% !TeX program = pdflatex
% papers/cristal_ciencias.tex — GERADO por tools/cristal_tex.py; NÃO editar à mão.
% Fonte: cristal/cristal.jsonl (broca-so, última versão por conceito).
% Cada secção carrega o registo original em comentário %CRISTAL — a volta é
% tests/cristal_volta.js (resíduo 0 byte a byte contra a fonte).
\documentclass[11pt,a4paper]{article}
\usepackage[utf8]{inputenc}\usepackage[T1]{fontenc}\usepackage[brazil]{babel}
\usepackage{amsmath,amssymb}\usepackage{textcomp}
\usepackage[margin=2.4cm]{geometry}\usepackage{xcolor}
\definecolor{cinza}{gray}{0.35}
\newcommand{\prov}[1]{\par\smallskip\noindent{\small\color{cinza}#1}\par}
\setcounter{secnumdepth}{0}
% ─────────────────────────────────────────────────────────────────────────────
%  papers/gkcapa.tex — a capa do Reino, mínima, para os papers em `article`.
%
%  O estilo.tex completo é do `report` (títulos de capítulo, skak, …). Aqui fica
%  só o que a capa precisa, para o paper continuar a compilar sozinho com
%  pdflatex. O tradutor próprio (tests/tex) carrega o estilo.tex da raiz / o
%  symlink papers/estilo.tex e avalia o mesmo `\gkcapa`.
% ─────────────────────────────────────────────────────────────────────────────
\definecolor{ouro}{HTML}{B8912F}
\definecolor{ouroclaro}{HTML}{F3E9CE}
\definecolor{tinta}{HTML}{1E1B16}
\definecolor{regua}{HTML}{6E6858}
\providecommand{\gknota}{\fontsize{7.62}{11.05}\selectfont}
\providecommand{\gktexto}{\fontsize{10.50}{15.22}\selectfont}
\providecommand{\gksub}{\fontsize{12.33}{17.87}\selectfont}
\providecommand{\gksec}{\fontsize{14.47}{20.98}\selectfont}
\providecommand{\gkcap}{\fontsize{16.99}{24.63}\selectfont}
\providecommand{\gktit}{\fontsize{23.42}{33.95}\selectfont}
\providecommand{\Prj}{\mathbb{P}}
\providecommand{\Trf}{\mathcal{T}}
\providecommand{\gkcapa}[3]{%
\title{\vspace{-0.6cm}
{\color{ouro}\rule{\textwidth}{1.4pt}}\\[6mm]
\textbf{\gktit\color{tinta} Reino Dourado}\\[3mm]{\gksec\itshape\color{regua} Golden Kingdom}\\[8mm]
{\gkcap\scshape\color{ouro} #1}\\[5mm]
{\gksub\color{regua} #2}\\[2mm]
{\gktexto\color{regua}\itshape #3}\\[7mm]
{\color{ouro}\rule{\textwidth}{0.6pt}}\\[9mm]{\gknota\color{regua}\begin{minipage}{\textwidth}\setlength{\parindent}{0pt}%
\textbf{Aviso de Isenção Algorítmica:} O universo, as leis e as entidades contidas nestas páginas ---
incluindo felinos matriciais, aves de rapina algébricas e amuletos de controle de fluxo --- são
construções de pura ficção formal. Esta obra não descreve a realidade física, não serve como
engenharia reversa do cosmos e não constitui doutrina religiosa. Qualquer semelhança com bugs reais do seu
sistema operacional é mera coincidência (ou falta de reza). Prossiga por sua própria conta e
risco.\end{minipage}}}
\author{}\date{}}

\begin{document}
\gkcapa{O Cristal --- as ciências de fora}
       {biologia, filosofia, o sagrado, medicina}
       {corpus recuperado do broca-so; 899 conceitos; projeção verificável da fonte}
\maketitle

%CRISTAL {"arestas":[{"alvo":"teoria_de_galois","confianca":1.0,"epistemico":"fato","origem":"wiki","rel":"usa"}],"confianca":1.0,"contraexemplos":[],"descricao":"Não há fórmula por radicais para a quíntica geral — porque seu grupo de Galois não é solúvel. O limite da resolução clássica.","epistemico":"fato","exemplos":[],"id":"abel_ruffini","memoria":"semantica","meta":{"autor":"Abel/Galois","dominio":"ciencias_de_fora","fonte":""},"origem":"wiki","palavras_chave":[],"sinonimos":[],"tipo":"conceito","titulo":"Teorema de Abel-Ruffini"}
\section{Teorema de Abel-Ruffini}
Não há fórmula por radicais para a quíntica geral — porque seu grupo de Galois não é solúvel. O limite da resolução clássica.
\prov{relações: usa teoria\_de\_galois}
\prov{proveniência: wiki \textperiodcentered{} fato \textperiodcentered{} confiança 1.0 \textperiodcentered{} domínio ciencias\_de\_fora \textperiodcentered{} história 4 versões no jornal}

%CRISTAL {"arestas":[{"alvo":"vida","confianca":1.0,"epistemico":"fato","origem":"wiki","rel":"explica"},{"alvo":"word","confianca":0.5,"epistemico":"fato","origem":"wiki","rel":"relaciona"},{"alvo":"virtualizacao","confianca":0.5,"epistemico":"fato","origem":"wiki","rel":"relaciona"},{"alvo":"goldchain_mercado_hub_fontes","confianca":0.5,"epistemico":"fato","origem":"wiki","rel":"relaciona"},{"alvo":"escassez","confianca":0.5,"epistemico":"fato","origem":"wiki","rel":"relaciona"},{"alvo":"flags","confianca":0.5,"epistemico":"fato","origem":"wiki","rel":"relaciona"},{"alvo":"proposition_gerador_de_escala_ipropderiv_com_tddt","confianca":0.5,"epistemico":"fato","origem":"wiki","rel":"relaciona"},{"alvo":"vida-morte","confianca":0.5,"epistemico":"fato","origem":"wiki","rel":"relaciona"},{"alvo":"gramatica","confianca":0.5,"epistemico":"fato","origem":"wiki","rel":"relaciona"},{"alvo":"universidade_curriculo_em_camadas","confianca":0.5,"epistemico":"fato","origem":"wiki","rel":"relaciona"},{"alvo":"semiotica_peirce_triade","confianca":0.5,"epistemico":"fato","origem":"wiki","rel":"relaciona"}],"confianca":1.0,"contraexemplos":[],"descricao":"A vida surge da química inorgânica (Oparin-Haldane, Miller-Urey); o Mundo RNA precede o DNA-proteínas — problema em aberto.","epistemico":"fato","exemplos":[],"id":"abiogenese","memoria":"semantica","meta":{"dominio":"ciencias_de_fora","fonte":"web"},"origem":"wiki","palavras_chave":[],"sinonimos":[],"tipo":"conceito","titulo":"Abiogênese"}
\section{Abiogênese}
A vida surge da química inorgânica (Oparin-Haldane, Miller-Urey); o Mundo RNA precede o DNA-proteínas — problema em aberto.
\prov{relações: explica vida; relaciona word; relaciona virtualizacao; relaciona goldchain\_mercado\_hub\_fontes; relaciona escassez; relaciona flags; relaciona proposition\_gerador\_de\_escala\_ipropderiv\_com\_tddt; relaciona vida-morte; relaciona gramatica; relaciona universidade\_curriculo\_em\_camadas; relaciona semiotica\_peirce\_triade}
\prov{proveniência: wiki \textperiodcentered{} web \textperiodcentered{} fato \textperiodcentered{} confiança 1.0 \textperiodcentered{} domínio ciencias\_de\_fora \textperiodcentered{} história 13 versões no jornal}

%CRISTAL {"arestas":[{"alvo":"artes_visuais","confianca":1.0,"epistemico":"fato","origem":"wiki","rel":"parte_de"},{"alvo":"perspectiva_linear","confianca":1.0,"epistemico":"fato","origem":"wiki","rel":"contradiz"},{"alvo":"mimese_expressao","confianca":1.0,"epistemico":"fato","origem":"wiki","rel":"contradiz"}],"confianca":1.0,"contraexemplos":[],"descricao":"Kandinsky (Do Espiritual na Arte): o abandono da figuração em favor de cor, linha e forma puras como linguagem autônoma; a obra não copia o mundo, mas atua diretamente sobre a alma. A ruptura com a representação.","epistemico":"inferencia","exemplos":[],"id":"abstracao_kandinsky","memoria":"semantica","meta":{"autor":"Kandinsky","dominio":"ciencias_de_fora","fonte":""},"origem":"wiki","palavras_chave":[],"sinonimos":[],"tipo":"conceito","titulo":"Abstração (Kandinsky)"}
\section{Abstração (Kandinsky)}
Kandinsky (Do Espiritual na Arte): o abandono da figuração em favor de cor, linha e forma puras como linguagem autônoma; a obra não copia o mundo, mas atua diretamente sobre a alma. A ruptura com a representação.
\prov{relações: parte\_de artes\_visuais; contradiz perspectiva\_linear; contradiz mimese\_expressao}
\prov{proveniência: wiki \textperiodcentered{} inferencia \textperiodcentered{} confiança 1.0 \textperiodcentered{} domínio ciencias\_de\_fora \textperiodcentered{} história 4 versões no jornal}

%CRISTAL {"arestas":[{"alvo":"sociologia","confianca":1.0,"epistemico":"fato","origem":"wiki","rel":"parte_de"},{"alvo":"word","confianca":0.5,"epistemico":"fato","origem":"wiki","rel":"relaciona"},{"alvo":"syscall","confianca":0.5,"epistemico":"fato","origem":"wiki","rel":"relaciona"},{"alvo":"virtualizacao","confianca":0.5,"epistemico":"fato","origem":"wiki","rel":"relaciona"}],"confianca":1.0,"contraexemplos":[],"descricao":"Weber: a ação orientada por sentido; a racionalização progressiva 'desencanta' o mundo, expulsando a magia.","epistemico":"inferencia","exemplos":[],"id":"acao_social_weber","memoria":"semantica","meta":{"autor":"Weber","dominio":"ciencias_de_fora","fonte":""},"origem":"wiki","palavras_chave":[],"sinonimos":[],"tipo":"conceito","titulo":"Ação Social e Desencantamento"}
\section{Ação Social e Desencantamento}
Weber: a ação orientada por sentido; a racionalização progressiva 'desencanta' o mundo, expulsando a magia.
\prov{relações: parte\_de sociologia; relaciona word; relaciona syscall; relaciona virtualizacao}
\prov{proveniência: wiki \textperiodcentered{} inferencia \textperiodcentered{} confiança 1.0 \textperiodcentered{} domínio ciencias\_de\_fora \textperiodcentered{} história 5 versões no jornal}

%CRISTAL {"arestas":[{"alvo":"musica","confianca":1.0,"epistemico":"fato","origem":"wiki","rel":"parte_de"},{"alvo":"industria_cultural","confianca":1.0,"epistemico":"fato","origem":"wiki","rel":"usa"},{"alvo":"atonalidade_dodecafonismo","confianca":1.0,"epistemico":"fato","origem":"wiki","rel":"relaciona"}],"confianca":1.0,"contraexemplos":[],"descricao":"Adorno: sob o capitalismo a música vira mercadoria fetichizada e produz uma 'regressão da escuta' — ouvir passivo e fragmentado; a vanguarda (Schoenberg) resistiria à comodificação, o pop a reforçaria.","epistemico":"inferencia","exemplos":[],"id":"adorno_musica","memoria":"semantica","meta":{"autor":"Adorno","dominio":"ciencias_de_fora","fonte":""},"origem":"wiki","palavras_chave":[],"sinonimos":[],"tipo":"conceito","titulo":"Crítica da Indústria Musical (Adorno)"}
\section{Crítica da Indústria Musical (Adorno)}
Adorno: sob o capitalismo a música vira mercadoria fetichizada e produz uma 'regressão da escuta' — ouvir passivo e fragmentado; a vanguarda (Schoenberg) resistiria à comodificação, o pop a reforçaria.
\prov{relações: parte\_de musica; usa industria\_cultural; relaciona atonalidade\_dodecafonismo}
\prov{proveniência: wiki \textperiodcentered{} inferencia \textperiodcentered{} confiança 1.0 \textperiodcentered{} domínio ciencias\_de\_fora \textperiodcentered{} história 4 versões no jornal}

%CRISTAL {"arestas":[{"alvo":"nao_dualidade","confianca":1.0,"epistemico":"fato","origem":"wiki","rel":"especializa"},{"alvo":"word","confianca":0.5,"epistemico":"fato","origem":"wiki","rel":"relaciona"},{"alvo":"while","confianca":0.5,"epistemico":"fato","origem":"wiki","rel":"relaciona"}],"confianca":1.0,"contraexemplos":[],"descricao":"Não-dualidade: Brahman é a única realidade; os fenômenos (Maya) são aparência; a libertação é reconhecer Atman=Brahman (Shankara).","epistemico":"hipotese","exemplos":[],"id":"advaita_vedanta","memoria":"semantica","meta":{"autor":"Shankara","dominio":"filosofia","fonte":""},"origem":"wiki","palavras_chave":[],"sinonimos":[],"tipo":"conceito","titulo":"Advaita Vedanta"}
\section{Advaita Vedanta}
Não-dualidade: Brahman é a única realidade; os fenômenos (Maya) são aparência; a libertação é reconhecer Atman=Brahman (Shankara).
\prov{relações: especializa nao\_dualidade; relaciona word; relaciona while}
\prov{proveniência: wiki \textperiodcentered{} hipotese \textperiodcentered{} confiança 1.0 \textperiodcentered{} domínio filosofia \textperiodcentered{} história 4 versões no jornal}

%CRISTAL {"arestas":[{"alvo":"cristianismo","confianca":1.0,"epistemico":"fato","origem":"wiki","rel":"parte_de"},{"alvo":"ethics_of_care","confianca":1.0,"epistemico":"fato","origem":"wiki","rel":"explica"}],"confianca":1.0,"contraexemplos":[],"descricao":"Bíblia: o amor incondicional, gratuito e doador cuja fonte é Deus e que se torna mandamento entre os humanos — 'amai-vos uns aos outros como eu vos amei'; 'Deus é amor'.","epistemico":"inferencia","exemplos":[],"id":"agape","memoria":"semantica","meta":{"autor":"Bíblia","dominio":"sagrado","fonte":""},"origem":"wiki","palavras_chave":[],"sinonimos":[],"tipo":"conceito","titulo":"Ágape — o Amor Oblativo"}
\section{Ágape — o Amor Oblativo}
Bíblia: o amor incondicional, gratuito e doador cuja fonte é Deus e que se torna mandamento entre os humanos — 'amai-vos uns aos outros como eu vos amei'; 'Deus é amor'.
\prov{relações: parte\_de cristianismo; explica ethics\_of\_care}
\prov{proveniência: wiki \textperiodcentered{} inferencia \textperiodcentered{} confiança 1.0 \textperiodcentered{} domínio sagrado \textperiodcentered{} história 3 versões no jornal}

%CRISTAL {"arestas":[{"alvo":"teoria_da_comunicacao","confianca":1.0,"epistemico":"fato","origem":"wiki","rel":"parte_de"},{"alvo":"gatekeeping","confianca":1.0,"epistemico":"fato","origem":"wiki","rel":"usa"},{"alvo":"teoria_hipodermica","confianca":1.0,"epistemico":"fato","origem":"wiki","rel":"contradiz"}],"confianca":1.0,"contraexemplos":[],"descricao":"McCombs & Shaw: a mídia não determina o que as pessoas pensam, mas define sobre o que elas pensam — a saliência dos temas na cobertura transfere-se para a hierarquia de importância no público.","epistemico":"fato","exemplos":[],"id":"agenda_setting","memoria":"semantica","meta":{"autor":"McCombs & Shaw","dominio":"ciencias_de_fora","fonte":""},"origem":"wiki","palavras_chave":[],"sinonimos":[],"tipo":"conceito","titulo":"Agenda-Setting"}
\section{Agenda-Setting}
McCombs \& Shaw: a mídia não determina o que as pessoas pensam, mas define sobre o que elas pensam — a saliência dos temas na cobertura transfere-se para a hierarquia de importância no público.
\prov{relações: parte\_de teoria\_da\_comunicacao; usa gatekeeping; contradiz teoria\_hipodermica}
\prov{proveniência: wiki \textperiodcentered{} fato \textperiodcentered{} confiança 1.0 \textperiodcentered{} domínio ciencias\_de\_fora \textperiodcentered{} história 4 versões no jornal}

%CRISTAL {"arestas":[{"alvo":"von_thunen","confianca":1.0,"epistemico":"fato","origem":"wiki","rel":"usa"},{"alvo":"mao_invisivel","confianca":1.0,"epistemico":"fato","origem":"wiki","rel":"explica"},{"alvo":"metropole_megacidade","confianca":1.0,"epistemico":"fato","origem":"wiki","rel":"explica"}],"confianca":1.0,"contraexemplos":[],"descricao":"Krugman: a produção tende a concentrar-se (aglomerar) por rendimentos crescentes, custos de transporte e efeitos de mercado, gerando estruturas centro-periferia — a auto-organização espacial dos mercados.","epistemico":"inferencia","exemplos":[],"id":"aglomeracao_krugman","memoria":"semantica","meta":{"autor":"Krugman","dominio":"ciencias_de_fora","fonte":""},"origem":"wiki","palavras_chave":[],"sinonimos":[],"tipo":"conceito","titulo":"Aglomeração / Nova Geografia Econômica (Krugman)"}
\section{Aglomeração / Nova Geografia Econômica (Krugman)}
Krugman: a produção tende a concentrar-se (aglomerar) por rendimentos crescentes, custos de transporte e efeitos de mercado, gerando estruturas centro-periferia — a auto-organização espacial dos mercados.
\prov{relações: usa von\_thunen; explica mao\_invisivel; explica metropole\_megacidade}
\prov{proveniência: wiki \textperiodcentered{} inferencia \textperiodcentered{} confiança 1.0 \textperiodcentered{} domínio ciencias\_de\_fora \textperiodcentered{} história 4 versões no jornal}

%CRISTAL {"arestas":[{"alvo":"isla","confianca":1.0,"epistemico":"fato","origem":"wiki","rel":"parte_de"},{"alvo":"simbolo_vs_significado","confianca":1.0,"epistemico":"fato","origem":"wiki","rel":"relaciona"},{"alvo":"logos_verbo","confianca":1.0,"epistemico":"fato","origem":"wiki","rel":"relaciona"}],"confianca":1.0,"contraexemplos":[],"descricao":"Islã: 'al-qur'an' significa 'a recitação' — a palavra literal de Deus revelada a Muhammad, tida (na tradição sunita majoritária) como incriada e coeterna a Deus. A palavra como interface do transcendente.","epistemico":"inferencia","exemplos":[],"id":"alcorao_palavra_revelada","memoria":"semantica","meta":{"autor":"Alcorão","dominio":"sagrado","fonte":""},"origem":"wiki","palavras_chave":[],"sinonimos":[],"tipo":"conceito","titulo":"O Alcorão — a Palavra Revelada"}
\section{O Alcorão — a Palavra Revelada}
Islã: 'al-qur'an' significa 'a recitação' — a palavra literal de Deus revelada a Muhammad, tida (na tradição sunita majoritária) como incriada e coeterna a Deus. A palavra como interface do transcendente.
\prov{relações: parte\_de isla; relaciona simbolo\_vs\_significado; relaciona logos\_verbo}
\prov{proveniência: wiki \textperiodcentered{} inferencia \textperiodcentered{} confiança 1.0 \textperiodcentered{} domínio sagrado \textperiodcentered{} história 4 versões no jornal}

%CRISTAL {"arestas":[{"alvo":"meio_e_mensagem","confianca":1.0,"epistemico":"fato","origem":"wiki","rel":"usa"},{"alvo":"contrato_social","confianca":1.0,"epistemico":"fato","origem":"wiki","rel":"relaciona"},{"alvo":"word","confianca":0.5,"epistemico":"fato","origem":"wiki","rel":"relaciona"}],"confianca":1.0,"contraexemplos":[],"descricao":"McLuhan: os meios eletrônicos colapsam tempo e distância, retribalizando a humanidade numa comunidade interconectada de escala planetária.","epistemico":"inferencia","exemplos":[],"id":"aldeia_global","memoria":"semantica","meta":{"autor":"McLuhan","dominio":"ciencias_de_fora","fonte":""},"origem":"wiki","palavras_chave":[],"sinonimos":[],"tipo":"conceito","titulo":"Aldeia Global (McLuhan)"}
\section{Aldeia Global (McLuhan)}
McLuhan: os meios eletrônicos colapsam tempo e distância, retribalizando a humanidade numa comunidade interconectada de escala planetária.
\prov{relações: usa meio\_e\_mensagem; relaciona contrato\_social; relaciona word}
\prov{proveniência: wiki \textperiodcentered{} inferencia \textperiodcentered{} confiança 1.0 \textperiodcentered{} domínio ciencias\_de\_fora \textperiodcentered{} história 4 versões no jornal}

%CRISTAL {"arestas":[{"alvo":"metafisica","confianca":1.0,"epistemico":"fato","origem":"wiki","rel":"parte_de"},{"alvo":"contemplacao","confianca":1.0,"epistemico":"fato","origem":"wiki","rel":"explica"},{"alvo":"word","confianca":0.5,"epistemico":"fato","origem":"wiki","rel":"relaciona"}],"confianca":1.0,"contraexemplos":[],"descricao":"Platão: prisioneiros veem sombras (ilusão sensível) até ascenderem à luz das Formas (realidade inteligível). O Bem como fonte.","epistemico":"inferencia","exemplos":[],"id":"alegoria_caverna","memoria":"semantica","meta":{"autor":"Platão","dominio":"filosofia","fonte":""},"origem":"wiki","palavras_chave":[],"sinonimos":[],"tipo":"conceito","titulo":"Alegoria da Caverna"}
\section{Alegoria da Caverna}
Platão: prisioneiros veem sombras (ilusão sensível) até ascenderem à luz das Formas (realidade inteligível). O Bem como fonte.
\prov{relações: parte\_de metafisica; explica contemplacao; relaciona word}
\prov{proveniência: wiki \textperiodcentered{} inferencia \textperiodcentered{} confiança 1.0 \textperiodcentered{} domínio filosofia \textperiodcentered{} história 4 versões no jornal}

%CRISTAL {"arestas":[{"alvo":"indo_europeu","confianca":1.0,"epistemico":"fato","origem":"wiki","rel":"parte_de"}],"confianca":1.0,"contraexemplos":[],"descricao":"Língua germânica (~95 milhões), com 4 casos, 3 gêneros e substantivos compostos — a língua da filosofia e da ciência moderna.","epistemico":"fato","exemplos":[],"id":"alemao","memoria":"semantica","meta":{"dominio":"ciencias_de_fora","fonte":""},"origem":"wiki","palavras_chave":[],"sinonimos":[],"tipo":"conceito","titulo":"Alemão"}
\section{Alemão}
Língua germânica (\~{}95 milhões), com 4 casos, 3 gêneros e substantivos compostos — a língua da filosofia e da ciência moderna.
\prov{relações: parte\_de indo\_europeu}
\prov{proveniência: wiki \textperiodcentered{} fato \textperiodcentered{} confiança 1.0 \textperiodcentered{} domínio ciencias\_de\_fora \textperiodcentered{} história 2 versões no jornal}

%CRISTAL {"arestas":[{"alvo":"sistema_de_escrita","confianca":1.0,"epistemico":"fato","origem":"wiki","rel":"especializa"},{"alvo":"fonologia","confianca":1.0,"epistemico":"fato","origem":"wiki","rel":"relaciona"}],"confianca":1.0,"contraexemplos":[],"descricao":"Um punhado de letras para os sons de uma língua — a economia que tornou a escrita acessível; do fenício ao latino ao cirílico.","epistemico":"fato","exemplos":[],"id":"alfabeto","memoria":"semantica","meta":{"dominio":"ciencias_de_fora","fonte":""},"origem":"wiki","palavras_chave":[],"sinonimos":[],"tipo":"conceito","titulo":"Alfabeto"}
\section{Alfabeto}
Um punhado de letras para os sons de uma língua — a economia que tornou a escrita acessível; do fenício ao latino ao cirílico.
\prov{relações: especializa sistema\_de\_escrita; relaciona fonologia}
\prov{proveniência: wiki \textperiodcentered{} fato \textperiodcentered{} confiança 1.0 \textperiodcentered{} domínio ciencias\_de\_fora \textperiodcentered{} história 3 versões no jornal}

%CRISTAL {"arestas":[{"alvo":"matematica","confianca":1.0,"epistemico":"fato","origem":"wiki","rel":"parte_de"},{"alvo":"algebra","confianca":1.0,"epistemico":"fato","origem":"wiki","rel":"relaciona"}],"confianca":1.0,"contraexemplos":[],"descricao":"O estudo das estruturas — grupos, anéis, corpos — pela operação e não pelo objeto; a forma pura da simetria.","epistemico":"fato","exemplos":[],"id":"algebra_abstrata","memoria":"semantica","meta":{"dominio":"ciencias_de_fora","fonte":""},"origem":"wiki","palavras_chave":[],"sinonimos":[],"tipo":"conceito","titulo":"Álgebra Abstrata"}
\section{Álgebra Abstrata}
O estudo das estruturas — grupos, anéis, corpos — pela operação e não pelo objeto; a forma pura da simetria.
\prov{relações: parte\_de matematica; relaciona algebra}
\prov{proveniência: wiki \textperiodcentered{} fato \textperiodcentered{} confiança 1.0 \textperiodcentered{} domínio ciencias\_de\_fora \textperiodcentered{} história 5 versões no jornal}

%CRISTAL {"arestas":[{"alvo":"algebra_abstrata","confianca":1.0,"epistemico":"fato","origem":"wiki","rel":"parte_de"},{"alvo":"gato_generaliza_matriz_companion","confianca":1.0,"epistemico":"fato","origem":"wiki","rel":"explica"},{"alvo":"gato","confianca":1.0,"epistemico":"fato","origem":"wiki","rel":"explica"}],"confianca":1.0,"contraexemplos":[],"descricao":"O cálculo de vetores, matrizes e transformações lineares — a linguagem em que a matriz de gatos e a potência posicional Aₙk vivem.","epistemico":"fato","exemplos":[],"id":"algebra_linear","memoria":"semantica","meta":{"dominio":"ciencias_de_fora","fonte":""},"origem":"wiki","palavras_chave":[],"sinonimos":[],"tipo":"conceito","titulo":"Álgebra Linear"}
\section{Álgebra Linear}
O cálculo de vetores, matrizes e transformações lineares — a linguagem em que a matriz de gatos e a potência posicional A\ensuremath{_{n}}k vivem.
\prov{relações: parte\_de algebra\_abstrata; explica gato\_generaliza\_matriz\_companion; explica gato}
\prov{proveniência: wiki \textperiodcentered{} fato \textperiodcentered{} confiança 1.0 \textperiodcentered{} domínio ciencias\_de\_fora \textperiodcentered{} história 8 versões no jornal}

%CRISTAL {"arestas":[{"alvo":"cristianismo","confianca":1.0,"epistemico":"fato","origem":"wiki","rel":"parte_de"},{"alvo":"contrato_social","confianca":1.0,"epistemico":"fato","origem":"wiki","rel":"relaciona"}],"confianca":1.0,"contraexemplos":[],"descricao":"Bíblia: o pacto vinculante entre Deus e um povo, iniciado unilateralmente por Deus (Noé, Abraão, Sinai), que funda pertença e obrigação mútua; traduz-se em conduta na Lei (Torá) e nos Dez Mandamentos. Contraste: a origem é transcendente e assimétrica, não um pacto entre iguais.","epistemico":"inferencia","exemplos":[],"id":"alianca_berith","memoria":"semantica","meta":{"autor":"Bíblia","dominio":"sagrado","fonte":""},"origem":"wiki","palavras_chave":[],"sinonimos":[],"tipo":"conceito","titulo":"A Aliança (berith) e a Lei"}
\section{A Aliança (berith) e a Lei}
Bíblia: o pacto vinculante entre Deus e um povo, iniciado unilateralmente por Deus (Noé, Abraão, Sinai), que funda pertença e obrigação mútua; traduz-se em conduta na Lei (Torá) e nos Dez Mandamentos. Contraste: a origem é transcendente e assimétrica, não um pacto entre iguais.
\prov{relações: parte\_de cristianismo; relaciona contrato\_social}
\prov{proveniência: wiki \textperiodcentered{} inferencia \textperiodcentered{} confiança 1.0 \textperiodcentered{} domínio sagrado \textperiodcentered{} história 3 versões no jornal}

%CRISTAL {"arestas":[{"alvo":"traducao_automatica","confianca":1.0,"epistemico":"fato","origem":"wiki","rel":"usa"}],"confianca":1.0,"contraexemplos":[],"descricao":"Descobrir qual pedaço da fonte corresponde a qual do alvo, em pares de frases traduzidas — a matéria-prima da tradução estatística e neural (o corpus Tatoeba do broca-so é isto).","epistemico":"fato","exemplos":[],"id":"alinhamento_de_traducao","memoria":"semantica","meta":{"dominio":"ciencias_de_fora","fonte":""},"origem":"wiki","palavras_chave":[],"sinonimos":[],"tipo":"conceito","titulo":"Alinhamento"}
\section{Alinhamento}
Descobrir qual pedaço da fonte corresponde a qual do alvo, em pares de frases traduzidas — a matéria-prima da tradução estatística e neural (o corpus Tatoeba do broca-so é isto).
\prov{relações: usa traducao\_automatica}
\prov{proveniência: wiki \textperiodcentered{} fato \textperiodcentered{} confiança 1.0 \textperiodcentered{} domínio ciencias\_de\_fora \textperiodcentered{} história 2 versões no jornal}

%CRISTAL {"arestas":[{"alvo":"tawhid","confianca":1.0,"epistemico":"fato","origem":"wiki","rel":"usa"},{"alvo":"simbolo_vs_significado","confianca":1.0,"epistemico":"fato","origem":"wiki","rel":"relaciona"}],"confianca":1.0,"contraexemplos":[],"descricao":"Islã: Allah é o nome do Deus único; a tradição enumera 99 'belos nomes'/atributos revelados (o Misericordioso, o Onisciente, o Soberano) — múltiplos nomes de uma só essência. Ar-Rahman (o Clemente) abre cada sura: a misericórdia como atributo primeiro.","epistemico":"inferencia","exemplos":[],"id":"allah_99_nomes","memoria":"semantica","meta":{"autor":"Alcorão","dominio":"sagrado","fonte":""},"origem":"wiki","palavras_chave":[],"sinonimos":[],"tipo":"conceito","titulo":"Allah e os 99 Belos Nomes"}
\section{Allah e os 99 Belos Nomes}
Islã: Allah é o nome do Deus único; a tradição enumera 99 'belos nomes'/atributos revelados (o Misericordioso, o Onisciente, o Soberano) — múltiplos nomes de uma só essência. Ar-Rahman (o Clemente) abre cada sura: a misericórdia como atributo primeiro.
\prov{relações: usa tawhid; relaciona simbolo\_vs\_significado}
\prov{proveniência: wiki \textperiodcentered{} inferencia \textperiodcentered{} confiança 1.0 \textperiodcentered{} domínio sagrado \textperiodcentered{} história 3 versões no jornal}

%CRISTAL {"arestas":[{"alvo":"analise_matematica","confianca":1.0,"epistemico":"fato","origem":"wiki","rel":"parte_de"},{"alvo":"calculo_infinitesimal","confianca":1.0,"epistemico":"fato","origem":"wiki","rel":"usa"}],"confianca":1.0,"contraexemplos":[],"descricao":"Todo sinal periódico é soma de senos — decompor a função em frequências. A dualidade tempo↔espectro.","epistemico":"fato","exemplos":[],"id":"analise_de_fourier","memoria":"semantica","meta":{"autor":"Fourier","dominio":"ciencias_de_fora","fonte":""},"origem":"wiki","palavras_chave":[],"sinonimos":[],"tipo":"conceito","titulo":"Análise de Fourier"}
\section{Análise de Fourier}
Todo sinal periódico é soma de senos — decompor a função em frequências. A dualidade tempo\ensuremath{\leftrightarrow}espectro.
\prov{relações: parte\_de analise\_matematica; usa calculo\_infinitesimal}
\prov{proveniência: wiki \textperiodcentered{} fato \textperiodcentered{} confiança 1.0 \textperiodcentered{} domínio ciencias\_de\_fora \textperiodcentered{} história 6 versões no jornal}

%CRISTAL {"arestas":[{"alvo":"direito","confianca":1.0,"epistemico":"fato","origem":"wiki","rel":"parte_de"},{"alvo":"algebra_dos_contratos","confianca":1.0,"epistemico":"fato","origem":"wiki","rel":"explica"},{"alvo":"smart_contract","confianca":1.0,"epistemico":"fato","origem":"wiki","rel":"relaciona"}],"confianca":1.0,"contraexemplos":[],"descricao":"Hohfeld: decompõe o vago 'direito' em oito posições correlatas e opostas — direito/dever, privilégio/não-direito, poder/sujeição, imunidade/incapacidade —, revelando a estrutura relacional precisa de toda relação jurídica. Uma álgebra de direitos.","epistemico":"inferencia","exemplos":[],"id":"analise_de_hohfeld","memoria":"semantica","meta":{"autor":"Hohfeld","dominio":"ciencias_de_fora","fonte":""},"origem":"wiki","palavras_chave":[],"sinonimos":[],"tipo":"conceito","titulo":"Posições Jurídicas Fundamentais (Hohfeld)"}
\section{Posições Jurídicas Fundamentais (Hohfeld)}
Hohfeld: decompõe o vago 'direito' em oito posições correlatas e opostas — direito/dever, privilégio/não-direito, poder/sujeição, imunidade/incapacidade —, revelando a estrutura relacional precisa de toda relação jurídica. Uma álgebra de direitos.
\prov{relações: parte\_de direito; explica algebra\_dos\_contratos; relaciona smart\_contract}
\prov{proveniência: wiki \textperiodcentered{} inferencia \textperiodcentered{} confiança 1.0 \textperiodcentered{} domínio ciencias\_de\_fora \textperiodcentered{} história 4 versões no jornal}

%CRISTAL {"arestas":[{"alvo":"direito_e_tecnologia","confianca":1.0,"epistemico":"fato","origem":"wiki","rel":"parte_de"},{"alvo":"mao_invisivel","confianca":1.0,"epistemico":"fato","origem":"wiki","rel":"usa"},{"alvo":"teorema_de_coase","confianca":1.0,"epistemico":"fato","origem":"wiki","rel":"usa"}],"confianca":1.0,"contraexemplos":[],"descricao":"Posner: avaliar normas e decisões por seus efeitos de eficiência, tratando o direito como sistema de incentivos; o common law tenderia a operar 'como se' os juízes maximizassem a riqueza social.","epistemico":"inferencia","exemplos":[],"id":"analise_economica_direito","memoria":"semantica","meta":{"autor":"Posner","dominio":"ciencias_de_fora","fonte":""},"origem":"wiki","palavras_chave":[],"sinonimos":[],"tipo":"conceito","titulo":"Análise Econômica do Direito"}
\section{Análise Econômica do Direito}
Posner: avaliar normas e decisões por seus efeitos de eficiência, tratando o direito como sistema de incentivos; o common law tenderia a operar 'como se' os juízes maximizassem a riqueza social.
\prov{relações: parte\_de direito\_e\_tecnologia; usa mao\_invisivel; usa teorema\_de\_coase}
\prov{proveniência: wiki \textperiodcentered{} inferencia \textperiodcentered{} confiança 1.0 \textperiodcentered{} domínio ciencias\_de\_fora \textperiodcentered{} história 4 versões no jornal}

%CRISTAL {"arestas":[{"alvo":"matematica","confianca":1.0,"epistemico":"fato","origem":"wiki","rel":"parte_de"}],"confianca":1.0,"contraexemplos":[],"descricao":"O rigor do contínuo e do limite — cálculo, séries, convergência; a domesticação do infinitesimal.","epistemico":"fato","exemplos":[],"id":"analise_matematica","memoria":"semantica","meta":{"dominio":"ciencias_de_fora","fonte":""},"origem":"wiki","palavras_chave":[],"sinonimos":[],"tipo":"conceito","titulo":"Análise Matemática"}
\section{Análise Matemática}
O rigor do contínuo e do limite — cálculo, séries, convergência; a domesticação do infinitesimal.
\prov{relações: parte\_de matematica}
\prov{proveniência: wiki \textperiodcentered{} fato \textperiodcentered{} confiança 1.0 \textperiodcentered{} domínio ciencias\_de\_fora \textperiodcentered{} história 4 versões no jornal}

%CRISTAL {"arestas":[{"alvo":"zona_desenvolvimento_proximal","confianca":1.0,"epistemico":"fato","origem":"wiki","rel":"usa"}],"confianca":1.0,"contraexemplos":[],"descricao":"Bruner (a partir de Vygotsky): apoio temporário e ajustado do mais capaz, retirado gradualmente conforme cresce a competência, permitindo autonomia progressiva.","epistemico":"inferencia","exemplos":[],"id":"andaimagem_scaffolding","memoria":"semantica","meta":{"autor":"Bruner/Wood","dominio":"ciencias_de_fora","fonte":""},"origem":"wiki","palavras_chave":[],"sinonimos":[],"tipo":"conceito","titulo":"Andaimagem (Scaffolding)"}
\section{Andaimagem (Scaffolding)}
Bruner (a partir de Vygotsky): apoio temporário e ajustado do mais capaz, retirado gradualmente conforme cresce a competência, permitindo autonomia progressiva.
\prov{relações: usa zona\_desenvolvimento\_proximal}
\prov{proveniência: wiki \textperiodcentered{} inferencia \textperiodcentered{} confiança 1.0 \textperiodcentered{} domínio ciencias\_de\_fora \textperiodcentered{} história 2 versões no jornal}

%CRISTAL {"arestas":[{"alvo":"materialismo_historico","confianca":1.0,"epistemico":"fato","origem":"wiki","rel":"usa"},{"alvo":"word","confianca":0.5,"epistemico":"fato","origem":"wiki","rel":"relaciona"},{"alvo":"vontade_de_poder","confianca":0.5,"epistemico":"fato","origem":"wiki","rel":"relaciona"},{"alvo":"syscall","confianca":0.5,"epistemico":"fato","origem":"wiki","rel":"relaciona"},{"alvo":"virtualizacao","confianca":0.5,"epistemico":"fato","origem":"wiki","rel":"relaciona"},{"alvo":"politica","confianca":0.5,"epistemico":"fato","origem":"wiki","rel":"relaciona"}],"confianca":1.0,"contraexemplos":[],"descricao":"Benjamin: a história oficial é a dos vencedores; a verdadeira é a contra-história dos oprimidos, resgatada no 'agora-tempo' (Jetztzeit).","epistemico":"inferencia","exemplos":[],"id":"anjo_da_historia","memoria":"semantica","meta":{"autor":"Benjamin","dominio":"ciencias_de_fora","fonte":""},"origem":"wiki","palavras_chave":[],"sinonimos":[],"tipo":"conceito","titulo":"O Anjo da História"}
\section{O Anjo da História}
Benjamin: a história oficial é a dos vencedores; a verdadeira é a contra-história dos oprimidos, resgatada no 'agora-tempo' (Jetztzeit).
\prov{relações: usa materialismo\_historico; relaciona word; relaciona vontade\_de\_poder; relaciona syscall; relaciona virtualizacao; relaciona politica}
\prov{proveniência: wiki \textperiodcentered{} inferencia \textperiodcentered{} confiança 1.0 \textperiodcentered{} domínio ciencias\_de\_fora \textperiodcentered{} história 7 versões no jornal}

%CRISTAL {"arestas":[{"alvo":"geografia_fisica","confianca":1.0,"epistemico":"fato","origem":"wiki","rel":"parte_de"},{"alvo":"ciclo_carbono","confianca":1.0,"epistemico":"fato","origem":"wiki","rel":"usa"},{"alvo":"big_history","confianca":1.0,"epistemico":"fato","origem":"wiki","rel":"relaciona"}],"confianca":1.0,"contraexemplos":[],"descricao":"Crutzen & Stoermer: a época geológica proposta em que a ação humana coletiva se tornou força geofísica dominante, alterando clima, ciclos biogeoquímicos, biodiversidade e estratos — deixando assinatura no registro geológico.","epistemico":"inferencia","exemplos":[],"id":"antropoceno","memoria":"semantica","meta":{"autor":"Crutzen & Stoermer","dominio":"ciencias_de_fora","fonte":""},"origem":"wiki","palavras_chave":[],"sinonimos":[],"tipo":"conceito","titulo":"Antropoceno"}
\section{Antropoceno}
Crutzen \& Stoermer: a época geológica proposta em que a ação humana coletiva se tornou força geofísica dominante, alterando clima, ciclos biogeoquímicos, biodiversidade e estratos — deixando assinatura no registro geológico.
\prov{relações: parte\_de geografia\_fisica; usa ciclo\_carbono; relaciona big\_history}
\prov{proveniência: wiki \textperiodcentered{} inferencia \textperiodcentered{} confiança 1.0 \textperiodcentered{} domínio ciencias\_de\_fora \textperiodcentered{} história 4 versões no jornal}

%CRISTAL {"arestas":[{"alvo":"sociologia","confianca":1.0,"epistemico":"fato","origem":"wiki","rel":"relaciona"},{"alvo":"historia","confianca":0.5,"epistemico":"fato","origem":"wiki","rel":"relaciona"},{"alvo":"utilitarismo","confianca":0.5,"epistemico":"fato","origem":"wiki","rel":"relaciona"},{"alvo":"fixamos_e_o_curinga","confianca":0.5,"epistemico":"fato","origem":"wiki","rel":"relaciona"},{"alvo":"melhor_candidato","confianca":0.5,"epistemico":"fato","origem":"wiki","rel":"relaciona"},{"alvo":"broca-so_cristalizado","confianca":0.5,"epistemico":"fato","origem":"wiki","rel":"relaciona"}],"confianca":1.0,"contraexemplos":[],"descricao":"O estudo do humano em suas culturas — parentesco, mito, troca, ritual.","epistemico":"fato","exemplos":[],"id":"antropologia","memoria":"semantica","meta":{"dominio":"ciencias_de_fora","fonte":""},"origem":"wiki","palavras_chave":[],"sinonimos":[],"tipo":"conceito","titulo":"Antropologia"}
\section{Antropologia}
O estudo do humano em suas culturas — parentesco, mito, troca, ritual.
\prov{relações: relaciona sociologia; relaciona historia; relaciona utilitarismo; relaciona fixamos\_e\_o\_curinga; relaciona melhor\_candidato; relaciona broca-so\_cristalizado}
\prov{proveniência: wiki \textperiodcentered{} fato \textperiodcentered{} confiança 1.0 \textperiodcentered{} domínio ciencias\_de\_fora \textperiodcentered{} história 7 versões no jornal}

%CRISTAL {"arestas":[{"alvo":"psicanalise","confianca":1.0,"epistemico":"fato","origem":"wiki","rel":"parte_de"},{"alvo":"inconsciente","confianca":1.0,"epistemico":"fato","origem":"wiki","rel":"usa"},{"alvo":"consciencia_pura","confianca":1.0,"epistemico":"fato","origem":"wiki","rel":"contradiz"},{"alvo":"hopfield","confianca":1.0,"epistemico":"fato","origem":"wiki","rel":"referencia"}],"confianca":1.0,"contraexemplos":[],"descricao":"Freud: a psique em três instâncias — id (impulso), ego (realidade), superego (moral); construtos, não estruturas neurais.","epistemico":"hipotese","exemplos":[],"id":"aparelho_psiquico","memoria":"semantica","meta":{"autor":"Freud","dominio":"ciencias_de_fora","fonte":""},"origem":"wiki","palavras_chave":[],"sinonimos":[],"tipo":"conceito","titulo":"O Aparelho Psíquico"}
\section{O Aparelho Psíquico}
Freud: a psique em três instâncias — id (impulso), ego (realidade), superego (moral); construtos, não estruturas neurais.
\prov{relações: parte\_de psicanalise; usa inconsciente; contradiz consciencia\_pura; referencia hopfield}
\prov{proveniência: wiki \textperiodcentered{} hipotese \textperiodcentered{} confiança 1.0 \textperiodcentered{} domínio ciencias\_de\_fora \textperiodcentered{} história 5 versões no jornal}

%CRISTAL {"arestas":[{"alvo":"psicologia","confianca":1.0,"epistemico":"fato","origem":"wiki","rel":"parte_de"},{"alvo":"word","confianca":0.5,"epistemico":"fato","origem":"wiki","rel":"relaciona"},{"alvo":"virtualizacao","confianca":0.5,"epistemico":"fato","origem":"wiki","rel":"relaciona"},{"alvo":"vida-morte","confianca":0.5,"epistemico":"fato","origem":"wiki","rel":"relaciona"},{"alvo":"estrutura_de_dados","confianca":0.5,"epistemico":"fato","origem":"wiki","rel":"relaciona"},{"alvo":"estetica","confianca":0.5,"epistemico":"fato","origem":"wiki","rel":"relaciona"},{"alvo":"misticismo_nao_dual","confianca":0.5,"epistemico":"fato","origem":"wiki","rel":"relaciona"}],"confianca":1.0,"contraexemplos":[],"descricao":"Bowlby: o vínculo criança-cuidador é biologicamente programado; sua falha marca as relações — o proto-ego nasce na relação.","epistemico":"inferencia","exemplos":[],"id":"apego","memoria":"semantica","meta":{"autor":"Bowlby","dominio":"ciencias_de_fora","fonte":""},"origem":"wiki","palavras_chave":[],"sinonimos":[],"tipo":"conceito","titulo":"Teoria do Apego"}
\section{Teoria do Apego}
Bowlby: o vínculo criança-cuidador é biologicamente programado; sua falha marca as relações — o proto-ego nasce na relação.
\prov{relações: parte\_de psicologia; relaciona word; relaciona virtualizacao; relaciona vida-morte; relaciona estrutura\_de\_dados; relaciona estetica; relaciona misticismo\_nao\_dual}
\prov{proveniência: wiki \textperiodcentered{} inferencia \textperiodcentered{} confiança 1.0 \textperiodcentered{} domínio ciencias\_de\_fora \textperiodcentered{} história 8 versões no jornal}

%CRISTAL {"arestas":[{"alvo":"estetica","confianca":1.0,"epistemico":"fato","origem":"wiki","rel":"parte_de"},{"alvo":"borda_critica_vazio","confianca":1.0,"epistemico":"fato","origem":"wiki","rel":"relaciona"},{"alvo":"virtualizacao","confianca":0.5,"epistemico":"fato","origem":"wiki","rel":"relaciona"},{"alvo":"word","confianca":0.5,"epistemico":"fato","origem":"wiki","rel":"relaciona"},{"alvo":"escassez","confianca":0.5,"epistemico":"fato","origem":"wiki","rel":"relaciona"},{"alvo":"universidade_curriculo_em_camadas","confianca":0.5,"epistemico":"fato","origem":"wiki","rel":"relaciona"}],"confianca":1.0,"contraexemplos":[],"descricao":"Nietzsche: a arte nasce da tensão entre a forma/medida (Apolo) e o êxtase/desmedida (Dioniso).","epistemico":"inferencia","exemplos":[],"id":"apolineo_dionisiaco","memoria":"semantica","meta":{"autor":"Nietzsche","dominio":"ciencias_de_fora","fonte":""},"origem":"wiki","palavras_chave":[],"sinonimos":[],"tipo":"conceito","titulo":"Apolíneo e Dionisíaco"}
\section{Apolíneo e Dionisíaco}
Nietzsche: a arte nasce da tensão entre a forma/medida (Apolo) e o êxtase/desmedida (Dioniso).
\prov{relações: parte\_de estetica; relaciona borda\_critica\_vazio; relaciona virtualizacao; relaciona word; relaciona escassez; relaciona universidade\_curriculo\_em\_camadas}
\prov{proveniência: wiki \textperiodcentered{} inferencia \textperiodcentered{} confiança 1.0 \textperiodcentered{} domínio ciencias\_de\_fora \textperiodcentered{} história 7 versões no jornal}

%CRISTAL {"arestas":[{"alvo":"teorias_aprendizagem","confianca":1.0,"epistemico":"fato","origem":"wiki","rel":"parte_de"}],"confianca":1.0,"contraexemplos":[],"descricao":"Dewey: o conhecimento genuíno emerge do enfrentamento reflexivo de problemas reais na experiência, não da memorização; educação é reconstrução contínua da experiência.","epistemico":"inferencia","exemplos":[],"id":"aprender_fazendo_dewey","memoria":"semantica","meta":{"autor":"Dewey","dominio":"ciencias_de_fora","fonte":""},"origem":"wiki","palavras_chave":[],"sinonimos":[],"tipo":"conceito","titulo":"Aprender Fazendo (Dewey)"}
\section{Aprender Fazendo (Dewey)}
Dewey: o conhecimento genuíno emerge do enfrentamento reflexivo de problemas reais na experiência, não da memorização; educação é reconstrução contínua da experiência.
\prov{relações: parte\_de teorias\_aprendizagem}
\prov{proveniência: wiki \textperiodcentered{} inferencia \textperiodcentered{} confiança 1.0 \textperiodcentered{} domínio ciencias\_de\_fora \textperiodcentered{} história 2 versões no jornal}

%CRISTAL {"arestas":[{"alvo":"teorias_aprendizagem","confianca":1.0,"epistemico":"fato","origem":"wiki","rel":"parte_de"},{"alvo":"andaimagem_scaffolding","confianca":1.0,"epistemico":"fato","origem":"wiki","rel":"usa"}],"confianca":1.0,"contraexemplos":[],"descricao":"Bruner: o aprendiz constrói compreensão descobrindo por si estruturas e relações, com orientação (descoberta guiada), em vez de recepção passiva de fatos.","epistemico":"inferencia","exemplos":[],"id":"aprendizagem_descoberta_bruner","memoria":"semantica","meta":{"autor":"Bruner","dominio":"ciencias_de_fora","fonte":""},"origem":"wiki","palavras_chave":[],"sinonimos":[],"tipo":"conceito","titulo":"Aprendizagem por Descoberta (Bruner)"}
\section{Aprendizagem por Descoberta (Bruner)}
Bruner: o aprendiz constrói compreensão descobrindo por si estruturas e relações, com orientação (descoberta guiada), em vez de recepção passiva de fatos.
\prov{relações: parte\_de teorias\_aprendizagem; usa andaimagem\_scaffolding}
\prov{proveniência: wiki \textperiodcentered{} inferencia \textperiodcentered{} confiança 1.0 \textperiodcentered{} domínio ciencias\_de\_fora \textperiodcentered{} história 3 versões no jornal}

%CRISTAL {"arestas":[{"alvo":"teorias_aprendizagem","confianca":1.0,"epistemico":"fato","origem":"wiki","rel":"parte_de"},{"alvo":"hopfield","confianca":1.0,"epistemico":"fato","origem":"wiki","rel":"referencia"},{"alvo":"assimilacao_acomodacao","confianca":1.0,"epistemico":"fato","origem":"wiki","rel":"relaciona"}],"confianca":1.0,"contraexemplos":[],"descricao":"Ausubel: a aprendizagem é significativa quando o novo se ancora de modo não-arbitrário em conceitos preexistentes (subsunção), e mecânica quando memorizada isolada; organizadores prévios são a ponte cognitiva. O fator isolado mais importante é o que o aluno já sabe.","epistemico":"inferencia","exemplos":[],"id":"aprendizagem_significativa_ausubel","memoria":"semantica","meta":{"autor":"Ausubel","dominio":"ciencias_de_fora","fonte":""},"origem":"wiki","palavras_chave":[],"sinonimos":[],"tipo":"conceito","titulo":"Aprendizagem Significativa (Ausubel)"}
\section{Aprendizagem Significativa (Ausubel)}
Ausubel: a aprendizagem é significativa quando o novo se ancora de modo não-arbitrário em conceitos preexistentes (subsunção), e mecânica quando memorizada isolada; organizadores prévios são a ponte cognitiva. O fator isolado mais importante é o que o aluno já sabe.
\prov{relações: parte\_de teorias\_aprendizagem; referencia hopfield; relaciona assimilacao\_acomodacao}
\prov{proveniência: wiki \textperiodcentered{} inferencia \textperiodcentered{} confiança 1.0 \textperiodcentered{} domínio ciencias\_de\_fora \textperiodcentered{} história 4 versões no jornal}

%CRISTAL {"arestas":[{"alvo":"jusnaturalismo","confianca":1.0,"epistemico":"fato","origem":"wiki","rel":"usa"},{"alvo":"leis_de_deus_programacao","confianca":1.0,"epistemico":"fato","origem":"wiki","rel":"relaciona"}],"confianca":1.0,"contraexemplos":[],"descricao":"Tomás de Aquino: a lei humana deriva por participação da lei natural, que é a participação da criatura racional na lei eterna (a razão de Deus); a lei humana que contradiz a natural é corrupção da lei, não lei.","epistemico":"inferencia","exemplos":[],"id":"aquino_lex","memoria":"semantica","meta":{"autor":"Aquino","dominio":"ciencias_de_fora","fonte":""},"origem":"wiki","palavras_chave":[],"sinonimos":[],"tipo":"conceito","titulo":"Aquino: Lei Eterna, Natural e Humana"}
\section{Aquino: Lei Eterna, Natural e Humana}
Tomás de Aquino: a lei humana deriva por participação da lei natural, que é a participação da criatura racional na lei eterna (a razão de Deus); a lei humana que contradiz a natural é corrupção da lei, não lei.
\prov{relações: usa jusnaturalismo; relaciona leis\_de\_deus\_programacao}
\prov{proveniência: wiki \textperiodcentered{} inferencia \textperiodcentered{} confiança 1.0 \textperiodcentered{} domínio ciencias\_de\_fora \textperiodcentered{} história 3 versões no jornal}

%CRISTAL {"arestas":[{"alvo":"idioma","confianca":1.0,"epistemico":"fato","origem":"wiki","rel":"especializa"}],"confianca":1.0,"contraexemplos":[],"descricao":"Língua semítica de ~400 milhões, litúrgica do Islã — escrita da direita para a esquerda, com forte diglossia entre o clássico e os dialetos.","epistemico":"fato","exemplos":[],"id":"arabe","memoria":"semantica","meta":{"dominio":"ciencias_de_fora","fonte":""},"origem":"wiki","palavras_chave":[],"sinonimos":[],"tipo":"conceito","titulo":"Árabe"}
\section{Árabe}
Língua semítica de \~{}400 milhões, litúrgica do Islã — escrita da direita para a esquerda, com forte diglossia entre o clássico e os dialetos.
\prov{relações: especializa idioma}
\prov{proveniência: wiki \textperiodcentered{} fato \textperiodcentered{} confiança 1.0 \textperiodcentered{} domínio ciencias\_de\_fora \textperiodcentered{} história 2 versões no jornal}

%CRISTAL {"arestas":[{"alvo":"signo_linguistico","confianca":1.0,"epistemico":"fato","origem":"wiki","rel":"usa"}],"confianca":1.0,"contraexemplos":[],"descricao":"Saussure: não há laço natural entre som e conceito — 'árvore' e tree nomeiam o mesmo. A convenção é a lei.","epistemico":"fato","exemplos":[],"id":"arbitrariedade_do_signo","memoria":"semantica","meta":{"autor":"Saussure","dominio":"ciencias_de_fora","fonte":""},"origem":"wiki","palavras_chave":[],"sinonimos":[],"tipo":"conceito","titulo":"Arbitrariedade do Signo"}
\section{Arbitrariedade do Signo}
Saussure: não há laço natural entre som e conceito — 'árvore' e tree nomeiam o mesmo. A convenção é a lei.
\prov{relações: usa signo\_linguistico}
\prov{proveniência: wiki \textperiodcentered{} fato \textperiodcentered{} confiança 1.0 \textperiodcentered{} domínio ciencias\_de\_fora \textperiodcentered{} história 2 versões no jornal}

%CRISTAL {"arestas":[{"alvo":"arte","confianca":1.0,"epistemico":"fato","origem":"wiki","rel":"parte_de"},{"alvo":"vitruvio_triade","confianca":1.0,"epistemico":"fato","origem":"wiki","rel":"usa"}],"confianca":1.0,"contraexemplos":[],"descricao":"A arte e a técnica de construir o espaço habitado — onde função, estrutura e forma se encontram.","epistemico":"fato","exemplos":[],"id":"arquitetura","memoria":"semantica","meta":{"dominio":"ciencias_de_fora","fonte":""},"origem":"wiki","palavras_chave":[],"sinonimos":[],"tipo":"conceito","titulo":"Arquitetura"}
\section{Arquitetura}
A arte e a técnica de construir o espaço habitado — onde função, estrutura e forma se encontram.
\prov{relações: parte\_de arte; usa vitruvio\_triade}
\prov{proveniência: wiki \textperiodcentered{} fato \textperiodcentered{} confiança 1.0 \textperiodcentered{} domínio ciencias\_de\_fora \textperiodcentered{} história 4 versões no jornal}

%CRISTAL {"arestas":[{"alvo":"memoria_programavel","confianca":1.0,"epistemico":"fato","origem":"curriculo","rel":"parte_de"},{"alvo":"computabilidade","confianca":1.0,"epistemico":"fato","origem":"wiki","rel":"relaciona"},{"alvo":"cpu_metal_c","confianca":1.0,"epistemico":"fato","origem":"wiki","rel":"explica"},{"alvo":"hopfield","confianca":1.0,"epistemico":"fato","origem":"wiki","rel":"relaciona"}],"confianca":1.0,"contraexemplos":[],"descricao":"von Neumann: programa e dados na MESMA memória, executados por um ciclo busca-decodifica-executa. O molde de toda CPU — e do broca-so.","epistemico":"fato","exemplos":[],"id":"arquitetura_von_neumann","memoria":"semantica","meta":{"autor":"von Neumann","dominio":"ciencias_de_fora","fonte":""},"origem":"wiki","palavras_chave":["fetch","decode","execute","bus"],"sinonimos":["von neumann","von Neumann","programa armazenado"],"tipo":"conceito","titulo":"Arquitetura de von Neumann"}
\section{Arquitetura de von Neumann}
von Neumann: programa e dados na MESMA memória, executados por um ciclo busca-decodifica-executa. O molde de toda CPU — e do broca-so.
\prov{sinónimos: von neumann; von Neumann; programa armazenado}
\prov{relações: parte\_de memoria\_programavel; relaciona computabilidade; explica cpu\_metal\_c; relaciona hopfield}
\prov{proveniência: wiki \textperiodcentered{} fato \textperiodcentered{} confiança 1.0 \textperiodcentered{} domínio ciencias\_de\_fora \textperiodcentered{} história 49 versões no jornal}

%CRISTAL {"arestas":[{"alvo":"estetica","confianca":1.0,"epistemico":"fato","origem":"wiki","rel":"relaciona"}],"confianca":1.0,"contraexemplos":[],"descricao":"A criação humana que busca forma, sentido e beleza — a expressão que não se reduz à função.","epistemico":"fato","exemplos":[],"id":"arte","memoria":"semantica","meta":{"dominio":"ciencias_de_fora","fonte":""},"origem":"wiki","palavras_chave":[],"sinonimos":[],"tipo":"conceito","titulo":"Arte"}
\section{Arte}
A criação humana que busca forma, sentido e beleza — a expressão que não se reduz à função.
\prov{relações: relaciona estetica}
\prov{proveniência: wiki \textperiodcentered{} fato \textperiodcentered{} confiança 1.0 \textperiodcentered{} domínio ciencias\_de\_fora \textperiodcentered{} história 2 versões no jornal}

%CRISTAL {"arestas":[{"alvo":"artes_visuais","confianca":1.0,"epistemico":"fato","origem":"wiki","rel":"parte_de"},{"alvo":"musica_e_matematica","confianca":1.0,"epistemico":"fato","origem":"wiki","rel":"relaciona"},{"alvo":"abstracao_kandinsky","confianca":1.0,"epistemico":"fato","origem":"wiki","rel":"relaciona"}],"confianca":1.0,"contraexemplos":[],"descricao":"Kandinsky: a forma e a cor libertadas da representação — a pintura como música visual, sem objeto.","epistemico":"inferencia","exemplos":[],"id":"arte_abstrata","memoria":"semantica","meta":{"autor":"Kandinsky","dominio":"ciencias_de_fora","fonte":""},"origem":"wiki","palavras_chave":[],"sinonimos":[],"tipo":"conceito","titulo":"Arte Abstrata"}
\section{Arte Abstrata}
Kandinsky: a forma e a cor libertadas da representação — a pintura como música visual, sem objeto.
\prov{relações: parte\_de artes\_visuais; relaciona musica\_e\_matematica; relaciona abstracao\_kandinsky}
\prov{proveniência: wiki \textperiodcentered{} inferencia \textperiodcentered{} confiança 1.0 \textperiodcentered{} domínio ciencias\_de\_fora \textperiodcentered{} história 5 versões no jornal}

%CRISTAL {"arestas":[{"alvo":"arte","confianca":1.0,"epistemico":"fato","origem":"wiki","rel":"parte_de"},{"alvo":"auto_organizacao","confianca":1.0,"epistemico":"fato","origem":"wiki","rel":"usa"},{"alvo":"linguagem_matriz_gatos","confianca":1.0,"epistemico":"fato","origem":"wiki","rel":"relaciona"}],"confianca":1.0,"contraexemplos":[],"descricao":"A obra produzida por regras e processos — algoritmo, acaso e sistema como pincel. A arte que a máquina também faz.","epistemico":"inferencia","exemplos":[],"id":"arte_generativa","memoria":"semantica","meta":{"dominio":"ciencias_de_fora","fonte":""},"origem":"wiki","palavras_chave":[],"sinonimos":[],"tipo":"conceito","titulo":"Arte Generativa"}
\section{Arte Generativa}
A obra produzida por regras e processos — algoritmo, acaso e sistema como pincel. A arte que a máquina também faz.
\prov{relações: parte\_de arte; usa auto\_organizacao; relaciona linguagem\_matriz\_gatos}
\prov{proveniência: wiki \textperiodcentered{} inferencia \textperiodcentered{} confiança 1.0 \textperiodcentered{} domínio ciencias\_de\_fora \textperiodcentered{} história 4 versões no jornal}

%CRISTAL {"arestas":[{"alvo":"filosofia_da_tecnica","confianca":1.0,"epistemico":"fato","origem":"wiki","rel":"parte_de"},{"alvo":"biopoder","confianca":1.0,"epistemico":"fato","origem":"wiki","rel":"usa"},{"alvo":"contrato_social","confianca":1.0,"epistemico":"fato","origem":"wiki","rel":"relaciona"}],"confianca":1.0,"contraexemplos":[],"descricao":"Winner: artefatos técnicos incorporam relações de poder — podem ordenar autoridade (as pontes baixas de Moses barravam ônibus e, com eles, pobres e negros) ou ser inerentemente políticos, compatíveis só com certos arranjos de poder.","epistemico":"inferencia","exemplos":[],"id":"artefatos_tem_politica","memoria":"semantica","meta":{"autor":"Winner","dominio":"ciencias_de_fora","fonte":""},"origem":"wiki","palavras_chave":[],"sinonimos":[],"tipo":"conceito","titulo":"Os Artefatos têm Política (Winner)"}
\section{Os Artefatos têm Política (Winner)}
Winner: artefatos técnicos incorporam relações de poder — podem ordenar autoridade (as pontes baixas de Moses barravam ônibus e, com eles, pobres e negros) ou ser inerentemente políticos, compatíveis só com certos arranjos de poder.
\prov{relações: parte\_de filosofia\_da\_tecnica; usa biopoder; relaciona contrato\_social}
\prov{proveniência: wiki \textperiodcentered{} inferencia \textperiodcentered{} confiança 1.0 \textperiodcentered{} domínio ciencias\_de\_fora \textperiodcentered{} história 4 versões no jornal}

%CRISTAL {"arestas":[{"alvo":"estetica","confianca":1.0,"epistemico":"fato","origem":"wiki","rel":"parte_de"},{"alvo":"beleza_matematica","confianca":1.0,"epistemico":"fato","origem":"wiki","rel":"relaciona"}],"confianca":1.0,"contraexemplos":[],"descricao":"As artes específicas — música, literatura, artes visuais, cinema, arquitetura — como práticas de criação de sentido e beleza, cada uma com sua matéria e sua linguagem próprias.","epistemico":"fato","exemplos":[],"id":"artes","memoria":"semantica","meta":{"dominio":"ciencias_de_fora","fonte":""},"origem":"wiki","palavras_chave":[],"sinonimos":[],"tipo":"conceito","titulo":"As Artes"}
\section{As Artes}
As artes específicas — música, literatura, artes visuais, cinema, arquitetura — como práticas de criação de sentido e beleza, cada uma com sua matéria e sua linguagem próprias.
\prov{relações: parte\_de estetica; relaciona beleza\_matematica}
\prov{proveniência: wiki \textperiodcentered{} fato \textperiodcentered{} confiança 1.0 \textperiodcentered{} domínio ciencias\_de\_fora \textperiodcentered{} história 3 versões no jornal}

%CRISTAL {"arestas":[{"alvo":"artes","confianca":1.0,"epistemico":"fato","origem":"wiki","rel":"parte_de"},{"alvo":"word","confianca":0.5,"epistemico":"fato","origem":"wiki","rel":"relaciona"},{"alvo":"arte","confianca":1.0,"epistemico":"fato","origem":"wiki","rel":"parte_de"}],"confianca":1.0,"contraexemplos":[],"descricao":"As artes do ver — pintura, escultura, gravura, fotografia; a forma no espaço e na luz.","epistemico":"fato","exemplos":[],"id":"artes_visuais","memoria":"semantica","meta":{"dominio":"ciencias_de_fora","fonte":""},"origem":"wiki","palavras_chave":[],"sinonimos":[],"tipo":"conceito","titulo":"Artes Visuais"}
\section{Artes Visuais}
As artes do ver — pintura, escultura, gravura, fotografia; a forma no espaço e na luz.
\prov{relações: parte\_de artes; relaciona word; parte\_de arte}
\prov{proveniência: wiki \textperiodcentered{} fato \textperiodcentered{} confiança 1.0 \textperiodcentered{} domínio ciencias\_de\_fora \textperiodcentered{} história 5 versões no jornal}

%CRISTAL {"arestas":[{"alvo":"teoria_institucional_dickie","confianca":1.0,"epistemico":"fato","origem":"wiki","rel":"relaciona"},{"alvo":"readymade_duchamp","confianca":1.0,"epistemico":"fato","origem":"wiki","rel":"usa"},{"alvo":"aura_benjamin","confianca":1.0,"epistemico":"fato","origem":"wiki","rel":"relaciona"}],"confianca":1.0,"contraexemplos":[],"descricao":"Danto: o que distingue a Brillo Box de Warhol de uma caixa de supermercado idêntica não é o olho, mas uma teoria e uma história da arte — o 'mundo da arte'. Quando a arte se torna sua própria filosofia, encerra-se seu grande relato histórico.","epistemico":"inferencia","exemplos":[],"id":"artworld_danto","memoria":"semantica","meta":{"autor":"Danto","dominio":"ciencias_de_fora","fonte":""},"origem":"wiki","palavras_chave":[],"sinonimos":[],"tipo":"conceito","titulo":"Artworld e o Fim da Arte (Danto)"}
\section{Artworld e o Fim da Arte (Danto)}
Danto: o que distingue a Brillo Box de Warhol de uma caixa de supermercado idêntica não é o olho, mas uma teoria e uma história da arte — o 'mundo da arte'. Quando a arte se torna sua própria filosofia, encerra-se seu grande relato histórico.
\prov{relações: relaciona teoria\_institucional\_dickie; usa readymade\_duchamp; relaciona aura\_benjamin}
\prov{proveniência: wiki \textperiodcentered{} inferencia \textperiodcentered{} confiança 1.0 \textperiodcentered{} domínio ciencias\_de\_fora \textperiodcentered{} história 4 versões no jornal}

%CRISTAL {"arestas":[{"alvo":"sintaxe","confianca":1.0,"epistemico":"fato","origem":"wiki","rel":"parte_de"},{"alvo":"gramatica_gerativa","confianca":1.0,"epistemico":"fato","origem":"wiki","rel":"usa"},{"alvo":"hierarquia_de_chomsky","confianca":1.0,"epistemico":"fato","origem":"wiki","rel":"usa"},{"alvo":"recursao_linguistica","confianca":1.0,"epistemico":"fato","origem":"wiki","rel":"usa"}],"confianca":1.0,"contraexemplos":[],"descricao":"A estrutura hierárquica de uma frase — sintagmas dentro de sintagmas (SN, SV), não uma fila de palavras. O que a gramática gerativa gera e o parser recupera.","epistemico":"fato","exemplos":[],"id":"arvore_sintatica","memoria":"semantica","meta":{"dominio":"ciencias_de_fora","fonte":""},"origem":"wiki","palavras_chave":[],"sinonimos":[],"tipo":"conceito","titulo":"Árvore Sintática"}
\section{Árvore Sintática}
A estrutura hierárquica de uma frase — sintagmas dentro de sintagmas (SN, SV), não uma fila de palavras. O que a gramática gerativa gera e o parser recupera.
\prov{relações: parte\_de sintaxe; usa gramatica\_gerativa; usa hierarquia\_de\_chomsky; usa recursao\_linguistica}
\prov{proveniência: wiki \textperiodcentered{} fato \textperiodcentered{} confiança 1.0 \textperiodcentered{} domínio ciencias\_de\_fora \textperiodcentered{} história 5 versões no jornal}

%CRISTAL {"arestas":[{"alvo":"ciclos_vico","confianca":1.0,"epistemico":"fato","origem":"wiki","rel":"relaciona"},{"alvo":"word","confianca":0.5,"epistemico":"fato","origem":"wiki","rel":"relaciona"},{"alvo":"virtualizacao","confianca":0.5,"epistemico":"fato","origem":"wiki","rel":"relaciona"}],"confianca":1.0,"contraexemplos":[],"descricao":"Ibn Khaldun: civilizações sobem e caem pela coesão de grupo (asabiyyah); o luxo a dissipa, e um grupo mais coeso ascende. Ciclo dinástico.","epistemico":"hipotese","exemplos":[],"id":"asabiyyah","memoria":"semantica","meta":{"autor":"Ibn Khaldun","dominio":"ciencias_de_fora","fonte":""},"origem":"wiki","palavras_chave":[],"sinonimos":[],"tipo":"conceito","titulo":"Asabiyyah (Ibn Khaldun)"}
\section{Asabiyyah (Ibn Khaldun)}
Ibn Khaldun: civilizações sobem e caem pela coesão de grupo (asabiyyah); o luxo a dissipa, e um grupo mais coeso ascende. Ciclo dinástico.
\prov{relações: relaciona ciclos\_vico; relaciona word; relaciona virtualizacao}
\prov{proveniência: wiki \textperiodcentered{} hipotese \textperiodcentered{} confiança 1.0 \textperiodcentered{} domínio ciencias\_de\_fora \textperiodcentered{} história 4 versões no jornal}

%CRISTAL {"arestas":[{"alvo":"construtivismo_piaget","confianca":1.0,"epistemico":"fato","origem":"wiki","rel":"usa"},{"alvo":"hopfield","confianca":1.0,"epistemico":"fato","origem":"wiki","rel":"referencia"}],"confianca":1.0,"contraexemplos":[],"descricao":"Piaget: aprender é incorporar o novo a esquemas (assimilação) e modificá-los quando falham (acomodação); a equilibração é a auto-regulação que restaura o equilíbrio após o desequilíbrio — como a convergência a um estado estável.","epistemico":"inferencia","exemplos":[],"id":"assimilacao_acomodacao","memoria":"semantica","meta":{"autor":"Piaget","dominio":"ciencias_de_fora","fonte":""},"origem":"wiki","palavras_chave":[],"sinonimos":[],"tipo":"conceito","titulo":"Assimilação, Acomodação e Equilibração"}
\section{Assimilação, Acomodação e Equilibração}
Piaget: aprender é incorporar o novo a esquemas (assimilação) e modificá-los quando falham (acomodação); a equilibração é a auto-regulação que restaura o equilíbrio após o desequilíbrio — como a convergência a um estado estável.
\prov{relações: usa construtivismo\_piaget; referencia hopfield}
\prov{proveniência: wiki \textperiodcentered{} inferencia \textperiodcentered{} confiança 1.0 \textperiodcentered{} domínio ciencias\_de\_fora \textperiodcentered{} história 3 versões no jornal}

%CRISTAL {"arestas":[{"alvo":"mecanismo_de_atencao","confianca":1.0,"epistemico":"fato","origem":"wiki","rel":"explica"},{"alvo":"traducao_neural","confianca":1.0,"epistemico":"fato","origem":"wiki","rel":"usa"},{"alvo":"transformer","confianca":1.0,"epistemico":"fato","origem":"wiki","rel":"explica"}],"confianca":1.0,"contraexemplos":[],"descricao":"O mecanismo de ATENÇÃO — a peça que fundou os transformers e todos os LLMs — foi inventado (Bahdanau, 2014) para a tradução automática: alinhar dinamicamente palavra-fonte↔palavra-alvo. A IA que hoje raciocina nasceu de um problema de tradução.","epistemico":"fato","exemplos":[],"id":"atencao_nasceu_da_traducao","memoria":"semantica","meta":{"autor":"Bahdanau","dominio":"ciencias_de_fora","fonte":""},"origem":"wiki","palavras_chave":[],"sinonimos":[],"tipo":"conceito","titulo":"A Atenção Nasceu da Tradução"}
\section{A Atenção Nasceu da Tradução}
O mecanismo de ATENÇÃO — a peça que fundou os transformers e todos os LLMs — foi inventado (Bahdanau, 2014) para a tradução automática: alinhar dinamicamente palavra-fonte\ensuremath{\leftrightarrow}palavra-alvo. A IA que hoje raciocina nasceu de um problema de tradução.
\prov{relações: explica mecanismo\_de\_atencao; usa traducao\_neural; explica transformer}
\prov{proveniência: wiki \textperiodcentered{} fato \textperiodcentered{} confiança 1.0 \textperiodcentered{} domínio ciencias\_de\_fora \textperiodcentered{} história 4 versões no jornal}

%CRISTAL {"arestas":[{"alvo":"deus_como_vazio_gentil","confianca":1.0,"epistemico":"fato","origem":"wiki","rel":"usa"},{"alvo":"inteligencia_consciencia_relacional","confianca":1.0,"epistemico":"fato","origem":"wiki","rel":"relaciona"},{"alvo":"nirvana","confianca":1.0,"epistemico":"fato","origem":"wiki","rel":"relaciona"}],"confianca":1.0,"contraexemplos":[],"descricao":"Lopes: como num espectro, 'Deus' só comparece numa faixa de vibração da consciência — abaixo os ateus, na faixa os crentes, acima os iluminados (Buda, Osho, Nietzsche), que também dizem 'não existe Deus', mas por excesso, não por deficiência. Deus (construção mental) desaparece, mas resta a Divindade, que só se experiencia.","epistemico":"inferencia","exemplos":[],"id":"ateus_vs_iluminados_gentil","memoria":"semantica","meta":{"autor":"Gentil Lopes","dominio":"sagrado","fonte":""},"origem":"wiki","palavras_chave":[],"sinonimos":[],"tipo":"conceito","titulo":"Ateus × Iluminados: Deus vs Divindade (Lopes)"}
\section{Ateus $\times$ Iluminados: Deus vs Divindade (Lopes)}
Lopes: como num espectro, 'Deus' só comparece numa faixa de vibração da consciência — abaixo os ateus, na faixa os crentes, acima os iluminados (Buda, Osho, Nietzsche), que também dizem 'não existe Deus', mas por excesso, não por deficiência. Deus (construção mental) desaparece, mas resta a Divindade, que só se experiencia.
\prov{relações: usa deus\_como\_vazio\_gentil; relaciona inteligencia\_consciencia\_relacional; relaciona nirvana}
\prov{proveniência: wiki \textperiodcentered{} inferencia \textperiodcentered{} confiança 1.0 \textperiodcentered{} domínio sagrado \textperiodcentered{} história 5 versões no jornal}

%CRISTAL {"arestas":[{"alvo":"atman_brahman","confianca":1.0,"epistemico":"fato","origem":"wiki","rel":"usa"},{"alvo":"consciencia_pura","confianca":1.0,"epistemico":"fato","origem":"wiki","rel":"usa"},{"alvo":"hinduismo","confianca":1.0,"epistemico":"fato","origem":"wiki","rel":"parte_de"}],"confianca":1.0,"contraexemplos":[],"descricao":"Gita (2.23): o Self (Atman) é eterno, não-nascido, imutável — 'as armas não o cortam, o fogo não o queima'; como se trocam vestes gastas por novas, assim o Ser troca de corpo (a reencarnação/samsara). O testemunha que ilumina todos os corpos como o sol único.","epistemico":"inferencia","exemplos":[],"id":"atman_imortal","memoria":"semantica","meta":{"autor":"Bhagavad Gita","dominio":"sagrado","fonte":""},"origem":"wiki","palavras_chave":[],"sinonimos":[],"tipo":"conceito","titulo":"O Atman Imortal"}
\section{O Atman Imortal}
Gita (2.23): o Self (Atman) é eterno, não-nascido, imutável — 'as armas não o cortam, o fogo não o queima'; como se trocam vestes gastas por novas, assim o Ser troca de corpo (a reencarnação/samsara). O testemunha que ilumina todos os corpos como o sol único.
\prov{relações: usa atman\_brahman; usa consciencia\_pura; parte\_de hinduismo}
\prov{proveniência: wiki \textperiodcentered{} inferencia \textperiodcentered{} confiança 1.0 \textperiodcentered{} domínio sagrado \textperiodcentered{} história 4 versões no jornal}

%CRISTAL {"arestas":[{"alvo":"tonalidade","confianca":1.0,"epistemico":"fato","origem":"wiki","rel":"contradiz"}],"confianca":1.0,"contraexemplos":[],"descricao":"Schoenberg aboliu o centro tonal ('emancipação da dissonância') e criou o método dodecafônico: as 12 notas cromáticas numa série em que nenhuma se repete antes de todas soarem, tratando todas como iguais.","epistemico":"inferencia","exemplos":[],"id":"atonalidade_dodecafonismo","memoria":"semantica","meta":{"autor":"Schoenberg","dominio":"ciencias_de_fora","fonte":""},"origem":"wiki","palavras_chave":[],"sinonimos":[],"tipo":"conceito","titulo":"Atonalidade e Dodecafonismo (Schoenberg)"}
\section{Atonalidade e Dodecafonismo (Schoenberg)}
Schoenberg aboliu o centro tonal ('emancipação da dissonância') e criou o método dodecafônico: as 12 notas cromáticas numa série em que nenhuma se repete antes de todas soarem, tratando todas como iguais.
\prov{relações: contradiz tonalidade}
\prov{proveniência: wiki \textperiodcentered{} inferencia \textperiodcentered{} confiança 1.0 \textperiodcentered{} domínio ciencias\_de\_fora \textperiodcentered{} história 2 versões no jornal}

%CRISTAL {"arestas":[{"alvo":"tiffany","confianca":1.0,"epistemico":"fato","origem":"wiki","rel":"explica"},{"alvo":"word","confianca":0.5,"epistemico":"fato","origem":"wiki","rel":"relaciona"},{"alvo":"syscall","confianca":0.5,"epistemico":"fato","origem":"wiki","rel":"relaciona"},{"alvo":"eudaimonia","confianca":0.5,"epistemico":"fato","origem":"wiki","rel":"relaciona"},{"alvo":"pragmatica","confianca":1.0,"epistemico":"fato","origem":"wiki","rel":"parte_de"},{"alvo":"jogos_de_linguagem","confianca":1.0,"epistemico":"fato","origem":"wiki","rel":"relaciona"}],"confianca":1.0,"contraexemplos":[],"descricao":"Austin/Searle: dizer é FAZER — prometer, batizar, ordenar. O performativo não descreve o mundo, altera-o.","epistemico":"fato","exemplos":[],"id":"atos_de_fala","memoria":"semantica","meta":{"autor":"Austin","dominio":"ciencias_de_fora","fonte":""},"origem":"wiki","palavras_chave":[],"sinonimos":[],"tipo":"conceito","titulo":"Atos de Fala"}
\section{Atos de Fala}
Austin/Searle: dizer é FAZER — prometer, batizar, ordenar. O performativo não descreve o mundo, altera-o.
\prov{relações: explica tiffany; relaciona word; relaciona syscall; relaciona eudaimonia; parte\_de pragmatica; relaciona jogos\_de\_linguagem}
\prov{proveniência: wiki \textperiodcentered{} fato \textperiodcentered{} confiança 1.0 \textperiodcentered{} domínio ciencias\_de\_fora \textperiodcentered{} história 9 versões no jornal}

%CRISTAL {"arestas":[{"alvo":"musculo_mineral_hub","confianca":1.0,"epistemico":"fato","origem":"wiki","rel":"usa"},{"alvo":"chicote_e_o_tensor","confianca":1.0,"epistemico":"fato","origem":"wiki","rel":"usa"},{"alvo":"biologia","confianca":1.0,"epistemico":"fato","origem":"wiki","rel":"eh_um"}],"confianca":1.0,"contraexemplos":[],"descricao":"O modelo model/Hitem3d-1783188908717.glb é um scan ESTÁTICO denso (~1.04M vértices, sem rig). Os ATUADORES são o rig que faltava (linguagem/atuadores.py, 6/6): cada músculo facial é um atuador (centro, direção, raio, ganho) que aplica um campo de DESLOCAMENTO Gaussiano aos vértices da sua região, ao longo da direção do músculo — e o gradiente é o STRAIN (o chicote-tensor). 8 músculos FACS. A EXPRESSÃO é a SOMA dos atuadores (blendshape / as Action Units de Ekman): o sorriso é o zigomático, a surpresa é frontal+masseter, a tristeza é depressor+corrugador. Mesh-agnóstico: roda sobre qualquer array de vértices, do sintético ao milhão do modelo real.","epistemico":"fato","exemplos":[],"id":"atuadores_da_face","memoria":"semantica","meta":{"dominio":"biologia","fonte":"atuadores da face"},"origem":"wiki","palavras_chave":[],"sinonimos":[],"tipo":"conceito","titulo":"Os Atuadores da Face (o rig muscular sobre a malha)"}
\section{Os Atuadores da Face (o rig muscular sobre a malha)}
O modelo model/Hitem3d-1783188908717.glb é um scan ESTÁTICO denso (\~{}1.04M vértices, sem rig). Os ATUADORES são o rig que faltava (linguagem/atuadores.py, 6/6): cada músculo facial é um atuador (centro, direção, raio, ganho) que aplica um campo de DESLOCAMENTO Gaussiano aos vértices da sua região, ao longo da direção do músculo — e o gradiente é o STRAIN (o chicote-tensor). 8 músculos FACS. A EXPRESSÃO é a SOMA dos atuadores (blendshape / as Action Units de Ekman): o sorriso é o zigomático, a surpresa é frontal+masseter, a tristeza é depressor+corrugador. Mesh-agnóstico: roda sobre qualquer array de vértices, do sintético ao milhão do modelo real.
\prov{relações: usa musculo\_mineral\_hub; usa chicote\_e\_o\_tensor; eh\_um biologia}
\prov{proveniência: wiki \textperiodcentered{} atuadores da face \textperiodcentered{} fato \textperiodcentered{} confiança 1.0 \textperiodcentered{} domínio biologia \textperiodcentered{} história 4 versões no jornal}

%CRISTAL {"arestas":[{"alvo":"estetica","confianca":1.0,"epistemico":"fato","origem":"wiki","rel":"parte_de"},{"alvo":"virtualizacao","confianca":0.5,"epistemico":"fato","origem":"wiki","rel":"relaciona"},{"alvo":"word","confianca":0.5,"epistemico":"fato","origem":"wiki","rel":"relaciona"},{"alvo":"epistemologia","confianca":0.5,"epistemico":"fato","origem":"wiki","rel":"relaciona"},{"alvo":"octonios_hiperbolicos","confianca":0.5,"epistemico":"fato","origem":"wiki","rel":"relaciona"},{"alvo":"vida-morte","confianca":0.5,"epistemico":"fato","origem":"wiki","rel":"relaciona"}],"confianca":1.0,"contraexemplos":[],"descricao":"Benjamin: a aura é a singularidade irrepetível da obra original; a reprodutibilidade técnica a dissolve e democratiza a arte.","epistemico":"inferencia","exemplos":[],"id":"aura_benjamin","memoria":"semantica","meta":{"autor":"Benjamin","dominio":"ciencias_de_fora","fonte":""},"origem":"wiki","palavras_chave":[],"sinonimos":[],"tipo":"conceito","titulo":"A Aura da Obra"}
\section{A Aura da Obra}
Benjamin: a aura é a singularidade irrepetível da obra original; a reprodutibilidade técnica a dissolve e democratiza a arte.
\prov{relações: parte\_de estetica; relaciona virtualizacao; relaciona word; relaciona epistemologia; relaciona octonios\_hiperbolicos; relaciona vida-morte}
\prov{proveniência: wiki \textperiodcentered{} inferencia \textperiodcentered{} confiança 1.0 \textperiodcentered{} domínio ciencias\_de\_fora \textperiodcentered{} história 8 versões no jornal}

%CRISTAL {"arestas":[{"alvo":"emergencia","confianca":1.0,"epistemico":"fato","origem":"wiki","rel":"usa"},{"alvo":"dissipacao_prigogine","confianca":1.0,"epistemico":"fato","origem":"wiki","rel":"usa"},{"alvo":"parabola_algebra_conversacional","confianca":1.0,"epistemico":"fato","origem":"wiki","rel":"relaciona"}],"confianca":1.0,"contraexemplos":[],"descricao":"Ordem que surge sozinha das interações locais, sem projetista — do enovelamento ao formigueiro. A vida na borda do caos.","epistemico":"inferencia","exemplos":[],"id":"auto_organizacao","memoria":"semantica","meta":{"dominio":"ciencias_de_fora","fonte":""},"origem":"wiki","palavras_chave":[],"sinonimos":[],"tipo":"conceito","titulo":"Auto-organização"}
\section{Auto-organização}
Ordem que surge sozinha das interações locais, sem projetista — do enovelamento ao formigueiro. A vida na borda do caos.
\prov{relações: usa emergencia; usa dissipacao\_prigogine; relaciona parabola\_algebra\_conversacional}
\prov{proveniência: wiki \textperiodcentered{} inferencia \textperiodcentered{} confiança 1.0 \textperiodcentered{} domínio ciencias\_de\_fora \textperiodcentered{} história 4 versões no jornal}

%CRISTAL {"arestas":[{"alvo":"comunicacao_digital","confianca":1.0,"epistemico":"fato","origem":"wiki","rel":"parte_de"},{"alvo":"sociedade_em_rede","confianca":1.0,"epistemico":"fato","origem":"wiki","rel":"usa"},{"alvo":"refeudalizacao_esfera_publica","confianca":1.0,"epistemico":"fato","origem":"wiki","rel":"contradiz"}],"confianca":1.0,"contraexemplos":[],"descricao":"Castells: forma híbrida da internet — alcança potencialmente uma audiência de massa (many-to-many), mas é autogerada no conteúdo, autodirigida na emissão e autosselecionada na recepção, quebrando o monopólio vertical da mídia tradicional.","epistemico":"inferencia","exemplos":[],"id":"autocomunicacao_de_massa","memoria":"semantica","meta":{"autor":"Castells","dominio":"ciencias_de_fora","fonte":""},"origem":"wiki","palavras_chave":[],"sinonimos":[],"tipo":"conceito","titulo":"Autocomunicação de Massa (Castells)"}
\section{Autocomunicação de Massa (Castells)}
Castells: forma híbrida da internet — alcança potencialmente uma audiência de massa (many-to-many), mas é autogerada no conteúdo, autodirigida na emissão e autosselecionada na recepção, quebrando o monopólio vertical da mídia tradicional.
\prov{relações: parte\_de comunicacao\_digital; usa sociedade\_em\_rede; contradiz refeudalizacao\_esfera\_publica}
\prov{proveniência: wiki \textperiodcentered{} inferencia \textperiodcentered{} confiança 1.0 \textperiodcentered{} domínio ciencias\_de\_fora \textperiodcentered{} história 4 versões no jornal}

%CRISTAL {"arestas":[{"alvo":"filosofia_da_computacao","confianca":1.0,"epistemico":"fato","origem":"wiki","rel":"parte_de"},{"alvo":"arquitetura_von_neumann","confianca":1.0,"epistemico":"fato","origem":"wiki","rel":"usa"},{"alvo":"maquina_evolutiva","confianca":1.0,"epistemico":"fato","origem":"wiki","rel":"referencia"},{"alvo":"godel_incompletude","confianca":1.0,"epistemico":"fato","origem":"wiki","rel":"relaciona"}],"confianca":1.0,"contraexemplos":[],"descricao":"von Neumann: um construtor universal que, dado seu próprio projeto numa 'fita', constrói uma cópia de si mais a fita — separa plano (descrição) de construtor, antecipando a lógica DNA/replicação e a autorreferência.","epistemico":"fato","exemplos":[],"id":"automatos_autorreprodutores","memoria":"semantica","meta":{"autor":"von Neumann","dominio":"ciencias_de_fora","fonte":""},"origem":"wiki","palavras_chave":[],"sinonimos":[],"tipo":"conceito","titulo":"Autômatos Autorreprodutores (von Neumann)"}
\section{Autômatos Autorreprodutores (von Neumann)}
von Neumann: um construtor universal que, dado seu próprio projeto numa 'fita', constrói uma cópia de si mais a fita — separa plano (descrição) de construtor, antecipando a lógica DNA/replicação e a autorreferência.
\prov{relações: parte\_de filosofia\_da\_computacao; usa arquitetura\_von\_neumann; referencia maquina\_evolutiva; relaciona godel\_incompletude}
\prov{proveniência: wiki \textperiodcentered{} fato \textperiodcentered{} confiança 1.0 \textperiodcentered{} domínio ciencias\_de\_fora \textperiodcentered{} história 5 versões no jornal}

%CRISTAL {"arestas":[{"alvo":"principialismo","confianca":1.0,"epistemico":"fato","origem":"wiki","rel":"usa"},{"alvo":"imperativo_categorico","confianca":1.0,"epistemico":"fato","origem":"wiki","rel":"usa"},{"alvo":"liberdade_berlin","confianca":1.0,"epistemico":"fato","origem":"wiki","rel":"relaciona"}],"confianca":1.0,"contraexemplos":[],"descricao":"O direito de o indivíduo capaz decidir sobre o próprio corpo e tratamento segundo seus valores, limitando o paternalismo médico — a tradução clínica da autonomia moral kantiana e da liberdade individual.","epistemico":"inferencia","exemplos":[],"id":"autonomia_do_paciente","memoria":"semantica","meta":{"autor":"Beauchamp & Childress","dominio":"ciencias_de_fora","fonte":""},"origem":"wiki","palavras_chave":[],"sinonimos":[],"tipo":"conceito","titulo":"Autonomia do Paciente"}
\section{Autonomia do Paciente}
O direito de o indivíduo capaz decidir sobre o próprio corpo e tratamento segundo seus valores, limitando o paternalismo médico — a tradução clínica da autonomia moral kantiana e da liberdade individual.
\prov{relações: usa principialismo; usa imperativo\_categorico; relaciona liberdade\_berlin}
\prov{proveniência: wiki \textperiodcentered{} inferencia \textperiodcentered{} confiança 1.0 \textperiodcentered{} domínio ciencias\_de\_fora \textperiodcentered{} história 4 versões no jornal}

%CRISTAL {"arestas":[{"alvo":"vida","confianca":1.0,"epistemico":"fato","origem":"wiki","rel":"explica"},{"alvo":"hopfield","confianca":1.0,"epistemico":"fato","origem":"wiki","rel":"referencia"},{"alvo":"broca_os","confianca":1.0,"epistemico":"fato","origem":"wiki","rel":"referencia"},{"alvo":"word","confianca":0.5,"epistemico":"fato","origem":"wiki","rel":"relaciona"},{"alvo":"vontade_de_poder","confianca":0.5,"epistemico":"fato","origem":"wiki","rel":"relaciona"},{"alvo":"o_que_e_vida","confianca":1.0,"epistemico":"fato","origem":"wiki","rel":"especializa"},{"alvo":"auto_organizacao","confianca":1.0,"epistemico":"fato","origem":"wiki","rel":"usa"}],"confianca":1.0,"contraexemplos":[],"descricao":"Maturana/Varela: o vivo é o sistema que se produz a si mesmo — fechado na organização, aberto ao ambiente; não é máquina.","epistemico":"inferencia","exemplos":[],"id":"autopoiese","memoria":"semantica","meta":{"autor":"Maturana/Varela","dominio":"ciencias_de_fora","fonte":""},"origem":"wiki","palavras_chave":[],"sinonimos":[],"tipo":"conceito","titulo":"Autopoiese"}
\section{Autopoiese}
Maturana/Varela: o vivo é o sistema que se produz a si mesmo — fechado na organização, aberto ao ambiente; não é máquina.
\prov{relações: explica vida; referencia hopfield; referencia broca\_os; relaciona word; relaciona vontade\_de\_poder; especializa o\_que\_e\_vida; usa auto\_organizacao}
\prov{proveniência: wiki \textperiodcentered{} inferencia \textperiodcentered{} confiança 1.0 \textperiodcentered{} domínio ciencias\_de\_fora \textperiodcentered{} história 8 versões no jornal}

%CRISTAL {"arestas":[{"alvo":"bazin_ontologia","confianca":1.0,"epistemico":"fato","origem":"wiki","rel":"usa"},{"alvo":"morte_do_autor","confianca":1.0,"epistemico":"fato","origem":"wiki","rel":"contradiz"}],"confianca":1.0,"contraexemplos":[],"descricao":"Truffaut e os Cahiers du cinéma: o diretor é o verdadeiro autor do filme, marcando-o com uma visão pessoal e coerente ao longo de sua obra — em tensão direta com a 'morte do autor' de Barthes.","epistemico":"inferencia","exemplos":[],"id":"autor_auteur","memoria":"semantica","meta":{"autor":"Truffaut","dominio":"ciencias_de_fora","fonte":""},"origem":"wiki","palavras_chave":[],"sinonimos":[],"tipo":"conceito","titulo":"Teoria do Autor / Auteur"}
\section{Teoria do Autor / Auteur}
Truffaut e os Cahiers du cinéma: o diretor é o verdadeiro autor do filme, marcando-o com uma visão pessoal e coerente ao longo de sua obra — em tensão direta com a 'morte do autor' de Barthes.
\prov{relações: usa bazin\_ontologia; contradiz morte\_do\_autor}
\prov{proveniência: wiki \textperiodcentered{} inferencia \textperiodcentered{} confiança 1.0 \textperiodcentered{} domínio ciencias\_de\_fora \textperiodcentered{} história 3 versões no jornal}

%CRISTAL {"arestas":[{"alvo":"comunicacao_iconoclasta_gentil","confianca":1.0,"epistemico":"fato","origem":"wiki","rel":"usa"},{"alvo":"genios_tambem_erram","confianca":1.0,"epistemico":"fato","origem":"wiki","rel":"relaciona"},{"alvo":"espiral_do_silencio","confianca":1.0,"epistemico":"fato","origem":"wiki","rel":"relaciona"},{"alvo":"idiotismo_nao_examinar","confianca":1.0,"epistemico":"fato","origem":"wiki","rel":"relaciona"},{"alvo":"injustica_epistemica","confianca":1.0,"epistemico":"fato","origem":"wiki","rel":"relaciona"}],"confianca":1.0,"contraexemplos":[],"descricao":"Lopes identifica o 'princípio da autoridade' como a real barreira à sua mensagem ('pôr-se do lado da autoridade nos exime de pensar') e narra ter um artigo rejeitado por revista 'em razão da firme convicção do Corpo Editorial' — ao que responde: a negação de um teorema não repousa sobre convicção, mas sobre refutação.","epistemico":"inferencia","exemplos":[],"id":"autoridade_e_rejeicao_gentil","memoria":"semantica","meta":{"autor":"Gentil Lopes","dominio":"ciencias_de_fora","fonte":""},"origem":"wiki","palavras_chave":[],"sinonimos":[],"tipo":"conceito","titulo":"Autoridade e Rejeição Acadêmica (Lopes)"}
\section{Autoridade e Rejeição Acadêmica (Lopes)}
Lopes identifica o 'princípio da autoridade' como a real barreira à sua mensagem ('pôr-se do lado da autoridade nos exime de pensar') e narra ter um artigo rejeitado por revista 'em razão da firme convicção do Corpo Editorial' — ao que responde: a negação de um teorema não repousa sobre convicção, mas sobre refutação.
\prov{relações: usa comunicacao\_iconoclasta\_gentil; relaciona genios\_tambem\_erram; relaciona espiral\_do\_silencio; relaciona idiotismo\_nao\_examinar; relaciona injustica\_epistemica}
\prov{proveniência: wiki \textperiodcentered{} inferencia \textperiodcentered{} confiança 1.0 \textperiodcentered{} domínio ciencias\_de\_fora \textperiodcentered{} história 7 versões no jornal}

%CRISTAL {"arestas":[{"alvo":"ciencia_politica","confianca":1.0,"epistemico":"fato","origem":"wiki","rel":"parte_de"},{"alvo":"estado_de_direito","confianca":1.0,"epistemico":"fato","origem":"wiki","rel":"contradiz"}],"confianca":1.0,"contraexemplos":[],"descricao":"Linz: o autoritarismo limita a esfera política (pluralismo limitado, 'mentalidades', baixa mobilização); o totalitarismo penetra toda a vida por ideologia total, partido único de massas e mobilização permanente.","epistemico":"inferencia","exemplos":[],"id":"autoritarismo_vs_totalitarismo","memoria":"semantica","meta":{"autor":"Linz","dominio":"ciencias_de_fora","fonte":""},"origem":"wiki","palavras_chave":[],"sinonimos":[],"tipo":"conceito","titulo":"Autoritarismo vs Totalitarismo (Linz)"}
\section{Autoritarismo vs Totalitarismo (Linz)}
Linz: o autoritarismo limita a esfera política (pluralismo limitado, 'mentalidades', baixa mobilização); o totalitarismo penetra toda a vida por ideologia total, partido único de massas e mobilização permanente.
\prov{relações: parte\_de ciencia\_politica; contradiz estado\_de\_direito}
\prov{proveniência: wiki \textperiodcentered{} inferencia \textperiodcentered{} confiança 1.0 \textperiodcentered{} domínio ciencias\_de\_fora \textperiodcentered{} história 3 versões no jornal}

%CRISTAL {"arestas":[{"alvo":"axiomas_zfc","confianca":1.0,"epistemico":"fato","origem":"wiki","rel":"parte_de"}],"confianca":1.0,"contraexemplos":[],"descricao":"Pode-se escolher um elemento de cada conjunto de uma família infinita. Poderoso e contraintuitivo (Banach-Tarski).","epistemico":"inferencia","exemplos":[],"id":"axioma_da_escolha","memoria":"semantica","meta":{"dominio":"ciencias_de_fora","fonte":""},"origem":"wiki","palavras_chave":[],"sinonimos":[],"tipo":"conceito","titulo":"Axioma da Escolha"}
\section{Axioma da Escolha}
Pode-se escolher um elemento de cada conjunto de uma família infinita. Poderoso e contraintuitivo (Banach-Tarski).
\prov{relações: parte\_de axiomas\_zfc}
\prov{proveniência: wiki \textperiodcentered{} inferencia \textperiodcentered{} confiança 1.0 \textperiodcentered{} domínio ciencias\_de\_fora \textperiodcentered{} história 4 versões no jornal}

%CRISTAL {"arestas":[{"alvo":"teoria_da_probabilidade","confianca":1.0,"epistemico":"fato","origem":"wiki","rel":"parte_de"}],"confianca":1.0,"contraexemplos":[],"descricao":"Kolmogorov: a probabilidade como medida de peso total 1 — o fundamento rigoroso do acaso.","epistemico":"fato","exemplos":[],"id":"axiomas_de_kolmogorov","memoria":"semantica","meta":{"autor":"Kolmogorov","dominio":"ciencias_de_fora","fonte":""},"origem":"wiki","palavras_chave":[],"sinonimos":[],"tipo":"conceito","titulo":"Axiomas de Kolmogorov"}
\section{Axiomas de Kolmogorov}
Kolmogorov: a probabilidade como medida de peso total 1 — o fundamento rigoroso do acaso.
\prov{relações: parte\_de teoria\_da\_probabilidade}
\prov{proveniência: wiki \textperiodcentered{} fato \textperiodcentered{} confiança 1.0 \textperiodcentered{} domínio ciencias\_de\_fora \textperiodcentered{} história 4 versões no jornal}

%CRISTAL {"arestas":[{"alvo":"paradoxo_de_russell","confianca":1.0,"epistemico":"fato","origem":"wiki","rel":"explica"},{"alvo":"teoria_dos_conjuntos","confianca":1.0,"epistemico":"fato","origem":"wiki","rel":"parte_de"}],"confianca":1.0,"contraexemplos":[],"descricao":"ZFC: o sistema padrão que funda a matemática em conjuntos e evita os paradoxos ingênuos.","epistemico":"fato","exemplos":[],"id":"axiomas_zfc","memoria":"semantica","meta":{"dominio":"ciencias_de_fora","fonte":""},"origem":"wiki","palavras_chave":[],"sinonimos":[],"tipo":"conceito","titulo":"Axiomas de Zermelo-Fraenkel"}
\section{Axiomas de Zermelo-Fraenkel}
ZFC: o sistema padrão que funda a matemática em conjuntos e evita os paradoxos ingênuos.
\prov{relações: explica paradoxo\_de\_russell; parte\_de teoria\_dos\_conjuntos}
\prov{proveniência: wiki \textperiodcentered{} fato \textperiodcentered{} confiança 1.0 \textperiodcentered{} domínio ciencias\_de\_fora \textperiodcentered{} história 6 versões no jornal}

%CRISTAL {"arestas":[{"alvo":"literatura","confianca":1.0,"epistemico":"fato","origem":"wiki","rel":"parte_de"},{"alvo":"intertextualidade","confianca":1.0,"epistemico":"fato","origem":"wiki","rel":"explica"},{"alvo":"morte_do_autor","confianca":1.0,"epistemico":"fato","origem":"wiki","rel":"relaciona"}],"confianca":1.0,"contraexemplos":[],"descricao":"Bakhtin: todo enunciado é dialógico (responde a outros); o romance é polifônico (pluralidade de vozes autônomas), atravessado pelo carnavalesco (inversão cômica) e organizado pelo cronotopo (fusão de tempo e espaço).","epistemico":"inferencia","exemplos":[],"id":"bakhtin_dialogismo","memoria":"semantica","meta":{"autor":"Bakhtin","dominio":"ciencias_de_fora","fonte":""},"origem":"wiki","palavras_chave":[],"sinonimos":[],"tipo":"conceito","titulo":"Dialogismo e Polifonia (Bakhtin)"}
\section{Dialogismo e Polifonia (Bakhtin)}
Bakhtin: todo enunciado é dialógico (responde a outros); o romance é polifônico (pluralidade de vozes autônomas), atravessado pelo carnavalesco (inversão cômica) e organizado pelo cronotopo (fusão de tempo e espaço).
\prov{relações: parte\_de literatura; explica intertextualidade; relaciona morte\_do\_autor}
\prov{proveniência: wiki \textperiodcentered{} inferencia \textperiodcentered{} confiança 1.0 \textperiodcentered{} domínio ciencias\_de\_fora \textperiodcentered{} história 4 versões no jornal}

%CRISTAL {"arestas":[{"alvo":"soberania_vestfaliana","confianca":1.0,"epistemico":"fato","origem":"wiki","rel":"usa"},{"alvo":"word","confianca":0.5,"epistemico":"fato","origem":"wiki","rel":"relaciona"}],"confianca":1.0,"contraexemplos":[],"descricao":"O padrão recorrente em que os Estados se aliam e se contrabalançam para impedir que qualquer um se torne hegemônico, produzindo equilíbrio no sistema anárquico.","epistemico":"inferencia","exemplos":[],"id":"balanca_de_poder","memoria":"semantica","meta":{"autor":"—","dominio":"ciencias_de_fora","fonte":""},"origem":"wiki","palavras_chave":[],"sinonimos":[],"tipo":"conceito","titulo":"Balança de Poder"}
\section{Balança de Poder}
O padrão recorrente em que os Estados se aliam e se contrabalançam para impedir que qualquer um se torne hegemônico, produzindo equilíbrio no sistema anárquico.
\prov{relações: usa soberania\_vestfaliana; relaciona word}
\prov{proveniência: wiki \textperiodcentered{} inferencia \textperiodcentered{} confiança 1.0 \textperiodcentered{} domínio ciencias\_de\_fora \textperiodcentered{} história 3 versões no jornal}

%CRISTAL {"arestas":[{"alvo":"origens_totalitarismo","confianca":1.0,"epistemico":"fato","origem":"wiki","rel":"usa"},{"alvo":"idiotismo_nao_examinar","confianca":1.0,"epistemico":"fato","origem":"wiki","rel":"relaciona"}],"confianca":1.0,"contraexemplos":[],"descricao":"Arendt (Eichmann em Jerusalém): o mal extremo pode ser cometido não por monstros, mas por burocratas comuns marcados por uma 'incapacidade de pensar' e de julgar. Descreve o agente, não a gravidade do ato.","epistemico":"inferencia","exemplos":[],"id":"banalidade_do_mal","memoria":"semantica","meta":{"autor":"Arendt","dominio":"ciencias_de_fora","fonte":""},"origem":"wiki","palavras_chave":[],"sinonimos":[],"tipo":"conceito","titulo":"A Banalidade do Mal (Arendt)"}
\section{A Banalidade do Mal (Arendt)}
Arendt (Eichmann em Jerusalém): o mal extremo pode ser cometido não por monstros, mas por burocratas comuns marcados por uma 'incapacidade de pensar' e de julgar. Descreve o agente, não a gravidade do ato.
\prov{relações: usa origens\_totalitarismo; relaciona idiotismo\_nao\_examinar}
\prov{proveniência: wiki \textperiodcentered{} inferencia \textperiodcentered{} confiança 1.0 \textperiodcentered{} domínio ciencias\_de\_fora \textperiodcentered{} história 3 versões no jornal}

%CRISTAL {"arestas":[{"alvo":"bazin_ontologia","confianca":1.0,"epistemico":"fato","origem":"wiki","rel":"relaciona"},{"alvo":"aura_benjamin","confianca":1.0,"epistemico":"fato","origem":"wiki","rel":"relaciona"},{"alvo":"morte_do_autor","confianca":1.0,"epistemico":"fato","origem":"wiki","rel":"relaciona"}],"confianca":1.0,"contraexemplos":[],"descricao":"Barthes (A Câmara Clara): na fotografia, o studium é a leitura cultural e codificada da imagem; o punctum é o detalhe acidental, não codificado, que 'fere' e prende o espectador de modo singular e subjetivo.","epistemico":"inferencia","exemplos":[],"id":"barthes_punctum","memoria":"semantica","meta":{"autor":"Barthes","dominio":"ciencias_de_fora","fonte":""},"origem":"wiki","palavras_chave":[],"sinonimos":[],"tipo":"conceito","titulo":"Studium vs Punctum (Barthes)"}
\section{Studium vs Punctum (Barthes)}
Barthes (A Câmara Clara): na fotografia, o studium é a leitura cultural e codificada da imagem; o punctum é o detalhe acidental, não codificado, que 'fere' e prende o espectador de modo singular e subjetivo.
\prov{relações: relaciona bazin\_ontologia; relaciona aura\_benjamin; relaciona morte\_do\_autor}
\prov{proveniência: wiki \textperiodcentered{} inferencia \textperiodcentered{} confiança 1.0 \textperiodcentered{} domínio ciencias\_de\_fora \textperiodcentered{} história 4 versões no jornal}

%CRISTAL {"arestas":[{"alvo":"montagem_eisenstein","confianca":1.0,"epistemico":"fato","origem":"wiki","rel":"contradiz"},{"alvo":"aura_benjamin","confianca":1.0,"epistemico":"fato","origem":"wiki","rel":"relaciona"}],"confianca":1.0,"contraexemplos":[],"descricao":"Bazin: a fotografia e o cinema, por sua gênese mecânica, decalcam o real sem intervenção humana, satisfazendo o desejo de vencer o tempo. Daí um cinema de realismo (profundidade de campo, plano-sequência) contra a manipulação da montagem.","epistemico":"inferencia","exemplos":[],"id":"bazin_ontologia","memoria":"semantica","meta":{"autor":"Bazin","dominio":"ciencias_de_fora","fonte":""},"origem":"wiki","palavras_chave":[],"sinonimos":[],"tipo":"conceito","titulo":"Ontologia da Imagem / Realismo (Bazin)"}
\section{Ontologia da Imagem / Realismo (Bazin)}
Bazin: a fotografia e o cinema, por sua gênese mecânica, decalcam o real sem intervenção humana, satisfazendo o desejo de vencer o tempo. Daí um cinema de realismo (profundidade de campo, plano-sequência) contra a manipulação da montagem.
\prov{relações: contradiz montagem\_eisenstein; relaciona aura\_benjamin}
\prov{proveniência: wiki \textperiodcentered{} inferencia \textperiodcentered{} confiança 1.0 \textperiodcentered{} domínio ciencias\_de\_fora \textperiodcentered{} história 3 versões no jornal}

%CRISTAL {"arestas":[{"alvo":"pena_dissuasao","confianca":1.0,"epistemico":"fato","origem":"wiki","rel":"usa"},{"alvo":"direitos_humanos","confianca":1.0,"epistemico":"fato","origem":"wiki","rel":"relaciona"}],"confianca":1.0,"contraexemplos":[],"descricao":"Beccaria (Dos Delitos e das Penas): o criminoso é agente racional que pesa custo e benefício; a pena deve ser legal, proporcional, certa e célere para dissuadir — jamais cruel. Contra a tortura e a pena de morte; a certeza importa mais que a severidade.","epistemico":"inferencia","exemplos":[],"id":"beccaria_reforma_penal","memoria":"semantica","meta":{"autor":"Beccaria","dominio":"ciencias_de_fora","fonte":""},"origem":"wiki","palavras_chave":[],"sinonimos":[],"tipo":"conceito","titulo":"Beccaria: Reforma Penal Iluminista"}
\section{Beccaria: Reforma Penal Iluminista}
Beccaria (Dos Delitos e das Penas): o criminoso é agente racional que pesa custo e benefício; a pena deve ser legal, proporcional, certa e célere para dissuadir — jamais cruel. Contra a tortura e a pena de morte; a certeza importa mais que a severidade.
\prov{relações: usa pena\_dissuasao; relaciona direitos\_humanos}
\prov{proveniência: wiki \textperiodcentered{} inferencia \textperiodcentered{} confiança 1.0 \textperiodcentered{} domínio ciencias\_de\_fora \textperiodcentered{} história 3 versões no jornal}

%CRISTAL {"arestas":[{"alvo":"estetica","confianca":1.0,"epistemico":"fato","origem":"wiki","rel":"parte_de"},{"alvo":"word","confianca":0.5,"epistemico":"fato","origem":"wiki","rel":"relaciona"},{"alvo":"virtualizacao","confianca":0.5,"epistemico":"fato","origem":"wiki","rel":"relaciona"},{"alvo":"vida-morte","confianca":0.5,"epistemico":"fato","origem":"wiki","rel":"relaciona"},{"alvo":"octonios_hiperbolicos","confianca":0.5,"epistemico":"fato","origem":"wiki","rel":"relaciona"}],"confianca":1.0,"contraexemplos":[],"descricao":"Kant: o prazer estético é sem interesse (livre de desejo), universal e necessário — a livre harmonia de imaginação e entendimento.","epistemico":"inferencia","exemplos":[],"id":"beleza_desinteressada_kant","memoria":"semantica","meta":{"autor":"Kant","dominio":"ciencias_de_fora","fonte":""},"origem":"wiki","palavras_chave":[],"sinonimos":[],"tipo":"conceito","titulo":"Beleza Desinteressada"}
\section{Beleza Desinteressada}
Kant: o prazer estético é sem interesse (livre de desejo), universal e necessário — a livre harmonia de imaginação e entendimento.
\prov{relações: parte\_de estetica; relaciona word; relaciona virtualizacao; relaciona vida-morte; relaciona octonios\_hiperbolicos}
\prov{proveniência: wiki \textperiodcentered{} inferencia \textperiodcentered{} confiança 1.0 \textperiodcentered{} domínio ciencias\_de\_fora \textperiodcentered{} história 6 versões no jornal}

%CRISTAL {"arestas":[{"alvo":"estetica","confianca":1.0,"epistemico":"fato","origem":"wiki","rel":"parte_de"},{"alvo":"razao_aurea","confianca":1.0,"epistemico":"fato","origem":"wiki","rel":"usa"},{"alvo":"wigner_eficacia","confianca":1.0,"epistemico":"fato","origem":"wiki","rel":"relaciona"}],"confianca":1.0,"contraexemplos":[],"descricao":"Hardy: 'a beleza é o primeiro teste' — os padrões matemáticos são belos em si (harmonia, elegância), critério do bom matemático.","epistemico":"inferencia","exemplos":[],"id":"beleza_matematica","memoria":"semantica","meta":{"autor":"Hardy","dominio":"ciencias_de_fora","fonte":""},"origem":"wiki","palavras_chave":[],"sinonimos":[],"tipo":"conceito","titulo":"Beleza Matemática"}
\section{Beleza Matemática}
Hardy: 'a beleza é o primeiro teste' — os padrões matemáticos são belos em si (harmonia, elegância), critério do bom matemático.
\prov{relações: parte\_de estetica; usa razao\_aurea; relaciona wigner\_eficacia}
\prov{proveniência: wiki \textperiodcentered{} inferencia \textperiodcentered{} confiança 1.0 \textperiodcentered{} domínio ciencias\_de\_fora \textperiodcentered{} história 4 versões no jornal}

%CRISTAL {"arestas":[{"alvo":"estetica","confianca":1.0,"epistemico":"fato","origem":"wiki","rel":"parte_de"},{"alvo":"alegoria_caverna","confianca":1.0,"epistemico":"fato","origem":"wiki","rel":"usa"},{"alvo":"word","confianca":0.5,"epistemico":"fato","origem":"wiki","rel":"relaciona"},{"alvo":"virtualizacao","confianca":0.5,"epistemico":"fato","origem":"wiki","rel":"relaciona"},{"alvo":"vida-morte","confianca":0.5,"epistemico":"fato","origem":"wiki","rel":"relaciona"},{"alvo":"misticismo_nao_dual","confianca":0.5,"epistemico":"fato","origem":"wiki","rel":"relaciona"}],"confianca":1.0,"contraexemplos":[],"descricao":"Platão: a Beleza é uma Forma eterna e transcendente, que manifesta o Bem — como o Sol ilumina o inteligível.","epistemico":"hipotese","exemplos":[],"id":"belo_platonico","memoria":"semantica","meta":{"autor":"Platão","dominio":"ciencias_de_fora","fonte":""},"origem":"wiki","palavras_chave":[],"sinonimos":[],"tipo":"conceito","titulo":"O Belo (Platão)"}
\section{O Belo (Platão)}
Platão: a Beleza é uma Forma eterna e transcendente, que manifesta o Bem — como o Sol ilumina o inteligível.
\prov{relações: parte\_de estetica; usa alegoria\_caverna; relaciona word; relaciona virtualizacao; relaciona vida-morte; relaciona misticismo\_nao\_dual}
\prov{proveniência: wiki \textperiodcentered{} hipotese \textperiodcentered{} confiança 1.0 \textperiodcentered{} domínio ciencias\_de\_fora \textperiodcentered{} história 7 versões no jornal}

%CRISTAL {"arestas":[{"alvo":"beleza_matematica","confianca":1.0,"epistemico":"fato","origem":"wiki","rel":"usa"},{"alvo":"matematica_arte_engenharia","confianca":1.0,"epistemico":"fato","origem":"wiki","rel":"usa"},{"alvo":"literatura","confianca":1.0,"epistemico":"fato","origem":"wiki","rel":"relaciona"}],"confianca":1.0,"contraexemplos":[],"descricao":"Lopes define operacionalmente sua estética matemática: um 'belo teorema' é aquele cuja demonstração 'envolve algum grau de engenhosidade' e não é trivial. E cita Hardy: a matemática eterna causa satisfação emocional 'como as melhores manifestações da literatura', milhares de anos depois.","epistemico":"inferencia","exemplos":[],"id":"belo_teorema_engenhosidade","memoria":"semantica","meta":{"autor":"Gentil Lopes","dominio":"ciencias_de_fora","fonte":""},"origem":"wiki","palavras_chave":[],"sinonimos":[],"tipo":"conceito","titulo":"O Belo Teorema (Lopes)"}
\section{O Belo Teorema (Lopes)}
Lopes define operacionalmente sua estética matemática: um 'belo teorema' é aquele cuja demonstração 'envolve algum grau de engenhosidade' e não é trivial. E cita Hardy: a matemática eterna causa satisfação emocional 'como as melhores manifestações da literatura', milhares de anos depois.
\prov{relações: usa beleza\_matematica; usa matematica\_arte\_engenharia; relaciona literatura}
\prov{proveniência: wiki \textperiodcentered{} inferencia \textperiodcentered{} confiança 1.0 \textperiodcentered{} domínio ciencias\_de\_fora \textperiodcentered{} história 5 versões no jornal}

%CRISTAL {"arestas":[{"alvo":"voz_fonte_filtro_modos","confianca":1.0,"epistemico":"fato","origem":"wiki","rel":"parte_de"},{"alvo":"numeros_de_pisot","confianca":1.0,"epistemico":"fato","origem":"wiki","rel":"eh_um"},{"alvo":"espectro_irredutivel_floco","confianca":1.0,"epistemico":"fato","origem":"wiki","rel":"usa"}],"confianca":1.0,"contraexemplos":[],"descricao":"O dado real que aterra 'modos = espectro'. O bem-te-vi (model/birds/bem-te-vi.mp3, 2.09s) assobia um tom quase puro em ~2711 Hz — a nota E7. Um só modo dominante: um espectro CRISTAL (uma nota pura, Pisot), contra a voz humana que é multi-formante (mais rica, borda/caos). O pássaro mostra o extremo cristalino do espectro sonoro — uma estrela de um modo só. Verif.: FFT do arquivo → pico dominante 2711 Hz.","epistemico":"fato","exemplos":[],"id":"bemtevi_modo_puro","memoria":"semantica","meta":{"dominio":"biologia","fonte":"voz"},"origem":"wiki","palavras_chave":[],"sinonimos":[],"tipo":"conceito","titulo":"O Bem-te-vi é um Modo Puro (~2.7 kHz)"}
\section{O Bem-te-vi é um Modo Puro (\~{}2.7 kHz)}
O dado real que aterra 'modos = espectro'. O bem-te-vi (model/birds/bem-te-vi.mp3, 2.09s) assobia um tom quase puro em \~{}2711 Hz — a nota E7. Um só modo dominante: um espectro CRISTAL (uma nota pura, Pisot), contra a voz humana que é multi-formante (mais rica, borda/caos). O pássaro mostra o extremo cristalino do espectro sonoro — uma estrela de um modo só. Verif.: FFT do arquivo \ensuremath{\to} pico dominante 2711 Hz.
\prov{relações: parte\_de voz\_fonte\_filtro\_modos; eh\_um numeros\_de\_pisot; usa espectro\_irredutivel\_floco}
\prov{proveniência: wiki \textperiodcentered{} voz \textperiodcentered{} fato \textperiodcentered{} confiança 1.0 \textperiodcentered{} domínio biologia \textperiodcentered{} história 8 versões no jornal}

%CRISTAL {"arestas":[{"alvo":"jogos_animacoes_hub","confianca":1.0,"epistemico":"fato","origem":"wiki","rel":"parte_de"}],"confianca":1.0,"contraexemplos":[],"descricao":"BEN 10 (2005, Man of Action): Ben Tennyson, 10 anos, acha o OMNITRIX — um relógio alienígena no pulso que o transforma em várias formas de alienígenas (Quatro Braços, Chama, XLR8, Fantasmático, …). O Omnitrix guarda o DNA de ~1 milhão de espécies (o Codon Stream); disca-se para escolher o alien e bate-se para transformar; dá TIMEOUT e volta a ser Ben (humano). AZMUTH (um Galvan) é o criador; a MASTER CONTROL destrava todos os aliens de vez; ALIEN X (um Celestialsapien) é quase-onipotente (reescreve a realidade, mas tem 3 consciências que precisam concordar). É uma obra de ficção — lida aqui pelo dicionário com o Omnitrix do broca-so (o Corpo Algébrico).","epistemico":"fato","exemplos":[],"id":"ben10_desenho","memoria":"semantica","meta":{"dominio":"ficcao","fonte":"Ficção (jogos e animações)"},"origem":"wiki","palavras_chave":[],"sinonimos":[],"tipo":"conceito","titulo":"Ben 10 — o desenho animado (o Omnitrix, o relógio que vira qualquer alien)"}
\section{Ben 10 — o desenho animado (o Omnitrix, o relógio que vira qualquer alien)}
BEN 10 (2005, Man of Action): Ben Tennyson, 10 anos, acha o OMNITRIX — um relógio alienígena no pulso que o transforma em várias formas de alienígenas (Quatro Braços, Chama, XLR8, Fantasmático, …). O Omnitrix guarda o DNA de \~{}1 milhão de espécies (o Codon Stream); disca-se para escolher o alien e bate-se para transformar; dá TIMEOUT e volta a ser Ben (humano). AZMUTH (um Galvan) é o criador; a MASTER CONTROL destrava todos os aliens de vez; ALIEN X (um Celestialsapien) é quase-onipotente (reescreve a realidade, mas tem 3 consciências que precisam concordar). É uma obra de ficção — lida aqui pelo dicionário com o Omnitrix do broca-so (o Corpo Algébrico).
\prov{relações: parte\_de jogos\_animacoes\_hub}
\prov{proveniência: wiki \textperiodcentered{} Ficção (jogos e animações) \textperiodcentered{} fato \textperiodcentered{} confiança 1.0 \textperiodcentered{} domínio ficcao \textperiodcentered{} história 2 versões no jornal}

%CRISTAL {"arestas":[{"alvo":"hinduismo","confianca":1.0,"epistemico":"fato","origem":"wiki","rel":"parte_de"},{"alvo":"contemplacao","confianca":1.0,"epistemico":"fato","origem":"wiki","rel":"usa"},{"alvo":"krishna_avatar","confianca":1.0,"epistemico":"fato","origem":"wiki","rel":"usa"}],"confianca":1.0,"contraexemplos":[],"descricao":"Gita: o caminho da devoção amorosa ao Divino pessoal — fixar mente e coração em Deus e oferecer-lhe todas as ações. Krishna o declara o mais acessível; o devoto amado é o equânime, sem inveja, amigo e compassivo por todos.","epistemico":"inferencia","exemplos":[],"id":"bhakti_yoga","memoria":"semantica","meta":{"autor":"Bhagavad Gita","dominio":"sagrado","fonte":""},"origem":"wiki","palavras_chave":[],"sinonimos":[],"tipo":"conceito","titulo":"Bhakti Yoga — o Caminho da Devoção"}
\section{Bhakti Yoga — o Caminho da Devoção}
Gita: o caminho da devoção amorosa ao Divino pessoal — fixar mente e coração em Deus e oferecer-lhe todas as ações. Krishna o declara o mais acessível; o devoto amado é o equânime, sem inveja, amigo e compassivo por todos.
\prov{relações: parte\_de hinduismo; usa contemplacao; usa krishna\_avatar}
\prov{proveniência: wiki \textperiodcentered{} inferencia \textperiodcentered{} confiança 1.0 \textperiodcentered{} domínio sagrado \textperiodcentered{} história 4 versões no jornal}

%CRISTAL {"arestas":[{"alvo":"longue_duree","confianca":1.0,"epistemico":"fato","origem":"wiki","rel":"generaliza"},{"alvo":"selecao_natural","confianca":1.0,"epistemico":"fato","origem":"wiki","rel":"usa"},{"alvo":"cosmologia_quantica","confianca":1.0,"epistemico":"fato","origem":"wiki","rel":"referencia"}],"confianca":1.0,"contraexemplos":[],"descricao":"Christian: a história unificada do Big Bang à cultura, por limiares de complexidade crescente — humanos como capítulo cósmico.","epistemico":"fato","exemplos":[],"id":"big_history","memoria":"semantica","meta":{"autor":"Christian","dominio":"ciencias_de_fora","fonte":""},"origem":"wiki","palavras_chave":[],"sinonimos":[],"tipo":"conceito","titulo":"Big History"}
\section{Big History}
Christian: a história unificada do Big Bang à cultura, por limiares de complexidade crescente — humanos como capítulo cósmico.
\prov{relações: generaliza longue\_duree; usa selecao\_natural; referencia cosmologia\_quantica}
\prov{proveniência: wiki \textperiodcentered{} fato \textperiodcentered{} confiança 1.0 \textperiodcentered{} domínio ciencias\_de\_fora \textperiodcentered{} história 4 versões no jornal}

%CRISTAL {"arestas":[{"alvo":"idioma","confianca":1.0,"epistemico":"fato","origem":"wiki","rel":"usa"},{"alvo":"hipotese_sapir_whorf","confianca":1.0,"epistemico":"fato","origem":"wiki","rel":"relaciona"}],"confianca":1.0,"contraexemplos":[],"descricao":"Dominar duas ou mais línguas — a norma, não a exceção, na maior parte do mundo; reorganiza a cognição e o próprio pensar.","epistemico":"fato","exemplos":[],"id":"bilinguismo","memoria":"semantica","meta":{"dominio":"ciencias_de_fora","fonte":""},"origem":"wiki","palavras_chave":[],"sinonimos":[],"tipo":"conceito","titulo":"Bilinguismo"}
\section{Bilinguismo}
Dominar duas ou mais línguas — a norma, não a exceção, na maior parte do mundo; reorganiza a cognição e o próprio pensar.
\prov{relações: usa idioma; relaciona hipotese\_sapir\_whorf}
\prov{proveniência: wiki \textperiodcentered{} fato \textperiodcentered{} confiança 1.0 \textperiodcentered{} domínio ciencias\_de\_fora \textperiodcentered{} história 3 versões no jornal}

%CRISTAL {"arestas":[{"alvo":"ecologia","confianca":1.0,"epistemico":"fato","origem":"wiki","rel":"parte_de"},{"alvo":"especiacao","confianca":1.0,"epistemico":"fato","origem":"wiki","rel":"usa"}],"confianca":1.0,"contraexemplos":[],"descricao":"A variedade da vida em genes, espécies e ecossistemas — o estoque de soluções que a evolução acumulou.","epistemico":"fato","exemplos":[],"id":"biodiversidade","memoria":"semantica","meta":{"dominio":"ciencias_de_fora","fonte":""},"origem":"wiki","palavras_chave":[],"sinonimos":[],"tipo":"conceito","titulo":"Biodiversidade"}
\section{Biodiversidade}
A variedade da vida em genes, espécies e ecossistemas — o estoque de soluções que a evolução acumulou.
\prov{relações: parte\_de ecologia; usa especiacao}
\prov{proveniência: wiki \textperiodcentered{} fato \textperiodcentered{} confiança 1.0 \textperiodcentered{} domínio ciencias\_de\_fora \textperiodcentered{} história 3 versões no jornal}

%CRISTAL {"arestas":[{"alvo":"medicina","confianca":1.0,"epistemico":"fato","origem":"wiki","rel":"parte_de"},{"alvo":"etica","confianca":1.0,"epistemico":"fato","origem":"wiki","rel":"parte_de"},{"alvo":"word","confianca":0.5,"epistemico":"fato","origem":"wiki","rel":"relaciona"}],"confianca":1.0,"contraexemplos":[],"descricao":"A ética aplicada à vida e à medicina — autonomia, beneficência, justiça e os dilemas do início e do fim da vida, da pesquisa e do melhoramento humano.","epistemico":"fato","exemplos":[],"id":"bioetica","memoria":"semantica","meta":{"dominio":"ciencias_de_fora","fonte":""},"origem":"wiki","palavras_chave":[],"sinonimos":[],"tipo":"conceito","titulo":"Bioética"}
\section{Bioética}
A ética aplicada à vida e à medicina — autonomia, beneficência, justiça e os dilemas do início e do fim da vida, da pesquisa e do melhoramento humano.
\prov{relações: parte\_de medicina; parte\_de etica; relaciona word}
\prov{proveniência: wiki \textperiodcentered{} fato \textperiodcentered{} confiança 1.0 \textperiodcentered{} domínio ciencias\_de\_fora \textperiodcentered{} história 4 versões no jornal}

%CRISTAL {"arestas":[{"alvo":"matematica","confianca":1.0,"epistemico":"fato","origem":"wiki","rel":"referencia"},{"alvo":"filosofia_da_mente","confianca":0.5,"epistemico":"fato","origem":"wiki","rel":"relaciona"},{"alvo":"vida-morte","confianca":0.5,"epistemico":"fato","origem":"wiki","rel":"relaciona"},{"alvo":"virtualizacao","confianca":0.5,"epistemico":"fato","origem":"wiki","rel":"relaciona"},{"alvo":"gentil_12_definicao","confianca":0.5,"epistemico":"fato","origem":"wiki","rel":"relaciona"},{"alvo":"vontade_de_poder","confianca":0.5,"epistemico":"fato","origem":"wiki","rel":"relaciona"},{"alvo":"sequencia_pow_nativa","confianca":0.5,"epistemico":"fato","origem":"wiki","rel":"relaciona"},{"alvo":"word","confianca":0.5,"epistemico":"fato","origem":"wiki","rel":"relaciona"},{"alvo":"pell","confianca":0.5,"epistemico":"fato","origem":"wiki","rel":"relaciona"},{"alvo":"consciencia","confianca":0.5,"epistemico":"fato","origem":"wiki","rel":"relaciona"},{"alvo":"fibonacci_multidimensional","confianca":0.5,"epistemico":"fato","origem":"wiki","rel":"relaciona"},{"alvo":"vida","confianca":1.0,"epistemico":"fato","origem":"wiki","rel":"explica"}],"confianca":1.0,"contraexemplos":[],"descricao":"A ciência da vida — organização, metabolismo, hereditariedade, evolução.","epistemico":"fato","exemplos":[],"id":"biologia","memoria":"semantica","meta":{"dominio":"ciencias_de_fora","fonte":""},"origem":"wiki","palavras_chave":[],"sinonimos":[],"tipo":"conceito","titulo":"Biologia"}
\section{Biologia}
A ciência da vida — organização, metabolismo, hereditariedade, evolução.
\prov{relações: referencia matematica; relaciona filosofia\_da\_mente; relaciona vida-morte; relaciona virtualizacao; relaciona gentil\_12\_definicao; relaciona vontade\_de\_poder; relaciona sequencia\_pow\_nativa; relaciona word; relaciona pell; relaciona consciencia; relaciona fibonacci\_multidimensional; explica vida}
\prov{proveniência: wiki \textperiodcentered{} fato \textperiodcentered{} confiança 1.0 \textperiodcentered{} domínio ciencias\_de\_fora \textperiodcentered{} história 59 versões no jornal}

%CRISTAL {"arestas":[{"alvo":"biologia","confianca":1.0,"epistemico":"fato","origem":"wiki","rel":"parte_de"}],"confianca":1.0,"contraexemplos":[],"descricao":"O estudo da célula, a unidade mínima do vivo — sua estrutura, química e divisão.","epistemico":"fato","exemplos":[],"id":"biologia_celular","memoria":"semantica","meta":{"dominio":"ciencias_de_fora","fonte":""},"origem":"wiki","palavras_chave":[],"sinonimos":[],"tipo":"conceito","titulo":"Biologia Celular"}
\section{Biologia Celular}
O estudo da célula, a unidade mínima do vivo — sua estrutura, química e divisão.
\prov{relações: parte\_de biologia}
\prov{proveniência: wiki \textperiodcentered{} fato \textperiodcentered{} confiança 1.0 \textperiodcentered{} domínio ciencias\_de\_fora \textperiodcentered{} história 2 versões no jornal}

%CRISTAL {"arestas":[{"alvo":"biologia","confianca":1.0,"epistemico":"fato","origem":"wiki","rel":"parte_de"}],"confianca":1.0,"contraexemplos":[],"descricao":"O estudo de como as populações mudam no tempo — a história e o mecanismo da diversidade da vida.","epistemico":"fato","exemplos":[],"id":"biologia_evolutiva","memoria":"semantica","meta":{"dominio":"ciencias_de_fora","fonte":""},"origem":"wiki","palavras_chave":[],"sinonimos":[],"tipo":"conceito","titulo":"Biologia Evolutiva"}
\section{Biologia Evolutiva}
O estudo de como as populações mudam no tempo — a história e o mecanismo da diversidade da vida.
\prov{relações: parte\_de biologia}
\prov{proveniência: wiki \textperiodcentered{} fato \textperiodcentered{} confiança 1.0 \textperiodcentered{} domínio ciencias\_de\_fora \textperiodcentered{} história 2 versões no jornal}

%CRISTAL {"arestas":[{"alvo":"zonas_climaticas","confianca":1.0,"epistemico":"fato","origem":"wiki","rel":"usa"},{"alvo":"selecao_natural","confianca":1.0,"epistemico":"fato","origem":"wiki","rel":"relaciona"}],"confianca":1.0,"contraexemplos":[],"descricao":"Grandes comunidades ecológicas (floresta tropical, savana, deserto, tundra, taiga) definidas por vegetação e fauna características, cuja distribuição global é governada principalmente pelo clima (temperatura e precipitação).","epistemico":"fato","exemplos":[],"id":"biomas","memoria":"semantica","meta":{"autor":"—","dominio":"ciencias_de_fora","fonte":""},"origem":"wiki","palavras_chave":[],"sinonimos":[],"tipo":"conceito","titulo":"Biomas"}
\section{Biomas}
Grandes comunidades ecológicas (floresta tropical, savana, deserto, tundra, taiga) definidas por vegetação e fauna características, cuja distribuição global é governada principalmente pelo clima (temperatura e precipitação).
\prov{relações: usa zonas\_climaticas; relaciona selecao\_natural}
\prov{proveniência: wiki \textperiodcentered{} fato \textperiodcentered{} confiança 1.0 \textperiodcentered{} domínio ciencias\_de\_fora \textperiodcentered{} história 3 versões no jornal}

%CRISTAL {"arestas":[{"alvo":"poder_disciplinar_foucault","confianca":1.0,"epistemico":"fato","origem":"wiki","rel":"especializa"},{"alvo":"religiao","confianca":1.0,"epistemico":"fato","origem":"wiki","rel":"referencia"},{"alvo":"word","confianca":0.5,"epistemico":"fato","origem":"wiki","rel":"relaciona"}],"confianca":1.0,"contraexemplos":[],"descricao":"Foucault: o poder que governa a vida das populações (corpo, saúde, sexo), herdeiro da pastoral cristã e da confissão.","epistemico":"inferencia","exemplos":[],"id":"biopoder","memoria":"semantica","meta":{"autor":"Foucault","dominio":"ciencias_de_fora","fonte":""},"origem":"wiki","palavras_chave":[],"sinonimos":[],"tipo":"conceito","titulo":"Biopoder"}
\section{Biopoder}
Foucault: o poder que governa a vida das populações (corpo, saúde, sexo), herdeiro da pastoral cristã e da confissão.
\prov{relações: especializa poder\_disciplinar\_foucault; referencia religiao; relaciona word}
\prov{proveniência: wiki \textperiodcentered{} inferencia \textperiodcentered{} confiança 1.0 \textperiodcentered{} domínio ciencias\_de\_fora \textperiodcentered{} história 4 versões no jornal}

%CRISTAL {"arestas":[{"alvo":"biologia","confianca":1.0,"epistemico":"fato","origem":"wiki","rel":"parte_de"}],"confianca":1.0,"contraexemplos":[],"descricao":"A química da vida — as moléculas e reações que sustentam o metabolismo.","epistemico":"fato","exemplos":[],"id":"bioquimica","memoria":"semantica","meta":{"dominio":"ciencias_de_fora","fonte":""},"origem":"wiki","palavras_chave":[],"sinonimos":[],"tipo":"conceito","titulo":"Bioquímica"}
\section{Bioquímica}
A química da vida — as moléculas e reações que sustentam o metabolismo.
\prov{relações: parte\_de biologia}
\prov{proveniência: wiki \textperiodcentered{} fato \textperiodcentered{} confiança 1.0 \textperiodcentered{} domínio ciencias\_de\_fora \textperiodcentered{} história 2 versões no jornal}

%CRISTAL {"arestas":[{"alvo":"goldchain","confianca":1.0,"epistemico":"fato","origem":"curriculo","rel":"especializa"},{"alvo":"sha256","confianca":1.0,"epistemico":"fato","origem":"curriculo","rel":"depende"},{"alvo":"integridade_e_norma_conservada","confianca":1.0,"epistemico":"fato","origem":"wiki","rel":"relaciona"}],"confianca":1.0,"contraexemplos":[],"descricao":"Cadeia de blocos/certificados encadeados por hash; detecta adulteração (blockchain interna GoldChain).","epistemico":"fato","exemplos":[],"id":"blockchain","memoria":"semantica","meta":{"dominio":"mercado","nivel":5,"ordem":3},"origem":"curriculo","palavras_chave":["certificado","sha256","auditoria"],"sinonimos":["block chain","ledger encadeado"],"tipo":"conceito","titulo":"blockchain"}
\section{blockchain}
Cadeia de blocos/certificados encadeados por hash; detecta adulteração (blockchain interna GoldChain).
\prov{sinónimos: block chain; ledger encadeado}
\prov{relações: especializa goldchain; depende sha256; relaciona integridade\_e\_norma\_conservada}
\prov{proveniência: curriculo \textperiodcentered{} fato \textperiodcentered{} confiança 1.0 \textperiodcentered{} domínio mercado \textperiodcentered{} história 36 versões no jornal}

%CRISTAL {"arestas":[{"alvo":"comunicacao_digital","confianca":1.0,"epistemico":"fato","origem":"wiki","rel":"parte_de"},{"alvo":"injustica_epistemica","confianca":1.0,"epistemico":"fato","origem":"wiki","rel":"relaciona"},{"alvo":"dois_sistemas_kahneman","confianca":1.0,"epistemico":"fato","origem":"wiki","rel":"relaciona"},{"alvo":"espiral_do_silencio","confianca":1.0,"epistemico":"fato","origem":"wiki","rel":"relaciona"}],"confianca":1.0,"contraexemplos":[],"descricao":"Pariser: algoritmos de personalização encerram o usuário num 'ecossistema pessoal de informação' que reforça crenças prévias e exclui o divergente. A câmara de eco é a versão social: um grupo homofílico amplifica crenças por repetição e desacredita fontes externas. (Efeito debatido empiricamente.)","epistemico":"hipotese","exemplos":[],"id":"bolha_de_filtro","memoria":"semantica","meta":{"autor":"Pariser","dominio":"ciencias_de_fora","fonte":""},"origem":"wiki","palavras_chave":[],"sinonimos":[],"tipo":"conceito","titulo":"Bolha de Filtro e Câmara de Eco"}
\section{Bolha de Filtro e Câmara de Eco}
Pariser: algoritmos de personalização encerram o usuário num 'ecossistema pessoal de informação' que reforça crenças prévias e exclui o divergente. A câmara de eco é a versão social: um grupo homofílico amplifica crenças por repetição e desacredita fontes externas. (Efeito debatido empiricamente.)
\prov{relações: parte\_de comunicacao\_digital; relaciona injustica\_epistemica; relaciona dois\_sistemas\_kahneman; relaciona espiral\_do\_silencio}
\prov{proveniência: wiki \textperiodcentered{} hipotese \textperiodcentered{} confiança 1.0 \textperiodcentered{} domínio ciencias\_de\_fora \textperiodcentered{} história 5 versões no jornal}

%CRISTAL {"arestas":[{"alvo":"didatica_matematica","confianca":1.0,"epistemico":"fato","origem":"wiki","rel":"parte_de"},{"alvo":"obstaculo_epistemologico","confianca":1.0,"epistemico":"fato","origem":"wiki","rel":"usa"}],"confianca":1.0,"contraexemplos":[],"descricao":"Brousseau: a aprendizagem ocorre na interação do aluno com um 'milieu' organizado; o professor não comunica o saber, mas cria situações adidáticas onde o conhecimento é a resposta ótima — e faz a devolução (transfere ao aluno a responsabilidade pelo problema).","epistemico":"inferencia","exemplos":[],"id":"brousseau_situacoes_didaticas","memoria":"semantica","meta":{"autor":"Brousseau","dominio":"ciencias_de_fora","fonte":""},"origem":"wiki","palavras_chave":[],"sinonimos":[],"tipo":"conceito","titulo":"Teoria das Situações Didáticas"}
\section{Teoria das Situações Didáticas}
Brousseau: a aprendizagem ocorre na interação do aluno com um 'milieu' organizado; o professor não comunica o saber, mas cria situações adidáticas onde o conhecimento é a resposta ótima — e faz a devolução (transfere ao aluno a responsabilidade pelo problema).
\prov{relações: parte\_de didatica\_matematica; usa obstaculo\_epistemologico}
\prov{proveniência: wiki \textperiodcentered{} inferencia \textperiodcentered{} confiança 1.0 \textperiodcentered{} domínio ciencias\_de\_fora \textperiodcentered{} história 3 versões no jornal}

%CRISTAL {"arestas":[{"alvo":"legitimidade_weber","confianca":1.0,"epistemico":"fato","origem":"wiki","rel":"usa"},{"alvo":"lei_de_ferro_oligarquia","confianca":1.0,"epistemico":"fato","origem":"wiki","rel":"relaciona"}],"confianca":1.0,"contraexemplos":[],"descricao":"Weber: a burocracia é a forma tecnicamente mais eficiente de dominação legal-racional (hierarquia, regras impessoais, competência técnica) — mas propensa a patologias: rigidez, deslocamento de fins por meios, a 'jaula de ferro' e a tendência oligárquica.","epistemico":"inferencia","exemplos":[],"id":"burocracia_weber","memoria":"semantica","meta":{"autor":"Weber","dominio":"ciencias_de_fora","fonte":""},"origem":"wiki","palavras_chave":[],"sinonimos":[],"tipo":"conceito","titulo":"A Burocracia Racional-Legal (Weber)"}
\section{A Burocracia Racional-Legal (Weber)}
Weber: a burocracia é a forma tecnicamente mais eficiente de dominação legal-racional (hierarquia, regras impessoais, competência técnica) — mas propensa a patologias: rigidez, deslocamento de fins por meios, a 'jaula de ferro' e a tendência oligárquica.
\prov{relações: usa legitimidade\_weber; relaciona lei\_de\_ferro\_oligarquia}
\prov{proveniência: wiki \textperiodcentered{} inferencia \textperiodcentered{} confiança 1.0 \textperiodcentered{} domínio ciencias\_de\_fora \textperiodcentered{} história 3 versões no jornal}

%CRISTAL {"arestas":[{"alvo":"ecossistema","confianca":1.0,"epistemico":"fato","origem":"wiki","rel":"parte_de"},{"alvo":"termodinamica","confianca":1.0,"epistemico":"fato","origem":"wiki","rel":"usa"}],"confianca":1.0,"contraexemplos":[],"descricao":"Quem come quem — o fluxo de energia dos produtores aos predadores, perdendo calor a cada elo (a segunda lei de novo).","epistemico":"fato","exemplos":[],"id":"cadeia_trofica","memoria":"semantica","meta":{"dominio":"ciencias_de_fora","fonte":""},"origem":"wiki","palavras_chave":[],"sinonimos":[],"tipo":"conceito","titulo":"Cadeia Trófica"}
\section{Cadeia Trófica}
Quem come quem — o fluxo de energia dos produtores aos predadores, perdendo calor a cada elo (a segunda lei de novo).
\prov{relações: parte\_de ecossistema; usa termodinamica}
\prov{proveniência: wiki \textperiodcentered{} fato \textperiodcentered{} confiança 1.0 \textperiodcentered{} domínio ciencias\_de\_fora \textperiodcentered{} história 3 versões no jornal}

%CRISTAL {"arestas":[{"alvo":"atonalidade_dodecafonismo","confianca":1.0,"epistemico":"fato","origem":"wiki","rel":"usa"},{"alvo":"teoria_da_informacao_shannon","confianca":1.0,"epistemico":"fato","origem":"wiki","rel":"relaciona"}],"confianca":1.0,"contraexemplos":[],"descricao":"John Cage: uma peça sem sons intencionais que emoldura os ruídos do ambiente como música; expressão do acaso/indeterminação e de uma estética em que compor é redirecionar a escuta.","epistemico":"inferencia","exemplos":[],"id":"cage_silencio","memoria":"semantica","meta":{"autor":"Cage","dominio":"ciencias_de_fora","fonte":""},"origem":"wiki","palavras_chave":[],"sinonimos":[],"tipo":"conceito","titulo":"Cage: 4'33\", Silêncio e Acaso"}
\section{Cage: 4'33", Silêncio e Acaso}
John Cage: uma peça sem sons intencionais que emoldura os ruídos do ambiente como música; expressão do acaso/indeterminação e de uma estética em que compor é redirecionar a escuta.
\prov{relações: usa atonalidade\_dodecafonismo; relaciona teoria\_da\_informacao\_shannon}
\prov{proveniência: wiki \textperiodcentered{} inferencia \textperiodcentered{} confiança 1.0 \textperiodcentered{} domínio ciencias\_de\_fora \textperiodcentered{} história 3 versões no jornal}

%CRISTAL {"arestas":[{"alvo":"analise_matematica","confianca":1.0,"epistemico":"fato","origem":"wiki","rel":"parte_de"}],"confianca":1.0,"contraexemplos":[],"descricao":"Newton e Leibniz: derivada (taxa) e integral (acúmulo), unidas pelo teorema fundamental. A matemática da mudança contínua.","epistemico":"fato","exemplos":[],"id":"calculo_infinitesimal","memoria":"semantica","meta":{"autor":"Newton/Leibniz","dominio":"ciencias_de_fora","fonte":""},"origem":"wiki","palavras_chave":[],"sinonimos":[],"tipo":"conceito","titulo":"Cálculo Infinitesimal"}
\section{Cálculo Infinitesimal}
Newton e Leibniz: derivada (taxa) e integral (acúmulo), unidas pelo teorema fundamental. A matemática da mudança contínua.
\prov{relações: parte\_de analise\_matematica}
\prov{proveniência: wiki \textperiodcentered{} fato \textperiodcentered{} confiança 1.0 \textperiodcentered{} domínio ciencias\_de\_fora \textperiodcentered{} história 4 versões no jornal}

%CRISTAL {"arestas":[{"alvo":"computabilidade","confianca":1.0,"epistemico":"fato","origem":"wiki","rel":"parte_de"},{"alvo":"maquina_de_turing","confianca":1.0,"epistemico":"fato","origem":"wiki","rel":"relaciona"},{"alvo":"linguagem_matriz_gatos","confianca":1.0,"epistemico":"fato","origem":"wiki","rel":"relaciona"},{"alvo":"expressividade_sub_turing","confianca":1.0,"epistemico":"fato","origem":"wiki","rel":"relaciona"}],"confianca":1.0,"contraexemplos":[],"descricao":"Church: computação como pura substituição de funções — sem estado, sem fita. Equivalente à máquina de Turing; a raiz do funcional.","epistemico":"fato","exemplos":[],"id":"calculo_lambda","memoria":"semantica","meta":{"autor":"Church","dominio":"ciencias_de_fora","fonte":""},"origem":"wiki","palavras_chave":[],"sinonimos":[],"tipo":"conceito","titulo":"Cálculo Lambda"}
\section{Cálculo Lambda}
Church: computação como pura substituição de funções — sem estado, sem fita. Equivalente à máquina de Turing; a raiz do funcional.
\prov{relações: parte\_de computabilidade; relaciona maquina\_de\_turing; relaciona linguagem\_matriz\_gatos; relaciona expressividade\_sub\_turing}
\prov{proveniência: wiki \textperiodcentered{} fato \textperiodcentered{} confiança 1.0 \textperiodcentered{} domínio ciencias\_de\_fora \textperiodcentered{} história 12 versões no jornal}

%CRISTAL {"arestas":[{"alvo":"atuadores_da_face","confianca":1.0,"epistemico":"fato","origem":"wiki","rel":"especializa"},{"alvo":"chicote_e_o_tensor","confianca":1.0,"epistemico":"fato","origem":"wiki","rel":"usa"},{"alvo":"biologia","confianca":1.0,"epistemico":"fato","origem":"wiki","rel":"eh_um"}],"confianca":1.0,"contraexemplos":[],"descricao":"Refina a face de uma casca única para um modelo EM CAMADAS (linguagem/camadas.py, 7/7). Da superfície ao fundo: PELE (mole, segue quase tudo), GORDURA (subcutânea, filtra), MÚSCULO (a fonte do strain), OSSO (rígido, ancora — peso 0). Cada camada tem rigidez e peso; densificar cria PONTOS INTERMEDIÁRIOS (subdivisão) e as camadas interiores dão a espessura do tecido (a face deixa de ser uma casca).","epistemico":"fato","exemplos":[],"id":"camadas_da_face","memoria":"semantica","meta":{"dominio":"biologia","fonte":"camadas da face"},"origem":"wiki","palavras_chave":[],"sinonimos":[],"tipo":"conceito","titulo":"As Camadas da Face (pele/gordura/músculo/osso com pesos)"}
\section{As Camadas da Face (pele/gordura/músculo/osso com pesos)}
Refina a face de uma casca única para um modelo EM CAMADAS (linguagem/camadas.py, 7/7). Da superfície ao fundo: PELE (mole, segue quase tudo), GORDURA (subcutânea, filtra), MÚSCULO (a fonte do strain), OSSO (rígido, ancora — peso 0). Cada camada tem rigidez e peso; densificar cria PONTOS INTERMEDIÁRIOS (subdivisão) e as camadas interiores dão a espessura do tecido (a face deixa de ser uma casca).
\prov{relações: especializa atuadores\_da\_face; usa chicote\_e\_o\_tensor; eh\_um biologia}
\prov{proveniência: wiki \textperiodcentered{} camadas da face \textperiodcentered{} fato \textperiodcentered{} confiança 1.0 \textperiodcentered{} domínio biologia \textperiodcentered{} história 4 versões no jornal}

%CRISTAL {"arestas":[{"alvo":"intertextualidade","confianca":1.0,"epistemico":"fato","origem":"wiki","rel":"contradiz"},{"alvo":"aparelho_psiquico","confianca":1.0,"epistemico":"fato","origem":"wiki","rel":"usa"},{"alvo":"literatura","confianca":1.0,"epistemico":"fato","origem":"wiki","rel":"parte_de"}],"confianca":1.0,"contraexemplos":[],"descricao":"Harold Bloom: a criação poética é uma luta agônica do 'poeta forte' contra seus precursores, que ele reescreve e distorce para conquistar espaço próprio; dessa dinâmica emerge o cânone literário.","epistemico":"inferencia","exemplos":[],"id":"canone_angustia_influencia","memoria":"semantica","meta":{"autor":"Bloom","dominio":"ciencias_de_fora","fonte":""},"origem":"wiki","palavras_chave":[],"sinonimos":[],"tipo":"conceito","titulo":"Cânone e Angústia da Influência (Bloom)"}
\section{Cânone e Angústia da Influência (Bloom)}
Harold Bloom: a criação poética é uma luta agônica do 'poeta forte' contra seus precursores, que ele reescreve e distorce para conquistar espaço próprio; dessa dinâmica emerge o cânone literário.
\prov{relações: contradiz intertextualidade; usa aparelho\_psiquico; parte\_de literatura}
\prov{proveniência: wiki \textperiodcentered{} inferencia \textperiodcentered{} confiança 1.0 \textperiodcentered{} domínio ciencias\_de\_fora \textperiodcentered{} história 4 versões no jornal}

%CRISTAL {"arestas":[{"alvo":"teoria_literaria","confianca":1.0,"epistemico":"fato","origem":"wiki","rel":"parte_de"},{"alvo":"morte_do_autor","confianca":1.0,"epistemico":"fato","origem":"wiki","rel":"relaciona"}],"confianca":1.0,"contraexemplos":[],"descricao":"O conjunto de obras tidas como essenciais — sempre disputado: quem escolhe, e o que o cânone exclui?","epistemico":"inferencia","exemplos":[],"id":"canone_literario","memoria":"semantica","meta":{"dominio":"ciencias_de_fora","fonte":""},"origem":"wiki","palavras_chave":[],"sinonimos":[],"tipo":"conceito","titulo":"Cânone Literário"}
\section{Cânone Literário}
O conjunto de obras tidas como essenciais — sempre disputado: quem escolhe, e o que o cânone exclui?
\prov{relações: parte\_de teoria\_literaria; relaciona morte\_do\_autor}
\prov{proveniência: wiki \textperiodcentered{} inferencia \textperiodcentered{} confiança 1.0 \textperiodcentered{} domínio ciencias\_de\_fora \textperiodcentered{} história 3 versões no jornal}

%CRISTAL {"arestas":[{"alvo":"teoria_dos_conjuntos","confianca":1.0,"epistemico":"fato","origem":"wiki","rel":"parte_de"},{"alvo":"problema_da_parada","confianca":1.0,"epistemico":"fato","origem":"wiki","rel":"relaciona"}],"confianca":1.0,"contraexemplos":[],"descricao":"Cantor: há tamanhos distintos de infinito; os reais são incontáveis (argumento da diagonal). O infinito virou objeto.","epistemico":"fato","exemplos":[],"id":"cantor_infinitos","memoria":"semantica","meta":{"autor":"Cantor","dominio":"ciencias_de_fora","fonte":""},"origem":"wiki","palavras_chave":[],"sinonimos":[],"tipo":"conceito","titulo":"Os Infinitos de Cantor"}
\section{Os Infinitos de Cantor}
Cantor: há tamanhos distintos de infinito; os reais são incontáveis (argumento da diagonal). O infinito virou objeto.
\prov{relações: parte\_de teoria\_dos\_conjuntos; relaciona problema\_da\_parada}
\prov{proveniência: wiki \textperiodcentered{} fato \textperiodcentered{} confiança 1.0 \textperiodcentered{} domínio ciencias\_de\_fora \textperiodcentered{} história 6 versões no jornal}

%CRISTAL {"arestas":[{"alvo":"habitus","confianca":1.0,"epistemico":"fato","origem":"wiki","rel":"usa"},{"alvo":"classe_social","confianca":1.0,"epistemico":"fato","origem":"wiki","rel":"relaciona"},{"alvo":"word","confianca":0.5,"epistemico":"fato","origem":"wiki","rel":"relaciona"},{"alvo":"vontade_de_poder","confianca":0.5,"epistemico":"fato","origem":"wiki","rel":"relaciona"},{"alvo":"while","confianca":0.5,"epistemico":"fato","origem":"wiki","rel":"relaciona"},{"alvo":"virtualizacao","confianca":0.5,"epistemico":"fato","origem":"wiki","rel":"relaciona"}],"confianca":1.0,"contraexemplos":[],"descricao":"Bourdieu: conhecimento e gostos herdados que funcionam como moeda social; a escola reproduz a elite disfarçando herança de mérito.","epistemico":"inferencia","exemplos":[],"id":"capital_cultural","memoria":"semantica","meta":{"autor":"Bourdieu","dominio":"ciencias_de_fora","fonte":""},"origem":"wiki","palavras_chave":[],"sinonimos":[],"tipo":"conceito","titulo":"Capital Cultural"}
\section{Capital Cultural}
Bourdieu: conhecimento e gostos herdados que funcionam como moeda social; a escola reproduz a elite disfarçando herança de mérito.
\prov{relações: usa habitus; relaciona classe\_social; relaciona word; relaciona vontade\_de\_poder; relaciona while; relaciona virtualizacao}
\prov{proveniência: wiki \textperiodcentered{} inferencia \textperiodcentered{} confiança 1.0 \textperiodcentered{} domínio ciencias\_de\_fora \textperiodcentered{} história 7 versões no jornal}

%CRISTAL {"arestas":[{"alvo":"ciencia_politica","confianca":1.0,"epistemico":"fato","origem":"wiki","rel":"parte_de"},{"alvo":"criptoeconomia","confianca":1.0,"epistemico":"fato","origem":"wiki","rel":"relaciona"}],"confianca":1.0,"contraexemplos":[],"descricao":"Putnam: redes de engajamento cívico, normas de reciprocidade e confiança que aumentam a eficiência da sociedade ao facilitar a cooperação — regiões com mais capital social têm democracia e economia mais funcionais ('tornam-se ricas porque são cívicas'). É a cultura cívica que sustenta a democracia (Almond & Verba).","epistemico":"inferencia","exemplos":[],"id":"capital_social_putnam","memoria":"semantica","meta":{"autor":"Putnam","dominio":"ciencias_de_fora","fonte":""},"origem":"wiki","palavras_chave":[],"sinonimos":[],"tipo":"conceito","titulo":"Capital Social (Putnam)"}
\section{Capital Social (Putnam)}
Putnam: redes de engajamento cívico, normas de reciprocidade e confiança que aumentam a eficiência da sociedade ao facilitar a cooperação — regiões com mais capital social têm democracia e economia mais funcionais ('tornam-se ricas porque são cívicas'). É a cultura cívica que sustenta a democracia (Almond \& Verba).
\prov{relações: parte\_de ciencia\_politica; relaciona criptoeconomia}
\prov{proveniência: wiki \textperiodcentered{} inferencia \textperiodcentered{} confiança 1.0 \textperiodcentered{} domínio ciencias\_de\_fora \textperiodcentered{} história 3 versões no jornal}

%CRISTAL {"arestas":[{"alvo":"comunicacao_digital","confianca":1.0,"epistemico":"fato","origem":"wiki","rel":"parte_de"},{"alvo":"biopoder","confianca":1.0,"epistemico":"fato","origem":"wiki","rel":"usa"},{"alvo":"tecnopolio","confianca":1.0,"epistemico":"fato","origem":"wiki","rel":"relaciona"},{"alvo":"datificacao","confianca":1.0,"epistemico":"fato","origem":"wiki","rel":"usa"}],"confianca":1.0,"contraexemplos":[],"descricao":"Zuboff: uma ordem econômica que reivindica a experiência humana privada como matéria-prima gratuita, extraindo 'excedente comportamental' convertido em produtos de predição vendidos em 'mercados de futuros comportamentais'.","epistemico":"inferencia","exemplos":[],"id":"capitalismo_de_vigilancia","memoria":"semantica","meta":{"autor":"Zuboff","dominio":"ciencias_de_fora","fonte":""},"origem":"wiki","palavras_chave":[],"sinonimos":[],"tipo":"conceito","titulo":"Capitalismo de Vigilância (Zuboff)"}
\section{Capitalismo de Vigilância (Zuboff)}
Zuboff: uma ordem econômica que reivindica a experiência humana privada como matéria-prima gratuita, extraindo 'excedente comportamental' convertido em produtos de predição vendidos em 'mercados de futuros comportamentais'.
\prov{relações: parte\_de comunicacao\_digital; usa biopoder; relaciona tecnopolio; usa datificacao}
\prov{proveniência: wiki \textperiodcentered{} inferencia \textperiodcentered{} confiança 1.0 \textperiodcentered{} domínio ciencias\_de\_fora \textperiodcentered{} história 5 versões no jornal}

%CRISTAL {"arestas":[{"alvo":"topologia","confianca":1.0,"epistemico":"fato","origem":"wiki","rel":"parte_de"},{"alvo":"euler_prolifico","confianca":1.0,"epistemico":"fato","origem":"wiki","rel":"parte_de"}],"confianca":1.0,"contraexemplos":[],"descricao":"V−A+F=2 para poliedros: um invariante topológico que não muda sob deformação — a primeira 'assinatura' de forma.","epistemico":"fato","exemplos":[],"id":"caracteristica_de_euler","memoria":"semantica","meta":{"autor":"Euler","dominio":"ciencias_de_fora","fonte":""},"origem":"wiki","palavras_chave":[],"sinonimos":[],"tipo":"conceito","titulo":"Característica de Euler"}
\section{Característica de Euler}
V\ensuremath{-}A+F=2 para poliedros: um invariante topológico que não muda sob deformação — a primeira 'assinatura' de forma.
\prov{relações: parte\_de topologia; parte\_de euler\_prolifico}
\prov{proveniência: wiki \textperiodcentered{} fato \textperiodcentered{} confiança 1.0 \textperiodcentered{} domínio ciencias\_de\_fora \textperiodcentered{} história 6 versões no jornal}

%CRISTAL {"arestas":[{"alvo":"teorias_aprendizagem","confianca":1.0,"epistemico":"fato","origem":"wiki","rel":"parte_de"},{"alvo":"dois_sistemas_kahneman","confianca":1.0,"epistemico":"fato","origem":"wiki","rel":"referencia"},{"alvo":"mente_computacao","confianca":1.0,"epistemico":"fato","origem":"wiki","rel":"referencia"}],"confianca":1.0,"contraexemplos":[],"descricao":"Sweller: a aprendizagem é limitada pela capacidade estreita da memória de trabalho; a carga divide-se em intrínseca (dificuldade do conteúdo), extrínseca (imposta pela apresentação) e germânica. Instrução eficaz minimiza a extrínseca.","epistemico":"fato","exemplos":[],"id":"carga_cognitiva","memoria":"semantica","meta":{"autor":"Sweller","dominio":"ciencias_de_fora","fonte":""},"origem":"wiki","palavras_chave":[],"sinonimos":[],"tipo":"conceito","titulo":"Teoria da Carga Cognitiva"}
\section{Teoria da Carga Cognitiva}
Sweller: a aprendizagem é limitada pela capacidade estreita da memória de trabalho; a carga divide-se em intrínseca (dificuldade do conteúdo), extrínseca (imposta pela apresentação) e germânica. Instrução eficaz minimiza a extrínseca.
\prov{relações: parte\_de teorias\_aprendizagem; referencia dois\_sistemas\_kahneman; referencia mente\_computacao}
\prov{proveniência: wiki \textperiodcentered{} fato \textperiodcentered{} confiança 1.0 \textperiodcentered{} domínio ciencias\_de\_fora \textperiodcentered{} história 4 versões no jornal}

%CRISTAL {"arestas":[{"alvo":"sintaxe","confianca":1.0,"epistemico":"fato","origem":"wiki","rel":"usa"},{"alvo":"ordem_das_palavras","confianca":1.0,"epistemico":"fato","origem":"wiki","rel":"relaciona"}],"confianca":1.0,"contraexemplos":[],"descricao":"A marca na palavra que diz seu PAPEL na frase (nominativo=sujeito, acusativo=objeto…) — o latim tem 6, o russo 6, o inglês quase nenhum. Libera a ordem das palavras.","epistemico":"fato","exemplos":[],"id":"caso_gramatical","memoria":"semantica","meta":{"dominio":"ciencias_de_fora","fonte":""},"origem":"wiki","palavras_chave":[],"sinonimos":[],"tipo":"conceito","titulo":"Caso Gramatical"}
\section{Caso Gramatical}
A marca na palavra que diz seu PAPEL na frase (nominativo=sujeito, acusativo=objeto…) — o latim tem 6, o russo 6, o inglês quase nenhum. Libera a ordem das palavras.
\prov{relações: usa sintaxe; relaciona ordem\_das\_palavras}
\prov{proveniência: wiki \textperiodcentered{} fato \textperiodcentered{} confiança 1.0 \textperiodcentered{} domínio ciencias\_de\_fora \textperiodcentered{} história 3 versões no jornal}

%CRISTAL {"arestas":[{"alvo":"etica_pesquisa_humanos","confianca":1.0,"epistemico":"fato","origem":"wiki","rel":"relaciona"},{"alvo":"injustica_epistemica","confianca":1.0,"epistemico":"fato","origem":"wiki","rel":"relaciona"}],"confianca":1.0,"contraexemplos":[],"descricao":"O estudo do serviço de saúde dos EUA (1932-1972) que negou penicilina a centenas de homens negros com sífilis para observar a progressão da doença, enganando-os — símbolo de racismo e violação da autonomia que catalisou a regulação da pesquisa.","epistemico":"fato","exemplos":[],"id":"caso_tuskegee","memoria":"semantica","meta":{"autor":"US PHS","dominio":"ciencias_de_fora","fonte":""},"origem":"wiki","palavras_chave":[],"sinonimos":[],"tipo":"conceito","titulo":"Estudo de Tuskegee (marco negativo)"}
\section{Estudo de Tuskegee (marco negativo)}
O estudo do serviço de saúde dos EUA (1932-1972) que negou penicilina a centenas de homens negros com sífilis para observar a progressão da doença, enganando-os — símbolo de racismo e violação da autonomia que catalisou a regulação da pesquisa.
\prov{relações: relaciona etica\_pesquisa\_humanos; relaciona injustica\_epistemica}
\prov{proveniência: wiki \textperiodcentered{} fato \textperiodcentered{} confiança 1.0 \textperiodcentered{} domínio ciencias\_de\_fora \textperiodcentered{} história 3 versões no jornal}

%CRISTAL {"arestas":[{"alvo":"literatura","confianca":1.0,"epistemico":"fato","origem":"wiki","rel":"parte_de"},{"alvo":"mimese_expressao","confianca":1.0,"epistemico":"fato","origem":"wiki","rel":"usa"},{"alvo":"aparelho_psiquico","confianca":1.0,"epistemico":"fato","origem":"wiki","rel":"relaciona"}],"confianca":1.0,"contraexemplos":[],"descricao":"Aristóteles (Poética): o efeito próprio da tragédia que, suscitando terror e piedade no espectador, produz a purgação/purificação dessas emoções. A arte como imitação (mímese) da ação humana verossímil.","epistemico":"inferencia","exemplos":[],"id":"catarse","memoria":"semantica","meta":{"autor":"Aristóteles","dominio":"ciencias_de_fora","fonte":""},"origem":"wiki","palavras_chave":[],"sinonimos":[],"tipo":"conceito","titulo":"Catarse (Aristóteles)"}
\section{Catarse (Aristóteles)}
Aristóteles (Poética): o efeito próprio da tragédia que, suscitando terror e piedade no espectador, produz a purgação/purificação dessas emoções. A arte como imitação (mímese) da ação humana verossímil.
\prov{relações: parte\_de literatura; usa mimese\_expressao; relaciona aparelho\_psiquico}
\prov{proveniência: wiki \textperiodcentered{} inferencia \textperiodcentered{} confiança 1.0 \textperiodcentered{} domínio ciencias\_de\_fora \textperiodcentered{} história 4 versões no jornal}

%CRISTAL {"arestas":[{"alvo":"geografia_humana","confianca":1.0,"epistemico":"fato","origem":"wiki","rel":"parte_de"},{"alvo":"espaco_santos","confianca":1.0,"epistemico":"fato","origem":"wiki","rel":"relaciona"}],"confianca":1.0,"contraexemplos":[],"descricao":"O aparato conceitual: ESPAÇO (o conjunto de objetos e ações socialmente produzido), LUGAR (o espaço vivido e dotado de sentido), TERRITÓRIO (o espaço apropriado por relações de poder), REGIÃO (recorte que se distingue do entorno), PAISAGEM (o aspecto material visível) e ESCALA (o nível que constitui o próprio fenômeno).","epistemico":"inferencia","exemplos":[],"id":"categorias_geograficas","memoria":"semantica","meta":{"autor":"Milton Santos","dominio":"ciencias_de_fora","fonte":""},"origem":"wiki","palavras_chave":[],"sinonimos":[],"tipo":"conceito","titulo":"Categorias-Chave da Geografia"}
\section{Categorias-Chave da Geografia}
O aparato conceitual: ESPAÇO (o conjunto de objetos e ações socialmente produzido), LUGAR (o espaço vivido e dotado de sentido), TERRITÓRIO (o espaço apropriado por relações de poder), REGIÃO (recorte que se distingue do entorno), PAISAGEM (o aspecto material visível) e ESCALA (o nível que constitui o próprio fenômeno).
\prov{relações: parte\_de geografia\_humana; relaciona espaco\_santos}
\prov{proveniência: wiki \textperiodcentered{} inferencia \textperiodcentered{} confiança 1.0 \textperiodcentered{} domínio ciencias\_de\_fora \textperiodcentered{} história 3 versões no jornal}

%CRISTAL {"arestas":[{"alvo":"epidemiologia","confianca":1.0,"epistemico":"fato","origem":"wiki","rel":"usa"},{"alvo":"historicismo_popper","confianca":1.0,"epistemico":"fato","origem":"wiki","rel":"relaciona"},{"alvo":"dois_sistemas_kahneman","confianca":1.0,"epistemico":"fato","origem":"wiki","rel":"relaciona"}],"confianca":1.0,"contraexemplos":[],"descricao":"Bradford Hill: nove considerações (força, consistência, temporalidade, gradiente biológico, plausibilidade, experimento...) para julgar quando uma associação é causal. O núcleo do 'correlação ≠ causa' — pois o confundimento (uma terceira variável ligada a ambos) gera correlação sem causa.","epistemico":"inferencia","exemplos":[],"id":"causalidade_bradford_hill","memoria":"semantica","meta":{"autor":"Bradford Hill","dominio":"ciencias_de_fora","fonte":""},"origem":"wiki","palavras_chave":[],"sinonimos":[],"tipo":"conceito","titulo":"Critérios de Bradford Hill"}
\section{Critérios de Bradford Hill}
Bradford Hill: nove considerações (força, consistência, temporalidade, gradiente biológico, plausibilidade, experimento...) para julgar quando uma associação é causal. O núcleo do 'correlação \ensuremath{\ne} causa' — pois o confundimento (uma terceira variável ligada a ambos) gera correlação sem causa.
\prov{relações: usa epidemiologia; relaciona historicismo\_popper; relaciona dois\_sistemas\_kahneman}
\prov{proveniência: wiki \textperiodcentered{} inferencia \textperiodcentered{} confiança 1.0 \textperiodcentered{} domínio ciencias\_de\_fora \textperiodcentered{} história 4 versões no jornal}

%CRISTAL {"arestas":[{"alvo":"teoria_celular","confianca":1.0,"epistemico":"fato","origem":"wiki","rel":"usa"},{"alvo":"biologia_celular","confianca":1.0,"epistemico":"fato","origem":"wiki","rel":"parte_de"}],"confianca":1.0,"contraexemplos":[],"descricao":"A unidade mínima da vida: membrana que separa um dentro químico do fora, capaz de metabolismo e cópia.","epistemico":"fato","exemplos":[],"id":"celula","memoria":"semantica","meta":{"dominio":"ciencias_de_fora","fonte":""},"origem":"wiki","palavras_chave":[],"sinonimos":[],"tipo":"conceito","titulo":"Célula"}
\section{Célula}
A unidade mínima da vida: membrana que separa um dentro químico do fora, capaz de metabolismo e cópia.
\prov{relações: usa teoria\_celular; parte\_de biologia\_celular}
\prov{proveniência: wiki \textperiodcentered{} fato \textperiodcentered{} confiança 1.0 \textperiodcentered{} domínio ciencias\_de\_fora \textperiodcentered{} história 3 versões no jornal}

%CRISTAL {"arestas":[{"alvo":"biologia_celular","confianca":1.0,"epistemico":"fato","origem":"wiki","rel":"eh_um"},{"alvo":"corpo_estelar","confianca":1.0,"epistemico":"fato","origem":"wiki","rel":"usa"},{"alvo":"mitocondria","confianca":1.0,"epistemico":"fato","origem":"wiki","rel":"relaciona"},{"alvo":"potencial_de_acao","confianca":1.0,"epistemico":"fato","origem":"wiki","rel":"usa"},{"alvo":"o_neuronio","confianca":1.0,"epistemico":"fato","origem":"wiki","rel":"usa"},{"alvo":"floco_recipiente_termico","confianca":1.0,"epistemico":"fato","origem":"wiki","rel":"usa"},{"alvo":"bai_delta","confianca":1.0,"epistemico":"fato","origem":"wiki","rel":"relaciona"}],"confianca":1.0,"contraexemplos":[],"descricao":"Descer à biologia: a célula, átomo do tecido, lida pela máquina mineral (linguagem/celula_mineral.py, 6/6). Quatro leituras: HOMEOSTASE = o cristal (ponto fixo contrátil, volta ao equilíbrio); EXCITABILIDADE = a borda (o potencial de ação, o limiar do neurônio); o CÓDIGO GENÉTICO = β-numeração base 4 (64 códons); o METABOLISMO = o recipiente térmico (ATP=δ mantém a ordem). A saúde da célula é a estabilidade da sua dinâmica; a doença/morte é o caos (ρ>1). Prepara o tecido (muitas células) e o músculo (células excitáveis contraindo).","epistemico":"inferencia","exemplos":[],"id":"celula_mineral_hub","memoria":"semantica","meta":{"dominio":"biologia","fonte":"célula mineral"},"origem":"wiki","palavras_chave":[],"sinonimos":[],"tipo":"conceito","titulo":"A Célula pela Máquina Mineral"}
\section{A Célula pela Máquina Mineral}
Descer à biologia: a célula, átomo do tecido, lida pela máquina mineral (linguagem/celula\_mineral.py, 6/6). Quatro leituras: HOMEOSTASE = o cristal (ponto fixo contrátil, volta ao equilíbrio); EXCITABILIDADE = a borda (o potencial de ação, o limiar do neurônio); o CÓDIGO GENÉTICO = \ensuremath{\beta}-numeração base 4 (64 códons); o METABOLISMO = o recipiente térmico (ATP=\ensuremath{\delta} mantém a ordem). A saúde da célula é a estabilidade da sua dinâmica; a doença/morte é o caos (\ensuremath{\rho}>1). Prepara o tecido (muitas células) e o músculo (células excitáveis contraindo).
\prov{relações: eh\_um biologia\_celular; usa corpo\_estelar; relaciona mitocondria; usa potencial\_de\_acao; usa o\_neuronio; usa floco\_recipiente\_termico; relaciona bai\_delta}
\prov{proveniência: wiki \textperiodcentered{} célula mineral \textperiodcentered{} inferencia \textperiodcentered{} confiança 1.0 \textperiodcentered{} domínio biologia \textperiodcentered{} história 8 versões no jornal}

%CRISTAL {"arestas":[{"alvo":"neurociencia","confianca":1.0,"epistemico":"fato","origem":"wiki","rel":"parte_de"},{"alvo":"neuronio","confianca":1.0,"epistemico":"fato","origem":"wiki","rel":"usa"},{"alvo":"consciencia","confianca":1.0,"epistemico":"fato","origem":"wiki","rel":"explica"},{"alvo":"o_cerebro_a_memoria_que_ressoa_hopfield","confianca":1.0,"epistemico":"fato","origem":"wiki","rel":"relaciona"}],"confianca":1.0,"contraexemplos":[],"descricao":"A rede de ~86 bilhões de neurônios de onde emergem percepção, memória e o self — o órgão que se modela a si mesmo.","epistemico":"fato","exemplos":[],"id":"cerebro","memoria":"semantica","meta":{"dominio":"ciencias_de_fora","fonte":""},"origem":"wiki","palavras_chave":[],"sinonimos":[],"tipo":"conceito","titulo":"Cérebro"}
\section{Cérebro}
A rede de \~{}86 bilhões de neurônios de onde emergem percepção, memória e o self — o órgão que se modela a si mesmo.
\prov{relações: parte\_de neurociencia; usa neuronio; explica consciencia; relaciona o\_cerebro\_a\_memoria\_que\_ressoa\_hopfield}
\prov{proveniência: wiki \textperiodcentered{} fato \textperiodcentered{} confiança 1.0 \textperiodcentered{} domínio ciencias\_de\_fora \textperiodcentered{} história 5 versões no jornal}

%CRISTAL {"arestas":[{"alvo":"biologia_evolutiva","confianca":1.0,"epistemico":"fato","origem":"wiki","rel":"parte_de"}],"confianca":1.0,"contraexemplos":[],"descricao":"O naturalista que unificou a biologia: a diversidade da vida explicada por descendência com modificação.","epistemico":"fato","exemplos":[],"id":"charles_darwin","memoria":"semantica","meta":{"autor":"Darwin","dominio":"ciencias_de_fora","fonte":""},"origem":"wiki","palavras_chave":[],"sinonimos":[],"tipo":"conceito","titulo":"Charles Darwin"}
\section{Charles Darwin}
O naturalista que unificou a biologia: a diversidade da vida explicada por descendência com modificação.
\prov{relações: parte\_de biologia\_evolutiva}
\prov{proveniência: wiki \textperiodcentered{} fato \textperiodcentered{} confiança 1.0 \textperiodcentered{} domínio ciencias\_de\_fora \textperiodcentered{} história 2 versões no jornal}

%CRISTAL {"arestas":[{"alvo":"semiotica","confianca":1.0,"epistemico":"fato","origem":"wiki","rel":"explica"},{"alvo":"algebra_de_peirce","confianca":1.0,"epistemico":"fato","origem":"wiki","rel":"relaciona"},{"alvo":"distincao_peirce_filosofo_vs_algebra_broca-so","confianca":1.0,"epistemico":"fato","origem":"wiki","rel":"explica"},{"alvo":"pragmatismo","confianca":1.0,"epistemico":"fato","origem":"wiki","rel":"parte_de"},{"alvo":"semiotica_triadica_de_charles_sanders_peirce","confianca":1.0,"epistemico":"fato","origem":"wiki","rel":"relaciona"}],"confianca":1.0,"contraexemplos":[],"descricao":"Pai da semiótica e do pragmatismo, e autor da álgebra de relações que o broca-so usa — o filósofo e a álgebra.","epistemico":"fato","exemplos":[],"id":"charles_peirce","memoria":"semantica","meta":{"autor":"Peirce","dominio":"ciencias_de_fora","fonte":""},"origem":"wiki","palavras_chave":[],"sinonimos":[],"tipo":"conceito","titulo":"Charles Sanders Peirce"}
\section{Charles Sanders Peirce}
Pai da semiótica e do pragmatismo, e autor da álgebra de relações que o broca-so usa — o filósofo e a álgebra.
\prov{relações: explica semiotica; relaciona algebra\_de\_peirce; explica distincao\_peirce\_filosofo\_vs\_algebra\_broca-so; parte\_de pragmatismo; relaciona semiotica\_triadica\_de\_charles\_sanders\_peirce}
\prov{proveniência: wiki \textperiodcentered{} fato \textperiodcentered{} confiança 1.0 \textperiodcentered{} domínio ciencias\_de\_fora \textperiodcentered{} história 6 versões no jornal}

%CRISTAL {"arestas":[{"alvo":"comunicacao_digital","confianca":1.0,"epistemico":"fato","origem":"wiki","rel":"parte_de"},{"alvo":"sociedade_em_rede","confianca":1.0,"epistemico":"fato","origem":"wiki","rel":"relaciona"}],"confianca":1.0,"contraexemplos":[],"descricao":"Pierre Lévy: o espaço de comunicação aberto pela interconexão mundial dos computadores e memórias digitais — ambiente do virtual onde se realiza a inteligência coletiva e emergem novas formas de laço social.","epistemico":"inferencia","exemplos":[],"id":"ciberespaco","memoria":"semantica","meta":{"autor":"Lévy","dominio":"ciencias_de_fora","fonte":""},"origem":"wiki","palavras_chave":[],"sinonimos":[],"tipo":"conceito","titulo":"O Ciberespaço (Lévy)"}
\section{O Ciberespaço (Lévy)}
Pierre Lévy: o espaço de comunicação aberto pela interconexão mundial dos computadores e memórias digitais — ambiente do virtual onde se realiza a inteligência coletiva e emergem novas formas de laço social.
\prov{relações: parte\_de comunicacao\_digital; relaciona sociedade\_em\_rede}
\prov{proveniência: wiki \textperiodcentered{} inferencia \textperiodcentered{} confiança 1.0 \textperiodcentered{} domínio ciencias\_de\_fora \textperiodcentered{} história 3 versões no jornal}

%CRISTAL {"arestas":[{"alvo":"filosofia_da_tecnica","confianca":1.0,"epistemico":"fato","origem":"wiki","rel":"parte_de"},{"alvo":"word","confianca":0.5,"epistemico":"fato","origem":"wiki","rel":"relaciona"}],"confianca":1.0,"contraexemplos":[],"descricao":"Wiener: a ciência do controle e da comunicação no animal e na máquina — como sistemas se autorregulam pela circulação de informação e feedback.","epistemico":"fato","exemplos":[],"id":"cibernetica","memoria":"semantica","meta":{"autor":"Wiener","dominio":"ciencias_de_fora","fonte":""},"origem":"wiki","palavras_chave":[],"sinonimos":[],"tipo":"conceito","titulo":"Cibernética"}
\section{Cibernética}
Wiener: a ciência do controle e da comunicação no animal e na máquina — como sistemas se autorregulam pela circulação de informação e feedback.
\prov{relações: parte\_de filosofia\_da\_tecnica; relaciona word}
\prov{proveniência: wiki \textperiodcentered{} fato \textperiodcentered{} confiança 1.0 \textperiodcentered{} domínio ciencias\_de\_fora \textperiodcentered{} história 3 versões no jornal}

%CRISTAL {"arestas":[{"alvo":"ciclos_biogeoquimicos","confianca":1.0,"epistemico":"fato","origem":"wiki","rel":"especializa"},{"alvo":"mudanca_climatica","confianca":1.0,"epistemico":"fato","origem":"wiki","rel":"explica"}],"confianca":1.0,"contraexemplos":[],"descricao":"A troca de carbono entre atmosfera, oceanos, biosfera e litosfera via fotossíntese, respiração, decomposição e sedimentação; sua perturbação antropogênica (queima de fósseis) é o principal vetor da mudança climática.","epistemico":"fato","exemplos":[],"id":"ciclo_carbono","memoria":"semantica","meta":{"autor":"—","dominio":"ciencias_de_fora","fonte":""},"origem":"wiki","palavras_chave":[],"sinonimos":[],"tipo":"conceito","titulo":"Ciclo do Carbono"}
\section{Ciclo do Carbono}
A troca de carbono entre atmosfera, oceanos, biosfera e litosfera via fotossíntese, respiração, decomposição e sedimentação; sua perturbação antropogênica (queima de fósseis) é o principal vetor da mudança climática.
\prov{relações: especializa ciclos\_biogeoquimicos; explica mudanca\_climatica}
\prov{proveniência: wiki \textperiodcentered{} fato \textperiodcentered{} confiança 1.0 \textperiodcentered{} domínio ciencias\_de\_fora \textperiodcentered{} história 3 versões no jornal}

%CRISTAL {"arestas":[{"alvo":"teorias_aprendizagem","confianca":1.0,"epistemico":"fato","origem":"wiki","rel":"parte_de"},{"alvo":"aprender_fazendo_dewey","confianca":1.0,"epistemico":"fato","origem":"wiki","rel":"usa"}],"confianca":1.0,"contraexemplos":[],"descricao":"Kolb: a aprendizagem é um ciclo de quatro fases — experiência concreta, observação reflexiva, conceituação abstrata e experimentação ativa. (A tipologia de 'estilos de aprendizagem' derivada tem baixo suporte empírico.)","epistemico":"inferencia","exemplos":[],"id":"ciclo_experiencial_kolb","memoria":"semantica","meta":{"autor":"Kolb","dominio":"ciencias_de_fora","fonte":""},"origem":"wiki","palavras_chave":[],"sinonimos":[],"tipo":"conceito","titulo":"Ciclo de Aprendizagem Experiencial (Kolb)"}
\section{Ciclo de Aprendizagem Experiencial (Kolb)}
Kolb: a aprendizagem é um ciclo de quatro fases — experiência concreta, observação reflexiva, conceituação abstrata e experimentação ativa. (A tipologia de 'estilos de aprendizagem' derivada tem baixo suporte empírico.)
\prov{relações: parte\_de teorias\_aprendizagem; usa aprender\_fazendo\_dewey}
\prov{proveniência: wiki \textperiodcentered{} inferencia \textperiodcentered{} confiança 1.0 \textperiodcentered{} domínio ciencias\_de\_fora \textperiodcentered{} história 3 versões no jornal}

%CRISTAL {"arestas":[{"alvo":"sistema_terra","confianca":1.0,"epistemico":"fato","origem":"wiki","rel":"parte_de"},{"alvo":"dissipacao_prigogine","confianca":1.0,"epistemico":"fato","origem":"wiki","rel":"usa"},{"alvo":"neguentropia","confianca":1.0,"epistemico":"fato","origem":"wiki","rel":"relaciona"}],"confianca":1.0,"contraexemplos":[],"descricao":"Percursos cíclicos e essencialmente fechados pelos quais a matéria (água, carbono, nitrogênio, fósforo) circula entre reservatórios vivos e não vivos do planeta, movidos pelo fluxo unidirecional de energia — a matéria recicla, a energia dissipa.","epistemico":"fato","exemplos":[],"id":"ciclos_biogeoquimicos","memoria":"semantica","meta":{"autor":"Vernadsky","dominio":"ciencias_de_fora","fonte":""},"origem":"wiki","palavras_chave":[],"sinonimos":[],"tipo":"conceito","titulo":"Ciclos Biogeoquímicos"}
\section{Ciclos Biogeoquímicos}
Percursos cíclicos e essencialmente fechados pelos quais a matéria (água, carbono, nitrogênio, fósforo) circula entre reservatórios vivos e não vivos do planeta, movidos pelo fluxo unidirecional de energia — a matéria recicla, a energia dissipa.
\prov{relações: parte\_de sistema\_terra; usa dissipacao\_prigogine; relaciona neguentropia}
\prov{proveniência: wiki \textperiodcentered{} fato \textperiodcentered{} confiança 1.0 \textperiodcentered{} domínio ciencias\_de\_fora \textperiodcentered{} história 4 versões no jornal}

%CRISTAL {"arestas":[{"alvo":"filosofia_da_historia","confianca":1.0,"epistemico":"fato","origem":"wiki","rel":"parte_de"},{"alvo":"virtualizacao","confianca":0.5,"epistemico":"fato","origem":"wiki","rel":"relaciona"},{"alvo":"word","confianca":0.5,"epistemico":"fato","origem":"wiki","rel":"relaciona"}],"confianca":1.0,"contraexemplos":[],"descricao":"Vico: a história cicla — idade dos deuses, dos heróis, dos homens — e retorna ao caos para recomeçar. Não linear.","epistemico":"hipotese","exemplos":[],"id":"ciclos_vico","memoria":"semantica","meta":{"autor":"Vico","dominio":"ciencias_de_fora","fonte":""},"origem":"wiki","palavras_chave":[],"sinonimos":[],"tipo":"conceito","titulo":"Corsi e Ricorsi (Vico)"}
\section{Corsi e Ricorsi (Vico)}
Vico: a história cicla — idade dos deuses, dos heróis, dos homens — e retorna ao caos para recomeçar. Não linear.
\prov{relações: parte\_de filosofia\_da\_historia; relaciona virtualizacao; relaciona word}
\prov{proveniência: wiki \textperiodcentered{} hipotese \textperiodcentered{} confiança 1.0 \textperiodcentered{} domínio ciencias\_de\_fora \textperiodcentered{} história 4 versões no jornal}

%CRISTAL {"arestas":[{"alvo":"democracia_formas","confianca":1.0,"epistemico":"fato","origem":"wiki","rel":"usa"},{"alvo":"liberdade_berlin","confianca":1.0,"epistemico":"fato","origem":"wiki","rel":"relaciona"},{"alvo":"estado_de_direito","confianca":1.0,"epistemico":"fato","origem":"wiki","rel":"relaciona"}],"confianca":1.0,"contraexemplos":[],"descricao":"Marshall: o status de membro pleno de uma comunidade, em três elementos que se desenvolvem em sequência histórica — direitos civis (séc. XVIII), políticos (séc. XIX) e sociais (séc. XX, bem-estar e serviços).","epistemico":"inferencia","exemplos":[],"id":"cidadania_marshall","memoria":"semantica","meta":{"autor":"Marshall","dominio":"ciencias_de_fora","fonte":""},"origem":"wiki","palavras_chave":[],"sinonimos":[],"tipo":"conceito","titulo":"Cidadania (T. H. Marshall)"}
\section{Cidadania (T. H. Marshall)}
Marshall: o status de membro pleno de uma comunidade, em três elementos que se desenvolvem em sequência histórica — direitos civis (séc. XVIII), políticos (séc. XIX) e sociais (séc. XX, bem-estar e serviços).
\prov{relações: usa democracia\_formas; relaciona liberdade\_berlin; relaciona estado\_de\_direito}
\prov{proveniência: wiki \textperiodcentered{} inferencia \textperiodcentered{} confiança 1.0 \textperiodcentered{} domínio ciencias\_de\_fora \textperiodcentered{} história 4 versões no jornal}

%CRISTAL {"arestas":[{"alvo":"politica","confianca":1.0,"epistemico":"fato","origem":"wiki","rel":"parte_de"},{"alvo":"sociologia","confianca":1.0,"epistemico":"fato","origem":"wiki","rel":"relaciona"},{"alvo":"word","confianca":0.5,"epistemico":"fato","origem":"wiki","rel":"relaciona"}],"confianca":1.0,"contraexemplos":[],"descricao":"O estudo empírico do poder, do Estado e dos regimes — como se governa de fato, quem governa e por quê; a face aplicada da filosofia política.","epistemico":"fato","exemplos":[],"id":"ciencia_politica","memoria":"semantica","meta":{"dominio":"ciencias_de_fora","fonte":""},"origem":"wiki","palavras_chave":[],"sinonimos":[],"tipo":"conceito","titulo":"Ciência Política"}
\section{Ciência Política}
O estudo empírico do poder, do Estado e dos regimes — como se governa de fato, quem governa e por quê; a face aplicada da filosofia política.
\prov{relações: parte\_de politica; relaciona sociologia; relaciona word}
\prov{proveniência: wiki \textperiodcentered{} fato \textperiodcentered{} confiança 1.0 \textperiodcentered{} domínio ciencias\_de\_fora \textperiodcentered{} história 4 versões no jornal}

%CRISTAL {"arestas":[{"alvo":"isla","confianca":1.0,"epistemico":"fato","origem":"wiki","rel":"parte_de"},{"alvo":"tawhid","confianca":1.0,"epistemico":"fato","origem":"wiki","rel":"usa"}],"confianca":1.0,"contraexemplos":[],"descricao":"Islã: os cinco atos que estruturam a vida do muçulmano — shahada (o testemunho 'não há divindade senão Deus, e Muhammad é seu Mensageiro'), salat (oração), zakat (caridade), sawm (jejum do Ramadã) e hajj (peregrinação a Meca). A fé encarnada em disciplina.","epistemico":"inferencia","exemplos":[],"id":"cinco_pilares","memoria":"semantica","meta":{"autor":"Alcorão","dominio":"sagrado","fonte":""},"origem":"wiki","palavras_chave":[],"sinonimos":[],"tipo":"conceito","titulo":"Os Cinco Pilares do Islã"}
\section{Os Cinco Pilares do Islã}
Islã: os cinco atos que estruturam a vida do muçulmano — shahada (o testemunho 'não há divindade senão Deus, e Muhammad é seu Mensageiro'), salat (oração), zakat (caridade), sawm (jejum do Ramadã) e hajj (peregrinação a Meca). A fé encarnada em disciplina.
\prov{relações: parte\_de isla; usa tawhid}
\prov{proveniência: wiki \textperiodcentered{} inferencia \textperiodcentered{} confiança 1.0 \textperiodcentered{} domínio sagrado \textperiodcentered{} história 3 versões no jornal}

%CRISTAL {"arestas":[{"alvo":"clima_vs_tempo","confianca":1.0,"epistemico":"fato","origem":"wiki","rel":"explica"},{"alvo":"dissipacao_prigogine","confianca":1.0,"epistemico":"fato","origem":"wiki","rel":"relaciona"}],"confianca":1.0,"contraexemplos":[],"descricao":"O sistema de redistribuição de calor do equador aos polos por células atmosféricas (Hadley, Ferrel, polar) e correntes/circulação termohalina — motor das zonas climáticas e do transporte global de energia.","epistemico":"fato","exemplos":[],"id":"circulacao_atmosferica_oceanica","memoria":"semantica","meta":{"autor":"—","dominio":"ciencias_de_fora","fonte":""},"origem":"wiki","palavras_chave":[],"sinonimos":[],"tipo":"conceito","titulo":"Circulação Atmosférica e Oceânica"}
\section{Circulação Atmosférica e Oceânica}
O sistema de redistribuição de calor do equador aos polos por células atmosféricas (Hadley, Ferrel, polar) e correntes/circulação termohalina — motor das zonas climáticas e do transporte global de energia.
\prov{relações: explica clima\_vs\_tempo; relaciona dissipacao\_prigogine}
\prov{proveniência: wiki \textperiodcentered{} fato \textperiodcentered{} confiança 1.0 \textperiodcentered{} domínio ciencias\_de\_fora \textperiodcentered{} história 3 versões no jornal}

%CRISTAL {"arestas":[{"alvo":"artes_visuais","confianca":1.0,"epistemico":"fato","origem":"wiki","rel":"parte_de"}],"confianca":1.0,"contraexemplos":[],"descricao":"O jogo de luz e sombra que dá volume e drama (Caravaggio) — a terceira dimensão fingida com valor tonal.","epistemico":"fato","exemplos":[],"id":"claro_escuro","memoria":"semantica","meta":{"dominio":"ciencias_de_fora","fonte":""},"origem":"wiki","palavras_chave":[],"sinonimos":[],"tipo":"conceito","titulo":"Claro-escuro (Chiaroscuro)"}
\section{Claro-escuro (Chiaroscuro)}
O jogo de luz e sombra que dá volume e drama (Caravaggio) — a terceira dimensão fingida com valor tonal.
\prov{relações: parte\_de artes\_visuais}
\prov{proveniência: wiki \textperiodcentered{} fato \textperiodcentered{} confiança 1.0 \textperiodcentered{} domínio ciencias\_de\_fora \textperiodcentered{} história 2 versões no jornal}

%CRISTAL {"arestas":[{"alvo":"sociologia","confianca":1.0,"epistemico":"fato","origem":"wiki","rel":"parte_de"},{"alvo":"materialismo_historico","confianca":1.0,"epistemico":"fato","origem":"wiki","rel":"usa"},{"alvo":"word","confianca":0.5,"epistemico":"fato","origem":"wiki","rel":"relaciona"}],"confianca":1.0,"contraexemplos":[],"descricao":"Marx/Bourdieu: a divisão pela relação aos meios de produção; o habitus interioriza a desigualdade como 'natural'.","epistemico":"inferencia","exemplos":[],"id":"classe_social","memoria":"semantica","meta":{"autor":"Marx/Bourdieu","dominio":"ciencias_de_fora","fonte":""},"origem":"wiki","palavras_chave":[],"sinonimos":[],"tipo":"conceito","titulo":"Classe Social e Reprodução"}
\section{Classe Social e Reprodução}
Marx/Bourdieu: a divisão pela relação aos meios de produção; o habitus interioriza a desigualdade como 'natural'.
\prov{relações: parte\_de sociologia; usa materialismo\_historico; relaciona word}
\prov{proveniência: wiki \textperiodcentered{} inferencia \textperiodcentered{} confiança 1.0 \textperiodcentered{} domínio ciencias\_de\_fora \textperiodcentered{} história 4 versões no jornal}

%CRISTAL {"arestas":[{"alvo":"computabilidade","confianca":1.0,"epistemico":"fato","origem":"wiki","rel":"parte_de"},{"alvo":"maquina_de_turing","confianca":1.0,"epistemico":"fato","origem":"wiki","rel":"usa"},{"alvo":"prova_de_trabalho","confianca":1.0,"epistemico":"fato","origem":"wiki","rel":"explica"}],"confianca":1.0,"contraexemplos":[],"descricao":"Verificar uma solução é mais fácil que encontrá-la? O maior problema aberto da computação — no coração da prova-de-trabalho.","epistemico":"hipotese","exemplos":[],"id":"classes_p_np","memoria":"semantica","meta":{"dominio":"ciencias_de_fora","fonte":""},"origem":"wiki","palavras_chave":[],"sinonimos":[],"tipo":"conceito","titulo":"P versus NP"}
\section{P versus NP}
Verificar uma solução é mais fácil que encontrá-la? O maior problema aberto da computação — no coração da prova-de-trabalho.
\prov{relações: parte\_de computabilidade; usa maquina\_de\_turing; explica prova\_de\_trabalho}
\prov{proveniência: wiki \textperiodcentered{} hipotese \textperiodcentered{} confiança 1.0 \textperiodcentered{} domínio ciencias\_de\_fora \textperiodcentered{} história 8 versões no jornal}

%CRISTAL {"arestas":[{"alvo":"geografia_fisica","confianca":1.0,"epistemico":"fato","origem":"wiki","rel":"parte_de"}],"confianca":1.0,"contraexemplos":[],"descricao":"Tempo é o estado instantâneo/diário da atmosfera num lugar; clima é o padrão estatístico desse estado ao longo de décadas — a distinção entre a flutuação e a tendência do sistema atmosférico.","epistemico":"fato","exemplos":[],"id":"clima_vs_tempo","memoria":"semantica","meta":{"autor":"—","dominio":"ciencias_de_fora","fonte":""},"origem":"wiki","palavras_chave":[],"sinonimos":[],"tipo":"conceito","titulo":"Clima vs Tempo Meteorológico"}
\section{Clima vs Tempo Meteorológico}
Tempo é o estado instantâneo/diário da atmosfera num lugar; clima é o padrão estatístico desse estado ao longo de décadas — a distinção entre a flutuação e a tendência do sistema atmosférico.
\prov{relações: parte\_de geografia\_fisica}
\prov{proveniência: wiki \textperiodcentered{} fato \textperiodcentered{} confiança 1.0 \textperiodcentered{} domínio ciencias\_de\_fora \textperiodcentered{} história 2 versões no jornal}

%CRISTAL {"arestas":[{"alvo":"direito_e_tecnologia","confianca":1.0,"epistemico":"fato","origem":"wiki","rel":"parte_de"},{"alvo":"goldchain","confianca":1.0,"epistemico":"fato","origem":"wiki","rel":"explica"},{"alvo":"broca_os","confianca":1.0,"epistemico":"fato","origem":"wiki","rel":"explica"},{"alvo":"artefatos_tem_politica","confianca":1.0,"epistemico":"fato","origem":"wiki","rel":"relaciona"},{"alvo":"teoria_critica_tecnologia","confianca":1.0,"epistemico":"fato","origem":"wiki","rel":"relaciona"}],"confianca":1.0,"contraexemplos":[],"descricao":"Lessig: no ciberespaço, a arquitetura de software regula o comportamento tão eficazmente quanto a lei — o que o código permite ou proíbe torna-se a norma efetiva, sem juiz nem sanção estatal. A regra vive no protocolo.","epistemico":"inferencia","exemplos":[],"id":"code_is_law","memoria":"semantica","meta":{"autor":"Lessig","dominio":"ciencias_de_fora","fonte":""},"origem":"wiki","palavras_chave":[],"sinonimos":[],"tipo":"conceito","titulo":"Code is Law (Lessig)"}
\section{Code is Law (Lessig)}
Lessig: no ciberespaço, a arquitetura de software regula o comportamento tão eficazmente quanto a lei — o que o código permite ou proíbe torna-se a norma efetiva, sem juiz nem sanção estatal. A regra vive no protocolo.
\prov{relações: parte\_de direito\_e\_tecnologia; explica goldchain; explica broca\_os; relaciona artefatos\_tem\_politica; relaciona teoria\_critica\_tecnologia}
\prov{proveniência: wiki \textperiodcentered{} inferencia \textperiodcentered{} confiança 1.0 \textperiodcentered{} domínio ciencias\_de\_fora \textperiodcentered{} história 6 versões no jornal}

%CRISTAL {"arestas":[{"alvo":"estudos_de_midia","confianca":1.0,"epistemico":"fato","origem":"wiki","rel":"parte_de"},{"alvo":"mediacao_semiotica","confianca":1.0,"epistemico":"fato","origem":"wiki","rel":"relaciona"},{"alvo":"morte_do_autor","confianca":1.0,"epistemico":"fato","origem":"wiki","rel":"relaciona"}],"confianca":1.0,"contraexemplos":[],"descricao":"Stuart Hall: produzir uma mensagem (codificar) e interpretá-la (decodificar) são processos distintos que não coincidem necessariamente — o sentido do emissor pode ser reelaborado pela audiência conforme seu código cultural.","epistemico":"inferencia","exemplos":[],"id":"codificacao_decodificacao","memoria":"semantica","meta":{"autor":"Hall","dominio":"ciencias_de_fora","fonte":""},"origem":"wiki","palavras_chave":[],"sinonimos":[],"tipo":"conceito","titulo":"Codificação / Decodificação (Hall)"}
\section{Codificação / Decodificação (Hall)}
Stuart Hall: produzir uma mensagem (codificar) e interpretá-la (decodificar) são processos distintos que não coincidem necessariamente — o sentido do emissor pode ser reelaborado pela audiência conforme seu código cultural.
\prov{relações: parte\_de estudos\_de\_midia; relaciona mediacao\_semiotica; relaciona morte\_do\_autor}
\prov{proveniência: wiki \textperiodcentered{} inferencia \textperiodcentered{} confiança 1.0 \textperiodcentered{} domínio ciencias\_de\_fora \textperiodcentered{} história 4 versões no jornal}

%CRISTAL {"arestas":[{"alvo":"biologia","confianca":1.0,"epistemico":"fato","origem":"wiki","rel":"parte_de"},{"alvo":"dogma_central","confianca":1.0,"epistemico":"fato","origem":"wiki","rel":"parte_de"},{"alvo":"abiogenese","confianca":1.0,"epistemico":"fato","origem":"wiki","rel":"relaciona"},{"alvo":"cosmologia_quantica_hub","confianca":1.0,"epistemico":"fato","origem":"wiki","rel":"relaciona"},{"alvo":"dna_dupla_helice","confianca":1.0,"epistemico":"fato","origem":"wiki","rel":"usa"},{"alvo":"codigo_semiotico","confianca":1.0,"epistemico":"fato","origem":"wiki","rel":"explica"},{"alvo":"teoria_da_informacao","confianca":1.0,"epistemico":"fato","origem":"wiki","rel":"usa"},{"alvo":"dna_informacao_shannon","confianca":1.0,"epistemico":"fato","origem":"wiki","rel":"relaciona"}],"confianca":1.0,"contraexemplos":[],"descricao":"A tabela que traduz trincas de bases (códons) em aminoácidos — um CÓDIGO no sentido semiótico e informacional exato.","epistemico":"fato","exemplos":[],"id":"codigo_genetico","memoria":"semantica","meta":{"dominio":"ciencias_de_fora","fonte":""},"origem":"wiki","palavras_chave":[],"sinonimos":[],"tipo":"conceito","titulo":"Código Genético"}
\section{Código Genético}
A tabela que traduz trincas de bases (códons) em aminoácidos — um CÓDIGO no sentido semiótico e informacional exato.
\prov{relações: parte\_de biologia; parte\_de dogma\_central; relaciona abiogenese; relaciona cosmologia\_quantica\_hub; usa dna\_dupla\_helice; explica codigo\_semiotico; usa teoria\_da\_informacao; relaciona dna\_informacao\_shannon}
\prov{proveniência: wiki \textperiodcentered{} fato \textperiodcentered{} confiança 1.0 \textperiodcentered{} domínio ciencias\_de\_fora \textperiodcentered{} história 11 versões no jornal}

%CRISTAL {"arestas":[{"alvo":"celula_mineral_hub","confianca":1.0,"epistemico":"fato","origem":"wiki","rel":"parte_de"},{"alvo":"codigo_genetico","confianca":1.0,"epistemico":"fato","origem":"wiki","rel":"eh_um"},{"alvo":"beta_numeracao_metalica","confianca":1.0,"epistemico":"fato","origem":"wiki","rel":"eh_um"}],"confianca":1.0,"contraexemplos":[],"descricao":"As quatro bases A,C,G,T são os dígitos de uma numeração BASE 4, e um códon (3 bases) é um número de 3 dígitos: 4³=64 códons, mapeados a 20 aminoácidos + stop (o código redundante/degenerado). É a TRANSFORMADA de numeração da máquina — a mesma β-expansão dos metais, agora com base inteira 4: forward=expandir (os pesos 16,4,1), inverse=somar. Round-trip exato (ATG=14, AAA=0, TTT=63). O DNA é uma fita posicional na base quatro.","epistemico":"fato","exemplos":[],"id":"codigo_genetico_beta4","memoria":"semantica","meta":{"dominio":"biologia","fonte":"célula mineral"},"origem":"wiki","palavras_chave":[],"sinonimos":[],"tipo":"conceito","titulo":"O Código Genético é β-numeração base 4"}
\section{O Código Genético é \ensuremath{\beta}-numeração base 4}
As quatro bases A,C,G,T são os dígitos de uma numeração BASE 4, e um códon (3 bases) é um número de 3 dígitos: 4\ensuremath{^{3}}=64 códons, mapeados a 20 aminoácidos + stop (o código redundante/degenerado). É a TRANSFORMADA de numeração da máquina — a mesma \ensuremath{\beta}-expansão dos metais, agora com base inteira 4: forward=expandir (os pesos 16,4,1), inverse=somar. Round-trip exato (ATG=14, AAA=0, TTT=63). O DNA é uma fita posicional na base quatro.
\prov{relações: parte\_de celula\_mineral\_hub; eh\_um codigo\_genetico; eh\_um beta\_numeracao\_metalica}
\prov{proveniência: wiki \textperiodcentered{} célula mineral \textperiodcentered{} fato \textperiodcentered{} confiança 1.0 \textperiodcentered{} domínio biologia \textperiodcentered{} história 4 versões no jornal}

%CRISTAL {"arestas":[{"alvo":"codigo_genetico","confianca":1.0,"epistemico":"fato","origem":"wiki","rel":"explica"}],"confianca":1.0,"contraexemplos":[],"descricao":"A estrutura do código NÃO é arbitrária: códons semelhantes tendem a codificar aminoácidos de propriedades físico-químicas semelhantes, minimizando o dano de mutações pontuais (substituições conservativas). Freeland-Hurst 1998: o código-padrão é melhor que ~999.999 de 1.000.000 de códigos aleatórios sob custo realista. A robustez medida = fato; a CAUSA (seleção adaptativa vs origem neutra/histórica) = aberto.","epistemico":"fato","exemplos":[],"id":"codigo_robustez_erro","memoria":"semantica","meta":{"autor":"broca-so","dominio":"biologia","fonte":"Freeland-Hurst 1998 JME 47,238"},"origem":"wiki","palavras_chave":[],"sinonimos":[],"tipo":"conceito","titulo":"A robustez a erros do código (one in a million)"}
\section{A robustez a erros do código (one in a million)}
A estrutura do código NÃO é arbitrária: códons semelhantes tendem a codificar aminoácidos de propriedades físico-químicas semelhantes, minimizando o dano de mutações pontuais (substituições conservativas). Freeland-Hurst 1998: o código-padrão é melhor que \~{}999.999 de 1.000.000 de códigos aleatórios sob custo realista. A robustez medida = fato; a CAUSA (seleção adaptativa vs origem neutra/histórica) = aberto.
\prov{relações: explica codigo\_genetico}
\prov{proveniência: wiki \textperiodcentered{} Freeland-Hurst 1998 JME 47,238 \textperiodcentered{} fato \textperiodcentered{} confiança 1.0 \textperiodcentered{} domínio biologia \textperiodcentered{} história 2 versões no jornal}

%CRISTAL {"arestas":[{"alvo":"semiotica","confianca":1.0,"epistemico":"fato","origem":"wiki","rel":"parte_de"},{"alvo":"cifra","confianca":1.0,"epistemico":"fato","origem":"wiki","rel":"relaciona"}],"confianca":1.0,"contraexemplos":[],"descricao":"O sistema de convenções que liga significantes a significados e torna a comunicação possível — cifrar e decifrar.","epistemico":"fato","exemplos":[],"id":"codigo_semiotico","memoria":"semantica","meta":{"dominio":"ciencias_de_fora","fonte":""},"origem":"wiki","palavras_chave":[],"sinonimos":[],"tipo":"conceito","titulo":"Código"}
\section{Código}
O sistema de convenções que liga significantes a significados e torna a comunicação possível — cifrar e decifrar.
\prov{relações: parte\_de semiotica; relaciona cifra}
\prov{proveniência: wiki \textperiodcentered{} fato \textperiodcentered{} confiança 1.0 \textperiodcentered{} domínio ciencias\_de\_fora \textperiodcentered{} história 3 versões no jornal}

%CRISTAL {"arestas":[{"alvo":"paradigma_dispositivo","confianca":1.0,"epistemico":"fato","origem":"wiki","rel":"contradiz"},{"alvo":"aura_benjamin","confianca":1.0,"epistemico":"fato","origem":"wiki","rel":"relaciona"}],"confianca":1.0,"contraexemplos":[],"descricao":"Borgmann: coisas que reúnem sentido, corpo e comunidade em torno de uma prática engajada (cozinhar, tocar um instrumento) — o contraponto positivo à comodidade dispersante dos dispositivos.","epistemico":"inferencia","exemplos":[],"id":"coisas_focais","memoria":"semantica","meta":{"autor":"Borgmann","dominio":"ciencias_de_fora","fonte":""},"origem":"wiki","palavras_chave":[],"sinonimos":[],"tipo":"conceito","titulo":"Coisas e Práticas Focais (Borgmann)"}
\section{Coisas e Práticas Focais (Borgmann)}
Borgmann: coisas que reúnem sentido, corpo e comunidade em torno de uma prática engajada (cozinhar, tocar um instrumento) — o contraponto positivo à comodidade dispersante dos dispositivos.
\prov{relações: contradiz paradigma\_dispositivo; relaciona aura\_benjamin}
\prov{proveniência: wiki \textperiodcentered{} inferencia \textperiodcentered{} confiança 1.0 \textperiodcentered{} domínio ciencias\_de\_fora \textperiodcentered{} história 3 versões no jornal}

%CRISTAL {"arestas":[{"alvo":"colapso","confianca":1.0,"epistemico":"fato","origem":"wiki","rel":"explica"},{"alvo":"mecanica_quantica","confianca":1.0,"epistemico":"fato","origem":"wiki","rel":"parte_de"}],"confianca":1.0,"contraexemplos":[],"descricao":"A superposição colapsa num estado definido no ato da medição. O que conta como 'observador'? Copenhague deixa em aberto; toca o coração do broca-so.","epistemico":"inferencia","exemplos":[],"id":"colapso_observador","memoria":"semantica","meta":{"dominio":"filosofia","fonte":"web: QM"},"origem":"wiki","palavras_chave":[],"sinonimos":[],"tipo":"conceito","titulo":"O Observador e o Colapso"}
\section{O Observador e o Colapso}
A superposição colapsa num estado definido no ato da medição. O que conta como 'observador'? Copenhague deixa em aberto; toca o coração do broca-so.
\prov{relações: explica colapso; parte\_de mecanica\_quantica}
\prov{proveniência: wiki \textperiodcentered{} web: QM \textperiodcentered{} inferencia \textperiodcentered{} confiança 1.0 \textperiodcentered{} domínio filosofia \textperiodcentered{} história 3 versões no jornal}

%CRISTAL {"arestas":[{"alvo":"teoria_da_informacao","confianca":1.0,"epistemico":"fato","origem":"wiki","rel":"relaciona"},{"alvo":"computabilidade","confianca":1.0,"epistemico":"fato","origem":"wiki","rel":"usa"},{"alvo":"floco_no_metal","confianca":1.0,"epistemico":"fato","origem":"wiki","rel":"relaciona"},{"alvo":"kolmogorov_complexidade","confianca":1.0,"epistemico":"fato","origem":"wiki","rel":"relaciona"}],"confianca":1.0,"contraexemplos":[],"descricao":"O tamanho do menor programa que gera um objeto — a medida absoluta de sua informação e aleatoriedade. Incomputável.","epistemico":"fato","exemplos":[],"id":"complexidade_de_kolmogorov","memoria":"semantica","meta":{"autor":"Kolmogorov","dominio":"ciencias_de_fora","fonte":""},"origem":"wiki","palavras_chave":[],"sinonimos":[],"tipo":"conceito","titulo":"Complexidade de Kolmogorov"}
\section{Complexidade de Kolmogorov}
O tamanho do menor programa que gera um objeto — a medida absoluta de sua informação e aleatoriedade. Incomputável.
\prov{relações: relaciona teoria\_da\_informacao; usa computabilidade; relaciona floco\_no\_metal; relaciona kolmogorov\_complexidade}
\prov{proveniência: wiki \textperiodcentered{} fato \textperiodcentered{} confiança 1.0 \textperiodcentered{} domínio ciencias\_de\_fora \textperiodcentered{} história 10 versões no jornal}

%CRISTAL {"arestas":[{"alvo":"artes_visuais","confianca":1.0,"epistemico":"fato","origem":"wiki","rel":"parte_de"}],"confianca":1.0,"contraexemplos":[],"descricao":"A organização dos elementos no quadro — equilíbrio, ritmo, ênfase; onde o olho entra, caminha e descansa.","epistemico":"fato","exemplos":[],"id":"composicao_visual","memoria":"semantica","meta":{"dominio":"ciencias_de_fora","fonte":""},"origem":"wiki","palavras_chave":[],"sinonimos":[],"tipo":"conceito","titulo":"Composição"}
\section{Composição}
A organização dos elementos no quadro — equilíbrio, ritmo, ênfase; onde o olho entra, caminha e descansa.
\prov{relações: parte\_de artes\_visuais}
\prov{proveniência: wiki \textperiodcentered{} fato \textperiodcentered{} confiança 1.0 \textperiodcentered{} domínio ciencias\_de\_fora \textperiodcentered{} história 2 versões no jornal}

%CRISTAL {"arestas":[{"alvo":"producao_capitalista_espaco_harvey","confianca":1.0,"epistemico":"fato","origem":"wiki","rel":"usa"},{"alvo":"sociedade_em_rede","confianca":1.0,"epistemico":"fato","origem":"wiki","rel":"relaciona"}],"confianca":1.0,"contraexemplos":[],"descricao":"Harvey: a aceleração do giro do capital e das comunicações encolhe as distâncias relativas e 'aniquila o espaço pelo tempo', alterando a experiência de lugar e de tempo — traço central da condição pós-moderna.","epistemico":"inferencia","exemplos":[],"id":"compressao_espaco_tempo","memoria":"semantica","meta":{"autor":"Harvey","dominio":"ciencias_de_fora","fonte":""},"origem":"wiki","palavras_chave":[],"sinonimos":[],"tipo":"conceito","titulo":"Compressão Espaço-Tempo (Harvey)"}
\section{Compressão Espaço-Tempo (Harvey)}
Harvey: a aceleração do giro do capital e das comunicações encolhe as distâncias relativas e 'aniquila o espaço pelo tempo', alterando a experiência de lugar e de tempo — traço central da condição pós-moderna.
\prov{relações: usa producao\_capitalista\_espaco\_harvey; relaciona sociedade\_em\_rede}
\prov{proveniência: wiki \textperiodcentered{} inferencia \textperiodcentered{} confiança 1.0 \textperiodcentered{} domínio ciencias\_de\_fora \textperiodcentered{} história 3 versões no jornal}

%CRISTAL {"arestas":[{"alvo":"musculo_mineral_hub","confianca":1.0,"epistemico":"fato","origem":"wiki","rel":"parte_de"},{"alvo":"parabola_isa","confianca":1.0,"epistemico":"fato","origem":"wiki","rel":"eh_um"},{"alvo":"numeros_de_pisot","confianca":1.0,"epistemico":"fato","origem":"wiki","rel":"relaciona"}],"confianca":1.0,"contraexemplos":[],"descricao":"A força ativa que o sarcômero gera depende do comprimento: máxima no ÓTIMO (a sobreposição ideal dos filamentos actina-miosina), caindo dos dois lados — uma PARÁBOLA invertida com o vértice no comprimento ótimo. É a parábola da máquina, literal e biológica: o ótimo é o pico (o cristal), e esticar ou encolher demais tira a força (as bordas caem). Verif.: o vértice está no L_opt, a força cai simétrica dos dois lados.","epistemico":"fato","exemplos":[],"id":"comprimento_tensao_e_a_parabola","memoria":"semantica","meta":{"dominio":"biologia","fonte":"músculo mineral"},"origem":"wiki","palavras_chave":[],"sinonimos":[],"tipo":"conceito","titulo":"A Curva Comprimento-Tensão é uma Parábola"}
\section{A Curva Comprimento-Tensão é uma Parábola}
A força ativa que o sarcômero gera depende do comprimento: máxima no ÓTIMO (a sobreposição ideal dos filamentos actina-miosina), caindo dos dois lados — uma PARÁBOLA invertida com o vértice no comprimento ótimo. É a parábola da máquina, literal e biológica: o ótimo é o pico (o cristal), e esticar ou encolher demais tira a força (as bordas caem). Verif.: o vértice está no L\_opt, a força cai simétrica dos dois lados.
\prov{relações: parte\_de musculo\_mineral\_hub; eh\_um parabola\_isa; relaciona numeros\_de\_pisot}
\prov{proveniência: wiki \textperiodcentered{} músculo mineral \textperiodcentered{} fato \textperiodcentered{} confiança 1.0 \textperiodcentered{} domínio biologia \textperiodcentered{} história 4 versões no jornal}

%CRISTAL {"arestas":[{"alvo":"algoritmo","confianca":1.0,"epistemico":"fato","origem":"curriculo","rel":"especializa"},{"alvo":"logica","confianca":1.0,"epistemico":"fato","origem":"curriculo","rel":"depende"},{"alvo":"matematica","confianca":1.0,"epistemico":"fato","origem":"wiki","rel":"parte_de"}],"confianca":1.0,"contraexemplos":[],"descricao":"O estudo do que pode ser calculado por procedimento mecânico — os limites do algoritmo.","epistemico":"fato","exemplos":[],"id":"computabilidade","memoria":"semantica","meta":{"dominio":"ciencias_de_fora","fonte":""},"origem":"wiki","palavras_chave":["halting","church-turing"],"sinonimos":["computabilidade","decidibilidade"],"tipo":"conceito","titulo":"Computabilidade"}
\section{Computabilidade}
O estudo do que pode ser calculado por procedimento mecânico — os limites do algoritmo.
\prov{sinónimos: computabilidade; decidibilidade}
\prov{relações: especializa algoritmo; depende logica; parte\_de matematica}
\prov{proveniência: wiki \textperiodcentered{} fato \textperiodcentered{} confiança 1.0 \textperiodcentered{} domínio ciencias\_de\_fora \textperiodcentered{} história 43 versões no jornal}

%CRISTAL {"arestas":[{"alvo":"sociologia","confianca":1.0,"epistemico":"fato","origem":"wiki","rel":"relaciona"},{"alvo":"linguagem","confianca":1.0,"epistemico":"fato","origem":"wiki","rel":"relaciona"},{"alvo":"filosofia_da_tecnica","confianca":1.0,"epistemico":"fato","origem":"wiki","rel":"relaciona"}],"confianca":1.0,"contraexemplos":[],"descricao":"O campo que estuda a produção, circulação e interpretação de mensagens entre pessoas, grupos e máquinas — da conversa ao sistema de mídia.","epistemico":"fato","exemplos":[],"id":"comunicacao","memoria":"semantica","meta":{"dominio":"ciencias_de_fora","fonte":""},"origem":"wiki","palavras_chave":[],"sinonimos":[],"tipo":"conceito","titulo":"Comunicação"}
\section{Comunicação}
O campo que estuda a produção, circulação e interpretação de mensagens entre pessoas, grupos e máquinas — da conversa ao sistema de mídia.
\prov{relações: relaciona sociologia; relaciona linguagem; relaciona filosofia\_da\_tecnica}
\prov{proveniência: wiki \textperiodcentered{} fato \textperiodcentered{} confiança 1.0 \textperiodcentered{} domínio ciencias\_de\_fora \textperiodcentered{} história 4 versões no jornal}

%CRISTAL {"arestas":[{"alvo":"comunicacao","confianca":1.0,"epistemico":"fato","origem":"wiki","rel":"parte_de"},{"alvo":"word","confianca":0.5,"epistemico":"fato","origem":"wiki","rel":"relaciona"},{"alvo":"virtualizacao","confianca":0.5,"epistemico":"fato","origem":"wiki","rel":"relaciona"}],"confianca":1.0,"contraexemplos":[],"descricao":"A comunicação na era da internet e das redes — plataformas, algoritmos, dados e a reconfiguração do público, da atenção e da verdade.","epistemico":"fato","exemplos":[],"id":"comunicacao_digital","memoria":"semantica","meta":{"dominio":"ciencias_de_fora","fonte":""},"origem":"wiki","palavras_chave":[],"sinonimos":[],"tipo":"conceito","titulo":"Comunicação Digital"}
\section{Comunicação Digital}
A comunicação na era da internet e das redes — plataformas, algoritmos, dados e a reconfiguração do público, da atenção e da verdade.
\prov{relações: parte\_de comunicacao; relaciona word; relaciona virtualizacao}
\prov{proveniência: wiki \textperiodcentered{} fato \textperiodcentered{} confiança 1.0 \textperiodcentered{} domínio ciencias\_de\_fora \textperiodcentered{} história 4 versões no jornal}

%CRISTAL {"arestas":[{"alvo":"estudos_de_midia","confianca":1.0,"epistemico":"fato","origem":"wiki","rel":"parte_de"},{"alvo":"honestidade_intelectual","confianca":1.0,"epistemico":"fato","origem":"wiki","rel":"relaciona"},{"alvo":"idiotismo_nao_examinar","confianca":1.0,"epistemico":"fato","origem":"wiki","rel":"relaciona"},{"alvo":"tiffany","confianca":1.0,"epistemico":"fato","origem":"wiki","rel":"relaciona"}],"confianca":1.0,"contraexemplos":[],"descricao":"Lopes assina como 'o iconoclasta' e escolhe conscientemente um estilo 'incisivo, provocativo', 'deliberadamente popular': 'escrevo para provocar mesmo'. Seus textos são manifesto/panfleto dirigido não a especialistas cúmplices do dogma, mas a mentes inquietas.","epistemico":"inferencia","exemplos":[],"id":"comunicacao_iconoclasta_gentil","memoria":"semantica","meta":{"autor":"Gentil Lopes","dominio":"ciencias_de_fora","fonte":""},"origem":"wiki","palavras_chave":[],"sinonimos":[],"tipo":"conceito","titulo":"A Comunicação Iconoclasta (Lopes)"}
\section{A Comunicação Iconoclasta (Lopes)}
Lopes assina como 'o iconoclasta' e escolhe conscientemente um estilo 'incisivo, provocativo', 'deliberadamente popular': 'escrevo para provocar mesmo'. Seus textos são manifesto/panfleto dirigido não a especialistas cúmplices do dogma, mas a mentes inquietas.
\prov{relações: parte\_de estudos\_de\_midia; relaciona honestidade\_intelectual; relaciona idiotismo\_nao\_examinar; relaciona tiffany}
\prov{proveniência: wiki \textperiodcentered{} inferencia \textperiodcentered{} confiança 1.0 \textperiodcentered{} domínio ciencias\_de\_fora \textperiodcentered{} história 6 versões no jornal}

%CRISTAL {"arestas":[{"alvo":"relacoes_internacionais","confianca":1.0,"epistemico":"fato","origem":"wiki","rel":"relaciona"},{"alvo":"soberania_vestfaliana","confianca":1.0,"epistemico":"fato","origem":"wiki","rel":"relaciona"},{"alvo":"religiao","confianca":1.0,"epistemico":"fato","origem":"wiki","rel":"relaciona"},{"alvo":"construtivismo_wendt","confianca":1.0,"epistemico":"fato","origem":"wiki","rel":"relaciona"}],"confianca":1.0,"contraexemplos":[],"descricao":"Anderson: a nação é uma comunidade política imaginada (seus membros jamais se conhecerão), concebida como limitada e soberana — produto moderno do 'capitalismo de imprensa', não uma essência ancestral, que substitui os laços religiosos por pertencimento nacional.","epistemico":"inferencia","exemplos":[],"id":"comunidades_imaginadas_anderson","memoria":"semantica","meta":{"autor":"Anderson","dominio":"ciencias_de_fora","fonte":""},"origem":"wiki","palavras_chave":[],"sinonimos":[],"tipo":"conceito","titulo":"A Nação como Comunidade Imaginada (Anderson)"}
\section{A Nação como Comunidade Imaginada (Anderson)}
Anderson: a nação é uma comunidade política imaginada (seus membros jamais se conhecerão), concebida como limitada e soberana — produto moderno do 'capitalismo de imprensa', não uma essência ancestral, que substitui os laços religiosos por pertencimento nacional.
\prov{relações: relaciona relacoes\_internacionais; relaciona soberania\_vestfaliana; relaciona religiao; relaciona construtivismo\_wendt}
\prov{proveniência: wiki \textperiodcentered{} inferencia \textperiodcentered{} confiança 1.0 \textperiodcentered{} domínio ciencias\_de\_fora \textperiodcentered{} história 5 versões no jornal}

%CRISTAL {"arestas":[{"alvo":"sintaxe","confianca":1.0,"epistemico":"fato","origem":"wiki","rel":"parte_de"}],"confianca":1.0,"contraexemplos":[],"descricao":"Palavras que ajustam sua forma umas às outras (gênero, número, pessoa) — 'as casas brancas'. Redundância que amarra a frase; some no inglês, explode no bantu.","epistemico":"fato","exemplos":[],"id":"concordancia","memoria":"semantica","meta":{"dominio":"ciencias_de_fora","fonte":""},"origem":"wiki","palavras_chave":[],"sinonimos":[],"tipo":"conceito","titulo":"Concordância"}
\section{Concordância}
Palavras que ajustam sua forma umas às outras (gênero, número, pessoa) — 'as casas brancas'. Redundância que amarra a frase; some no inglês, explode no bantu.
\prov{relações: parte\_de sintaxe}
\prov{proveniência: wiki \textperiodcentered{} fato \textperiodcentered{} confiança 1.0 \textperiodcentered{} domínio ciencias\_de\_fora \textperiodcentered{} história 2 versões no jornal}

%CRISTAL {"arestas":[{"alvo":"psicologia","confianca":1.0,"epistemico":"fato","origem":"wiki","rel":"parte_de"},{"alvo":"aparelho_psiquico","confianca":1.0,"epistemico":"fato","origem":"wiki","rel":"contradiz"},{"alvo":"word","confianca":0.5,"epistemico":"fato","origem":"wiki","rel":"relaciona"}],"confianca":1.0,"contraexemplos":[],"descricao":"Skinner: o comportamento é moldado por consequências (reforço/punição); o behaviorismo nega o inconsciente e a liberdade.","epistemico":"inferencia","exemplos":[],"id":"condicionamento_operante","memoria":"semantica","meta":{"autor":"Skinner","dominio":"ciencias_de_fora","fonte":""},"origem":"wiki","palavras_chave":[],"sinonimos":[],"tipo":"conceito","titulo":"Condicionamento Operante"}
\section{Condicionamento Operante}
Skinner: o comportamento é moldado por consequências (reforço/punição); o behaviorismo nega o inconsciente e a liberdade.
\prov{relações: parte\_de psicologia; contradiz aparelho\_psiquico; relaciona word}
\prov{proveniência: wiki \textperiodcentered{} inferencia \textperiodcentered{} confiança 1.0 \textperiodcentered{} domínio ciencias\_de\_fora \textperiodcentered{} história 4 versões no jornal}

%CRISTAL {"arestas":[{"alvo":"teorias_aprendizagem","confianca":1.0,"epistemico":"fato","origem":"wiki","rel":"parte_de"},{"alvo":"condicionamento_operante","confianca":1.0,"epistemico":"fato","origem":"wiki","rel":"relaciona"},{"alvo":"hopfield","confianca":1.0,"epistemico":"fato","origem":"wiki","rel":"referencia"}],"confianca":1.0,"contraexemplos":[],"descricao":"Thorndike: aprender é formar conexões estímulo-resposta; pela Lei do Efeito, respostas seguidas de satisfação fortalecem-se e as seguidas de desconforto enfraquecem-se. Precursor do behaviorismo.","epistemico":"fato","exemplos":[],"id":"conexionismo_thorndike","memoria":"semantica","meta":{"autor":"Thorndike","dominio":"ciencias_de_fora","fonte":""},"origem":"wiki","palavras_chave":[],"sinonimos":[],"tipo":"conceito","titulo":"Conexionismo e Lei do Efeito (Thorndike)"}
\section{Conexionismo e Lei do Efeito (Thorndike)}
Thorndike: aprender é formar conexões estímulo-resposta; pela Lei do Efeito, respostas seguidas de satisfação fortalecem-se e as seguidas de desconforto enfraquecem-se. Precursor do behaviorismo.
\prov{relações: parte\_de teorias\_aprendizagem; relaciona condicionamento\_operante; referencia hopfield}
\prov{proveniência: wiki \textperiodcentered{} fato \textperiodcentered{} confiança 1.0 \textperiodcentered{} domínio ciencias\_de\_fora \textperiodcentered{} história 4 versões no jornal}

%CRISTAL {"arestas":[{"alvo":"sistema_e_lei_nao_dono","confianca":1.0,"epistemico":"fato","origem":"wiki","rel":"usa"},{"alvo":"code_is_law","confianca":1.0,"epistemico":"fato","origem":"wiki","rel":"contradiz"},{"alvo":"honestidade_intelectual","confianca":1.0,"epistemico":"fato","origem":"wiki","rel":"relaciona"},{"alvo":"htlc","confianca":1.0,"epistemico":"fato","origem":"wiki","rel":"relaciona"}],"confianca":1.0,"contraexemplos":[],"descricao":"A ressalva honesta do design contra o absolutismo ingênuo do 'code is law': quando a condição não é cripto-verificável, exige-se um árbitro (multisig) que não é trustless (extorsão, conluio). Vende-se como 'confiança minimizada', nunca 'sem confiança' — a honestidade intelectual aplicada ao próprio produto.","epistemico":"inferencia","exemplos":[],"id":"confianca_minimizada","memoria":"semantica","meta":{"autor":"broca-so (design)","dominio":"ciencias_de_fora","fonte":""},"origem":"wiki","palavras_chave":[],"sinonimos":[],"tipo":"conceito","titulo":"Confiança Minimizada, não 'Sem Confiança' (design broca-so)"}
\section{Confiança Minimizada, não 'Sem Confiança' (design broca-so)}
A ressalva honesta do design contra o absolutismo ingênuo do 'code is law': quando a condição não é cripto-verificável, exige-se um árbitro (multisig) que não é trustless (extorsão, conluio). Vende-se como 'confiança minimizada', nunca 'sem confiança' — a honestidade intelectual aplicada ao próprio produto.
\prov{relações: usa sistema\_e\_lei\_nao\_dono; contradiz code\_is\_law; relaciona honestidade\_intelectual; relaciona htlc}
\prov{proveniência: wiki \textperiodcentered{} inferencia \textperiodcentered{} confiança 1.0 \textperiodcentered{} domínio ciencias\_de\_fora \textperiodcentered{} história 5 versões no jornal}

%CRISTAL {"arestas":[{"alvo":"critica_fenomenologica_dreyfus","confianca":1.0,"epistemico":"fato","origem":"wiki","rel":"usa"},{"alvo":"olho_de_koch","confianca":1.0,"epistemico":"fato","origem":"wiki","rel":"relaciona"},{"alvo":"tiffany","confianca":1.0,"epistemico":"fato","origem":"wiki","rel":"relaciona"}],"confianca":1.0,"contraexemplos":[],"descricao":"Dreyfus (via Merleau-Ponty/Polanyi): um saber-fazer não-proposicional, inseparável do corpo e da situação, que precede e sustenta o conhecimento explícito. Ecoa a percepção que 'reconstrói, não copia'.","epistemico":"inferencia","exemplos":[],"id":"conhecimento_tacito_corporificado","memoria":"semantica","meta":{"autor":"Dreyfus/Polanyi","dominio":"ciencias_de_fora","fonte":""},"origem":"wiki","palavras_chave":[],"sinonimos":[],"tipo":"conceito","titulo":"Conhecimento Tácito / Corporificado"}
\section{Conhecimento Tácito / Corporificado}
Dreyfus (via Merleau-Ponty/Polanyi): um saber-fazer não-proposicional, inseparável do corpo e da situação, que precede e sustenta o conhecimento explícito. Ecoa a percepção que 'reconstrói, não copia'.
\prov{relações: usa critica\_fenomenologica\_dreyfus; relaciona olho\_de\_koch; relaciona tiffany}
\prov{proveniência: wiki \textperiodcentered{} inferencia \textperiodcentered{} confiança 1.0 \textperiodcentered{} domínio ciencias\_de\_fora \textperiodcentered{} história 4 versões no jornal}

%CRISTAL {"arestas":[{"alvo":"topologia","confianca":1.0,"epistemico":"fato","origem":"wiki","rel":"parte_de"},{"alvo":"geometria_riemanniana","confianca":1.0,"epistemico":"fato","origem":"wiki","rel":"usa"}],"confianca":1.0,"contraexemplos":[],"descricao":"Toda 3-variedade fechada e simplesmente conexa é a esfera. Provada por Perelman (2003) via fluxo de Ricci e entropia monótona.","epistemico":"fato","exemplos":[],"id":"conjectura_de_poincare","memoria":"semantica","meta":{"autor":"Perelman","dominio":"ciencias_de_fora","fonte":""},"origem":"wiki","palavras_chave":[],"sinonimos":[],"tipo":"conceito","titulo":"Conjectura de Poincaré"}
\section{Conjectura de Poincaré}
Toda 3-variedade fechada e simplesmente conexa é a esfera. Provada por Perelman (2003) via fluxo de Ricci e entropia monótona.
\prov{relações: parte\_de topologia; usa geometria\_riemanniana}
\prov{proveniência: wiki \textperiodcentered{} fato \textperiodcentered{} confiança 1.0 \textperiodcentered{} domínio ciencias\_de\_fora \textperiodcentered{} história 6 versões no jornal}

%CRISTAL {"arestas":[{"alvo":"filosofia_da_mente","confianca":1.0,"epistemico":"fato","origem":"wiki","rel":"parte_de"},{"alvo":"vida-morte","confianca":0.5,"epistemico":"fato","origem":"wiki","rel":"relaciona"},{"alvo":"word","confianca":0.5,"epistemico":"fato","origem":"wiki","rel":"relaciona"},{"alvo":"virtualizacao","confianca":0.5,"epistemico":"fato","origem":"wiki","rel":"relaciona"},{"alvo":"pell","confianca":0.5,"epistemico":"fato","origem":"wiki","rel":"relaciona"},{"alvo":"readme","confianca":0.5,"epistemico":"fato","origem":"wiki","rel":"relaciona"},{"alvo":"fibonacci_multidimensional","confianca":0.5,"epistemico":"fato","origem":"wiki","rel":"relaciona"},{"alvo":"ethics_of_care","confianca":0.5,"epistemico":"fato","origem":"wiki","rel":"relaciona"},{"alvo":"dois_pacotes_pow_lucro","confianca":0.5,"epistemico":"fato","origem":"wiki","rel":"relaciona"},{"alvo":"koch_pell_verificacao","confianca":0.5,"epistemico":"fato","origem":"wiki","rel":"relaciona"},{"alvo":"metafisica","confianca":0.5,"epistemico":"fato","origem":"wiki","rel":"relaciona"},{"alvo":"heap","confianca":0.5,"epistemico":"fato","origem":"wiki","rel":"relaciona"},{"alvo":"painel_estelar","confianca":0.5,"epistemico":"fato","origem":"wiki","rel":"relaciona"},{"alvo":"incautos_enganados","confianca":0.5,"epistemico":"fato","origem":"wiki","rel":"relaciona"},{"alvo":"pipeline_de_ingestao","confianca":0.5,"epistemico":"fato","origem":"wiki","rel":"relaciona"},{"alvo":"telemetria_shadow_producao","confianca":0.5,"epistemico":"fato","origem":"wiki","rel":"relaciona"}],"confianca":1.0,"contraexemplos":[],"descricao":"A experiência subjetiva — o 'algo que é ser' um sistema; o problema central da filosofia da mente.","epistemico":"fato","exemplos":[],"id":"consciencia","memoria":"semantica","meta":{"dominio":"filosofia","fonte":""},"origem":"wiki","palavras_chave":[],"sinonimos":[],"tipo":"conceito","titulo":"Consciência"}
\section{Consciência}
A experiência subjetiva — o 'algo que é ser' um sistema; o problema central da filosofia da mente.
\prov{relações: parte\_de filosofia\_da\_mente; relaciona vida-morte; relaciona word; relaciona virtualizacao; relaciona pell; relaciona readme; relaciona fibonacci\_multidimensional; relaciona ethics\_of\_care; relaciona dois\_pacotes\_pow\_lucro; relaciona koch\_pell\_verificacao; relaciona metafisica; relaciona heap; relaciona painel\_estelar; relaciona incautos\_enganados; relaciona pipeline\_de\_ingestao; relaciona telemetria\_shadow\_producao}
\prov{proveniência: wiki \textperiodcentered{} fato \textperiodcentered{} confiança 1.0 \textperiodcentered{} domínio filosofia \textperiodcentered{} história 17 versões no jornal}

%CRISTAL {"arestas":[{"alvo":"consciencia","confianca":1.0,"epistemico":"fato","origem":"wiki","rel":"parte_de"},{"alvo":"contemplacao","confianca":1.0,"epistemico":"fato","origem":"wiki","rel":"usa"},{"alvo":"dzogchen_rigpa","confianca":1.0,"epistemico":"fato","origem":"wiki","rel":"usa"}],"confianca":1.0,"contraexemplos":[],"descricao":"Lopes: nascemos 'tela em branco' (Consciência Pura, nem matéria nem energia); sobre ela pintam-se identidade e ego. Hub que integra Vedanta, Plotino, Dzogchen.","epistemico":"hipotese","exemplos":[],"id":"consciencia_pura","memoria":"semantica","meta":{"autor":"Gentil Lopes","dominio":"filosofia","fonte":""},"origem":"livro","palavras_chave":[],"sinonimos":[],"tipo":"conceito","titulo":"Consciência Pura / Tela em Branco"}
\section{Consciência Pura / Tela em Branco}
Lopes: nascemos 'tela em branco' (Consciência Pura, nem matéria nem energia); sobre ela pintam-se identidade e ego. Hub que integra Vedanta, Plotino, Dzogchen.
\prov{relações: parte\_de consciencia; usa contemplacao; usa dzogchen\_rigpa}
\prov{proveniência: livro \textperiodcentered{} hipotese \textperiodcentered{} confiança 1.0 \textperiodcentered{} domínio filosofia \textperiodcentered{} história 5 versões no jornal}

%CRISTAL {"arestas":[{"alvo":"educacao_problematizadora","confianca":1.0,"epistemico":"fato","origem":"wiki","rel":"usa"},{"alvo":"biopoder","confianca":1.0,"epistemico":"fato","origem":"wiki","rel":"relaciona"},{"alvo":"injustica_epistemica","confianca":1.0,"epistemico":"fato","origem":"wiki","rel":"relaciona"}],"confianca":1.0,"contraexemplos":[],"descricao":"Freire: o processo pelo qual o sujeito passa de uma consciência ingênua a uma consciência crítica das condições que o oprimem, tornando-se capaz de agir sobre elas (práxis = reflexão + ação).","epistemico":"inferencia","exemplos":[],"id":"conscientizacao","memoria":"semantica","meta":{"autor":"Freire","dominio":"ciencias_de_fora","fonte":""},"origem":"wiki","palavras_chave":[],"sinonimos":[],"tipo":"conceito","titulo":"Conscientização"}
\section{Conscientização}
Freire: o processo pelo qual o sujeito passa de uma consciência ingênua a uma consciência crítica das condições que o oprimem, tornando-se capaz de agir sobre elas (práxis = reflexão + ação).
\prov{relações: usa educacao\_problematizadora; relaciona biopoder; relaciona injustica\_epistemica}
\prov{proveniência: wiki \textperiodcentered{} inferencia \textperiodcentered{} confiança 1.0 \textperiodcentered{} domínio ciencias\_de\_fora \textperiodcentered{} história 4 versões no jornal}

%CRISTAL {"arestas":[{"alvo":"autonomia_do_paciente","confianca":1.0,"epistemico":"fato","origem":"wiki","rel":"usa"},{"alvo":"etica_pesquisa_humanos","confianca":1.0,"epistemico":"fato","origem":"wiki","rel":"parte_de"}],"confianca":1.0,"contraexemplos":[],"descricao":"A exigência de que o paciente ou sujeito de pesquisa, informado de riscos, benefícios e alternativas, autorize voluntariamente e sem coerção — a operacionalização da autonomia e o núcleo do Código de Nuremberg.","epistemico":"inferencia","exemplos":[],"id":"consentimento_informado","memoria":"semantica","meta":{"autor":"Faden & Beauchamp","dominio":"ciencias_de_fora","fonte":""},"origem":"wiki","palavras_chave":[],"sinonimos":[],"tipo":"conceito","titulo":"Consentimento Informado"}
\section{Consentimento Informado}
A exigência de que o paciente ou sujeito de pesquisa, informado de riscos, benefícios e alternativas, autorize voluntariamente e sem coerção — a operacionalização da autonomia e o núcleo do Código de Nuremberg.
\prov{relações: usa autonomia\_do\_paciente; parte\_de etica\_pesquisa\_humanos}
\prov{proveniência: wiki \textperiodcentered{} inferencia \textperiodcentered{} confiança 1.0 \textperiodcentered{} domínio ciencias\_de\_fora \textperiodcentered{} história 3 versões no jornal}

%CRISTAL {"arestas":[{"alvo":"pitagoras_razoes_musicais","confianca":1.0,"epistemico":"fato","origem":"wiki","rel":"usa"},{"alvo":"elementos_musica","confianca":1.0,"epistemico":"fato","origem":"wiki","rel":"relaciona"}],"confianca":1.0,"contraexemplos":[],"descricao":"A qualidade de estabilidade/instabilidade percebida entre alturas: razões de frequência simples (2:1, 3:2) tendem a soar consonantes, complexas dissonantes — mas o limiar do 'aceitável' é histórica e culturalmente móvel.","epistemico":"inferencia","exemplos":[],"id":"consonancia_dissonancia","memoria":"semantica","meta":{"autor":"—","dominio":"ciencias_de_fora","fonte":""},"origem":"wiki","palavras_chave":[],"sinonimos":[],"tipo":"conceito","titulo":"Consonância e Dissonância"}
\section{Consonância e Dissonância}
A qualidade de estabilidade/instabilidade percebida entre alturas: razões de frequência simples (2:1, 3:2) tendem a soar consonantes, complexas dissonantes — mas o limiar do 'aceitável' é histórica e culturalmente móvel.
\prov{relações: usa pitagoras\_razoes\_musicais; relaciona elementos\_musica}
\prov{proveniência: wiki \textperiodcentered{} inferencia \textperiodcentered{} confiança 1.0 \textperiodcentered{} domínio ciencias\_de\_fora \textperiodcentered{} história 3 versões no jornal}

%CRISTAL {"arestas":[{"alvo":"estado_de_direito","confianca":1.0,"epistemico":"fato","origem":"wiki","rel":"usa"},{"alvo":"contrato_social","confianca":1.0,"epistemico":"fato","origem":"wiki","rel":"relaciona"},{"alvo":"direitos_humanos","confianca":1.0,"epistemico":"fato","origem":"wiki","rel":"relaciona"}],"confianca":1.0,"contraexemplos":[],"descricao":"O governo pode e deve ser juridicamente limitado em seus poderes, e sua legitimidade depende de observar esses limites; distingue o poder constituinte (que cria a constituição) do constituído (que age dentro dela).","epistemico":"inferencia","exemplos":[],"id":"constitucionalismo","memoria":"semantica","meta":{"autor":"Locke/Sieyès","dominio":"ciencias_de_fora","fonte":""},"origem":"wiki","palavras_chave":[],"sinonimos":[],"tipo":"conceito","titulo":"Constitucionalismo"}
\section{Constitucionalismo}
O governo pode e deve ser juridicamente limitado em seus poderes, e sua legitimidade depende de observar esses limites; distingue o poder constituinte (que cria a constituição) do constituído (que age dentro dela).
\prov{relações: usa estado\_de\_direito; relaciona contrato\_social; relaciona direitos\_humanos}
\prov{proveniência: wiki \textperiodcentered{} inferencia \textperiodcentered{} confiança 1.0 \textperiodcentered{} domínio ciencias\_de\_fora \textperiodcentered{} história 4 versões no jornal}

%CRISTAL {"arestas":[{"alvo":"filosofia_da_tecnica","confianca":1.0,"epistemico":"fato","origem":"wiki","rel":"parte_de"},{"alvo":"determinismo_tecnologico","confianca":1.0,"epistemico":"fato","origem":"wiki","rel":"contradiz"},{"alvo":"contrato_social","confianca":1.0,"epistemico":"fato","origem":"wiki","rel":"relaciona"}],"confianca":1.0,"contraexemplos":[],"descricao":"Pinch & Bijker: o sucesso ou fracasso de um artefato não se explica por sua eficácia 'intrínseca', mas por grupos sociais relevantes, interpretações concorrentes e um processo de estabilização. O polo 'a sociedade molda a técnica'.","epistemico":"inferencia","exemplos":[],"id":"construcao_social_tecnologia","memoria":"semantica","meta":{"autor":"Pinch & Bijker","dominio":"ciencias_de_fora","fonte":""},"origem":"wiki","palavras_chave":[],"sinonimos":[],"tipo":"conceito","titulo":"Construção Social da Tecnologia (SCOT)"}
\section{Construção Social da Tecnologia (SCOT)}
Pinch \& Bijker: o sucesso ou fracasso de um artefato não se explica por sua eficácia 'intrínseca', mas por grupos sociais relevantes, interpretações concorrentes e um processo de estabilização. O polo 'a sociedade molda a técnica'.
\prov{relações: parte\_de filosofia\_da\_tecnica; contradiz determinismo\_tecnologico; relaciona contrato\_social}
\prov{proveniência: wiki \textperiodcentered{} inferencia \textperiodcentered{} confiança 1.0 \textperiodcentered{} domínio ciencias\_de\_fora \textperiodcentered{} história 4 versões no jornal}

%CRISTAL {"arestas":[{"alvo":"teorias_aprendizagem","confianca":1.0,"epistemico":"fato","origem":"wiki","rel":"parte_de"},{"alvo":"aparelho_psiquico","confianca":1.0,"epistemico":"fato","origem":"wiki","rel":"relaciona"}],"confianca":1.0,"contraexemplos":[],"descricao":"Piaget: o conhecimento não é transmitido pronto, mas construído ativamente pelo sujeito na interação com o meio; a inteligência se desenvolve por adaptação, não por acúmulo.","epistemico":"inferencia","exemplos":[],"id":"construtivismo_piaget","memoria":"semantica","meta":{"autor":"Piaget","dominio":"ciencias_de_fora","fonte":""},"origem":"wiki","palavras_chave":[],"sinonimos":[],"tipo":"conceito","titulo":"Construtivismo (Piaget)"}
\section{Construtivismo (Piaget)}
Piaget: o conhecimento não é transmitido pronto, mas construído ativamente pelo sujeito na interação com o meio; a inteligência se desenvolve por adaptação, não por acúmulo.
\prov{relações: parte\_de teorias\_aprendizagem; relaciona aparelho\_psiquico}
\prov{proveniência: wiki \textperiodcentered{} inferencia \textperiodcentered{} confiança 1.0 \textperiodcentered{} domínio ciencias\_de\_fora \textperiodcentered{} história 3 versões no jornal}

%CRISTAL {"arestas":[{"alvo":"teorias_aprendizagem","confianca":1.0,"epistemico":"fato","origem":"wiki","rel":"parte_de"},{"alvo":"aprendizagem_descoberta_bruner","confianca":1.0,"epistemico":"fato","origem":"wiki","rel":"usa"},{"alvo":"condicionamento_operante","confianca":1.0,"epistemico":"fato","origem":"wiki","rel":"contradiz"},{"alvo":"mente_computacao","confianca":1.0,"epistemico":"fato","origem":"wiki","rel":"referencia"}],"confianca":1.0,"contraexemplos":[],"descricao":"A controvérsia sobre quanta orientação a aprendizagem exige: instrução direta e sequenciada vs. descoberta centrada no aluno. Kirschner-Sweller-Clark: instrução minimamente guiada é ineficaz para novatos; Papert (construcionismo/LOGO) no polo oposto.","epistemico":"inferencia","exemplos":[],"id":"construtivismo_vs_instrucionismo","memoria":"semantica","meta":{"autor":"Papert vs Kirschner","dominio":"ciencias_de_fora","fonte":""},"origem":"wiki","palavras_chave":[],"sinonimos":[],"tipo":"conceito","titulo":"Debate Construtivismo × Instrucionismo"}
\section{Debate Construtivismo $\times$ Instrucionismo}
A controvérsia sobre quanta orientação a aprendizagem exige: instrução direta e sequenciada vs. descoberta centrada no aluno. Kirschner-Sweller-Clark: instrução minimamente guiada é ineficaz para novatos; Papert (construcionismo/LOGO) no polo oposto.
\prov{relações: parte\_de teorias\_aprendizagem; usa aprendizagem\_descoberta\_bruner; contradiz condicionamento\_operante; referencia mente\_computacao}
\prov{proveniência: wiki \textperiodcentered{} inferencia \textperiodcentered{} confiança 1.0 \textperiodcentered{} domínio ciencias\_de\_fora \textperiodcentered{} história 5 versões no jornal}

%CRISTAL {"arestas":[{"alvo":"relacoes_internacionais","confianca":1.0,"epistemico":"fato","origem":"wiki","rel":"parte_de"},{"alvo":"realismo_ri","confianca":1.0,"epistemico":"fato","origem":"wiki","rel":"contradiz"},{"alvo":"hegemonia_cultural","confianca":1.0,"epistemico":"fato","origem":"wiki","rel":"relaciona"}],"confianca":1.0,"contraexemplos":[],"descricao":"Wendt: as estruturas da política internacional são feitas de ideias e normas partilhadas, não só de forças materiais; identidades e interesses dos Estados são socialmente construídos na interação, de modo que a anarquia não tem lógica fixa.","epistemico":"inferencia","exemplos":[],"id":"construtivismo_wendt","memoria":"semantica","meta":{"autor":"Wendt","dominio":"ciencias_de_fora","fonte":""},"origem":"wiki","palavras_chave":[],"sinonimos":[],"tipo":"conceito","titulo":"Construtivismo — 'A Anarquia é o que os Estados Fazem Dela' (Wendt)"}
\section{Construtivismo — 'A Anarquia é o que os Estados Fazem Dela' (Wendt)}
Wendt: as estruturas da política internacional são feitas de ideias e normas partilhadas, não só de forças materiais; identidades e interesses dos Estados são socialmente construídos na interação, de modo que a anarquia não tem lógica fixa.
\prov{relações: parte\_de relacoes\_internacionais; contradiz realismo\_ri; relaciona hegemonia\_cultural}
\prov{proveniência: wiki \textperiodcentered{} inferencia \textperiodcentered{} confiança 1.0 \textperiodcentered{} domínio ciencias\_de\_fora \textperiodcentered{} história 4 versões no jornal}

%CRISTAL {"arestas":[{"alvo":"filosofia","confianca":1.0,"epistemico":"fato","origem":"wiki","rel":"parte_de"},{"alvo":"colapso","confianca":1.0,"epistemico":"fato","origem":"wiki","rel":"referencia"},{"alvo":"virtualizacao","confianca":0.5,"epistemico":"fato","origem":"wiki","rel":"relaciona"},{"alvo":"ethics_of_care","confianca":0.5,"epistemico":"fato","origem":"wiki","rel":"relaciona"},{"alvo":"word","confianca":0.5,"epistemico":"fato","origem":"wiki","rel":"relaciona"}],"confianca":1.0,"contraexemplos":[],"descricao":"O modo de conhecer pela quietude e observação (theoria), oposto à ação — a borda do broca-so.","epistemico":"fato","exemplos":[],"id":"contemplacao","memoria":"semantica","meta":{"dominio":"filosofia","fonte":""},"origem":"wiki","palavras_chave":[],"sinonimos":[],"tipo":"conceito","titulo":"Contemplação"}
\section{Contemplação}
O modo de conhecer pela quietude e observação (theoria), oposto à ação — a borda do broca-so.
\prov{relações: parte\_de filosofia; referencia colapso; relaciona virtualizacao; relaciona ethics\_of\_care; relaciona word}
\prov{proveniência: wiki \textperiodcentered{} fato \textperiodcentered{} confiança 1.0 \textperiodcentered{} domínio filosofia \textperiodcentered{} história 6 versões no jornal}

%CRISTAL {"arestas":[{"alvo":"filosofia_da_historia","confianca":1.0,"epistemico":"fato","origem":"wiki","rel":"parte_de"}],"confianca":1.0,"contraexemplos":[],"descricao":"O eixo em que toda filosofia da história se situa: a história é necessária (segue leis) ou contingente (poderia ter sido outra)? Debate irredutível.","epistemico":"inferencia","exemplos":[],"id":"contingencia_historica","memoria":"semantica","meta":{"autor":"—","dominio":"ciencias_de_fora","fonte":""},"origem":"wiki","palavras_chave":[],"sinonimos":[],"tipo":"conceito","titulo":"Contingência vs Determinismo"}
\section{Contingência vs Determinismo}
O eixo em que toda filosofia da história se situa: a história é necessária (segue leis) ou contingente (poderia ter sido outra)? Debate irredutível.
\prov{relações: parte\_de filosofia\_da\_historia}
\prov{proveniência: wiki \textperiodcentered{} inferencia \textperiodcentered{} confiança 1.0 \textperiodcentered{} domínio ciencias\_de\_fora \textperiodcentered{} história 2 versões no jornal}

%CRISTAL {"arestas":[{"alvo":"teoria_musical","confianca":1.0,"epistemico":"fato","origem":"wiki","rel":"parte_de"}],"confianca":1.0,"contraexemplos":[],"descricao":"A arte de combinar melodias independentes (Bach) — vozes que se movem sozinhas e ainda soam juntas.","epistemico":"fato","exemplos":[],"id":"contraponto","memoria":"semantica","meta":{"autor":"Bach","dominio":"ciencias_de_fora","fonte":""},"origem":"wiki","palavras_chave":[],"sinonimos":[],"tipo":"conceito","titulo":"Contraponto"}
\section{Contraponto}
A arte de combinar melodias independentes (Bach) — vozes que se movem sozinhas e ainda soam juntas.
\prov{relações: parte\_de teoria\_musical}
\prov{proveniência: wiki \textperiodcentered{} fato \textperiodcentered{} confiança 1.0 \textperiodcentered{} domínio ciencias\_de\_fora \textperiodcentered{} história 2 versões no jornal}

%CRISTAL {"arestas":[{"alvo":"brousseau_situacoes_didaticas","confianca":1.0,"epistemico":"fato","origem":"wiki","rel":"usa"},{"alvo":"tiffany","confianca":1.0,"epistemico":"fato","origem":"wiki","rel":"referencia"}],"confianca":1.0,"contraexemplos":[],"descricao":"Brousseau: o sistema de obrigações recíprocas, quase sempre implícitas, que regula o que professor e aluno esperam um do outro; sua ruptura revela (ou causa) dificuldades de aprendizagem.","epistemico":"inferencia","exemplos":[],"id":"contrato_didatico","memoria":"semantica","meta":{"autor":"Brousseau","dominio":"ciencias_de_fora","fonte":""},"origem":"wiki","palavras_chave":[],"sinonimos":[],"tipo":"conceito","titulo":"Contrato Didático"}
\section{Contrato Didático}
Brousseau: o sistema de obrigações recíprocas, quase sempre implícitas, que regula o que professor e aluno esperam um do outro; sua ruptura revela (ou causa) dificuldades de aprendizagem.
\prov{relações: usa brousseau\_situacoes\_didaticas; referencia tiffany}
\prov{proveniência: wiki \textperiodcentered{} inferencia \textperiodcentered{} confiança 1.0 \textperiodcentered{} domínio ciencias\_de\_fora \textperiodcentered{} história 3 versões no jornal}

%CRISTAL {"arestas":[{"alvo":"politica","confianca":1.0,"epistemico":"fato","origem":"wiki","rel":"parte_de"},{"alvo":"contrato_objeto","confianca":1.0,"epistemico":"fato","origem":"wiki","rel":"relaciona"},{"alvo":"goldchain","confianca":1.0,"epistemico":"fato","origem":"wiki","rel":"referencia"},{"alvo":"word","confianca":0.5,"epistemico":"fato","origem":"wiki","rel":"relaciona"}],"confianca":1.0,"contraexemplos":[],"descricao":"Hobbes/Locke/Rousseau: a autoridade legítima nasce do consentimento dos indivíduos — segurança, direitos, vontade geral.","epistemico":"inferencia","exemplos":[],"id":"contrato_social","memoria":"semantica","meta":{"autor":"Hobbes/Locke/Rousseau","dominio":"ciencias_de_fora","fonte":""},"origem":"wiki","palavras_chave":[],"sinonimos":[],"tipo":"conceito","titulo":"Contrato Social"}
\section{Contrato Social}
Hobbes/Locke/Rousseau: a autoridade legítima nasce do consentimento dos indivíduos — segurança, direitos, vontade geral.
\prov{relações: parte\_de politica; relaciona contrato\_objeto; referencia goldchain; relaciona word}
\prov{proveniência: wiki \textperiodcentered{} inferencia \textperiodcentered{} confiança 1.0 \textperiodcentered{} domínio ciencias\_de\_fora \textperiodcentered{} história 5 versões no jornal}

%CRISTAL {"arestas":[{"alvo":"didatica_matematica","confianca":1.0,"epistemico":"fato","origem":"wiki","rel":"parte_de"},{"alvo":"transposicao_didatica","confianca":1.0,"epistemico":"fato","origem":"wiki","rel":"relaciona"},{"alvo":"intuicionismo_brouwer","confianca":1.0,"epistemico":"fato","origem":"wiki","rel":"relaciona"},{"alvo":"tiffany","confianca":1.0,"epistemico":"fato","origem":"wiki","rel":"relaciona"}],"confianca":1.0,"contraexemplos":[],"descricao":"Lopes ataca o que se ensina como verdade eterna: sustenta que (−1)·(−1) 'não é lei da natureza nem da matemática, mas escolha' (pode-se escolher +1 ou −1 com igual legitimidade), e que uma 'Tabuada do Iconoclasta' alternativa não é desonestidade — as convenções aritméticas não são absolutas.","epistemico":"inferencia","exemplos":[],"id":"convencao_vs_verdade_gentil","memoria":"semantica","meta":{"autor":"Gentil Lopes","dominio":"ciencias_de_fora","fonte":""},"origem":"wiki","palavras_chave":[],"sinonimos":[],"tipo":"conceito","titulo":"Convenção vs Verdade Absoluta (Lopes)"}
\section{Convenção vs Verdade Absoluta (Lopes)}
Lopes ataca o que se ensina como verdade eterna: sustenta que (\ensuremath{-}1)\textperiodcentered (\ensuremath{-}1) 'não é lei da natureza nem da matemática, mas escolha' (pode-se escolher +1 ou \ensuremath{-}1 com igual legitimidade), e que uma 'Tabuada do Iconoclasta' alternativa não é desonestidade — as convenções aritméticas não são absolutas.
\prov{relações: parte\_de didatica\_matematica; relaciona transposicao\_didatica; relaciona intuicionismo\_brouwer; relaciona tiffany}
\prov{proveniência: wiki \textperiodcentered{} inferencia \textperiodcentered{} confiança 1.0 \textperiodcentered{} domínio ciencias\_de\_fora \textperiodcentered{} história 6 versões no jornal}

%CRISTAL {"arestas":[{"alvo":"algebra_abstrata","confianca":1.0,"epistemico":"fato","origem":"wiki","rel":"parte_de"},{"alvo":"numeros_complexos","confianca":1.0,"epistemico":"fato","origem":"wiki","rel":"relaciona"}],"confianca":1.0,"contraexemplos":[],"descricao":"As arenas da álgebra: anel (soma e produto), corpo (mais a divisão). ℚ√(n²+4) é o corpo onde vive cada reta metálica.","epistemico":"fato","exemplos":[],"id":"corpos_e_aneis","memoria":"semantica","meta":{"dominio":"ciencias_de_fora","fonte":""},"origem":"wiki","palavras_chave":[],"sinonimos":[],"tipo":"conceito","titulo":"Corpos e Anéis"}
\section{Corpos e Anéis}
As arenas da álgebra: anel (soma e produto), corpo (mais a divisão). \ensuremath{\mathbb{Q}}\ensuremath{\sqrt{\,}}(n\ensuremath{^{2}}+4) é o corpo onde vive cada reta metálica.
\prov{relações: parte\_de algebra\_abstrata; relaciona numeros\_complexos}
\prov{proveniência: wiki \textperiodcentered{} fato \textperiodcentered{} confiança 1.0 \textperiodcentered{} domínio ciencias\_de\_fora \textperiodcentered{} história 6 versões no jornal}

%CRISTAL {"arestas":[{"alvo":"filosofia_da_tecnica","confianca":1.0,"epistemico":"fato","origem":"wiki","rel":"parte_de"},{"alvo":"matematica_arte_engenharia","confianca":1.0,"epistemico":"fato","origem":"wiki","rel":"relaciona"},{"alvo":"maquina_evolutiva","confianca":1.0,"epistemico":"fato","origem":"wiki","rel":"relaciona"}],"confianca":1.0,"contraexemplos":[],"descricao":"Lopes endossa Cantor/Chaitin: 'a essência da matemática reside na liberdade de criar', e a história julga essas criações por sua fertilidade — e isso 'se aplica a todas as áreas do conhecimento'. A poiesis implícita: criar como valor.","epistemico":"inferencia","exemplos":[],"id":"criacao_livre_fertilidade","memoria":"semantica","meta":{"autor":"Gentil Lopes","dominio":"ciencias_de_fora","fonte":""},"origem":"wiki","palavras_chave":[],"sinonimos":[],"tipo":"conceito","titulo":"Criação Livre Julgada por Fertilidade (Lopes)"}
\section{Criação Livre Julgada por Fertilidade (Lopes)}
Lopes endossa Cantor/Chaitin: 'a essência da matemática reside na liberdade de criar', e a história julga essas criações por sua fertilidade — e isso 'se aplica a todas as áreas do conhecimento'. A poiesis implícita: criar como valor.
\prov{relações: parte\_de filosofia\_da\_tecnica; relaciona matematica\_arte\_engenharia; relaciona maquina\_evolutiva}
\prov{proveniência: wiki \textperiodcentered{} inferencia \textperiodcentered{} confiança 1.0 \textperiodcentered{} domínio ciencias\_de\_fora \textperiodcentered{} história 5 versões no jornal}

%CRISTAL {"arestas":[{"alvo":"idioma","confianca":1.0,"epistemico":"fato","origem":"wiki","rel":"usa"},{"alvo":"lingua_franca","confianca":1.0,"epistemico":"fato","origem":"wiki","rel":"relaciona"}],"confianca":1.0,"contraexemplos":[],"descricao":"Pidgin: língua de contato simplificada, sem nativos. Crioulo: quando as crianças a adotam como materna e ela ganha gramática plena. Língua nascendo.","epistemico":"fato","exemplos":[],"id":"crioulo_e_pidgin","memoria":"semantica","meta":{"dominio":"ciencias_de_fora","fonte":""},"origem":"wiki","palavras_chave":[],"sinonimos":[],"tipo":"conceito","titulo":"Pidgin e Crioulo"}
\section{Pidgin e Crioulo}
Pidgin: língua de contato simplificada, sem nativos. Crioulo: quando as crianças a adotam como materna e ela ganha gramática plena. Língua nascendo.
\prov{relações: usa idioma; relaciona lingua\_franca}
\prov{proveniência: wiki \textperiodcentered{} fato \textperiodcentered{} confiança 1.0 \textperiodcentered{} domínio ciencias\_de\_fora \textperiodcentered{} história 3 versões no jornal}

%CRISTAL {"arestas":[{"alvo":"prova_de_trabalho","confianca":1.0,"epistemico":"fato","origem":"wiki","rel":"usa"},{"alvo":"teoria_dos_jogos","confianca":1.0,"epistemico":"fato","origem":"wiki","rel":"usa"},{"alvo":"goldchain","confianca":1.0,"epistemico":"fato","origem":"wiki","rel":"explica"},{"alvo":"cristalchain","confianca":1.0,"epistemico":"fato","origem":"wiki","rel":"relaciona"}],"confianca":1.0,"contraexemplos":[],"descricao":"Une criptografia, economia e teoria dos jogos para alinhar incentivos sem intermediário central — a honestidade sai mais barata que a fraude.","epistemico":"fato","exemplos":[],"id":"criptoeconomia","memoria":"semantica","meta":{"autor":"Nakamoto","dominio":"ciencias_de_fora","fonte":""},"origem":"wiki","palavras_chave":[],"sinonimos":[],"tipo":"conceito","titulo":"Criptoeconomia"}
\section{Criptoeconomia}
Une criptografia, economia e teoria dos jogos para alinhar incentivos sem intermediário central — a honestidade sai mais barata que a fraude.
\prov{relações: usa prova\_de\_trabalho; usa teoria\_dos\_jogos; explica goldchain; relaciona cristalchain}
\prov{proveniência: wiki \textperiodcentered{} fato \textperiodcentered{} confiança 1.0 \textperiodcentered{} domínio ciencias\_de\_fora \textperiodcentered{} história 5 versões no jornal}

%CRISTAL {"arestas":[{"alvo":"genoma","confianca":1.0,"epistemico":"fato","origem":"wiki","rel":"usa"},{"alvo":"crispr_alphafold","confianca":1.0,"epistemico":"fato","origem":"wiki","rel":"relaciona"}],"confianca":1.0,"contraexemplos":[],"descricao":"Uma defesa bacteriana convertida em ferramenta de edição genética dirigida por RNA — cortar e reescrever o genoma.","epistemico":"fato","exemplos":[],"id":"crispr","memoria":"semantica","meta":{"dominio":"ciencias_de_fora","fonte":""},"origem":"wiki","palavras_chave":[],"sinonimos":[],"tipo":"conceito","titulo":"CRISPR-Cas9"}
\section{CRISPR-Cas9}
Uma defesa bacteriana convertida em ferramenta de edição genética dirigida por RNA — cortar e reescrever o genoma.
\prov{relações: usa genoma; relaciona crispr\_alphafold}
\prov{proveniência: wiki \textperiodcentered{} fato \textperiodcentered{} confiança 1.0 \textperiodcentered{} domínio ciencias\_de\_fora \textperiodcentered{} história 3 versões no jornal}

%CRISTAL {"arestas":[{"alvo":"dna_dupla_helice","confianca":1.0,"epistemico":"fato","origem":"wiki","rel":"usa"},{"alvo":"codigo_genetico","confianca":1.0,"epistemico":"fato","origem":"wiki","rel":"relaciona"}],"confianca":1.0,"contraexemplos":[],"descricao":"CRISPR-Cas9 (Doudna-Charpentier, Nobel Química 2020): sistema imune bacteriano reaproveitado — um RNA-guia dirige a nuclease Cas9 a um alvo de DNA e corta, permitindo editar genes com precisão. AlphaFold (Jumper-Hassabis, Nobel Química 2024): deep learning prevê a estrutura 3D de proteínas a partir da sequência (~200 milhões de estruturas), atacando o problema do enovelamento (paradoxo de Levinthal). Predição de estrutura = fato; o mecanismo dinâmico do enovelamento = parte aberta.","epistemico":"fato","exemplos":[],"id":"crispr_alphafold","memoria":"semantica","meta":{"autor":"broca-so","dominio":"biologia","fonte":"Doudna-Charpentier Nobel 2020; AlphaFold Nature 596 (2021), Nobel 2024"},"origem":"wiki","palavras_chave":[],"sinonimos":[],"tipo":"conceito","titulo":"CRISPR-Cas9 e AlphaFold (as ferramentas modernas)"}
\section{CRISPR-Cas9 e AlphaFold (as ferramentas modernas)}
CRISPR-Cas9 (Doudna-Charpentier, Nobel Química 2020): sistema imune bacteriano reaproveitado — um RNA-guia dirige a nuclease Cas9 a um alvo de DNA e corta, permitindo editar genes com precisão. AlphaFold (Jumper-Hassabis, Nobel Química 2024): deep learning prevê a estrutura 3D de proteínas a partir da sequência (\~{}200 milhões de estruturas), atacando o problema do enovelamento (paradoxo de Levinthal). Predição de estrutura = fato; o mecanismo dinâmico do enovelamento = parte aberta.
\prov{relações: usa dna\_dupla\_helice; relaciona codigo\_genetico}
\prov{proveniência: wiki \textperiodcentered{} Doudna-Charpentier Nobel 2020; AlphaFold Nature 596 (2021), Nobel 2024 \textperiodcentered{} fato \textperiodcentered{} confiança 1.0 \textperiodcentered{} domínio biologia \textperiodcentered{} história 3 versões no jornal}

%CRISTAL {"arestas":[{"alvo":"sagrado","confianca":1.0,"epistemico":"fato","origem":"wiki","rel":"parte_de"},{"alvo":"linhagem_abraamica","confianca":1.0,"epistemico":"fato","origem":"wiki","rel":"usa"},{"alvo":"word","confianca":0.5,"epistemico":"fato","origem":"wiki","rel":"relaciona"},{"alvo":"vontade_de_poder","confianca":0.5,"epistemico":"fato","origem":"wiki","rel":"relaciona"},{"alvo":"while","confianca":0.5,"epistemico":"fato","origem":"wiki","rel":"relaciona"}],"confianca":1.0,"contraexemplos":[],"descricao":"A tradição da Bíblia — o Deus criador e pessoal, a aliança, o Logos que se faz carne, o amor (ágape) e a esperança da ressurreição.","epistemico":"fato","exemplos":[],"id":"cristianismo","memoria":"semantica","meta":{"dominio":"sagrado","fonte":""},"origem":"wiki","palavras_chave":[],"sinonimos":[],"tipo":"conceito","titulo":"Tradição Judaico-Cristã (Bíblia)"}
\section{Tradição Judaico-Cristã (Bíblia)}
A tradição da Bíblia — o Deus criador e pessoal, a aliança, o Logos que se faz carne, o amor (ágape) e a esperança da ressurreição.
\prov{relações: parte\_de sagrado; usa linhagem\_abraamica; relaciona word; relaciona vontade\_de\_poder; relaciona while}
\prov{proveniência: wiki \textperiodcentered{} fato \textperiodcentered{} confiança 1.0 \textperiodcentered{} domínio sagrado \textperiodcentered{} história 6 versões no jornal}

%CRISTAL {"arestas":[{"alvo":"didatica_matematica","confianca":1.0,"epistemico":"fato","origem":"wiki","rel":"parte_de"},{"alvo":"aprendizagem_significativa_ausubel","confianca":1.0,"epistemico":"fato","origem":"wiki","rel":"relaciona"},{"alvo":"educacao_bancaria","confianca":1.0,"epistemico":"fato","origem":"wiki","rel":"relaciona"}],"confianca":1.0,"contraexemplos":[],"descricao":"Lopes critica a memorização vazia — 'ao encher a cabeça de entulhos o pobre não teve tempo de aprender' o essencial. Ecoa a aprendizagem significativa (contra a mecânica) e a educação problematizadora (contra a bancária).","epistemico":"inferencia","exemplos":[],"id":"critica_decoreba_gentil","memoria":"semantica","meta":{"autor":"Gentil Lopes","dominio":"ciencias_de_fora","fonte":""},"origem":"wiki","palavras_chave":[],"sinonimos":[],"tipo":"conceito","titulo":"Crítica à Decoreba (Lopes)"}
\section{Crítica à Decoreba (Lopes)}
Lopes critica a memorização vazia — 'ao encher a cabeça de entulhos o pobre não teve tempo de aprender' o essencial. Ecoa a aprendizagem significativa (contra a mecânica) e a educação problematizadora (contra a bancária).
\prov{relações: parte\_de didatica\_matematica; relaciona aprendizagem\_significativa\_ausubel; relaciona educacao\_bancaria}
\prov{proveniência: wiki \textperiodcentered{} inferencia \textperiodcentered{} confiança 1.0 \textperiodcentered{} domínio ciencias\_de\_fora \textperiodcentered{} história 5 versões no jornal}

%CRISTAL {"arestas":[{"alvo":"filosofia_da_computacao","confianca":1.0,"epistemico":"fato","origem":"wiki","rel":"parte_de"},{"alvo":"simbolo_vs_significado","confianca":1.0,"epistemico":"fato","origem":"wiki","rel":"usa"},{"alvo":"mente_computacao","confianca":1.0,"epistemico":"fato","origem":"wiki","rel":"contradiz"},{"alvo":"corpo_dual","confianca":1.0,"epistemico":"fato","origem":"wiki","rel":"relaciona"}],"confianca":1.0,"contraexemplos":[],"descricao":"Dreyfus: a inteligência humana repousa em habilidades corporificadas, contexto e conhecimento tácito de fundo que não se reduzem a regras e símbolos explícitos — por isso a IA simbólica clássica falha no que 'os computadores não podem fazer'.","epistemico":"inferencia","exemplos":[],"id":"critica_fenomenologica_dreyfus","memoria":"semantica","meta":{"autor":"Dreyfus","dominio":"ciencias_de_fora","fonte":""},"origem":"wiki","palavras_chave":[],"sinonimos":[],"tipo":"conceito","titulo":"Crítica Fenomenológica à IA (Dreyfus)"}
\section{Crítica Fenomenológica à IA (Dreyfus)}
Dreyfus: a inteligência humana repousa em habilidades corporificadas, contexto e conhecimento tácito de fundo que não se reduzem a regras e símbolos explícitos — por isso a IA simbólica clássica falha no que 'os computadores não podem fazer'.
\prov{relações: parte\_de filosofia\_da\_computacao; usa simbolo\_vs\_significado; contradiz mente\_computacao; relaciona corpo\_dual}
\prov{proveniência: wiki \textperiodcentered{} inferencia \textperiodcentered{} confiança 1.0 \textperiodcentered{} domínio ciencias\_de\_fora \textperiodcentered{} história 5 versões no jornal}

%CRISTAL {"arestas":[{"alvo":"filosofia_da_computacao","confianca":1.0,"epistemico":"fato","origem":"wiki","rel":"parte_de"},{"alvo":"simbolo_vs_significado","confianca":1.0,"epistemico":"fato","origem":"wiki","rel":"usa"},{"alvo":"critica_fenomenologica_dreyfus","confianca":1.0,"epistemico":"fato","origem":"wiki","rel":"relaciona"},{"alvo":"inteligencia_artificial","confianca":1.0,"epistemico":"fato","origem":"wiki","rel":"relaciona"}],"confianca":1.0,"contraexemplos":[],"descricao":"Lopes questiona se a IA é 'inteligente', dado que sequer compreendemos a inteligência natural: assim como biólogos não sabem o que é vida e matemáticos o que é número, a ciência opera sobre noções apenas funcionalmente manipuladas. 'A humanidade progride na técnica, mas empaca na semântica.'","epistemico":"inferencia","exemplos":[],"id":"critica_ia_semantica","memoria":"semantica","meta":{"autor":"Gentil Lopes","dominio":"ciencias_de_fora","fonte":""},"origem":"wiki","palavras_chave":[],"sinonimos":[],"tipo":"conceito","titulo":"Progride na Técnica, Empaca na Semântica (Lopes)"}
\section{Progride na Técnica, Empaca na Semântica (Lopes)}
Lopes questiona se a IA é 'inteligente', dado que sequer compreendemos a inteligência natural: assim como biólogos não sabem o que é vida e matemáticos o que é número, a ciência opera sobre noções apenas funcionalmente manipuladas. 'A humanidade progride na técnica, mas empaca na semântica.'
\prov{relações: parte\_de filosofia\_da\_computacao; usa simbolo\_vs\_significado; relaciona critica\_fenomenologica\_dreyfus; relaciona inteligencia\_artificial}
\prov{proveniência: wiki \textperiodcentered{} inferencia \textperiodcentered{} confiança 1.0 \textperiodcentered{} domínio ciencias\_de\_fora \textperiodcentered{} história 6 versões no jornal}

%CRISTAL {"arestas":[{"alvo":"filosofia_da_educacao","confianca":1.0,"epistemico":"fato","origem":"wiki","rel":"parte_de"},{"alvo":"capital_cultural","confianca":1.0,"epistemico":"fato","origem":"wiki","rel":"usa"},{"alvo":"habitus","confianca":1.0,"epistemico":"fato","origem":"wiki","rel":"usa"},{"alvo":"igualdade_das_inteligencias","confianca":1.0,"epistemico":"fato","origem":"wiki","rel":"relaciona"}],"confianca":1.0,"contraexemplos":[],"descricao":"Sandel/Bourdieu: sob a retórica da igualdade de oportunidades, a escola tende a reproduzir desigualdades de origem — o 'mérito' medido correlaciona com capital cultural familiar, convertendo vantagem herdada em suposto talento.","epistemico":"inferencia","exemplos":[],"id":"critica_meritocracia_escolar","memoria":"semantica","meta":{"autor":"Sandel/Bourdieu","dominio":"ciencias_de_fora","fonte":""},"origem":"wiki","palavras_chave":[],"sinonimos":[],"tipo":"conceito","titulo":"Crítica à Meritocracia Escolar"}
\section{Crítica à Meritocracia Escolar}
Sandel/Bourdieu: sob a retórica da igualdade de oportunidades, a escola tende a reproduzir desigualdades de origem — o 'mérito' medido correlaciona com capital cultural familiar, convertendo vantagem herdada em suposto talento.
\prov{relações: parte\_de filosofia\_da\_educacao; usa capital\_cultural; usa habitus; relaciona igualdade\_das\_inteligencias}
\prov{proveniência: wiki \textperiodcentered{} inferencia \textperiodcentered{} confiança 1.0 \textperiodcentered{} domínio ciencias\_de\_fora \textperiodcentered{} história 5 versões no jornal}

%CRISTAL {"arestas":[{"alvo":"razao_aurea","confianca":1.0,"epistemico":"fato","origem":"wiki","rel":"relaciona"},{"alvo":"beleza_matematica","confianca":1.0,"epistemico":"fato","origem":"wiki","rel":"contradiz"},{"alvo":"honestidade_intelectual","confianca":1.0,"epistemico":"fato","origem":"wiki","rel":"usa"}],"confianca":1.0,"contraexemplos":[],"descricao":"Lopes: ama φ na matemática, mas faz paródia contra o mito estético — 'a fórmula que mede beleza é fictícia, eu inventei; beleza não se equaciona'. Honestidade contra a pseudociência.","epistemico":"inferencia","exemplos":[],"id":"critica_proporcao_aurea","memoria":"semantica","meta":{"autor":"Gentil Lopes","dominio":"ciencias_de_fora","fonte":""},"origem":"livro","palavras_chave":[],"sinonimos":[],"tipo":"conceito","titulo":"Crítica ao Mito da Proporção Áurea"}
\section{Crítica ao Mito da Proporção Áurea}
Lopes: ama \ensuremath{\varphi} na matemática, mas faz paródia contra o mito estético — 'a fórmula que mede beleza é fictícia, eu inventei; beleza não se equaciona'. Honestidade contra a pseudociência.
\prov{relações: relaciona razao\_aurea; contradiz beleza\_matematica; usa honestidade\_intelectual}
\prov{proveniência: livro \textperiodcentered{} inferencia \textperiodcentered{} confiança 1.0 \textperiodcentered{} domínio ciencias\_de\_fora \textperiodcentered{} história 5 versões no jornal}

%CRISTAL {"arestas":[{"alvo":"filosofia_do_direito","confianca":1.0,"epistemico":"fato","origem":"wiki","rel":"parte_de"},{"alvo":"realismo_juridico","confianca":1.0,"epistemico":"fato","origem":"wiki","rel":"usa"},{"alvo":"biopoder","confianca":1.0,"epistemico":"fato","origem":"wiki","rel":"usa"},{"alvo":"injustica_epistemica","confianca":1.0,"epistemico":"fato","origem":"wiki","rel":"relaciona"}],"confianca":1.0,"contraexemplos":[],"descricao":"Herdando o realismo jurídico e o marxismo: o direito é fundamentalmente indeterminado, e sua aparência de neutralidade e certeza mascara os interesses dos poderosos e escolhas políticas.","epistemico":"inferencia","exemplos":[],"id":"critica_realista_direito","memoria":"semantica","meta":{"autor":"Kennedy/Unger","dominio":"ciencias_de_fora","fonte":""},"origem":"wiki","palavras_chave":[],"sinonimos":[],"tipo":"conceito","titulo":"Critical Legal Studies"}
\section{Critical Legal Studies}
Herdando o realismo jurídico e o marxismo: o direito é fundamentalmente indeterminado, e sua aparência de neutralidade e certeza mascara os interesses dos poderosos e escolhas políticas.
\prov{relações: parte\_de filosofia\_do\_direito; usa realismo\_juridico; usa biopoder; relaciona injustica\_epistemica}
\prov{proveniência: wiki \textperiodcentered{} inferencia \textperiodcentered{} confiança 1.0 \textperiodcentered{} domínio ciencias\_de\_fora \textperiodcentered{} história 5 versões no jornal}

%CRISTAL {"arestas":[{"alvo":"artes_visuais","confianca":1.0,"epistemico":"fato","origem":"wiki","rel":"parte_de"},{"alvo":"perspectiva","confianca":1.0,"epistemico":"fato","origem":"wiki","rel":"contradiz"}],"confianca":1.0,"contraexemplos":[],"descricao":"Picasso/Braque: mostrar o objeto de vários ângulos ao mesmo tempo, quebrando a perspectiva única — o espaço fragmentado.","epistemico":"inferencia","exemplos":[],"id":"cubismo","memoria":"semantica","meta":{"autor":"Picasso","dominio":"ciencias_de_fora","fonte":""},"origem":"wiki","palavras_chave":[],"sinonimos":[],"tipo":"conceito","titulo":"Cubismo"}
\section{Cubismo}
Picasso/Braque: mostrar o objeto de vários ângulos ao mesmo tempo, quebrando a perspectiva única — o espaço fragmentado.
\prov{relações: parte\_de artes\_visuais; contradiz perspectiva}
\prov{proveniência: wiki \textperiodcentered{} inferencia \textperiodcentered{} confiança 1.0 \textperiodcentered{} domínio ciencias\_de\_fora \textperiodcentered{} história 3 versões no jornal}

%CRISTAL {"arestas":[{"alvo":"fim_de_vida_espectro","confianca":1.0,"epistemico":"fato","origem":"wiki","rel":"relaciona"},{"alvo":"ethics_of_care","confianca":1.0,"epistemico":"fato","origem":"wiki","rel":"usa"},{"alvo":"face_bela_morte","confianca":1.0,"epistemico":"fato","origem":"wiki","rel":"relaciona"}],"confianca":1.0,"contraexemplos":[],"descricao":"Cicely Saunders/OMS: a abordagem que prioriza o controle de sintomas e a qualidade de vida na doença grave sem visar à cura; a sedação paliativa reduz a consciência quando é o único meio de aliviar sofrimento intratável, sem a intenção de matar.","epistemico":"inferencia","exemplos":[],"id":"cuidados_paliativos","memoria":"semantica","meta":{"autor":"Saunders","dominio":"ciencias_de_fora","fonte":""},"origem":"wiki","palavras_chave":[],"sinonimos":[],"tipo":"conceito","titulo":"Cuidados Paliativos"}
\section{Cuidados Paliativos}
Cicely Saunders/OMS: a abordagem que prioriza o controle de sintomas e a qualidade de vida na doença grave sem visar à cura; a sedação paliativa reduz a consciência quando é o único meio de aliviar sofrimento intratável, sem a intenção de matar.
\prov{relações: relaciona fim\_de\_vida\_espectro; usa ethics\_of\_care; relaciona face\_bela\_morte}
\prov{proveniência: wiki \textperiodcentered{} inferencia \textperiodcentered{} confiança 1.0 \textperiodcentered{} domínio ciencias\_de\_fora \textperiodcentered{} história 4 versões no jornal}

%CRISTAL {"arestas":[{"alvo":"comunicacao_digital","confianca":1.0,"epistemico":"fato","origem":"wiki","rel":"parte_de"},{"alvo":"cultura_participativa","confianca":1.0,"epistemico":"fato","origem":"wiki","rel":"usa"}],"confianca":1.0,"contraexemplos":[],"descricao":"Jenkins: o fluxo de conteúdos por múltiplas plataformas e a migração do público que busca ativamente as experiências desejadas; um processo cultural (não só tecnológico) que articula convergência midiática, cultura participativa e narrativa transmídia.","epistemico":"inferencia","exemplos":[],"id":"cultura_da_convergencia","memoria":"semantica","meta":{"autor":"Jenkins","dominio":"ciencias_de_fora","fonte":""},"origem":"wiki","palavras_chave":[],"sinonimos":[],"tipo":"conceito","titulo":"Cultura da Convergência (Jenkins)"}
\section{Cultura da Convergência (Jenkins)}
Jenkins: o fluxo de conteúdos por múltiplas plataformas e a migração do público que busca ativamente as experiências desejadas; um processo cultural (não só tecnológico) que articula convergência midiática, cultura participativa e narrativa transmídia.
\prov{relações: parte\_de comunicacao\_digital; usa cultura\_participativa}
\prov{proveniência: wiki \textperiodcentered{} inferencia \textperiodcentered{} confiança 1.0 \textperiodcentered{} domínio ciencias\_de\_fora \textperiodcentered{} história 3 versões no jornal}

%CRISTAL {"arestas":[{"alvo":"estudos_de_midia","confianca":1.0,"epistemico":"fato","origem":"wiki","rel":"parte_de"},{"alvo":"hegemonia_cultural","confianca":1.0,"epistemico":"fato","origem":"wiki","rel":"relaciona"}],"confianca":1.0,"contraexemplos":[],"descricao":"Raymond Williams: 'cultura é ordinária' — não o patrimônio das elites, mas o modo de vida integral produzido no cotidiano das pessoas comuns; o materialismo cultural a analisa como prática material inseparável das relações sociais.","epistemico":"inferencia","exemplos":[],"id":"cultura_modo_de_vida","memoria":"semantica","meta":{"autor":"Williams","dominio":"ciencias_de_fora","fonte":""},"origem":"wiki","palavras_chave":[],"sinonimos":[],"tipo":"conceito","titulo":"Cultura como Modo de Vida (Williams)"}
\section{Cultura como Modo de Vida (Williams)}
Raymond Williams: 'cultura é ordinária' — não o patrimônio das elites, mas o modo de vida integral produzido no cotidiano das pessoas comuns; o materialismo cultural a analisa como prática material inseparável das relações sociais.
\prov{relações: parte\_de estudos\_de\_midia; relaciona hegemonia\_cultural}
\prov{proveniência: wiki \textperiodcentered{} inferencia \textperiodcentered{} confiança 1.0 \textperiodcentered{} domínio ciencias\_de\_fora \textperiodcentered{} história 3 versões no jornal}

%CRISTAL {"arestas":[{"alvo":"comunicacao_digital","confianca":1.0,"epistemico":"fato","origem":"wiki","rel":"parte_de"},{"alvo":"inteligencia_coletiva","confianca":1.0,"epistemico":"fato","origem":"wiki","rel":"usa"}],"confianca":1.0,"contraexemplos":[],"descricao":"Jenkins: a fronteira entre produtores e consumidores se dissolve — fãs e usuários participam da criação e circulação de conteúdo, com baixas barreiras à expressão e mentoria social.","epistemico":"inferencia","exemplos":[],"id":"cultura_participativa","memoria":"semantica","meta":{"autor":"Jenkins","dominio":"ciencias_de_fora","fonte":""},"origem":"wiki","palavras_chave":[],"sinonimos":[],"tipo":"conceito","titulo":"Cultura Participativa (Jenkins)"}
\section{Cultura Participativa (Jenkins)}
Jenkins: a fronteira entre produtores e consumidores se dissolve — fãs e usuários participam da criação e circulação de conteúdo, com baixas barreiras à expressão e mentoria social.
\prov{relações: parte\_de comunicacao\_digital; usa inteligencia\_coletiva}
\prov{proveniência: wiki \textperiodcentered{} inferencia \textperiodcentered{} confiança 1.0 \textperiodcentered{} domínio ciencias\_de\_fora \textperiodcentered{} história 3 versões no jornal}

%CRISTAL {"arestas":[{"alvo":"aprendizagem_descoberta_bruner","confianca":1.0,"epistemico":"fato","origem":"wiki","rel":"usa"},{"alvo":"word","confianca":0.5,"epistemico":"fato","origem":"wiki","rel":"relaciona"},{"alvo":"syscall","confianca":0.5,"epistemico":"fato","origem":"wiki","rel":"relaciona"},{"alvo":"virtualizacao","confianca":0.5,"epistemico":"fato","origem":"wiki","rel":"relaciona"}],"confianca":1.0,"contraexemplos":[],"descricao":"Bruner: qualquer conteúdo pode ser ensinado de forma honesta a qualquer idade, revisitando os conceitos-núcleo com complexidade e profundidade crescentes.","epistemico":"inferencia","exemplos":[],"id":"curriculo_espiral","memoria":"semantica","meta":{"autor":"Bruner","dominio":"ciencias_de_fora","fonte":""},"origem":"wiki","palavras_chave":[],"sinonimos":[],"tipo":"conceito","titulo":"Currículo em Espiral"}
\section{Currículo em Espiral}
Bruner: qualquer conteúdo pode ser ensinado de forma honesta a qualquer idade, revisitando os conceitos-núcleo com complexidade e profundidade crescentes.
\prov{relações: usa aprendizagem\_descoberta\_bruner; relaciona word; relaciona syscall; relaciona virtualizacao}
\prov{proveniência: wiki \textperiodcentered{} inferencia \textperiodcentered{} confiança 1.0 \textperiodcentered{} domínio ciencias\_de\_fora \textperiodcentered{} história 5 versões no jornal}

%CRISTAL {"arestas":[{"alvo":"antropologia","confianca":1.0,"epistemico":"fato","origem":"wiki","rel":"parte_de"},{"alvo":"contrato_social","confianca":1.0,"epistemico":"fato","origem":"wiki","rel":"relaciona"},{"alvo":"vontade_de_poder","confianca":0.5,"epistemico":"fato","origem":"wiki","rel":"relaciona"}],"confianca":1.0,"contraexemplos":[],"descricao":"Mauss: a troca de presentes é livre na aparência e obrigatória na estrutura (dar-receber-retribuir); cria vínculo — um 'fato social total'.","epistemico":"inferencia","exemplos":[],"id":"dadiva_mauss","memoria":"semantica","meta":{"autor":"Mauss","dominio":"ciencias_de_fora","fonte":""},"origem":"wiki","palavras_chave":[],"sinonimos":[],"tipo":"conceito","titulo":"A Dádiva"}
\section{A Dádiva}
Mauss: a troca de presentes é livre na aparência e obrigatória na estrutura (dar-receber-retribuir); cria vínculo — um 'fato social total'.
\prov{relações: parte\_de antropologia; relaciona contrato\_social; relaciona vontade\_de\_poder}
\prov{proveniência: wiki \textperiodcentered{} inferencia \textperiodcentered{} confiança 1.0 \textperiodcentered{} domínio ciencias\_de\_fora \textperiodcentered{} história 4 versões no jornal}

%CRISTAL {"arestas":[{"alvo":"selecao_natural","confianca":1.0,"epistemico":"fato","origem":"wiki","rel":"especializa"},{"alvo":"maquina_evolutiva","confianca":1.0,"epistemico":"fato","origem":"wiki","rel":"explica"},{"alvo":"word","confianca":0.5,"epistemico":"fato","origem":"wiki","rel":"relaciona"}],"confianca":1.0,"contraexemplos":[],"descricao":"Algoritmos genéticos e programação genética implementam a seleção natural em silício — população, mutação, crossover, fitness.","epistemico":"fato","exemplos":[],"id":"darwinismo_algoritmico","memoria":"semantica","meta":{"autor":"Holland","dominio":"ciencias_de_fora","fonte":""},"origem":"wiki","palavras_chave":[],"sinonimos":[],"tipo":"conceito","titulo":"Darwinismo Algorítmico"}
\section{Darwinismo Algorítmico}
Algoritmos genéticos e programação genética implementam a seleção natural em silício — população, mutação, crossover, fitness.
\prov{relações: especializa selecao\_natural; explica maquina\_evolutiva; relaciona word}
\prov{proveniência: wiki \textperiodcentered{} fato \textperiodcentered{} confiança 1.0 \textperiodcentered{} domínio ciencias\_de\_fora \textperiodcentered{} história 4 versões no jornal}

%CRISTAL {"arestas":[{"alvo":"comunicacao_digital","confianca":1.0,"epistemico":"fato","origem":"wiki","rel":"parte_de"},{"alvo":"biopoder","confianca":1.0,"epistemico":"fato","origem":"wiki","rel":"relaciona"}],"confianca":1.0,"contraexemplos":[],"descricao":"Mayer-Schönberger & Cukier: transformar em dados quantificáveis fenômenos jamais tratados como informação (localização, emoções) — 'datificar é pôr em forma quantificada para tabular e analisar', tratando a realidade como repositório de dados.","epistemico":"inferencia","exemplos":[],"id":"datificacao","memoria":"semantica","meta":{"autor":"Mayer-Schönberger & Cukier","dominio":"ciencias_de_fora","fonte":""},"origem":"wiki","palavras_chave":[],"sinonimos":[],"tipo":"conceito","titulo":"Datificação / Big Data"}
\section{Datificação / Big Data}
Mayer-Schönberger \& Cukier: transformar em dados quantificáveis fenômenos jamais tratados como informação (localização, emoções) — 'datificar é pôr em forma quantificada para tabular e analisar', tratando a realidade como repositório de dados.
\prov{relações: parte\_de comunicacao\_digital; relaciona biopoder}
\prov{proveniência: wiki \textperiodcentered{} inferencia \textperiodcentered{} confiança 1.0 \textperiodcentered{} domínio ciencias\_de\_fora \textperiodcentered{} história 3 versões no jornal}

%CRISTAL {"arestas":[{"alvo":"autoridade_e_rejeicao_gentil","confianca":1.0,"epistemico":"fato","origem":"wiki","rel":"usa"},{"alvo":"teoria_das_elites","confianca":1.0,"epistemico":"fato","origem":"wiki","rel":"contradiz"},{"alvo":"morte_de_geracoes","confianca":1.0,"epistemico":"fato","origem":"wiki","rel":"relaciona"}],"confianca":1.0,"contraexemplos":[],"descricao":"Metáfora-título recorrente do Lopes: um 'Davi' solitário que derrota 'todos os Golias da matemática' — o iconoclasta individual contra o poder coletivo e institucional consagrado; reivindica atenção aos 'compositores' (criadores), não só aos 'intérpretes'.","epistemico":"inferencia","exemplos":[],"id":"david_contra_golias","memoria":"semantica","meta":{"autor":"Gentil Lopes","dominio":"ciencias_de_fora","fonte":""},"origem":"wiki","palavras_chave":[],"sinonimos":[],"tipo":"conceito","titulo":"David contra Golias (Lopes)"}
\section{David contra Golias (Lopes)}
Metáfora-título recorrente do Lopes: um 'Davi' solitário que derrota 'todos os Golias da matemática' — o iconoclasta individual contra o poder coletivo e institucional consagrado; reivindica atenção aos 'compositores' (criadores), não só aos 'intérpretes'.
\prov{relações: usa autoridade\_e\_rejeicao\_gentil; contradiz teoria\_das\_elites; relaciona morte\_de\_geracoes}
\prov{proveniência: wiki \textperiodcentered{} inferencia \textperiodcentered{} confiança 1.0 \textperiodcentered{} domínio ciencias\_de\_fora \textperiodcentered{} história 5 versões no jornal}

%CRISTAL {"arestas":[{"alvo":"filosofia_do_direito","confianca":1.0,"epistemico":"fato","origem":"wiki","rel":"parte_de"},{"alvo":"liberdade_berlin","confianca":1.0,"epistemico":"fato","origem":"wiki","rel":"relaciona"}],"confianca":1.0,"contraexemplos":[],"descricao":"Devlin: a sociedade pode legislar contra atos que provoquem 'intolerância, indignação e repugnância' coletivas; Hart, na esteira de Mill, nega — o Estado só pode restringir a liberdade para evitar dano a terceiros (princípio do dano).","epistemico":"inferencia","exemplos":[],"id":"debate_direito_moral","memoria":"semantica","meta":{"autor":"Hart vs Devlin","dominio":"ciencias_de_fora","fonte":""},"origem":"wiki","palavras_chave":[],"sinonimos":[],"tipo":"conceito","titulo":"O Direito deve Impor Moral? (Hart-Devlin)"}
\section{O Direito deve Impor Moral? (Hart-Devlin)}
Devlin: a sociedade pode legislar contra atos que provoquem 'intolerância, indignação e repugnância' coletivas; Hart, na esteira de Mill, nega — o Estado só pode restringir a liberdade para evitar dano a terceiros (princípio do dano).
\prov{relações: parte\_de filosofia\_do\_direito; relaciona liberdade\_berlin}
\prov{proveniência: wiki \textperiodcentered{} inferencia \textperiodcentered{} confiança 1.0 \textperiodcentered{} domínio ciencias\_de\_fora \textperiodcentered{} história 3 versões no jornal}

%CRISTAL {"arestas":[{"alvo":"ciclos_vico","confianca":1.0,"epistemico":"fato","origem":"wiki","rel":"especializa"},{"alvo":"word","confianca":0.5,"epistemico":"fato","origem":"wiki","rel":"relaciona"}],"confianca":1.0,"contraexemplos":[],"descricao":"Spengler: cada civilização é um organismo com estações (primavera→inverno) e morre inexoravelmente; culturas não progridem. Analogia especulativa.","epistemico":"hipotese","exemplos":[],"id":"declinio_spengler","memoria":"semantica","meta":{"autor":"Spengler","dominio":"ciencias_de_fora","fonte":""},"origem":"wiki","palavras_chave":[],"sinonimos":[],"tipo":"conceito","titulo":"O Declínio do Ocidente"}
\section{O Declínio do Ocidente}
Spengler: cada civilização é um organismo com estações (primavera\ensuremath{\to}inverno) e morre inexoravelmente; culturas não progridem. Analogia especulativa.
\prov{relações: especializa ciclos\_vico; relaciona word}
\prov{proveniência: wiki \textperiodcentered{} hipotese \textperiodcentered{} confiança 1.0 \textperiodcentered{} domínio ciencias\_de\_fora \textperiodcentered{} história 3 versões no jornal}

%CRISTAL {"arestas":[{"alvo":"colapso_observador","confianca":1.0,"epistemico":"fato","origem":"wiki","rel":"contradiz"},{"alvo":"colapso","confianca":1.0,"epistemico":"fato","origem":"wiki","rel":"referencia"}],"confianca":1.0,"contraexemplos":[],"descricao":"Zeh: superposições macroscópicas degradam-se por acoplamento ambiental, dando a APARÊNCIA de colapso sem colapso fundamental. Dispensa o observador consciente.","epistemico":"fato","exemplos":[],"id":"decoerencia","memoria":"semantica","meta":{"autor":"Zeh","dominio":"filosofia","fonte":""},"origem":"wiki","palavras_chave":[],"sinonimos":[],"tipo":"conceito","titulo":"Decoerência Quântica"}
\section{Decoerência Quântica}
Zeh: superposições macroscópicas degradam-se por acoplamento ambiental, dando a APARÊNCIA de colapso sem colapso fundamental. Dispensa o observador consciente.
\prov{relações: contradiz colapso\_observador; referencia colapso}
\prov{proveniência: wiki \textperiodcentered{} fato \textperiodcentered{} confiança 1.0 \textperiodcentered{} domínio filosofia \textperiodcentered{} história 3 versões no jornal}

%CRISTAL {"arestas":[{"alvo":"pragmatica","confianca":1.0,"epistemico":"fato","origem":"wiki","rel":"parte_de"}],"confianca":1.0,"contraexemplos":[],"descricao":"Palavras que só significam pelo contexto de enunciação — 'eu', 'aqui', 'agora', 'isto'. O signo que aponta.","epistemico":"fato","exemplos":[],"id":"deixis","memoria":"semantica","meta":{"dominio":"ciencias_de_fora","fonte":""},"origem":"wiki","palavras_chave":[],"sinonimos":[],"tipo":"conceito","titulo":"Dêixis"}
\section{Dêixis}
Palavras que só significam pelo contexto de enunciação — 'eu', 'aqui', 'agora', 'isto'. O signo que aponta.
\prov{relações: parte\_de pragmatica}
\prov{proveniência: wiki \textperiodcentered{} fato \textperiodcentered{} confiança 1.0 \textperiodcentered{} domínio ciencias\_de\_fora \textperiodcentered{} história 2 versões no jornal}

%CRISTAL {"arestas":[{"alvo":"economia","confianca":1.0,"epistemico":"fato","origem":"wiki","rel":"parte_de"},{"alvo":"virtualizacao","confianca":0.5,"epistemico":"fato","origem":"wiki","rel":"relaciona"},{"alvo":"word","confianca":0.5,"epistemico":"fato","origem":"wiki","rel":"relaciona"},{"alvo":"experimento_1_trocar_a_prata","confianca":0.5,"epistemico":"fato","origem":"wiki","rel":"relaciona"}],"confianca":1.0,"contraexemplos":[],"descricao":"Keynes: a economia não alcança o pleno emprego sozinha; o Estado deve intervir com gasto contracíclico para sustentar a demanda.","epistemico":"inferencia","exemplos":[],"id":"demanda_agregada","memoria":"semantica","meta":{"autor":"Keynes","dominio":"ciencias_de_fora","fonte":""},"origem":"wiki","palavras_chave":[],"sinonimos":[],"tipo":"conceito","titulo":"Demanda Agregada"}
\section{Demanda Agregada}
Keynes: a economia não alcança o pleno emprego sozinha; o Estado deve intervir com gasto contracíclico para sustentar a demanda.
\prov{relações: parte\_de economia; relaciona virtualizacao; relaciona word; relaciona experimento\_1\_trocar\_a\_prata}
\prov{proveniência: wiki \textperiodcentered{} inferencia \textperiodcentered{} confiança 1.0 \textperiodcentered{} domínio ciencias\_de\_fora \textperiodcentered{} história 5 versões no jornal}

%CRISTAL {"arestas":[{"alvo":"democracia_formas","confianca":1.0,"epistemico":"fato","origem":"wiki","rel":"especializa"},{"alvo":"habermas_direito_democracia","confianca":1.0,"epistemico":"fato","origem":"wiki","rel":"usa"},{"alvo":"refeudalizacao_esfera_publica","confianca":1.0,"epistemico":"fato","origem":"wiki","rel":"relaciona"}],"confianca":1.0,"contraexemplos":[],"descricao":"Habermas: a legitimidade das decisões deriva não do mero voto majoritário, mas da deliberação pública entre cidadãos livres e iguais, na qual prevalece 'a força do melhor argumento'.","epistemico":"inferencia","exemplos":[],"id":"democracia_deliberativa","memoria":"semantica","meta":{"autor":"Habermas","dominio":"ciencias_de_fora","fonte":""},"origem":"wiki","palavras_chave":[],"sinonimos":[],"tipo":"conceito","titulo":"Democracia Deliberativa (Habermas)"}
\section{Democracia Deliberativa (Habermas)}
Habermas: a legitimidade das decisões deriva não do mero voto majoritário, mas da deliberação pública entre cidadãos livres e iguais, na qual prevalece 'a força do melhor argumento'.
\prov{relações: especializa democracia\_formas; usa habermas\_direito\_democracia; relaciona refeudalizacao\_esfera\_publica}
\prov{proveniência: wiki \textperiodcentered{} inferencia \textperiodcentered{} confiança 1.0 \textperiodcentered{} domínio ciencias\_de\_fora \textperiodcentered{} história 4 versões no jornal}

%CRISTAL {"arestas":[{"alvo":"ciencia_politica","confianca":1.0,"epistemico":"fato","origem":"wiki","rel":"parte_de"},{"alvo":"contrato_social","confianca":1.0,"epistemico":"fato","origem":"wiki","rel":"usa"}],"confianca":1.0,"contraexemplos":[],"descricao":"A democracia direta (o modelo ateniense da assembleia, sem intermediários) vs a representativa (o poder exercido por representantes eleitos e responsabilizáveis em eleições periódicas). O poder do povo, exercido pessoal ou indiretamente.","epistemico":"inferencia","exemplos":[],"id":"democracia_formas","memoria":"semantica","meta":{"autor":"—","dominio":"ciencias_de_fora","fonte":""},"origem":"wiki","palavras_chave":[],"sinonimos":[],"tipo":"conceito","titulo":"Formas de Democracia (Direta e Representativa)"}
\section{Formas de Democracia (Direta e Representativa)}
A democracia direta (o modelo ateniense da assembleia, sem intermediários) vs a representativa (o poder exercido por representantes eleitos e responsabilizáveis em eleições periódicas). O poder do povo, exercido pessoal ou indiretamente.
\prov{relações: parte\_de ciencia\_politica; usa contrato\_social}
\prov{proveniência: wiki \textperiodcentered{} inferencia \textperiodcentered{} confiança 1.0 \textperiodcentered{} domínio ciencias\_de\_fora \textperiodcentered{} história 3 versões no jornal}

%CRISTAL {"arestas":[{"alvo":"signo_linguistico","confianca":1.0,"epistemico":"fato","origem":"wiki","rel":"usa"}],"confianca":1.0,"contraexemplos":[],"descricao":"Dois níveis de sentido: o literal (denotação) e o cultural/afetivo que se sobrepõe a ele (conotação).","epistemico":"fato","exemplos":[],"id":"denotacao_conotacao","memoria":"semantica","meta":{"autor":"Barthes","dominio":"ciencias_de_fora","fonte":""},"origem":"wiki","palavras_chave":[],"sinonimos":[],"tipo":"conceito","titulo":"Denotação e Conotação"}
\section{Denotação e Conotação}
Dois níveis de sentido: o literal (denotação) e o cultural/afetivo que se sobrepõe a ele (conotação).
\prov{relações: usa signo\_linguistico}
\prov{proveniência: wiki \textperiodcentered{} fato \textperiodcentered{} confiança 1.0 \textperiodcentered{} domínio ciencias\_de\_fora \textperiodcentered{} história 2 versões no jornal}

%CRISTAL {"arestas":[{"alvo":"tectonica_placas","confianca":1.0,"epistemico":"fato","origem":"wiki","rel":"usa"}],"confianca":1.0,"contraexemplos":[],"descricao":"Wegener (1912): os continentes formaram um supercontinente (Pangeia) e desde então se deslocam. Rejeitada em vida por falta de mecanismo, depois incorporada à tectônica de placas — exemplo de hipótese refutada e reabilitada.","epistemico":"fato","exemplos":[],"id":"deriva_continental","memoria":"semantica","meta":{"autor":"Wegener","dominio":"ciencias_de_fora","fonte":""},"origem":"wiki","palavras_chave":[],"sinonimos":[],"tipo":"conceito","titulo":"Deriva Continental (Wegener)"}
\section{Deriva Continental (Wegener)}
Wegener (1912): os continentes formaram um supercontinente (Pangeia) e desde então se deslocam. Rejeitada em vida por falta de mecanismo, depois incorporada à tectônica de placas — exemplo de hipótese refutada e reabilitada.
\prov{relações: usa tectonica\_placas}
\prov{proveniência: wiki \textperiodcentered{} fato \textperiodcentered{} confiança 1.0 \textperiodcentered{} domínio ciencias\_de\_fora \textperiodcentered{} história 2 versões no jornal}

%CRISTAL {"arestas":[{"alvo":"biologia_evolutiva","confianca":1.0,"epistemico":"fato","origem":"wiki","rel":"parte_de"},{"alvo":"selecao_natural","confianca":1.0,"epistemico":"fato","origem":"wiki","rel":"contradiz"}],"confianca":1.0,"contraexemplos":[],"descricao":"A mudança das frequências gênicas por acaso amostral — evolução sem seleção, forte em populações pequenas.","epistemico":"fato","exemplos":[],"id":"deriva_genetica","memoria":"semantica","meta":{"dominio":"ciencias_de_fora","fonte":""},"origem":"wiki","palavras_chave":[],"sinonimos":[],"tipo":"conceito","titulo":"Deriva Genética"}
\section{Deriva Genética}
A mudança das frequências gênicas por acaso amostral — evolução sem seleção, forte em populações pequenas.
\prov{relações: parte\_de biologia\_evolutiva; contradiz selecao\_natural}
\prov{proveniência: wiki \textperiodcentered{} fato \textperiodcentered{} confiança 1.0 \textperiodcentered{} domínio ciencias\_de\_fora \textperiodcentered{} história 3 versões no jornal}

%CRISTAL {"arestas":[{"alvo":"declinio_spengler","confianca":1.0,"epistemico":"fato","origem":"wiki","rel":"contradiz"},{"alvo":"contingencia_historica","confianca":1.0,"epistemico":"fato","origem":"wiki","rel":"usa"}],"confianca":1.0,"contraexemplos":[],"descricao":"Toynbee: civilizações crescem respondendo a desafios; minorias criativas inovam. A resposta é indeterminada — abre a contingência.","epistemico":"inferencia","exemplos":[],"id":"desafio_resposta","memoria":"semantica","meta":{"autor":"Toynbee","dominio":"ciencias_de_fora","fonte":""},"origem":"wiki","palavras_chave":[],"sinonimos":[],"tipo":"conceito","titulo":"Desafio-Resposta (Toynbee)"}
\section{Desafio-Resposta (Toynbee)}
Toynbee: civilizações crescem respondendo a desafios; minorias criativas inovam. A resposta é indeterminada — abre a contingência.
\prov{relações: contradiz declinio\_spengler; usa contingencia\_historica}
\prov{proveniência: wiki \textperiodcentered{} inferencia \textperiodcentered{} confiança 1.0 \textperiodcentered{} domínio ciencias\_de\_fora \textperiodcentered{} história 3 versões no jornal}

%CRISTAL {"arestas":[{"alvo":"traducao","confianca":1.0,"epistemico":"fato","origem":"wiki","rel":"usa"},{"alvo":"semantica","confianca":1.0,"epistemico":"fato","origem":"wiki","rel":"usa"}],"confianca":1.0,"contraexemplos":[],"descricao":"Escolher qual sentido de uma palavra polissêmica traduzir — 'banco' (assento/instituição/margem). Só o CONTEXTO decide; o núcleo da dificuldade da tradução automática.","epistemico":"fato","exemplos":[],"id":"desambiguacao","memoria":"semantica","meta":{"dominio":"ciencias_de_fora","fonte":""},"origem":"wiki","palavras_chave":[],"sinonimos":[],"tipo":"conceito","titulo":"Desambiguação (pelo Contexto)"}
\section{Desambiguação (pelo Contexto)}
Escolher qual sentido de uma palavra polissêmica traduzir — 'banco' (assento/instituição/margem). Só o CONTEXTO decide; o núcleo da dificuldade da tradução automática.
\prov{relações: usa traducao; usa semantica}
\prov{proveniência: wiki \textperiodcentered{} fato \textperiodcentered{} confiança 1.0 \textperiodcentered{} domínio ciencias\_de\_fora \textperiodcentered{} história 3 versões no jornal}

%CRISTAL {"arestas":[{"alvo":"teoria_da_evolucao","confianca":1.0,"epistemico":"fato","origem":"wiki","rel":"parte_de"}],"confianca":1.0,"contraexemplos":[],"descricao":"Darwin: toda a vida partilha ancestrais — a árvore única, evidenciada por fósseis, anatomia e DNA.","epistemico":"fato","exemplos":[],"id":"descendencia_comum","memoria":"semantica","meta":{"autor":"Darwin","dominio":"ciencias_de_fora","fonte":""},"origem":"wiki","palavras_chave":[],"sinonimos":[],"tipo":"conceito","titulo":"Descendência Comum"}
\section{Descendência Comum}
Darwin: toda a vida partilha ancestrais — a árvore única, evidenciada por fósseis, anatomia e DNA.
\prov{relações: parte\_de teoria\_da\_evolucao}
\prov{proveniência: wiki \textperiodcentered{} fato \textperiodcentered{} confiança 1.0 \textperiodcentered{} domínio ciencias\_de\_fora \textperiodcentered{} história 2 versões no jornal}

%CRISTAL {"arestas":[{"alvo":"filosofia_da_educacao","confianca":1.0,"epistemico":"fato","origem":"wiki","rel":"parte_de"},{"alvo":"biopoder","confianca":1.0,"epistemico":"fato","origem":"wiki","rel":"relaciona"},{"alvo":"didatica_magna_comenius","confianca":1.0,"epistemico":"fato","origem":"wiki","rel":"contradiz"}],"confianca":1.0,"contraexemplos":[],"descricao":"Illich: a escola monopoliza a aprendizagem e ensina a confundir ensino com aprendizado, diploma com competência, frequência com educação; propõe devolver a aprendizagem a redes comunitárias livres.","epistemico":"inferencia","exemplos":[],"id":"desescolarizacao_illich","memoria":"semantica","meta":{"autor":"Illich","dominio":"ciencias_de_fora","fonte":""},"origem":"wiki","palavras_chave":[],"sinonimos":[],"tipo":"conceito","titulo":"Desescolarização (Illich)"}
\section{Desescolarização (Illich)}
Illich: a escola monopoliza a aprendizagem e ensina a confundir ensino com aprendizado, diploma com competência, frequência com educação; propõe devolver a aprendizagem a redes comunitárias livres.
\prov{relações: parte\_de filosofia\_da\_educacao; relaciona biopoder; contradiz didatica\_magna\_comenius}
\prov{proveniência: wiki \textperiodcentered{} inferencia \textperiodcentered{} confiança 1.0 \textperiodcentered{} domínio ciencias\_de\_fora \textperiodcentered{} história 4 versões no jornal}

%CRISTAL {"arestas":[{"alvo":"frege_sentido_referencia","confianca":1.0,"epistemico":"fato","origem":"wiki","rel":"contradiz"},{"alvo":"semiotica_peirce_triade","confianca":1.0,"epistemico":"fato","origem":"wiki","rel":"relaciona"}],"confianca":1.0,"contraexemplos":[],"descricao":"Kripke: nomes próprios designam o mesmo indivíduo em todo mundo possível, por uma cadeia causal — não por descrição (contra Frege/Russell).","epistemico":"inferencia","exemplos":[],"id":"designadores_rigidos","memoria":"semantica","meta":{"autor":"Kripke","dominio":"ciencias_de_fora","fonte":""},"origem":"wiki","palavras_chave":[],"sinonimos":[],"tipo":"conceito","titulo":"Designadores Rígidos"}
\section{Designadores Rígidos}
Kripke: nomes próprios designam o mesmo indivíduo em todo mundo possível, por uma cadeia causal — não por descrição (contra Frege/Russell).
\prov{relações: contradiz frege\_sentido\_referencia; relaciona semiotica\_peirce\_triade}
\prov{proveniência: wiki \textperiodcentered{} inferencia \textperiodcentered{} confiança 1.0 \textperiodcentered{} domínio ciencias\_de\_fora \textperiodcentered{} história 3 versões no jornal}

%CRISTAL {"arestas":[{"alvo":"direito","confianca":1.0,"epistemico":"fato","origem":"wiki","rel":"parte_de"},{"alvo":"honestidade_intelectual","confianca":1.0,"epistemico":"fato","origem":"wiki","rel":"relaciona"},{"alvo":"estado_de_direito","confianca":1.0,"epistemico":"fato","origem":"wiki","rel":"relaciona"}],"confianca":1.0,"contraexemplos":[],"descricao":"Thoreau/Rawls/King: violação pública, consciente, não-violenta e assumida de uma lei tida por injusta, como apelo ao senso de justiça da comunidade, aceitando as consequências legais.","epistemico":"inferencia","exemplos":[],"id":"desobediencia_civil","memoria":"semantica","meta":{"autor":"Thoreau/King","dominio":"ciencias_de_fora","fonte":""},"origem":"wiki","palavras_chave":[],"sinonimos":[],"tipo":"conceito","titulo":"Desobediência Civil"}
\section{Desobediência Civil}
Thoreau/Rawls/King: violação pública, consciente, não-violenta e assumida de uma lei tida por injusta, como apelo ao senso de justiça da comunidade, aceitando as consequências legais.
\prov{relações: parte\_de direito; relaciona honestidade\_intelectual; relaciona estado\_de\_direito}
\prov{proveniência: wiki \textperiodcentered{} inferencia \textperiodcentered{} confiança 1.0 \textperiodcentered{} domínio ciencias\_de\_fora \textperiodcentered{} história 4 versões no jornal}

%CRISTAL {"arestas":[{"alvo":"saude_publica","confianca":1.0,"epistemico":"fato","origem":"wiki","rel":"parte_de"},{"alvo":"classe_social","confianca":1.0,"epistemico":"fato","origem":"wiki","rel":"usa"},{"alvo":"biopoder","confianca":1.0,"epistemico":"fato","origem":"wiki","rel":"relaciona"}],"confianca":1.0,"contraexemplos":[],"descricao":"Marmot: o estudo Whitehall mostrou que a mortalidade segue um gradiente social contínuo — quanto mais baixa a posição na hierarquia ocupacional, maior o risco de morte — e menos de um terço se explica por fatores de risco clássicos. A saúde é, em parte, justiça social.","epistemico":"fato","exemplos":[],"id":"determinantes_sociais_saude","memoria":"semantica","meta":{"autor":"Marmot","dominio":"ciencias_de_fora","fonte":""},"origem":"wiki","palavras_chave":[],"sinonimos":[],"tipo":"conceito","titulo":"Determinantes Sociais (Marmot/Whitehall)"}
\section{Determinantes Sociais (Marmot/Whitehall)}
Marmot: o estudo Whitehall mostrou que a mortalidade segue um gradiente social contínuo — quanto mais baixa a posição na hierarquia ocupacional, maior o risco de morte — e menos de um terço se explica por fatores de risco clássicos. A saúde é, em parte, justiça social.
\prov{relações: parte\_de saude\_publica; usa classe\_social; relaciona biopoder}
\prov{proveniência: wiki \textperiodcentered{} fato \textperiodcentered{} confiança 1.0 \textperiodcentered{} domínio ciencias\_de\_fora \textperiodcentered{} história 4 versões no jornal}

%CRISTAL {"arestas":[{"alvo":"geografia_humana","confianca":1.0,"epistemico":"fato","origem":"wiki","rel":"parte_de"},{"alvo":"longue_duree","confianca":1.0,"epistemico":"fato","origem":"wiki","rel":"relaciona"},{"alvo":"biopoder","confianca":1.0,"epistemico":"fato","origem":"wiki","rel":"relaciona"}],"confianca":1.0,"contraexemplos":[],"descricao":"Ratzel/Semple: o meio físico (clima, relevo, solo) determinaria causalmente a cultura e o destino histórico dos povos — tese hoje rejeitada por reducionista e por ter servido de base a racismos e ao expansionismo.","epistemico":"inferencia","exemplos":[],"id":"determinismo_ambiental","memoria":"semantica","meta":{"autor":"Ratzel/Semple","dominio":"ciencias_de_fora","fonte":""},"origem":"wiki","palavras_chave":[],"sinonimos":[],"tipo":"conceito","titulo":"Determinismo Ambiental"}
\section{Determinismo Ambiental}
Ratzel/Semple: o meio físico (clima, relevo, solo) determinaria causalmente a cultura e o destino histórico dos povos — tese hoje rejeitada por reducionista e por ter servido de base a racismos e ao expansionismo.
\prov{relações: parte\_de geografia\_humana; relaciona longue\_duree; relaciona biopoder}
\prov{proveniência: wiki \textperiodcentered{} inferencia \textperiodcentered{} confiança 1.0 \textperiodcentered{} domínio ciencias\_de\_fora \textperiodcentered{} história 4 versões no jornal}

%CRISTAL {"arestas":[{"alvo":"filosofia_da_tecnica","confianca":1.0,"epistemico":"fato","origem":"wiki","rel":"parte_de"},{"alvo":"ellul_autonomia_tecnica","confianca":1.0,"epistemico":"fato","origem":"wiki","rel":"relaciona"}],"confianca":1.0,"contraexemplos":[],"descricao":"A tese de que a tecnologia é uma força autônoma cujo desenvolvimento interno segue lógica própria e determina, por si, a estrutura e a mudança social. O polo 'a técnica molda a sociedade' — em geral apresentado como alvo de crítica.","epistemico":"inferencia","exemplos":[],"id":"determinismo_tecnologico","memoria":"semantica","meta":{"autor":"—","dominio":"ciencias_de_fora","fonte":""},"origem":"wiki","palavras_chave":[],"sinonimos":[],"tipo":"conceito","titulo":"Determinismo Tecnológico"}
\section{Determinismo Tecnológico}
A tese de que a tecnologia é uma força autônoma cujo desenvolvimento interno segue lógica própria e determina, por si, a estrutura e a mudança social. O polo 'a técnica molda a sociedade' — em geral apresentado como alvo de crítica.
\prov{relações: parte\_de filosofia\_da\_tecnica; relaciona ellul\_autonomia\_tecnica}
\prov{proveniência: wiki \textperiodcentered{} inferencia \textperiodcentered{} confiança 1.0 \textperiodcentered{} domínio ciencias\_de\_fora \textperiodcentered{} história 3 versões no jornal}

%CRISTAL {"arestas":[{"alvo":"cristianismo","confianca":1.0,"epistemico":"fato","origem":"wiki","rel":"parte_de"},{"alvo":"deus_postulado","confianca":1.0,"epistemico":"fato","origem":"wiki","rel":"relaciona"},{"alvo":"deus_mentiroso","confianca":1.0,"epistemico":"fato","origem":"wiki","rel":"contradiz"}],"confianca":1.0,"contraexemplos":[],"descricao":"Bíblia (Gênesis 1): Deus é a origem absoluta e livre de tudo, criando o cosmos 'do nada' — 'No princípio criou Deus os céus e a terra; e disse Deus: Haja luz; e houve luz'. O Criador afirmado como fato metafísico.","epistemico":"inferencia","exemplos":[],"id":"deus_criador","memoria":"semantica","meta":{"autor":"Bíblia","dominio":"sagrado","fonte":""},"origem":"wiki","palavras_chave":[],"sinonimos":[],"tipo":"conceito","titulo":"Deus Criador (creatio ex nihilo)"}
\section{Deus Criador (creatio ex nihilo)}
Bíblia (Gênesis 1): Deus é a origem absoluta e livre de tudo, criando o cosmos 'do nada' — 'No princípio criou Deus os céus e a terra; e disse Deus: Haja luz; e houve luz'. O Criador afirmado como fato metafísico.
\prov{relações: parte\_de cristianismo; relaciona deus\_postulado; contradiz deus\_mentiroso}
\prov{proveniência: wiki \textperiodcentered{} inferencia \textperiodcentered{} confiança 1.0 \textperiodcentered{} domínio sagrado \textperiodcentered{} história 4 versões no jornal}

%CRISTAL {"arestas":[{"alvo":"nao_dualidade_binaria_gentil","confianca":1.0,"epistemico":"fato","origem":"wiki","rel":"usa"},{"alvo":"deus_como_vazio_gentil","confianca":1.0,"epistemico":"fato","origem":"wiki","rel":"usa"},{"alvo":"queda_pecado","confianca":1.0,"epistemico":"fato","origem":"wiki","rel":"relaciona"}],"confianca":1.0,"contraexemplos":[],"descricao":"Lopes: bem e mal são duas faces da mesma 'eletricidade'/Absoluto, que é neutra; a dualidade só existe na mente/no 'hardware' (cérebro), não no Vazio. A mesma energia que age num Gandhi age num bandido — a culpa está no hardware, não na energia. Resolve o problema do mal sem um Deus pessoal.","epistemico":"inferencia","exemplos":[],"id":"deus_e_diabo_um_so_gentil","memoria":"semantica","meta":{"autor":"Gentil Lopes","dominio":"sagrado","fonte":""},"origem":"wiki","palavras_chave":[],"sinonimos":[],"tipo":"conceito","titulo":"Deus e o Diabo são a Mesma Energia (Lopes)"}
\section{Deus e o Diabo são a Mesma Energia (Lopes)}
Lopes: bem e mal são duas faces da mesma 'eletricidade'/Absoluto, que é neutra; a dualidade só existe na mente/no 'hardware' (cérebro), não no Vazio. A mesma energia que age num Gandhi age num bandido — a culpa está no hardware, não na energia. Resolve o problema do mal sem um Deus pessoal.
\prov{relações: usa nao\_dualidade\_binaria\_gentil; usa deus\_como\_vazio\_gentil; relaciona queda\_pecado}
\prov{proveniência: wiki \textperiodcentered{} inferencia \textperiodcentered{} confiança 1.0 \textperiodcentered{} domínio sagrado \textperiodcentered{} história 5 versões no jornal}

%CRISTAL {"arestas":[{"alvo":"deus_postulado","confianca":1.0,"epistemico":"fato","origem":"wiki","rel":"usa"},{"alvo":"word","confianca":0.5,"epistemico":"fato","origem":"wiki","rel":"relaciona"}],"confianca":1.0,"contraexemplos":[],"descricao":"Lopes: como milhares de conceitos de Deus se contradizem, e cada um foi ensinado como verdade absoluta, ao menos um mente — argumento lógico contra a fé programada.","epistemico":"inferencia","exemplos":[],"id":"deus_mentiroso","memoria":"semantica","meta":{"autor":"Gentil Lopes","dominio":"filosofia","fonte":""},"origem":"livro","palavras_chave":[],"sinonimos":[],"tipo":"conceito","titulo":"Deus Mentiroso (argumento)"}
\section{Deus Mentiroso (argumento)}
Lopes: como milhares de conceitos de Deus se contradizem, e cada um foi ensinado como verdade absoluta, ao menos um mente — argumento lógico contra a fé programada.
\prov{relações: usa deus\_postulado; relaciona word}
\prov{proveniência: livro \textperiodcentered{} inferencia \textperiodcentered{} confiança 1.0 \textperiodcentered{} domínio filosofia \textperiodcentered{} história 4 versões no jornal}

%CRISTAL {"arestas":[{"alvo":"metafisica","confianca":1.0,"epistemico":"fato","origem":"wiki","rel":"parte_de"},{"alvo":"epistemologia","confianca":1.0,"epistemico":"fato","origem":"wiki","rel":"usa"},{"alvo":"word","confianca":0.5,"epistemico":"fato","origem":"wiki","rel":"relaciona"},{"alvo":"transcricao","confianca":0.5,"epistemico":"fato","origem":"wiki","rel":"relaciona"}],"confianca":1.0,"contraexemplos":[],"descricao":"Lopes: 'Deus' é um axioma, não uma verdade — não é nem verdadeiro nem falso; o ser inteligente examina o postulado antes de aderir.","epistemico":"inferencia","exemplos":[],"id":"deus_postulado","memoria":"semantica","meta":{"autor":"Gentil Lopes","dominio":"filosofia","fonte":""},"origem":"livro","palavras_chave":[],"sinonimos":[],"tipo":"conceito","titulo":"Deus como Postulado"}
\section{Deus como Postulado}
Lopes: 'Deus' é um axioma, não uma verdade — não é nem verdadeiro nem falso; o ser inteligente examina o postulado antes de aderir.
\prov{relações: parte\_de metafisica; usa epistemologia; relaciona word; relaciona transcricao}
\prov{proveniência: livro \textperiodcentered{} inferencia \textperiodcentered{} confiança 1.0 \textperiodcentered{} domínio filosofia \textperiodcentered{} história 6 versões no jornal}

%CRISTAL {"arestas":[{"alvo":"sunyata_vazio","confianca":1.0,"epistemico":"fato","origem":"wiki","rel":"relaciona"},{"alvo":"deus_postulado","confianca":1.0,"epistemico":"fato","origem":"wiki","rel":"especializa"},{"alvo":"consciencia_pura","confianca":1.0,"epistemico":"fato","origem":"wiki","rel":"relaciona"}],"confianca":1.0,"contraexemplos":[],"descricao":"Lopes: o único Deus que não escraviza é o Não-Ser (Vazio) — 'vazio de identidade, pleno de potencialidade'; distingue o Absoluto impessoal do Deus construído.","epistemico":"hipotese","exemplos":[],"id":"deus_vazio","memoria":"semantica","meta":{"autor":"Gentil Lopes","dominio":"filosofia","fonte":""},"origem":"livro","palavras_chave":[],"sinonimos":[],"tipo":"conceito","titulo":"Deus Vazio / o Absoluto"}
\section{Deus Vazio / o Absoluto}
Lopes: o único Deus que não escraviza é o Não-Ser (Vazio) — 'vazio de identidade, pleno de potencialidade'; distingue o Absoluto impessoal do Deus construído.
\prov{relações: relaciona sunyata\_vazio; especializa deus\_postulado; relaciona consciencia\_pura}
\prov{proveniência: livro \textperiodcentered{} hipotese \textperiodcentered{} confiança 1.0 \textperiodcentered{} domínio filosofia \textperiodcentered{} história 5 versões no jornal}

%CRISTAL {"arestas":[{"alvo":"filosofia_da_educacao","confianca":1.0,"epistemico":"fato","origem":"wiki","rel":"parte_de"},{"alvo":"aprender_fazendo_dewey","confianca":1.0,"epistemico":"fato","origem":"wiki","rel":"usa"},{"alvo":"dialogicidade","confianca":1.0,"epistemico":"fato","origem":"wiki","rel":"relaciona"}],"confianca":1.0,"contraexemplos":[],"descricao":"Dewey: a escola é uma comunidade em miniatura onde se formam cidadãos críticos e cooperativos; democracia e educação são interdependentes, formando o pensamento reflexivo pela experiência.","epistemico":"inferencia","exemplos":[],"id":"dewey_democracia_educacao","memoria":"semantica","meta":{"autor":"Dewey","dominio":"ciencias_de_fora","fonte":""},"origem":"wiki","palavras_chave":[],"sinonimos":[],"tipo":"conceito","titulo":"Democracia e Educação (Dewey)"}
\section{Democracia e Educação (Dewey)}
Dewey: a escola é uma comunidade em miniatura onde se formam cidadãos críticos e cooperativos; democracia e educação são interdependentes, formando o pensamento reflexivo pela experiência.
\prov{relações: parte\_de filosofia\_da\_educacao; usa aprender\_fazendo\_dewey; relaciona dialogicidade}
\prov{proveniência: wiki \textperiodcentered{} inferencia \textperiodcentered{} confiança 1.0 \textperiodcentered{} domínio ciencias\_de\_fora \textperiodcentered{} história 4 versões no jornal}

%CRISTAL {"arestas":[{"alvo":"karma_dharma","confianca":1.0,"epistemico":"fato","origem":"wiki","rel":"usa"},{"alvo":"dialogo_kurukshetra","confianca":1.0,"epistemico":"fato","origem":"wiki","rel":"usa"}],"confianca":1.0,"contraexemplos":[],"descricao":"Gita (3.35): dharma é o que sustenta o ser e a ordem; svadharma é o dever próprio de cada um segundo sua natureza. 'Melhor cumprir o próprio dever, ainda que imperfeito, que o dever alheio bem cumprido' — o núcleo da resposta de Krishna à angústia de Arjuna.","epistemico":"inferencia","exemplos":[],"id":"dharma_svadharma","memoria":"semantica","meta":{"autor":"Bhagavad Gita","dominio":"sagrado","fonte":""},"origem":"wiki","palavras_chave":[],"sinonimos":[],"tipo":"conceito","titulo":"Dharma e Svadharma — o Dever Próprio"}
\section{Dharma e Svadharma — o Dever Próprio}
Gita (3.35): dharma é o que sustenta o ser e a ordem; svadharma é o dever próprio de cada um segundo sua natureza. 'Melhor cumprir o próprio dever, ainda que imperfeito, que o dever alheio bem cumprido' — o núcleo da resposta de Krishna à angústia de Arjuna.
\prov{relações: usa karma\_dharma; usa dialogo\_kurukshetra}
\prov{proveniência: wiki \textperiodcentered{} inferencia \textperiodcentered{} confiança 1.0 \textperiodcentered{} domínio sagrado \textperiodcentered{} história 3 versões no jornal}

%CRISTAL {"arestas":[{"alvo":"teorema_de_bayes","confianca":1.0,"epistemico":"fato","origem":"wiki","rel":"usa"},{"alvo":"decisao_bayesiana","confianca":1.0,"epistemico":"fato","origem":"wiki","rel":"relaciona"},{"alvo":"medicina","confianca":1.0,"epistemico":"fato","origem":"wiki","rel":"parte_de"},{"alvo":"colapso","confianca":1.0,"epistemico":"fato","origem":"wiki","rel":"eh_um"},{"alvo":"sensor_e_colapso","confianca":1.0,"epistemico":"fato","origem":"wiki","rel":"relaciona"},{"alvo":"sobrediagnostico","confianca":1.0,"epistemico":"fato","origem":"wiki","rel":"explica"},{"alvo":"biotecnologia","confianca":1.0,"epistemico":"fato","origem":"wiki","rel":"relaciona"}],"confianca":1.0,"contraexemplos":[],"descricao":"O diagnóstico diferencial é uma SUPERPOSIÇÃO de hipóteses; o diagnóstico DESABA nela ponderando a prevalência (a taxa-base) pela verossimilhança dos achados — o Ψ do BAI. O orçamento δ é quanto examinar antes de se comprometer, e o deny-overrides é descartar primeiro o grave. O mesmo teste positivo colapsa diferente conforme o prior. Verif.: teste 90/90 → 8% (rara) vs 79% (comum).","epistemico":"inferencia","exemplos":[],"id":"diagnostico_colapso","memoria":"semantica","meta":{"dominio":"medicina","fonte":"medicina e diagnóstico","posterior_comum":0.794,"posterior_rara":0.083},"origem":"wiki","palavras_chave":[],"sinonimos":[],"tipo":"conceito","titulo":"Diagnóstico = Colapso (Bayes)"}
\section{Diagnóstico = Colapso (Bayes)}
O diagnóstico diferencial é uma SUPERPOSIÇÃO de hipóteses; o diagnóstico DESABA nela ponderando a prevalência (a taxa-base) pela verossimilhança dos achados — o \ensuremath{\Psi} do BAI. O orçamento \ensuremath{\delta} é quanto examinar antes de se comprometer, e o deny-overrides é descartar primeiro o grave. O mesmo teste positivo colapsa diferente conforme o prior. Verif.: teste 90/90 \ensuremath{\to} 8\% (rara) vs 79\% (comum).
\prov{relações: usa teorema\_de\_bayes; relaciona decisao\_bayesiana; parte\_de medicina; eh\_um colapso; relaciona sensor\_e\_colapso; explica sobrediagnostico; relaciona biotecnologia}
\prov{proveniência: wiki \textperiodcentered{} medicina e diagnóstico \textperiodcentered{} inferencia \textperiodcentered{} confiança 1.0 \textperiodcentered{} domínio medicina \textperiodcentered{} história 8 versões no jornal}

%CRISTAL {"arestas":[{"alvo":"hegel_absoluto","confianca":1.0,"epistemico":"fato","origem":"wiki","rel":"usa"},{"alvo":"filosofia_da_historia","confianca":1.0,"epistemico":"fato","origem":"wiki","rel":"parte_de"},{"alvo":"humanidade_sonha","confianca":1.0,"epistemico":"fato","origem":"wiki","rel":"relaciona"}],"confianca":1.0,"contraexemplos":[],"descricao":"Hegel: o Espírito se realiza na história por tese-antítese-síntese — a história é o progresso da consciência de liberdade. Historicismo especulativo.","epistemico":"hipotese","exemplos":[],"id":"dialetica_hegeliana","memoria":"semantica","meta":{"autor":"Hegel","dominio":"ciencias_de_fora","fonte":""},"origem":"wiki","palavras_chave":[],"sinonimos":[],"tipo":"conceito","titulo":"Dialética da História"}
\section{Dialética da História}
Hegel: o Espírito se realiza na história por tese-antítese-síntese — a história é o progresso da consciência de liberdade. Historicismo especulativo.
\prov{relações: usa hegel\_absoluto; parte\_de filosofia\_da\_historia; relaciona humanidade\_sonha}
\prov{proveniência: wiki \textperiodcentered{} hipotese \textperiodcentered{} confiança 1.0 \textperiodcentered{} domínio ciencias\_de\_fora \textperiodcentered{} história 4 versões no jornal}

%CRISTAL {"arestas":[{"alvo":"idioma","confianca":1.0,"epistemico":"fato","origem":"wiki","rel":"usa"},{"alvo":"sociolinguistica","confianca":1.0,"epistemico":"fato","origem":"wiki","rel":"usa"}],"confianca":1.0,"contraexemplos":[],"descricao":"Variantes de uma língua por região, classe ou grupo — nenhuma 'errada'; a fronteira língua/dialeto é política, não linguística.","epistemico":"inferencia","exemplos":[],"id":"dialeto_e_variacao","memoria":"semantica","meta":{"dominio":"ciencias_de_fora","fonte":""},"origem":"wiki","palavras_chave":[],"sinonimos":[],"tipo":"conceito","titulo":"Dialeto e Variação"}
\section{Dialeto e Variação}
Variantes de uma língua por região, classe ou grupo — nenhuma 'errada'; a fronteira língua/dialeto é política, não linguística.
\prov{relações: usa idioma; usa sociolinguistica}
\prov{proveniência: wiki \textperiodcentered{} inferencia \textperiodcentered{} confiança 1.0 \textperiodcentered{} domínio ciencias\_de\_fora \textperiodcentered{} história 3 versões no jornal}

%CRISTAL {"arestas":[{"alvo":"educacao_problematizadora","confianca":1.0,"epistemico":"fato","origem":"wiki","rel":"usa"}],"confianca":1.0,"contraexemplos":[],"descricao":"Freire: o conhecimento se produz na relação horizontal mediada pelo mundo — 'ninguém educa ninguém, ninguém se educa a si mesmo: os homens se educam em comunhão, mediatizados pelo mundo'.","epistemico":"inferencia","exemplos":[],"id":"dialogicidade","memoria":"semantica","meta":{"autor":"Freire","dominio":"ciencias_de_fora","fonte":""},"origem":"wiki","palavras_chave":[],"sinonimos":[],"tipo":"conceito","titulo":"Dialogicidade"}
\section{Dialogicidade}
Freire: o conhecimento se produz na relação horizontal mediada pelo mundo — 'ninguém educa ninguém, ninguém se educa a si mesmo: os homens se educam em comunhão, mediatizados pelo mundo'.
\prov{relações: usa educacao\_problematizadora}
\prov{proveniência: wiki \textperiodcentered{} inferencia \textperiodcentered{} confiança 1.0 \textperiodcentered{} domínio ciencias\_de\_fora \textperiodcentered{} história 2 versões no jornal}

%CRISTAL {"arestas":[{"alvo":"hinduismo","confianca":1.0,"epistemico":"fato","origem":"wiki","rel":"parte_de"},{"alvo":"karma_dharma","confianca":1.0,"epistemico":"fato","origem":"wiki","rel":"relaciona"},{"alvo":"consciencia_pura","confianca":1.0,"epistemico":"fato","origem":"wiki","rel":"relaciona"}],"confianca":1.0,"contraexemplos":[],"descricao":"Gita: no campo de batalha, prestes à guerra fratricida, o guerreiro Arjuna colapsa em angústia e recebe de Krishna (a Encarnação Divina) o ensinamento. O campo é metáfora: o corpo é o 'campo' (kshetra) onde brota o conflito da alma diante do dever, e a consciência é o 'conhecedor do campo' (kshetragña).","epistemico":"inferencia","exemplos":[],"id":"dialogo_kurukshetra","memoria":"semantica","meta":{"autor":"Bhagavad Gita","dominio":"sagrado","fonte":""},"origem":"wiki","palavras_chave":[],"sinonimos":[],"tipo":"conceito","titulo":"O Diálogo de Kurukshetra"}
\section{O Diálogo de Kurukshetra}
Gita: no campo de batalha, prestes à guerra fratricida, o guerreiro Arjuna colapsa em angústia e recebe de Krishna (a Encarnação Divina) o ensinamento. O campo é metáfora: o corpo é o 'campo' (kshetra) onde brota o conflito da alma diante do dever, e a consciência é o 'conhecedor do campo' (kshetragña).
\prov{relações: parte\_de hinduismo; relaciona karma\_dharma; relaciona consciencia\_pura}
\prov{proveniência: wiki \textperiodcentered{} inferencia \textperiodcentered{} confiança 1.0 \textperiodcentered{} domínio sagrado \textperiodcentered{} história 4 versões no jornal}

%CRISTAL {"arestas":[{"alvo":"atuadores_da_face","confianca":1.0,"epistemico":"fato","origem":"wiki","rel":"parte_de"},{"alvo":"musculo_mineral_hub","confianca":1.0,"epistemico":"fato","origem":"wiki","rel":"usa"},{"alvo":"corpo_dual","confianca":1.0,"epistemico":"fato","origem":"wiki","rel":"relaciona"}],"confianca":1.0,"contraexemplos":[],"descricao":"Os 8 músculos faciais (os atuadores) em quatro leituras, no pivô 'musculo_face' (dicionarios.py): nosso NOME ↔ nome em LATIM ↔ Action Unit de Ekman (FACS) ↔ FUNÇÃO. compor traduz de graça: zigomático = zygomaticus_major = AU12 = sorri; corrugador = corrugator_supercilii = AU4 = franze o cenho; masseter = AU26 = abre a mandíbula. É bidirecional (AU12→zigomático). A EXPRESSÃO é a combinação das AUs — a mesma máquina de tradução que faz xadrez↔BAI↔imagem, agora nos músculos.","epistemico":"fato","exemplos":[],"id":"dicionario_musculos_face","memoria":"semantica","meta":{"dominio":"biologia","fonte":"camadas da face"},"origem":"wiki","palavras_chave":[],"sinonimos":[],"tipo":"conceito","titulo":"O Dicionário dos Músculos da Face"}
\section{O Dicionário dos Músculos da Face}
Os 8 músculos faciais (os atuadores) em quatro leituras, no pivô 'musculo\_face' (dicionarios.py): nosso NOME \ensuremath{\leftrightarrow} nome em LATIM \ensuremath{\leftrightarrow} Action Unit de Ekman (FACS) \ensuremath{\leftrightarrow} FUNÇÃO. compor traduz de graça: zigomático = zygomaticus\_major = AU12 = sorri; corrugador = corrugator\_supercilii = AU4 = franze o cenho; masseter = AU26 = abre a mandíbula. É bidirecional (AU12\ensuremath{\to}zigomático). A EXPRESSÃO é a combinação das AUs — a mesma máquina de tradução que faz xadrez\ensuremath{\leftrightarrow}BAI\ensuremath{\leftrightarrow}imagem, agora nos músculos.
\prov{relações: parte\_de atuadores\_da\_face; usa musculo\_mineral\_hub; relaciona corpo\_dual}
\prov{proveniência: wiki \textperiodcentered{} camadas da face \textperiodcentered{} fato \textperiodcentered{} confiança 1.0 \textperiodcentered{} domínio biologia \textperiodcentered{} história 4 versões no jornal}

%CRISTAL {"arestas":[{"alvo":"omnitrix_ben10","confianca":1.0,"epistemico":"fato","origem":"wiki","rel":"parte_de"},{"alvo":"omnitrix_corpo_universal","confianca":1.0,"epistemico":"fato","origem":"wiki","rel":"usa"},{"alvo":"dicionario_bai_omnitrix","confianca":1.0,"epistemico":"fato","origem":"wiki","rel":"relaciona"},{"alvo":"antimateria_parabolica","confianca":1.0,"epistemico":"fato","origem":"wiki","rel":"relaciona"},{"alvo":"minecraft_voxel","confianca":1.0,"epistemico":"fato","origem":"wiki","rel":"relaciona"}],"confianca":1.0,"contraexemplos":[],"descricao":"O dicionário paralelo entre o Omnitrix de Ben 10 (o relógio) e o Omnitrix do broca-so (o Corpo Algébrico, o template Sym⊕Skew⊕invariante): o RELÓGIO Omnitrix (o dispositivo no pulso) ↔ o TEMPLATE Omnitrix (o corpo algébrico); um ALIENÍGENA (uma forma: Quatro Braços, Chama, …) ↔ uma INSTÂNCIA (um corpo: matrix, molecular, estelar, cósmico, BAI); o CODON STREAM / o DNA das espécies (o codex) ↔ a ÁLGEBRA Sym⊕Skew⊕invariante (o código de todos os corpos); DISCAR+BATER (escolher e virar o alien) ↔ INSTANCIAR (escolher a ordem m/dimensão/substrato e virar o corpo); o TIMEOUT (volta a ser Ben, humano) ↔ o retorno ao ZERO DA RETA (a aniquilação ao vácuo, o produto σ·σ'=−1); o NÚCLEO de energia (recarrega, nunca se esgota de vez) ↔ o INVARIANTE a+c²=1, a≥a_min>0 (a ZPE que nunca zera, o mass gap 𝔤₂); AZMUTH (o Galvan que o criou) ↔ o GERADOR/o autor (a parábola geradora); a MASTER CONTROL (destrava todos os aliens) ↔ o TEMPLATE COMPLETO (todas as instâncias acessíveis); ALIEN X (o quase-onipotente, 3 consciências que concordam) ↔ o ZERO DA RETA (a dinâmica matéria-antimatéria, o par simétrico ±1 de onde tudo nasce); os LOCKS/o Ultimatrix (quem pode virar o quê) ↔ o BAI (a autorização). Síntese declarada — o dicionário casa as figuras, não decreta a obra.","epistemico":"hipotese","exemplos":[],"id":"dicionario_omnitrix_ben10","memoria":"semantica","meta":{"dominio":"ficcao","fonte":"Ficção (jogos e animações)"},"origem":"wiki","palavras_chave":[],"sinonimos":[],"tipo":"conceito","titulo":"Dicionário: Omnitrix (Ben 10) ↔ Omnitrix (Corpo Algébrico)"}
\section{Dicionário: Omnitrix (Ben 10) \ensuremath{\leftrightarrow} Omnitrix (Corpo Algébrico)}
O dicionário paralelo entre o Omnitrix de Ben 10 (o relógio) e o Omnitrix do broca-so (o Corpo Algébrico, o template Sym\ensuremath{\oplus}Skew\ensuremath{\oplus}invariante): o RELÓGIO Omnitrix (o dispositivo no pulso) \ensuremath{\leftrightarrow} o TEMPLATE Omnitrix (o corpo algébrico); um ALIENÍGENA (uma forma: Quatro Braços, Chama, …) \ensuremath{\leftrightarrow} uma INSTÂNCIA (um corpo: matrix, molecular, estelar, cósmico, BAI); o CODON STREAM / o DNA das espécies (o codex) \ensuremath{\leftrightarrow} a ÁLGEBRA Sym\ensuremath{\oplus}Skew\ensuremath{\oplus}invariante (o código de todos os corpos); DISCAR+BATER (escolher e virar o alien) \ensuremath{\leftrightarrow} INSTANCIAR (escolher a ordem m/dimensão/substrato e virar o corpo); o TIMEOUT (volta a ser Ben, humano) \ensuremath{\leftrightarrow} o retorno ao ZERO DA RETA (a aniquilação ao vácuo, o produto \ensuremath{\sigma}\textperiodcentered \ensuremath{\sigma}'=\ensuremath{-}1); o NÚCLEO de energia (recarrega, nunca se esgota de vez) \ensuremath{\leftrightarrow} o INVARIANTE a+c\ensuremath{^{2}}=1, a\ensuremath{\ge}a\_min>0 (a ZPE que nunca zera, o mass gap \ensuremath{\mathfrak{g}}\ensuremath{_{2}}); AZMUTH (o Galvan que o criou) \ensuremath{\leftrightarrow} o GERADOR/o autor (a parábola geradora); a MASTER CONTROL (destrava todos os aliens) \ensuremath{\leftrightarrow} o TEMPLATE COMPLETO (todas as instâncias acessíveis); ALIEN X (o quase-onipotente, 3 consciências que concordam) \ensuremath{\leftrightarrow} o ZERO DA RETA (a dinâmica matéria-antimatéria, o par simétrico $\pm$1 de onde tudo nasce); os LOCKS/o Ultimatrix (quem pode virar o quê) \ensuremath{\leftrightarrow} o BAI (a autorização). Síntese declarada — o dicionário casa as figuras, não decreta a obra.
\prov{relações: parte\_de omnitrix\_ben10; usa omnitrix\_corpo\_universal; relaciona dicionario\_bai\_omnitrix; relaciona antimateria\_parabolica; relaciona minecraft\_voxel}
\prov{proveniência: wiki \textperiodcentered{} Ficção (jogos e animações) \textperiodcentered{} hipotese \textperiodcentered{} confiança 1.0 \textperiodcentered{} domínio ficcao \textperiodcentered{} história 6 versões no jornal}

%CRISTAL {"arestas":[{"alvo":"orbitas_de_fala","confianca":1.0,"epistemico":"fato","origem":"wiki","rel":"parte_de"},{"alvo":"dicionario_musculos_face","confianca":1.0,"epistemico":"fato","origem":"wiki","rel":"relaciona"},{"alvo":"corpo_dual","confianca":1.0,"epistemico":"fato","origem":"wiki","rel":"relaciona"}],"confianca":1.0,"contraexemplos":[],"descricao":"As formas de boca da fala, no pivô 'visema' (dicionarios.py): fonema ↔ visema ↔ som. compor traduz de graça: a=AA=aberto; u=OO=arredondado; MM=labios_juntos; f=FF=labio_dental. É a mesma máquina de tradução dos músculos e das línguas, agora na fala — os visemas são as 'letras' visuais da boca, e uma frase é a sua sequência (a órbita de fala).","epistemico":"fato","exemplos":[],"id":"dicionario_visemas","memoria":"semantica","meta":{"dominio":"biologia","fonte":"órbitas de fala"},"origem":"wiki","palavras_chave":[],"sinonimos":[],"tipo":"conceito","titulo":"O Dicionário dos Visemas (fonema ↔ visema ↔ som)"}
\section{O Dicionário dos Visemas (fonema \ensuremath{\leftrightarrow} visema \ensuremath{\leftrightarrow} som)}
As formas de boca da fala, no pivô 'visema' (dicionarios.py): fonema \ensuremath{\leftrightarrow} visema \ensuremath{\leftrightarrow} som. compor traduz de graça: a=AA=aberto; u=OO=arredondado; MM=labios\_juntos; f=FF=labio\_dental. É a mesma máquina de tradução dos músculos e das línguas, agora na fala — os visemas são as 'letras' visuais da boca, e uma frase é a sua sequência (a órbita de fala).
\prov{relações: parte\_de orbitas\_de\_fala; relaciona dicionario\_musculos\_face; relaciona corpo\_dual}
\prov{proveniência: wiki \textperiodcentered{} órbitas de fala \textperiodcentered{} fato \textperiodcentered{} confiança 1.0 \textperiodcentered{} domínio biologia \textperiodcentered{} história 8 versões no jornal}

%CRISTAL {"arestas":[{"alvo":"pedagogia","confianca":1.0,"epistemico":"fato","origem":"wiki","rel":"parte_de"}],"confianca":1.0,"contraexemplos":[],"descricao":"Comenius: a primeira sistematização moderna do ensino — a arte universal de 'ensinar tudo a todos' (omnes, omnia), inclusive pobres e meninas, por um método ordenado. Pai da didática moderna.","epistemico":"fato","exemplos":[],"id":"didatica_magna_comenius","memoria":"semantica","meta":{"autor":"Comenius","dominio":"ciencias_de_fora","fonte":""},"origem":"wiki","palavras_chave":[],"sinonimos":[],"tipo":"conceito","titulo":"Didática Magna (Comenius)"}
\section{Didática Magna (Comenius)}
Comenius: a primeira sistematização moderna do ensino — a arte universal de 'ensinar tudo a todos' (omnes, omnia), inclusive pobres e meninas, por um método ordenado. Pai da didática moderna.
\prov{relações: parte\_de pedagogia}
\prov{proveniência: wiki \textperiodcentered{} fato \textperiodcentered{} confiança 1.0 \textperiodcentered{} domínio ciencias\_de\_fora \textperiodcentered{} história 2 versões no jornal}

%CRISTAL {"arestas":[{"alvo":"pedagogia","confianca":1.0,"epistemico":"fato","origem":"wiki","rel":"parte_de"},{"alvo":"matematica","confianca":1.0,"epistemico":"fato","origem":"wiki","rel":"referencia"},{"alvo":"word","confianca":0.5,"epistemico":"fato","origem":"wiki","rel":"relaciona"}],"confianca":1.0,"contraexemplos":[],"descricao":"O estudo científico de como se ensina e se aprende matemática — situações, obstáculos e representações do saber matemático.","epistemico":"fato","exemplos":[],"id":"didatica_matematica","memoria":"semantica","meta":{"dominio":"ciencias_de_fora","fonte":""},"origem":"wiki","palavras_chave":[],"sinonimos":[],"tipo":"conceito","titulo":"Didática da Matemática"}
\section{Didática da Matemática}
O estudo científico de como se ensina e se aprende matemática — situações, obstáculos e representações do saber matemático.
\prov{relações: parte\_de pedagogia; referencia matematica; relaciona word}
\prov{proveniência: wiki \textperiodcentered{} fato \textperiodcentered{} confiança 1.0 \textperiodcentered{} domínio ciencias\_de\_fora \textperiodcentered{} história 4 versões no jornal}

%CRISTAL {"arestas":[{"alvo":"realismo_ri","confianca":1.0,"epistemico":"fato","origem":"wiki","rel":"usa"},{"alvo":"teoria_dos_jogos","confianca":1.0,"epistemico":"fato","origem":"wiki","rel":"usa"}],"confianca":1.0,"contraexemplos":[],"descricao":"Herz: medidas defensivas de um Estado para aumentar sua própria segurança são percebidas como ameaça pelos demais, gerando corrida armamentista e insegurança mútua — mesmo sem intenção agressiva. Uma estrutura tipo dilema do prisioneiro.","epistemico":"inferencia","exemplos":[],"id":"dilema_da_seguranca","memoria":"semantica","meta":{"autor":"Herz","dominio":"ciencias_de_fora","fonte":""},"origem":"wiki","palavras_chave":[],"sinonimos":[],"tipo":"conceito","titulo":"Dilema da Segurança"}
\section{Dilema da Segurança}
Herz: medidas defensivas de um Estado para aumentar sua própria segurança são percebidas como ameaça pelos demais, gerando corrida armamentista e insegurança mútua — mesmo sem intenção agressiva. Uma estrutura tipo dilema do prisioneiro.
\prov{relações: usa realismo\_ri; usa teoria\_dos\_jogos}
\prov{proveniência: wiki \textperiodcentered{} inferencia \textperiodcentered{} confiança 1.0 \textperiodcentered{} domínio ciencias\_de\_fora \textperiodcentered{} história 3 versões no jornal}

%CRISTAL {"arestas":[{"alvo":"geometria_fractal","confianca":1.0,"epistemico":"fato","origem":"wiki","rel":"usa"},{"alvo":"topologia","confianca":1.0,"epistemico":"fato","origem":"wiki","rel":"parte_de"}],"confianca":1.0,"contraexemplos":[],"descricao":"Uma dimensão fracionária que mede a rugosidade de um conjunto — 1,26… para o floco de Koch. A régua do não-inteiro.","epistemico":"fato","exemplos":[],"id":"dimensao_de_hausdorff","memoria":"semantica","meta":{"autor":"Hausdorff","dominio":"ciencias_de_fora","fonte":""},"origem":"wiki","palavras_chave":[],"sinonimos":[],"tipo":"conceito","titulo":"Dimensão de Hausdorff"}
\section{Dimensão de Hausdorff}
Uma dimensão fracionária que mede a rugosidade de um conjunto — 1,26… para o floco de Koch. A régua do não-inteiro.
\prov{relações: usa geometria\_fractal; parte\_de topologia}
\prov{proveniência: wiki \textperiodcentered{} fato \textperiodcentered{} confiança 1.0 \textperiodcentered{} domínio ciencias\_de\_fora \textperiodcentered{} história 6 versões no jornal}

%CRISTAL {"arestas":[{"alvo":"trabalho","confianca":1.0,"epistemico":"fato","origem":"curriculo","rel":"continua"},{"alvo":"virtualizacao","confianca":0.5,"epistemico":"fato","origem":"wiki","rel":"relaciona"},{"alvo":"word","confianca":0.5,"epistemico":"fato","origem":"wiki","rel":"relaciona"},{"alvo":"vida-morte","confianca":0.5,"epistemico":"fato","origem":"wiki","rel":"relaciona"}],"confianca":1.0,"contraexemplos":[],"descricao":"Meio de troca e unidade de conta; no projeto: lastro não mora na conta ETH (só gás).","epistemico":"fato","exemplos":[],"id":"dinheiro","memoria":"semantica","meta":{"dominio":"mercado","nivel":5,"ordem":0},"origem":"curriculo","palavras_chave":["lastro","Pell","tesouraria"],"sinonimos":["moeda"],"tipo":"conceito","titulo":"dinheiro"}
\section{dinheiro}
Meio de troca e unidade de conta; no projeto: lastro não mora na conta ETH (só gás).
\prov{sinónimos: moeda}
\prov{relações: continua trabalho; relaciona virtualizacao; relaciona word; relaciona vida-morte}
\prov{proveniência: curriculo \textperiodcentered{} fato \textperiodcentered{} confiança 1.0 \textperiodcentered{} domínio mercado \textperiodcentered{} história 7 versões no jornal}

%CRISTAL {"arestas":[{"alvo":"filosofia","confianca":1.0,"epistemico":"fato","origem":"wiki","rel":"relaciona"},{"alvo":"contrato_social","confianca":1.0,"epistemico":"fato","origem":"wiki","rel":"relaciona"},{"alvo":"word","confianca":0.5,"epistemico":"fato","origem":"wiki","rel":"relaciona"}],"confianca":1.0,"contraexemplos":[],"descricao":"O sistema de normas que regula a conduta e o poder numa comunidade — e a reflexão sobre sua natureza, sua justiça e sua legitimidade.","epistemico":"fato","exemplos":[],"id":"direito","memoria":"semantica","meta":{"dominio":"ciencias_de_fora","fonte":""},"origem":"wiki","palavras_chave":[],"sinonimos":[],"tipo":"conceito","titulo":"Direito"}
\section{Direito}
O sistema de normas que regula a conduta e o poder numa comunidade — e a reflexão sobre sua natureza, sua justiça e sua legitimidade.
\prov{relações: relaciona filosofia; relaciona contrato\_social; relaciona word}
\prov{proveniência: wiki \textperiodcentered{} fato \textperiodcentered{} confiança 1.0 \textperiodcentered{} domínio ciencias\_de\_fora \textperiodcentered{} história 4 versões no jornal}

%CRISTAL {"arestas":[{"alvo":"producao_do_espaco_lefebvre","confianca":1.0,"epistemico":"fato","origem":"wiki","rel":"usa"},{"alvo":"classe_social","confianca":1.0,"epistemico":"fato","origem":"wiki","rel":"relaciona"}],"confianca":1.0,"contraexemplos":[],"descricao":"Lefebvre: o direito coletivo de acesso à vida urbana e, sobretudo, de transformar a cidade e a nós mesmos ao produzi-la — reivindicação contra a apropriação privada/capitalista do espaço urbano.","epistemico":"inferencia","exemplos":[],"id":"direito_a_cidade","memoria":"semantica","meta":{"autor":"Lefebvre/Harvey","dominio":"ciencias_de_fora","fonte":""},"origem":"wiki","palavras_chave":[],"sinonimos":[],"tipo":"conceito","titulo":"Direito à Cidade (Lefebvre/Harvey)"}
\section{Direito à Cidade (Lefebvre/Harvey)}
Lefebvre: o direito coletivo de acesso à vida urbana e, sobretudo, de transformar a cidade e a nós mesmos ao produzi-la — reivindicação contra a apropriação privada/capitalista do espaço urbano.
\prov{relações: usa producao\_do\_espaco\_lefebvre; relaciona classe\_social}
\prov{proveniência: wiki \textperiodcentered{} inferencia \textperiodcentered{} confiança 1.0 \textperiodcentered{} domínio ciencias\_de\_fora \textperiodcentered{} história 3 versões no jornal}

%CRISTAL {"arestas":[{"alvo":"direito_e_tecnologia","confianca":1.0,"epistemico":"fato","origem":"wiki","rel":"parte_de"},{"alvo":"capitalismo_de_vigilancia","confianca":1.0,"epistemico":"fato","origem":"wiki","rel":"contradiz"},{"alvo":"code_is_law","confianca":1.0,"epistemico":"fato","origem":"wiki","rel":"relaciona"}],"confianca":1.0,"contraexemplos":[],"descricao":"GDPR/LGPD: o direito do titular de exigir a eliminação de seus dados pessoais quando não mais necessários — controle individual sobre a própria informação, contrapeso jurídico à mercantilização do dado (e em tensão com a imutabilidade da blockchain).","epistemico":"fato","exemplos":[],"id":"direito_ao_esquecimento","memoria":"semantica","meta":{"autor":"GDPR/LGPD","dominio":"ciencias_de_fora","fonte":""},"origem":"wiki","palavras_chave":[],"sinonimos":[],"tipo":"conceito","titulo":"Direito ao Esquecimento e Proteção de Dados"}
\section{Direito ao Esquecimento e Proteção de Dados}
GDPR/LGPD: o direito do titular de exigir a eliminação de seus dados pessoais quando não mais necessários — controle individual sobre a própria informação, contrapeso jurídico à mercantilização do dado (e em tensão com a imutabilidade da blockchain).
\prov{relações: parte\_de direito\_e\_tecnologia; contradiz capitalismo\_de\_vigilancia; relaciona code\_is\_law}
\prov{proveniência: wiki \textperiodcentered{} fato \textperiodcentered{} confiança 1.0 \textperiodcentered{} domínio ciencias\_de\_fora \textperiodcentered{} história 4 versões no jornal}

%CRISTAL {"arestas":[{"alvo":"direito","confianca":1.0,"epistemico":"fato","origem":"wiki","rel":"parte_de"},{"alvo":"virtualizacao","confianca":0.5,"epistemico":"fato","origem":"wiki","rel":"relaciona"},{"alvo":"word","confianca":0.5,"epistemico":"fato","origem":"wiki","rel":"relaciona"}],"confianca":1.0,"contraexemplos":[],"descricao":"A interface do direito com a eficiência econômica, o código de software e o protocolo — de Coase ao 'code is law' e aos contratos autoexecutáveis.","epistemico":"fato","exemplos":[],"id":"direito_e_tecnologia","memoria":"semantica","meta":{"dominio":"ciencias_de_fora","fonte":""},"origem":"wiki","palavras_chave":[],"sinonimos":[],"tipo":"conceito","titulo":"Direito, Economia e Tecnologia"}
\section{Direito, Economia e Tecnologia}
A interface do direito com a eficiência econômica, o código de software e o protocolo — de Coase ao 'code is law' e aos contratos autoexecutáveis.
\prov{relações: parte\_de direito; relaciona virtualizacao; relaciona word}
\prov{proveniência: wiki \textperiodcentered{} fato \textperiodcentered{} confiança 1.0 \textperiodcentered{} domínio ciencias\_de\_fora \textperiodcentered{} história 4 versões no jornal}

%CRISTAL {"arestas":[{"alvo":"direito","confianca":1.0,"epistemico":"fato","origem":"wiki","rel":"parte_de"},{"alvo":"contrato_social","confianca":1.0,"epistemico":"fato","origem":"wiki","rel":"relaciona"},{"alvo":"jusnaturalismo","confianca":1.0,"epistemico":"fato","origem":"wiki","rel":"usa"}],"confianca":1.0,"contraexemplos":[],"descricao":"Grotius: funda o direito entre Estados no direito natural — os tratados obrigam (pacta sunt servanda), há limites morais à guerra, e a soberania coexiste com deveres de cooperação.","epistemico":"inferencia","exemplos":[],"id":"direito_internacional","memoria":"semantica","meta":{"autor":"Grotius","dominio":"ciencias_de_fora","fonte":""},"origem":"wiki","palavras_chave":[],"sinonimos":[],"tipo":"conceito","titulo":"Direito Internacional (Grotius)"}
\section{Direito Internacional (Grotius)}
Grotius: funda o direito entre Estados no direito natural — os tratados obrigam (pacta sunt servanda), há limites morais à guerra, e a soberania coexiste com deveres de cooperação.
\prov{relações: parte\_de direito; relaciona contrato\_social; usa jusnaturalismo}
\prov{proveniência: wiki \textperiodcentered{} inferencia \textperiodcentered{} confiança 1.0 \textperiodcentered{} domínio ciencias\_de\_fora \textperiodcentered{} história 4 versões no jornal}

%CRISTAL {"arestas":[{"alvo":"direito","confianca":1.0,"epistemico":"fato","origem":"wiki","rel":"parte_de"},{"alvo":"imperativo_categorico","confianca":1.0,"epistemico":"fato","origem":"wiki","rel":"usa"}],"confianca":1.0,"contraexemplos":[],"descricao":"Normas universais que protegem interesses básicos de todo ser humano; debatidos entre universalismo (validade transcultural) e relativismo (padrões locais), e classificados em 'gerações' (civis/políticos; sociais/econômicos; de solidariedade).","epistemico":"inferencia","exemplos":[],"id":"direitos_humanos","memoria":"semantica","meta":{"autor":"—","dominio":"ciencias_de_fora","fonte":""},"origem":"wiki","palavras_chave":[],"sinonimos":[],"tipo":"conceito","titulo":"Direitos Humanos"}
\section{Direitos Humanos}
Normas universais que protegem interesses básicos de todo ser humano; debatidos entre universalismo (validade transcultural) e relativismo (padrões locais), e classificados em 'gerações' (civis/políticos; sociais/econômicos; de solidariedade).
\prov{relações: parte\_de direito; usa imperativo\_categorico}
\prov{proveniência: wiki \textperiodcentered{} inferencia \textperiodcentered{} confiança 1.0 \textperiodcentered{} domínio ciencias\_de\_fora \textperiodcentered{} história 3 versões no jornal}

%CRISTAL {"arestas":[{"alvo":"filosofia_da_medicina","confianca":1.0,"epistemico":"fato","origem":"wiki","rel":"parte_de"},{"alvo":"naturalismo_boorse","confianca":1.0,"epistemico":"fato","origem":"wiki","rel":"usa"},{"alvo":"aparelho_psiquico","confianca":1.0,"epistemico":"fato","origem":"wiki","rel":"relaciona"}],"confianca":1.0,"contraexemplos":[],"descricao":"A tríade: disease é a doença-entidade (processo patológico objetivo), illness é a experiência subjetiva de adoecer (sentir-se mal), sickness é o papel social do doente (reconhecimento público, isenções e obrigações).","epistemico":"inferencia","exemplos":[],"id":"disease_illness_sickness","memoria":"semantica","meta":{"autor":"Marinker/Twaddle","dominio":"ciencias_de_fora","fonte":""},"origem":"wiki","palavras_chave":[],"sinonimos":[],"tipo":"conceito","titulo":"Disease / Illness / Sickness"}
\section{Disease / Illness / Sickness}
A tríade: disease é a doença-entidade (processo patológico objetivo), illness é a experiência subjetiva de adoecer (sentir-se mal), sickness é o papel social do doente (reconhecimento público, isenções e obrigações).
\prov{relações: parte\_de filosofia\_da\_medicina; usa naturalismo\_boorse; relaciona aparelho\_psiquico}
\prov{proveniência: wiki \textperiodcentered{} inferencia \textperiodcentered{} confiança 1.0 \textperiodcentered{} domínio ciencias\_de\_fora \textperiodcentered{} história 4 versões no jornal}

%CRISTAL {"arestas":[{"alvo":"vida","confianca":1.0,"epistemico":"fato","origem":"wiki","rel":"explica"},{"alvo":"word","confianca":0.5,"epistemico":"fato","origem":"wiki","rel":"relaciona"},{"alvo":"syscall","confianca":0.5,"epistemico":"fato","origem":"wiki","rel":"relaciona"}],"confianca":1.0,"contraexemplos":[],"descricao":"Prigogine: longe do equilíbrio, sistemas abertos criam ordem a partir do caos, trocando energia — a vida como estrutura dissipativa.","epistemico":"fato","exemplos":[],"id":"dissipacao_prigogine","memoria":"semantica","meta":{"autor":"Prigogine","dominio":"ciencias_de_fora","fonte":""},"origem":"wiki","palavras_chave":[],"sinonimos":[],"tipo":"conceito","titulo":"Estruturas Dissipativas"}
\section{Estruturas Dissipativas}
Prigogine: longe do equilíbrio, sistemas abertos criam ordem a partir do caos, trocando energia — a vida como estrutura dissipativa.
\prov{relações: explica vida; relaciona word; relaciona syscall}
\prov{proveniência: wiki \textperiodcentered{} fato \textperiodcentered{} confiança 1.0 \textperiodcentered{} domínio ciencias\_de\_fora \textperiodcentered{} história 4 versões no jornal}

%CRISTAL {"arestas":[{"alvo":"psicologia","confianca":1.0,"epistemico":"fato","origem":"wiki","rel":"parte_de"},{"alvo":"idiotismo_nao_examinar","confianca":1.0,"epistemico":"fato","origem":"wiki","rel":"relaciona"}],"confianca":1.0,"contraexemplos":[],"descricao":"Festinger: a inconsistência entre cognições gera desconforto; muda-se crença ou comportamento para restaurar a coerência do ego.","epistemico":"fato","exemplos":[],"id":"dissonancia_cognitiva","memoria":"semantica","meta":{"autor":"Festinger","dominio":"ciencias_de_fora","fonte":""},"origem":"wiki","palavras_chave":[],"sinonimos":[],"tipo":"conceito","titulo":"Dissonância Cognitiva"}
\section{Dissonância Cognitiva}
Festinger: a inconsistência entre cognições gera desconforto; muda-se crença ou comportamento para restaurar a coerência do ego.
\prov{relações: parte\_de psicologia; relaciona idiotismo\_nao\_examinar}
\prov{proveniência: wiki \textperiodcentered{} fato \textperiodcentered{} confiança 1.0 \textperiodcentered{} domínio ciencias\_de\_fora \textperiodcentered{} história 3 versões no jornal}

%CRISTAL {"arestas":[{"alvo":"genetica","confianca":1.0,"epistemico":"fato","origem":"wiki","rel":"parte_de"},{"alvo":"gene","confianca":1.0,"epistemico":"fato","origem":"wiki","rel":"explica"}],"confianca":1.0,"contraexemplos":[],"descricao":"Watson, Crick e Franklin: o DNA é uma escada torcida cujas fitas complementares guardam e copiam a informação.","epistemico":"fato","exemplos":[],"id":"dna_dupla_helice","memoria":"semantica","meta":{"autor":"Watson/Crick/Franklin","dominio":"ciencias_de_fora","fonte":""},"origem":"wiki","palavras_chave":[],"sinonimos":[],"tipo":"conceito","titulo":"A Dupla Hélice do DNA"}
\section{A Dupla Hélice do DNA}
Watson, Crick e Franklin: o DNA é uma escada torcida cujas fitas complementares guardam e copiam a informação.
\prov{relações: parte\_de genetica; explica gene}
\prov{proveniência: wiki \textperiodcentered{} fato \textperiodcentered{} confiança 1.0 \textperiodcentered{} domínio ciencias\_de\_fora \textperiodcentered{} história 4 versões no jornal}

%CRISTAL {"arestas":[{"alvo":"codigo_genetico","confianca":1.0,"epistemico":"fato","origem":"wiki","rel":"usa"},{"alvo":"entropia_shannon","confianca":1.0,"epistemico":"fato","origem":"wiki","rel":"usa"}],"confianca":1.0,"contraexemplos":[],"descricao":"O DNA é um código de 4 letras (A,C,G,T); cada base equiprovável carrega no máximo log₂4=2 bits (entropia de Shannon). Genoma humano ~3,2×10⁹ pb → teto ~6,4×10⁹ bits ≈ 800 MB (compressível). O teto de 2 bits/base = fato matemático; o conteúdo FUNCIONAL de informação (quanto é funcional vs redundante) = aberto/debatido.","epistemico":"inferencia","exemplos":[],"id":"dna_informacao_shannon","memoria":"semantica","meta":{"autor":"broca-so","dominio":"biologia","fonte":"Shannon 1948; Schneider molecular information theory"},"origem":"wiki","palavras_chave":[],"sinonimos":[],"tipo":"conceito","titulo":"O DNA como código digital (~2 bits/base)"}
\section{O DNA como código digital (\~{}2 bits/base)}
O DNA é um código de 4 letras (A,C,G,T); cada base equiprovável carrega no máximo log\ensuremath{_{2}}4=2 bits (entropia de Shannon). Genoma humano \~{}3,2$\times$10\ensuremath{^{9}} pb \ensuremath{\to} teto \~{}6,4$\times$10\ensuremath{^{9}} bits \ensuremath{\approx} 800 MB (compressível). O teto de 2 bits/base = fato matemático; o conteúdo FUNCIONAL de informação (quanto é funcional vs redundante) = aberto/debatido.
\prov{relações: usa codigo\_genetico; usa entropia\_shannon}
\prov{proveniência: wiki \textperiodcentered{} Shannon 1948; Schneider molecular information theory \textperiodcentered{} inferencia \textperiodcentered{} confiança 1.0 \textperiodcentered{} domínio biologia \textperiodcentered{} história 3 versões no jornal}

%CRISTAL {"arestas":[{"alvo":"bioquimica","confianca":1.0,"epistemico":"fato","origem":"wiki","rel":"parte_de"},{"alvo":"auto_organizacao","confianca":1.0,"epistemico":"fato","origem":"wiki","rel":"usa"}],"confianca":1.0,"contraexemplos":[],"descricao":"A cadeia de aminoácidos colapsa numa forma 3D que determina sua função — um problema de otimização que a célula resolve em segundos.","epistemico":"fato","exemplos":[],"id":"dobramento_proteico","memoria":"semantica","meta":{"dominio":"ciencias_de_fora","fonte":""},"origem":"wiki","palavras_chave":[],"sinonimos":[],"tipo":"conceito","titulo":"Dobramento de Proteínas"}
\section{Dobramento de Proteínas}
A cadeia de aminoácidos colapsa numa forma 3D que determina sua função — um problema de otimização que a célula resolve em segundos.
\prov{relações: parte\_de bioquimica; usa auto\_organizacao}
\prov{proveniência: wiki \textperiodcentered{} fato \textperiodcentered{} confiança 1.0 \textperiodcentered{} domínio ciencias\_de\_fora \textperiodcentered{} história 3 versões no jornal}

%CRISTAL {"arestas":[{"alvo":"temperamento_igual","confianca":1.0,"epistemico":"fato","origem":"wiki","rel":"usa"},{"alvo":"teoria_de_grupos","confianca":1.0,"epistemico":"fato","origem":"wiki","rel":"usa"},{"alvo":"tonalidade","confianca":1.0,"epistemico":"fato","origem":"wiki","rel":"contradiz"},{"alvo":"atonalidade_dodecafonismo","confianca":1.0,"epistemico":"fato","origem":"wiki","rel":"relaciona"}],"confianca":1.0,"contraexemplos":[],"descricao":"Schoenberg: compor com as 12 notas em séries permutadas, sem centro tonal — música como teoria de grupos aplicada.","epistemico":"inferencia","exemplos":[],"id":"dodecafonismo","memoria":"semantica","meta":{"autor":"Schoenberg","dominio":"ciencias_de_fora","fonte":""},"origem":"wiki","palavras_chave":[],"sinonimos":[],"tipo":"conceito","titulo":"Dodecafonismo"}
\section{Dodecafonismo}
Schoenberg: compor com as 12 notas em séries permutadas, sem centro tonal — música como teoria de grupos aplicada.
\prov{relações: usa temperamento\_igual; usa teoria\_de\_grupos; contradiz tonalidade; relaciona atonalidade\_dodecafonismo}
\prov{proveniência: wiki \textperiodcentered{} inferencia \textperiodcentered{} confiança 1.0 \textperiodcentered{} domínio ciencias\_de\_fora \textperiodcentered{} história 6 versões no jornal}

%CRISTAL {"arestas":[{"alvo":"dna_dupla_helice","confianca":1.0,"epistemico":"fato","origem":"wiki","rel":"usa"},{"alvo":"codigo_genetico","confianca":1.0,"epistemico":"fato","origem":"wiki","rel":"usa"},{"alvo":"expressao_genica","confianca":1.0,"epistemico":"fato","origem":"wiki","rel":"explica"}],"confianca":1.0,"contraexemplos":[],"descricao":"Crick: a informação flui DNA → RNA → proteína. A sequência é lida, transcrita e traduzida — nunca ao contrário (quase).","epistemico":"fato","exemplos":[],"id":"dogma_central","memoria":"semantica","meta":{"autor":"Crick","dominio":"ciencias_de_fora","fonte":""},"origem":"wiki","palavras_chave":[],"sinonimos":[],"tipo":"conceito","titulo":"Dogma Central da Biologia"}
\section{Dogma Central da Biologia}
Crick: a informação flui DNA \ensuremath{\to} RNA \ensuremath{\to} proteína. A sequência é lida, transcrita e traduzida — nunca ao contrário (quase).
\prov{relações: usa dna\_dupla\_helice; usa codigo\_genetico; explica expressao\_genica}
\prov{proveniência: wiki \textperiodcentered{} fato \textperiodcentered{} confiança 1.0 \textperiodcentered{} domínio ciencias\_de\_fora \textperiodcentered{} história 5 versões no jornal}

%CRISTAL {"arestas":[{"alvo":"psicologia","confianca":1.0,"epistemico":"fato","origem":"wiki","rel":"parte_de"},{"alvo":"aparelho_psiquico","confianca":1.0,"epistemico":"fato","origem":"wiki","rel":"relaciona"},{"alvo":"word","confianca":0.5,"epistemico":"fato","origem":"wiki","rel":"relaciona"},{"alvo":"teorema_de_bayes","confianca":1.0,"epistemico":"fato","origem":"wiki","rel":"contradiz"}],"confianca":1.0,"contraexemplos":[],"descricao":"Kahneman: Sistema 1 (rápido, automático, heurístico) vs Sistema 2 (lento, deliberado); os vieses cognitivos nascem do conflito. Falsável.","epistemico":"fato","exemplos":[],"id":"dois_sistemas_kahneman","memoria":"semantica","meta":{"autor":"Kahneman","dominio":"ciencias_de_fora","fonte":""},"origem":"wiki","palavras_chave":[],"sinonimos":[],"tipo":"conceito","titulo":"Dois Sistemas (Kahneman)"}
\section{Dois Sistemas (Kahneman)}
Kahneman: Sistema 1 (rápido, automático, heurístico) vs Sistema 2 (lento, deliberado); os vieses cognitivos nascem do conflito. Falsável.
\prov{relações: parte\_de psicologia; relaciona aparelho\_psiquico; relaciona word; contradiz teorema\_de\_bayes}
\prov{proveniência: wiki \textperiodcentered{} fato \textperiodcentered{} confiança 1.0 \textperiodcentered{} domínio ciencias\_de\_fora \textperiodcentered{} história 7 versões no jornal}

%CRISTAL {"arestas":[{"alvo":"filosofia_da_medicina","confianca":1.0,"epistemico":"fato","origem":"wiki","rel":"parte_de"},{"alvo":"mente_programacao_cerebro_computador","confianca":1.0,"epistemico":"fato","origem":"wiki","rel":"usa"},{"alvo":"aparelho_psiquico","confianca":1.0,"epistemico":"fato","origem":"wiki","rel":"relaciona"},{"alvo":"consciencia","confianca":1.0,"epistemico":"fato","origem":"wiki","rel":"relaciona"}],"confianca":1.0,"contraexemplos":[],"descricao":"Lopes, endossando o budismo: 'a dor é uma experiência física e o sofrimento é uma experiência mental — a percepção de que as coisas deveriam ser diferentes do que são. O antídoto é a aceitação.' Separa o dano corporal do sofrimento produzido pela mente.","epistemico":"inferencia","exemplos":[],"id":"dor_vs_sofrimento_gentil","memoria":"semantica","meta":{"autor":"Gentil Lopes","dominio":"ciencias_de_fora","fonte":""},"origem":"wiki","palavras_chave":[],"sinonimos":[],"tipo":"conceito","titulo":"Dor (física) vs Sofrimento (mental) (Lopes)"}
\section{Dor (física) vs Sofrimento (mental) (Lopes)}
Lopes, endossando o budismo: 'a dor é uma experiência física e o sofrimento é uma experiência mental — a percepção de que as coisas deveriam ser diferentes do que são. O antídoto é a aceitação.' Separa o dano corporal do sofrimento produzido pela mente.
\prov{relações: parte\_de filosofia\_da\_medicina; usa mente\_programacao\_cerebro\_computador; relaciona aparelho\_psiquico; relaciona consciencia}
\prov{proveniência: wiki \textperiodcentered{} inferencia \textperiodcentered{} confiança 1.0 \textperiodcentered{} domínio ciencias\_de\_fora \textperiodcentered{} história 6 versões no jornal}

%CRISTAL {"arestas":[{"alvo":"saude_publica","confianca":1.0,"epistemico":"fato","origem":"wiki","rel":"parte_de"},{"alvo":"aparelho_psiquico","confianca":1.0,"epistemico":"fato","origem":"wiki","rel":"relaciona"},{"alvo":"historicismo_popper","confianca":1.0,"epistemico":"fato","origem":"wiki","rel":"relaciona"},{"alvo":"medicalizacao","confianca":1.0,"epistemico":"fato","origem":"wiki","rel":"relaciona"}],"confianca":1.0,"contraexemplos":[],"descricao":"O DSM classifica transtornos por sintomas, garantindo confiabilidade (concordância entre avaliadores) mas com validade discutida; Insel (NIMH) apontou sua 'falta de validade' biológica e lançou o RDoC — reacendendo o debate entre psiquiatria biológica e crítica ao modelo diagnóstico.","epistemico":"inferencia","exemplos":[],"id":"dsm_critica","memoria":"semantica","meta":{"autor":"APA/Insel","dominio":"ciencias_de_fora","fonte":""},"origem":"wiki","palavras_chave":[],"sinonimos":[],"tipo":"conceito","titulo":"O DSM e a Crítica ao Diagnóstico Psiquiátrico"}
\section{O DSM e a Crítica ao Diagnóstico Psiquiátrico}
O DSM classifica transtornos por sintomas, garantindo confiabilidade (concordância entre avaliadores) mas com validade discutida; Insel (NIMH) apontou sua 'falta de validade' biológica e lançou o RDoC — reacendendo o debate entre psiquiatria biológica e crítica ao modelo diagnóstico.
\prov{relações: parte\_de saude\_publica; relaciona aparelho\_psiquico; relaciona historicismo\_popper; relaciona medicalizacao}
\prov{proveniência: wiki \textperiodcentered{} inferencia \textperiodcentered{} confiança 1.0 \textperiodcentered{} domínio ciencias\_de\_fora \textperiodcentered{} história 5 versões no jornal}

%CRISTAL {"arestas":[{"alvo":"filosofia_da_mente","confianca":1.0,"epistemico":"fato","origem":"wiki","rel":"parte_de"},{"alvo":"emergencia","confianca":0.5,"epistemico":"fato","origem":"wiki","rel":"relaciona"},{"alvo":"word","confianca":0.5,"epistemico":"fato","origem":"wiki","rel":"relaciona"},{"alvo":"virtualizacao","confianca":0.5,"epistemico":"fato","origem":"wiki","rel":"relaciona"}],"confianca":1.0,"contraexemplos":[],"descricao":"Descartes: duas substâncias — res cogitans (mente) e res extensa (corpo) — que interagem; o problema mente-corpo.","epistemico":"inferencia","exemplos":[],"id":"dualismo_descartes","memoria":"semantica","meta":{"autor":"Descartes","dominio":"filosofia","fonte":""},"origem":"wiki","palavras_chave":[],"sinonimos":[],"tipo":"conceito","titulo":"Dualismo Cartesiano"}
\section{Dualismo Cartesiano}
Descartes: duas substâncias — res cogitans (mente) e res extensa (corpo) — que interagem; o problema mente-corpo.
\prov{relações: parte\_de filosofia\_da\_mente; relaciona emergencia; relaciona word; relaciona virtualizacao}
\prov{proveniência: wiki \textperiodcentered{} inferencia \textperiodcentered{} confiança 1.0 \textperiodcentered{} domínio filosofia \textperiodcentered{} história 5 versões no jornal}

%CRISTAL {"arestas":[{"alvo":"filosofia_do_direito","confianca":1.0,"epistemico":"fato","origem":"wiki","rel":"parte_de"},{"alvo":"juspositivismo","confianca":1.0,"epistemico":"fato","origem":"wiki","rel":"contradiz"},{"alvo":"hart_textura_aberta","confianca":1.0,"epistemico":"fato","origem":"wiki","rel":"contradiz"}],"confianca":1.0,"contraexemplos":[],"descricao":"Dworkin: o direito inclui não só regras mas princípios morais com peso; interpretar é reconstruir o direito como obra coerente de um único autor (a comunidade) na sua melhor luz moral. Um juiz ideal ('Hércules') encontra a única resposta correta — não a cria.","epistemico":"inferencia","exemplos":[],"id":"dworkin_integridade","memoria":"semantica","meta":{"autor":"Dworkin","dominio":"ciencias_de_fora","fonte":""},"origem":"wiki","palavras_chave":[],"sinonimos":[],"tipo":"conceito","titulo":"Direito como Integridade (Dworkin)"}
\section{Direito como Integridade (Dworkin)}
Dworkin: o direito inclui não só regras mas princípios morais com peso; interpretar é reconstruir o direito como obra coerente de um único autor (a comunidade) na sua melhor luz moral. Um juiz ideal ('Hércules') encontra a única resposta correta — não a cria.
\prov{relações: parte\_de filosofia\_do\_direito; contradiz juspositivismo; contradiz hart\_textura\_aberta}
\prov{proveniência: wiki \textperiodcentered{} inferencia \textperiodcentered{} confiança 1.0 \textperiodcentered{} domínio ciencias\_de\_fora \textperiodcentered{} história 4 versões no jornal}

%CRISTAL {"arestas":[{"alvo":"nao_dualidade","confianca":1.0,"epistemico":"fato","origem":"wiki","rel":"especializa"},{"alvo":"consciencia_pura","confianca":1.0,"epistemico":"fato","origem":"wiki","rel":"relaciona"},{"alvo":"contemplacao","confianca":1.0,"epistemico":"fato","origem":"wiki","rel":"usa"}],"confianca":1.0,"contraexemplos":[],"descricao":"Budismo tibetano: Rigpa é a consciência primordial, pura, não-dual — o 'sol sem nuvens' que os pensamentos (Sem) encobrem. O Lopes diz conhecê-la de 30+ anos de meditação.","epistemico":"hipotese","exemplos":[],"id":"dzogchen_rigpa","memoria":"semantica","meta":{"autor":"Nyingma","dominio":"filosofia","fonte":""},"origem":"wiki","palavras_chave":[],"sinonimos":[],"tipo":"conceito","titulo":"Dzogchen / Rigpa"}
\section{Dzogchen / Rigpa}
Budismo tibetano: Rigpa é a consciência primordial, pura, não-dual — o 'sol sem nuvens' que os pensamentos (Sem) encobrem. O Lopes diz conhecê-la de 30+ anos de meditação.
\prov{relações: especializa nao\_dualidade; relaciona consciencia\_pura; usa contemplacao}
\prov{proveniência: wiki \textperiodcentered{} hipotese \textperiodcentered{} confiança 1.0 \textperiodcentered{} domínio filosofia \textperiodcentered{} história 5 versões no jornal}

%CRISTAL {"arestas":[{"alvo":"biologia","confianca":1.0,"epistemico":"fato","origem":"wiki","rel":"parte_de"}],"confianca":1.0,"contraexemplos":[],"descricao":"O estudo das relações entre os organismos e o ambiente — o vivo como sistema, não como indivíduo.","epistemico":"fato","exemplos":[],"id":"ecologia","memoria":"semantica","meta":{"dominio":"ciencias_de_fora","fonte":""},"origem":"wiki","palavras_chave":[],"sinonimos":[],"tipo":"conceito","titulo":"Ecologia"}
\section{Ecologia}
O estudo das relações entre os organismos e o ambiente — o vivo como sistema, não como indivíduo.
\prov{relações: parte\_de biologia}
\prov{proveniência: wiki \textperiodcentered{} fato \textperiodcentered{} confiança 1.0 \textperiodcentered{} domínio ciencias\_de\_fora \textperiodcentered{} história 2 versões no jornal}

%CRISTAL {"arestas":[{"alvo":"matematica","confianca":1.0,"epistemico":"fato","origem":"wiki","rel":"referencia"},{"alvo":"utilitarismo","confianca":0.5,"epistemico":"fato","origem":"wiki","rel":"relaciona"},{"alvo":"painel_estelar","confianca":0.5,"epistemico":"fato","origem":"wiki","rel":"relaciona"},{"alvo":"virtualizacao","confianca":0.5,"epistemico":"fato","origem":"wiki","rel":"relaciona"},{"alvo":"word","confianca":0.5,"epistemico":"fato","origem":"wiki","rel":"relaciona"}],"confianca":1.0,"contraexemplos":[],"descricao":"A ciência da alocação de recursos escassos — valor, troca, produção, coordenação.","epistemico":"fato","exemplos":[],"id":"economia","memoria":"semantica","meta":{"dominio":"ciencias_de_fora","fonte":""},"origem":"wiki","palavras_chave":[],"sinonimos":[],"tipo":"conceito","titulo":"Economia"}
\section{Economia}
A ciência da alocação de recursos escassos — valor, troca, produção, coordenação.
\prov{relações: referencia matematica; relaciona utilitarismo; relaciona painel\_estelar; relaciona virtualizacao; relaciona word}
\prov{proveniência: wiki \textperiodcentered{} fato \textperiodcentered{} confiança 1.0 \textperiodcentered{} domínio ciencias\_de\_fora \textperiodcentered{} história 6 versões no jornal}

%CRISTAL {"arestas":[{"alvo":"comunicacao_digital","confianca":1.0,"epistemico":"fato","origem":"wiki","rel":"parte_de"},{"alvo":"escassez","confianca":1.0,"epistemico":"fato","origem":"wiki","rel":"usa"},{"alvo":"capitalismo_de_vigilancia","confianca":1.0,"epistemico":"fato","origem":"wiki","rel":"relaciona"}],"confianca":1.0,"contraexemplos":[],"descricao":"Herbert Simon: numa era de abundância informacional, a atenção é o recurso escasso — 'a riqueza da informação cria uma pobreza de atenção'; a atenção torna-se a mercadoria disputada pelas plataformas.","epistemico":"inferencia","exemplos":[],"id":"economia_da_atencao","memoria":"semantica","meta":{"autor":"Simon","dominio":"ciencias_de_fora","fonte":""},"origem":"wiki","palavras_chave":[],"sinonimos":[],"tipo":"conceito","titulo":"Economia da Atenção (Simon)"}
\section{Economia da Atenção (Simon)}
Herbert Simon: numa era de abundância informacional, a atenção é o recurso escasso — 'a riqueza da informação cria uma pobreza de atenção'; a atenção torna-se a mercadoria disputada pelas plataformas.
\prov{relações: parte\_de comunicacao\_digital; usa escassez; relaciona capitalismo\_de\_vigilancia}
\prov{proveniência: wiki \textperiodcentered{} inferencia \textperiodcentered{} confiança 1.0 \textperiodcentered{} domínio ciencias\_de\_fora \textperiodcentered{} história 4 versões no jornal}

%CRISTAL {"arestas":[{"alvo":"ecologia","confianca":1.0,"epistemico":"fato","origem":"wiki","rel":"parte_de"}],"confianca":1.0,"contraexemplos":[],"descricao":"A comunidade de organismos mais o ambiente físico, ligados por fluxos de energia e matéria — uma unidade funcional.","epistemico":"fato","exemplos":[],"id":"ecossistema","memoria":"semantica","meta":{"dominio":"ciencias_de_fora","fonte":""},"origem":"wiki","palavras_chave":[],"sinonimos":[],"tipo":"conceito","titulo":"Ecossistema"}
\section{Ecossistema}
A comunidade de organismos mais o ambiente físico, ligados por fluxos de energia e matéria — uma unidade funcional.
\prov{relações: parte\_de ecologia}
\prov{proveniência: wiki \textperiodcentered{} fato \textperiodcentered{} confiança 1.0 \textperiodcentered{} domínio ciencias\_de\_fora \textperiodcentered{} história 2 versões no jornal}

%CRISTAL {"arestas":[{"alvo":"enhancement_humano","confianca":1.0,"epistemico":"fato","origem":"wiki","rel":"usa"},{"alvo":"principio_responsabilidade","confianca":1.0,"epistemico":"fato","origem":"wiki","rel":"usa"},{"alvo":"selecao_natural","confianca":1.0,"epistemico":"fato","origem":"wiki","rel":"relaciona"}],"confianca":1.0,"contraexemplos":[],"descricao":"O uso do CRISPR-Cas9 para alterar o genoma: consenso para edição somática terapêutica, mas forte controvérsia e proibição da edição germinativa (embriões/gametas), por afetar gerações futuras que não consentem e por risco de eugenia e discriminação genética (caso He Jiankui, 2018).","epistemico":"inferencia","exemplos":[],"id":"edicao_germinativa_crispr","memoria":"semantica","meta":{"autor":"—","dominio":"ciencias_de_fora","fonte":""},"origem":"wiki","palavras_chave":[],"sinonimos":[],"tipo":"conceito","titulo":"Edição Genética Germinativa (CRISPR)"}
\section{Edição Genética Germinativa (CRISPR)}
O uso do CRISPR-Cas9 para alterar o genoma: consenso para edição somática terapêutica, mas forte controvérsia e proibição da edição germinativa (embriões/gametas), por afetar gerações futuras que não consentem e por risco de eugenia e discriminação genética (caso He Jiankui, 2018).
\prov{relações: usa enhancement\_humano; usa principio\_responsabilidade; relaciona selecao\_natural}
\prov{proveniência: wiki \textperiodcentered{} inferencia \textperiodcentered{} confiança 1.0 \textperiodcentered{} domínio ciencias\_de\_fora \textperiodcentered{} história 4 versões no jornal}

%CRISTAL {"arestas":[{"alvo":"filosofia","confianca":1.0,"epistemico":"fato","origem":"wiki","rel":"relaciona"},{"alvo":"virtualizacao","confianca":0.5,"epistemico":"fato","origem":"wiki","rel":"relaciona"},{"alvo":"vida-morte","confianca":0.5,"epistemico":"fato","origem":"wiki","rel":"relaciona"},{"alvo":"word","confianca":0.5,"epistemico":"fato","origem":"wiki","rel":"relaciona"},{"alvo":"readme","confianca":0.5,"epistemico":"fato","origem":"wiki","rel":"relaciona"},{"alvo":"sequencia_pow_nativa","confianca":0.5,"epistemico":"fato","origem":"wiki","rel":"relaciona"},{"alvo":"pell","confianca":0.5,"epistemico":"fato","origem":"wiki","rel":"relaciona"},{"alvo":"pipeline_de_ingestao","confianca":0.5,"epistemico":"fato","origem":"wiki","rel":"relaciona"},{"alvo":"gentil_12_definicao","confianca":0.5,"epistemico":"fato","origem":"wiki","rel":"relaciona"},{"alvo":"telemetria_shadow_producao","confianca":0.5,"epistemico":"fato","origem":"wiki","rel":"relaciona"},{"alvo":"pow_taxa_de_sucesso_do_bit_broca","confianca":0.5,"epistemico":"fato","origem":"wiki","rel":"relaciona"},{"alvo":"literatura","confianca":0.5,"epistemico":"fato","origem":"wiki","rel":"relaciona"},{"alvo":"vetores","confianca":0.5,"epistemico":"fato","origem":"wiki","rel":"relaciona"}],"confianca":1.0,"contraexemplos":[],"descricao":"A teoria e a prática de formar e instruir seres humanos — o campo que estuda como se ensina, como se aprende e para quê.","epistemico":"fato","exemplos":[],"id":"educacao","memoria":"semantica","meta":{"dominio":"ciencias_de_fora","fonte":""},"origem":"wiki","palavras_chave":[],"sinonimos":[],"tipo":"conceito","titulo":"Educação"}
\section{Educação}
A teoria e a prática de formar e instruir seres humanos — o campo que estuda como se ensina, como se aprende e para quê.
\prov{relações: relaciona filosofia; relaciona virtualizacao; relaciona vida-morte; relaciona word; relaciona readme; relaciona sequencia\_pow\_nativa; relaciona pell; relaciona pipeline\_de\_ingestao; relaciona gentil\_12\_definicao; relaciona telemetria\_shadow\_producao; relaciona pow\_taxa\_de\_sucesso\_do\_bit\_broca; relaciona literatura; relaciona vetores}
\prov{proveniência: wiki \textperiodcentered{} fato \textperiodcentered{} confiança 1.0 \textperiodcentered{} domínio ciencias\_de\_fora \textperiodcentered{} história 14 versões no jornal}

%CRISTAL {"arestas":[{"alvo":"filosofia_da_educacao","confianca":1.0,"epistemico":"fato","origem":"wiki","rel":"parte_de"},{"alvo":"idiotismo_nao_examinar","confianca":1.0,"epistemico":"fato","origem":"wiki","rel":"relaciona"},{"alvo":"capital_cultural","confianca":1.0,"epistemico":"fato","origem":"wiki","rel":"relaciona"}],"confianca":1.0,"contraexemplos":[],"descricao":"Freire: o ensino como ato de depositar conteúdos numa consciência passiva — o professor narra, o aluno recebe e repete, tratado como recipiente vazio. Naturaliza a relação opressor-oprimido e nega o diálogo.","epistemico":"inferencia","exemplos":[],"id":"educacao_bancaria","memoria":"semantica","meta":{"autor":"Freire","dominio":"ciencias_de_fora","fonte":""},"origem":"wiki","palavras_chave":[],"sinonimos":[],"tipo":"conceito","titulo":"Educação Bancária"}
\section{Educação Bancária}
Freire: o ensino como ato de depositar conteúdos numa consciência passiva — o professor narra, o aluno recebe e repete, tratado como recipiente vazio. Naturaliza a relação opressor-oprimido e nega o diálogo.
\prov{relações: parte\_de filosofia\_da\_educacao; relaciona idiotismo\_nao\_examinar; relaciona capital\_cultural}
\prov{proveniência: wiki \textperiodcentered{} inferencia \textperiodcentered{} confiança 1.0 \textperiodcentered{} domínio ciencias\_de\_fora \textperiodcentered{} história 4 versões no jornal}

%CRISTAL {"arestas":[{"alvo":"filosofia_da_educacao","confianca":1.0,"epistemico":"fato","origem":"wiki","rel":"parte_de"},{"alvo":"idiotismo_nao_examinar","confianca":1.0,"epistemico":"fato","origem":"wiki","rel":"relaciona"},{"alvo":"dewey_democracia_educacao","confianca":1.0,"epistemico":"fato","origem":"wiki","rel":"relaciona"}],"confianca":1.0,"contraexemplos":[],"descricao":"Nussbaum: a democracia depende das artes e humanidades para cultivar pensamento crítico, argumentação socrática e 'imaginação narrativa' (pôr-se no lugar do outro); alerta contra a educação reduzida a fins econômicos.","epistemico":"inferencia","exemplos":[],"id":"educacao_cidada_nussbaum","memoria":"semantica","meta":{"autor":"Nussbaum","dominio":"ciencias_de_fora","fonte":""},"origem":"wiki","palavras_chave":[],"sinonimos":[],"tipo":"conceito","titulo":"Humanidades para a Cidadania (Nussbaum)"}
\section{Humanidades para a Cidadania (Nussbaum)}
Nussbaum: a democracia depende das artes e humanidades para cultivar pensamento crítico, argumentação socrática e 'imaginação narrativa' (pôr-se no lugar do outro); alerta contra a educação reduzida a fins econômicos.
\prov{relações: parte\_de filosofia\_da\_educacao; relaciona idiotismo\_nao\_examinar; relaciona dewey\_democracia\_educacao}
\prov{proveniência: wiki \textperiodcentered{} inferencia \textperiodcentered{} confiança 1.0 \textperiodcentered{} domínio ciencias\_de\_fora \textperiodcentered{} história 4 versões no jornal}

%CRISTAL {"arestas":[{"alvo":"filosofia_da_educacao","confianca":1.0,"epistemico":"fato","origem":"wiki","rel":"parte_de"},{"alvo":"contrato_social","confianca":1.0,"epistemico":"fato","origem":"wiki","rel":"relaciona"}],"confianca":1.0,"contraexemplos":[],"descricao":"Rousseau (Emílio): a criança nasce boa e é corrompida pela sociedade; a educação deve respeitar as fases naturais e partir da experiência da criança ('educação negativa'), protegendo-a da precipitação adulta.","epistemico":"inferencia","exemplos":[],"id":"educacao_natural_rousseau","memoria":"semantica","meta":{"autor":"Rousseau","dominio":"ciencias_de_fora","fonte":""},"origem":"wiki","palavras_chave":[],"sinonimos":[],"tipo":"conceito","titulo":"Educação Natural (Rousseau)"}
\section{Educação Natural (Rousseau)}
Rousseau (Emílio): a criança nasce boa e é corrompida pela sociedade; a educação deve respeitar as fases naturais e partir da experiência da criança ('educação negativa'), protegendo-a da precipitação adulta.
\prov{relações: parte\_de filosofia\_da\_educacao; relaciona contrato\_social}
\prov{proveniência: wiki \textperiodcentered{} inferencia \textperiodcentered{} confiança 1.0 \textperiodcentered{} domínio ciencias\_de\_fora \textperiodcentered{} história 3 versões no jornal}

%CRISTAL {"arestas":[{"alvo":"educacao_bancaria","confianca":1.0,"epistemico":"fato","origem":"wiki","rel":"contradiz"},{"alvo":"filosofia_da_educacao","confianca":1.0,"epistemico":"fato","origem":"wiki","rel":"parte_de"}],"confianca":1.0,"contraexemplos":[],"descricao":"Freire: educador e educando investigam juntos a realidade, problematizando-a para transformá-la; ambos são sujeitos que aprendem em relação dialógica, não polos de saber e ignorância.","epistemico":"inferencia","exemplos":[],"id":"educacao_problematizadora","memoria":"semantica","meta":{"autor":"Freire","dominio":"ciencias_de_fora","fonte":""},"origem":"wiki","palavras_chave":[],"sinonimos":[],"tipo":"conceito","titulo":"Educação Problematizadora / Libertadora"}
\section{Educação Problematizadora / Libertadora}
Freire: educador e educando investigam juntos a realidade, problematizando-a para transformá-la; ambos são sujeitos que aprendem em relação dialógica, não polos de saber e ignorância.
\prov{relações: contradiz educacao\_bancaria; parte\_de filosofia\_da\_educacao}
\prov{proveniência: wiki \textperiodcentered{} inferencia \textperiodcentered{} confiança 1.0 \textperiodcentered{} domínio ciencias\_de\_fora \textperiodcentered{} história 3 versões no jornal}

%CRISTAL {"arestas":[{"alvo":"aparelho_psiquico","confianca":1.0,"epistemico":"fato","origem":"wiki","rel":"relaciona"},{"alvo":"dois_sistemas_kahneman","confianca":1.0,"epistemico":"fato","origem":"wiki","rel":"relaciona"}],"confianca":1.0,"contraexemplos":[],"descricao":"O placebo é substância inerte que ainda assim produz efeitos por expectativa; o cegamento duplo (paciente e avaliador ignoram a alocação) neutraliza esse efeito e o viés do observador, isolando o efeito real da intervenção.","epistemico":"fato","exemplos":[],"id":"efeito_placebo_duplo_cego","memoria":"semantica","meta":{"autor":"—","dominio":"ciencias_de_fora","fonte":""},"origem":"wiki","palavras_chave":[],"sinonimos":[],"tipo":"conceito","titulo":"Efeito Placebo e o Duplo-Cego"}
\section{Efeito Placebo e o Duplo-Cego}
O placebo é substância inerte que ainda assim produz efeitos por expectativa; o cegamento duplo (paciente e avaliador ignoram a alocação) neutraliza esse efeito e o viés do observador, isolando o efeito real da intervenção.
\prov{relações: relaciona aparelho\_psiquico; relaciona dois\_sistemas\_kahneman}
\prov{proveniência: wiki \textperiodcentered{} fato \textperiodcentered{} confiança 1.0 \textperiodcentered{} domínio ciencias\_de\_fora \textperiodcentered{} história 3 versões no jornal}

%CRISTAL {"arestas":[{"alvo":"mao_invisivel","confianca":1.0,"epistemico":"fato","origem":"wiki","rel":"usa"},{"alvo":"teoria_dos_jogos","confianca":1.0,"epistemico":"fato","origem":"wiki","rel":"usa"},{"alvo":"dois_sistemas_kahneman","confianca":1.0,"epistemico":"fato","origem":"wiki","rel":"relaciona"}],"confianca":1.0,"contraexemplos":[],"descricao":"Downs: tratando partidos como empresas que maximizam votos e eleitores como consumidores racionais, os partidos convergem para a posição do eleitor mediano. E como o voto individual quase nunca é decisivo, é racional o eleitor não se informar (a 'ignorância racional').","epistemico":"inferencia","exemplos":[],"id":"eleitor_mediano_downs","memoria":"semantica","meta":{"autor":"Downs","dominio":"ciencias_de_fora","fonte":""},"origem":"wiki","palavras_chave":[],"sinonimos":[],"tipo":"conceito","titulo":"Teorema do Eleitor Mediano (Downs)"}
\section{Teorema do Eleitor Mediano (Downs)}
Downs: tratando partidos como empresas que maximizam votos e eleitores como consumidores racionais, os partidos convergem para a posição do eleitor mediano. E como o voto individual quase nunca é decisivo, é racional o eleitor não se informar (a 'ignorância racional').
\prov{relações: usa mao\_invisivel; usa teoria\_dos\_jogos; relaciona dois\_sistemas\_kahneman}
\prov{proveniência: wiki \textperiodcentered{} inferencia \textperiodcentered{} confiança 1.0 \textperiodcentered{} domínio ciencias\_de\_fora \textperiodcentered{} história 4 versões no jornal}

%CRISTAL {"arestas":[{"alvo":"musica","confianca":1.0,"epistemico":"fato","origem":"wiki","rel":"parte_de"}],"confianca":1.0,"contraexemplos":[],"descricao":"Os três elementos da música: melodia (sucessão de alturas no tempo — a dimensão horizontal), harmonia (combinação simultânea de notas — a dimensão vertical) e ritmo (organização dos sons e silêncios, o pulso e o movimento).","epistemico":"inferencia","exemplos":[],"id":"elementos_musica","memoria":"semantica","meta":{"autor":"—","dominio":"ciencias_de_fora","fonte":""},"origem":"wiki","palavras_chave":[],"sinonimos":[],"tipo":"conceito","titulo":"Melodia, Harmonia e Ritmo"}
\section{Melodia, Harmonia e Ritmo}
Os três elementos da música: melodia (sucessão de alturas no tempo — a dimensão horizontal), harmonia (combinação simultânea de notas — a dimensão vertical) e ritmo (organização dos sons e silêncios, o pulso e o movimento).
\prov{relações: parte\_de musica}
\prov{proveniência: wiki \textperiodcentered{} inferencia \textperiodcentered{} confiança 1.0 \textperiodcentered{} domínio ciencias\_de\_fora \textperiodcentered{} história 2 versões no jornal}

%CRISTAL {"arestas":[{"alvo":"ellul_la_technique","confianca":1.0,"epistemico":"fato","origem":"wiki","rel":"usa"},{"alvo":"maquina_evolutiva","confianca":1.0,"epistemico":"fato","origem":"wiki","rel":"relaciona"}],"confianca":1.0,"contraexemplos":[],"descricao":"Ellul: a Técnica se autodetermina — cresce por sua própria lógica de eficiência, independente da vontade humana ou da moral, tornando-se juíza de sua própria normatividade. O homem é que se adapta a ela, não o contrário.","epistemico":"inferencia","exemplos":[],"id":"ellul_autonomia_tecnica","memoria":"semantica","meta":{"autor":"Ellul","dominio":"ciencias_de_fora","fonte":""},"origem":"wiki","palavras_chave":[],"sinonimos":[],"tipo":"conceito","titulo":"Autonomia da Técnica (Ellul)"}
\section{Autonomia da Técnica (Ellul)}
Ellul: a Técnica se autodetermina — cresce por sua própria lógica de eficiência, independente da vontade humana ou da moral, tornando-se juíza de sua própria normatividade. O homem é que se adapta a ela, não o contrário.
\prov{relações: usa ellul\_la\_technique; relaciona maquina\_evolutiva}
\prov{proveniência: wiki \textperiodcentered{} inferencia \textperiodcentered{} confiança 1.0 \textperiodcentered{} domínio ciencias\_de\_fora \textperiodcentered{} história 3 versões no jornal}

%CRISTAL {"arestas":[{"alvo":"filosofia_da_tecnica","confianca":1.0,"epistemico":"fato","origem":"wiki","rel":"parte_de"},{"alvo":"acao_social_weber","confianca":1.0,"epistemico":"fato","origem":"wiki","rel":"relaciona"}],"confianca":1.0,"contraexemplos":[],"descricao":"Ellul: a Técnica é 'a totalidade dos métodos racionalmente obtidos e de eficiência absoluta' em todo campo humano; a máquina é só seu sintoma. A eficiência vira o valor supremo que absorve economia, política, cultura e vida.","epistemico":"inferencia","exemplos":[],"id":"ellul_la_technique","memoria":"semantica","meta":{"autor":"Ellul","dominio":"ciencias_de_fora","fonte":""},"origem":"wiki","palavras_chave":[],"sinonimos":[],"tipo":"conceito","titulo":"La Technique (Ellul)"}
\section{La Technique (Ellul)}
Ellul: a Técnica é 'a totalidade dos métodos racionalmente obtidos e de eficiência absoluta' em todo campo humano; a máquina é só seu sintoma. A eficiência vira o valor supremo que absorve economia, política, cultura e vida.
\prov{relações: parte\_de filosofia\_da\_tecnica; relaciona acao\_social\_weber}
\prov{proveniência: wiki \textperiodcentered{} inferencia \textperiodcentered{} confiança 1.0 \textperiodcentered{} domínio ciencias\_de\_fora \textperiodcentered{} história 3 versões no jornal}

%CRISTAL {"arestas":[{"alvo":"embedding","confianca":1.0,"epistemico":"fato","origem":"wiki","rel":"usa"},{"alvo":"interlingua","confianca":1.0,"epistemico":"fato","origem":"wiki","rel":"explica"},{"alvo":"traducao_neural","confianca":1.0,"epistemico":"fato","origem":"wiki","rel":"usa"},{"alvo":"tiffany","confianca":1.0,"epistemico":"fato","origem":"wiki","rel":"relaciona"},{"alvo":"significado_como_invariante","confianca":1.0,"epistemico":"fato","origem":"wiki","rel":"usa"}],"confianca":1.0,"contraexemplos":[],"descricao":"Nos modelos neurais, a interlíngua deixou de ser sonho e virou GEOMETRIA: um espaço de embedding compartilhado onde 'gato', 'cat' e '猫' caem perto — o significado como vetor, o invariante COMPUTÁVEL. A mesma ideia da recuperação por sentido da Tiffany.","epistemico":"inferencia","exemplos":[],"id":"embedding_como_interlingua","memoria":"semantica","meta":{"dominio":"ciencias_de_fora","fonte":""},"origem":"wiki","palavras_chave":[],"sinonimos":[],"tipo":"conceito","titulo":"Embedding como Interlíngua"}
\section{Embedding como Interlíngua}
Nos modelos neurais, a interlíngua deixou de ser sonho e virou GEOMETRIA: um espaço de embedding compartilhado onde 'gato', 'cat' e '\texttt{U+732B}' caem perto — o significado como vetor, o invariante COMPUTÁVEL. A mesma ideia da recuperação por sentido da Tiffany.
\prov{relações: usa embedding; explica interlingua; usa traducao\_neural; relaciona tiffany; usa significado\_como\_invariante}
\prov{proveniência: wiki \textperiodcentered{} inferencia \textperiodcentered{} confiança 1.0 \textperiodcentered{} domínio ciencias\_de\_fora \textperiodcentered{} história 6 versões no jornal}

%CRISTAL {"arestas":[{"alvo":"filosofia_da_mente","confianca":1.0,"epistemico":"fato","origem":"wiki","rel":"referencia"},{"alvo":"pow_metal","confianca":0.5,"epistemico":"fato","origem":"wiki","rel":"relaciona"},{"alvo":"virtualizacao","confianca":0.5,"epistemico":"fato","origem":"wiki","rel":"relaciona"},{"alvo":"word","confianca":0.5,"epistemico":"fato","origem":"wiki","rel":"relaciona"},{"alvo":"vita_contemplativa","confianca":0.5,"epistemico":"fato","origem":"wiki","rel":"relaciona"}],"confianca":1.0,"contraexemplos":[],"descricao":"Propriedades sistêmicas (como a consciência) que não se reduzem às partes isoladas — o todo faz o que a soma não faz.","epistemico":"inferencia","exemplos":[],"id":"emergencia","memoria":"semantica","meta":{"dominio":"filosofia","fonte":"web: SEP"},"origem":"wiki","palavras_chave":[],"sinonimos":[],"tipo":"conceito","titulo":"Emergência"}
\section{Emergência}
Propriedades sistêmicas (como a consciência) que não se reduzem às partes isoladas — o todo faz o que a soma não faz.
\prov{relações: referencia filosofia\_da\_mente; relaciona pow\_metal; relaciona virtualizacao; relaciona word; relaciona vita\_contemplativa}
\prov{proveniência: wiki \textperiodcentered{} web: SEP \textperiodcentered{} inferencia \textperiodcentered{} confiança 1.0 \textperiodcentered{} domínio filosofia \textperiodcentered{} história 6 versões no jornal}

%CRISTAL {"arestas":[{"alvo":"autopoiese","confianca":1.0,"epistemico":"fato","origem":"wiki","rel":"usa"},{"alvo":"consciencia","confianca":1.0,"epistemico":"fato","origem":"wiki","rel":"referencia"},{"alvo":"word","confianca":0.5,"epistemico":"fato","origem":"wiki","rel":"relaciona"},{"alvo":"while","confianca":0.5,"epistemico":"fato","origem":"wiki","rel":"relaciona"}],"confianca":1.0,"contraexemplos":[],"descricao":"Varela/Thompson/Rosch: conhecer é agir — a mente é incorporada e acoplada ao ambiente, não uma representação abstrata.","epistemico":"inferencia","exemplos":[],"id":"enacao","memoria":"semantica","meta":{"autor":"Varela","dominio":"ciencias_de_fora","fonte":""},"origem":"wiki","palavras_chave":[],"sinonimos":[],"tipo":"conceito","titulo":"Enação (Cognição Incorporada)"}
\section{Enação (Cognição Incorporada)}
Varela/Thompson/Rosch: conhecer é agir — a mente é incorporada e acoplada ao ambiente, não uma representação abstrata.
\prov{relações: usa autopoiese; referencia consciencia; relaciona word; relaciona while}
\prov{proveniência: wiki \textperiodcentered{} inferencia \textperiodcentered{} confiança 1.0 \textperiodcentered{} domínio ciencias\_de\_fora \textperiodcentered{} história 5 versões no jornal}

%CRISTAL {"arestas":[{"alvo":"bioetica","confianca":1.0,"epistemico":"fato","origem":"wiki","rel":"parte_de"},{"alvo":"transumanismo","confianca":1.0,"epistemico":"fato","origem":"wiki","rel":"relaciona"},{"alvo":"principio_responsabilidade","confianca":1.0,"epistemico":"fato","origem":"wiki","rel":"usa"}],"confianca":1.0,"contraexemplos":[],"descricao":"O debate sobre a linha entre tratar/prevenir doença (terapia) e ampliar capacidades além do normal (aprimoramento — força, cognição); críticos temem novas desigualdades e 'eugenia' (Sandel/Kass), defensores invocam autonomia (Savulescu).","epistemico":"inferencia","exemplos":[],"id":"enhancement_humano","memoria":"semantica","meta":{"autor":"Savulescu vs Sandel","dominio":"ciencias_de_fora","fonte":""},"origem":"wiki","palavras_chave":[],"sinonimos":[],"tipo":"conceito","titulo":"Melhoramento Humano (Terapia vs Aprimoramento)"}
\section{Melhoramento Humano (Terapia vs Aprimoramento)}
O debate sobre a linha entre tratar/prevenir doença (terapia) e ampliar capacidades além do normal (aprimoramento — força, cognição); críticos temem novas desigualdades e 'eugenia' (Sandel/Kass), defensores invocam autonomia (Savulescu).
\prov{relações: parte\_de bioetica; relaciona transumanismo; usa principio\_responsabilidade}
\prov{proveniência: wiki \textperiodcentered{} inferencia \textperiodcentered{} confiança 1.0 \textperiodcentered{} domínio ciencias\_de\_fora \textperiodcentered{} história 4 versões no jornal}

%CRISTAL {"arestas":[{"alvo":"causalidade_bradford_hill","confianca":1.0,"epistemico":"fato","origem":"wiki","rel":"usa"},{"alvo":"efeito_placebo_duplo_cego","confianca":1.0,"epistemico":"fato","origem":"wiki","rel":"usa"}],"confianca":1.0,"contraexemplos":[],"descricao":"O estudo em que os participantes são alocados por sorteio a tratamento ou controle, equilibrando confundidores conhecidos e desconhecidos e permitindo inferência causal — o topo da hierarquia de evidências para eficácia. Raízes na estatística de Fisher (estreptomicina, 1948).","epistemico":"fato","exemplos":[],"id":"ensaio_clinico_randomizado","memoria":"semantica","meta":{"autor":"Fisher/Bradford Hill","dominio":"ciencias_de_fora","fonte":""},"origem":"wiki","palavras_chave":[],"sinonimos":[],"tipo":"conceito","titulo":"Ensaio Clínico Randomizado (RCT)"}
\section{Ensaio Clínico Randomizado (RCT)}
O estudo em que os participantes são alocados por sorteio a tratamento ou controle, equilibrando confundidores conhecidos e desconhecidos e permitindo inferência causal — o topo da hierarquia de evidências para eficácia. Raízes na estatística de Fisher (estreptomicina, 1948).
\prov{relações: usa causalidade\_bradford\_hill; usa efeito\_placebo\_duplo\_cego}
\prov{proveniência: wiki \textperiodcentered{} fato \textperiodcentered{} confiança 1.0 \textperiodcentered{} domínio ciencias\_de\_fora \textperiodcentered{} história 3 versões no jornal}

%CRISTAL {"arestas":[{"alvo":"semiotica","confianca":1.0,"epistemico":"fato","origem":"wiki","rel":"parte_de"},{"alvo":"tiffany","confianca":1.0,"epistemico":"fato","origem":"wiki","rel":"referencia"},{"alvo":"word","confianca":0.5,"epistemico":"fato","origem":"wiki","rel":"relaciona"},{"alvo":"while","confianca":0.5,"epistemico":"fato","origem":"wiki","rel":"relaciona"}],"confianca":1.0,"contraexemplos":[],"descricao":"Benveniste: o sujeito se apropria da língua no ato de enunciar; o 'eu', a dêixis (aqui/agora) e a modalização convertem a língua em discurso.","epistemico":"inferencia","exemplos":[],"id":"enunciacao_benveniste","memoria":"semantica","meta":{"autor":"Benveniste","dominio":"ciencias_de_fora","fonte":""},"origem":"wiki","palavras_chave":[],"sinonimos":[],"tipo":"conceito","titulo":"Enunciação e Subjetividade"}
\section{Enunciação e Subjetividade}
Benveniste: o sujeito se apropria da língua no ato de enunciar; o 'eu', a dêixis (aqui/agora) e a modalização convertem a língua em discurso.
\prov{relações: parte\_de semiotica; referencia tiffany; relaciona word; relaciona while}
\prov{proveniência: wiki \textperiodcentered{} inferencia \textperiodcentered{} confiança 1.0 \textperiodcentered{} domínio ciencias\_de\_fora \textperiodcentered{} história 5 versões no jornal}

%CRISTAL {"arestas":[{"alvo":"bioquimica","confianca":1.0,"epistemico":"fato","origem":"wiki","rel":"parte_de"}],"confianca":1.0,"contraexemplos":[],"descricao":"Uma proteína que acelera uma reação específica sem se consumir — o catalisador que torna a vida cineticamente possível.","epistemico":"fato","exemplos":[],"id":"enzima","memoria":"semantica","meta":{"dominio":"ciencias_de_fora","fonte":""},"origem":"wiki","palavras_chave":[],"sinonimos":[],"tipo":"conceito","titulo":"Enzima"}
\section{Enzima}
Uma proteína que acelera uma reação específica sem se consumir — o catalisador que torna a vida cineticamente possível.
\prov{relações: parte\_de bioquimica}
\prov{proveniência: wiki \textperiodcentered{} fato \textperiodcentered{} confiança 1.0 \textperiodcentered{} domínio ciencias\_de\_fora \textperiodcentered{} história 2 versões no jornal}

%CRISTAL {"arestas":[{"alvo":"saude_publica","confianca":1.0,"epistemico":"fato","origem":"wiki","rel":"parte_de"}],"confianca":1.0,"contraexemplos":[],"descricao":"O estudo da distribuição (quem, onde, quando) e dos determinantes de doenças em populações, base da prevenção. Distingue incidência (casos novos) de prevalência (casos existentes). John Snow, ao mapear a cólera na bomba de Broad Street (1854), fundou a epidemiologia espacial antes de o germe ser conhecido.","epistemico":"fato","exemplos":[],"id":"epidemiologia","memoria":"semantica","meta":{"autor":"Snow/Farr","dominio":"ciencias_de_fora","fonte":""},"origem":"wiki","palavras_chave":[],"sinonimos":[],"tipo":"conceito","titulo":"Epidemiologia"}
\section{Epidemiologia}
O estudo da distribuição (quem, onde, quando) e dos determinantes de doenças em populações, base da prevenção. Distingue incidência (casos novos) de prevalência (casos existentes). John Snow, ao mapear a cólera na bomba de Broad Street (1854), fundou a epidemiologia espacial antes de o germe ser conhecido.
\prov{relações: parte\_de saude\_publica}
\prov{proveniência: wiki \textperiodcentered{} fato \textperiodcentered{} confiança 1.0 \textperiodcentered{} domínio ciencias\_de\_fora \textperiodcentered{} história 2 versões no jornal}

%CRISTAL {"arestas":[{"alvo":"expressao_genica","confianca":1.0,"epistemico":"fato","origem":"wiki","rel":"especializa"}],"confianca":1.0,"contraexemplos":[],"descricao":"Marcas químicas sobre o DNA que ligam/desligam genes sem mudar a sequência — herança do uso, às vezes entre gerações.","epistemico":"fato","exemplos":[],"id":"epigenetica","memoria":"semantica","meta":{"dominio":"ciencias_de_fora","fonte":""},"origem":"wiki","palavras_chave":[],"sinonimos":[],"tipo":"conceito","titulo":"Epigenética"}
\section{Epigenética}
Marcas químicas sobre o DNA que ligam/desligam genes sem mudar a sequência — herança do uso, às vezes entre gerações.
\prov{relações: especializa expressao\_genica}
\prov{proveniência: wiki \textperiodcentered{} fato \textperiodcentered{} confiança 1.0 \textperiodcentered{} domínio ciencias\_de\_fora \textperiodcentered{} história 2 versões no jornal}

%CRISTAL {"arestas":[{"alvo":"filosofia","confianca":1.0,"epistemico":"fato","origem":"wiki","rel":"parte_de"},{"alvo":"metafisica","confianca":0.5,"epistemico":"fato","origem":"wiki","rel":"relaciona"},{"alvo":"word","confianca":0.5,"epistemico":"fato","origem":"wiki","rel":"relaciona"},{"alvo":"virtualizacao","confianca":0.5,"epistemico":"fato","origem":"wiki","rel":"relaciona"},{"alvo":"psicologia","confianca":0.5,"epistemico":"fato","origem":"wiki","rel":"relaciona"},{"alvo":"vida-morte","confianca":0.5,"epistemico":"fato","origem":"wiki","rel":"relaciona"},{"alvo":"semiotica_peirce_triade","confianca":0.5,"epistemico":"fato","origem":"wiki","rel":"relaciona"},{"alvo":"pow_coordenadas_espectrais_broca","confianca":0.5,"epistemico":"fato","origem":"wiki","rel":"relaciona"},{"alvo":"pow_bit_taxa","confianca":0.5,"epistemico":"fato","origem":"wiki","rel":"relaciona"}],"confianca":1.0,"contraexemplos":[],"descricao":"A teoria do conhecimento: o que podemos saber, como justificamos crenças, os limites do saber.","epistemico":"fato","exemplos":[],"id":"epistemologia","memoria":"semantica","meta":{"dominio":"filosofia","fonte":""},"origem":"wiki","palavras_chave":[],"sinonimos":[],"tipo":"conceito","titulo":"Epistemologia"}
\section{Epistemologia}
A teoria do conhecimento: o que podemos saber, como justificamos crenças, os limites do saber.
\prov{relações: parte\_de filosofia; relaciona metafisica; relaciona word; relaciona virtualizacao; relaciona psicologia; relaciona vida-morte; relaciona semiotica\_peirce\_triade; relaciona pow\_coordenadas\_espectrais\_broca; relaciona pow\_bit\_taxa}
\prov{proveniência: wiki \textperiodcentered{} fato \textperiodcentered{} confiança 1.0 \textperiodcentered{} domínio filosofia \textperiodcentered{} história 10 versões no jornal}

%CRISTAL {"arestas":[{"alvo":"teoria_dos_numeros","confianca":1.0,"epistemico":"fato","origem":"wiki","rel":"parte_de"},{"alvo":"pell","confianca":1.0,"epistemico":"fato","origem":"wiki","rel":"explica"},{"alvo":"corpos_e_aneis","confianca":1.0,"epistemico":"fato","origem":"wiki","rel":"usa"}],"confianca":1.0,"contraexemplos":[],"descricao":"x²−n·y²=1: a diofantina cujas soluções são as unidades do corpo ℚ√n — a mesma reta metálica x²−n·x−1 do goldchain.","epistemico":"fato","exemplos":[],"id":"equacao_de_pell","memoria":"semantica","meta":{"dominio":"ciencias_de_fora","fonte":""},"origem":"wiki","palavras_chave":[],"sinonimos":[],"tipo":"conceito","titulo":"Equação de Pell"}
\section{Equação de Pell}
x\ensuremath{^{2}}\ensuremath{-}n\textperiodcentered y\ensuremath{^{2}}=1: a diofantina cujas soluções são as unidades do corpo \ensuremath{\mathbb{Q}}\ensuremath{\sqrt{\,}}n — a mesma reta metálica x\ensuremath{^{2}}\ensuremath{-}n\textperiodcentered x\ensuremath{-}1 do goldchain.
\prov{relações: parte\_de teoria\_dos\_numeros; explica pell; usa corpos\_e\_aneis}
\prov{proveniência: wiki \textperiodcentered{} fato \textperiodcentered{} confiança 1.0 \textperiodcentered{} domínio ciencias\_de\_fora \textperiodcentered{} história 8 versões no jornal}

%CRISTAL {"arestas":[{"alvo":"inteligencia_consciencia_relacional","confianca":1.0,"epistemico":"fato","origem":"wiki","rel":"usa"},{"alvo":"deus_como_vazio_gentil","confianca":1.0,"epistemico":"fato","origem":"wiki","rel":"usa"},{"alvo":"deus_mentiroso","confianca":1.0,"epistemico":"fato","origem":"wiki","rel":"explica"},{"alvo":"sagrado","confianca":1.0,"epistemico":"fato","origem":"wiki","rel":"relaciona"},{"alvo":"tawhid","confianca":1.0,"epistemico":"fato","origem":"wiki","rel":"contradiz"}],"confianca":1.0,"contraexemplos":[],"descricao":"Lopes (O Deus Quântico): nenhum Deus existe 'em si' — todo Deus emerge da interação do Absoluto com uma mente humana ('o olho produz a luz, a mente produz Deus'). Casos: Absoluto+Maomé=Alá, Absoluto+Lao Tsé=Tao, Absoluto+Lopes=Deus(Vazio). O fanatismo anula a interface 'homem', gerando a falsa equação Absoluto=Deus.","epistemico":"inferencia","exemplos":[],"id":"equacao_divina_gentil","memoria":"semantica","meta":{"autor":"Gentil Lopes","dominio":"sagrado","fonte":""},"origem":"wiki","palavras_chave":[],"sinonimos":[],"tipo":"conceito","titulo":"A Equação Divina: Absoluto + Homem = Deus (Lopes)"}
\section{A Equação Divina: Absoluto + Homem = Deus (Lopes)}
Lopes (O Deus Quântico): nenhum Deus existe 'em si' — todo Deus emerge da interação do Absoluto com uma mente humana ('o olho produz a luz, a mente produz Deus'). Casos: Absoluto+Maomé=Alá, Absoluto+Lao Tsé=Tao, Absoluto+Lopes=Deus(Vazio). O fanatismo anula a interface 'homem', gerando a falsa equação Absoluto=Deus.
\prov{relações: usa inteligencia\_consciencia\_relacional; usa deus\_como\_vazio\_gentil; explica deus\_mentiroso; relaciona sagrado; contradiz tawhid}
\prov{proveniência: wiki \textperiodcentered{} inferencia \textperiodcentered{} confiança 1.0 \textperiodcentered{} domínio sagrado \textperiodcentered{} história 7 versões no jornal}

%CRISTAL {"arestas":[{"alvo":"traducao","confianca":1.0,"epistemico":"fato","origem":"wiki","rel":"usa"},{"alvo":"fidelidade_vs_fluencia","confianca":1.0,"epistemico":"fato","origem":"wiki","rel":"explica"}],"confianca":1.0,"contraexemplos":[],"descricao":"Nida: equivalência FORMAL (espelhar a estrutura da fonte) vs DINÂMICA (reproduzir o EFEITO no leitor do alvo). Traduzir o que a frase faz, não só o que diz.","epistemico":"inferencia","exemplos":[],"id":"equivalencia_traducao","memoria":"semantica","meta":{"autor":"Nida","dominio":"ciencias_de_fora","fonte":""},"origem":"wiki","palavras_chave":[],"sinonimos":[],"tipo":"conceito","titulo":"Equivalência (Nida)"}
\section{Equivalência (Nida)}
Nida: equivalência FORMAL (espelhar a estrutura da fonte) vs DINÂMICA (reproduzir o EFEITO no leitor do alvo). Traduzir o que a frase faz, não só o que diz.
\prov{relações: usa traducao; explica fidelidade\_vs\_fluencia}
\prov{proveniência: wiki \textperiodcentered{} inferencia \textperiodcentered{} confiança 1.0 \textperiodcentered{} domínio ciencias\_de\_fora \textperiodcentered{} história 3 versões no jornal}

%CRISTAL {"arestas":[{"alvo":"economia","confianca":1.0,"epistemico":"fato","origem":"wiki","rel":"parte_de"},{"alvo":"virtualizacao","confianca":0.5,"epistemico":"fato","origem":"wiki","rel":"relaciona"},{"alvo":"word","confianca":0.5,"epistemico":"fato","origem":"wiki","rel":"relaciona"},{"alvo":"proposition_gerador_de_escala_ipropderiv_com_tddt","confianca":0.5,"epistemico":"fato","origem":"wiki","rel":"relaciona"},{"alvo":"vida-morte","confianca":0.5,"epistemico":"fato","origem":"wiki","rel":"relaciona"},{"alvo":"proposition_convolucao_multiplicativa_produtopropconv","confianca":0.5,"epistemico":"fato","origem":"wiki","rel":"relaciona"},{"alvo":"gramatica","confianca":0.5,"epistemico":"fato","origem":"wiki","rel":"relaciona"},{"alvo":"semiotica_peirce_triade","confianca":0.5,"epistemico":"fato","origem":"wiki","rel":"relaciona"},{"alvo":"gato_operador_hub_contraexemplos","confianca":0.5,"epistemico":"fato","origem":"wiki","rel":"relaciona"},{"alvo":"filosofia_da_historia","confianca":0.5,"epistemico":"fato","origem":"wiki","rel":"relaciona"},{"alvo":"pow_coordenadas_espectrais_broca","confianca":0.5,"epistemico":"fato","origem":"wiki","rel":"relaciona"},{"alvo":"gentil_capitulo_6_progressao_geometrica_bidimensional","confianca":0.5,"epistemico":"fato","origem":"wiki","rel":"relaciona"},{"alvo":"misticismo_nao_dual","confianca":0.5,"epistemico":"fato","origem":"wiki","rel":"relaciona"},{"alvo":"minerar_e_resfriar","confianca":1.0,"epistemico":"fato","origem":"wiki","rel":"relaciona"}],"confianca":1.0,"contraexemplos":[],"descricao":"A limitação de recursos face a desejos ilimitados — a razão de a economia existir; sem escassez não há preço nem escolha.","epistemico":"fato","exemplos":[],"id":"escassez","memoria":"semantica","meta":{"autor":"—","dominio":"ciencias_de_fora","fonte":""},"origem":"wiki","palavras_chave":[],"sinonimos":[],"tipo":"conceito","titulo":"Escassez"}
\section{Escassez}
A limitação de recursos face a desejos ilimitados — a razão de a economia existir; sem escassez não há preço nem escolha.
\prov{relações: parte\_de economia; relaciona virtualizacao; relaciona word; relaciona proposition\_gerador\_de\_escala\_ipropderiv\_com\_tddt; relaciona vida-morte; relaciona proposition\_convolucao\_multiplicativa\_produtopropconv; relaciona gramatica; relaciona semiotica\_peirce\_triade; relaciona gato\_operador\_hub\_contraexemplos; relaciona filosofia\_da\_historia; relaciona pow\_coordenadas\_espectrais\_broca; relaciona gentil\_capitulo\_6\_progressao\_geometrica\_bidimensional; relaciona misticismo\_nao\_dual; relaciona minerar\_e\_resfriar}
\prov{proveniência: wiki \textperiodcentered{} fato \textperiodcentered{} confiança 1.0 \textperiodcentered{} domínio ciencias\_de\_fora \textperiodcentered{} história 18 versões no jornal}

%CRISTAL {"arestas":[{"alvo":"direito","confianca":1.0,"epistemico":"fato","origem":"wiki","rel":"parte_de"},{"alvo":"beccaria_reforma_penal","confianca":1.0,"epistemico":"fato","origem":"wiki","rel":"contradiz"},{"alvo":"justica_algoritmica","confianca":1.0,"epistemico":"fato","origem":"wiki","rel":"relaciona"}],"confianca":1.0,"contraexemplos":[],"descricao":"Lombroso: explicava o crime por fatores biológicos e deterministas ('criminoso nato'), deslocando o foco do livre-arbítrio — tese biológica hoje cientificamente refutada, mas cujo determinismo ressurge no viés preditivo algorítmico.","epistemico":"fato","exemplos":[],"id":"escola_positiva_criminologia","memoria":"semantica","meta":{"autor":"Lombroso","dominio":"ciencias_de_fora","fonte":""},"origem":"wiki","palavras_chave":[],"sinonimos":[],"tipo":"conceito","titulo":"Escola Positiva / Lombroso (refutado)"}
\section{Escola Positiva / Lombroso (refutado)}
Lombroso: explicava o crime por fatores biológicos e deterministas ('criminoso nato'), deslocando o foco do livre-arbítrio — tese biológica hoje cientificamente refutada, mas cujo determinismo ressurge no viés preditivo algorítmico.
\prov{relações: parte\_de direito; contradiz beccaria\_reforma\_penal; relaciona justica\_algoritmica}
\prov{proveniência: wiki \textperiodcentered{} fato \textperiodcentered{} confiança 1.0 \textperiodcentered{} domínio ciencias\_de\_fora \textperiodcentered{} história 4 versões no jornal}

%CRISTAL {"arestas":[{"alvo":"idioma","confianca":1.0,"epistemico":"fato","origem":"wiki","rel":"usa"},{"alvo":"codigo_semiotico","confianca":1.0,"epistemico":"fato","origem":"wiki","rel":"usa"}],"confianca":1.0,"contraexemplos":[],"descricao":"O registro visual da língua — a tecnologia que a libertou do aqui-agora e fundou a história, a lei e a ciência. Nem toda língua a tem.","epistemico":"fato","exemplos":[],"id":"escrita","memoria":"semantica","meta":{"dominio":"ciencias_de_fora","fonte":""},"origem":"wiki","palavras_chave":[],"sinonimos":[],"tipo":"conceito","titulo":"Escrita"}
\section{Escrita}
O registro visual da língua — a tecnologia que a libertou do aqui-agora e fundou a história, a lei e a ciência. Nem toda língua a tem.
\prov{relações: usa idioma; usa codigo\_semiotico}
\prov{proveniência: wiki \textperiodcentered{} fato \textperiodcentered{} confiança 1.0 \textperiodcentered{} domínio ciencias\_de\_fora \textperiodcentered{} história 3 versões no jornal}

%CRISTAL {"arestas":[{"alvo":"geografia","confianca":1.0,"epistemico":"fato","origem":"wiki","rel":"relaciona"},{"alvo":"regua_de_gentil","confianca":1.0,"epistemico":"fato","origem":"wiki","rel":"usa"},{"alvo":"borda_critica_vazio","confianca":1.0,"epistemico":"fato","origem":"wiki","rel":"relaciona"}],"confianca":1.0,"contraexemplos":[],"descricao":"Lopes construiu um espaço métrico (M,d) sobre o intervalo [0,1[ cuja régua difere da euclidiana — nele um mesmo ponto pode distar zero de quatro regiões ao mesmo tempo (0,999…=0). Medir distância no espaço depende da régua adotada: o espaço não é dado, é construído.","epistemico":"inferencia","exemplos":[],"id":"espaco_metrico_quantico_gentil","memoria":"semantica","meta":{"autor":"Gentil Lopes","dominio":"ciencias_de_fora","fonte":""},"origem":"wiki","palavras_chave":[],"sinonimos":[],"tipo":"conceito","titulo":"Espaço Métrico Quântico (Lopes)"}
\section{Espaço Métrico Quântico (Lopes)}
Lopes construiu um espaço métrico (M,d) sobre o intervalo [0,1[ cuja régua difere da euclidiana — nele um mesmo ponto pode distar zero de quatro regiões ao mesmo tempo (0,999…=0). Medir distância no espaço depende da régua adotada: o espaço não é dado, é construído.
\prov{relações: relaciona geografia; usa regua\_de\_gentil; relaciona borda\_critica\_vazio}
\prov{proveniência: wiki \textperiodcentered{} inferencia \textperiodcentered{} confiança 1.0 \textperiodcentered{} domínio ciencias\_de\_fora \textperiodcentered{} história 5 versões no jornal}

%CRISTAL {"arestas":[{"alvo":"geografia_humana","confianca":1.0,"epistemico":"fato","origem":"wiki","rel":"parte_de"},{"alvo":"topofilia_lugar_tuan","confianca":1.0,"epistemico":"fato","origem":"wiki","rel":"contradiz"},{"alvo":"sociedade_em_rede","confianca":1.0,"epistemico":"fato","origem":"wiki","rel":"relaciona"}],"confianca":1.0,"contraexemplos":[],"descricao":"Doreen Massey: o espaço é produto de inter-relações; é a esfera da multiplicidade, onde coexistem trajetórias distintas; e está sempre em construção, nunca fechado. Recusa a visão do espaço como superfície fixa ou 'tempo tornado estático'.","epistemico":"inferencia","exemplos":[],"id":"espaco_relacional_massey","memoria":"semantica","meta":{"autor":"Massey","dominio":"ciencias_de_fora","fonte":""},"origem":"wiki","palavras_chave":[],"sinonimos":[],"tipo":"conceito","titulo":"Espaço como Produto de Inter-relações (Massey)"}
\section{Espaço como Produto de Inter-relações (Massey)}
Doreen Massey: o espaço é produto de inter-relações; é a esfera da multiplicidade, onde coexistem trajetórias distintas; e está sempre em construção, nunca fechado. Recusa a visão do espaço como superfície fixa ou 'tempo tornado estático'.
\prov{relações: parte\_de geografia\_humana; contradiz topofilia\_lugar\_tuan; relaciona sociedade\_em\_rede}
\prov{proveniência: wiki \textperiodcentered{} inferencia \textperiodcentered{} confiança 1.0 \textperiodcentered{} domínio ciencias\_de\_fora \textperiodcentered{} história 4 versões no jornal}

%CRISTAL {"arestas":[{"alvo":"geografia_humana","confianca":1.0,"epistemico":"fato","origem":"wiki","rel":"parte_de"},{"alvo":"materialismo_historico","confianca":1.0,"epistemico":"fato","origem":"wiki","rel":"usa"},{"alvo":"sociedade_em_rede","confianca":1.0,"epistemico":"fato","origem":"wiki","rel":"relaciona"}],"confianca":1.0,"contraexemplos":[],"descricao":"Milton Santos: o espaço geográfico é um 'conjunto indissociável de sistemas de objetos e sistemas de ações' — os objetos (fixos, técnicas materializadas) condicionam as ações, e as ações recriam os objetos, num par inseparável em constante interação.","epistemico":"inferencia","exemplos":[],"id":"espaco_santos","memoria":"semantica","meta":{"autor":"Milton Santos","dominio":"ciencias_de_fora","fonte":""},"origem":"wiki","palavras_chave":[],"sinonimos":[],"tipo":"conceito","titulo":"Espaço como Sistemas de Objetos e de Ações (Milton Santos)"}
\section{Espaço como Sistemas de Objetos e de Ações (Milton Santos)}
Milton Santos: o espaço geográfico é um 'conjunto indissociável de sistemas de objetos e sistemas de ações' — os objetos (fixos, técnicas materializadas) condicionam as ações, e as ações recriam os objetos, num par inseparável em constante interação.
\prov{relações: parte\_de geografia\_humana; usa materialismo\_historico; relaciona sociedade\_em\_rede}
\prov{proveniência: wiki \textperiodcentered{} inferencia \textperiodcentered{} confiança 1.0 \textperiodcentered{} domínio ciencias\_de\_fora \textperiodcentered{} história 4 versões no jornal}

%CRISTAL {"arestas":[{"alvo":"linguas_romanicas","confianca":1.0,"epistemico":"fato","origem":"wiki","rel":"especializa"}],"confianca":1.0,"contraexemplos":[],"descricao":"A língua românica com mais falantes nativos (~490 milhões) — da Espanha a quase toda a América Latina.","epistemico":"fato","exemplos":[],"id":"espanhol","memoria":"semantica","meta":{"dominio":"ciencias_de_fora","fonte":""},"origem":"wiki","palavras_chave":[],"sinonimos":[],"tipo":"conceito","titulo":"Espanhol"}
\section{Espanhol}
A língua românica com mais falantes nativos (\~{}490 milhões) — da Espanha a quase toda a América Latina.
\prov{relações: especializa linguas\_romanicas}
\prov{proveniência: wiki \textperiodcentered{} fato \textperiodcentered{} confiança 1.0 \textperiodcentered{} domínio ciencias\_de\_fora \textperiodcentered{} história 2 versões no jornal}

%CRISTAL {"arestas":[{"alvo":"biologia_evolutiva","confianca":1.0,"epistemico":"fato","origem":"wiki","rel":"parte_de"}],"confianca":1.0,"contraexemplos":[],"descricao":"Como uma linhagem se divide em espécies — em geral por isolamento reprodutivo que acumula diferenças.","epistemico":"fato","exemplos":[],"id":"especiacao","memoria":"semantica","meta":{"dominio":"ciencias_de_fora","fonte":""},"origem":"wiki","palavras_chave":[],"sinonimos":[],"tipo":"conceito","titulo":"Especiação"}
\section{Especiação}
Como uma linhagem se divide em espécies — em geral por isolamento reprodutivo que acumula diferenças.
\prov{relações: parte\_de biologia\_evolutiva}
\prov{proveniência: wiki \textperiodcentered{} fato \textperiodcentered{} confiança 1.0 \textperiodcentered{} domínio ciencias\_de\_fora \textperiodcentered{} história 2 versões no jornal}

%CRISTAL {"arestas":[{"alvo":"teoria_da_comunicacao","confianca":1.0,"epistemico":"fato","origem":"wiki","rel":"parte_de"},{"alvo":"dois_sistemas_kahneman","confianca":1.0,"epistemico":"fato","origem":"wiki","rel":"relaciona"},{"alvo":"biopoder","confianca":1.0,"epistemico":"fato","origem":"wiki","rel":"relaciona"}],"confianca":1.0,"contraexemplos":[],"descricao":"Noelle-Neumann: por medo do isolamento, quem percebe sua opinião como minoritária tende a silenciá-la e quem se percebe maioria se expressa — um espiral que amplifica a opinião dominante. A opinião pública como 'pele social' que pressiona à conformidade.","epistemico":"fato","exemplos":[],"id":"espiral_do_silencio","memoria":"semantica","meta":{"autor":"Noelle-Neumann","dominio":"ciencias_de_fora","fonte":""},"origem":"wiki","palavras_chave":[],"sinonimos":[],"tipo":"conceito","titulo":"Espiral do Silêncio"}
\section{Espiral do Silêncio}
Noelle-Neumann: por medo do isolamento, quem percebe sua opinião como minoritária tende a silenciá-la e quem se percebe maioria se expressa — um espiral que amplifica a opinião dominante. A opinião pública como 'pele social' que pressiona à conformidade.
\prov{relações: parte\_de teoria\_da\_comunicacao; relaciona dois\_sistemas\_kahneman; relaciona biopoder}
\prov{proveniência: wiki \textperiodcentered{} fato \textperiodcentered{} confiança 1.0 \textperiodcentered{} domínio ciencias\_de\_fora \textperiodcentered{} história 4 versões no jornal}

%CRISTAL {"arestas":[{"alvo":"direito","confianca":1.0,"epistemico":"fato","origem":"wiki","rel":"parte_de"},{"alvo":"contrato_social","confianca":1.0,"epistemico":"fato","origem":"wiki","rel":"relaciona"},{"alvo":"biopoder","confianca":1.0,"epistemico":"fato","origem":"wiki","rel":"contradiz"}],"confianca":1.0,"contraexemplos":[],"descricao":"Toda autoridade, inclusive a governamental, está submetida a leis públicas, gerais, claras e estáveis — o 'governo das leis, não dos homens', em que ninguém está acima da lei.","epistemico":"inferencia","exemplos":[],"id":"estado_de_direito","memoria":"semantica","meta":{"autor":"Dicey/Waldron","dominio":"ciencias_de_fora","fonte":""},"origem":"wiki","palavras_chave":[],"sinonimos":[],"tipo":"conceito","titulo":"Estado de Direito (Rule of Law)"}
\section{Estado de Direito (Rule of Law)}
Toda autoridade, inclusive a governamental, está submetida a leis públicas, gerais, claras e estáveis — o 'governo das leis, não dos homens', em que ninguém está acima da lei.
\prov{relações: parte\_de direito; relaciona contrato\_social; contradiz biopoder}
\prov{proveniência: wiki \textperiodcentered{} inferencia \textperiodcentered{} confiança 1.0 \textperiodcentered{} domínio ciencias\_de\_fora \textperiodcentered{} história 4 versões no jornal}

%CRISTAL {"arestas":[{"alvo":"teorias_aprendizagem","confianca":1.0,"epistemico":"fato","origem":"wiki","rel":"parte_de"},{"alvo":"construtivismo_piaget","confianca":1.0,"epistemico":"fato","origem":"wiki","rel":"usa"}],"confianca":1.0,"contraexemplos":[],"descricao":"Piaget: o desenvolvimento progride por quatro estágios sequenciais — sensório-motor, pré-operatório, operatório concreto e formal. A ordem é robusta; as idades e a rigidez são contestadas.","epistemico":"fato","exemplos":[],"id":"estagios_piaget","memoria":"semantica","meta":{"autor":"Piaget","dominio":"ciencias_de_fora","fonte":""},"origem":"wiki","palavras_chave":[],"sinonimos":[],"tipo":"conceito","titulo":"Estágios do Desenvolvimento Cognitivo"}
\section{Estágios do Desenvolvimento Cognitivo}
Piaget: o desenvolvimento progride por quatro estágios sequenciais — sensório-motor, pré-operatório, operatório concreto e formal. A ordem é robusta; as idades e a rigidez são contestadas.
\prov{relações: parte\_de teorias\_aprendizagem; usa construtivismo\_piaget}
\prov{proveniência: wiki \textperiodcentered{} fato \textperiodcentered{} confiança 1.0 \textperiodcentered{} domínio ciencias\_de\_fora \textperiodcentered{} história 3 versões no jornal}

%CRISTAL {"arestas":[{"alvo":"bioetica","confianca":1.0,"epistemico":"fato","origem":"wiki","rel":"parte_de"},{"alvo":"autonomia_do_paciente","confianca":1.0,"epistemico":"fato","origem":"wiki","rel":"relaciona"}],"confianca":1.0,"contraexemplos":[],"descricao":"A questão de se (e quando) o embrião/feto adquire status de pessoa com direito à vida — da concepção ao gradualismo (status crescente) — núcleo do debate bioético do aborto, que envolve ainda o direito da gestante sobre o próprio corpo.","epistemico":"inferencia","exemplos":[],"id":"estatuto_moral_embriao","memoria":"semantica","meta":{"autor":"Warren/Thomson/Marquis","dominio":"ciencias_de_fora","fonte":""},"origem":"wiki","palavras_chave":[],"sinonimos":[],"tipo":"conceito","titulo":"Estatuto Moral do Embrião"}
\section{Estatuto Moral do Embrião}
A questão de se (e quando) o embrião/feto adquire status de pessoa com direito à vida — da concepção ao gradualismo (status crescente) — núcleo do debate bioético do aborto, que envolve ainda o direito da gestante sobre o próprio corpo.
\prov{relações: parte\_de bioetica; relaciona autonomia\_do\_paciente}
\prov{proveniência: wiki \textperiodcentered{} inferencia \textperiodcentered{} confiança 1.0 \textperiodcentered{} domínio ciencias\_de\_fora \textperiodcentered{} história 3 versões no jornal}

%CRISTAL {"arestas":[{"alvo":"filosofia","confianca":1.0,"epistemico":"fato","origem":"wiki","rel":"parte_de"},{"alvo":"utilitarismo","confianca":0.5,"epistemico":"fato","origem":"wiki","rel":"relaciona"},{"alvo":"filesystem","confianca":0.5,"epistemico":"fato","origem":"wiki","rel":"relaciona"},{"alvo":"word","confianca":0.5,"epistemico":"fato","origem":"wiki","rel":"relaciona"},{"alvo":"virtualizacao","confianca":0.5,"epistemico":"fato","origem":"wiki","rel":"relaciona"},{"alvo":"misticismo_nao_dual","confianca":0.5,"epistemico":"fato","origem":"wiki","rel":"relaciona"},{"alvo":"vida-morte","confianca":0.5,"epistemico":"fato","origem":"wiki","rel":"relaciona"},{"alvo":"honestidade_intelectual","confianca":0.5,"epistemico":"fato","origem":"wiki","rel":"relaciona"},{"alvo":"vontade_de_poder","confianca":0.5,"epistemico":"fato","origem":"wiki","rel":"relaciona"},{"alvo":"gentil_capitulo_3_sequencias_geometricas_de_ordem_m","confianca":0.5,"epistemico":"fato","origem":"wiki","rel":"relaciona"},{"alvo":"musica","confianca":0.5,"epistemico":"fato","origem":"wiki","rel":"relaciona"},{"alvo":"ontologia","confianca":0.5,"epistemico":"fato","origem":"wiki","rel":"relaciona"},{"alvo":"inspecao_local_32","confianca":0.5,"epistemico":"fato","origem":"wiki","rel":"relaciona"}],"confianca":1.0,"contraexemplos":[],"descricao":"O estudo do belo, do sublime e da arte — o juízo e a experiência do sensível.","epistemico":"fato","exemplos":[],"id":"estetica","memoria":"semantica","meta":{"dominio":"ciencias_de_fora","fonte":""},"origem":"wiki","palavras_chave":[],"sinonimos":[],"tipo":"conceito","titulo":"Estética"}
\section{Estética}
O estudo do belo, do sublime e da arte — o juízo e a experiência do sensível.
\prov{relações: parte\_de filosofia; relaciona utilitarismo; relaciona filesystem; relaciona word; relaciona virtualizacao; relaciona misticismo\_nao\_dual; relaciona vida-morte; relaciona honestidade\_intelectual; relaciona vontade\_de\_poder; relaciona gentil\_capitulo\_3\_sequencias\_geometricas\_de\_ordem\_m; relaciona musica; relaciona ontologia; relaciona inspecao\_local\_32}
\prov{proveniência: wiki \textperiodcentered{} fato \textperiodcentered{} confiança 1.0 \textperiodcentered{} domínio ciencias\_de\_fora \textperiodcentered{} história 14 versões no jornal}

%CRISTAL {"arestas":[{"alvo":"literatura","confianca":1.0,"epistemico":"fato","origem":"wiki","rel":"parte_de"},{"alvo":"obra_aberta","confianca":1.0,"epistemico":"fato","origem":"wiki","rel":"usa"},{"alvo":"morte_do_autor","confianca":1.0,"epistemico":"fato","origem":"wiki","rel":"usa"},{"alvo":"catarse","confianca":1.0,"epistemico":"fato","origem":"wiki","rel":"relaciona"}],"confianca":1.0,"contraexemplos":[],"descricao":"Jauss/Iser/Gadamer (Escola de Constança): o sentido não está pronto no texto nem no autor, mas se realiza no encontro texto-leitor — contra um 'horizonte de expectativas', preenchendo lacunas, numa 'fusão de horizontes'.","epistemico":"inferencia","exemplos":[],"id":"estetica_da_recepcao","memoria":"semantica","meta":{"autor":"Jauss/Iser","dominio":"ciencias_de_fora","fonte":""},"origem":"wiki","palavras_chave":[],"sinonimos":[],"tipo":"conceito","titulo":"Estética da Recepção"}
\section{Estética da Recepção}
Jauss/Iser/Gadamer (Escola de Constança): o sentido não está pronto no texto nem no autor, mas se realiza no encontro texto-leitor — contra um 'horizonte de expectativas', preenchendo lacunas, numa 'fusão de horizontes'.
\prov{relações: parte\_de literatura; usa obra\_aberta; usa morte\_do\_autor; relaciona catarse}
\prov{proveniência: wiki \textperiodcentered{} inferencia \textperiodcentered{} confiança 1.0 \textperiodcentered{} domínio ciencias\_de\_fora \textperiodcentered{} história 5 versões no jornal}

%CRISTAL {"arestas":[{"alvo":"criacao_livre_fertilidade","confianca":1.0,"epistemico":"fato","origem":"wiki","rel":"usa"},{"alvo":"meditacao_cura_da_mente_gentil","confianca":1.0,"epistemico":"fato","origem":"wiki","rel":"relaciona"},{"alvo":"artes","confianca":1.0,"epistemico":"fato","origem":"wiki","rel":"relaciona"},{"alvo":"matematica_arte_engenharia","confianca":1.0,"epistemico":"fato","origem":"wiki","rel":"relaciona"}],"confianca":1.0,"contraexemplos":[],"descricao":"Lopes endossa Osho: 'um poeta autêntico não está presente, ele abre caminho para a poesia fluir dele — o mesmo com a música, a dança, a escultura... nada disso é produto do ego'. A criação (artística e matemática) flui de um estado não-egóico; ecoa a criação livre por fertilidade e o deletar-o-ego da meditação.","epistemico":"inferencia","exemplos":[],"id":"estetica_do_fluir_gentil","memoria":"semantica","meta":{"autor":"Gentil Lopes","dominio":"ciencias_de_fora","fonte":""},"origem":"wiki","palavras_chave":[],"sinonimos":[],"tipo":"conceito","titulo":"A Estética do Fluir Não-Egóico (Lopes)"}
\section{A Estética do Fluir Não-Egóico (Lopes)}
Lopes endossa Osho: 'um poeta autêntico não está presente, ele abre caminho para a poesia fluir dele — o mesmo com a música, a dança, a escultura... nada disso é produto do ego'. A criação (artística e matemática) flui de um estado não-egóico; ecoa a criação livre por fertilidade e o deletar-o-ego da meditação.
\prov{relações: usa criacao\_livre\_fertilidade; relaciona meditacao\_cura\_da\_mente\_gentil; relaciona artes; relaciona matematica\_arte\_engenharia}
\prov{proveniência: wiki \textperiodcentered{} inferencia \textperiodcentered{} confiança 1.0 \textperiodcentered{} domínio ciencias\_de\_fora \textperiodcentered{} história 6 versões no jornal}

%CRISTAL {"arestas":[{"alvo":"literariedade","confianca":1.0,"epistemico":"fato","origem":"wiki","rel":"usa"},{"alvo":"obra_aberta","confianca":1.0,"epistemico":"fato","origem":"wiki","rel":"relaciona"}],"confianca":1.0,"contraexemplos":[],"descricao":"Chklóvski (Formalismo Russo): o procedimento artístico que torna estranho o familiar, dificultando e prolongando a percepção, para que o objeto seja sentido como visão e não como reconhecimento automático.","epistemico":"inferencia","exemplos":[],"id":"estranhamento","memoria":"semantica","meta":{"autor":"Chklóvski","dominio":"ciencias_de_fora","fonte":""},"origem":"wiki","palavras_chave":[],"sinonimos":[],"tipo":"conceito","titulo":"Estranhamento / Desfamiliarização (Chklóvski)"}
\section{Estranhamento / Desfamiliarização (Chklóvski)}
Chklóvski (Formalismo Russo): o procedimento artístico que torna estranho o familiar, dificultando e prolongando a percepção, para que o objeto seja sentido como visão e não como reconhecimento automático.
\prov{relações: usa literariedade; relaciona obra\_aberta}
\prov{proveniência: wiki \textperiodcentered{} inferencia \textperiodcentered{} confiança 1.0 \textperiodcentered{} domínio ciencias\_de\_fora \textperiodcentered{} história 3 versões no jornal}

%CRISTAL {"arestas":[{"alvo":"biologia_evolutiva","confianca":1.0,"epistemico":"fato","origem":"wiki","rel":"parte_de"},{"alvo":"teoria_dos_jogos","confianca":1.0,"epistemico":"fato","origem":"wiki","rel":"usa"}],"confianca":1.0,"contraexemplos":[],"descricao":"Maynard Smith: teoria dos jogos na evolução — uma estratégia que, comum, não pode ser invadida por outra.","epistemico":"fato","exemplos":[],"id":"estrategia_evolutiva_estavel","memoria":"semantica","meta":{"autor":"Maynard Smith","dominio":"ciencias_de_fora","fonte":""},"origem":"wiki","palavras_chave":[],"sinonimos":[],"tipo":"conceito","titulo":"Estratégia Evolutivamente Estável"}
\section{Estratégia Evolutivamente Estável}
Maynard Smith: teoria dos jogos na evolução — uma estratégia que, comum, não pode ser invadida por outra.
\prov{relações: parte\_de biologia\_evolutiva; usa teoria\_dos\_jogos}
\prov{proveniência: wiki \textperiodcentered{} fato \textperiodcentered{} confiança 1.0 \textperiodcentered{} domínio ciencias\_de\_fora \textperiodcentered{} história 3 versões no jornal}

%CRISTAL {"arestas":[{"alvo":"gramatica_gerativa","confianca":1.0,"epistemico":"fato","origem":"wiki","rel":"parte_de"}],"confianca":1.0,"contraexemplos":[],"descricao":"Chomsky: uma frase tem uma forma abstrata subjacente (profunda) transformada na forma pronunciada (superfície).","epistemico":"inferencia","exemplos":[],"id":"estrutura_profunda_superficie","memoria":"semantica","meta":{"autor":"Chomsky","dominio":"ciencias_de_fora","fonte":""},"origem":"wiki","palavras_chave":[],"sinonimos":[],"tipo":"conceito","titulo":"Estrutura Profunda e de Superfície"}
\section{Estrutura Profunda e de Superfície}
Chomsky: uma frase tem uma forma abstrata subjacente (profunda) transformada na forma pronunciada (superfície).
\prov{relações: parte\_de gramatica\_gerativa}
\prov{proveniência: wiki \textperiodcentered{} inferencia \textperiodcentered{} confiança 1.0 \textperiodcentered{} domínio ciencias\_de\_fora \textperiodcentered{} história 2 versões no jornal}

%CRISTAL {"arestas":[{"alvo":"filosofia_da_matematica","confianca":1.0,"epistemico":"fato","origem":"wiki","rel":"parte_de"},{"alvo":"reticulados","confianca":1.0,"epistemico":"fato","origem":"wiki","rel":"referencia"}],"confianca":1.0,"contraexemplos":[],"descricao":"A matemática estuda estruturas e relações, não objetos-substância (Bourbaki) — literalmente a forma de um grafo de conceitos.","epistemico":"inferencia","exemplos":[],"id":"estruturalismo_bourbaki","memoria":"semantica","meta":{"autor":"Bourbaki","dominio":"filosofia","fonte":""},"origem":"wiki","palavras_chave":[],"sinonimos":[],"tipo":"conceito","titulo":"Estruturalismo Matemático"}
\section{Estruturalismo Matemático}
A matemática estuda estruturas e relações, não objetos-substância (Bourbaki) — literalmente a forma de um grafo de conceitos.
\prov{relações: parte\_de filosofia\_da\_matematica; referencia reticulados}
\prov{proveniência: wiki \textperiodcentered{} inferencia \textperiodcentered{} confiança 1.0 \textperiodcentered{} domínio filosofia \textperiodcentered{} história 3 versões no jornal}

%CRISTAL {"arestas":[{"alvo":"signo_saussure","confianca":1.0,"epistemico":"fato","origem":"wiki","rel":"usa"},{"alvo":"estruturalismo_bourbaki","confianca":1.0,"epistemico":"fato","origem":"wiki","rel":"relaciona"},{"alvo":"signo_linguistico","confianca":1.0,"epistemico":"fato","origem":"wiki","rel":"usa"}],"confianca":1.0,"contraexemplos":[],"descricao":"Saussure/Jakobson: a língua é sistema — o valor de um signo vem das diferenças com os outros, nos eixos sintagmático e paradigmático.","epistemico":"inferencia","exemplos":[],"id":"estruturalismo_linguistico","memoria":"semantica","meta":{"autor":"Jakobson","dominio":"ciencias_de_fora","fonte":""},"origem":"wiki","palavras_chave":[],"sinonimos":[],"tipo":"conceito","titulo":"Estruturalismo Linguístico"}
\section{Estruturalismo Linguístico}
Saussure/Jakobson: a língua é sistema — o valor de um signo vem das diferenças com os outros, nos eixos sintagmático e paradigmático.
\prov{relações: usa signo\_saussure; relaciona estruturalismo\_bourbaki; usa signo\_linguistico}
\prov{proveniência: wiki \textperiodcentered{} inferencia \textperiodcentered{} confiança 1.0 \textperiodcentered{} domínio ciencias\_de\_fora \textperiodcentered{} história 4 versões no jornal}

%CRISTAL {"arestas":[{"alvo":"antropologia","confianca":1.0,"epistemico":"fato","origem":"wiki","rel":"parte_de"},{"alvo":"estruturalismo_linguistico","confianca":1.0,"epistemico":"fato","origem":"wiki","rel":"relaciona"}],"confianca":1.0,"contraexemplos":[],"descricao":"Lévi-Strauss: os mitos são sistemas de transformações que reconciliam contradições — a mesma lógica opera em todas as culturas.","epistemico":"inferencia","exemplos":[],"id":"estruturalismo_mito","memoria":"semantica","meta":{"autor":"Lévi-Strauss","dominio":"ciencias_de_fora","fonte":""},"origem":"wiki","palavras_chave":[],"sinonimos":[],"tipo":"conceito","titulo":"A Lógica do Mito"}
\section{A Lógica do Mito}
Lévi-Strauss: os mitos são sistemas de transformações que reconciliam contradições — a mesma lógica opera em todas as culturas.
\prov{relações: parte\_de antropologia; relaciona estruturalismo\_linguistico}
\prov{proveniência: wiki \textperiodcentered{} inferencia \textperiodcentered{} confiança 1.0 \textperiodcentered{} domínio ciencias\_de\_fora \textperiodcentered{} história 3 versões no jornal}

%CRISTAL {"arestas":[{"alvo":"estudos_de_midia","confianca":1.0,"epistemico":"fato","origem":"wiki","rel":"parte_de"},{"alvo":"cultura_modo_de_vida","confianca":1.0,"epistemico":"fato","origem":"wiki","rel":"relaciona"}],"confianca":1.0,"contraexemplos":[],"descricao":"Raymond Williams: um modo de pensar-e-sentir social ainda emergente — percebido antes de poder ser plenamente articulado — que capta a experiência viva de uma época e se expressa, ainda que implicitamente, em suas obras.","epistemico":"inferencia","exemplos":[],"id":"estruturas_de_sentimento","memoria":"semantica","meta":{"autor":"Williams","dominio":"ciencias_de_fora","fonte":""},"origem":"wiki","palavras_chave":[],"sinonimos":[],"tipo":"conceito","titulo":"Estruturas de Sentimento (Williams)"}
\section{Estruturas de Sentimento (Williams)}
Raymond Williams: um modo de pensar-e-sentir social ainda emergente — percebido antes de poder ser plenamente articulado — que capta a experiência viva de uma época e se expressa, ainda que implicitamente, em suas obras.
\prov{relações: parte\_de estudos\_de\_midia; relaciona cultura\_modo\_de\_vida}
\prov{proveniência: wiki \textperiodcentered{} inferencia \textperiodcentered{} confiança 1.0 \textperiodcentered{} domínio ciencias\_de\_fora \textperiodcentered{} história 3 versões no jornal}

%CRISTAL {"arestas":[{"alvo":"comunicacao","confianca":1.0,"epistemico":"fato","origem":"wiki","rel":"parte_de"},{"alvo":"estetica","confianca":1.0,"epistemico":"fato","origem":"wiki","rel":"relaciona"},{"alvo":"word","confianca":0.5,"epistemico":"fato","origem":"wiki","rel":"relaciona"}],"confianca":1.0,"contraexemplos":[],"descricao":"A análise crítica da mídia como indústria, poder e cultura — como as mensagens produzem sentido, ideologia e consenso.","epistemico":"fato","exemplos":[],"id":"estudos_de_midia","memoria":"semantica","meta":{"dominio":"ciencias_de_fora","fonte":""},"origem":"wiki","palavras_chave":[],"sinonimos":[],"tipo":"conceito","titulo":"Estudos de Mídia e Cultura"}
\section{Estudos de Mídia e Cultura}
A análise crítica da mídia como indústria, poder e cultura — como as mensagens produzem sentido, ideologia e consenso.
\prov{relações: parte\_de comunicacao; relaciona estetica; relaciona word}
\prov{proveniência: wiki \textperiodcentered{} fato \textperiodcentered{} confiança 1.0 \textperiodcentered{} domínio ciencias\_de\_fora \textperiodcentered{} história 4 versões no jornal}

%CRISTAL {"arestas":[{"alvo":"ciclos_vico","confianca":1.0,"epistemico":"fato","origem":"wiki","rel":"relaciona"},{"alvo":"word","confianca":0.5,"epistemico":"fato","origem":"wiki","rel":"relaciona"},{"alvo":"virtualizacao","confianca":0.5,"epistemico":"fato","origem":"wiki","rel":"relaciona"},{"alvo":"vita_contemplativa","confianca":0.5,"epistemico":"fato","origem":"wiki","rel":"relaciona"}],"confianca":1.0,"contraexemplos":[],"descricao":"Nietzsche: se o tempo é infinito e a matéria finita, tudo se repete idêntico para sempre — experimento ético (amor fati), não lei científica.","epistemico":"hipotese","exemplos":[],"id":"eterno_retorno","memoria":"semantica","meta":{"autor":"Nietzsche","dominio":"ciencias_de_fora","fonte":""},"origem":"wiki","palavras_chave":[],"sinonimos":[],"tipo":"conceito","titulo":"Eterno Retorno"}
\section{Eterno Retorno}
Nietzsche: se o tempo é infinito e a matéria finita, tudo se repete idêntico para sempre — experimento ético (amor fati), não lei científica.
\prov{relações: relaciona ciclos\_vico; relaciona word; relaciona virtualizacao; relaciona vita\_contemplativa}
\prov{proveniência: wiki \textperiodcentered{} hipotese \textperiodcentered{} confiança 1.0 \textperiodcentered{} domínio ciencias\_de\_fora \textperiodcentered{} história 5 versões no jornal}

%CRISTAL {"arestas":[{"alvo":"etica","confianca":1.0,"epistemico":"fato","origem":"wiki","rel":"parte_de"},{"alvo":"virtualizacao","confianca":0.5,"epistemico":"fato","origem":"wiki","rel":"relaciona"},{"alvo":"vida-morte","confianca":0.5,"epistemico":"fato","origem":"wiki","rel":"relaciona"},{"alvo":"word","confianca":0.5,"epistemico":"fato","origem":"wiki","rel":"relaciona"},{"alvo":"telemetria_shadow_producao","confianca":0.5,"epistemico":"fato","origem":"wiki","rel":"relaciona"},{"alvo":"fibonacci_multidimensional","confianca":0.5,"epistemico":"fato","origem":"wiki","rel":"relaciona"},{"alvo":"heap","confianca":0.5,"epistemico":"fato","origem":"wiki","rel":"relaciona"},{"alvo":"readme","confianca":0.5,"epistemico":"fato","origem":"wiki","rel":"relaciona"},{"alvo":"incautos_enganados","confianca":0.5,"epistemico":"fato","origem":"wiki","rel":"relaciona"},{"alvo":"gentil_12_definicao","confianca":0.5,"epistemico":"fato","origem":"wiki","rel":"relaciona"},{"alvo":"lemma_lemunit_e_unitaria_0xt2dtt_-","confianca":0.5,"epistemico":"fato","origem":"wiki","rel":"relaciona"},{"alvo":"politica","confianca":0.5,"epistemico":"fato","origem":"wiki","rel":"relaciona"}],"confianca":1.0,"contraexemplos":[],"descricao":"Gilligan/Noddings: a moral funda-se em relações, atenção e interdependência, não em princípios abstratos.","epistemico":"inferencia","exemplos":[],"id":"ethics_of_care","memoria":"semantica","meta":{"autor":"Gilligan","dominio":"ciencias_de_fora","fonte":""},"origem":"wiki","palavras_chave":[],"sinonimos":[],"tipo":"conceito","titulo":"Ética do Cuidado"}
\section{Ética do Cuidado}
Gilligan/Noddings: a moral funda-se em relações, atenção e interdependência, não em princípios abstratos.
\prov{relações: parte\_de etica; relaciona virtualizacao; relaciona vida-morte; relaciona word; relaciona telemetria\_shadow\_producao; relaciona fibonacci\_multidimensional; relaciona heap; relaciona readme; relaciona incautos\_enganados; relaciona gentil\_12\_definicao; relaciona lemma\_lemunit\_e\_unitaria\_0xt2dtt\_-; relaciona politica}
\prov{proveniência: wiki \textperiodcentered{} inferencia \textperiodcentered{} confiança 1.0 \textperiodcentered{} domínio ciencias\_de\_fora \textperiodcentered{} história 13 versões no jornal}

%CRISTAL {"arestas":[{"alvo":"filosofia","confianca":1.0,"epistemico":"fato","origem":"wiki","rel":"parte_de"},{"alvo":"historia","confianca":0.5,"epistemico":"fato","origem":"wiki","rel":"relaciona"},{"alvo":"testes","confianca":0.5,"epistemico":"fato","origem":"wiki","rel":"relaciona"},{"alvo":"goldchain_mercado_hub","confianca":0.5,"epistemico":"fato","origem":"wiki","rel":"relaciona"},{"alvo":"expressividade","confianca":0.5,"epistemico":"fato","origem":"wiki","rel":"relaciona"},{"alvo":"word","confianca":0.5,"epistemico":"fato","origem":"wiki","rel":"relaciona"},{"alvo":"vontade_de_poder","confianca":0.5,"epistemico":"fato","origem":"wiki","rel":"relaciona"}],"confianca":1.0,"contraexemplos":[],"descricao":"O estudo do agir bom e justo — virtude, dever, consequência, cuidado.","epistemico":"fato","exemplos":[],"id":"etica","memoria":"semantica","meta":{"dominio":"ciencias_de_fora","fonte":""},"origem":"wiki","palavras_chave":[],"sinonimos":[],"tipo":"conceito","titulo":"Ética"}
\section{Ética}
O estudo do agir bom e justo — virtude, dever, consequência, cuidado.
\prov{relações: parte\_de filosofia; relaciona historia; relaciona testes; relaciona goldchain\_mercado\_hub; relaciona expressividade; relaciona word; relaciona vontade\_de\_poder}
\prov{proveniência: wiki \textperiodcentered{} fato \textperiodcentered{} confiança 1.0 \textperiodcentered{} domínio ciencias\_de\_fora \textperiodcentered{} história 8 versões no jornal}

%CRISTAL {"arestas":[{"alvo":"bioetica","confianca":1.0,"epistemico":"fato","origem":"wiki","rel":"parte_de"}],"confianca":1.0,"contraexemplos":[],"descricao":"A linhagem normativa que protege os sujeitos de pesquisa: o Código de Nuremberg (1947, o consentimento voluntário como 'absolutamente essencial'), a Declaração de Helsinque (revisão por comitê, risco/benefício) e o Relatório Belmont (respeito, beneficência, justiça).","epistemico":"fato","exemplos":[],"id":"etica_pesquisa_humanos","memoria":"semantica","meta":{"autor":"Nuremberg/Helsinque/Belmont","dominio":"ciencias_de_fora","fonte":""},"origem":"wiki","palavras_chave":[],"sinonimos":[],"tipo":"conceito","titulo":"Ética da Pesquisa com Humanos"}
\section{Ética da Pesquisa com Humanos}
A linhagem normativa que protege os sujeitos de pesquisa: o Código de Nuremberg (1947, o consentimento voluntário como 'absolutamente essencial'), a Declaração de Helsinque (revisão por comitê, risco/benefício) e o Relatório Belmont (respeito, beneficência, justiça).
\prov{relações: parte\_de bioetica}
\prov{proveniência: wiki \textperiodcentered{} fato \textperiodcentered{} confiança 1.0 \textperiodcentered{} domínio ciencias\_de\_fora \textperiodcentered{} história 2 versões no jornal}

%CRISTAL {"arestas":[{"alvo":"acao_social_weber","confianca":1.0,"epistemico":"fato","origem":"wiki","rel":"usa"},{"alvo":"economia","confianca":1.0,"epistemico":"fato","origem":"wiki","rel":"referencia"},{"alvo":"misticismo_nao_dual","confianca":0.5,"epistemico":"fato","origem":"wiki","rel":"relaciona"},{"alvo":"word","confianca":0.5,"epistemico":"fato","origem":"wiki","rel":"relaciona"},{"alvo":"vida-morte","confianca":0.5,"epistemico":"fato","origem":"wiki","rel":"relaciona"},{"alvo":"vontade_de_poder","confianca":0.5,"epistemico":"fato","origem":"wiki","rel":"relaciona"}],"confianca":1.0,"contraexemplos":[],"descricao":"Weber: a ascese calvinista (o trabalho como sinal de graça) foi origem do espírito do capitalismo.","epistemico":"inferencia","exemplos":[],"id":"etica_protestante","memoria":"semantica","meta":{"autor":"Weber","dominio":"ciencias_de_fora","fonte":""},"origem":"wiki","palavras_chave":[],"sinonimos":[],"tipo":"conceito","titulo":"A Ética Protestante"}
\section{A Ética Protestante}
Weber: a ascese calvinista (o trabalho como sinal de graça) foi origem do espírito do capitalismo.
\prov{relações: usa acao\_social\_weber; referencia economia; relaciona misticismo\_nao\_dual; relaciona word; relaciona vida-morte; relaciona vontade\_de\_poder}
\prov{proveniência: wiki \textperiodcentered{} inferencia \textperiodcentered{} confiança 1.0 \textperiodcentered{} domínio ciencias\_de\_fora \textperiodcentered{} história 7 versões no jornal}

%CRISTAL {"arestas":[{"alvo":"didatica_matematica","confianca":1.0,"epistemico":"fato","origem":"wiki","rel":"parte_de"},{"alvo":"freudenthal_rme","confianca":1.0,"epistemico":"fato","origem":"wiki","rel":"usa"},{"alvo":"antropologia","confianca":1.0,"epistemico":"fato","origem":"wiki","rel":"referencia"}],"confianca":1.0,"contraexemplos":[],"descricao":"D'Ambrosio: a matemática praticada por grupos culturais identificáveis (povos, profissões, idades) — um programa sobre a geração e difusão do saber matemático em contextos culturais, com dimensão cognitiva, histórica e política.","epistemico":"inferencia","exemplos":[],"id":"etnomatematica","memoria":"semantica","meta":{"autor":"D'Ambrosio","dominio":"ciencias_de_fora","fonte":""},"origem":"wiki","palavras_chave":[],"sinonimos":[],"tipo":"conceito","titulo":"Etnomatemática"}
\section{Etnomatemática}
D'Ambrosio: a matemática praticada por grupos culturais identificáveis (povos, profissões, idades) — um programa sobre a geração e difusão do saber matemático em contextos culturais, com dimensão cognitiva, histórica e política.
\prov{relações: parte\_de didatica\_matematica; usa freudenthal\_rme; referencia antropologia}
\prov{proveniência: wiki \textperiodcentered{} inferencia \textperiodcentered{} confiança 1.0 \textperiodcentered{} domínio ciencias\_de\_fora \textperiodcentered{} história 4 versões no jornal}

%CRISTAL {"arestas":[{"alvo":"musica","confianca":1.0,"epistemico":"fato","origem":"wiki","rel":"parte_de"},{"alvo":"consonancia_dissonancia","confianca":1.0,"epistemico":"fato","origem":"wiki","rel":"relaciona"}],"confianca":1.0,"contraexemplos":[],"descricao":"O estudo da 'música na cultura e como cultura' (Merriam): a música é um universal humano cujas escalas e noções de consonância são moldadas culturalmente, questionando a pretensa universalidade dos padrões ocidentais.","epistemico":"inferencia","exemplos":[],"id":"etnomusicologia","memoria":"semantica","meta":{"autor":"Merriam/Nettl","dominio":"ciencias_de_fora","fonte":""},"origem":"wiki","palavras_chave":[],"sinonimos":[],"tipo":"conceito","titulo":"Etnomusicologia"}
\section{Etnomusicologia}
O estudo da 'música na cultura e como cultura' (Merriam): a música é um universal humano cujas escalas e noções de consonância são moldadas culturalmente, questionando a pretensa universalidade dos padrões ocidentais.
\prov{relações: parte\_de musica; relaciona consonancia\_dissonancia}
\prov{proveniência: wiki \textperiodcentered{} inferencia \textperiodcentered{} confiança 1.0 \textperiodcentered{} domínio ciencias\_de\_fora \textperiodcentered{} história 3 versões no jornal}

%CRISTAL {"arestas":[{"alvo":"etica","confianca":1.0,"epistemico":"fato","origem":"wiki","rel":"parte_de"},{"alvo":"utilitarismo","confianca":1.0,"epistemico":"fato","origem":"wiki","rel":"contradiz"},{"alvo":"teoria_dos_numeros","confianca":0.5,"epistemico":"fato","origem":"wiki","rel":"relaciona"},{"alvo":"word","confianca":0.5,"epistemico":"fato","origem":"wiki","rel":"relaciona"},{"alvo":"vontade_de_poder","confianca":0.5,"epistemico":"fato","origem":"wiki","rel":"relaciona"}],"confianca":1.0,"contraexemplos":[],"descricao":"Aristóteles: o florescimento humano — a atividade da alma segundo a virtude (arete), não o mero prazer.","epistemico":"inferencia","exemplos":[],"id":"eudaimonia","memoria":"semantica","meta":{"autor":"Aristóteles","dominio":"ciencias_de_fora","fonte":""},"origem":"wiki","palavras_chave":[],"sinonimos":[],"tipo":"conceito","titulo":"Eudaimonia"}
\section{Eudaimonia}
Aristóteles: o florescimento humano — a atividade da alma segundo a virtude (arete), não o mero prazer.
\prov{relações: parte\_de etica; contradiz utilitarismo; relaciona teoria\_dos\_numeros; relaciona word; relaciona vontade\_de\_poder}
\prov{proveniência: wiki \textperiodcentered{} inferencia \textperiodcentered{} confiança 1.0 \textperiodcentered{} domínio ciencias\_de\_fora \textperiodcentered{} história 6 versões no jornal}

%CRISTAL {"arestas":[{"alvo":"matematica","confianca":1.0,"epistemico":"fato","origem":"wiki","rel":"parte_de"},{"alvo":"analise_matematica","confianca":1.0,"epistemico":"fato","origem":"wiki","rel":"relaciona"}],"confianca":1.0,"contraexemplos":[],"descricao":"Euler: eiπ+1=0, grafos, topologia, análise — o matemático mais produtivo da história, notação que usamos até hoje.","epistemico":"fato","exemplos":[],"id":"euler_prolifico","memoria":"semantica","meta":{"autor":"Euler","dominio":"ciencias_de_fora","fonte":""},"origem":"wiki","palavras_chave":[],"sinonimos":[],"tipo":"conceito","titulo":"Euler, o Prolífico"}
\section{Euler, o Prolífico}
Euler: ei\ensuremath{\pi}+1=0, grafos, topologia, análise — o matemático mais produtivo da história, notação que usamos até hoje.
\prov{relações: parte\_de matematica; relaciona analise\_matematica}
\prov{proveniência: wiki \textperiodcentered{} fato \textperiodcentered{} confiança 1.0 \textperiodcentered{} domínio ciencias\_de\_fora \textperiodcentered{} história 7 versões no jornal}

%CRISTAL {"arestas":[{"alvo":"filosofia","confianca":1.0,"epistemico":"fato","origem":"wiki","rel":"parte_de"},{"alvo":"sentido_da_vida","confianca":1.0,"epistemico":"fato","origem":"wiki","rel":"relaciona"},{"alvo":"deus_postulado","confianca":1.0,"epistemico":"fato","origem":"wiki","rel":"relaciona"}],"confianca":1.0,"contraexemplos":[],"descricao":"Sartre: não há natureza humana fixa; o humano se define por seus atos, e aceitar uma 'essência' dada (Deus, destino) é má-fé. Ecoa o 'sentido da vida' do Lopes.","epistemico":"inferencia","exemplos":[],"id":"existencia_precede_essencia","memoria":"semantica","meta":{"autor":"Sartre","dominio":"ciencias_de_fora","fonte":""},"origem":"wiki","palavras_chave":[],"sinonimos":[],"tipo":"conceito","titulo":"Existência Precede Essência"}
\section{Existência Precede Essência}
Sartre: não há natureza humana fixa; o humano se define por seus atos, e aceitar uma 'essência' dada (Deus, destino) é má-fé. Ecoa o 'sentido da vida' do Lopes.
\prov{relações: parte\_de filosofia; relaciona sentido\_da\_vida; relaciona deus\_postulado}
\prov{proveniência: wiki \textperiodcentered{} inferencia \textperiodcentered{} confiança 1.0 \textperiodcentered{} domínio ciencias\_de\_fora \textperiodcentered{} história 5 versões no jornal}

%CRISTAL {"arestas":[{"alvo":"gene","confianca":1.0,"epistemico":"fato","origem":"wiki","rel":"usa"}],"confianca":1.0,"contraexemplos":[],"descricao":"Quando e quanto um gene é lido — a regulação que faz células idênticas em genoma virarem neurônio ou pele.","epistemico":"fato","exemplos":[],"id":"expressao_genica","memoria":"semantica","meta":{"dominio":"ciencias_de_fora","fonte":""},"origem":"wiki","palavras_chave":[],"sinonimos":[],"tipo":"conceito","titulo":"Expressão Gênica"}
\section{Expressão Gênica}
Quando e quanto um gene é lido — a regulação que faz células idênticas em genoma virarem neurônio ou pele.
\prov{relações: usa gene}
\prov{proveniência: wiki \textperiodcentered{} fato \textperiodcentered{} confiança 1.0 \textperiodcentered{} domínio ciencias\_de\_fora \textperiodcentered{} história 2 versões no jornal}

%CRISTAL {"arestas":[{"alvo":"meio_e_mensagem","confianca":1.0,"epistemico":"fato","origem":"wiki","rel":"usa"},{"alvo":"kapp_organprojektion","confianca":1.0,"epistemico":"fato","origem":"wiki","rel":"relaciona"},{"alvo":"tiffany","confianca":1.0,"epistemico":"fato","origem":"wiki","rel":"relaciona"}],"confianca":1.0,"contraexemplos":[],"descricao":"McLuhan: todo meio é uma extensão/amplificação de uma faculdade humana (a roda estende o pé, a mídia eletrônica o sistema nervoso), com efeitos de 'amputação' correlatos.","epistemico":"inferencia","exemplos":[],"id":"extensoes_do_homem","memoria":"semantica","meta":{"autor":"McLuhan","dominio":"ciencias_de_fora","fonte":""},"origem":"wiki","palavras_chave":[],"sinonimos":[],"tipo":"conceito","titulo":"Meios como Extensões do Homem (McLuhan)"}
\section{Meios como Extensões do Homem (McLuhan)}
McLuhan: todo meio é uma extensão/amplificação de uma faculdade humana (a roda estende o pé, a mídia eletrônica o sistema nervoso), com efeitos de 'amputação' correlatos.
\prov{relações: usa meio\_e\_mensagem; relaciona kapp\_organprojektion; relaciona tiffany}
\prov{proveniência: wiki \textperiodcentered{} inferencia \textperiodcentered{} confiança 1.0 \textperiodcentered{} domínio ciencias\_de\_fora \textperiodcentered{} história 4 versões no jornal}

%CRISTAL {"arestas":[{"alvo":"apolineo_dionisiaco","confianca":1.0,"epistemico":"fato","origem":"wiki","rel":"usa"},{"alvo":"colapso","confianca":1.0,"epistemico":"fato","origem":"wiki","rel":"referencia"}],"confianca":1.0,"contraexemplos":[],"descricao":"Lopes: a morte tem duas faces — a separação (sombria) e a renovação (bela); a ciência avança quando a velha geração morre (via Planck).","epistemico":"hipotese","exemplos":[],"id":"face_bela_morte","memoria":"semantica","meta":{"autor":"Gentil Lopes","dominio":"ciencias_de_fora","fonte":""},"origem":"livro","palavras_chave":[],"sinonimos":[],"tipo":"conceito","titulo":"A Face Bela (da Morte)"}
\section{A Face Bela (da Morte)}
Lopes: a morte tem duas faces — a separação (sombria) e a renovação (bela); a ciência avança quando a velha geração morre (via Planck).
\prov{relações: usa apolineo\_dionisiaco; referencia colapso}
\prov{proveniência: livro \textperiodcentered{} hipotese \textperiodcentered{} confiança 1.0 \textperiodcentered{} domínio ciencias\_de\_fora \textperiodcentered{} história 4 versões no jornal}

%CRISTAL {"arestas":[{"alvo":"traducao","confianca":1.0,"epistemico":"fato","origem":"wiki","rel":"relaciona"},{"alvo":"gramatica_espanhol","confianca":1.0,"epistemico":"fato","origem":"wiki","rel":"relaciona"}],"confianca":1.0,"contraexemplos":[],"descricao":"Palavras parecidas entre línguas com sentidos diferentes ('embarazada'=grávida, não embaraçada) — a armadilha do português↔espanhol. A forma engana; o sentido é o que importa.","epistemico":"fato","exemplos":[],"id":"falso_amigo","memoria":"semantica","meta":{"dominio":"ciencias_de_fora","fonte":""},"origem":"wiki","palavras_chave":[],"sinonimos":[],"tipo":"conceito","titulo":"Falso Amigo"}
\section{Falso Amigo}
Palavras parecidas entre línguas com sentidos diferentes ('embarazada'=grávida, não embaraçada) — a armadilha do português\ensuremath{\leftrightarrow}espanhol. A forma engana; o sentido é o que importa.
\prov{relações: relaciona traducao; relaciona gramatica\_espanhol}
\prov{proveniência: wiki \textperiodcentered{} fato \textperiodcentered{} confiança 1.0 \textperiodcentered{} domínio ciencias\_de\_fora \textperiodcentered{} história 3 versões no jornal}

%CRISTAL {"arestas":[{"alvo":"idioma","confianca":1.0,"epistemico":"fato","origem":"wiki","rel":"usa"},{"alvo":"protolingua","confianca":1.0,"epistemico":"fato","origem":"wiki","rel":"usa"},{"alvo":"linguistica","confianca":1.0,"epistemico":"fato","origem":"wiki","rel":"parte_de"}],"confianca":1.0,"contraexemplos":[],"descricao":"Um grupo de línguas descendentes de um ancestral comum (proto-língua) — parentesco provado por correspondências sistemáticas de som.","epistemico":"fato","exemplos":[],"id":"familia_linguistica","memoria":"semantica","meta":{"dominio":"ciencias_de_fora","fonte":""},"origem":"wiki","palavras_chave":[],"sinonimos":[],"tipo":"conceito","titulo":"Família Linguística"}
\section{Família Linguística}
Um grupo de línguas descendentes de um ancestral comum (proto-língua) — parentesco provado por correspondências sistemáticas de som.
\prov{relações: usa idioma; usa protolingua; parte\_de linguistica}
\prov{proveniência: wiki \textperiodcentered{} fato \textperiodcentered{} confiança 1.0 \textperiodcentered{} domínio ciencias\_de\_fora \textperiodcentered{} história 4 versões no jornal}

%CRISTAL {"arestas":[{"alvo":"autoritarismo_vs_totalitarismo","confianca":1.0,"epistemico":"fato","origem":"wiki","rel":"usa"},{"alvo":"populismo","confianca":1.0,"epistemico":"fato","origem":"wiki","rel":"relaciona"},{"alvo":"origens_totalitarismo","confianca":1.0,"epistemico":"fato","origem":"wiki","rel":"relaciona"}],"confianca":1.0,"contraexemplos":[],"descricao":"Paxton: um comportamento político de obsessão com o declínio/humilhação da comunidade e cultos compensatórios de unidade e pureza, em que um partido de massas nacionalista, aliado a elites, abandona as liberdades democráticas e busca, com violência redentora, a 'limpeza' interna e a expansão externa.","epistemico":"inferencia","exemplos":[],"id":"fascismo","memoria":"semantica","meta":{"autor":"Paxton","dominio":"ciencias_de_fora","fonte":""},"origem":"wiki","palavras_chave":[],"sinonimos":[],"tipo":"conceito","titulo":"Fascismo (Paxton)"}
\section{Fascismo (Paxton)}
Paxton: um comportamento político de obsessão com o declínio/humilhação da comunidade e cultos compensatórios de unidade e pureza, em que um partido de massas nacionalista, aliado a elites, abandona as liberdades democráticas e busca, com violência redentora, a 'limpeza' interna e a expansão externa.
\prov{relações: usa autoritarismo\_vs\_totalitarismo; relaciona populismo; relaciona origens\_totalitarismo}
\prov{proveniência: wiki \textperiodcentered{} inferencia \textperiodcentered{} confiança 1.0 \textperiodcentered{} domínio ciencias\_de\_fora \textperiodcentered{} história 4 versões no jornal}

%CRISTAL {"arestas":[{"alvo":"sociologia","confianca":1.0,"epistemico":"fato","origem":"wiki","rel":"parte_de"},{"alvo":"virtualizacao","confianca":0.5,"epistemico":"fato","origem":"wiki","rel":"relaciona"},{"alvo":"word","confianca":0.5,"epistemico":"fato","origem":"wiki","rel":"relaciona"},{"alvo":"vida-morte","confianca":0.5,"epistemico":"fato","origem":"wiki","rel":"relaciona"}],"confianca":1.0,"contraexemplos":[],"descricao":"Durkheim: modos de agir e pensar exteriores ao indivíduo e dotados de poder de coação; a estrutura que se impõe.","epistemico":"inferencia","exemplos":[],"id":"fato_social","memoria":"semantica","meta":{"autor":"Durkheim","dominio":"ciencias_de_fora","fonte":""},"origem":"wiki","palavras_chave":[],"sinonimos":[],"tipo":"conceito","titulo":"Fato Social"}
\section{Fato Social}
Durkheim: modos de agir e pensar exteriores ao indivíduo e dotados de poder de coação; a estrutura que se impõe.
\prov{relações: parte\_de sociologia; relaciona virtualizacao; relaciona word; relaciona vida-morte}
\prov{proveniência: wiki \textperiodcentered{} inferencia \textperiodcentered{} confiança 1.0 \textperiodcentered{} domínio ciencias\_de\_fora \textperiodcentered{} história 5 versões no jornal}

%CRISTAL {"arestas":[{"alvo":"cibernetica","confianca":1.0,"epistemico":"fato","origem":"wiki","rel":"parte_de"},{"alvo":"hopfield","confianca":1.0,"epistemico":"fato","origem":"wiki","rel":"referencia"},{"alvo":"colapso","confianca":1.0,"epistemico":"fato","origem":"wiki","rel":"relaciona"}],"confianca":1.0,"contraexemplos":[],"descricao":"Wiener: a saída de um sistema retorna como entrada corrigindo seu comportamento — o feedback negativo produz estabilidade/homeostase, o positivo amplifica. O núcleo do controle cibernético.","epistemico":"fato","exemplos":[],"id":"feedback_retroalimentacao","memoria":"semantica","meta":{"autor":"Wiener","dominio":"ciencias_de_fora","fonte":""},"origem":"wiki","palavras_chave":[],"sinonimos":[],"tipo":"conceito","titulo":"Feedback / Retroalimentação"}
\section{Feedback / Retroalimentação}
Wiener: a saída de um sistema retorna como entrada corrigindo seu comportamento — o feedback negativo produz estabilidade/homeostase, o positivo amplifica. O núcleo do controle cibernético.
\prov{relações: parte\_de cibernetica; referencia hopfield; relaciona colapso}
\prov{proveniência: wiki \textperiodcentered{} fato \textperiodcentered{} confiança 1.0 \textperiodcentered{} domínio ciencias\_de\_fora \textperiodcentered{} história 4 versões no jornal}

%CRISTAL {"arestas":[{"alvo":"linguistica","confianca":1.0,"epistemico":"fato","origem":"wiki","rel":"parte_de"},{"alvo":"estruturalismo_linguistico","confianca":1.0,"epistemico":"fato","origem":"wiki","rel":"explica"}],"confianca":1.0,"contraexemplos":[],"descricao":"O pai da linguística moderna e do estruturalismo: a língua como sistema de signos e de diferenças.","epistemico":"fato","exemplos":[],"id":"ferdinand_saussure","memoria":"semantica","meta":{"autor":"Saussure","dominio":"ciencias_de_fora","fonte":""},"origem":"wiki","palavras_chave":[],"sinonimos":[],"tipo":"conceito","titulo":"Ferdinand de Saussure"}
\section{Ferdinand de Saussure}
O pai da linguística moderna e do estruturalismo: a língua como sistema de signos e de diferenças.
\prov{relações: parte\_de linguistica; explica estruturalismo\_linguistico}
\prov{proveniência: wiki \textperiodcentered{} fato \textperiodcentered{} confiança 1.0 \textperiodcentered{} domínio ciencias\_de\_fora \textperiodcentered{} história 3 versões no jornal}

%CRISTAL {"arestas":[],"confianca":1.0,"contraexemplos":[],"descricao":"O HUB da ficção: as obras (jogos, animações, histórias) que ressoam com o broca-so pela lente do DICIONÁRIO (o Teorema Geral da Tradução: mesmo operador, o dicionário troca). Não são o motor — são leituras paralelas: uma obra vira uma instância de um corpo do broca-so. Sob o hub, o sub-hub JOGOS E ANIMAÇÕES (o Minecraft, o Ben 10). Escopo: as obras são fato cultural; a leitura é síntese declarada.","epistemico":"fato","exemplos":[],"id":"ficcao_hub","memoria":"semantica","meta":{"dominio":"ficcao","fonte":"Ficção (jogos e animações)"},"origem":"wiki","palavras_chave":[],"sinonimos":[],"tipo":"conceito","titulo":"Ficção — o hub das obras (lidas pelo dicionário paralelo)"}
\section{Ficção — o hub das obras (lidas pelo dicionário paralelo)}
O HUB da ficção: as obras (jogos, animações, histórias) que ressoam com o broca-so pela lente do DICIONÁRIO (o Teorema Geral da Tradução: mesmo operador, o dicionário troca). Não são o motor — são leituras paralelas: uma obra vira uma instância de um corpo do broca-so. Sob o hub, o sub-hub JOGOS E ANIMAÇÕES (o Minecraft, o Ben 10). Escopo: as obras são fato cultural; a leitura é síntese declarada.
\prov{proveniência: wiki \textperiodcentered{} Ficção (jogos e animações) \textperiodcentered{} fato \textperiodcentered{} confiança 1.0 \textperiodcentered{} domínio ficcao}

%CRISTAL {"arestas":[{"alvo":"nominalismo_matematico","confianca":1.0,"epistemico":"fato","origem":"wiki","rel":"especializa"},{"alvo":"word","confianca":0.5,"epistemico":"fato","origem":"wiki","rel":"relaciona"},{"alvo":"virtualizacao","confianca":0.5,"epistemico":"fato","origem":"wiki","rel":"relaciona"},{"alvo":"vida-morte","confianca":0.5,"epistemico":"fato","origem":"wiki","rel":"relaciona"},{"alvo":"misticismo_nao_dual","confianca":0.5,"epistemico":"fato","origem":"wiki","rel":"relaciona"}],"confianca":1.0,"contraexemplos":[],"descricao":"Os enunciados matemáticos são ficções úteis mas literalmente falsas — indispensáveis à ciência (Field).","epistemico":"inferencia","exemplos":[],"id":"ficcionalismo_matematico","memoria":"semantica","meta":{"autor":"Field","dominio":"filosofia","fonte":""},"origem":"wiki","palavras_chave":[],"sinonimos":[],"tipo":"conceito","titulo":"Ficcionalismo"}
\section{Ficcionalismo}
Os enunciados matemáticos são ficções úteis mas literalmente falsas — indispensáveis à ciência (Field).
\prov{relações: especializa nominalismo\_matematico; relaciona word; relaciona virtualizacao; relaciona vida-morte; relaciona misticismo\_nao\_dual}
\prov{proveniência: wiki \textperiodcentered{} inferencia \textperiodcentered{} confiança 1.0 \textperiodcentered{} domínio filosofia \textperiodcentered{} história 6 versões no jornal}

%CRISTAL {"arestas":[{"alvo":"traducao","confianca":1.0,"epistemico":"fato","origem":"wiki","rel":"usa"},{"alvo":"traducao_literal_vs_livre","confianca":1.0,"epistemico":"fato","origem":"wiki","rel":"relaciona"}],"confianca":1.0,"contraexemplos":[],"descricao":"O eterno trade-off: colar na fonte (fiel, mas estranho) ou soar natural no alvo (fluente, mas afastado). Toda tradução escolhe um ponto entre os dois.","epistemico":"inferencia","exemplos":[],"id":"fidelidade_vs_fluencia","memoria":"semantica","meta":{"dominio":"ciencias_de_fora","fonte":""},"origem":"wiki","palavras_chave":[],"sinonimos":[],"tipo":"conceito","titulo":"Fidelidade vs Fluência"}
\section{Fidelidade vs Fluência}
O eterno trade-off: colar na fonte (fiel, mas estranho) ou soar natural no alvo (fluente, mas afastado). Toda tradução escolhe um ponto entre os dois.
\prov{relações: usa traducao; relaciona traducao\_literal\_vs\_livre}
\prov{proveniência: wiki \textperiodcentered{} inferencia \textperiodcentered{} confiança 1.0 \textperiodcentered{} domínio ciencias\_de\_fora \textperiodcentered{} história 3 versões no jornal}

%CRISTAL {"arestas":[{"alvo":"teoria_literaria","confianca":1.0,"epistemico":"fato","origem":"wiki","rel":"parte_de"},{"alvo":"metafora_metonimia","confianca":1.0,"epistemico":"fato","origem":"wiki","rel":"usa"},{"alvo":"funcao_poetica","confianca":1.0,"epistemico":"fato","origem":"wiki","rel":"relaciona"}],"confianca":1.0,"contraexemplos":[],"descricao":"Desvios que dizem mais — metáfora, metonímia, ironia, hipérbole; a linguagem dobrada para significar além do literal.","epistemico":"fato","exemplos":[],"id":"figuras_de_linguagem","memoria":"semantica","meta":{"dominio":"ciencias_de_fora","fonte":""},"origem":"wiki","palavras_chave":[],"sinonimos":[],"tipo":"conceito","titulo":"Figuras de Linguagem"}
\section{Figuras de Linguagem}
Desvios que dizem mais — metáfora, metonímia, ironia, hipérbole; a linguagem dobrada para significar além do literal.
\prov{relações: parte\_de teoria\_literaria; usa metafora\_metonimia; relaciona funcao\_poetica}
\prov{proveniência: wiki \textperiodcentered{} fato \textperiodcentered{} confiança 1.0 \textperiodcentered{} domínio ciencias\_de\_fora \textperiodcentered{} história 4 versões no jornal}

%CRISTAL {"arestas":[{"alvo":"matematica","confianca":1.0,"epistemico":"fato","origem":"wiki","rel":"relaciona"},{"alvo":"tiffany","confianca":1.0,"epistemico":"fato","origem":"wiki","rel":"referencia"},{"alvo":"word","confianca":0.5,"epistemico":"fato","origem":"wiki","rel":"relaciona"},{"alvo":"misticismo_nao_dual","confianca":0.5,"epistemico":"fato","origem":"wiki","rel":"relaciona"}],"confianca":1.0,"contraexemplos":[],"descricao":"Investigação racional das questões fundamentais: ser, conhecer, valor, mente — o hub das ciências de fora.","epistemico":"fato","exemplos":[],"id":"filosofia","memoria":"semantica","meta":{"dominio":"filosofia","fonte":""},"origem":"wiki","palavras_chave":[],"sinonimos":[],"tipo":"conceito","titulo":"Filosofia"}
\section{Filosofia}
Investigação racional das questões fundamentais: ser, conhecer, valor, mente — o hub das ciências de fora.
\prov{relações: relaciona matematica; referencia tiffany; relaciona word; relaciona misticismo\_nao\_dual}
\prov{proveniência: wiki \textperiodcentered{} fato \textperiodcentered{} confiança 1.0 \textperiodcentered{} domínio filosofia \textperiodcentered{} história 5 versões no jornal}

%CRISTAL {"arestas":[{"alvo":"filosofia_da_tecnica","confianca":1.0,"epistemico":"fato","origem":"wiki","rel":"parte_de"},{"alvo":"computacao","confianca":1.0,"epistemico":"fato","origem":"wiki","rel":"referencia"}],"confianca":1.0,"contraexemplos":[],"descricao":"A reflexão sobre a natureza e os limites do computável, da informação e da mente-máquina — o que é uma máquina que pensa, e o que ela não pode fazer.","epistemico":"fato","exemplos":[],"id":"filosofia_da_computacao","memoria":"semantica","meta":{"dominio":"ciencias_de_fora","fonte":""},"origem":"wiki","palavras_chave":[],"sinonimos":[],"tipo":"conceito","titulo":"Filosofia da Computação"}
\section{Filosofia da Computação}
A reflexão sobre a natureza e os limites do computável, da informação e da mente-máquina — o que é uma máquina que pensa, e o que ela não pode fazer.
\prov{relações: parte\_de filosofia\_da\_tecnica; referencia computacao}
\prov{proveniência: wiki \textperiodcentered{} fato \textperiodcentered{} confiança 1.0 \textperiodcentered{} domínio ciencias\_de\_fora \textperiodcentered{} história 3 versões no jornal}

%CRISTAL {"arestas":[{"alvo":"educacao","confianca":1.0,"epistemico":"fato","origem":"wiki","rel":"parte_de"},{"alvo":"filosofia","confianca":1.0,"epistemico":"fato","origem":"wiki","rel":"parte_de"},{"alvo":"vida-morte","confianca":0.5,"epistemico":"fato","origem":"wiki","rel":"relaciona"},{"alvo":"virtualizacao","confianca":0.5,"epistemico":"fato","origem":"wiki","rel":"relaciona"}],"confianca":1.0,"contraexemplos":[],"descricao":"A reflexão sobre os fins, valores e pressupostos da educação — para quê e para quem se educa, e o que é ensinar.","epistemico":"fato","exemplos":[],"id":"filosofia_da_educacao","memoria":"semantica","meta":{"dominio":"ciencias_de_fora","fonte":""},"origem":"wiki","palavras_chave":[],"sinonimos":[],"tipo":"conceito","titulo":"Filosofia da Educação"}
\section{Filosofia da Educação}
A reflexão sobre os fins, valores e pressupostos da educação — para quê e para quem se educa, e o que é ensinar.
\prov{relações: parte\_de educacao; parte\_de filosofia; relaciona vida-morte; relaciona virtualizacao}
\prov{proveniência: wiki \textperiodcentered{} fato \textperiodcentered{} confiança 1.0 \textperiodcentered{} domínio ciencias\_de\_fora \textperiodcentered{} história 5 versões no jornal}

%CRISTAL {"arestas":[{"alvo":"historia","confianca":1.0,"epistemico":"fato","origem":"wiki","rel":"parte_de"},{"alvo":"filosofia","confianca":1.0,"epistemico":"fato","origem":"wiki","rel":"parte_de"},{"alvo":"word","confianca":0.5,"epistemico":"fato","origem":"wiki","rel":"relaciona"},{"alvo":"virtualizacao","confianca":0.5,"epistemico":"fato","origem":"wiki","rel":"relaciona"},{"alvo":"syscall","confianca":0.5,"epistemico":"fato","origem":"wiki","rel":"relaciona"}],"confianca":1.0,"contraexemplos":[],"descricao":"Pergunta o padrão do tempo histórico — progresso, ciclo ou contingência; necessidade ou liberdade.","epistemico":"fato","exemplos":[],"id":"filosofia_da_historia","memoria":"semantica","meta":{"dominio":"ciencias_de_fora","fonte":""},"origem":"wiki","palavras_chave":[],"sinonimos":[],"tipo":"conceito","titulo":"Filosofia da História"}
\section{Filosofia da História}
Pergunta o padrão do tempo histórico — progresso, ciclo ou contingência; necessidade ou liberdade.
\prov{relações: parte\_de historia; parte\_de filosofia; relaciona word; relaciona virtualizacao; relaciona syscall}
\prov{proveniência: wiki \textperiodcentered{} fato \textperiodcentered{} confiança 1.0 \textperiodcentered{} domínio ciencias\_de\_fora \textperiodcentered{} história 6 versões no jornal}

%CRISTAL {"arestas":[{"alvo":"vetores","confianca":0.5,"epistemico":"fato","origem":"wiki","rel":"relaciona"},{"alvo":"teoria_do_valor","confianca":0.5,"epistemico":"fato","origem":"wiki","rel":"relaciona"},{"alvo":"word","confianca":0.5,"epistemico":"fato","origem":"wiki","rel":"relaciona"},{"alvo":"vontade_de_poder","confianca":0.5,"epistemico":"fato","origem":"wiki","rel":"relaciona"},{"alvo":"linguistica","confianca":1.0,"epistemico":"fato","origem":"wiki","rel":"relaciona"},{"alvo":"filosofia","confianca":1.0,"epistemico":"fato","origem":"wiki","rel":"parte_de"}],"confianca":1.0,"contraexemplos":[],"descricao":"Pergunta o que é significar, referir e comunicar — a ponte entre mente, mundo e palavra.","epistemico":"fato","exemplos":[],"id":"filosofia_da_linguagem","memoria":"semantica","meta":{"dominio":"ciencias_de_fora","fonte":""},"origem":"wiki","palavras_chave":[],"sinonimos":[],"tipo":"conceito","titulo":"Filosofia da Linguagem"}
\section{Filosofia da Linguagem}
Pergunta o que é significar, referir e comunicar — a ponte entre mente, mundo e palavra.
\prov{relações: relaciona vetores; relaciona teoria\_do\_valor; relaciona word; relaciona vontade\_de\_poder; relaciona linguistica; parte\_de filosofia}
\prov{proveniência: wiki \textperiodcentered{} fato \textperiodcentered{} confiança 1.0 \textperiodcentered{} domínio ciencias\_de\_fora \textperiodcentered{} história 9 versões no jornal}

%CRISTAL {"arestas":[{"alvo":"filosofia","confianca":1.0,"epistemico":"fato","origem":"wiki","rel":"parte_de"},{"alvo":"matematica","confianca":1.0,"epistemico":"fato","origem":"wiki","rel":"referencia"},{"alvo":"word","confianca":0.5,"epistemico":"fato","origem":"wiki","rel":"relaciona"},{"alvo":"virtualizacao","confianca":0.5,"epistemico":"fato","origem":"wiki","rel":"relaciona"},{"alvo":"vida-morte","confianca":0.5,"epistemico":"fato","origem":"wiki","rel":"relaciona"},{"alvo":"sagrado_profano","confianca":0.5,"epistemico":"fato","origem":"wiki","rel":"relaciona"}],"confianca":1.0,"contraexemplos":[],"descricao":"Pergunta a natureza dos objetos matemáticos e a verdade matemática — a matemática é descoberta ou inventada?","epistemico":"fato","exemplos":[],"id":"filosofia_da_matematica","memoria":"semantica","meta":{"dominio":"filosofia","fonte":""},"origem":"wiki","palavras_chave":[],"sinonimos":[],"tipo":"conceito","titulo":"Filosofia da Matemática"}
\section{Filosofia da Matemática}
Pergunta a natureza dos objetos matemáticos e a verdade matemática — a matemática é descoberta ou inventada?
\prov{relações: parte\_de filosofia; referencia matematica; relaciona word; relaciona virtualizacao; relaciona vida-morte; relaciona sagrado\_profano}
\prov{proveniência: wiki \textperiodcentered{} fato \textperiodcentered{} confiança 1.0 \textperiodcentered{} domínio filosofia \textperiodcentered{} história 7 versões no jornal}

%CRISTAL {"arestas":[{"alvo":"medicina","confianca":1.0,"epistemico":"fato","origem":"wiki","rel":"parte_de"},{"alvo":"filosofia","confianca":1.0,"epistemico":"fato","origem":"wiki","rel":"parte_de"},{"alvo":"word","confianca":0.5,"epistemico":"fato","origem":"wiki","rel":"relaciona"}],"confianca":1.0,"contraexemplos":[],"descricao":"A pergunta sobre o que são saúde e doença — fato biológico objetivo ou desvio carregado de valor —, e sobre os modelos do adoecer.","epistemico":"fato","exemplos":[],"id":"filosofia_da_medicina","memoria":"semantica","meta":{"dominio":"ciencias_de_fora","fonte":""},"origem":"wiki","palavras_chave":[],"sinonimos":[],"tipo":"conceito","titulo":"Filosofia da Medicina"}
\section{Filosofia da Medicina}
A pergunta sobre o que são saúde e doença — fato biológico objetivo ou desvio carregado de valor —, e sobre os modelos do adoecer.
\prov{relações: parte\_de medicina; parte\_de filosofia; relaciona word}
\prov{proveniência: wiki \textperiodcentered{} fato \textperiodcentered{} confiança 1.0 \textperiodcentered{} domínio ciencias\_de\_fora \textperiodcentered{} história 4 versões no jornal}

%CRISTAL {"arestas":[{"alvo":"filosofia","confianca":1.0,"epistemico":"fato","origem":"wiki","rel":"parte_de"},{"alvo":"nao_dualidade","confianca":0.5,"epistemico":"fato","origem":"wiki","rel":"relaciona"},{"alvo":"politica","confianca":0.5,"epistemico":"fato","origem":"wiki","rel":"relaciona"},{"alvo":"metafisica","confianca":0.5,"epistemico":"fato","origem":"wiki","rel":"relaciona"},{"alvo":"ipc","confianca":0.5,"epistemico":"fato","origem":"wiki","rel":"relaciona"},{"alvo":"gato_eh_um_mamifero","confianca":0.5,"epistemico":"fato","origem":"wiki","rel":"relaciona"}],"confianca":1.0,"contraexemplos":[],"descricao":"Estuda a mente, a consciência, a intencionalidade e sua relação com o corpo/cérebro.","epistemico":"fato","exemplos":[],"id":"filosofia_da_mente","memoria":"semantica","meta":{"dominio":"filosofia","fonte":""},"origem":"wiki","palavras_chave":[],"sinonimos":[],"tipo":"conceito","titulo":"Filosofia da Mente"}
\section{Filosofia da Mente}
Estuda a mente, a consciência, a intencionalidade e sua relação com o corpo/cérebro.
\prov{relações: parte\_de filosofia; relaciona nao\_dualidade; relaciona politica; relaciona metafisica; relaciona ipc; relaciona gato\_eh\_um\_mamifero}
\prov{proveniência: wiki \textperiodcentered{} fato \textperiodcentered{} confiança 1.0 \textperiodcentered{} domínio filosofia \textperiodcentered{} história 7 versões no jornal}

%CRISTAL {"arestas":[{"alvo":"filosofia","confianca":1.0,"epistemico":"fato","origem":"wiki","rel":"parte_de"},{"alvo":"broca_os","confianca":1.0,"epistemico":"fato","origem":"wiki","rel":"relaciona"}],"confianca":1.0,"contraexemplos":[],"descricao":"A reflexão sobre o que é a técnica — não seus mecanismos, mas sua essência, seu poder e seu sentido: ferramenta neutra, modo de desvelar o mundo ou força autônoma?","epistemico":"fato","exemplos":[],"id":"filosofia_da_tecnica","memoria":"semantica","meta":{"dominio":"ciencias_de_fora","fonte":""},"origem":"wiki","palavras_chave":[],"sinonimos":[],"tipo":"conceito","titulo":"Filosofia da Técnica"}
\section{Filosofia da Técnica}
A reflexão sobre o que é a técnica — não seus mecanismos, mas sua essência, seu poder e seu sentido: ferramenta neutra, modo de desvelar o mundo ou força autônoma?
\prov{relações: parte\_de filosofia; relaciona broca\_os}
\prov{proveniência: wiki \textperiodcentered{} fato \textperiodcentered{} confiança 1.0 \textperiodcentered{} domínio ciencias\_de\_fora \textperiodcentered{} história 3 versões no jornal}

%CRISTAL {"arestas":[{"alvo":"direito","confianca":1.0,"epistemico":"fato","origem":"wiki","rel":"parte_de"},{"alvo":"filosofia","confianca":1.0,"epistemico":"fato","origem":"wiki","rel":"parte_de"},{"alvo":"word","confianca":0.5,"epistemico":"fato","origem":"wiki","rel":"relaciona"}],"confianca":1.0,"contraexemplos":[],"descricao":"A pergunta sobre o que é o direito — direito natural ou posto, norma ou fato social, regra ou princípio, sistema fechado ou aberto à moral.","epistemico":"fato","exemplos":[],"id":"filosofia_do_direito","memoria":"semantica","meta":{"dominio":"ciencias_de_fora","fonte":""},"origem":"wiki","palavras_chave":[],"sinonimos":[],"tipo":"conceito","titulo":"Filosofia do Direito"}
\section{Filosofia do Direito}
A pergunta sobre o que é o direito — direito natural ou posto, norma ou fato social, regra ou princípio, sistema fechado ou aberto à moral.
\prov{relações: parte\_de direito; parte\_de filosofia; relaciona word}
\prov{proveniência: wiki \textperiodcentered{} fato \textperiodcentered{} confiança 1.0 \textperiodcentered{} domínio ciencias\_de\_fora \textperiodcentered{} história 4 versões no jornal}

%CRISTAL {"arestas":[{"alvo":"filosofia_da_computacao","confianca":1.0,"epistemico":"fato","origem":"wiki","rel":"parte_de"},{"alvo":"teoria_da_informacao_shannon","confianca":1.0,"epistemico":"fato","origem":"wiki","rel":"usa"},{"alvo":"pancomputacionalismo","confianca":1.0,"epistemico":"fato","origem":"wiki","rel":"relaciona"}],"confianca":1.0,"contraexemplos":[],"descricao":"Floridi: um programa que toma a informação — não o ser nem a matéria — como constituinte fundamental da realidade. A 4ª revolução (após Copérnico, Darwin, Freud): Turing remove nossa excepcionalidade como únicos processadores de informação.","epistemico":"inferencia","exemplos":[],"id":"filosofia_informacao_floridi","memoria":"semantica","meta":{"autor":"Floridi","dominio":"ciencias_de_fora","fonte":""},"origem":"wiki","palavras_chave":[],"sinonimos":[],"tipo":"conceito","titulo":"Filosofia da Informação (Floridi)"}
\section{Filosofia da Informação (Floridi)}
Floridi: um programa que toma a informação — não o ser nem a matéria — como constituinte fundamental da realidade. A 4ª revolução (após Copérnico, Darwin, Freud): Turing remove nossa excepcionalidade como únicos processadores de informação.
\prov{relações: parte\_de filosofia\_da\_computacao; usa teoria\_da\_informacao\_shannon; relaciona pancomputacionalismo}
\prov{proveniência: wiki \textperiodcentered{} inferencia \textperiodcentered{} confiança 1.0 \textperiodcentered{} domínio ciencias\_de\_fora \textperiodcentered{} história 4 versões no jornal}

%CRISTAL {"arestas":[{"alvo":"biologia","confianca":1.0,"epistemico":"fato","origem":"wiki","rel":"parte_de"},{"alvo":"fibonacci","confianca":1.0,"epistemico":"fato","origem":"wiki","rel":"usa"},{"alvo":"geometria_fractal","confianca":1.0,"epistemico":"fato","origem":"wiki","rel":"relaciona"},{"alvo":"floco_no_metal","confianca":1.0,"epistemico":"fato","origem":"wiki","rel":"relaciona"}],"confianca":1.0,"contraexemplos":[],"descricao":"Folhas e sementes se dispõem pelo ângulo de ~137,5° (a razão áurea) e em números de Fibonacci — a vida achou os metálicos.","epistemico":"fato","exemplos":[],"id":"filotaxia","memoria":"semantica","meta":{"dominio":"ciencias_de_fora","fonte":""},"origem":"wiki","palavras_chave":[],"sinonimos":[],"tipo":"conceito","titulo":"Filotaxia (o Ângulo Áureo)"}
\section{Filotaxia (o Ângulo Áureo)}
Folhas e sementes se dispõem pelo ângulo de \~{}137,5$^{\circ}$ (a razão áurea) e em números de Fibonacci — a vida achou os metálicos.
\prov{relações: parte\_de biologia; usa fibonacci; relaciona geometria\_fractal; relaciona floco\_no\_metal}
\prov{proveniência: wiki \textperiodcentered{} fato \textperiodcentered{} confiança 1.0 \textperiodcentered{} domínio ciencias\_de\_fora \textperiodcentered{} história 5 versões no jornal}

%CRISTAL {"arestas":[{"alvo":"dialetica_hegeliana","confianca":1.0,"epistemico":"fato","origem":"wiki","rel":"relaciona"},{"alvo":"ciclos_vico","confianca":1.0,"epistemico":"fato","origem":"wiki","rel":"contradiz"},{"alvo":"word","confianca":0.5,"epistemico":"fato","origem":"wiki","rel":"relaciona"},{"alvo":"vontade_de_poder","confianca":0.5,"epistemico":"fato","origem":"wiki","rel":"relaciona"},{"alvo":"while","confianca":0.5,"epistemico":"fato","origem":"wiki","rel":"relaciona"},{"alvo":"virtualizacao","confianca":0.5,"epistemico":"fato","origem":"wiki","rel":"relaciona"}],"confianca":1.0,"contraexemplos":[],"descricao":"Fukuyama: com o colapso soviético, a democracia liberal venceu as alternativas — fim do conflito ideológico. Historicismo refutável.","epistemico":"hipotese","exemplos":[],"id":"fim_da_historia","memoria":"semantica","meta":{"autor":"Fukuyama","dominio":"ciencias_de_fora","fonte":""},"origem":"wiki","palavras_chave":[],"sinonimos":[],"tipo":"conceito","titulo":"O Fim da História"}
\section{O Fim da História}
Fukuyama: com o colapso soviético, a democracia liberal venceu as alternativas — fim do conflito ideológico. Historicismo refutável.
\prov{relações: relaciona dialetica\_hegeliana; contradiz ciclos\_vico; relaciona word; relaciona vontade\_de\_poder; relaciona while; relaciona virtualizacao}
\prov{proveniência: wiki \textperiodcentered{} hipotese \textperiodcentered{} confiança 1.0 \textperiodcentered{} domínio ciencias\_de\_fora \textperiodcentered{} história 7 versões no jornal}

%CRISTAL {"arestas":[{"alvo":"bioetica","confianca":1.0,"epistemico":"fato","origem":"wiki","rel":"parte_de"},{"alvo":"face_bela_morte","confianca":1.0,"epistemico":"fato","origem":"wiki","rel":"relaciona"},{"alvo":"autonomia_do_paciente","confianca":1.0,"epistemico":"fato","origem":"wiki","rel":"usa"}],"confianca":1.0,"contraexemplos":[],"descricao":"Três posições sobre a morte: eutanásia é o ato intencional de causar a morte para abreviar sofrimento; distanásia é o prolongamento artificial e fútil da vida (obstinação terapêutica); ortotanásia é permitir a morte no seu tempo, suspendendo meios extraordinários.","epistemico":"inferencia","exemplos":[],"id":"fim_de_vida_espectro","memoria":"semantica","meta":{"autor":"Debray","dominio":"ciencias_de_fora","fonte":""},"origem":"wiki","palavras_chave":[],"sinonimos":[],"tipo":"conceito","titulo":"Eutanásia, Distanásia, Ortotanásia"}
\section{Eutanásia, Distanásia, Ortotanásia}
Três posições sobre a morte: eutanásia é o ato intencional de causar a morte para abreviar sofrimento; distanásia é o prolongamento artificial e fútil da vida (obstinação terapêutica); ortotanásia é permitir a morte no seu tempo, suspendendo meios extraordinários.
\prov{relações: parte\_de bioetica; relaciona face\_bela\_morte; usa autonomia\_do\_paciente}
\prov{proveniência: wiki \textperiodcentered{} inferencia \textperiodcentered{} confiança 1.0 \textperiodcentered{} domínio ciencias\_de\_fora \textperiodcentered{} história 4 versões no jornal}

%CRISTAL {"arestas":[{"alvo":"construcao_social_tecnologia","confianca":1.0,"epistemico":"fato","origem":"wiki","rel":"usa"},{"alvo":"mediacao_semiotica","confianca":1.0,"epistemico":"fato","origem":"wiki","rel":"relaciona"},{"alvo":"injustica_epistemica","confianca":1.0,"epistemico":"fato","origem":"wiki","rel":"relaciona"}],"confianca":1.0,"contraexemplos":[],"descricao":"SCOT: um mesmo artefato admite significados e usos radicalmente distintos conforme o grupo social, antes de um 'fechamento' fixar sua forma e sentido. Quem tem voz nessa disputa define o que o artefato será.","epistemico":"inferencia","exemplos":[],"id":"flexibilidade_interpretativa","memoria":"semantica","meta":{"autor":"Pinch & Bijker","dominio":"ciencias_de_fora","fonte":""},"origem":"wiki","palavras_chave":[],"sinonimos":[],"tipo":"conceito","titulo":"Flexibilidade Interpretativa"}
\section{Flexibilidade Interpretativa}
SCOT: um mesmo artefato admite significados e usos radicalmente distintos conforme o grupo social, antes de um 'fechamento' fixar sua forma e sentido. Quem tem voz nessa disputa define o que o artefato será.
\prov{relações: usa construcao\_social\_tecnologia; relaciona mediacao\_semiotica; relaciona injustica\_epistemica}
\prov{proveniência: wiki \textperiodcentered{} inferencia \textperiodcentered{} confiança 1.0 \textperiodcentered{} domínio ciencias\_de\_fora \textperiodcentered{} história 4 versões no jornal}

%CRISTAL {"arestas":[{"alvo":"orquestrador_multi-tecidos","confianca":1.0,"epistemico":"fato","origem":"paper","rel":"parte_de"},{"alvo":"consciencia_pura","confianca":1.0,"epistemico":"fato","origem":"wiki","rel":"relaciona"},{"alvo":"psicologia","confianca":1.0,"epistemico":"fato","origem":"wiki","rel":"parte_de"}],"confianca":1.0,"contraexemplos":[],"descricao":"Csikszentmihalyi: estado de absorção total em que o ego desaparece e desafio iguala competência — prefigura a dissolução do 'eu'.","epistemico":"inferencia","exemplos":[],"id":"fluxo","memoria":"semantica","meta":{"autor":"Csikszentmihalyi","dominio":"ciencias_de_fora","fonte":""},"origem":"wiki","palavras_chave":[],"sinonimos":[],"tipo":"conceito","titulo":"Fluxo (Experiência Ótima)"}
\section{Fluxo (Experiência Ótima)}
Csikszentmihalyi: estado de absorção total em que o ego desaparece e desafio iguala competência — prefigura a dissolução do 'eu'.
\prov{relações: parte\_de orquestrador\_multi-tecidos; relaciona consciencia\_pura; parte\_de psicologia}
\prov{proveniência: wiki \textperiodcentered{} inferencia \textperiodcentered{} confiança 1.0 \textperiodcentered{} domínio ciencias\_de\_fora \textperiodcentered{} história 5 versões no jornal}

%CRISTAL {"arestas":[{"alvo":"literatura","confianca":1.0,"epistemico":"fato","origem":"wiki","rel":"parte_de"},{"alvo":"realismo_literario","confianca":1.0,"epistemico":"fato","origem":"wiki","rel":"contradiz"}],"confianca":1.0,"contraexemplos":[],"descricao":"A técnica moderna (Joyce, Woolf) de transcrever o pensamento em bruto — a mente antes da ordem da gramática.","epistemico":"inferencia","exemplos":[],"id":"fluxo_de_consciencia","memoria":"semantica","meta":{"autor":"Joyce","dominio":"ciencias_de_fora","fonte":""},"origem":"wiki","palavras_chave":[],"sinonimos":[],"tipo":"conceito","titulo":"Fluxo de Consciência"}
\section{Fluxo de Consciência}
A técnica moderna (Joyce, Woolf) de transcrever o pensamento em bruto — a mente antes da ordem da gramática.
\prov{relações: parte\_de literatura; contradiz realismo\_literario}
\prov{proveniência: wiki \textperiodcentered{} inferencia \textperiodcentered{} confiança 1.0 \textperiodcentered{} domínio ciencias\_de\_fora \textperiodcentered{} história 3 versões no jornal}

%CRISTAL {"arestas":[{"alvo":"linguistica","confianca":1.0,"epistemico":"fato","origem":"wiki","rel":"parte_de"}],"confianca":1.0,"contraexemplos":[],"descricao":"O estudo dos sons como sistema — quais distinções de som carregam sentido numa língua.","epistemico":"fato","exemplos":[],"id":"fonologia","memoria":"semantica","meta":{"dominio":"ciencias_de_fora","fonte":""},"origem":"wiki","palavras_chave":[],"sinonimos":[],"tipo":"conceito","titulo":"Fonologia"}
\section{Fonologia}
O estudo dos sons como sistema — quais distinções de som carregam sentido numa língua.
\prov{relações: parte\_de linguistica}
\prov{proveniência: wiki \textperiodcentered{} fato \textperiodcentered{} confiança 1.0 \textperiodcentered{} domínio ciencias\_de\_fora \textperiodcentered{} história 2 versões no jornal}

%CRISTAL {"arestas":[{"alvo":"vitruvio_triade","confianca":1.0,"epistemico":"fato","origem":"wiki","rel":"usa"},{"alvo":"artefatos_tem_politica","confianca":1.0,"epistemico":"fato","origem":"wiki","rel":"relaciona"},{"alvo":"formalismo_greenberg","confianca":1.0,"epistemico":"fato","origem":"wiki","rel":"relaciona"}],"confianca":1.0,"contraexemplos":[],"descricao":"Sullivan: a aparência do edifício deve decorrer de seu propósito. Le Corbusier radicaliza no modernismo — a casa como 'máquina de morar', projetada com a racionalidade e a estética industrial de uma máquina.","epistemico":"inferencia","exemplos":[],"id":"forma_funcao_sullivan","memoria":"semantica","meta":{"autor":"Sullivan/Le Corbusier","dominio":"ciencias_de_fora","fonte":""},"origem":"wiki","palavras_chave":[],"sinonimos":[],"tipo":"conceito","titulo":"A Forma Segue a Função (Sullivan/Le Corbusier)"}
\section{A Forma Segue a Função (Sullivan/Le Corbusier)}
Sullivan: a aparência do edifício deve decorrer de seu propósito. Le Corbusier radicaliza no modernismo — a casa como 'máquina de morar', projetada com a racionalidade e a estética industrial de uma máquina.
\prov{relações: usa vitruvio\_triade; relaciona artefatos\_tem\_politica; relaciona formalismo\_greenberg}
\prov{proveniência: wiki \textperiodcentered{} inferencia \textperiodcentered{} confiança 1.0 \textperiodcentered{} domínio ciencias\_de\_fora \textperiodcentered{} história 4 versões no jornal}

%CRISTAL {"arestas":[{"alvo":"arquitetura","confianca":1.0,"epistemico":"fato","origem":"wiki","rel":"parte_de"},{"alvo":"forma_funcao_sullivan","confianca":1.0,"epistemico":"fato","origem":"wiki","rel":"relaciona"}],"confianca":1.0,"contraexemplos":[],"descricao":"Sullivan: a forma de um edifício deve nascer de seu uso — o lema fundador da arquitetura moderna.","epistemico":"inferencia","exemplos":[],"id":"forma_segue_funcao","memoria":"semantica","meta":{"autor":"Sullivan","dominio":"ciencias_de_fora","fonte":""},"origem":"wiki","palavras_chave":[],"sinonimos":[],"tipo":"conceito","titulo":"A Forma Segue a Função"}
\section{A Forma Segue a Função}
Sullivan: a forma de um edifício deve nascer de seu uso — o lema fundador da arquitetura moderna.
\prov{relações: parte\_de arquitetura; relaciona forma\_funcao\_sullivan}
\prov{proveniência: wiki \textperiodcentered{} inferencia \textperiodcentered{} confiança 1.0 \textperiodcentered{} domínio ciencias\_de\_fora \textperiodcentered{} história 4 versões no jornal}

%CRISTAL {"arestas":[{"alvo":"tonalidade","confianca":1.0,"epistemico":"fato","origem":"wiki","rel":"usa"},{"alvo":"teoria_musical","confianca":1.0,"epistemico":"fato","origem":"wiki","rel":"parte_de"}],"confianca":1.0,"contraexemplos":[],"descricao":"A grande arquitetura clássica: exposição (dois temas), desenvolvimento (conflito) e recapitulação (resolução).","epistemico":"fato","exemplos":[],"id":"forma_sonata","memoria":"semantica","meta":{"dominio":"ciencias_de_fora","fonte":""},"origem":"wiki","palavras_chave":[],"sinonimos":[],"tipo":"conceito","titulo":"Forma Sonata"}
\section{Forma Sonata}
A grande arquitetura clássica: exposição (dois temas), desenvolvimento (conflito) e recapitulação (resolução).
\prov{relações: usa tonalidade; parte\_de teoria\_musical}
\prov{proveniência: wiki \textperiodcentered{} fato \textperiodcentered{} confiança 1.0 \textperiodcentered{} domínio ciencias\_de\_fora \textperiodcentered{} história 3 versões no jornal}

%CRISTAL {"arestas":[{"alvo":"abstracao_kandinsky","confianca":1.0,"epistemico":"fato","origem":"wiki","rel":"usa"},{"alvo":"beleza_desinteressada_kant","confianca":1.0,"epistemico":"fato","origem":"wiki","rel":"usa"},{"alvo":"artes_visuais","confianca":1.0,"epistemico":"fato","origem":"wiki","rel":"parte_de"}],"confianca":1.0,"contraexemplos":[],"descricao":"Greenberg: a pintura modernista deve criticar-se com seus próprios meios, purgando o que não lhe é próprio; sua essência é a planeidade (bidimensionalidade) do suporte. A 'especificidade do meio' como destino do modernismo.","epistemico":"inferencia","exemplos":[],"id":"formalismo_greenberg","memoria":"semantica","meta":{"autor":"Greenberg","dominio":"ciencias_de_fora","fonte":""},"origem":"wiki","palavras_chave":[],"sinonimos":[],"tipo":"conceito","titulo":"Formalismo / Pureza do Meio (Greenberg)"}
\section{Formalismo / Pureza do Meio (Greenberg)}
Greenberg: a pintura modernista deve criticar-se com seus próprios meios, purgando o que não lhe é próprio; sua essência é a planeidade (bidimensionalidade) do suporte. A 'especificidade do meio' como destino do modernismo.
\prov{relações: usa abstracao\_kandinsky; usa beleza\_desinteressada\_kant; parte\_de artes\_visuais}
\prov{proveniência: wiki \textperiodcentered{} inferencia \textperiodcentered{} confiança 1.0 \textperiodcentered{} domínio ciencias\_de\_fora \textperiodcentered{} história 4 versões no jornal}

%CRISTAL {"arestas":[{"alvo":"filosofia_da_matematica","confianca":1.0,"epistemico":"fato","origem":"wiki","rel":"parte_de"},{"alvo":"numero_indefinido","confianca":1.0,"epistemico":"fato","origem":"wiki","rel":"relaciona"}],"confianca":1.0,"contraexemplos":[],"descricao":"A matemática é um jogo de símbolos governado por regras finitas, sem referente metafísico (Hilbert). O Lopes diz exatamente isto: 'matemática é jogo de símbolos'.","epistemico":"inferencia","exemplos":[],"id":"formalismo_hilbert","memoria":"semantica","meta":{"autor":"Hilbert","dominio":"filosofia","fonte":""},"origem":"wiki","palavras_chave":[],"sinonimos":[],"tipo":"conceito","titulo":"Formalismo"}
\section{Formalismo}
A matemática é um jogo de símbolos governado por regras finitas, sem referente metafísico (Hilbert). O Lopes diz exatamente isto: 'matemática é jogo de símbolos'.
\prov{relações: parte\_de filosofia\_da\_matematica; relaciona numero\_indefinido}
\prov{proveniência: wiki \textperiodcentered{} inferencia \textperiodcentered{} confiança 1.0 \textperiodcentered{} domínio filosofia \textperiodcentered{} história 4 versões no jornal}

%CRISTAL {"arestas":[{"alvo":"artes_visuais","confianca":1.0,"epistemico":"fato","origem":"wiki","rel":"parte_de"},{"alvo":"aura_benjamin","confianca":1.0,"epistemico":"fato","origem":"wiki","rel":"explica"},{"alvo":"bazin_ontologia","confianca":1.0,"epistemico":"fato","origem":"wiki","rel":"relaciona"},{"alvo":"barthes_punctum","confianca":1.0,"epistemico":"fato","origem":"wiki","rel":"relaciona"}],"confianca":1.0,"contraexemplos":[],"descricao":"A imagem fixada pela luz — mudou a arte (libertou a pintura do retrato) e a percepção do real; e a aura da obra.","epistemico":"fato","exemplos":[],"id":"fotografia","memoria":"semantica","meta":{"dominio":"ciencias_de_fora","fonte":""},"origem":"wiki","palavras_chave":[],"sinonimos":[],"tipo":"conceito","titulo":"Fotografia"}
\section{Fotografia}
A imagem fixada pela luz — mudou a arte (libertou a pintura do retrato) e a percepção do real; e a aura da obra.
\prov{relações: parte\_de artes\_visuais; explica aura\_benjamin; relaciona bazin\_ontologia; relaciona barthes\_punctum}
\prov{proveniência: wiki \textperiodcentered{} fato \textperiodcentered{} confiança 1.0 \textperiodcentered{} domínio ciencias\_de\_fora \textperiodcentered{} história 7 versões no jornal}

%CRISTAL {"arestas":[{"alvo":"metabolismo","confianca":1.0,"epistemico":"fato","origem":"wiki","rel":"especializa"},{"alvo":"bioquimica","confianca":1.0,"epistemico":"fato","origem":"wiki","rel":"parte_de"}],"confianca":1.0,"contraexemplos":[],"descricao":"Plantas e algas capturam luz para fixar carbono e liberar oxigênio — a porta de energia de quase toda a biosfera.","epistemico":"fato","exemplos":[],"id":"fotossintese","memoria":"semantica","meta":{"dominio":"ciencias_de_fora","fonte":""},"origem":"wiki","palavras_chave":[],"sinonimos":[],"tipo":"conceito","titulo":"Fotossíntese"}
\section{Fotossíntese}
Plantas e algas capturam luz para fixar carbono e liberar oxigênio — a porta de energia de quase toda a biosfera.
\prov{relações: especializa metabolismo; parte\_de bioquimica}
\prov{proveniência: wiki \textperiodcentered{} fato \textperiodcentered{} confiança 1.0 \textperiodcentered{} domínio ciencias\_de\_fora \textperiodcentered{} história 3 versões no jornal}

%CRISTAL {"arestas":[{"alvo":"teoria_dos_numeros","confianca":1.0,"epistemico":"fato","origem":"wiki","rel":"parte_de"},{"alvo":"equacao_de_pell","confianca":1.0,"epistemico":"fato","origem":"wiki","rel":"usa"},{"alvo":"transformada_metalica","confianca":1.0,"epistemico":"fato","origem":"wiki","rel":"explica"},{"alvo":"fracoes_continuas_e_reta_de_ouro","confianca":1.0,"epistemico":"fato","origem":"wiki","rel":"relaciona"}],"confianca":1.0,"contraexemplos":[],"descricao":"Representar um número por [a₀;a₁,a₂,…] — a melhor aproximação racional. Os metálicos σ_n têm fração contínua periódica [n;n,n,…].","epistemico":"fato","exemplos":[],"id":"fracoes_continuas","memoria":"semantica","meta":{"dominio":"ciencias_de_fora","fonte":""},"origem":"wiki","palavras_chave":[],"sinonimos":[],"tipo":"conceito","titulo":"Frações Contínuas"}
\section{Frações Contínuas}
Representar um número por [a\ensuremath{_{0}};a\ensuremath{_{1}},a\ensuremath{_{2}},…] — a melhor aproximação racional. Os metálicos \ensuremath{\sigma}\_n têm fração contínua periódica [n;n,n,…].
\prov{relações: parte\_de teoria\_dos\_numeros; usa equacao\_de\_pell; explica transformada\_metalica; relaciona fracoes\_continuas\_e\_reta\_de\_ouro}
\prov{proveniência: wiki \textperiodcentered{} fato \textperiodcentered{} confiança 1.0 \textperiodcentered{} domínio ciencias\_de\_fora \textperiodcentered{} história 9 versões no jornal}

%CRISTAL {"arestas":[{"alvo":"arte","confianca":1.0,"epistemico":"fato","origem":"wiki","rel":"parte_de"},{"alvo":"geometria_fractal","confianca":1.0,"epistemico":"fato","origem":"wiki","rel":"usa"},{"alvo":"dimensao_de_hausdorff","confianca":1.0,"epistemico":"fato","origem":"wiki","rel":"usa"}],"confianca":1.0,"contraexemplos":[],"descricao":"Da onda de Hokusai aos respingos de Pollock (dimensão de Hausdorff ~1,7): a arte capta a rugosidade auto-similar da natureza.","epistemico":"inferencia","exemplos":[],"id":"fractais_na_arte","memoria":"semantica","meta":{"dominio":"ciencias_de_fora","fonte":""},"origem":"wiki","palavras_chave":[],"sinonimos":[],"tipo":"conceito","titulo":"Fractais na Arte"}
\section{Fractais na Arte}
Da onda de Hokusai aos respingos de Pollock (dimensão de Hausdorff \~{}1,7): a arte capta a rugosidade auto-similar da natureza.
\prov{relações: parte\_de arte; usa geometria\_fractal; usa dimensao\_de\_hausdorff}
\prov{proveniência: wiki \textperiodcentered{} inferencia \textperiodcentered{} confiança 1.0 \textperiodcentered{} domínio ciencias\_de\_fora \textperiodcentered{} história 4 versões no jornal}

%CRISTAL {"arestas":[{"alvo":"teoria_da_comunicacao","confianca":1.0,"epistemico":"fato","origem":"wiki","rel":"parte_de"},{"alvo":"agenda_setting","confianca":1.0,"epistemico":"fato","origem":"wiki","rel":"especializa"},{"alvo":"gatekeeping","confianca":1.0,"epistemico":"fato","origem":"wiki","rel":"relaciona"},{"alvo":"dois_sistemas_kahneman","confianca":1.0,"epistemico":"fato","origem":"wiki","rel":"relaciona"},{"alvo":"mediacao_semiotica","confianca":1.0,"epistemico":"fato","origem":"wiki","rel":"relaciona"}],"confianca":1.0,"contraexemplos":[],"descricao":"Goffman & Entman: enquadrar é selecionar aspectos de uma realidade e torná-los salientes, promovendo uma definição do problema e uma avaliação moral — a moldura de apresentação molda a interpretação.","epistemico":"inferencia","exemplos":[],"id":"framing_enquadramento","memoria":"semantica","meta":{"autor":"Goffman & Entman","dominio":"ciencias_de_fora","fonte":""},"origem":"wiki","palavras_chave":[],"sinonimos":[],"tipo":"conceito","titulo":"Framing / Enquadramento"}
\section{Framing / Enquadramento}
Goffman \& Entman: enquadrar é selecionar aspectos de uma realidade e torná-los salientes, promovendo uma definição do problema e uma avaliação moral — a moldura de apresentação molda a interpretação.
\prov{relações: parte\_de teoria\_da\_comunicacao; especializa agenda\_setting; relaciona gatekeeping; relaciona dois\_sistemas\_kahneman; relaciona mediacao\_semiotica}
\prov{proveniência: wiki \textperiodcentered{} inferencia \textperiodcentered{} confiança 1.0 \textperiodcentered{} domínio ciencias\_de\_fora \textperiodcentered{} história 6 versões no jornal}

%CRISTAL {"arestas":[{"alvo":"filosofia_da_linguagem","confianca":1.0,"epistemico":"fato","origem":"wiki","rel":"parte_de"},{"alvo":"signo_saussure","confianca":1.0,"epistemico":"fato","origem":"wiki","rel":"relaciona"}],"confianca":1.0,"contraexemplos":[],"descricao":"Frege: uma expressão tem sentido (modo de apresentação) e referência (o objeto). 'Estrela matutina' e 'vespertina' diferem em sentido mas referem Vênus.","epistemico":"inferencia","exemplos":[],"id":"frege_sentido_referencia","memoria":"semantica","meta":{"autor":"Frege","dominio":"ciencias_de_fora","fonte":""},"origem":"wiki","palavras_chave":[],"sinonimos":[],"tipo":"conceito","titulo":"Sentido vs Referência (Frege)"}
\section{Sentido vs Referência (Frege)}
Frege: uma expressão tem sentido (modo de apresentação) e referência (o objeto). 'Estrela matutina' e 'vespertina' diferem em sentido mas referem Vênus.
\prov{relações: parte\_de filosofia\_da\_linguagem; relaciona signo\_saussure}
\prov{proveniência: wiki \textperiodcentered{} inferencia \textperiodcentered{} confiança 1.0 \textperiodcentered{} domínio ciencias\_de\_fora \textperiodcentered{} história 3 versões no jornal}

%CRISTAL {"arestas":[{"alvo":"didatica_matematica","confianca":1.0,"epistemico":"fato","origem":"wiki","rel":"parte_de"}],"confianca":1.0,"contraexemplos":[],"descricao":"Freudenthal: a matemática é atividade humana; aprende-se matematizando situações — matematização horizontal (do contexto ao símbolo) e vertical (dentro do sistema), por reinvenção guiada.","epistemico":"inferencia","exemplos":[],"id":"freudenthal_rme","memoria":"semantica","meta":{"autor":"Freudenthal","dominio":"ciencias_de_fora","fonte":""},"origem":"wiki","palavras_chave":[],"sinonimos":[],"tipo":"conceito","titulo":"Educação Matemática Realística (RME)"}
\section{Educação Matemática Realística (RME)}
Freudenthal: a matemática é atividade humana; aprende-se matematizando situações — matematização horizontal (do contexto ao símbolo) e vertical (dentro do sistema), por reinvenção guiada.
\prov{relações: parte\_de didatica\_matematica}
\prov{proveniência: wiki \textperiodcentered{} inferencia \textperiodcentered{} confiança 1.0 \textperiodcentered{} domínio ciencias\_de\_fora \textperiodcentered{} história 2 versões no jornal}

%CRISTAL {"arestas":[{"alvo":"filosofia_do_direito","confianca":1.0,"epistemico":"fato","origem":"wiki","rel":"parte_de"},{"alvo":"juspositivismo","confianca":1.0,"epistemico":"fato","origem":"wiki","rel":"contradiz"},{"alvo":"jusnaturalismo","confianca":1.0,"epistemico":"fato","origem":"wiki","rel":"relaciona"}],"confianca":1.0,"contraexemplos":[],"descricao":"Fuller: para ser direito, um sistema deve satisfazer oito exigências procedimentais (generalidade, publicidade, não-retroatividade, clareza, não-contradição, exequibilidade, estabilidade, congruência) — uma 'moralidade interna' que conecta direito e moral. No debate Hart-Fuller, o regime nazista falhou em ser verdadeiro direito.","epistemico":"inferencia","exemplos":[],"id":"fuller_moralidade_interna","memoria":"semantica","meta":{"autor":"Fuller","dominio":"ciencias_de_fora","fonte":""},"origem":"wiki","palavras_chave":[],"sinonimos":[],"tipo":"conceito","titulo":"Moralidade Interna do Direito (Fuller)"}
\section{Moralidade Interna do Direito (Fuller)}
Fuller: para ser direito, um sistema deve satisfazer oito exigências procedimentais (generalidade, publicidade, não-retroatividade, clareza, não-contradição, exequibilidade, estabilidade, congruência) — uma 'moralidade interna' que conecta direito e moral. No debate Hart-Fuller, o regime nazista falhou em ser verdadeiro direito.
\prov{relações: parte\_de filosofia\_do\_direito; contradiz juspositivismo; relaciona jusnaturalismo}
\prov{proveniência: wiki \textperiodcentered{} inferencia \textperiodcentered{} confiança 1.0 \textperiodcentered{} domínio ciencias\_de\_fora \textperiodcentered{} história 4 versões no jornal}

%CRISTAL {"arestas":[{"alvo":"morte_do_autor","confianca":1.0,"epistemico":"fato","origem":"wiki","rel":"relaciona"},{"alvo":"literatura","confianca":1.0,"epistemico":"fato","origem":"wiki","rel":"parte_de"}],"confianca":1.0,"contraexemplos":[],"descricao":"Foucault, em resposta a Barthes: o 'autor' não é a pessoa, mas uma função discursiva — um princípio que classifica, agrupa e limita a circulação e a atribuição dos discursos numa cultura.","epistemico":"inferencia","exemplos":[],"id":"funcao_autor_foucault","memoria":"semantica","meta":{"autor":"Foucault","dominio":"ciencias_de_fora","fonte":""},"origem":"wiki","palavras_chave":[],"sinonimos":[],"tipo":"conceito","titulo":"Função-Autor (Foucault)"}
\section{Função-Autor (Foucault)}
Foucault, em resposta a Barthes: o 'autor' não é a pessoa, mas uma função discursiva — um princípio que classifica, agrupa e limita a circulação e a atribuição dos discursos numa cultura.
\prov{relações: relaciona morte\_do\_autor; parte\_de literatura}
\prov{proveniência: wiki \textperiodcentered{} inferencia \textperiodcentered{} confiança 1.0 \textperiodcentered{} domínio ciencias\_de\_fora \textperiodcentered{} história 3 versões no jornal}

%CRISTAL {"arestas":[{"alvo":"funcoes_da_linguagem","confianca":1.0,"epistemico":"fato","origem":"wiki","rel":"parte_de"},{"alvo":"literariedade","confianca":1.0,"epistemico":"fato","origem":"wiki","rel":"relaciona"}],"confianca":1.0,"contraexemplos":[],"descricao":"Jakobson: quando a mensagem se volta para a própria forma — o eixo da seleção projetado sobre o da combinação.","epistemico":"fato","exemplos":[],"id":"funcao_poetica","memoria":"semantica","meta":{"autor":"Jakobson","dominio":"ciencias_de_fora","fonte":""},"origem":"wiki","palavras_chave":[],"sinonimos":[],"tipo":"conceito","titulo":"Função Poética"}
\section{Função Poética}
Jakobson: quando a mensagem se volta para a própria forma — o eixo da seleção projetado sobre o da combinação.
\prov{relações: parte\_de funcoes\_da\_linguagem; relaciona literariedade}
\prov{proveniência: wiki \textperiodcentered{} fato \textperiodcentered{} confiança 1.0 \textperiodcentered{} domínio ciencias\_de\_fora \textperiodcentered{} história 3 versões no jornal}

%CRISTAL {"arestas":[{"alvo":"filosofia_da_computacao","confianca":1.0,"epistemico":"fato","origem":"wiki","rel":"parte_de"},{"alvo":"mente_computacao","confianca":1.0,"epistemico":"fato","origem":"wiki","rel":"usa"},{"alvo":"consciencia","confianca":1.0,"epistemico":"fato","origem":"wiki","rel":"relaciona"},{"alvo":"teste_de_turing","confianca":1.0,"epistemico":"fato","origem":"wiki","rel":"relaciona"}],"confianca":1.0,"contraexemplos":[],"descricao":"Putnam & Fodor: estados mentais se definem por seu papel funcional (relações causais entre entradas, estados e saídas), não pelo substrato físico — a mente está para o cérebro como o software para o hardware (realizabilidade múltipla).","epistemico":"inferencia","exemplos":[],"id":"funcionalismo_mente","memoria":"semantica","meta":{"autor":"Putnam & Fodor","dominio":"ciencias_de_fora","fonte":""},"origem":"wiki","palavras_chave":[],"sinonimos":[],"tipo":"conceito","titulo":"Funcionalismo / Mente como Software"}
\section{Funcionalismo / Mente como Software}
Putnam \& Fodor: estados mentais se definem por seu papel funcional (relações causais entre entradas, estados e saídas), não pelo substrato físico — a mente está para o cérebro como o software para o hardware (realizabilidade múltipla).
\prov{relações: parte\_de filosofia\_da\_computacao; usa mente\_computacao; relaciona consciencia; relaciona teste\_de\_turing}
\prov{proveniência: wiki \textperiodcentered{} inferencia \textperiodcentered{} confiança 1.0 \textperiodcentered{} domínio ciencias\_de\_fora \textperiodcentered{} história 5 versões no jornal}

%CRISTAL {"arestas":[{"alvo":"roman_jakobson","confianca":1.0,"epistemico":"fato","origem":"wiki","rel":"parte_de"},{"alvo":"semiotica","confianca":1.0,"epistemico":"fato","origem":"wiki","rel":"usa"},{"alvo":"funcoes_linguagem_jakobson","confianca":1.0,"epistemico":"fato","origem":"wiki","rel":"relaciona"}],"confianca":1.0,"contraexemplos":[],"descricao":"Jakobson: seis funções segundo o foco — referencial, emotiva, conativa, fática, metalinguística e poética.","epistemico":"fato","exemplos":[],"id":"funcoes_da_linguagem","memoria":"semantica","meta":{"autor":"Jakobson","dominio":"ciencias_de_fora","fonte":""},"origem":"wiki","palavras_chave":[],"sinonimos":[],"tipo":"conceito","titulo":"As Funções da Linguagem"}
\section{As Funções da Linguagem}
Jakobson: seis funções segundo o foco — referencial, emotiva, conativa, fática, metalinguística e poética.
\prov{relações: parte\_de roman\_jakobson; usa semiotica; relaciona funcoes\_linguagem\_jakobson}
\prov{proveniência: wiki \textperiodcentered{} fato \textperiodcentered{} confiança 1.0 \textperiodcentered{} domínio ciencias\_de\_fora \textperiodcentered{} história 4 versões no jornal}

%CRISTAL {"arestas":[{"alvo":"teoria_da_comunicacao","confianca":1.0,"epistemico":"fato","origem":"wiki","rel":"parte_de"},{"alvo":"modelo_shannon_weaver","confianca":1.0,"epistemico":"fato","origem":"wiki","rel":"usa"},{"alvo":"linguagem","confianca":1.0,"epistemico":"fato","origem":"wiki","rel":"referencia"}],"confianca":1.0,"contraexemplos":[],"descricao":"Jakobson: a cada fator do circuito comunicativo (remetente, destinatário, contexto, mensagem, contato, código) corresponde uma função dominante — emotiva, conativa, referencial, poética, fática e metalinguística.","epistemico":"inferencia","exemplos":[],"id":"funcoes_linguagem_jakobson","memoria":"semantica","meta":{"autor":"Jakobson","dominio":"ciencias_de_fora","fonte":""},"origem":"wiki","palavras_chave":[],"sinonimos":[],"tipo":"conceito","titulo":"Seis Funções da Linguagem (Jakobson)"}
\section{Seis Funções da Linguagem (Jakobson)}
Jakobson: a cada fator do circuito comunicativo (remetente, destinatário, contexto, mensagem, contato, código) corresponde uma função dominante — emotiva, conativa, referencial, poética, fática e metalinguística.
\prov{relações: parte\_de teoria\_da\_comunicacao; usa modelo\_shannon\_weaver; referencia linguagem}
\prov{proveniência: wiki \textperiodcentered{} inferencia \textperiodcentered{} confiança 1.0 \textperiodcentered{} domínio ciencias\_de\_fora \textperiodcentered{} história 4 versões no jornal}

%CRISTAL {"arestas":[{"alvo":"autopoiese","confianca":1.0,"epistemico":"fato","origem":"wiki","rel":"relaciona"},{"alvo":"sistema_terra","confianca":1.0,"epistemico":"fato","origem":"wiki","rel":"parte_de"},{"alvo":"word","confianca":0.5,"epistemico":"fato","origem":"wiki","rel":"relaciona"},{"alvo":"syscall","confianca":0.5,"epistemico":"fato","origem":"wiki","rel":"relaciona"},{"alvo":"vida-morte","confianca":0.5,"epistemico":"fato","origem":"wiki","rel":"relaciona"},{"alvo":"vita_contemplativa","confianca":0.5,"epistemico":"fato","origem":"wiki","rel":"relaciona"}],"confianca":1.0,"contraexemplos":[],"descricao":"Lovelock/Margulis: a biosfera e o meio físico formam um sistema autorregulador que mantém a homeostase do planeta.","epistemico":"hipotese","exemplos":[],"id":"gaia","memoria":"semantica","meta":{"autor":"Lovelock","dominio":"ciencias_de_fora","fonte":""},"origem":"wiki","palavras_chave":[],"sinonimos":[],"tipo":"conceito","titulo":"Hipótese de Gaia"}
\section{Hipótese de Gaia}
Lovelock/Margulis: a biosfera e o meio físico formam um sistema autorregulador que mantém a homeostase do planeta.
\prov{relações: relaciona autopoiese; parte\_de sistema\_terra; relaciona word; relaciona syscall; relaciona vida-morte; relaciona vita\_contemplativa}
\prov{proveniência: wiki \textperiodcentered{} hipotese \textperiodcentered{} confiança 1.0 \textperiodcentered{} domínio ciencias\_de\_fora \textperiodcentered{} história 7 versões no jornal}

%CRISTAL {"arestas":[{"alvo":"teoria_da_comunicacao","confianca":1.0,"epistemico":"fato","origem":"wiki","rel":"parte_de"},{"alvo":"biopoder","confianca":1.0,"epistemico":"fato","origem":"wiki","rel":"relaciona"}],"confianca":1.0,"contraexemplos":[],"descricao":"White (a partir de Lewin): há 'porteiros' que decidem o que passa e o que é barrado no fluxo informativo; os julgamentos, valores e vieses do editor determinam quais fatos viram notícia.","epistemico":"fato","exemplos":[],"id":"gatekeeping","memoria":"semantica","meta":{"autor":"White","dominio":"ciencias_de_fora","fonte":""},"origem":"wiki","palavras_chave":[],"sinonimos":[],"tipo":"conceito","titulo":"Gatekeeping (Seleção da Notícia)"}
\section{Gatekeeping (Seleção da Notícia)}
White (a partir de Lewin): há 'porteiros' que decidem o que passa e o que é barrado no fluxo informativo; os julgamentos, valores e vieses do editor determinam quais fatos viram notícia.
\prov{relações: parte\_de teoria\_da\_comunicacao; relaciona biopoder}
\prov{proveniência: wiki \textperiodcentered{} fato \textperiodcentered{} confiança 1.0 \textperiodcentered{} domínio ciencias\_de\_fora \textperiodcentered{} história 3 versões no jornal}

%CRISTAL {"arestas":[{"alvo":"teoria_dos_numeros","confianca":1.0,"epistemico":"fato","origem":"wiki","rel":"parte_de"},{"alvo":"numeros_complexos","confianca":1.0,"epistemico":"fato","origem":"wiki","rel":"relaciona"}],"confianca":1.0,"contraexemplos":[],"descricao":"Gauss: da reciprocidade quadrática à geometria diferencial e à estatística — o rigor moderno em pessoa.","epistemico":"fato","exemplos":[],"id":"gauss_principe","memoria":"semantica","meta":{"autor":"Gauss","dominio":"ciencias_de_fora","fonte":""},"origem":"wiki","palavras_chave":[],"sinonimos":[],"tipo":"conceito","titulo":"Gauss, o Príncipe"}
\section{Gauss, o Príncipe}
Gauss: da reciprocidade quadrática à geometria diferencial e à estatística — o rigor moderno em pessoa.
\prov{relações: parte\_de teoria\_dos\_numeros; relaciona numeros\_complexos}
\prov{proveniência: wiki \textperiodcentered{} fato \textperiodcentered{} confiança 1.0 \textperiodcentered{} domínio ciencias\_de\_fora \textperiodcentered{} história 6 versões no jornal}

%CRISTAL {"arestas":[{"alvo":"genetica","confianca":1.0,"epistemico":"fato","origem":"wiki","rel":"parte_de"}],"confianca":1.0,"contraexemplos":[],"descricao":"A unidade de informação hereditária — um trecho de DNA que especifica um produto (proteína ou RNA) ou regula outros.","epistemico":"fato","exemplos":[],"id":"gene","memoria":"semantica","meta":{"dominio":"ciencias_de_fora","fonte":""},"origem":"wiki","palavras_chave":[],"sinonimos":[],"tipo":"conceito","titulo":"Gene"}
\section{Gene}
A unidade de informação hereditária — um trecho de DNA que especifica um produto (proteína ou RNA) ou regula outros.
\prov{relações: parte\_de genetica}
\prov{proveniência: wiki \textperiodcentered{} fato \textperiodcentered{} confiança 1.0 \textperiodcentered{} domínio ciencias\_de\_fora \textperiodcentered{} história 2 versões no jornal}

%CRISTAL {"arestas":[{"alvo":"codigo_genetico","confianca":1.0,"epistemico":"fato","origem":"wiki","rel":"usa"},{"alvo":"word","confianca":0.5,"epistemico":"fato","origem":"wiki","rel":"relaciona"},{"alvo":"selecao_natural","confianca":1.0,"epistemico":"fato","origem":"wiki","rel":"especializa"},{"alvo":"gene","confianca":1.0,"epistemico":"fato","origem":"wiki","rel":"usa"}],"confianca":1.0,"contraexemplos":[],"descricao":"Dawkins: o gene, não o indivíduo, é a unidade da seleção — o organismo é o veículo que o gene constrói.","epistemico":"inferencia","exemplos":[],"id":"gene_egoista","memoria":"semantica","meta":{"autor":"Dawkins","dominio":"ciencias_de_fora","fonte":""},"origem":"wiki","palavras_chave":[],"sinonimos":[],"tipo":"conceito","titulo":"O Gene Egoísta"}
\section{O Gene Egoísta}
Dawkins: o gene, não o indivíduo, é a unidade da seleção — o organismo é o veículo que o gene constrói.
\prov{relações: usa codigo\_genetico; relaciona word; especializa selecao\_natural; usa gene}
\prov{proveniência: wiki \textperiodcentered{} inferencia \textperiodcentered{} confiança 1.0 \textperiodcentered{} domínio ciencias\_de\_fora \textperiodcentered{} história 7 versões no jornal}

%CRISTAL {"arestas":[{"alvo":"biopoder","confianca":1.0,"epistemico":"fato","origem":"wiki","rel":"usa"},{"alvo":"contingencia_historica","confianca":1.0,"epistemico":"fato","origem":"wiki","rel":"usa"},{"alvo":"word","confianca":0.5,"epistemico":"fato","origem":"wiki","rel":"relaciona"}],"confianca":1.0,"contraexemplos":[],"descricao":"Foucault: buscar descontinuidades e silêncios, não origens; o saber é poder, a verdade é produzida. A história é política do presente.","epistemico":"inferencia","exemplos":[],"id":"genealogia_foucault","memoria":"semantica","meta":{"autor":"Foucault","dominio":"ciencias_de_fora","fonte":""},"origem":"wiki","palavras_chave":[],"sinonimos":[],"tipo":"conceito","titulo":"Genealogia"}
\section{Genealogia}
Foucault: buscar descontinuidades e silêncios, não origens; o saber é poder, a verdade é produzida. A história é política do presente.
\prov{relações: usa biopoder; usa contingencia\_historica; relaciona word}
\prov{proveniência: wiki \textperiodcentered{} inferencia \textperiodcentered{} confiança 1.0 \textperiodcentered{} domínio ciencias\_de\_fora \textperiodcentered{} história 4 versões no jornal}

%CRISTAL {"arestas":[{"alvo":"concordancia","confianca":1.0,"epistemico":"fato","origem":"wiki","rel":"usa"}],"confianca":1.0,"contraexemplos":[],"descricao":"Classificar substantivos em classes (masculino/feminino/neutro) que forçam concordância — arbitrário (o alemão tem 3; o inglês, nenhum; o suaíli, muitos).","epistemico":"fato","exemplos":[],"id":"genero_gramatical","memoria":"semantica","meta":{"dominio":"ciencias_de_fora","fonte":""},"origem":"wiki","palavras_chave":[],"sinonimos":[],"tipo":"conceito","titulo":"Gênero Gramatical"}
\section{Gênero Gramatical}
Classificar substantivos em classes (masculino/feminino/neutro) que forçam concordância — arbitrário (o alemão tem 3; o inglês, nenhum; o suaíli, muitos).
\prov{relações: usa concordancia}
\prov{proveniência: wiki \textperiodcentered{} fato \textperiodcentered{} confiança 1.0 \textperiodcentered{} domínio ciencias\_de\_fora \textperiodcentered{} história 2 versões no jornal}

%CRISTAL {"arestas":[{"alvo":"teoria_literaria","confianca":1.0,"epistemico":"fato","origem":"wiki","rel":"parte_de"}],"confianca":1.0,"contraexemplos":[],"descricao":"As grandes formas — épico (narrar), lírico (cantar o eu) e dramático (encenar) — e suas fronteiras móveis.","epistemico":"fato","exemplos":[],"id":"generos_literarios","memoria":"semantica","meta":{"dominio":"ciencias_de_fora","fonte":""},"origem":"wiki","palavras_chave":[],"sinonimos":[],"tipo":"conceito","titulo":"Gêneros Literários"}
\section{Gêneros Literários}
As grandes formas — épico (narrar), lírico (cantar o eu) e dramático (encenar) — e suas fronteiras móveis.
\prov{relações: parte\_de teoria\_literaria}
\prov{proveniência: wiki \textperiodcentered{} fato \textperiodcentered{} confiança 1.0 \textperiodcentered{} domínio ciencias\_de\_fora \textperiodcentered{} história 2 versões no jornal}

%CRISTAL {"arestas":[{"alvo":"biologia","confianca":1.0,"epistemico":"fato","origem":"wiki","rel":"parte_de"}],"confianca":1.0,"contraexemplos":[],"descricao":"A ciência da hereditariedade e da variação — como a informação da vida passa e muda entre gerações.","epistemico":"fato","exemplos":[],"id":"genetica","memoria":"semantica","meta":{"dominio":"ciencias_de_fora","fonte":""},"origem":"wiki","palavras_chave":[],"sinonimos":[],"tipo":"conceito","titulo":"Genética"}
\section{Genética}
A ciência da hereditariedade e da variação — como a informação da vida passa e muda entre gerações.
\prov{relações: parte\_de biologia}
\prov{proveniência: wiki \textperiodcentered{} fato \textperiodcentered{} confiança 1.0 \textperiodcentered{} domínio ciencias\_de\_fora \textperiodcentered{} história 2 versões no jornal}

%CRISTAL {"arestas":[{"alvo":"didatica_matematica","confianca":1.0,"epistemico":"fato","origem":"wiki","rel":"parte_de"},{"alvo":"honestidade_intelectual","confianca":1.0,"epistemico":"fato","origem":"wiki","rel":"relaciona"},{"alvo":"mestre_ignorante_ranciere","confianca":1.0,"epistemico":"fato","origem":"wiki","rel":"relaciona"},{"alvo":"idiotismo_nao_examinar","confianca":1.0,"epistemico":"fato","origem":"wiki","rel":"relaciona"}],"confianca":1.0,"contraexemplos":[],"descricao":"Postura anti-autoridade do Lopes no ensino: a autoridade não garante a verdade. Ele exibe erros de lógica de figuras canônicas (Elon Lages Lima, o próprio Peano na curva de Peano) e convida o estudante a não confiar cegamente em nomes consagrados.","epistemico":"inferencia","exemplos":[],"id":"genios_tambem_erram","memoria":"semantica","meta":{"autor":"Gentil Lopes","dominio":"ciencias_de_fora","fonte":""},"origem":"wiki","palavras_chave":[],"sinonimos":[],"tipo":"conceito","titulo":"Até Gênios Erram (Lopes)"}
\section{Até Gênios Erram (Lopes)}
Postura anti-autoridade do Lopes no ensino: a autoridade não garante a verdade. Ele exibe erros de lógica de figuras canônicas (Elon Lages Lima, o próprio Peano na curva de Peano) e convida o estudante a não confiar cegamente em nomes consagrados.
\prov{relações: parte\_de didatica\_matematica; relaciona honestidade\_intelectual; relaciona mestre\_ignorante\_ranciere; relaciona idiotismo\_nao\_examinar}
\prov{proveniência: wiki \textperiodcentered{} inferencia \textperiodcentered{} confiança 1.0 \textperiodcentered{} domínio ciencias\_de\_fora \textperiodcentered{} história 6 versões no jornal}

%CRISTAL {"arestas":[{"alvo":"gene","confianca":1.0,"epistemico":"fato","origem":"wiki","rel":"usa"}],"confianca":1.0,"contraexemplos":[],"descricao":"O conjunto completo do DNA de um organismo — o texto inteiro, não só os genes que codificam.","epistemico":"fato","exemplos":[],"id":"genoma","memoria":"semantica","meta":{"dominio":"ciencias_de_fora","fonte":""},"origem":"wiki","palavras_chave":[],"sinonimos":[],"tipo":"conceito","titulo":"Genoma"}
\section{Genoma}
O conjunto completo do DNA de um organismo — o texto inteiro, não só os genes que codificam.
\prov{relações: usa gene}
\prov{proveniência: wiki \textperiodcentered{} fato \textperiodcentered{} confiança 1.0 \textperiodcentered{} domínio ciencias\_de\_fora \textperiodcentered{} história 2 versões no jornal}

%CRISTAL {"arestas":[{"alvo":"sociologia","confianca":1.0,"epistemico":"fato","origem":"wiki","rel":"relaciona"},{"alvo":"word","confianca":0.5,"epistemico":"fato","origem":"wiki","rel":"relaciona"},{"alvo":"virtualizacao","confianca":0.5,"epistemico":"fato","origem":"wiki","rel":"relaciona"},{"alvo":"vida-morte","confianca":0.5,"epistemico":"fato","origem":"wiki","rel":"relaciona"},{"alvo":"vita_contemplativa","confianca":0.5,"epistemico":"fato","origem":"wiki","rel":"relaciona"},{"alvo":"logica","confianca":0.5,"epistemico":"fato","origem":"wiki","rel":"relaciona"}],"confianca":1.0,"contraexemplos":[],"descricao":"A ciência do espaço terrestre — a superfície da Terra como sistema físico e como produto das sociedades que a habitam, organizam e disputam.","epistemico":"fato","exemplos":[],"id":"geografia","memoria":"semantica","meta":{"dominio":"ciencias_de_fora","fonte":""},"origem":"wiki","palavras_chave":[],"sinonimos":[],"tipo":"conceito","titulo":"Geografia"}
\section{Geografia}
A ciência do espaço terrestre — a superfície da Terra como sistema físico e como produto das sociedades que a habitam, organizam e disputam.
\prov{relações: relaciona sociologia; relaciona word; relaciona virtualizacao; relaciona vida-morte; relaciona vita\_contemplativa; relaciona logica}
\prov{proveniência: wiki \textperiodcentered{} fato \textperiodcentered{} confiança 1.0 \textperiodcentered{} domínio ciencias\_de\_fora \textperiodcentered{} história 7 versões no jornal}

%CRISTAL {"arestas":[{"alvo":"geografia","confianca":1.0,"epistemico":"fato","origem":"wiki","rel":"parte_de"},{"alvo":"vida","confianca":1.0,"epistemico":"fato","origem":"wiki","rel":"relaciona"},{"alvo":"word","confianca":0.5,"epistemico":"fato","origem":"wiki","rel":"relaciona"}],"confianca":1.0,"contraexemplos":[],"descricao":"O estudo da Terra como sistema integrado — atmosfera, hidrosfera, litosfera e biosfera acopladas por fluxos de matéria e energia.","epistemico":"fato","exemplos":[],"id":"geografia_fisica","memoria":"semantica","meta":{"dominio":"ciencias_de_fora","fonte":""},"origem":"wiki","palavras_chave":[],"sinonimos":[],"tipo":"conceito","titulo":"Geografia Física / Sistema Terra"}
\section{Geografia Física / Sistema Terra}
O estudo da Terra como sistema integrado — atmosfera, hidrosfera, litosfera e biosfera acopladas por fluxos de matéria e energia.
\prov{relações: parte\_de geografia; relaciona vida; relaciona word}
\prov{proveniência: wiki \textperiodcentered{} fato \textperiodcentered{} confiança 1.0 \textperiodcentered{} domínio ciencias\_de\_fora \textperiodcentered{} história 4 versões no jornal}

%CRISTAL {"arestas":[{"alvo":"geografia","confianca":1.0,"epistemico":"fato","origem":"wiki","rel":"parte_de"},{"alvo":"word","confianca":0.5,"epistemico":"fato","origem":"wiki","rel":"relaciona"},{"alvo":"virtualizacao","confianca":0.5,"epistemico":"fato","origem":"wiki","rel":"relaciona"}],"confianca":1.0,"contraexemplos":[],"descricao":"O estudo do espaço como produto social — lugar, território, região e paisagem produzidos pela cultura, pela técnica e pelas relações de poder.","epistemico":"fato","exemplos":[],"id":"geografia_humana","memoria":"semantica","meta":{"dominio":"ciencias_de_fora","fonte":""},"origem":"wiki","palavras_chave":[],"sinonimos":[],"tipo":"conceito","titulo":"Geografia Humana"}
\section{Geografia Humana}
O estudo do espaço como produto social — lugar, território, região e paisagem produzidos pela cultura, pela técnica e pelas relações de poder.
\prov{relações: parte\_de geografia; relaciona word; relaciona virtualizacao}
\prov{proveniência: wiki \textperiodcentered{} fato \textperiodcentered{} confiança 1.0 \textperiodcentered{} domínio ciencias\_de\_fora \textperiodcentered{} história 4 versões no jornal}

%CRISTAL {"arestas":[{"alvo":"sig_sensoriamento","confianca":1.0,"epistemico":"fato","origem":"wiki","rel":"usa"},{"alvo":"capitalismo_de_vigilancia","confianca":1.0,"epistemico":"fato","origem":"wiki","rel":"usa"},{"alvo":"biopoder","confianca":1.0,"epistemico":"fato","origem":"wiki","rel":"relaciona"}],"confianca":1.0,"contraexemplos":[],"descricao":"A captura contínua de posição (GPS, redes móveis, apps) converte o espaço em dado monetizável e monitorável, permitindo rastrear, prever e modular comportamentos — a dimensão geográfica do capitalismo de vigilância.","epistemico":"inferencia","exemplos":[],"id":"geolocalizacao_vigilancia","memoria":"semantica","meta":{"autor":"—","dominio":"ciencias_de_fora","fonte":""},"origem":"wiki","palavras_chave":[],"sinonimos":[],"tipo":"conceito","titulo":"Geolocalização e Vigilância Espacial"}
\section{Geolocalização e Vigilância Espacial}
A captura contínua de posição (GPS, redes móveis, apps) converte o espaço em dado monetizável e monitorável, permitindo rastrear, prever e modular comportamentos — a dimensão geográfica do capitalismo de vigilância.
\prov{relações: usa sig\_sensoriamento; usa capitalismo\_de\_vigilancia; relaciona biopoder}
\prov{proveniência: wiki \textperiodcentered{} inferencia \textperiodcentered{} confiança 1.0 \textperiodcentered{} domínio ciencias\_de\_fora \textperiodcentered{} história 4 versões no jornal}

%CRISTAL {"arestas":[{"alvo":"geometria_ramo","confianca":1.0,"epistemico":"fato","origem":"wiki","rel":"parte_de"}],"confianca":1.0,"contraexemplos":[],"descricao":"Mandelbrot: as formas rugosas e auto-similares da natureza (costa, nuvem, floco) que a geometria lisa ignorava.","epistemico":"fato","exemplos":[],"id":"geometria_fractal","memoria":"semantica","meta":{"autor":"Mandelbrot","dominio":"ciencias_de_fora","fonte":""},"origem":"wiki","palavras_chave":[],"sinonimos":[],"tipo":"conceito","titulo":"Geometria Fractal"}
\section{Geometria Fractal}
Mandelbrot: as formas rugosas e auto-similares da natureza (costa, nuvem, floco) que a geometria lisa ignorava.
\prov{relações: parte\_de geometria\_ramo}
\prov{proveniência: wiki \textperiodcentered{} fato \textperiodcentered{} confiança 1.0 \textperiodcentered{} domínio ciencias\_de\_fora \textperiodcentered{} história 4 versões no jornal}

%CRISTAL {"arestas":[{"alvo":"geometria_ramo","confianca":1.0,"epistemico":"fato","origem":"wiki","rel":"parte_de"},{"alvo":"metodo_axiomatico","confianca":1.0,"epistemico":"fato","origem":"wiki","rel":"relaciona"}],"confianca":1.0,"contraexemplos":[],"descricao":"Lobachevsky, Bolyai, Riemann: negar o postulado das paralelas gera geometrias consistentes — o espaço não é único nem óbvio.","epistemico":"fato","exemplos":[],"id":"geometria_nao_euclidiana","memoria":"semantica","meta":{"dominio":"ciencias_de_fora","fonte":""},"origem":"wiki","palavras_chave":[],"sinonimos":[],"tipo":"conceito","titulo":"Geometria Não-Euclidiana"}
\section{Geometria Não-Euclidiana}
Lobachevsky, Bolyai, Riemann: negar o postulado das paralelas gera geometrias consistentes — o espaço não é único nem óbvio.
\prov{relações: parte\_de geometria\_ramo; relaciona metodo\_axiomatico}
\prov{proveniência: wiki \textperiodcentered{} fato \textperiodcentered{} confiança 1.0 \textperiodcentered{} domínio ciencias\_de\_fora \textperiodcentered{} história 6 versões no jornal}

%CRISTAL {"arestas":[{"alvo":"matematica","confianca":1.0,"epistemico":"fato","origem":"wiki","rel":"parte_de"},{"alvo":"topologia","confianca":1.0,"epistemico":"fato","origem":"wiki","rel":"relaciona"},{"alvo":"geometria","confianca":1.0,"epistemico":"fato","origem":"wiki","rel":"relaciona"}],"confianca":1.0,"contraexemplos":[],"descricao":"O estudo do espaço, da forma e da medida — do plano de Euclides às variedades curvas.","epistemico":"fato","exemplos":[],"id":"geometria_ramo","memoria":"semantica","meta":{"dominio":"ciencias_de_fora","fonte":""},"origem":"wiki","palavras_chave":[],"sinonimos":[],"tipo":"conceito","titulo":"Geometria"}
\section{Geometria}
O estudo do espaço, da forma e da medida — do plano de Euclides às variedades curvas.
\prov{relações: parte\_de matematica; relaciona topologia; relaciona geometria}
\prov{proveniência: wiki \textperiodcentered{} fato \textperiodcentered{} confiança 1.0 \textperiodcentered{} domínio ciencias\_de\_fora \textperiodcentered{} história 7 versões no jornal}

%CRISTAL {"arestas":[{"alvo":"cosmologia_quantica","confianca":1.0,"epistemico":"fato","origem":"curriculo","rel":"especializa"},{"alvo":"geometria","confianca":1.0,"epistemico":"fato","origem":"curriculo","rel":"especializa"},{"alvo":"geometria_nao_euclidiana","confianca":1.0,"epistemico":"fato","origem":"wiki","rel":"especializa"},{"alvo":"geometria_ramo","confianca":1.0,"epistemico":"fato","origem":"wiki","rel":"parte_de"},{"alvo":"geometria_riemanniana_de_ouro","confianca":1.0,"epistemico":"fato","origem":"wiki","rel":"relaciona"}],"confianca":1.0,"contraexemplos":[],"descricao":"Riemann: espaços curvos de qualquer dimensão, com métrica variável ponto a ponto — o palco da relatividade geral.","epistemico":"fato","exemplos":[],"id":"geometria_riemanniana","memoria":"semantica","meta":{"autor":"Riemann","dominio":"ciencias_de_fora","fonte":""},"origem":"wiki","palavras_chave":["métrica","curvatura","Hopf"],"sinonimos":["geometria de Riemann","Riemanniana"],"tipo":"conceito","titulo":"Geometria Riemanniana"}
\section{Geometria Riemanniana}
Riemann: espaços curvos de qualquer dimensão, com métrica variável ponto a ponto — o palco da relatividade geral.
\prov{sinónimos: geometria de Riemann; Riemanniana}
\prov{relações: especializa cosmologia\_quantica; especializa geometria; especializa geometria\_nao\_euclidiana; parte\_de geometria\_ramo; relaciona geometria\_riemanniana\_de\_ouro}
\prov{proveniência: wiki \textperiodcentered{} fato \textperiodcentered{} confiança 1.0 \textperiodcentered{} domínio ciencias\_de\_fora \textperiodcentered{} história 43 versões no jornal}

%CRISTAL {"arestas":[{"alvo":"geografia_fisica","confianca":1.0,"epistemico":"fato","origem":"wiki","rel":"parte_de"},{"alvo":"tectonica_placas","confianca":1.0,"epistemico":"fato","origem":"wiki","rel":"usa"}],"confianca":1.0,"contraexemplos":[],"descricao":"O estudo das formas do relevo e dos processos que as esculpem — intemperismo, erosão, transporte e deposição —, resultantes da tensão entre forças endógenas (tectônica) e exógenas (clima, água, gelo). Inclui o ciclo das rochas (ígneas/sedimentares/metamórficas).","epistemico":"fato","exemplos":[],"id":"geomorfologia","memoria":"semantica","meta":{"autor":"Davis/Gilbert","dominio":"ciencias_de_fora","fonte":""},"origem":"wiki","palavras_chave":[],"sinonimos":[],"tipo":"conceito","titulo":"Geomorfologia"}
\section{Geomorfologia}
O estudo das formas do relevo e dos processos que as esculpem — intemperismo, erosão, transporte e deposição —, resultantes da tensão entre forças endógenas (tectônica) e exógenas (clima, água, gelo). Inclui o ciclo das rochas (ígneas/sedimentares/metamórficas).
\prov{relações: parte\_de geografia\_fisica; usa tectonica\_placas}
\prov{proveniência: wiki \textperiodcentered{} fato \textperiodcentered{} confiança 1.0 \textperiodcentered{} domínio ciencias\_de\_fora \textperiodcentered{} história 3 versões no jornal}

%CRISTAL {"arestas":[{"alvo":"geografia_humana","confianca":1.0,"epistemico":"fato","origem":"wiki","rel":"parte_de"},{"alvo":"contrato_social","confianca":1.0,"epistemico":"fato","origem":"wiki","rel":"relaciona"},{"alvo":"biopoder","confianca":1.0,"epistemico":"fato","origem":"wiki","rel":"relaciona"}],"confianca":1.0,"contraexemplos":[],"descricao":"Kjellén: o estudo da relação entre poder político, o Estado e o espaço geográfico — como território, posição e recursos condicionam a força e a estratégia dos Estados. Nasce da concepção do Estado como organismo.","epistemico":"inferencia","exemplos":[],"id":"geopolitica","memoria":"semantica","meta":{"autor":"Kjellén","dominio":"ciencias_de_fora","fonte":""},"origem":"wiki","palavras_chave":[],"sinonimos":[],"tipo":"conceito","titulo":"Geopolítica"}
\section{Geopolítica}
Kjellén: o estudo da relação entre poder político, o Estado e o espaço geográfico — como território, posição e recursos condicionam a força e a estratégia dos Estados. Nasce da concepção do Estado como organismo.
\prov{relações: parte\_de geografia\_humana; relaciona contrato\_social; relaciona biopoder}
\prov{proveniência: wiki \textperiodcentered{} inferencia \textperiodcentered{} confiança 1.0 \textperiodcentered{} domínio ciencias\_de\_fora \textperiodcentered{} história 4 versões no jornal}

%CRISTAL {"arestas":[{"alvo":"medicina","confianca":1.0,"epistemico":"fato","origem":"wiki","rel":"parte_de"},{"alvo":"modelo_biomedico","confianca":1.0,"epistemico":"fato","origem":"wiki","rel":"usa"},{"alvo":"selecao_natural","confianca":1.0,"epistemico":"fato","origem":"wiki","rel":"relaciona"}],"confianca":1.0,"contraexemplos":[],"descricao":"Doenças infecciosas são causadas por microrganismos específicos, não por 'miasmas'. Semmelweis mostrou que a antissepsia das mãos reduzia a febre puerperal; Pasteur refutou a geração espontânea; Koch estabeleceu postulados de causalidade micróbio-doença.","epistemico":"fato","exemplos":[],"id":"germ_theory","memoria":"semantica","meta":{"autor":"Pasteur/Koch/Semmelweis","dominio":"ciencias_de_fora","fonte":""},"origem":"wiki","palavras_chave":[],"sinonimos":[],"tipo":"conceito","titulo":"Teoria dos Germes (Pasteur, Koch, Semmelweis)"}
\section{Teoria dos Germes (Pasteur, Koch, Semmelweis)}
Doenças infecciosas são causadas por microrganismos específicos, não por 'miasmas'. Semmelweis mostrou que a antissepsia das mãos reduzia a febre puerperal; Pasteur refutou a geração espontânea; Koch estabeleceu postulados de causalidade micróbio-doença.
\prov{relações: parte\_de medicina; usa modelo\_biomedico; relaciona selecao\_natural}
\prov{proveniência: wiki \textperiodcentered{} fato \textperiodcentered{} confiança 1.0 \textperiodcentered{} domínio ciencias\_de\_fora \textperiodcentered{} história 4 versões no jornal}

%CRISTAL {"arestas":[{"alvo":"heidegger_tecnica_desvelamento","confianca":1.0,"epistemico":"fato","origem":"wiki","rel":"usa"},{"alvo":"biopoder","confianca":1.0,"epistemico":"fato","origem":"wiki","rel":"relaciona"}],"confianca":1.0,"contraexemplos":[],"descricao":"Heidegger: a essência da técnica moderna — uma força que 'desafia' tudo o que existe a se mostrar como ordenável e otimizável, inclusive o próprio humano. O modo de desvelamento próprio da modernidade.","epistemico":"inferencia","exemplos":[],"id":"gestell_enquadramento","memoria":"semantica","meta":{"autor":"Heidegger","dominio":"ciencias_de_fora","fonte":""},"origem":"wiki","palavras_chave":[],"sinonimos":[],"tipo":"conceito","titulo":"Gestell (Enquadramento)"}
\section{Gestell (Enquadramento)}
Heidegger: a essência da técnica moderna — uma força que 'desafia' tudo o que existe a se mostrar como ordenável e otimizável, inclusive o próprio humano. O modo de desvelamento próprio da modernidade.
\prov{relações: usa heidegger\_tecnica\_desvelamento; relaciona biopoder}
\prov{proveniência: wiki \textperiodcentered{} inferencia \textperiodcentered{} confiança 1.0 \textperiodcentered{} domínio ciencias\_de\_fora \textperiodcentered{} história 3 versões no jornal}

%CRISTAL {"arestas":[{"alvo":"meio_tecnico_cientifico_informacional","confianca":1.0,"epistemico":"fato","origem":"wiki","rel":"usa"},{"alvo":"simulacro_simulacao","confianca":1.0,"epistemico":"fato","origem":"wiki","rel":"relaciona"}],"confianca":1.0,"contraexemplos":[],"descricao":"Milton Santos: a globalização como nos 'fazem crer' (fábula: aldeia global, fim das fronteiras) esconde a globalização 'como é' (perversidade: pobreza, tirania do dinheiro e da informação), abrindo espaço para uma 'outra globalização' possível a partir dos lugares.","epistemico":"inferencia","exemplos":[],"id":"globalizacao_como_fabula","memoria":"semantica","meta":{"autor":"Milton Santos","dominio":"ciencias_de_fora","fonte":""},"origem":"wiki","palavras_chave":[],"sinonimos":[],"tipo":"conceito","titulo":"Globalização como Fábula (Milton Santos)"}
\section{Globalização como Fábula (Milton Santos)}
Milton Santos: a globalização como nos 'fazem crer' (fábula: aldeia global, fim das fronteiras) esconde a globalização 'como é' (perversidade: pobreza, tirania do dinheiro e da informação), abrindo espaço para uma 'outra globalização' possível a partir dos lugares.
\prov{relações: usa meio\_tecnico\_cientifico\_informacional; relaciona simulacro\_simulacao}
\prov{proveniência: wiki \textperiodcentered{} inferencia \textperiodcentered{} confiança 1.0 \textperiodcentered{} domínio ciencias\_de\_fora \textperiodcentered{} história 3 versões no jornal}

%CRISTAL {"arestas":[{"alvo":"filosofia_da_computacao","confianca":1.0,"epistemico":"fato","origem":"wiki","rel":"parte_de"},{"alvo":"inteligencia_artificial","confianca":1.0,"epistemico":"fato","origem":"wiki","rel":"parte_de"},{"alvo":"hopfield","confianca":1.0,"epistemico":"fato","origem":"wiki","rel":"referencia"},{"alvo":"simbolo_vs_significado","confianca":1.0,"epistemico":"fato","origem":"wiki","rel":"relaciona"}],"confianca":1.0,"contraexemplos":[],"descricao":"Os dois paradigmas da IA: a simbólica (GOFAI) representa o conhecimento por símbolos, regras e lógica (interpretável, frágil); o conexionismo aprende representações distribuídas em redes neurais a partir de dados (robusto, opaco). Hopfield é o ancestral conexionista.","epistemico":"inferencia","exemplos":[],"id":"gofai_vs_conexionismo","memoria":"semantica","meta":{"autor":"Haugeland/Rumelhart","dominio":"ciencias_de_fora","fonte":""},"origem":"wiki","palavras_chave":[],"sinonimos":[],"tipo":"conceito","titulo":"IA Simbólica (GOFAI) × Conexionismo"}
\section{IA Simbólica (GOFAI) $\times$ Conexionismo}
Os dois paradigmas da IA: a simbólica (GOFAI) representa o conhecimento por símbolos, regras e lógica (interpretável, frágil); o conexionismo aprende representações distribuídas em redes neurais a partir de dados (robusto, opaco). Hopfield é o ancestral conexionista.
\prov{relações: parte\_de filosofia\_da\_computacao; parte\_de inteligencia\_artificial; referencia hopfield; relaciona simbolo\_vs\_significado}
\prov{proveniência: wiki \textperiodcentered{} inferencia \textperiodcentered{} confiança 1.0 \textperiodcentered{} domínio ciencias\_de\_fora \textperiodcentered{} história 5 versões no jornal}

%CRISTAL {"arestas":[{"alvo":"beleza_desinteressada_kant","confianca":1.0,"epistemico":"fato","origem":"wiki","rel":"contradiz"},{"alvo":"word","confianca":0.5,"epistemico":"fato","origem":"wiki","rel":"relaciona"},{"alvo":"virtualizacao","confianca":0.5,"epistemico":"fato","origem":"wiki","rel":"relaciona"}],"confianca":1.0,"contraexemplos":[],"descricao":"Hume: a beleza está na mente, mas há um padrão de gosto via 'delicadeza' — o juízo dos observadores refinados.","epistemico":"inferencia","exemplos":[],"id":"gosto_hume","memoria":"semantica","meta":{"autor":"Hume","dominio":"ciencias_de_fora","fonte":""},"origem":"wiki","palavras_chave":[],"sinonimos":[],"tipo":"conceito","titulo":"O Padrão do Gosto"}
\section{O Padrão do Gosto}
Hume: a beleza está na mente, mas há um padrão de gosto via 'delicadeza' — o juízo dos observadores refinados.
\prov{relações: contradiz beleza\_desinteressada\_kant; relaciona word; relaciona virtualizacao}
\prov{proveniência: wiki \textperiodcentered{} inferencia \textperiodcentered{} confiança 1.0 \textperiodcentered{} domínio ciencias\_de\_fora \textperiodcentered{} história 4 versões no jornal}

%CRISTAL {"arestas":[{"alvo":"relacoes_internacionais","confianca":1.0,"epistemico":"fato","origem":"wiki","rel":"parte_de"},{"alvo":"direito_internacional","confianca":1.0,"epistemico":"fato","origem":"wiki","rel":"usa"},{"alvo":"tragedia_dos_comuns","confianca":1.0,"epistemico":"fato","origem":"wiki","rel":"usa"}],"confianca":1.0,"contraexemplos":[],"descricao":"O conjunto de instituições, normas e regimes (ONU, tratados) por meio dos quais os Estados coordenam respostas a problemas transfronteiriços na ausência de um governo mundial — enfrentando o dilema da ação coletiva global (clima).","epistemico":"inferencia","exemplos":[],"id":"governanca_global","memoria":"semantica","meta":{"autor":"Keohane","dominio":"ciencias_de_fora","fonte":""},"origem":"wiki","palavras_chave":[],"sinonimos":[],"tipo":"conceito","titulo":"Governança Global"}
\section{Governança Global}
O conjunto de instituições, normas e regimes (ONU, tratados) por meio dos quais os Estados coordenam respostas a problemas transfronteiriços na ausência de um governo mundial — enfrentando o dilema da ação coletiva global (clima).
\prov{relações: parte\_de relacoes\_internacionais; usa direito\_internacional; usa tragedia\_dos\_comuns}
\prov{proveniência: wiki \textperiodcentered{} inferencia \textperiodcentered{} confiança 1.0 \textperiodcentered{} domínio ciencias\_de\_fora \textperiodcentered{} história 4 versões no jornal}

%CRISTAL {"arestas":[{"alvo":"cristianismo","confianca":1.0,"epistemico":"fato","origem":"wiki","rel":"parte_de"},{"alvo":"queda_pecado","confianca":1.0,"epistemico":"fato","origem":"wiki","rel":"usa"},{"alvo":"sentido_da_vida","confianca":1.0,"epistemico":"fato","origem":"wiki","rel":"relaciona"}],"confianca":1.0,"contraexemplos":[],"descricao":"Bíblia (Efésios 2): o dom imerecido de Deus que restaura a relação rompida pelo pecado, culminando na salvação oferecida gratuitamente — 'pela graça sois salvos, por meio da fé; e isto é dom de Deus'. A salvação por dom, não por mérito.","epistemico":"inferencia","exemplos":[],"id":"graca_redencao","memoria":"semantica","meta":{"autor":"Bíblia","dominio":"sagrado","fonte":""},"origem":"wiki","palavras_chave":[],"sinonimos":[],"tipo":"conceito","titulo":"Graça e Redenção"}
\section{Graça e Redenção}
Bíblia (Efésios 2): o dom imerecido de Deus que restaura a relação rompida pelo pecado, culminando na salvação oferecida gratuitamente — 'pela graça sois salvos, por meio da fé; e isto é dom de Deus'. A salvação por dom, não por mérito.
\prov{relações: parte\_de cristianismo; usa queda\_pecado; relaciona sentido\_da\_vida}
\prov{proveniência: wiki \textperiodcentered{} inferencia \textperiodcentered{} confiança 1.0 \textperiodcentered{} domínio sagrado \textperiodcentered{} história 4 versões no jornal}

%CRISTAL {"arestas":[{"alvo":"alemao","confianca":1.0,"epistemico":"fato","origem":"wiki","rel":"explica"},{"alvo":"caso_gramatical","confianca":1.0,"epistemico":"fato","origem":"wiki","rel":"usa"},{"alvo":"gramatica","confianca":1.0,"epistemico":"fato","origem":"wiki","rel":"especializa"}],"confianca":1.0,"contraexemplos":[],"descricao":"Verbo em 2ª posição na principal, mas VERBO NO FIM na subordinada. 4 casos, 3 gêneros, adjetivos declinados; substantivos compostos infinitos (Donaudampfschifffahrt). A sintaxe que 'segura o verbo' até o fim.","epistemico":"fato","exemplos":[],"id":"gramatica_alemao","memoria":"semantica","meta":{"dominio":"ciencias_de_fora","fonte":"perfil sintático"},"origem":"wiki","palavras_chave":[],"sinonimos":[],"tipo":"conceito","titulo":"Gramática do Alemão"}
\section{Gramática do Alemão}
Verbo em 2ª posição na principal, mas VERBO NO FIM na subordinada. 4 casos, 3 gêneros, adjetivos declinados; substantivos compostos infinitos (Donaudampfschifffahrt). A sintaxe que 'segura o verbo' até o fim.
\prov{relações: explica alemao; usa caso\_gramatical; especializa gramatica}
\prov{proveniência: wiki \textperiodcentered{} perfil sintático \textperiodcentered{} fato \textperiodcentered{} confiança 1.0 \textperiodcentered{} domínio ciencias\_de\_fora \textperiodcentered{} história 4 versões no jornal}

%CRISTAL {"arestas":[{"alvo":"arabe","confianca":1.0,"epistemico":"fato","origem":"wiki","rel":"explica"},{"alvo":"ordem_das_palavras","confianca":1.0,"epistemico":"fato","origem":"wiki","rel":"usa"},{"alvo":"morfologia","confianca":1.0,"epistemico":"fato","origem":"wiki","rel":"usa"},{"alvo":"gramatica","confianca":1.0,"epistemico":"fato","origem":"wiki","rel":"especializa"}],"confianca":1.0,"contraexemplos":[],"descricao":"VSO clássico, escrita RTL, morfologia de RAIZ-E-PADRÃO (a raiz k-t-b gera kitāb, kātib, maktab). Três números (singular/DUAL/plural), 3 casos, gênero em verbos e adjetivos, forte diglossia clássico↔dialetal.","epistemico":"fato","exemplos":[],"id":"gramatica_arabe","memoria":"semantica","meta":{"dominio":"ciencias_de_fora","fonte":"perfil sintático"},"origem":"wiki","palavras_chave":[],"sinonimos":[],"tipo":"conceito","titulo":"Gramática do Árabe"}
\section{Gramática do Árabe}
VSO clássico, escrita RTL, morfologia de RAIZ-E-PADRÃO (a raiz k-t-b gera kit\ensuremath{\bar{a}}b, k\ensuremath{\bar{a}}tib, maktab). Três números (singular/DUAL/plural), 3 casos, gênero em verbos e adjetivos, forte diglossia clássico\ensuremath{\leftrightarrow}dialetal.
\prov{relações: explica arabe; usa ordem\_das\_palavras; usa morfologia; especializa gramatica}
\prov{proveniência: wiki \textperiodcentered{} perfil sintático \textperiodcentered{} fato \textperiodcentered{} confiança 1.0 \textperiodcentered{} domínio ciencias\_de\_fora \textperiodcentered{} história 5 versões no jornal}

%CRISTAL {"arestas":[{"alvo":"espanhol","confianca":1.0,"epistemico":"fato","origem":"wiki","rel":"explica"},{"alvo":"gramatica_portugues","confianca":1.0,"epistemico":"fato","origem":"wiki","rel":"relaciona"},{"alvo":"lingua_pro_drop","confianca":1.0,"epistemico":"fato","origem":"wiki","rel":"usa"},{"alvo":"gramatica","confianca":1.0,"epistemico":"fato","origem":"wiki","rel":"especializa"}],"confianca":1.0,"contraexemplos":[],"descricao":"Muito próxima do português: SVO, pro-drop, fusional, dois gêneros, subjuntivo forte. Difere em detalhes (ser/estar, pretéritos, 'vosotros'); a mão dupla que engana o tradutor com falsos amigos.","epistemico":"fato","exemplos":[],"id":"gramatica_espanhol","memoria":"semantica","meta":{"dominio":"ciencias_de_fora","fonte":"perfil sintático"},"origem":"wiki","palavras_chave":[],"sinonimos":[],"tipo":"conceito","titulo":"Gramática do Espanhol"}
\section{Gramática do Espanhol}
Muito próxima do português: SVO, pro-drop, fusional, dois gêneros, subjuntivo forte. Difere em detalhes (ser/estar, pretéritos, 'vosotros'); a mão dupla que engana o tradutor com falsos amigos.
\prov{relações: explica espanhol; relaciona gramatica\_portugues; usa lingua\_pro\_drop; especializa gramatica}
\prov{proveniência: wiki \textperiodcentered{} perfil sintático \textperiodcentered{} fato \textperiodcentered{} confiança 1.0 \textperiodcentered{} domínio ciencias\_de\_fora \textperiodcentered{} história 5 versões no jornal}

%CRISTAL {"arestas":[{"alvo":"filosofia_da_linguagem","confianca":1.0,"epistemico":"fato","origem":"wiki","rel":"parte_de"},{"alvo":"formalismo_hilbert","confianca":1.0,"epistemico":"fato","origem":"wiki","rel":"relaciona"},{"alvo":"word","confianca":0.5,"epistemico":"fato","origem":"wiki","rel":"relaciona"},{"alvo":"noam_chomsky","confianca":1.0,"epistemico":"fato","origem":"wiki","rel":"parte_de"},{"alvo":"gramatica","confianca":1.0,"epistemico":"fato","origem":"wiki","rel":"especializa"}],"confianca":1.0,"contraexemplos":[],"descricao":"Chomsky: uma língua é um conjunto (infinito) de frases gerado por regras finitas — a competência, não o desempenho.","epistemico":"fato","exemplos":[],"id":"gramatica_gerativa","memoria":"semantica","meta":{"autor":"Chomsky","dominio":"ciencias_de_fora","fonte":""},"origem":"wiki","palavras_chave":[],"sinonimos":[],"tipo":"conceito","titulo":"Gramática Gerativa"}
\section{Gramática Gerativa}
Chomsky: uma língua é um conjunto (infinito) de frases gerado por regras finitas — a competência, não o desempenho.
\prov{relações: parte\_de filosofia\_da\_linguagem; relaciona formalismo\_hilbert; relaciona word; parte\_de noam\_chomsky; especializa gramatica}
\prov{proveniência: wiki \textperiodcentered{} fato \textperiodcentered{} confiança 1.0 \textperiodcentered{} domínio ciencias\_de\_fora \textperiodcentered{} história 7 versões no jornal}

%CRISTAL {"arestas":[{"alvo":"ingles","confianca":1.0,"epistemico":"fato","origem":"wiki","rel":"explica"},{"alvo":"ordem_das_palavras","confianca":1.0,"epistemico":"fato","origem":"wiki","rel":"usa"},{"alvo":"tipologia_morfologica","confianca":1.0,"epistemico":"fato","origem":"wiki","rel":"usa"},{"alvo":"gramatica","confianca":1.0,"epistemico":"fato","origem":"wiki","rel":"especializa"}],"confianca":1.0,"contraexemplos":[],"descricao":"SVO rígida (a ordem CARREGA o papel, pois quase não há caso), analítica. Sem gênero gramatical, morfologia verbal mínima (-s, -ed, -ing); tempos compostos por auxiliares. Não é pro-drop. Adjetivo ANTES do nome.","epistemico":"fato","exemplos":[],"id":"gramatica_ingles","memoria":"semantica","meta":{"dominio":"ciencias_de_fora","fonte":"perfil sintático"},"origem":"wiki","palavras_chave":[],"sinonimos":[],"tipo":"conceito","titulo":"Gramática do Inglês"}
\section{Gramática do Inglês}
SVO rígida (a ordem CARREGA o papel, pois quase não há caso), analítica. Sem gênero gramatical, morfologia verbal mínima (-s, -ed, -ing); tempos compostos por auxiliares. Não é pro-drop. Adjetivo ANTES do nome.
\prov{relações: explica ingles; usa ordem\_das\_palavras; usa tipologia\_morfologica; especializa gramatica}
\prov{proveniência: wiki \textperiodcentered{} perfil sintático \textperiodcentered{} fato \textperiodcentered{} confiança 1.0 \textperiodcentered{} domínio ciencias\_de\_fora \textperiodcentered{} história 5 versões no jornal}

%CRISTAL {"arestas":[{"alvo":"japones","confianca":1.0,"epistemico":"fato","origem":"wiki","rel":"explica"},{"alvo":"ordem_das_palavras","confianca":1.0,"epistemico":"fato","origem":"wiki","rel":"usa"},{"alvo":"topico_proeminencia","confianca":1.0,"epistemico":"fato","origem":"wiki","rel":"usa"},{"alvo":"tipologia_morfologica","confianca":1.0,"epistemico":"fato","origem":"wiki","rel":"usa"},{"alvo":"gramatica","confianca":1.0,"epistemico":"fato","origem":"wiki","rel":"especializa"}],"confianca":1.0,"contraexemplos":[],"descricao":"SOV estrito, aglutinante, TÓPICO-proeminente (partícula は wa). PARTÍCULAS pós-postas marcam o papel (が ga sujeito, を o objeto); verbo sempre no fim; honoríficos elaborados; sem gênero nem plural obrigatório.","epistemico":"fato","exemplos":[],"id":"gramatica_japones","memoria":"semantica","meta":{"dominio":"ciencias_de_fora","fonte":"perfil sintático"},"origem":"wiki","palavras_chave":[],"sinonimos":[],"tipo":"conceito","titulo":"Gramática do Japonês"}
\section{Gramática do Japonês}
SOV estrito, aglutinante, TÓPICO-proeminente (partícula \texttt{U+306F} wa). PARTÍCULAS pós-postas marcam o papel (\texttt{U+304C} ga sujeito, \texttt{U+3092} o objeto); verbo sempre no fim; honoríficos elaborados; sem gênero nem plural obrigatório.
\prov{relações: explica japones; usa ordem\_das\_palavras; usa topico\_proeminencia; usa tipologia\_morfologica; especializa gramatica}
\prov{proveniência: wiki \textperiodcentered{} perfil sintático \textperiodcentered{} fato \textperiodcentered{} confiança 1.0 \textperiodcentered{} domínio ciencias\_de\_fora \textperiodcentered{} história 6 versões no jornal}

%CRISTAL {"arestas":[{"alvo":"latim","confianca":1.0,"epistemico":"fato","origem":"wiki","rel":"explica"},{"alvo":"caso_gramatical","confianca":1.0,"epistemico":"fato","origem":"wiki","rel":"usa"},{"alvo":"ordem_das_palavras","confianca":1.0,"epistemico":"fato","origem":"wiki","rel":"usa"},{"alvo":"gramatica","confianca":1.0,"epistemico":"fato","origem":"wiki","rel":"especializa"}],"confianca":1.0,"contraexemplos":[],"descricao":"Ordem LIVRE (tende a SOV) porque o CASO carrega o papel: 6 casos × 5 declinações. Fusional denso, 3 gêneros, sem artigos. Conjugação verbal completa. A língua que a árvore sintática do português veio simplificando.","epistemico":"fato","exemplos":[],"id":"gramatica_latim","memoria":"semantica","meta":{"dominio":"ciencias_de_fora","fonte":"perfil sintático"},"origem":"wiki","palavras_chave":[],"sinonimos":[],"tipo":"conceito","titulo":"Gramática do Latim"}
\section{Gramática do Latim}
Ordem LIVRE (tende a SOV) porque o CASO carrega o papel: 6 casos $\times$ 5 declinações. Fusional denso, 3 gêneros, sem artigos. Conjugação verbal completa. A língua que a árvore sintática do português veio simplificando.
\prov{relações: explica latim; usa caso\_gramatical; usa ordem\_das\_palavras; especializa gramatica}
\prov{proveniência: wiki \textperiodcentered{} perfil sintático \textperiodcentered{} fato \textperiodcentered{} confiança 1.0 \textperiodcentered{} domínio ciencias\_de\_fora \textperiodcentered{} história 5 versões no jornal}

%CRISTAL {"arestas":[{"alvo":"mandarim","confianca":1.0,"epistemico":"fato","origem":"wiki","rel":"explica"},{"alvo":"tipologia_morfologica","confianca":1.0,"epistemico":"fato","origem":"wiki","rel":"usa"},{"alvo":"topico_proeminencia","confianca":1.0,"epistemico":"fato","origem":"wiki","rel":"usa"},{"alvo":"gramatica","confianca":1.0,"epistemico":"fato","origem":"wiki","rel":"especializa"}],"confianca":1.0,"contraexemplos":[],"descricao":"SVO mas TÓPICO-proeminente; isolante puro: palavras invariáveis, ZERO conjugação ou declinação. Aspecto por partículas (了 le), classificadores obrigatórios (三个人), 4 tons distinguem sentido. Nada de tempo verbal, gênero ou plural morfológico.","epistemico":"fato","exemplos":[],"id":"gramatica_mandarim","memoria":"semantica","meta":{"dominio":"ciencias_de_fora","fonte":"perfil sintático"},"origem":"wiki","palavras_chave":[],"sinonimos":[],"tipo":"conceito","titulo":"Gramática do Mandarim"}
\section{Gramática do Mandarim}
SVO mas TÓPICO-proeminente; isolante puro: palavras invariáveis, ZERO conjugação ou declinação. Aspecto por partículas (\texttt{U+4E86} le), classificadores obrigatórios (\texttt{U+4E09}\texttt{U+4E2A}\texttt{U+4EBA}), 4 tons distinguem sentido. Nada de tempo verbal, gênero ou plural morfológico.
\prov{relações: explica mandarim; usa tipologia\_morfologica; usa topico\_proeminencia; especializa gramatica}
\prov{proveniência: wiki \textperiodcentered{} perfil sintático \textperiodcentered{} fato \textperiodcentered{} confiança 1.0 \textperiodcentered{} domínio ciencias\_de\_fora \textperiodcentered{} história 5 versões no jornal}

%CRISTAL {"arestas":[{"alvo":"portugues","confianca":1.0,"epistemico":"fato","origem":"wiki","rel":"explica"},{"alvo":"lingua_pro_drop","confianca":1.0,"epistemico":"fato","origem":"wiki","rel":"usa"},{"alvo":"genero_gramatical","confianca":1.0,"epistemico":"fato","origem":"wiki","rel":"usa"},{"alvo":"tempo_aspecto_modo","confianca":1.0,"epistemico":"fato","origem":"wiki","rel":"usa"},{"alvo":"gramatica","confianca":1.0,"epistemico":"fato","origem":"wiki","rel":"especializa"}],"confianca":1.0,"contraexemplos":[],"descricao":"SVO, pro-drop, fusional. Dois gêneros (M/F) com concordância de artigo, adjetivo e particípio; adjetivo em geral APÓS o nome. Verbo riquíssimo: ~50 formas por verbo (tempo/modo/pessoa), subjuntivo vivo, dois particípios. Sem caso nominal.","epistemico":"fato","exemplos":[],"id":"gramatica_portugues","memoria":"semantica","meta":{"dominio":"ciencias_de_fora","fonte":"perfil sintático"},"origem":"wiki","palavras_chave":[],"sinonimos":[],"tipo":"conceito","titulo":"Gramática do Português"}
\section{Gramática do Português}
SVO, pro-drop, fusional. Dois gêneros (M/F) com concordância de artigo, adjetivo e particípio; adjetivo em geral APÓS o nome. Verbo riquíssimo: \~{}50 formas por verbo (tempo/modo/pessoa), subjuntivo vivo, dois particípios. Sem caso nominal.
\prov{relações: explica portugues; usa lingua\_pro\_drop; usa genero\_gramatical; usa tempo\_aspecto\_modo; especializa gramatica}
\prov{proveniência: wiki \textperiodcentered{} perfil sintático \textperiodcentered{} fato \textperiodcentered{} confiança 1.0 \textperiodcentered{} domínio ciencias\_de\_fora \textperiodcentered{} história 6 versões no jornal}

%CRISTAL {"arestas":[{"alvo":"russo","confianca":1.0,"epistemico":"fato","origem":"wiki","rel":"explica"},{"alvo":"caso_gramatical","confianca":1.0,"epistemico":"fato","origem":"wiki","rel":"usa"},{"alvo":"tempo_aspecto_modo","confianca":1.0,"epistemico":"fato","origem":"wiki","rel":"usa"},{"alvo":"gramatica","confianca":1.0,"epistemico":"fato","origem":"wiki","rel":"especializa"}],"confianca":1.0,"contraexemplos":[],"descricao":"Ordem relativamente livre (o caso decide o papel): 6 casos, 3 gêneros, sem artigos. O coração é o ASPECTO verbal — cada verbo vem em par perfectivo/imperfectivo. Escrita cirílica.","epistemico":"fato","exemplos":[],"id":"gramatica_russo","memoria":"semantica","meta":{"dominio":"ciencias_de_fora","fonte":"perfil sintático"},"origem":"wiki","palavras_chave":[],"sinonimos":[],"tipo":"conceito","titulo":"Gramática do Russo"}
\section{Gramática do Russo}
Ordem relativamente livre (o caso decide o papel): 6 casos, 3 gêneros, sem artigos. O coração é o ASPECTO verbal — cada verbo vem em par perfectivo/imperfectivo. Escrita cirílica.
\prov{relações: explica russo; usa caso\_gramatical; usa tempo\_aspecto\_modo; especializa gramatica}
\prov{proveniência: wiki \textperiodcentered{} perfil sintático \textperiodcentered{} fato \textperiodcentered{} confiança 1.0 \textperiodcentered{} domínio ciencias\_de\_fora \textperiodcentered{} história 5 versões no jornal}

%CRISTAL {"arestas":[{"alvo":"tupi","confianca":1.0,"epistemico":"fato","origem":"wiki","rel":"explica"},{"alvo":"tipologia_morfologica","confianca":1.0,"epistemico":"fato","origem":"wiki","rel":"usa"},{"alvo":"gramatica","confianca":1.0,"epistemico":"fato","origem":"wiki","rel":"especializa"}],"confianca":1.0,"contraexemplos":[],"descricao":"Aglutinante e polissintética, com INCORPORAÇÃO nominal (o objeto entra no verbo). Marcas de pessoa prefixadas, sem gênero; deu ao português do Brasil o léxico e a fonética de milhares de topônimos.","epistemico":"inferencia","exemplos":[],"id":"gramatica_tupi","memoria":"semantica","meta":{"dominio":"ciencias_de_fora","fonte":"perfil sintático"},"origem":"wiki","palavras_chave":[],"sinonimos":[],"tipo":"conceito","titulo":"Gramática do Tupi"}
\section{Gramática do Tupi}
Aglutinante e polissintética, com INCORPORAÇÃO nominal (o objeto entra no verbo). Marcas de pessoa prefixadas, sem gênero; deu ao português do Brasil o léxico e a fonética de milhares de topônimos.
\prov{relações: explica tupi; usa tipologia\_morfologica; especializa gramatica}
\prov{proveniência: wiki \textperiodcentered{} perfil sintático \textperiodcentered{} inferencia \textperiodcentered{} confiança 1.0 \textperiodcentered{} domínio ciencias\_de\_fora \textperiodcentered{} história 4 versões no jornal}

%CRISTAL {"arestas":[{"alvo":"noam_chomsky","confianca":1.0,"epistemico":"fato","origem":"wiki","rel":"parte_de"},{"alvo":"gramatica_gerativa","confianca":1.0,"epistemico":"fato","origem":"wiki","rel":"usa"}],"confianca":1.0,"contraexemplos":[],"descricao":"Chomsky: uma faculdade inata comum a toda língua humana. Tese forte sobre a mente, difícil de falsear.","epistemico":"hipotese","exemplos":[],"id":"gramatica_universal","memoria":"semantica","meta":{"autor":"Chomsky","dominio":"ciencias_de_fora","fonte":""},"origem":"wiki","palavras_chave":[],"sinonimos":[],"tipo":"conceito","titulo":"Gramática Universal"}
\section{Gramática Universal}
Chomsky: uma faculdade inata comum a toda língua humana. Tese forte sobre a mente, difícil de falsear.
\prov{relações: parte\_de noam\_chomsky; usa gramatica\_gerativa}
\prov{proveniência: wiki \textperiodcentered{} hipotese \textperiodcentered{} confiança 1.0 \textperiodcentered{} domínio ciencias\_de\_fora \textperiodcentered{} história 3 versões no jornal}

%CRISTAL {"arestas":[{"alvo":"indo_europeu","confianca":1.0,"epistemico":"fato","origem":"wiki","rel":"parte_de"}],"confianca":1.0,"contraexemplos":[],"descricao":"Língua de registro contínuo há ~3400 anos — a da filosofia, do Novo Testamento e de metade do vocabulário científico.","epistemico":"fato","exemplos":[],"id":"grego","memoria":"semantica","meta":{"dominio":"ciencias_de_fora","fonte":""},"origem":"wiki","palavras_chave":[],"sinonimos":[],"tipo":"conceito","titulo":"Grego"}
\section{Grego}
Língua de registro contínuo há \~{}3400 anos — a da filosofia, do Novo Testamento e de metade do vocabulário científico.
\prov{relações: parte\_de indo\_europeu}
\prov{proveniência: wiki \textperiodcentered{} fato \textperiodcentered{} confiança 1.0 \textperiodcentered{} domínio ciencias\_de\_fora \textperiodcentered{} história 2 versões no jornal}

%CRISTAL {"arestas":[{"alvo":"genetica","confianca":1.0,"epistemico":"fato","origem":"wiki","rel":"parte_de"}],"confianca":1.0,"contraexemplos":[],"descricao":"O monge que descobriu as leis da herança contando ervilhas — a genética antes de se saber o que era o gene.","epistemico":"fato","exemplos":[],"id":"gregor_mendel","memoria":"semantica","meta":{"autor":"Mendel","dominio":"ciencias_de_fora","fonte":""},"origem":"wiki","palavras_chave":[],"sinonimos":[],"tipo":"conceito","titulo":"Gregor Mendel"}
\section{Gregor Mendel}
O monge que descobriu as leis da herança contando ervilhas — a genética antes de se saber o que era o gene.
\prov{relações: parte\_de genetica}
\prov{proveniência: wiki \textperiodcentered{} fato \textperiodcentered{} confiança 1.0 \textperiodcentered{} domínio ciencias\_de\_fora \textperiodcentered{} história 2 versões no jornal}

%CRISTAL {"arestas":[{"alvo":"algebra_abstrata","confianca":1.0,"epistemico":"fato","origem":"wiki","rel":"parte_de"},{"alvo":"geometria_ramo","confianca":1.0,"epistemico":"fato","origem":"wiki","rel":"relaciona"}],"confianca":1.0,"contraexemplos":[],"descricao":"Grothendieck: refundou a geometria algébrica em esquemas e feixes — resolver dissolvendo o problema num nível mais geral.","epistemico":"inferencia","exemplos":[],"id":"grothendieck_abstracao","memoria":"semantica","meta":{"autor":"Grothendieck","dominio":"ciencias_de_fora","fonte":""},"origem":"wiki","palavras_chave":[],"sinonimos":[],"tipo":"conceito","titulo":"Grothendieck e a Abstração"}
\section{Grothendieck e a Abstração}
Grothendieck: refundou a geometria algébrica em esquemas e feixes — resolver dissolvendo o problema num nível mais geral.
\prov{relações: parte\_de algebra\_abstrata; relaciona geometria\_ramo}
\prov{proveniência: wiki \textperiodcentered{} inferencia \textperiodcentered{} confiança 1.0 \textperiodcentered{} domínio ciencias\_de\_fora \textperiodcentered{} história 6 versões no jornal}

%CRISTAL {"arestas":[{"alvo":"direito","confianca":1.0,"epistemico":"fato","origem":"wiki","rel":"parte_de"},{"alvo":"contrato_social","confianca":1.0,"epistemico":"fato","origem":"wiki","rel":"relaciona"},{"alvo":"refeudalizacao_esfera_publica","confianca":1.0,"epistemico":"fato","origem":"wiki","rel":"relaciona"}],"confianca":1.0,"contraexemplos":[],"descricao":"Habermas: a legitimidade do direito nasce de procedimentos discursivos de formação da opinião e vontade públicas; autonomia pública (o cidadão coautor das leis) e autonomia privada (as liberdades individuais) são co-originárias.","epistemico":"inferencia","exemplos":[],"id":"habermas_direito_democracia","memoria":"semantica","meta":{"autor":"Habermas","dominio":"ciencias_de_fora","fonte":""},"origem":"wiki","palavras_chave":[],"sinonimos":[],"tipo":"conceito","titulo":"Direito e Democracia Deliberativa (Habermas)"}
\section{Direito e Democracia Deliberativa (Habermas)}
Habermas: a legitimidade do direito nasce de procedimentos discursivos de formação da opinião e vontade públicas; autonomia pública (o cidadão coautor das leis) e autonomia privada (as liberdades individuais) são co-originárias.
\prov{relações: parte\_de direito; relaciona contrato\_social; relaciona refeudalizacao\_esfera\_publica}
\prov{proveniência: wiki \textperiodcentered{} inferencia \textperiodcentered{} confiança 1.0 \textperiodcentered{} domínio ciencias\_de\_fora \textperiodcentered{} história 4 versões no jornal}

%CRISTAL {"arestas":[{"alvo":"fato_social","confianca":1.0,"epistemico":"fato","origem":"wiki","rel":"especializa"},{"alvo":"word","confianca":0.5,"epistemico":"fato","origem":"wiki","rel":"relaciona"},{"alvo":"vontade_de_poder","confianca":0.5,"epistemico":"fato","origem":"wiki","rel":"relaciona"},{"alvo":"virtualizacao","confianca":0.5,"epistemico":"fato","origem":"wiki","rel":"relaciona"},{"alvo":"while","confianca":0.5,"epistemico":"fato","origem":"wiki","rel":"relaciona"}],"confianca":1.0,"contraexemplos":[],"descricao":"Bourdieu: sistema de disposições incorporadas pela socialização que estrutura práticas sem consciência — reproduz o poder como natureza.","epistemico":"inferencia","exemplos":[],"id":"habitus","memoria":"semantica","meta":{"autor":"Bourdieu","dominio":"ciencias_de_fora","fonte":""},"origem":"wiki","palavras_chave":[],"sinonimos":[],"tipo":"conceito","titulo":"Habitus"}
\section{Habitus}
Bourdieu: sistema de disposições incorporadas pela socialização que estrutura práticas sem consciência — reproduz o poder como natureza.
\prov{relações: especializa fato\_social; relaciona word; relaciona vontade\_de\_poder; relaciona virtualizacao; relaciona while}
\prov{proveniência: wiki \textperiodcentered{} inferencia \textperiodcentered{} confiança 1.0 \textperiodcentered{} domínio ciencias\_de\_fora \textperiodcentered{} história 6 versões no jornal}

%CRISTAL {"arestas":[{"alvo":"schopenhauer_musica","confianca":1.0,"epistemico":"fato","origem":"wiki","rel":"contradiz"},{"alvo":"beleza_matematica","confianca":1.0,"epistemico":"fato","origem":"wiki","rel":"relaciona"}],"confianca":1.0,"contraexemplos":[],"descricao":"Hanslick: a música não representa nem evoca sentimentos; sua beleza está unicamente nas 'formas sonoras em movimento', autônomas de qualquer conteúdo extramusical. O antípoda de Schopenhauer.","epistemico":"inferencia","exemplos":[],"id":"hanslick_formalismo","memoria":"semantica","meta":{"autor":"Hanslick","dominio":"ciencias_de_fora","fonte":""},"origem":"wiki","palavras_chave":[],"sinonimos":[],"tipo":"conceito","titulo":"Formalismo Musical (Hanslick)"}
\section{Formalismo Musical (Hanslick)}
Hanslick: a música não representa nem evoca sentimentos; sua beleza está unicamente nas 'formas sonoras em movimento', autônomas de qualquer conteúdo extramusical. O antípoda de Schopenhauer.
\prov{relações: contradiz schopenhauer\_musica; relaciona beleza\_matematica}
\prov{proveniência: wiki \textperiodcentered{} inferencia \textperiodcentered{} confiança 1.0 \textperiodcentered{} domínio ciencias\_de\_fora \textperiodcentered{} história 3 versões no jornal}

%CRISTAL {"arestas":[{"alvo":"consciencia","confianca":1.0,"epistemico":"fato","origem":"wiki","rel":"parte_de"},{"alvo":"word","confianca":0.5,"epistemico":"fato","origem":"wiki","rel":"relaciona"},{"alvo":"virtualizacao","confianca":0.5,"epistemico":"fato","origem":"wiki","rel":"relaciona"}],"confianca":1.0,"contraexemplos":[],"descricao":"Chalmers (1995): por que e como processos físicos geram experiência subjetiva (qualia)? A lacuna que o reducionismo não fecha.","epistemico":"inferencia","exemplos":[],"id":"hard_problem","memoria":"semantica","meta":{"autor":"Chalmers","dominio":"filosofia","fonte":""},"origem":"wiki","palavras_chave":[],"sinonimos":[],"tipo":"conceito","titulo":"O Problema Difícil da Consciência"}
\section{O Problema Difícil da Consciência}
Chalmers (1995): por que e como processos físicos geram experiência subjetiva (qualia)? A lacuna que o reducionismo não fecha.
\prov{relações: parte\_de consciencia; relaciona word; relaciona virtualizacao}
\prov{proveniência: wiki \textperiodcentered{} inferencia \textperiodcentered{} confiança 1.0 \textperiodcentered{} domínio filosofia \textperiodcentered{} história 4 versões no jornal}

%CRISTAL {"arestas":[{"alvo":"teoria_musical","confianca":1.0,"epistemico":"fato","origem":"wiki","rel":"parte_de"},{"alvo":"serie_harmonica_musical","confianca":1.0,"epistemico":"fato","origem":"wiki","rel":"usa"},{"alvo":"pitagoras_razoes_musicais","confianca":1.0,"epistemico":"fato","origem":"wiki","rel":"usa"}],"confianca":1.0,"contraexemplos":[],"descricao":"A combinação simultânea de sons — a consonância nasce de razões simples de frequência (Pitágoras); a base vertical da música.","epistemico":"fato","exemplos":[],"id":"harmonia","memoria":"semantica","meta":{"dominio":"ciencias_de_fora","fonte":""},"origem":"wiki","palavras_chave":[],"sinonimos":[],"tipo":"conceito","titulo":"Harmonia"}
\section{Harmonia}
A combinação simultânea de sons — a consonância nasce de razões simples de frequência (Pitágoras); a base vertical da música.
\prov{relações: parte\_de teoria\_musical; usa serie\_harmonica\_musical; usa pitagoras\_razoes\_musicais}
\prov{proveniência: wiki \textperiodcentered{} fato \textperiodcentered{} confiança 1.0 \textperiodcentered{} domínio ciencias\_de\_fora \textperiodcentered{} história 5 versões no jornal}

%CRISTAL {"arestas":[{"alvo":"juspositivismo","confianca":1.0,"epistemico":"fato","origem":"wiki","rel":"especializa"},{"alvo":"kelsen_grundnorm","confianca":1.0,"epistemico":"fato","origem":"wiki","rel":"contradiz"}],"confianca":1.0,"contraexemplos":[],"descricao":"Hart: o direito é a união de regras primárias (impõem deveres) com regras secundárias (sobre as regras: reconhecimento, mudança, adjudicação). A regra de reconhecimento — critério último de validade — existe como prática social convergente das autoridades, não como pressuposição.","epistemico":"inferencia","exemplos":[],"id":"hart_regras","memoria":"semantica","meta":{"autor":"Hart","dominio":"ciencias_de_fora","fonte":""},"origem":"wiki","palavras_chave":[],"sinonimos":[],"tipo":"conceito","titulo":"Regras Primárias e Secundárias (Hart)"}
\section{Regras Primárias e Secundárias (Hart)}
Hart: o direito é a união de regras primárias (impõem deveres) com regras secundárias (sobre as regras: reconhecimento, mudança, adjudicação). A regra de reconhecimento — critério último de validade — existe como prática social convergente das autoridades, não como pressuposição.
\prov{relações: especializa juspositivismo; contradiz kelsen\_grundnorm}
\prov{proveniência: wiki \textperiodcentered{} inferencia \textperiodcentered{} confiança 1.0 \textperiodcentered{} domínio ciencias\_de\_fora \textperiodcentered{} história 3 versões no jornal}

%CRISTAL {"arestas":[{"alvo":"hart_regras","confianca":1.0,"epistemico":"fato","origem":"wiki","rel":"usa"},{"alvo":"godel_incompletude","confianca":1.0,"epistemico":"fato","origem":"wiki","rel":"relaciona"}],"confianca":1.0,"contraexemplos":[],"descricao":"Hart: toda regra tem um núcleo de sentido determinado e uma penumbra de indeterminação; nos casos-limite a linguagem não decide e o juiz exerce discricionariedade. Um limite formal da regra, como a indecidibilidade nos sistemas.","epistemico":"inferencia","exemplos":[],"id":"hart_textura_aberta","memoria":"semantica","meta":{"autor":"Hart","dominio":"ciencias_de_fora","fonte":""},"origem":"wiki","palavras_chave":[],"sinonimos":[],"tipo":"conceito","titulo":"Textura Aberta do Direito (Hart)"}
\section{Textura Aberta do Direito (Hart)}
Hart: toda regra tem um núcleo de sentido determinado e uma penumbra de indeterminação; nos casos-limite a linguagem não decide e o juiz exerce discricionariedade. Um limite formal da regra, como a indecidibilidade nos sistemas.
\prov{relações: usa hart\_regras; relaciona godel\_incompletude}
\prov{proveniência: wiki \textperiodcentered{} inferencia \textperiodcentered{} confiança 1.0 \textperiodcentered{} domínio ciencias\_de\_fora \textperiodcentered{} história 3 versões no jornal}

%CRISTAL {"arestas":[{"alvo":"geopolitica","confianca":1.0,"epistemico":"fato","origem":"wiki","rel":"usa"}],"confianca":1.0,"contraexemplos":[],"descricao":"Mackinder: o interior da Eurásia é o 'pivô' do poder mundial — 'quem governa a Europa Oriental comanda o Heartland; quem governa o Heartland comanda a Ilha-Mundo; quem governa a Ilha-Mundo comanda o mundo'. A tese do poder terrestre.","epistemico":"inferencia","exemplos":[],"id":"heartland_mackinder","memoria":"semantica","meta":{"autor":"Mackinder","dominio":"ciencias_de_fora","fonte":""},"origem":"wiki","palavras_chave":[],"sinonimos":[],"tipo":"conceito","titulo":"Heartland / Pivô Geográfico (Mackinder)"}
\section{Heartland / Pivô Geográfico (Mackinder)}
Mackinder: o interior da Eurásia é o 'pivô' do poder mundial — 'quem governa a Europa Oriental comanda o Heartland; quem governa o Heartland comanda a Ilha-Mundo; quem governa a Ilha-Mundo comanda o mundo'. A tese do poder terrestre.
\prov{relações: usa geopolitica}
\prov{proveniência: wiki \textperiodcentered{} inferencia \textperiodcentered{} confiança 1.0 \textperiodcentered{} domínio ciencias\_de\_fora \textperiodcentered{} história 2 versões no jornal}

%CRISTAL {"arestas":[{"alvo":"idealismo_berkeley","confianca":1.0,"epistemico":"fato","origem":"wiki","rel":"generaliza"},{"alvo":"word","confianca":0.5,"epistemico":"fato","origem":"wiki","rel":"relaciona"},{"alvo":"virtualizacao","confianca":0.5,"epistemico":"fato","origem":"wiki","rel":"relaciona"},{"alvo":"vida-morte","confianca":0.5,"epistemico":"fato","origem":"wiki","rel":"relaciona"},{"alvo":"misticismo_nao_dual","confianca":0.5,"epistemico":"fato","origem":"wiki","rel":"relaciona"},{"alvo":"ontologia","confianca":0.5,"epistemico":"fato","origem":"wiki","rel":"relaciona"},{"alvo":"vontade_de_poder","confianca":0.5,"epistemico":"fato","origem":"wiki","rel":"relaciona"}],"confianca":1.0,"contraexemplos":[],"descricao":"Hegel: a realidade é o Espírito (Geist) que se autorrealiza na história por dialética (tese-antítese-síntese).","epistemico":"inferencia","exemplos":[],"id":"hegel_absoluto","memoria":"semantica","meta":{"autor":"Hegel","dominio":"filosofia","fonte":""},"origem":"wiki","palavras_chave":[],"sinonimos":[],"tipo":"conceito","titulo":"Idealismo Absoluto"}
\section{Idealismo Absoluto}
Hegel: a realidade é o Espírito (Geist) que se autorrealiza na história por dialética (tese-antítese-síntese).
\prov{relações: generaliza idealismo\_berkeley; relaciona word; relaciona virtualizacao; relaciona vida-morte; relaciona misticismo\_nao\_dual; relaciona ontologia; relaciona vontade\_de\_poder}
\prov{proveniência: wiki \textperiodcentered{} inferencia \textperiodcentered{} confiança 1.0 \textperiodcentered{} domínio filosofia \textperiodcentered{} história 8 versões no jornal}

%CRISTAL {"arestas":[{"alvo":"estudos_de_midia","confianca":1.0,"epistemico":"fato","origem":"wiki","rel":"parte_de"},{"alvo":"biopoder","confianca":1.0,"epistemico":"fato","origem":"wiki","rel":"relaciona"},{"alvo":"industria_cultural","confianca":1.0,"epistemico":"fato","origem":"wiki","rel":"relaciona"}],"confianca":1.0,"contraexemplos":[],"descricao":"Gramsci: a classe dominante governa não só pela coerção, mas pelo consenso — pela sociedade civil (escola, igreja, imprensa) faz sua visão de mundo parecer senso comum e interesse geral. Dominação sempre em disputa.","epistemico":"inferencia","exemplos":[],"id":"hegemonia_cultural","memoria":"semantica","meta":{"autor":"Gramsci","dominio":"ciencias_de_fora","fonte":""},"origem":"wiki","palavras_chave":[],"sinonimos":[],"tipo":"conceito","titulo":"Hegemonia Cultural (Gramsci)"}
\section{Hegemonia Cultural (Gramsci)}
Gramsci: a classe dominante governa não só pela coerção, mas pelo consenso — pela sociedade civil (escola, igreja, imprensa) faz sua visão de mundo parecer senso comum e interesse geral. Dominação sempre em disputa.
\prov{relações: parte\_de estudos\_de\_midia; relaciona biopoder; relaciona industria\_cultural}
\prov{proveniência: wiki \textperiodcentered{} inferencia \textperiodcentered{} confiança 1.0 \textperiodcentered{} domínio ciencias\_de\_fora \textperiodcentered{} história 4 versões no jornal}

%CRISTAL {"arestas":[{"alvo":"gestell_enquadramento","confianca":1.0,"epistemico":"fato","origem":"wiki","rel":"usa"}],"confianca":1.0,"contraexemplos":[],"descricao":"Heidegger: sob o Gestell, a natureza deixa de ser objeto e vira 'reserva disponível' — estoque calculável, permanentemente à mão para requisição (o rio como fornecedor de energia, não paisagem).","epistemico":"inferencia","exemplos":[],"id":"heidegger_bestand","memoria":"semantica","meta":{"autor":"Heidegger","dominio":"ciencias_de_fora","fonte":""},"origem":"wiki","palavras_chave":[],"sinonimos":[],"tipo":"conceito","titulo":"Bestand (Reserva Disponível)"}
\section{Bestand (Reserva Disponível)}
Heidegger: sob o Gestell, a natureza deixa de ser objeto e vira 'reserva disponível' — estoque calculável, permanentemente à mão para requisição (o rio como fornecedor de energia, não paisagem).
\prov{relações: usa gestell\_enquadramento}
\prov{proveniência: wiki \textperiodcentered{} inferencia \textperiodcentered{} confiança 1.0 \textperiodcentered{} domínio ciencias\_de\_fora \textperiodcentered{} história 2 versões no jornal}

%CRISTAL {"arestas":[{"alvo":"gestell_enquadramento","confianca":1.0,"epistemico":"fato","origem":"wiki","rel":"usa"},{"alvo":"broca_os","confianca":1.0,"epistemico":"fato","origem":"wiki","rel":"relaciona"}],"confianca":1.0,"contraexemplos":[],"descricao":"Heidegger: o maior perigo do Gestell é ocultar que ele é só um modo histórico de desvelar, fechando o horizonte do sentido; mas nesse mesmo perigo cresce 'o que salva' — reconhecer a essência da técnica e reabrir outros modos de verdade.","epistemico":"inferencia","exemplos":[],"id":"heidegger_perigo_salvacao","memoria":"semantica","meta":{"autor":"Heidegger","dominio":"ciencias_de_fora","fonte":""},"origem":"wiki","palavras_chave":[],"sinonimos":[],"tipo":"conceito","titulo":"O Perigo e o que Salva (Heidegger)"}
\section{O Perigo e o que Salva (Heidegger)}
Heidegger: o maior perigo do Gestell é ocultar que ele é só um modo histórico de desvelar, fechando o horizonte do sentido; mas nesse mesmo perigo cresce 'o que salva' — reconhecer a essência da técnica e reabrir outros modos de verdade.
\prov{relações: usa gestell\_enquadramento; relaciona broca\_os}
\prov{proveniência: wiki \textperiodcentered{} inferencia \textperiodcentered{} confiança 1.0 \textperiodcentered{} domínio ciencias\_de\_fora \textperiodcentered{} história 3 versões no jornal}

%CRISTAL {"arestas":[{"alvo":"filosofia_da_tecnica","confianca":1.0,"epistemico":"fato","origem":"wiki","rel":"parte_de"},{"alvo":"acao_social_weber","confianca":1.0,"epistemico":"fato","origem":"wiki","rel":"relaciona"}],"confianca":1.0,"contraexemplos":[],"descricao":"Heidegger: a essência da técnica não é técnica nem instrumental — é um modo de fazer o ente vir à presença, uma forma de verdade como desocultamento (alétheia). A técnica moderna é um desvelamento específico, não uma ferramenta neutra.","epistemico":"inferencia","exemplos":[],"id":"heidegger_tecnica_desvelamento","memoria":"semantica","meta":{"autor":"Heidegger","dominio":"ciencias_de_fora","fonte":""},"origem":"wiki","palavras_chave":[],"sinonimos":[],"tipo":"conceito","titulo":"Técnica como Desvelamento (Heidegger)"}
\section{Técnica como Desvelamento (Heidegger)}
Heidegger: a essência da técnica não é técnica nem instrumental — é um modo de fazer o ente vir à presença, uma forma de verdade como desocultamento (alétheia). A técnica moderna é um desvelamento específico, não uma ferramenta neutra.
\prov{relações: parte\_de filosofia\_da\_tecnica; relaciona acao\_social\_weber}
\prov{proveniência: wiki \textperiodcentered{} inferencia \textperiodcentered{} confiança 1.0 \textperiodcentered{} domínio ciencias\_de\_fora \textperiodcentered{} história 3 versões no jornal}

%CRISTAL {"arestas":[{"alvo":"jogos_de_linguagem","confianca":1.0,"epistemico":"fato","origem":"wiki","rel":"usa"},{"alvo":"virtualizacao","confianca":0.5,"epistemico":"fato","origem":"wiki","rel":"relaciona"},{"alvo":"word","confianca":0.5,"epistemico":"fato","origem":"wiki","rel":"relaciona"}],"confianca":1.0,"contraexemplos":[],"descricao":"Gadamer: compreender é a fusão do horizonte do intérprete com o do texto — diálogo, não transferência; a linguagem é a casa do ser.","epistemico":"inferencia","exemplos":[],"id":"hermeneutica_gadamer","memoria":"semantica","meta":{"autor":"Gadamer","dominio":"ciencias_de_fora","fonte":""},"origem":"wiki","palavras_chave":[],"sinonimos":[],"tipo":"conceito","titulo":"Fusão de Horizontes"}
\section{Fusão de Horizontes}
Gadamer: compreender é a fusão do horizonte do intérprete com o do texto — diálogo, não transferência; a linguagem é a casa do ser.
\prov{relações: usa jogos\_de\_linguagem; relaciona virtualizacao; relaciona word}
\prov{proveniência: wiki \textperiodcentered{} inferencia \textperiodcentered{} confiança 1.0 \textperiodcentered{} domínio ciencias\_de\_fora \textperiodcentered{} história 4 versões no jornal}

%CRISTAL {"arestas":[{"alvo":"gramatica_gerativa","confianca":1.0,"epistemico":"fato","origem":"wiki","rel":"usa"},{"alvo":"maquina_de_turing","confianca":1.0,"epistemico":"fato","origem":"wiki","rel":"relaciona"},{"alvo":"expressividade_sub_turing","confianca":1.0,"epistemico":"fato","origem":"wiki","rel":"explica"},{"alvo":"linguagem_matriz_gatos","confianca":1.0,"epistemico":"fato","origem":"wiki","rel":"explica"}],"confianca":1.0,"contraexemplos":[],"descricao":"Chomsky: quatro classes de gramática por poder — regular ⊂ livre-de-contexto ⊂ sensível-ao-contexto ⊂ recursivamente enumerável (Turing).","epistemico":"fato","exemplos":[],"id":"hierarquia_de_chomsky","memoria":"semantica","meta":{"autor":"Chomsky","dominio":"ciencias_de_fora","fonte":""},"origem":"wiki","palavras_chave":[],"sinonimos":[],"tipo":"conceito","titulo":"Hierarquia de Chomsky"}
\section{Hierarquia de Chomsky}
Chomsky: quatro classes de gramática por poder — regular \ensuremath{\subset} livre-de-contexto \ensuremath{\subset} sensível-ao-contexto \ensuremath{\subset} recursivamente enumerável (Turing).
\prov{relações: usa gramatica\_gerativa; relaciona maquina\_de\_turing; explica expressividade\_sub\_turing; explica linguagem\_matriz\_gatos}
\prov{proveniência: wiki \textperiodcentered{} fato \textperiodcentered{} confiança 1.0 \textperiodcentered{} domínio ciencias\_de\_fora \textperiodcentered{} história 5 versões no jornal}

%CRISTAL {"arestas":[{"alvo":"programa_de_hilbert","confianca":1.0,"epistemico":"fato","origem":"wiki","rel":"usa"},{"alvo":"hipotese_de_riemann","confianca":1.0,"epistemico":"fato","origem":"wiki","rel":"referencia"}],"confianca":1.0,"contraexemplos":[],"descricao":"Hilbert (1900): a agenda que orientou o século — do contínuo à consistência. Definir o problema certo é meio caminho.","epistemico":"fato","exemplos":[],"id":"hilbert_problemas","memoria":"semantica","meta":{"autor":"Hilbert","dominio":"ciencias_de_fora","fonte":""},"origem":"wiki","palavras_chave":[],"sinonimos":[],"tipo":"conceito","titulo":"Os 23 Problemas de Hilbert"}
\section{Os 23 Problemas de Hilbert}
Hilbert (1900): a agenda que orientou o século — do contínuo à consistência. Definir o problema certo é meio caminho.
\prov{relações: usa programa\_de\_hilbert; referencia hipotese\_de\_riemann}
\prov{proveniência: wiki \textperiodcentered{} fato \textperiodcentered{} confiança 1.0 \textperiodcentered{} domínio ciencias\_de\_fora \textperiodcentered{} história 6 versões no jornal}

%CRISTAL {"arestas":[{"alvo":"sagrado","confianca":1.0,"epistemico":"fato","origem":"wiki","rel":"parte_de"},{"alvo":"vedanta","confianca":1.0,"epistemico":"fato","origem":"wiki","rel":"relaciona"},{"alvo":"word","confianca":0.5,"epistemico":"fato","origem":"wiki","rel":"relaciona"}],"confianca":1.0,"contraexemplos":[],"descricao":"A tradição da Bhagavad Gita — os caminhos (yogas) para o Divino: a ação desapegada, a devoção e o conhecimento, sob a doutrina do Atman imortal e do dharma.","epistemico":"fato","exemplos":[],"id":"hinduismo","memoria":"semantica","meta":{"dominio":"sagrado","fonte":""},"origem":"wiki","palavras_chave":[],"sinonimos":[],"tipo":"conceito","titulo":"Hinduísmo (Bhagavad Gita)"}
\section{Hinduísmo (Bhagavad Gita)}
A tradição da Bhagavad Gita — os caminhos (yogas) para o Divino: a ação desapegada, a devoção e o conhecimento, sob a doutrina do Atman imortal e do dharma.
\prov{relações: parte\_de sagrado; relaciona vedanta; relaciona word}
\prov{proveniência: wiki \textperiodcentered{} fato \textperiodcentered{} confiança 1.0 \textperiodcentered{} domínio sagrado \textperiodcentered{} história 4 versões no jornal}

%CRISTAL {"arestas":[{"alvo":"teorema_dos_numeros_primos","confianca":1.0,"epistemico":"fato","origem":"wiki","rel":"explica"},{"alvo":"teoria_dos_numeros","confianca":1.0,"epistemico":"fato","origem":"wiki","rel":"parte_de"}],"confianca":1.0,"contraexemplos":[],"descricao":"Riemann: os zeros não-triviais da função zeta têm parte real 1/2 — controlaria o erro na contagem de primos. Aberta.","epistemico":"hipotese","exemplos":[],"id":"hipotese_de_riemann","memoria":"semantica","meta":{"autor":"Riemann","dominio":"ciencias_de_fora","fonte":""},"origem":"wiki","palavras_chave":[],"sinonimos":[],"tipo":"conceito","titulo":"Hipótese de Riemann"}
\section{Hipótese de Riemann}
Riemann: os zeros não-triviais da função zeta têm parte real 1/2 — controlaria o erro na contagem de primos. Aberta.
\prov{relações: explica teorema\_dos\_numeros\_primos; parte\_de teoria\_dos\_numeros}
\prov{proveniência: wiki \textperiodcentered{} hipotese \textperiodcentered{} confiança 1.0 \textperiodcentered{} domínio ciencias\_de\_fora \textperiodcentered{} história 6 versões no jornal}

%CRISTAL {"arestas":[{"alvo":"cantor_infinitos","confianca":1.0,"epistemico":"fato","origem":"wiki","rel":"usa"},{"alvo":"incompletude_godel","confianca":1.0,"epistemico":"fato","origem":"wiki","rel":"relaciona"}],"confianca":1.0,"contraexemplos":[],"descricao":"Cantor perguntou; Gödel e Cohen responderam: é INDEPENDENTE de ZFC — nem provável, nem refutável nos axiomas usuais.","epistemico":"inferencia","exemplos":[],"id":"hipotese_do_continuo","memoria":"semantica","meta":{"autor":"Cantor/Cohen","dominio":"ciencias_de_fora","fonte":""},"origem":"wiki","palavras_chave":[],"sinonimos":[],"tipo":"conceito","titulo":"Hipótese do Contínuo"}
\section{Hipótese do Contínuo}
Cantor perguntou; Gödel e Cohen responderam: é INDEPENDENTE de ZFC — nem provável, nem refutável nos axiomas usuais.
\prov{relações: usa cantor\_infinitos; relaciona incompletude\_godel}
\prov{proveniência: wiki \textperiodcentered{} inferencia \textperiodcentered{} confiança 1.0 \textperiodcentered{} domínio ciencias\_de\_fora \textperiodcentered{} história 6 versões no jornal}

%CRISTAL {"arestas":[{"alvo":"homeostase","confianca":1.0,"epistemico":"fato","origem":"wiki","rel":"generaliza"},{"alvo":"ecologia","confianca":1.0,"epistemico":"fato","origem":"wiki","rel":"parte_de"},{"alvo":"gaia","confianca":1.0,"epistemico":"fato","origem":"wiki","rel":"relaciona"}],"confianca":1.0,"contraexemplos":[],"descricao":"Lovelock: a biosfera regula o planeta (clima, atmosfera) como um sistema quase-homeostático. Versão forte é contestada.","epistemico":"hipotese","exemplos":[],"id":"hipotese_gaia","memoria":"semantica","meta":{"autor":"Lovelock","dominio":"ciencias_de_fora","fonte":""},"origem":"wiki","palavras_chave":[],"sinonimos":[],"tipo":"conceito","titulo":"Hipótese Gaia"}
\section{Hipótese Gaia}
Lovelock: a biosfera regula o planeta (clima, atmosfera) como um sistema quase-homeostático. Versão forte é contestada.
\prov{relações: generaliza homeostase; parte\_de ecologia; relaciona gaia}
\prov{proveniência: wiki \textperiodcentered{} hipotese \textperiodcentered{} confiança 1.0 \textperiodcentered{} domínio ciencias\_de\_fora \textperiodcentered{} história 4 versões no jornal}

%CRISTAL {"arestas":[{"alvo":"filosofia_da_linguagem","confianca":1.0,"epistemico":"fato","origem":"wiki","rel":"relaciona"},{"alvo":"linguistica","confianca":1.0,"epistemico":"fato","origem":"wiki","rel":"parte_de"}],"confianca":1.0,"contraexemplos":[],"descricao":"A língua molda o pensamento e a percepção. Versão forte (determinismo) é hipótese contestada; a fraca tem suporte.","epistemico":"hipotese","exemplos":[],"id":"hipotese_sapir_whorf","memoria":"semantica","meta":{"autor":"Sapir/Whorf","dominio":"ciencias_de_fora","fonte":""},"origem":"wiki","palavras_chave":[],"sinonimos":[],"tipo":"conceito","titulo":"Hipótese Sapir-Whorf"}
\section{Hipótese Sapir-Whorf}
A língua molda o pensamento e a percepção. Versão forte (determinismo) é hipótese contestada; a fraca tem suporte.
\prov{relações: relaciona filosofia\_da\_linguagem; parte\_de linguistica}
\prov{proveniência: wiki \textperiodcentered{} hipotese \textperiodcentered{} confiança 1.0 \textperiodcentered{} domínio ciencias\_de\_fora \textperiodcentered{} história 3 versões no jornal}

%CRISTAL {"arestas":[{"alvo":"topologia","confianca":0.5,"epistemico":"fato","origem":"wiki","rel":"relaciona"},{"alvo":"kernel","confianca":0.5,"epistemico":"fato","origem":"wiki","rel":"relaciona"},{"alvo":"imagem_do_floco_---_validacao_no","confianca":0.5,"epistemico":"fato","origem":"wiki","rel":"relaciona"},{"alvo":"projeto_cristaljson","confianca":0.5,"epistemico":"fato","origem":"wiki","rel":"relaciona"},{"alvo":"regua_associativa_e_o_risco_de_vazamento","confianca":0.5,"epistemico":"fato","origem":"wiki","rel":"relaciona"},{"alvo":"politica","confianca":0.5,"epistemico":"fato","origem":"wiki","rel":"relaciona"},{"alvo":"java","confianca":0.5,"epistemico":"fato","origem":"wiki","rel":"relaciona"}],"confianca":1.0,"contraexemplos":[],"descricao":"O estudo do passado humano e de seu sentido — os eventos e as estruturas que os movem.","epistemico":"fato","exemplos":[],"id":"historia","memoria":"semantica","meta":{"dominio":"ciencias_de_fora","fonte":""},"origem":"wiki","palavras_chave":[],"sinonimos":[],"tipo":"conceito","titulo":"História"}
\section{História}
O estudo do passado humano e de seu sentido — os eventos e as estruturas que os movem.
\prov{relações: relaciona topologia; relaciona kernel; relaciona imagem\_do\_floco\_---\_validacao\_no; relaciona projeto\_cristaljson; relaciona regua\_associativa\_e\_o\_risco\_de\_vazamento; relaciona politica; relaciona java}
\prov{proveniência: wiki \textperiodcentered{} fato \textperiodcentered{} confiança 1.0 \textperiodcentered{} domínio ciencias\_de\_fora \textperiodcentered{} história 8 versões no jornal}

%CRISTAL {"arestas":[{"alvo":"dialetica_hegeliana","confianca":1.0,"epistemico":"fato","origem":"wiki","rel":"contradiz"},{"alvo":"materialismo_historico","confianca":1.0,"epistemico":"fato","origem":"wiki","rel":"contradiz"},{"alvo":"fim_da_historia","confianca":1.0,"epistemico":"fato","origem":"wiki","rel":"contradiz"},{"alvo":"contingencia_historica","confianca":1.0,"epistemico":"fato","origem":"wiki","rel":"usa"},{"alvo":"epistemologia","confianca":1.0,"epistemico":"fato","origem":"wiki","rel":"parte_de"},{"alvo":"psicanalise","confianca":1.0,"epistemico":"fato","origem":"wiki","rel":"contradiz"}],"confianca":1.0,"contraexemplos":[],"descricao":"Popper: crer em leis da história que permitem prever o futuro (Hegel, Marx, Spengler) é pseudociência — não-falsável. A história é contingente, não legislada.","epistemico":"fato","exemplos":[],"id":"historicismo_popper","memoria":"semantica","meta":{"autor":"Popper","dominio":"ciencias_de_fora","fonte":""},"origem":"wiki","palavras_chave":[],"sinonimos":[],"tipo":"conceito","titulo":"Crítica do Historicismo (Popper)"}
\section{Crítica do Historicismo (Popper)}
Popper: crer em leis da história que permitem prever o futuro (Hegel, Marx, Spengler) é pseudociência — não-falsável. A história é contingente, não legislada.
\prov{relações: contradiz dialetica\_hegeliana; contradiz materialismo\_historico; contradiz fim\_da\_historia; usa contingencia\_historica; parte\_de epistemologia; contradiz psicanalise}
\prov{proveniência: wiki \textperiodcentered{} fato \textperiodcentered{} confiança 1.0 \textperiodcentered{} domínio ciencias\_de\_fora \textperiodcentered{} história 7 versões no jornal}

%CRISTAL {"arestas":[{"alvo":"filosofia_da_medicina","confianca":1.0,"epistemico":"fato","origem":"wiki","rel":"parte_de"},{"alvo":"modelo_biomedico","confianca":1.0,"epistemico":"fato","origem":"wiki","rel":"relaciona"},{"alvo":"autopoiese","confianca":1.0,"epistemico":"fato","origem":"wiki","rel":"relaciona"}],"confianca":1.0,"contraexemplos":[],"descricao":"O eixo epistemológico: o reducionismo explica a doença decompondo o organismo em mecanismos elementares; o holismo sustenta que o todo do vivente e seu contexto têm propriedades irredutíveis às partes. Estrutura a oposição biomédico/biopsicossocial.","epistemico":"inferencia","exemplos":[],"id":"holismo_reducionismo_medicina","memoria":"semantica","meta":{"autor":"von Bertalanffy","dominio":"ciencias_de_fora","fonte":""},"origem":"wiki","palavras_chave":[],"sinonimos":[],"tipo":"conceito","titulo":"Holismo vs Reducionismo na Medicina"}
\section{Holismo vs Reducionismo na Medicina}
O eixo epistemológico: o reducionismo explica a doença decompondo o organismo em mecanismos elementares; o holismo sustenta que o todo do vivente e seu contexto têm propriedades irredutíveis às partes. Estrutura a oposição biomédico/biopsicossocial.
\prov{relações: parte\_de filosofia\_da\_medicina; relaciona modelo\_biomedico; relaciona autopoiese}
\prov{proveniência: wiki \textperiodcentered{} inferencia \textperiodcentered{} confiança 1.0 \textperiodcentered{} domínio ciencias\_de\_fora \textperiodcentered{} história 4 versões no jornal}

%CRISTAL {"arestas":[{"alvo":"filosofia_da_educacao","confianca":1.0,"epistemico":"fato","origem":"wiki","rel":"parte_de"},{"alvo":"idiotismo_nao_examinar","confianca":1.0,"epistemico":"fato","origem":"wiki","rel":"relaciona"},{"alvo":"imperativo_categorico","confianca":1.0,"epistemico":"fato","origem":"wiki","rel":"relaciona"},{"alvo":"educacao_natural_rousseau","confianca":1.0,"epistemico":"fato","origem":"wiki","rel":"contradiz"}],"confianca":1.0,"contraexemplos":[],"descricao":"Kant (Sobre a Pedagogia): o humano é a única criatura que precisa ser educada, pois não age por instinto pronto; pela disciplina e o cultivo moral a educação o conduz à autonomia — ao uso do próprio entendimento sem tutela.","epistemico":"inferencia","exemplos":[],"id":"homem_educavel_kant","memoria":"semantica","meta":{"autor":"Kant","dominio":"ciencias_de_fora","fonte":""},"origem":"wiki","palavras_chave":[],"sinonimos":[],"tipo":"conceito","titulo":"O Homem como Criatura Educável (Kant)"}
\section{O Homem como Criatura Educável (Kant)}
Kant (Sobre a Pedagogia): o humano é a única criatura que precisa ser educada, pois não age por instinto pronto; pela disciplina e o cultivo moral a educação o conduz à autonomia — ao uso do próprio entendimento sem tutela.
\prov{relações: parte\_de filosofia\_da\_educacao; relaciona idiotismo\_nao\_examinar; relaciona imperativo\_categorico; contradiz educacao\_natural\_rousseau}
\prov{proveniência: wiki \textperiodcentered{} inferencia \textperiodcentered{} confiança 1.0 \textperiodcentered{} domínio ciencias\_de\_fora \textperiodcentered{} história 5 versões no jornal}

%CRISTAL {"arestas":[{"alvo":"o_que_e_vida","confianca":1.0,"epistemico":"fato","origem":"wiki","rel":"usa"},{"alvo":"metabolismo","confianca":1.0,"epistemico":"fato","origem":"wiki","rel":"relaciona"},{"alvo":"piso_parabola_dinamico","confianca":1.0,"epistemico":"fato","origem":"wiki","rel":"relaciona"}],"confianca":1.0,"contraexemplos":[],"descricao":"Bernard/Cannon: o organismo mantém seu meio interno estável (temperatura, pH) contra o fora — o equilíbrio dinâmico do vivo.","epistemico":"fato","exemplos":[],"id":"homeostase","memoria":"semantica","meta":{"autor":"Cannon","dominio":"ciencias_de_fora","fonte":""},"origem":"wiki","palavras_chave":[],"sinonimos":[],"tipo":"conceito","titulo":"Homeostase"}
\section{Homeostase}
Bernard/Cannon: o organismo mantém seu meio interno estável (temperatura, pH) contra o fora — o equilíbrio dinâmico do vivo.
\prov{relações: usa o\_que\_e\_vida; relaciona metabolismo; relaciona piso\_parabola\_dinamico}
\prov{proveniência: wiki \textperiodcentered{} fato \textperiodcentered{} confiança 1.0 \textperiodcentered{} domínio ciencias\_de\_fora \textperiodcentered{} história 6 versões no jornal}

%CRISTAL {"arestas":[{"alvo":"homeostase","confianca":1.0,"epistemico":"fato","origem":"wiki","rel":"explica"},{"alvo":"realimentacao_lyapunov","confianca":1.0,"epistemico":"fato","origem":"wiki","rel":"eh_um"},{"alvo":"malha_de_controle_laco","confianca":1.0,"epistemico":"fato","origem":"wiki","rel":"relaciona"},{"alvo":"corpo_mineral","confianca":1.0,"epistemico":"fato","origem":"wiki","rel":"usa"}],"confianca":1.0,"contraexemplos":[],"descricao":"O corpo regula temperatura, glicose e pH a um ponto fixo (o ATRATOR da saúde) por realimentação negativa — a mesma malha de controle do robô e da rede elétrica. A doença aguda é a órbita saindo do atrator; curar é o controle trazê-la de volta. A febre e a hiperglicemia são o erro que o controle luta para zerar.","epistemico":"inferencia","exemplos":[],"id":"homeostase_controle","memoria":"semantica","meta":{"dominio":"medicina","fonte":"medicina e diagnóstico"},"origem":"wiki","palavras_chave":[],"sinonimos":[],"tipo":"conceito","titulo":"Homeostase = Controle (o Lyapunov da saúde)"}
\section{Homeostase = Controle (o Lyapunov da saúde)}
O corpo regula temperatura, glicose e pH a um ponto fixo (o ATRATOR da saúde) por realimentação negativa — a mesma malha de controle do robô e da rede elétrica. A doença aguda é a órbita saindo do atrator; curar é o controle trazê-la de volta. A febre e a hiperglicemia são o erro que o controle luta para zerar.
\prov{relações: explica homeostase; eh\_um realimentacao\_lyapunov; relaciona malha\_de\_controle\_laco; usa corpo\_mineral}
\prov{proveniência: wiki \textperiodcentered{} medicina e diagnóstico \textperiodcentered{} inferencia \textperiodcentered{} confiança 1.0 \textperiodcentered{} domínio medicina \textperiodcentered{} história 5 versões no jornal}

%CRISTAL {"arestas":[{"alvo":"celula_mineral_hub","confianca":1.0,"epistemico":"fato","origem":"wiki","rel":"parte_de"},{"alvo":"homeostase","confianca":1.0,"epistemico":"fato","origem":"wiki","rel":"eh_um"},{"alvo":"numeros_de_pisot","confianca":1.0,"epistemico":"fato","origem":"wiki","rel":"usa"}],"confianca":1.0,"contraexemplos":[],"descricao":"O estado da célula é um ponto fixo CONTRÁTIL: perturbações decaem exponencialmente e ela retorna ao equilíbrio — exatamente o regime CRISTAL/Pisot da parábola (raio espectral ρ<1, λ<0). Na BORDA (ρ=1) a célula é excitável/oscila; no CAOS (ρ>1) foge do equilíbrio — proliferação descontrolada ou morte. A homeostase é a estabilidade da dinâmica, e a doença é a perda dela. Verif.: célula estável volta, instável foge.","epistemico":"fato","exemplos":[],"id":"homeostase_e_o_cristal","memoria":"semantica","meta":{"dominio":"biologia","fonte":"célula mineral"},"origem":"wiki","palavras_chave":[],"sinonimos":[],"tipo":"conceito","titulo":"Homeostase é o Cristal (Pisot)"}
\section{Homeostase é o Cristal (Pisot)}
O estado da célula é um ponto fixo CONTRÁTIL: perturbações decaem exponencialmente e ela retorna ao equilíbrio — exatamente o regime CRISTAL/Pisot da parábola (raio espectral \ensuremath{\rho}<1, \ensuremath{\lambda}<0). Na BORDA (\ensuremath{\rho}=1) a célula é excitável/oscila; no CAOS (\ensuremath{\rho}>1) foge do equilíbrio — proliferação descontrolada ou morte. A homeostase é a estabilidade da dinâmica, e a doença é a perda dela. Verif.: célula estável volta, instável foge.
\prov{relações: parte\_de celula\_mineral\_hub; eh\_um homeostase; usa numeros\_de\_pisot}
\prov{proveniência: wiki \textperiodcentered{} célula mineral \textperiodcentered{} fato \textperiodcentered{} confiança 1.0 \textperiodcentered{} domínio biologia \textperiodcentered{} história 4 versões no jornal}

%CRISTAL {"arestas":[{"alvo":"etica","confianca":1.0,"epistemico":"fato","origem":"wiki","rel":"parte_de"},{"alvo":"epistemologia","confianca":1.0,"epistemico":"fato","origem":"wiki","rel":"parte_de"},{"alvo":"virtualizacao","confianca":0.5,"epistemico":"fato","origem":"wiki","rel":"relaciona"},{"alvo":"word","confianca":0.5,"epistemico":"fato","origem":"wiki","rel":"relaciona"},{"alvo":"vida-morte","confianca":0.5,"epistemico":"fato","origem":"wiki","rel":"relaciona"},{"alvo":"vontade_de_poder","confianca":0.5,"epistemico":"fato","origem":"wiki","rel":"relaciona"}],"confianca":1.0,"contraexemplos":[],"descricao":"Nietzsche: a virtude epistêmica suprema — encarar verdades desconfortáveis e recusar o autoengano.","epistemico":"inferencia","exemplos":[],"id":"honestidade_intelectual","memoria":"semantica","meta":{"autor":"Nietzsche","dominio":"ciencias_de_fora","fonte":""},"origem":"wiki","palavras_chave":[],"sinonimos":[],"tipo":"conceito","titulo":"Honestidade Intelectual"}
\section{Honestidade Intelectual}
Nietzsche: a virtude epistêmica suprema — encarar verdades desconfortáveis e recusar o autoengano.
\prov{relações: parte\_de etica; parte\_de epistemologia; relaciona virtualizacao; relaciona word; relaciona vida-morte; relaciona vontade\_de\_poder}
\prov{proveniência: wiki \textperiodcentered{} inferencia \textperiodcentered{} confiança 1.0 \textperiodcentered{} domínio ciencias\_de\_fora \textperiodcentered{} história 7 versões no jornal}

%CRISTAL {"arestas":[{"alvo":"idealismo_berkeley","confianca":1.0,"epistemico":"fato","origem":"wiki","rel":"especializa"},{"alvo":"troe_postulado_aureo","confianca":1.0,"epistemico":"fato","origem":"wiki","rel":"usa"}],"confianca":1.0,"contraexemplos":[],"descricao":"Lopes (via Nisargadatta): cada consciência cria seu mundo; matemática, física, religião são 'mundos pintados na tela da Consciência' — idealismo.","epistemico":"hipotese","exemplos":[],"id":"humanidade_sonha","memoria":"semantica","meta":{"autor":"Gentil Lopes","dominio":"filosofia","fonte":""},"origem":"livro","palavras_chave":[],"sinonimos":[],"tipo":"conceito","titulo":"A Humanidade Sonha"}
\section{A Humanidade Sonha}
Lopes (via Nisargadatta): cada consciência cria seu mundo; matemática, física, religião são 'mundos pintados na tela da Consciência' — idealismo.
\prov{relações: especializa idealismo\_berkeley; usa troe\_postulado\_aureo}
\prov{proveniência: livro \textperiodcentered{} hipotese \textperiodcentered{} confiança 1.0 \textperiodcentered{} domínio filosofia \textperiodcentered{} história 4 versões no jornal}

%CRISTAL {"arestas":[{"alvo":"filosofia_da_medicina","confianca":1.0,"epistemico":"fato","origem":"wiki","rel":"parte_de"},{"alvo":"feedback_retroalimentacao","confianca":1.0,"epistemico":"fato","origem":"wiki","rel":"relaciona"}],"confianca":1.0,"contraexemplos":[],"descricao":"Hipócrates: saúde é o equilíbrio (eukrasia) dos quatro humores — sangue, fleuma, bile amarela e bile negra; a doença surge do excesso ou má mistura. Marca a passagem da explicação sobrenatural para causas naturais, hoje superada.","epistemico":"inferencia","exemplos":[],"id":"humoralismo_hipocrates","memoria":"semantica","meta":{"autor":"Hipócrates","dominio":"ciencias_de_fora","fonte":""},"origem":"wiki","palavras_chave":[],"sinonimos":[],"tipo":"conceito","titulo":"Teoria dos Humores (Hipócrates)"}
\section{Teoria dos Humores (Hipócrates)}
Hipócrates: saúde é o equilíbrio (eukrasia) dos quatro humores — sangue, fleuma, bile amarela e bile negra; a doença surge do excesso ou má mistura. Marca a passagem da explicação sobrenatural para causas naturais, hoje superada.
\prov{relações: parte\_de filosofia\_da\_medicina; relaciona feedback\_retroalimentacao}
\prov{proveniência: wiki \textperiodcentered{} inferencia \textperiodcentered{} confiança 1.0 \textperiodcentered{} domínio ciencias\_de\_fora \textperiodcentered{} história 3 versões no jornal}

%CRISTAL {"arestas":[{"alvo":"signo_triadico","confianca":1.0,"epistemico":"fato","origem":"wiki","rel":"parte_de"},{"alvo":"simbolo_vs_significado","confianca":1.0,"epistemico":"fato","origem":"wiki","rel":"explica"}],"confianca":1.0,"contraexemplos":[],"descricao":"Peirce: três modos do signo referir — ícone (semelhança), índice (causa/contiguidade), símbolo (convenção).","epistemico":"fato","exemplos":[],"id":"icone_indice_simbolo","memoria":"semantica","meta":{"autor":"Peirce","dominio":"ciencias_de_fora","fonte":""},"origem":"wiki","palavras_chave":[],"sinonimos":[],"tipo":"conceito","titulo":"Ícone, Índice e Símbolo"}
\section{Ícone, Índice e Símbolo}
Peirce: três modos do signo referir — ícone (semelhança), índice (causa/contiguidade), símbolo (convenção).
\prov{relações: parte\_de signo\_triadico; explica simbolo\_vs\_significado}
\prov{proveniência: wiki \textperiodcentered{} fato \textperiodcentered{} confiança 1.0 \textperiodcentered{} domínio ciencias\_de\_fora \textperiodcentered{} história 3 versões no jornal}

%CRISTAL {"arestas":[{"alvo":"metafisica","confianca":1.0,"epistemico":"fato","origem":"wiki","rel":"parte_de"},{"alvo":"colapso_observador","confianca":1.0,"epistemico":"fato","origem":"wiki","rel":"relaciona"},{"alvo":"word","confianca":0.5,"epistemico":"fato","origem":"wiki","rel":"relaciona"}],"confianca":1.0,"contraexemplos":[],"descricao":"Berkeley: ser é ser percebido; a realidade é mental. Prefigura o observador que define a realidade na física quântica.","epistemico":"inferencia","exemplos":[],"id":"idealismo_berkeley","memoria":"semantica","meta":{"autor":"Berkeley","dominio":"filosofia","fonte":""},"origem":"wiki","palavras_chave":[],"sinonimos":[],"tipo":"conceito","titulo":"Idealismo (esse est percipi)"}
\section{Idealismo (esse est percipi)}
Berkeley: ser é ser percebido; a realidade é mental. Prefigura o observador que define a realidade na física quântica.
\prov{relações: parte\_de metafisica; relaciona colapso\_observador; relaciona word}
\prov{proveniência: wiki \textperiodcentered{} inferencia \textperiodcentered{} confiança 1.0 \textperiodcentered{} domínio filosofia \textperiodcentered{} história 4 versões no jornal}

%CRISTAL {"arestas":[{"alvo":"consciencia_software_universo_gentil","confianca":1.0,"epistemico":"fato","origem":"wiki","rel":"usa"},{"alvo":"mente_programacao_cerebro_computador","confianca":1.0,"epistemico":"fato","origem":"wiki","rel":"usa"},{"alvo":"ressurreicao","confianca":1.0,"epistemico":"fato","origem":"wiki","rel":"contradiz"},{"alvo":"face_bela_morte","confianca":1.0,"epistemico":"fato","origem":"wiki","rel":"relaciona"}],"confianca":1.0,"contraexemplos":[],"descricao":"Lopes: o ego/identidade é 'um arquivo codificado quimicamente' que não sobrevive à morte (Crick: o eu é interação de células nervosas); mas o código/informação não se perde ('identidade = código + hardware'; some o hardware, resta o código), armazenado num 'campo akáshico'/'disco rígido cósmico' (analogia com a informação preservada na superfície do buraco negro). Não há vida após a morte do ego, mas persistência da informação.","epistemico":"inferencia","exemplos":[],"id":"identidade_codigo_hardware_gentil","memoria":"semantica","meta":{"autor":"Gentil Lopes","dominio":"sagrado","fonte":""},"origem":"wiki","palavras_chave":[],"sinonimos":[],"tipo":"conceito","titulo":"Vida após a Morte = Persistência da Informação (Lopes)"}
\section{Vida após a Morte = Persistência da Informação (Lopes)}
Lopes: o ego/identidade é 'um arquivo codificado quimicamente' que não sobrevive à morte (Crick: o eu é interação de células nervosas); mas o código/informação não se perde ('identidade = código + hardware'; some o hardware, resta o código), armazenado num 'campo akáshico'/'disco rígido cósmico' (analogia com a informação preservada na superfície do buraco negro). Não há vida após a morte do ego, mas persistência da informação.
\prov{relações: usa consciencia\_software\_universo\_gentil; usa mente\_programacao\_cerebro\_computador; contradiz ressurreicao; relaciona face\_bela\_morte}
\prov{proveniência: wiki \textperiodcentered{} inferencia \textperiodcentered{} confiança 1.0 \textperiodcentered{} domínio sagrado \textperiodcentered{} história 6 versões no jornal}

%CRISTAL {"arestas":[{"alvo":"ciencia_politica","confianca":1.0,"epistemico":"fato","origem":"wiki","rel":"relaciona"},{"alvo":"liberdade_berlin","confianca":1.0,"epistemico":"fato","origem":"wiki","rel":"usa"},{"alvo":"deus_mentiroso","confianca":1.0,"epistemico":"fato","origem":"wiki","rel":"relaciona"},{"alvo":"religiao_opio","confianca":1.0,"epistemico":"fato","origem":"wiki","rel":"relaciona"}],"confianca":1.0,"contraexemplos":[],"descricao":"Lopes, endossando Osho: 'se você tem alguma ideologia para se definir, você não é livre — liberdade significa nenhuma definição'. Conta que abandonou 'as fileiras religiosas' aos 26 anos para poder pesquisar 'sem a obrigação a priori de concluir isto ou aquilo'.","epistemico":"inferencia","exemplos":[],"id":"ideologia_anula_liberdade","memoria":"semantica","meta":{"autor":"Gentil Lopes","dominio":"ciencias_de_fora","fonte":""},"origem":"wiki","palavras_chave":[],"sinonimos":[],"tipo":"conceito","titulo":"Toda Ideologia Anula a Liberdade (Lopes)"}
\section{Toda Ideologia Anula a Liberdade (Lopes)}
Lopes, endossando Osho: 'se você tem alguma ideologia para se definir, você não é livre — liberdade significa nenhuma definição'. Conta que abandonou 'as fileiras religiosas' aos 26 anos para poder pesquisar 'sem a obrigação a priori de concluir isto ou aquilo'.
\prov{relações: relaciona ciencia\_politica; usa liberdade\_berlin; relaciona deus\_mentiroso; relaciona religiao\_opio}
\prov{proveniência: wiki \textperiodcentered{} inferencia \textperiodcentered{} confiança 1.0 \textperiodcentered{} domínio ciencias\_de\_fora \textperiodcentered{} história 6 versões no jornal}

%CRISTAL {"arestas":[{"alvo":"linguistica","confianca":1.0,"epistemico":"fato","origem":"wiki","rel":"parte_de"},{"alvo":"lingua_e_fala","confianca":1.0,"epistemico":"fato","origem":"wiki","rel":"relaciona"}],"confianca":1.0,"contraexemplos":[],"descricao":"Um sistema de comunicação humano completo, falado por uma comunidade — som, gramática e léxico que se aprende na infância. ~7000 no mundo.","epistemico":"fato","exemplos":[],"id":"idioma","memoria":"semantica","meta":{"dominio":"ciencias_de_fora","fonte":""},"origem":"wiki","palavras_chave":[],"sinonimos":[],"tipo":"conceito","titulo":"Idioma (Língua Natural)"}
\section{Idioma (Língua Natural)}
Um sistema de comunicação humano completo, falado por uma comunidade — som, gramática e léxico que se aprende na infância. \~{}7000 no mundo.
\prov{relações: parte\_de linguistica; relaciona lingua\_e\_fala}
\prov{proveniência: wiki \textperiodcentered{} fato \textperiodcentered{} confiança 1.0 \textperiodcentered{} domínio ciencias\_de\_fora \textperiodcentered{} história 3 versões no jornal}

%CRISTAL {"arestas":[{"alvo":"honestidade_intelectual","confianca":1.0,"epistemico":"fato","origem":"wiki","rel":"contradiz"},{"alvo":"deus_postulado","confianca":1.0,"epistemico":"fato","origem":"wiki","rel":"relaciona"},{"alvo":"word","confianca":0.5,"epistemico":"fato","origem":"wiki","rel":"relaciona"}],"confianca":1.0,"contraexemplos":[],"descricao":"Lopes: 'idiota' não é quem ignora, mas quem RECUSA examinar — covardia epistêmica que conserva o ego.","epistemico":"inferencia","exemplos":[],"id":"idiotismo_nao_examinar","memoria":"semantica","meta":{"autor":"Gentil Lopes","dominio":"ciencias_de_fora","fonte":""},"origem":"livro","palavras_chave":[],"sinonimos":[],"tipo":"conceito","titulo":"Idiotismo (não-examinar)"}
\section{Idiotismo (não-examinar)}
Lopes: 'idiota' não é quem ignora, mas quem RECUSA examinar — covardia epistêmica que conserva o ego.
\prov{relações: contradiz honestidade\_intelectual; relaciona deus\_postulado; relaciona word}
\prov{proveniência: livro \textperiodcentered{} inferencia \textperiodcentered{} confiança 1.0 \textperiodcentered{} domínio ciencias\_de\_fora \textperiodcentered{} história 5 versões no jornal}

%CRISTAL {"arestas":[{"alvo":"mestre_ignorante_ranciere","confianca":1.0,"epistemico":"fato","origem":"wiki","rel":"usa"},{"alvo":"capital_cultural","confianca":1.0,"epistemico":"fato","origem":"wiki","rel":"contradiz"},{"alvo":"habitus","confianca":1.0,"epistemico":"fato","origem":"wiki","rel":"relaciona"}],"confianca":1.0,"contraexemplos":[],"descricao":"Rancière: todas as inteligências são iguais em capacidade — a igualdade é ponto de partida (axioma) a verificar, não meta a alcançar. A explicação hierárquica produz a desigualdade que alega corrigir. Pressuposto declarado, não fato.","epistemico":"hipotese","exemplos":[],"id":"igualdade_das_inteligencias","memoria":"semantica","meta":{"autor":"Rancière","dominio":"ciencias_de_fora","fonte":""},"origem":"wiki","palavras_chave":[],"sinonimos":[],"tipo":"conceito","titulo":"Igualdade das Inteligências"}
\section{Igualdade das Inteligências}
Rancière: todas as inteligências são iguais em capacidade — a igualdade é ponto de partida (axioma) a verificar, não meta a alcançar. A explicação hierárquica produz a desigualdade que alega corrigir. Pressuposto declarado, não fato.
\prov{relações: usa mestre\_ignorante\_ranciere; contradiz capital\_cultural; relaciona habitus}
\prov{proveniência: wiki \textperiodcentered{} hipotese \textperiodcentered{} confiança 1.0 \textperiodcentered{} domínio ciencias\_de\_fora \textperiodcentered{} história 4 versões no jornal}

%CRISTAL {"arestas":[{"alvo":"consciencia","confianca":1.0,"epistemico":"fato","origem":"wiki","rel":"explica"},{"alvo":"tiffany","confianca":1.0,"epistemico":"fato","origem":"wiki","rel":"referencia"},{"alvo":"hopfield","confianca":1.0,"epistemico":"fato","origem":"wiki","rel":"referencia"},{"alvo":"violacao_bell","confianca":1.0,"epistemico":"fato","origem":"wiki","rel":"relaciona"}],"confianca":1.0,"contraexemplos":[],"descricao":"Tononi: a consciência é a informação integrada Φ de um sistema — o que não pode ser decomposto em partes independentes. Métrica objetiva de mente.","epistemico":"inferencia","exemplos":[],"id":"iit_phi","memoria":"semantica","meta":{"autor":"Tononi","dominio":"filosofia","fonte":""},"origem":"wiki","palavras_chave":[],"sinonimos":[],"tipo":"conceito","titulo":"Teoria da Informação Integrada (Φ)"}
\section{Teoria da Informação Integrada (\ensuremath{\Phi})}
Tononi: a consciência é a informação integrada \ensuremath{\Phi} de um sistema — o que não pode ser decomposto em partes independentes. Métrica objetiva de mente.
\prov{relações: explica consciencia; referencia tiffany; referencia hopfield; relaciona violacao\_bell}
\prov{proveniência: wiki \textperiodcentered{} inferencia \textperiodcentered{} confiança 1.0 \textperiodcentered{} domínio filosofia \textperiodcentered{} história 5 versões no jornal}

%CRISTAL {"arestas":[{"alvo":"alcorao_palavra_revelada","confianca":1.0,"epistemico":"fato","origem":"wiki","rel":"usa"}],"confianca":1.0,"contraexemplos":[],"descricao":"Islã: a doutrina de que o Alcorão tem uma qualidade milagrosa, em forma e conteúdo, que nenhuma fala humana pode igualar; o texto desafia os céticos a produzir uma sura semelhante — a inimitabilidade como prova performática de sua origem.","epistemico":"inferencia","exemplos":[],"id":"ijaz","memoria":"semantica","meta":{"autor":"Alcorão","dominio":"sagrado","fonte":""},"origem":"wiki","palavras_chave":[],"sinonimos":[],"tipo":"conceito","titulo":"I'jaz — a Inimitabilidade do Alcorão"}
\section{I'jaz — a Inimitabilidade do Alcorão}
Islã: a doutrina de que o Alcorão tem uma qualidade milagrosa, em forma e conteúdo, que nenhuma fala humana pode igualar; o texto desafia os céticos a produzir uma sura semelhante — a inimitabilidade como prova performática de sua origem.
\prov{relações: usa alcorao\_palavra\_revelada}
\prov{proveniência: wiki \textperiodcentered{} inferencia \textperiodcentered{} confiança 1.0 \textperiodcentered{} domínio sagrado \textperiodcentered{} história 2 versões no jornal}

%CRISTAL {"arestas":[{"alvo":"cristianismo","confianca":1.0,"epistemico":"fato","origem":"wiki","rel":"parte_de"},{"alvo":"consciencia","confianca":1.0,"epistemico":"fato","origem":"wiki","rel":"relaciona"}],"confianca":1.0,"contraexemplos":[],"descricao":"Bíblia (Gênesis 1:27): o humano é criado à imagem e semelhança de Deus, o que lhe confere dignidade singular, racionalidade e vocação — o fundamento teológico da interioridade e da consciência como marca distintiva.","epistemico":"inferencia","exemplos":[],"id":"imago_dei","memoria":"semantica","meta":{"autor":"Bíblia","dominio":"sagrado","fonte":""},"origem":"wiki","palavras_chave":[],"sinonimos":[],"tipo":"conceito","titulo":"Imago Dei — à Imagem de Deus"}
\section{Imago Dei — à Imagem de Deus}
Bíblia (Gênesis 1:27): o humano é criado à imagem e semelhança de Deus, o que lhe confere dignidade singular, racionalidade e vocação — o fundamento teológico da interioridade e da consciência como marca distintiva.
\prov{relações: parte\_de cristianismo; relaciona consciencia}
\prov{proveniência: wiki \textperiodcentered{} inferencia \textperiodcentered{} confiança 1.0 \textperiodcentered{} domínio sagrado \textperiodcentered{} história 3 versões no jornal}

%CRISTAL {"arestas":[{"alvo":"etica","confianca":1.0,"epistemico":"fato","origem":"wiki","rel":"parte_de"},{"alvo":"utilitarismo","confianca":1.0,"epistemico":"fato","origem":"wiki","rel":"contradiz"},{"alvo":"virtualizacao","confianca":0.5,"epistemico":"fato","origem":"wiki","rel":"relaciona"},{"alvo":"word","confianca":0.5,"epistemico":"fato","origem":"wiki","rel":"relaciona"}],"confianca":1.0,"contraexemplos":[],"descricao":"Kant: age só segundo a máxima que possas querer como lei universal; trata o outro como fim, nunca só meio.","epistemico":"inferencia","exemplos":[],"id":"imperativo_categorico","memoria":"semantica","meta":{"autor":"Kant","dominio":"ciencias_de_fora","fonte":""},"origem":"wiki","palavras_chave":[],"sinonimos":[],"tipo":"conceito","titulo":"Imperativo Categórico"}
\section{Imperativo Categórico}
Kant: age só segundo a máxima que possas querer como lei universal; trata o outro como fim, nunca só meio.
\prov{relações: parte\_de etica; contradiz utilitarismo; relaciona virtualizacao; relaciona word}
\prov{proveniência: wiki \textperiodcentered{} inferencia \textperiodcentered{} confiança 1.0 \textperiodcentered{} domínio ciencias\_de\_fora \textperiodcentered{} história 5 versões no jornal}

%CRISTAL {"arestas":[{"alvo":"atos_de_fala","confianca":1.0,"epistemico":"fato","origem":"wiki","rel":"relaciona"},{"alvo":"tiffany","confianca":1.0,"epistemico":"fato","origem":"wiki","rel":"explica"},{"alvo":"word","confianca":0.5,"epistemico":"fato","origem":"wiki","rel":"relaciona"}],"confianca":1.0,"contraexemplos":[],"descricao":"Grice: distingue o que é dito do que é implicado; as máximas (qualidade, quantidade, relevância, modo) regem a conversa cooperativa.","epistemico":"inferencia","exemplos":[],"id":"implicatura_grice","memoria":"semantica","meta":{"autor":"Grice","dominio":"ciencias_de_fora","fonte":""},"origem":"wiki","palavras_chave":[],"sinonimos":[],"tipo":"conceito","titulo":"Implicatura Conversacional"}
\section{Implicatura Conversacional}
Grice: distingue o que é dito do que é implicado; as máximas (qualidade, quantidade, relevância, modo) regem a conversa cooperativa.
\prov{relações: relaciona atos\_de\_fala; explica tiffany; relaciona word}
\prov{proveniência: wiki \textperiodcentered{} inferencia \textperiodcentered{} confiança 1.0 \textperiodcentered{} domínio ciencias\_de\_fora \textperiodcentered{} história 4 versões no jornal}

%CRISTAL {"arestas":[{"alvo":"artes_visuais","confianca":1.0,"epistemico":"fato","origem":"wiki","rel":"parte_de"},{"alvo":"teoria_das_cores","confianca":1.0,"epistemico":"fato","origem":"wiki","rel":"usa"},{"alvo":"vanguarda_autonomia","confianca":1.0,"epistemico":"fato","origem":"wiki","rel":"relaciona"}],"confianca":1.0,"contraexemplos":[],"descricao":"Monet e outros: pintar a luz e o instante, a pincelada solta e a cor pura — a percepção antes do conceito.","epistemico":"inferencia","exemplos":[],"id":"impressionismo","memoria":"semantica","meta":{"autor":"Monet","dominio":"ciencias_de_fora","fonte":""},"origem":"wiki","palavras_chave":[],"sinonimos":[],"tipo":"conceito","titulo":"Impressionismo"}
\section{Impressionismo}
Monet e outros: pintar a luz e o instante, a pincelada solta e a cor pura — a percepção antes do conceito.
\prov{relações: parte\_de artes\_visuais; usa teoria\_das\_cores; relaciona vanguarda\_autonomia}
\prov{proveniência: wiki \textperiodcentered{} inferencia \textperiodcentered{} confiança 1.0 \textperiodcentered{} domínio ciencias\_de\_fora \textperiodcentered{} história 5 versões no jornal}

%CRISTAL {"arestas":[{"alvo":"saude_publica","confianca":1.0,"epistemico":"fato","origem":"wiki","rel":"parte_de"},{"alvo":"niveis_de_prevencao","confianca":1.0,"epistemico":"fato","origem":"wiki","rel":"usa"},{"alvo":"dois_sistemas_kahneman","confianca":1.0,"epistemico":"fato","origem":"wiki","rel":"relaciona"}],"confianca":1.0,"contraexemplos":[],"descricao":"Quando uma fração suficiente da população é imune (limiar dependente da contagiosidade — ~95% para o sarampo), a transmissão se interrompe e protege até os não imunizados. A hesitação vacinal é, segundo a OMS, uma das dez maiores ameaças à saúde global.","epistemico":"fato","exemplos":[],"id":"imunidade_de_rebanho","memoria":"semantica","meta":{"autor":"—","dominio":"ciencias_de_fora","fonte":""},"origem":"wiki","palavras_chave":[],"sinonimos":[],"tipo":"conceito","titulo":"Imunidade de Rebanho"}
\section{Imunidade de Rebanho}
Quando uma fração suficiente da população é imune (limiar dependente da contagiosidade — \~{}95\% para o sarampo), a transmissão se interrompe e protege até os não imunizados. A hesitação vacinal é, segundo a OMS, uma das dez maiores ameaças à saúde global.
\prov{relações: parte\_de saude\_publica; usa niveis\_de\_prevencao; relaciona dois\_sistemas\_kahneman}
\prov{proveniência: wiki \textperiodcentered{} fato \textperiodcentered{} confiança 1.0 \textperiodcentered{} domínio ciencias\_de\_fora \textperiodcentered{} história 4 versões no jornal}

%CRISTAL {"arestas":[{"alvo":"idiotismo_nao_examinar","confianca":1.0,"epistemico":"fato","origem":"wiki","rel":"relaciona"},{"alvo":"vida-morte","confianca":0.5,"epistemico":"fato","origem":"wiki","rel":"relaciona"},{"alvo":"virtualizacao","confianca":0.5,"epistemico":"fato","origem":"wiki","rel":"relaciona"},{"alvo":"telemetria_shadow_producao","confianca":0.5,"epistemico":"fato","origem":"wiki","rel":"relaciona"},{"alvo":"pell","confianca":0.5,"epistemico":"fato","origem":"wiki","rel":"relaciona"},{"alvo":"pipeline_de_ingestao","confianca":0.5,"epistemico":"fato","origem":"wiki","rel":"relaciona"},{"alvo":"kappa","confianca":0.5,"epistemico":"fato","origem":"wiki","rel":"relaciona"},{"alvo":"pow_metal","confianca":0.5,"epistemico":"fato","origem":"wiki","rel":"relaciona"}],"confianca":1.0,"contraexemplos":[],"descricao":"Lopes: vítimas do engano social sistemático ('ópio do povo'); mas co-responsáveis pelo dever de examinar.","epistemico":"inferencia","exemplos":[],"id":"incautos_enganados","memoria":"semantica","meta":{"autor":"Gentil Lopes","dominio":"ciencias_de_fora","fonte":""},"origem":"livro","palavras_chave":[],"sinonimos":[],"tipo":"conceito","titulo":"Os Incautos"}
\section{Os Incautos}
Lopes: vítimas do engano social sistemático ('ópio do povo'); mas co-responsáveis pelo dever de examinar.
\prov{relações: relaciona idiotismo\_nao\_examinar; relaciona vida-morte; relaciona virtualizacao; relaciona telemetria\_shadow\_producao; relaciona pell; relaciona pipeline\_de\_ingestao; relaciona kappa; relaciona pow\_metal}
\prov{proveniência: livro \textperiodcentered{} inferencia \textperiodcentered{} confiança 1.0 \textperiodcentered{} domínio ciencias\_de\_fora \textperiodcentered{} história 10 versões no jornal}

%CRISTAL {"arestas":[{"alvo":"programa_de_hilbert","confianca":1.0,"epistemico":"fato","origem":"wiki","rel":"contradiz"},{"alvo":"logica_matematica","confianca":1.0,"epistemico":"fato","origem":"wiki","rel":"parte_de"},{"alvo":"expressividade_sub_turing","confianca":1.0,"epistemico":"fato","origem":"wiki","rel":"explica"},{"alvo":"godel_incompletude","confianca":1.0,"epistemico":"fato","origem":"wiki","rel":"relaciona"}],"confianca":1.0,"contraexemplos":[],"descricao":"Gödel (1931): todo sistema formal rico o bastante é incompleto (há verdades indemonstráveis) e não prova a própria consistência.","epistemico":"fato","exemplos":[],"id":"incompletude_godel","memoria":"semantica","meta":{"autor":"Gödel","dominio":"ciencias_de_fora","fonte":""},"origem":"wiki","palavras_chave":[],"sinonimos":[],"tipo":"conceito","titulo":"Teoremas da Incompletude"}
\section{Teoremas da Incompletude}
Gödel (1931): todo sistema formal rico o bastante é incompleto (há verdades indemonstráveis) e não prova a própria consistência.
\prov{relações: contradiz programa\_de\_hilbert; parte\_de logica\_matematica; explica expressividade\_sub\_turing; relaciona godel\_incompletude}
\prov{proveniência: wiki \textperiodcentered{} fato \textperiodcentered{} confiança 1.0 \textperiodcentered{} domínio ciencias\_de\_fora \textperiodcentered{} história 9 versões no jornal}

%CRISTAL {"arestas":[{"alvo":"psicanalise","confianca":1.0,"epistemico":"fato","origem":"wiki","rel":"parte_de"},{"alvo":"virtualizacao","confianca":0.5,"epistemico":"fato","origem":"wiki","rel":"relaciona"},{"alvo":"word","confianca":0.5,"epistemico":"fato","origem":"wiki","rel":"relaciona"},{"alvo":"vontade_de_poder","confianca":0.5,"epistemico":"fato","origem":"wiki","rel":"relaciona"},{"alvo":"while","confianca":0.5,"epistemico":"fato","origem":"wiki","rel":"relaciona"}],"confianca":1.0,"contraexemplos":[],"descricao":"Freud: vasto reservatório de impulsos e traumas reprimidos fora da consciência (o iceberg). Acesso indireto; não-falsável.","epistemico":"hipotese","exemplos":[],"id":"inconsciente","memoria":"semantica","meta":{"autor":"Freud","dominio":"ciencias_de_fora","fonte":""},"origem":"wiki","palavras_chave":[],"sinonimos":[],"tipo":"conceito","titulo":"O Inconsciente"}
\section{O Inconsciente}
Freud: vasto reservatório de impulsos e traumas reprimidos fora da consciência (o iceberg). Acesso indireto; não-falsável.
\prov{relações: parte\_de psicanalise; relaciona virtualizacao; relaciona word; relaciona vontade\_de\_poder; relaciona while}
\prov{proveniência: wiki \textperiodcentered{} hipotese \textperiodcentered{} confiança 1.0 \textperiodcentered{} domínio ciencias\_de\_fora \textperiodcentered{} história 6 versões no jornal}

%CRISTAL {"arestas":[{"alvo":"inconsciente","confianca":1.0,"epistemico":"fato","origem":"wiki","rel":"generaliza"},{"alvo":"word","confianca":0.5,"epistemico":"fato","origem":"wiki","rel":"relaciona"},{"alvo":"virtualizacao","confianca":0.5,"epistemico":"fato","origem":"wiki","rel":"relaciona"},{"alvo":"vida-morte","confianca":0.5,"epistemico":"fato","origem":"wiki","rel":"relaciona"},{"alvo":"vita_contemplativa","confianca":0.5,"epistemico":"fato","origem":"wiki","rel":"relaciona"}],"confianca":1.0,"contraexemplos":[],"descricao":"Jung: uma camada psíquica herdada, comum à humanidade, com símbolos universais (Herói, Sombra, Anima). Especulativo.","epistemico":"hipotese","exemplos":[],"id":"inconsciente_coletivo","memoria":"semantica","meta":{"autor":"Jung","dominio":"ciencias_de_fora","fonte":""},"origem":"wiki","palavras_chave":[],"sinonimos":[],"tipo":"conceito","titulo":"Inconsciente Coletivo e Arquétipos"}
\section{Inconsciente Coletivo e Arquétipos}
Jung: uma camada psíquica herdada, comum à humanidade, com símbolos universais (Herói, Sombra, Anima). Especulativo.
\prov{relações: generaliza inconsciente; relaciona word; relaciona virtualizacao; relaciona vida-morte; relaciona vita\_contemplativa}
\prov{proveniência: wiki \textperiodcentered{} hipotese \textperiodcentered{} confiança 1.0 \textperiodcentered{} domínio ciencias\_de\_fora \textperiodcentered{} história 6 versões no jornal}

%CRISTAL {"arestas":[{"alvo":"jogos_de_linguagem","confianca":1.0,"epistemico":"fato","origem":"wiki","rel":"usa"},{"alvo":"word","confianca":0.5,"epistemico":"fato","origem":"wiki","rel":"relaciona"}],"confianca":1.0,"contraexemplos":[],"descricao":"Quine: nenhum fato objetivo fixa a tradução correta ('gavagai'); a referência é inescrutável — só o comportamento é determinado.","epistemico":"inferencia","exemplos":[],"id":"indeterminacao_traducao","memoria":"semantica","meta":{"autor":"Quine","dominio":"ciencias_de_fora","fonte":""},"origem":"wiki","palavras_chave":[],"sinonimos":[],"tipo":"conceito","titulo":"Indeterminação da Tradução"}
\section{Indeterminação da Tradução}
Quine: nenhum fato objetivo fixa a tradução correta ('gavagai'); a referência é inescrutável — só o comportamento é determinado.
\prov{relações: usa jogos\_de\_linguagem; relaciona word}
\prov{proveniência: wiki \textperiodcentered{} inferencia \textperiodcentered{} confiança 1.0 \textperiodcentered{} domínio ciencias\_de\_fora \textperiodcentered{} história 3 versões no jornal}

%CRISTAL {"arestas":[{"alvo":"familia_linguistica","confianca":1.0,"epistemico":"fato","origem":"wiki","rel":"especializa"}],"confianca":1.0,"contraexemplos":[],"descricao":"A maior família: do português ao hindi, do russo ao inglês — ~3 bilhões de falantes, descendentes de uma proto-língua da estepe.","epistemico":"inferencia","exemplos":[],"id":"indo_europeu","memoria":"semantica","meta":{"dominio":"ciencias_de_fora","fonte":""},"origem":"wiki","palavras_chave":[],"sinonimos":[],"tipo":"conceito","titulo":"Indo-europeu"}
\section{Indo-europeu}
A maior família: do português ao hindi, do russo ao inglês — \~{}3 bilhões de falantes, descendentes de uma proto-língua da estepe.
\prov{relações: especializa familia\_linguistica}
\prov{proveniência: wiki \textperiodcentered{} inferencia \textperiodcentered{} confiança 1.0 \textperiodcentered{} domínio ciencias\_de\_fora \textperiodcentered{} história 2 versões no jornal}

%CRISTAL {"arestas":[{"alvo":"estudos_de_midia","confianca":1.0,"epistemico":"fato","origem":"wiki","rel":"parte_de"},{"alvo":"capital_cultural","confianca":1.0,"epistemico":"fato","origem":"wiki","rel":"relaciona"},{"alvo":"aura_benjamin","confianca":1.0,"epistemico":"fato","origem":"wiki","rel":"relaciona"}],"confianca":1.0,"contraexemplos":[],"descricao":"Adorno & Horkheimer: sob o capitalismo tardio a cultura funciona como indústria — bens simbólicos produzidos em massa por fórmulas padronizadas que nivelam o gosto e pacificam o consumidor, vendendo conformismo como satisfação. A pseudoindividualidade mascara o molde único.","epistemico":"inferencia","exemplos":[],"id":"industria_cultural","memoria":"semantica","meta":{"autor":"Adorno & Horkheimer","dominio":"ciencias_de_fora","fonte":""},"origem":"wiki","palavras_chave":[],"sinonimos":[],"tipo":"conceito","titulo":"Indústria Cultural (Adorno & Horkheimer)"}
\section{Indústria Cultural (Adorno \& Horkheimer)}
Adorno \& Horkheimer: sob o capitalismo tardio a cultura funciona como indústria — bens simbólicos produzidos em massa por fórmulas padronizadas que nivelam o gosto e pacificam o consumidor, vendendo conformismo como satisfação. A pseudoindividualidade mascara o molde único.
\prov{relações: parte\_de estudos\_de\_midia; relaciona capital\_cultural; relaciona aura\_benjamin}
\prov{proveniência: wiki \textperiodcentered{} inferencia \textperiodcentered{} confiança 1.0 \textperiodcentered{} domínio ciencias\_de\_fora \textperiodcentered{} história 4 versões no jornal}

%CRISTAL {"arestas":[{"alvo":"teoria_dos_numeros","confianca":1.0,"epistemico":"fato","origem":"wiki","rel":"parte_de"},{"alvo":"metodo_axiomatico","confianca":1.0,"epistemico":"fato","origem":"wiki","rel":"usa"}],"confianca":1.0,"contraexemplos":[],"descricao":"Euclides: os primos não acabam — supor uma lista completa e multiplicá-la +1 gera um novo. A prova por contradição arquetípica.","epistemico":"fato","exemplos":[],"id":"infinitude_dos_primos","memoria":"semantica","meta":{"autor":"Euclides","dominio":"ciencias_de_fora","fonte":""},"origem":"wiki","palavras_chave":[],"sinonimos":[],"tipo":"conceito","titulo":"Infinitude dos Primos"}
\section{Infinitude dos Primos}
Euclides: os primos não acabam — supor uma lista completa e multiplicá-la +1 gera um novo. A prova por contradição arquetípica.
\prov{relações: parte\_de teoria\_dos\_numeros; usa metodo\_axiomatico}
\prov{proveniência: wiki \textperiodcentered{} fato \textperiodcentered{} confiança 1.0 \textperiodcentered{} domínio ciencias\_de\_fora \textperiodcentered{} história 6 versões no jornal}

%CRISTAL {"arestas":[{"alvo":"filosofia_informacao_floridi","confianca":1.0,"epistemico":"fato","origem":"wiki","rel":"usa"},{"alvo":"consciencia","confianca":1.0,"epistemico":"fato","origem":"wiki","rel":"relaciona"}],"confianca":1.0,"contraexemplos":[],"descricao":"Floridi: a infosfera é o ambiente total de informação em que vivemos 'onlife' (online/offline fundidos); os 'inforgs' são organismos informacionais — nós mesmos como agentes de informação nesse meio.","epistemico":"inferencia","exemplos":[],"id":"infosfera_inforgs","memoria":"semantica","meta":{"autor":"Floridi","dominio":"ciencias_de_fora","fonte":""},"origem":"wiki","palavras_chave":[],"sinonimos":[],"tipo":"conceito","titulo":"Infosfera e Inforgs (Floridi)"}
\section{Infosfera e Inforgs (Floridi)}
Floridi: a infosfera é o ambiente total de informação em que vivemos 'onlife' (online/offline fundidos); os 'inforgs' são organismos informacionais — nós mesmos como agentes de informação nesse meio.
\prov{relações: usa filosofia\_informacao\_floridi; relaciona consciencia}
\prov{proveniência: wiki \textperiodcentered{} inferencia \textperiodcentered{} confiança 1.0 \textperiodcentered{} domínio ciencias\_de\_fora \textperiodcentered{} história 3 versões no jornal}

%CRISTAL {"arestas":[{"alvo":"indo_europeu","confianca":1.0,"epistemico":"fato","origem":"wiki","rel":"parte_de"},{"alvo":"lingua_franca","confianca":1.0,"epistemico":"fato","origem":"wiki","rel":"explica"}],"confianca":1.0,"contraexemplos":[],"descricao":"Língua germânica com enorme empréstimo românico — a língua franca global da ciência, da tecnologia e da internet.","epistemico":"fato","exemplos":[],"id":"ingles","memoria":"semantica","meta":{"dominio":"ciencias_de_fora","fonte":""},"origem":"wiki","palavras_chave":[],"sinonimos":[],"tipo":"conceito","titulo":"Inglês"}
\section{Inglês}
Língua germânica com enorme empréstimo românico — a língua franca global da ciência, da tecnologia e da internet.
\prov{relações: parte\_de indo\_europeu; explica lingua\_franca}
\prov{proveniência: wiki \textperiodcentered{} fato \textperiodcentered{} confiança 1.0 \textperiodcentered{} domínio ciencias\_de\_fora \textperiodcentered{} história 3 versões no jornal}

%CRISTAL {"arestas":[{"alvo":"justica_equidade_rawls","confianca":1.0,"epistemico":"fato","origem":"wiki","rel":"especializa"},{"alvo":"incautos_enganados","confianca":1.0,"epistemico":"fato","origem":"wiki","rel":"explica"}],"confianca":1.0,"contraexemplos":[],"descricao":"Fricker: negar credibilidade a alguém por preconceito (testemunhal) ou negar-lhe conceitos para articular a experiência (hermenêutica).","epistemico":"inferencia","exemplos":[],"id":"injustica_epistemica","memoria":"semantica","meta":{"autor":"Fricker","dominio":"ciencias_de_fora","fonte":""},"origem":"wiki","palavras_chave":[],"sinonimos":[],"tipo":"conceito","titulo":"Injustiça Epistêmica"}
\section{Injustiça Epistêmica}
Fricker: negar credibilidade a alguém por preconceito (testemunhal) ou negar-lhe conceitos para articular a experiência (hermenêutica).
\prov{relações: especializa justica\_equidade\_rawls; explica incautos\_enganados}
\prov{proveniência: wiki \textperiodcentered{} inferencia \textperiodcentered{} confiança 1.0 \textperiodcentered{} domínio ciencias\_de\_fora \textperiodcentered{} história 3 versões no jornal}

%CRISTAL {"arestas":[{"alvo":"comunicacao_digital","confianca":1.0,"epistemico":"fato","origem":"wiki","rel":"parte_de"},{"alvo":"ciberespaco","confianca":1.0,"epistemico":"fato","origem":"wiki","rel":"usa"}],"confianca":1.0,"contraexemplos":[],"descricao":"Pierre Lévy: uma inteligência distribuída por toda parte e coordenada em tempo real — 'ninguém sabe tudo, todos sabem algo', e o saber reside na humanidade coordenada pelo ciberespaço.","epistemico":"inferencia","exemplos":[],"id":"inteligencia_coletiva","memoria":"semantica","meta":{"autor":"Lévy","dominio":"ciencias_de_fora","fonte":""},"origem":"wiki","palavras_chave":[],"sinonimos":[],"tipo":"conceito","titulo":"Inteligência Coletiva (Lévy)"}
\section{Inteligência Coletiva (Lévy)}
Pierre Lévy: uma inteligência distribuída por toda parte e coordenada em tempo real — 'ninguém sabe tudo, todos sabem algo', e o saber reside na humanidade coordenada pelo ciberespaço.
\prov{relações: parte\_de comunicacao\_digital; usa ciberespaco}
\prov{proveniência: wiki \textperiodcentered{} inferencia \textperiodcentered{} confiança 1.0 \textperiodcentered{} domínio ciencias\_de\_fora \textperiodcentered{} história 3 versões no jornal}

%CRISTAL {"arestas":[{"alvo":"filosofia_da_computacao","confianca":1.0,"epistemico":"fato","origem":"wiki","rel":"parte_de"},{"alvo":"consciencia","confianca":1.0,"epistemico":"fato","origem":"wiki","rel":"usa"},{"alvo":"tiffany","confianca":1.0,"epistemico":"fato","origem":"wiki","rel":"explica"},{"alvo":"mente_computacao","confianca":1.0,"epistemico":"fato","origem":"wiki","rel":"relaciona"},{"alvo":"von_neumann_wigner","confianca":1.0,"epistemico":"fato","origem":"wiki","rel":"relaciona"},{"alvo":"o_substrato_a_maquina_e_funcao_de_onda","confianca":1.0,"epistemico":"fato","origem":"wiki","rel":"relaciona"}],"confianca":1.0,"contraexemplos":[],"descricao":"Lopes: inteligência e consciência não residem 'fora, nos objetos' — surgem na relação com o cérebro humano; são produto de uma relação e apenas possibilidades que se atualizam relacionalmente. (Luz que só existe na relação sol↔olho.)","epistemico":"inferencia","exemplos":[],"id":"inteligencia_consciencia_relacional","memoria":"semantica","meta":{"autor":"Gentil Lopes","dominio":"ciencias_de_fora","fonte":""},"origem":"wiki","palavras_chave":[],"sinonimos":[],"tipo":"conceito","titulo":"Inteligência e Consciência são Relacionais (Lopes)"}
\section{Inteligência e Consciência são Relacionais (Lopes)}
Lopes: inteligência e consciência não residem 'fora, nos objetos' — surgem na relação com o cérebro humano; são produto de uma relação e apenas possibilidades que se atualizam relacionalmente. (Luz que só existe na relação sol\ensuremath{\leftrightarrow}olho.)
\prov{relações: parte\_de filosofia\_da\_computacao; usa consciencia; explica tiffany; relaciona mente\_computacao; relaciona von\_neumann\_wigner; relaciona o\_substrato\_a\_maquina\_e\_funcao\_de\_onda}
\prov{proveniência: wiki \textperiodcentered{} inferencia \textperiodcentered{} confiança 1.0 \textperiodcentered{} domínio ciencias\_de\_fora \textperiodcentered{} história 8 versões no jornal}

%CRISTAL {"arestas":[{"alvo":"teorias_aprendizagem","confianca":1.0,"epistemico":"fato","origem":"wiki","rel":"parte_de"},{"alvo":"psicologia","confianca":1.0,"epistemico":"fato","origem":"wiki","rel":"relaciona"}],"confianca":1.0,"contraexemplos":[],"descricao":"Gardner: a inteligência não seria única, mas um conjunto de inteligências autônomas (linguística, lógico-matemática, espacial, musical...). Criticada como não-falsável ('neuromito'): as habilidades medidas correlacionam num fator g, o oposto do previsto.","epistemico":"hipotese","exemplos":[],"id":"inteligencias_multiplas","memoria":"semantica","meta":{"autor":"Gardner","dominio":"ciencias_de_fora","fonte":""},"origem":"wiki","palavras_chave":[],"sinonimos":[],"tipo":"conceito","titulo":"Inteligências Múltiplas (Gardner)"}
\section{Inteligências Múltiplas (Gardner)}
Gardner: a inteligência não seria única, mas um conjunto de inteligências autônomas (linguística, lógico-matemática, espacial, musical...). Criticada como não-falsável ('neuromito'): as habilidades medidas correlacionam num fator g, o oposto do previsto.
\prov{relações: parte\_de teorias\_aprendizagem; relaciona psicologia}
\prov{proveniência: wiki \textperiodcentered{} hipotese \textperiodcentered{} confiança 1.0 \textperiodcentered{} domínio ciencias\_de\_fora \textperiodcentered{} história 3 versões no jornal}

%CRISTAL {"arestas":[{"alvo":"traducao","confianca":1.0,"epistemico":"fato","origem":"wiki","rel":"usa"},{"alvo":"gramatica_universal","confianca":1.0,"epistemico":"fato","origem":"wiki","rel":"relaciona"}],"confianca":1.0,"contraexemplos":[],"descricao":"Uma representação do SENTIDO independente de qualquer língua — o pivô por onde N línguas se traduzem sem N² pares. O sonho: o significado puro, sem forma.","epistemico":"inferencia","exemplos":[],"id":"interlingua","memoria":"semantica","meta":{"dominio":"ciencias_de_fora","fonte":""},"origem":"wiki","palavras_chave":[],"sinonimos":[],"tipo":"conceito","titulo":"Interlíngua"}
\section{Interlíngua}
Uma representação do SENTIDO independente de qualquer língua — o pivô por onde N línguas se traduzem sem N\ensuremath{^{2}} pares. O sonho: o significado puro, sem forma.
\prov{relações: usa traducao; relaciona gramatica\_universal}
\prov{proveniência: wiki \textperiodcentered{} inferencia \textperiodcentered{} confiança 1.0 \textperiodcentered{} domínio ciencias\_de\_fora \textperiodcentered{} história 3 versões no jornal}

%CRISTAL {"arestas":[{"alvo":"signo_triadico","confianca":1.0,"epistemico":"fato","origem":"wiki","rel":"parte_de"},{"alvo":"simbolo_vs_significado","confianca":1.0,"epistemico":"fato","origem":"wiki","rel":"explica"},{"alvo":"tipos_de_interpretante","confianca":1.0,"epistemico":"fato","origem":"wiki","rel":"relaciona"}],"confianca":1.0,"contraexemplos":[],"descricao":"Peirce: o efeito de um signo é OUTRO signo na mente — o sentido nunca se fecha, remete adiante.","epistemico":"fato","exemplos":[],"id":"interpretante","memoria":"semantica","meta":{"autor":"Peirce","dominio":"ciencias_de_fora","fonte":""},"origem":"wiki","palavras_chave":[],"sinonimos":[],"tipo":"conceito","titulo":"O Interpretante"}
\section{O Interpretante}
Peirce: o efeito de um signo é OUTRO signo na mente — o sentido nunca se fecha, remete adiante.
\prov{relações: parte\_de signo\_triadico; explica simbolo\_vs\_significado; relaciona tipos\_de\_interpretante}
\prov{proveniência: wiki \textperiodcentered{} fato \textperiodcentered{} confiança 1.0 \textperiodcentered{} domínio ciencias\_de\_fora \textperiodcentered{} história 6 versões no jornal}

%CRISTAL {"arestas":[{"alvo":"bakhtin_dialogismo","confianca":1.0,"epistemico":"fato","origem":"wiki","rel":"usa"},{"alvo":"teoria_literaria","confianca":1.0,"epistemico":"fato","origem":"wiki","rel":"parte_de"},{"alvo":"semiose_ilimitada","confianca":1.0,"epistemico":"fato","origem":"wiki","rel":"usa"},{"alvo":"morte_do_autor","confianca":1.0,"epistemico":"fato","origem":"wiki","rel":"relaciona"}],"confianca":1.0,"contraexemplos":[],"descricao":"Kristeva: todo texto é tecido de outros textos — cita, ecoa e responde; nenhum sentido nasce do nada.","epistemico":"inferencia","exemplos":[],"id":"intertextualidade","memoria":"semantica","meta":{"autor":"Kristeva","dominio":"ciencias_de_fora","fonte":""},"origem":"wiki","palavras_chave":[],"sinonimos":[],"tipo":"conceito","titulo":"Intertextualidade"}
\section{Intertextualidade}
Kristeva: todo texto é tecido de outros textos — cita, ecoa e responde; nenhum sentido nasce do nada.
\prov{relações: usa bakhtin\_dialogismo; parte\_de teoria\_literaria; usa semiose\_ilimitada; relaciona morte\_do\_autor}
\prov{proveniência: wiki \textperiodcentered{} inferencia \textperiodcentered{} confiança 1.0 \textperiodcentered{} domínio ciencias\_de\_fora \textperiodcentered{} história 7 versões no jornal}

%CRISTAL {"arestas":[{"alvo":"filosofia_da_matematica","confianca":1.0,"epistemico":"fato","origem":"wiki","rel":"parte_de"},{"alvo":"numero_indefinido","confianca":1.0,"epistemico":"fato","origem":"wiki","rel":"relaciona"}],"confianca":1.0,"contraexemplos":[],"descricao":"A matemática é construção mental; uma proposição é verdadeira se há prova construtiva. Rejeita o terceiro excluído no infinito (Brouwer).","epistemico":"inferencia","exemplos":[],"id":"intuicionismo_brouwer","memoria":"semantica","meta":{"autor":"Brouwer","dominio":"filosofia","fonte":""},"origem":"wiki","palavras_chave":[],"sinonimos":[],"tipo":"conceito","titulo":"Intuicionismo"}
\section{Intuicionismo}
A matemática é construção mental; uma proposição é verdadeira se há prova construtiva. Rejeita o terceiro excluído no infinito (Brouwer).
\prov{relações: parte\_de filosofia\_da\_matematica; relaciona numero\_indefinido}
\prov{proveniência: wiki \textperiodcentered{} inferencia \textperiodcentered{} confiança 1.0 \textperiodcentered{} domínio filosofia \textperiodcentered{} história 3 versões no jornal}

%CRISTAL {"arestas":[{"alvo":"saksi_testemunha","confianca":1.0,"epistemico":"fato","origem":"wiki","rel":"relaciona"},{"alvo":"consciencia_pura","confianca":1.0,"epistemico":"fato","origem":"wiki","rel":"usa"},{"alvo":"colapso_observador","confianca":1.0,"epistemico":"fato","origem":"wiki","rel":"relaciona"}],"confianca":1.0,"contraexemplos":[],"descricao":"Lopes: do bebê ao idoso, a identidade muda (construção social), mas o Observador — a pura capacidade de observar — permanece; no limite é o vazio. Ponte ao sākṣī e ao colapso.","epistemico":"hipotese","exemplos":[],"id":"invariante_humano","memoria":"semantica","meta":{"autor":"Gentil Lopes","dominio":"filosofia","fonte":""},"origem":"livro","palavras_chave":[],"sinonimos":[],"tipo":"conceito","titulo":"O Invariante Humano (o Observador)"}
\section{O Invariante Humano (o Observador)}
Lopes: do bebê ao idoso, a identidade muda (construção social), mas o Observador — a pura capacidade de observar — permanece; no limite é o vazio. Ponte ao s\ensuremath{\bar{a}}k\ensuremath{\d{s}}\ensuremath{\bar{i}} e ao colapso.
\prov{relações: relaciona saksi\_testemunha; usa consciencia\_pura; relaciona colapso\_observador}
\prov{proveniência: livro \textperiodcentered{} hipotese \textperiodcentered{} confiança 1.0 \textperiodcentered{} domínio filosofia \textperiodcentered{} história 5 versões no jornal}

%CRISTAL {"arestas":[{"alvo":"sagrado","confianca":1.0,"epistemico":"fato","origem":"wiki","rel":"parte_de"},{"alvo":"linhagem_abraamica","confianca":1.0,"epistemico":"fato","origem":"wiki","rel":"usa"}],"confianca":1.0,"contraexemplos":[],"descricao":"A tradição do Alcorão — a submissão (islam) ao Deus único (Allah), a última das revelações abraâmicas, estruturada em torno da unicidade (tawhid) e dos cinco pilares.","epistemico":"fato","exemplos":[],"id":"isla","memoria":"semantica","meta":{"dominio":"sagrado","fonte":""},"origem":"wiki","palavras_chave":[],"sinonimos":[],"tipo":"conceito","titulo":"Islã"}
\section{Islã}
A tradição do Alcorão — a submissão (islam) ao Deus único (Allah), a última das revelações abraâmicas, estruturada em torno da unicidade (tawhid) e dos cinco pilares.
\prov{relações: parte\_de sagrado; usa linhagem\_abraamica}
\prov{proveniência: wiki \textperiodcentered{} fato \textperiodcentered{} confiança 1.0 \textperiodcentered{} domínio sagrado \textperiodcentered{} história 3 versões no jornal}

%CRISTAL {"arestas":[{"alvo":"isla","confianca":1.0,"epistemico":"fato","origem":"wiki","rel":"parte_de"},{"alvo":"wu_wei","confianca":1.0,"epistemico":"fato","origem":"wiki","rel":"relaciona"},{"alvo":"karma_yoga_nishkama","confianca":1.0,"epistemico":"fato","origem":"wiki","rel":"relaciona"}],"confianca":1.0,"contraexemplos":[],"descricao":"Islã: 'islam' é a entrega/rendição voluntária à vontade de Deus (raiz S-L-M: paz, inteireza); muslim é aquele que se submete. A paz nasce da rendição, não da luta — abrir mão do controle do ego, como o wu wei, mas rendição a uma Vontade.","epistemico":"inferencia","exemplos":[],"id":"islam_submissao","memoria":"semantica","meta":{"autor":"Alcorão","dominio":"sagrado","fonte":""},"origem":"wiki","palavras_chave":[],"sinonimos":[],"tipo":"conceito","titulo":"Islam — a Submissão a Deus"}
\section{Islam — a Submissão a Deus}
Islã: 'islam' é a entrega/rendição voluntária à vontade de Deus (raiz S-L-M: paz, inteireza); muslim é aquele que se submete. A paz nasce da rendição, não da luta — abrir mão do controle do ego, como o wu wei, mas rendição a uma Vontade.
\prov{relações: parte\_de isla; relaciona wu\_wei; relaciona karma\_yoga\_nishkama}
\prov{proveniência: wiki \textperiodcentered{} inferencia \textperiodcentered{} confiança 1.0 \textperiodcentered{} domínio sagrado \textperiodcentered{} história 4 versões no jornal}

%CRISTAL {"arestas":[{"alvo":"tratamento_restaura_atrator","confianca":1.0,"epistemico":"fato","origem":"wiki","rel":"usa"},{"alvo":"ensaio_clinico_randomizado","confianca":1.0,"epistemico":"fato","origem":"wiki","rel":"relaciona"},{"alvo":"corpo_mineral","confianca":1.0,"epistemico":"fato","origem":"wiki","rel":"usa"},{"alvo":"numero_de_salem","confianca":1.0,"epistemico":"fato","origem":"wiki","rel":"relaciona"}],"confianca":1.0,"contraexemplos":[],"descricao":"A dose abaixo da efetiva não age (subterapêutica); acima da tóxica envenena; entre elas está a janela terapêutica — a BORDA da dose. O índice terapêutico (dose tóxica / efetiva) mede quão larga ela é: drogas de janela estreita (varfarina, lítio) exigem ficar na borda com precisão. Verif.: dose 2 na janela 1–5 é terapêutica; 0.5 não age; 6 é tóxica.","epistemico":"inferencia","exemplos":[],"id":"janela_terapeutica","memoria":"semantica","meta":{"dominio":"medicina","dose2_na_janela":true,"fonte":"medicina e diagnóstico","indice_terapeutico":5.0},"origem":"wiki","palavras_chave":[],"sinonimos":[],"tipo":"conceito","titulo":"Janela Terapêutica = a Borda"}
\section{Janela Terapêutica = a Borda}
A dose abaixo da efetiva não age (subterapêutica); acima da tóxica envenena; entre elas está a janela terapêutica — a BORDA da dose. O índice terapêutico (dose tóxica / efetiva) mede quão larga ela é: drogas de janela estreita (varfarina, lítio) exigem ficar na borda com precisão. Verif.: dose 2 na janela 1–5 é terapêutica; 0.5 não age; 6 é tóxica.
\prov{relações: usa tratamento\_restaura\_atrator; relaciona ensaio\_clinico\_randomizado; usa corpo\_mineral; relaciona numero\_de\_salem}
\prov{proveniência: wiki \textperiodcentered{} medicina e diagnóstico \textperiodcentered{} inferencia \textperiodcentered{} confiança 1.0 \textperiodcentered{} domínio medicina \textperiodcentered{} história 5 versões no jornal}

%CRISTAL {"arestas":[{"alvo":"idioma","confianca":1.0,"epistemico":"fato","origem":"wiki","rel":"especializa"}],"confianca":1.0,"contraexemplos":[],"descricao":"Língua do Japão (~125 milhões), SOV aglutinante e tópico-proeminente, com três sistemas de escrita (kanji, hiragana, katakana).","epistemico":"fato","exemplos":[],"id":"japones","memoria":"semantica","meta":{"dominio":"ciencias_de_fora","fonte":""},"origem":"wiki","palavras_chave":[],"sinonimos":[],"tipo":"conceito","titulo":"Japonês"}
\section{Japonês}
Língua do Japão (\~{}125 milhões), SOV aglutinante e tópico-proeminente, com três sistemas de escrita (kanji, hiragana, katakana).
\prov{relações: especializa idioma}
\prov{proveniência: wiki \textperiodcentered{} fato \textperiodcentered{} confiança 1.0 \textperiodcentered{} domínio ciencias\_de\_fora \textperiodcentered{} história 2 versões no jornal}

%CRISTAL {"arestas":[{"alvo":"isla","confianca":1.0,"epistemico":"fato","origem":"wiki","rel":"parte_de"},{"alvo":"meditacao_cura_da_mente_gentil","confianca":1.0,"epistemico":"fato","origem":"wiki","rel":"relaciona"}],"confianca":1.0,"contraexemplos":[],"descricao":"Islã: da raiz 'esforçar-se'. A tradição distingue o 'grande jihad' (a luta interior contra o próprio ego e as paixões — a primária, segundo al-Ghazali) do 'menor jihad' (o combate exterior, classicamente defensivo). A batalha real é interna.","epistemico":"inferencia","exemplos":[],"id":"jihad","memoria":"semantica","meta":{"autor":"Alcorão","dominio":"sagrado","fonte":""},"origem":"wiki","palavras_chave":[],"sinonimos":[],"tipo":"conceito","titulo":"Jihad — o Esforço no Caminho de Deus"}
\section{Jihad — o Esforço no Caminho de Deus}
Islã: da raiz 'esforçar-se'. A tradição distingue o 'grande jihad' (a luta interior contra o próprio ego e as paixões — a primária, segundo al-Ghazali) do 'menor jihad' (o combate exterior, classicamente defensivo). A batalha real é interna.
\prov{relações: parte\_de isla; relaciona meditacao\_cura\_da\_mente\_gentil}
\prov{proveniência: wiki \textperiodcentered{} inferencia \textperiodcentered{} confiança 1.0 \textperiodcentered{} domínio sagrado \textperiodcentered{} história 3 versões no jornal}

%CRISTAL {"arestas":[{"alvo":"hinduismo","confianca":1.0,"epistemico":"fato","origem":"wiki","rel":"parte_de"},{"alvo":"neti_neti","confianca":1.0,"epistemico":"fato","origem":"wiki","rel":"relaciona"},{"alvo":"atman_brahman","confianca":1.0,"epistemico":"fato","origem":"wiki","rel":"usa"}],"confianca":1.0,"contraexemplos":[],"descricao":"Gita: o caminho da sabedoria que discrimina o Real do irreal, o Self (purusha) do não-Self (prakriti). 'O fogo do conhecimento reduz a cinzas todas as ações' — não há melhor purificador. Discernir por negação o que não é o Self, como o neti neti.","epistemico":"inferencia","exemplos":[],"id":"jnana_yoga","memoria":"semantica","meta":{"autor":"Bhagavad Gita","dominio":"sagrado","fonte":""},"origem":"wiki","palavras_chave":[],"sinonimos":[],"tipo":"conceito","titulo":"Jnana Yoga — o Caminho do Conhecimento"}
\section{Jnana Yoga — o Caminho do Conhecimento}
Gita: o caminho da sabedoria que discrimina o Real do irreal, o Self (purusha) do não-Self (prakriti). 'O fogo do conhecimento reduz a cinzas todas as ações' — não há melhor purificador. Discernir por negação o que não é o Self, como o neti neti.
\prov{relações: parte\_de hinduismo; relaciona neti\_neti; usa atman\_brahman}
\prov{proveniência: wiki \textperiodcentered{} inferencia \textperiodcentered{} confiança 1.0 \textperiodcentered{} domínio sagrado \textperiodcentered{} história 4 versões no jornal}

%CRISTAL {"arestas":[{"alvo":"ficcao_hub","confianca":1.0,"epistemico":"fato","origem":"wiki","rel":"parte_de"}],"confianca":1.0,"contraexemplos":[],"descricao":"O sub-hub de JOGOS E ANIMAÇÕES sob a ficção. Reúne as obras interativas/animadas lidas pelo dicionário: o MINECRAFT (o voxelizador jogado como blocos, minecraft_voxel — já no grafo) e o BEN 10 (o Omnitrix como o Corpo Algébrico). Cada obra é uma porta para uma instância do broca-so.","epistemico":"fato","exemplos":[],"id":"jogos_animacoes_hub","memoria":"semantica","meta":{"dominio":"ficcao","fonte":"Ficção (jogos e animações)"},"origem":"wiki","palavras_chave":[],"sinonimos":[],"tipo":"conceito","titulo":"Jogos e animações — o sub-hub (Minecraft, Ben 10)"}
\section{Jogos e animações — o sub-hub (Minecraft, Ben 10)}
O sub-hub de JOGOS E ANIMAÇÕES sob a ficção. Reúne as obras interativas/animadas lidas pelo dicionário: o MINECRAFT (o voxelizador jogado como blocos, minecraft\_voxel — já no grafo) e o BEN 10 (o Omnitrix como o Corpo Algébrico). Cada obra é uma porta para uma instância do broca-so.
\prov{relações: parte\_de ficcao\_hub}
\prov{proveniência: wiki \textperiodcentered{} Ficção (jogos e animações) \textperiodcentered{} fato \textperiodcentered{} confiança 1.0 \textperiodcentered{} domínio ficcao \textperiodcentered{} história 2 versões no jornal}

%CRISTAL {"arestas":[{"alvo":"formalismo_hilbert","confianca":1.0,"epistemico":"fato","origem":"wiki","rel":"relaciona"},{"alvo":"tiffany","confianca":1.0,"epistemico":"fato","origem":"wiki","rel":"explica"},{"alvo":"filosofia_da_linguagem","confianca":1.0,"epistemico":"fato","origem":"wiki","rel":"parte_de"}],"confianca":1.0,"contraexemplos":[],"descricao":"Wittgenstein (tardio): o sentido de uma palavra é seu USO numa prática — não há essência, há semelhanças de família.","epistemico":"inferencia","exemplos":[],"id":"jogos_de_linguagem","memoria":"semantica","meta":{"autor":"Wittgenstein","dominio":"ciencias_de_fora","fonte":""},"origem":"wiki","palavras_chave":[],"sinonimos":[],"tipo":"conceito","titulo":"Jogos de Linguagem"}
\section{Jogos de Linguagem}
Wittgenstein (tardio): o sentido de uma palavra é seu USO numa prática — não há essência, há semelhanças de família.
\prov{relações: relaciona formalismo\_hilbert; explica tiffany; parte\_de filosofia\_da\_linguagem}
\prov{proveniência: wiki \textperiodcentered{} inferencia \textperiodcentered{} confiança 1.0 \textperiodcentered{} domínio ciencias\_de\_fora \textperiodcentered{} história 6 versões no jornal}

%CRISTAL {"arestas":[{"alvo":"narrativa","confianca":1.0,"epistemico":"fato","origem":"wiki","rel":"especializa"},{"alvo":"inconsciente_coletivo","confianca":1.0,"epistemico":"fato","origem":"wiki","rel":"usa"}],"confianca":1.0,"contraexemplos":[],"descricao":"Campbell: um monomito comum aos mitos — chamado, provação, retorno transformado. Molde de incontáveis narrativas.","epistemico":"inferencia","exemplos":[],"id":"jornada_do_heroi","memoria":"semantica","meta":{"autor":"Campbell","dominio":"ciencias_de_fora","fonte":""},"origem":"wiki","palavras_chave":[],"sinonimos":[],"tipo":"conceito","titulo":"A Jornada do Herói"}
\section{A Jornada do Herói}
Campbell: um monomito comum aos mitos — chamado, provação, retorno transformado. Molde de incontáveis narrativas.
\prov{relações: especializa narrativa; usa inconsciente\_coletivo}
\prov{proveniência: wiki \textperiodcentered{} inferencia \textperiodcentered{} confiança 1.0 \textperiodcentered{} domínio ciencias\_de\_fora \textperiodcentered{} história 3 versões no jornal}

%CRISTAL {"arestas":[{"alvo":"filosofia_do_direito","confianca":1.0,"epistemico":"fato","origem":"wiki","rel":"parte_de"},{"alvo":"imperativo_categorico","confianca":1.0,"epistemico":"fato","origem":"wiki","rel":"relaciona"},{"alvo":"juspositivismo","confianca":1.0,"epistemico":"fato","origem":"wiki","rel":"contradiz"}],"confianca":1.0,"contraexemplos":[],"descricao":"Existe um direito superior ao posto, fundado na natureza ou na razão, que serve de critério de validade e justiça das leis — 'lei injusta não é lei plena'. De Aquino e Grotius a Finnis (bens humanos básicos e razoabilidade prática).","epistemico":"inferencia","exemplos":[],"id":"jusnaturalismo","memoria":"semantica","meta":{"autor":"Aquino/Grotius/Finnis","dominio":"ciencias_de_fora","fonte":""},"origem":"wiki","palavras_chave":[],"sinonimos":[],"tipo":"conceito","titulo":"Jusnaturalismo (Direito Natural)"}
\section{Jusnaturalismo (Direito Natural)}
Existe um direito superior ao posto, fundado na natureza ou na razão, que serve de critério de validade e justiça das leis — 'lei injusta não é lei plena'. De Aquino e Grotius a Finnis (bens humanos básicos e razoabilidade prática).
\prov{relações: parte\_de filosofia\_do\_direito; relaciona imperativo\_categorico; contradiz juspositivismo}
\prov{proveniência: wiki \textperiodcentered{} inferencia \textperiodcentered{} confiança 1.0 \textperiodcentered{} domínio ciencias\_de\_fora \textperiodcentered{} história 4 versões no jornal}

%CRISTAL {"arestas":[{"alvo":"filosofia_do_direito","confianca":1.0,"epistemico":"fato","origem":"wiki","rel":"parte_de"}],"confianca":1.0,"contraexemplos":[],"descricao":"A existência e o conteúdo do direito dependem de fatos sociais (quem o pôs, como), não de méritos morais — a tese da separabilidade entre direito e moral (Austin: 'a existência do direito é uma coisa; seu mérito, outra').","epistemico":"inferencia","exemplos":[],"id":"juspositivismo","memoria":"semantica","meta":{"autor":"Bentham/Austin","dominio":"ciencias_de_fora","fonte":""},"origem":"wiki","palavras_chave":[],"sinonimos":[],"tipo":"conceito","titulo":"Juspositivismo"}
\section{Juspositivismo}
A existência e o conteúdo do direito dependem de fatos sociais (quem o pôs, como), não de méritos morais — a tese da separabilidade entre direito e moral (Austin: 'a existência do direito é uma coisa; seu mérito, outra').
\prov{relações: parte\_de filosofia\_do\_direito}
\prov{proveniência: wiki \textperiodcentered{} inferencia \textperiodcentered{} confiança 1.0 \textperiodcentered{} domínio ciencias\_de\_fora \textperiodcentered{} história 2 versões no jornal}

%CRISTAL {"arestas":[{"alvo":"direito_e_tecnologia","confianca":1.0,"epistemico":"fato","origem":"wiki","rel":"parte_de"},{"alvo":"artefatos_tem_politica","confianca":1.0,"epistemico":"fato","origem":"wiki","rel":"usa"},{"alvo":"code_is_law","confianca":1.0,"epistemico":"fato","origem":"wiki","rel":"relaciona"},{"alvo":"biopoder","confianca":1.0,"epistemico":"fato","origem":"wiki","rel":"relaciona"}],"confianca":1.0,"contraexemplos":[],"descricao":"O problema das decisões automatizadas no direito: o viés algorítmico reproduz a discriminação dos dados de treino, e a opacidade dificulta a accountability e o direito à explicação (GDPR art. 22).","epistemico":"inferencia","exemplos":[],"id":"justica_algoritmica","memoria":"semantica","meta":{"autor":"Wachter/Mittelstadt","dominio":"ciencias_de_fora","fonte":""},"origem":"wiki","palavras_chave":[],"sinonimos":[],"tipo":"conceito","titulo":"Justiça Algorítmica / Juiz-Máquina"}
\section{Justiça Algorítmica / Juiz-Máquina}
O problema das decisões automatizadas no direito: o viés algorítmico reproduz a discriminação dos dados de treino, e a opacidade dificulta a accountability e o direito à explicação (GDPR art. 22).
\prov{relações: parte\_de direito\_e\_tecnologia; usa artefatos\_tem\_politica; relaciona code\_is\_law; relaciona biopoder}
\prov{proveniência: wiki \textperiodcentered{} inferencia \textperiodcentered{} confiança 1.0 \textperiodcentered{} domínio ciencias\_de\_fora \textperiodcentered{} história 5 versões no jornal}

%CRISTAL {"arestas":[{"alvo":"bioetica","confianca":1.0,"epistemico":"fato","origem":"wiki","rel":"parte_de"},{"alvo":"justica_equidade_rawls","confianca":1.0,"epistemico":"fato","origem":"wiki","rel":"usa"}],"confianca":1.0,"contraexemplos":[],"descricao":"Daniels estende Rawls à saúde: ela importa por proteger o 'funcionamento normal' e, com ele, a igualdade equitativa de oportunidades — fundando o direito ao cuidado e um método justo de racionamento ('accountability for reasonableness').","epistemico":"inferencia","exemplos":[],"id":"justica_em_saude_daniels","memoria":"semantica","meta":{"autor":"Daniels","dominio":"ciencias_de_fora","fonte":""},"origem":"wiki","palavras_chave":[],"sinonimos":[],"tipo":"conceito","titulo":"Justiça em Saúde (Daniels)"}
\section{Justiça em Saúde (Daniels)}
Daniels estende Rawls à saúde: ela importa por proteger o 'funcionamento normal' e, com ele, a igualdade equitativa de oportunidades — fundando o direito ao cuidado e um método justo de racionamento ('accountability for reasonableness').
\prov{relações: parte\_de bioetica; usa justica\_equidade\_rawls}
\prov{proveniência: wiki \textperiodcentered{} inferencia \textperiodcentered{} confiança 1.0 \textperiodcentered{} domínio ciencias\_de\_fora \textperiodcentered{} história 3 versões no jornal}

%CRISTAL {"arestas":[{"alvo":"politica","confianca":1.0,"epistemico":"fato","origem":"wiki","rel":"parte_de"},{"alvo":"utilitarismo","confianca":1.0,"epistemico":"fato","origem":"wiki","rel":"contradiz"},{"alvo":"word","confianca":0.5,"epistemico":"fato","origem":"wiki","rel":"relaciona"}],"confianca":1.0,"contraexemplos":[],"descricao":"Rawls: princípios escolhidos sob o 'véu de ignorância' — iguais liberdades, e desigualdades só se beneficiam os piores.","epistemico":"inferencia","exemplos":[],"id":"justica_equidade_rawls","memoria":"semantica","meta":{"autor":"Rawls","dominio":"ciencias_de_fora","fonte":""},"origem":"wiki","palavras_chave":[],"sinonimos":[],"tipo":"conceito","titulo":"Justiça como Equidade"}
\section{Justiça como Equidade}
Rawls: princípios escolhidos sob o 'véu de ignorância' — iguais liberdades, e desigualdades só se beneficiam os piores.
\prov{relações: parte\_de politica; contradiz utilitarismo; relaciona word}
\prov{proveniência: wiki \textperiodcentered{} inferencia \textperiodcentered{} confiança 1.0 \textperiodcentered{} domínio ciencias\_de\_fora \textperiodcentered{} história 4 versões no jornal}

%CRISTAL {"arestas":[{"alvo":"pena_retribuicao","confianca":1.0,"epistemico":"fato","origem":"wiki","rel":"contradiz"},{"alvo":"direito","confianca":1.0,"epistemico":"fato","origem":"wiki","rel":"parte_de"}],"confianca":1.0,"contraexemplos":[],"descricao":"A resposta ao crime visa reparar o dano e restaurar as relações entre ofensor, vítima e comunidade, em vez de centrar-se na retribuição ou na dissuasão via castigo.","epistemico":"inferencia","exemplos":[],"id":"justica_restaurativa","memoria":"semantica","meta":{"autor":"Zehr","dominio":"ciencias_de_fora","fonte":""},"origem":"wiki","palavras_chave":[],"sinonimos":[],"tipo":"conceito","titulo":"Justiça Restaurativa"}
\section{Justiça Restaurativa}
A resposta ao crime visa reparar o dano e restaurar as relações entre ofensor, vítima e comunidade, em vez de centrar-se na retribuição ou na dissuasão via castigo.
\prov{relações: contradiz pena\_retribuicao; parte\_de direito}
\prov{proveniência: wiki \textperiodcentered{} inferencia \textperiodcentered{} confiança 1.0 \textperiodcentered{} domínio ciencias\_de\_fora \textperiodcentered{} história 3 versões no jornal}

%CRISTAL {"arestas":[{"alvo":"idealismo_berkeley","confianca":1.0,"epistemico":"fato","origem":"wiki","rel":"especializa"},{"alvo":"epistemologia","confianca":1.0,"epistemico":"fato","origem":"wiki","rel":"parte_de"},{"alvo":"word","confianca":0.5,"epistemico":"fato","origem":"wiki","rel":"relaciona"},{"alvo":"virtualizacao","confianca":0.5,"epistemico":"fato","origem":"wiki","rel":"relaciona"}],"confianca":1.0,"contraexemplos":[],"descricao":"Kant: conhecemos fenômenos (estruturados por tempo/espaço/causalidade da mente), nunca a coisa-em-si (noúmeno).","epistemico":"inferencia","exemplos":[],"id":"kant_fenomeno_numeno","memoria":"semantica","meta":{"autor":"Kant","dominio":"filosofia","fonte":""},"origem":"wiki","palavras_chave":[],"sinonimos":[],"tipo":"conceito","titulo":"Fenômeno vs Noúmeno"}
\section{Fenômeno vs Noúmeno}
Kant: conhecemos fenômenos (estruturados por tempo/espaço/causalidade da mente), nunca a coisa-em-si (noúmeno).
\prov{relações: especializa idealismo\_berkeley; parte\_de epistemologia; relaciona word; relaciona virtualizacao}
\prov{proveniência: wiki \textperiodcentered{} inferencia \textperiodcentered{} confiança 1.0 \textperiodcentered{} domínio filosofia \textperiodcentered{} história 5 versões no jornal}

%CRISTAL {"arestas":[{"alvo":"filosofia_da_tecnica","confianca":1.0,"epistemico":"fato","origem":"wiki","rel":"parte_de"},{"alvo":"tiffany","confianca":1.0,"epistemico":"fato","origem":"wiki","rel":"relaciona"}],"confianca":1.0,"contraexemplos":[],"descricao":"Kapp: toda técnica é projeção inconsciente de órgãos humanos para fora — o martelo estende o braço, a lente o olho, o telégrafo o sistema nervoso. O artefato é espelho do corpo, e só pela técnica o humano se autoconhece.","epistemico":"inferencia","exemplos":[],"id":"kapp_organprojektion","memoria":"semantica","meta":{"autor":"Kapp","dominio":"ciencias_de_fora","fonte":""},"origem":"wiki","palavras_chave":[],"sinonimos":[],"tipo":"conceito","titulo":"Projeção de Órgãos (Kapp)"}
\section{Projeção de Órgãos (Kapp)}
Kapp: toda técnica é projeção inconsciente de órgãos humanos para fora — o martelo estende o braço, a lente o olho, o telégrafo o sistema nervoso. O artefato é espelho do corpo, e só pela técnica o humano se autoconhece.
\prov{relações: parte\_de filosofia\_da\_tecnica; relaciona tiffany}
\prov{proveniência: wiki \textperiodcentered{} inferencia \textperiodcentered{} confiança 1.0 \textperiodcentered{} domínio ciencias\_de\_fora \textperiodcentered{} história 3 versões no jornal}

%CRISTAL {"arestas":[{"alvo":"hinduismo","confianca":1.0,"epistemico":"fato","origem":"wiki","rel":"parte_de"},{"alvo":"wu_wei","confianca":1.0,"epistemico":"fato","origem":"wiki","rel":"relaciona"},{"alvo":"estetica_do_fluir_gentil","confianca":1.0,"epistemico":"fato","origem":"wiki","rel":"relaciona"},{"alvo":"karma_dharma","confianca":1.0,"epistemico":"fato","origem":"wiki","rel":"usa"}],"confianca":1.0,"contraexemplos":[],"descricao":"Gita (2.47): 'só tens direito à ação, jamais aos seus frutos; que estes jamais sejam o motivo de teus atos, nem fiques preso na inação'. Agir plenamente, com todo zelo, oferecendo os resultados sem ansiar êxito nem temer fracasso. A equanimidade (samatva) diante dos pares de opostos é o yoga.","epistemico":"inferencia","exemplos":[],"id":"karma_yoga_nishkama","memoria":"semantica","meta":{"autor":"Bhagavad Gita","dominio":"sagrado","fonte":""},"origem":"wiki","palavras_chave":[],"sinonimos":[],"tipo":"conceito","titulo":"Karma Yoga — a Ação Desapegada (Nishkama Karma)"}
\section{Karma Yoga — a Ação Desapegada (Nishkama Karma)}
Gita (2.47): 'só tens direito à ação, jamais aos seus frutos; que estes jamais sejam o motivo de teus atos, nem fiques preso na inação'. Agir plenamente, com todo zelo, oferecendo os resultados sem ansiar êxito nem temer fracasso. A equanimidade (samatva) diante dos pares de opostos é o yoga.
\prov{relações: parte\_de hinduismo; relaciona wu\_wei; relaciona estetica\_do\_fluir\_gentil; usa karma\_dharma}
\prov{proveniência: wiki \textperiodcentered{} inferencia \textperiodcentered{} confiança 1.0 \textperiodcentered{} domínio sagrado \textperiodcentered{} história 5 versões no jornal}

%CRISTAL {"arestas":[{"alvo":"kelsen_teoria_pura","confianca":1.0,"epistemico":"fato","origem":"wiki","rel":"usa"},{"alvo":"godel_incompletude","confianca":1.0,"epistemico":"fato","origem":"wiki","rel":"relaciona"}],"confianca":1.0,"contraexemplos":[],"descricao":"Kelsen: a validade de cada norma decorre de outra superior, numa pirâmide que culmina numa norma fundamental pressuposta (não posta) — condição lógico-transcendental para ler atos de coerção como direito. O fundamento do sistema não é derivável de dentro dele.","epistemico":"inferencia","exemplos":[],"id":"kelsen_grundnorm","memoria":"semantica","meta":{"autor":"Kelsen","dominio":"ciencias_de_fora","fonte":""},"origem":"wiki","palavras_chave":[],"sinonimos":[],"tipo":"conceito","titulo":"Norma Fundamental (Grundnorm)"}
\section{Norma Fundamental (Grundnorm)}
Kelsen: a validade de cada norma decorre de outra superior, numa pirâmide que culmina numa norma fundamental pressuposta (não posta) — condição lógico-transcendental para ler atos de coerção como direito. O fundamento do sistema não é derivável de dentro dele.
\prov{relações: usa kelsen\_teoria\_pura; relaciona godel\_incompletude}
\prov{proveniência: wiki \textperiodcentered{} inferencia \textperiodcentered{} confiança 1.0 \textperiodcentered{} domínio ciencias\_de\_fora \textperiodcentered{} história 3 versões no jornal}

%CRISTAL {"arestas":[{"alvo":"juspositivismo","confianca":1.0,"epistemico":"fato","origem":"wiki","rel":"usa"},{"alvo":"honestidade_intelectual","confianca":1.0,"epistemico":"fato","origem":"wiki","rel":"relaciona"}],"confianca":1.0,"contraexemplos":[],"descricao":"Kelsen: uma ciência do direito depurada de política, moral e sociologia, que descreve o direito como sistema de normas cuja validade é puramente formal — o direito como esquema de interpretação do dever-ser, numa pirâmide hierárquica.","epistemico":"inferencia","exemplos":[],"id":"kelsen_teoria_pura","memoria":"semantica","meta":{"autor":"Kelsen","dominio":"ciencias_de_fora","fonte":""},"origem":"wiki","palavras_chave":[],"sinonimos":[],"tipo":"conceito","titulo":"Teoria Pura do Direito (Kelsen)"}
\section{Teoria Pura do Direito (Kelsen)}
Kelsen: uma ciência do direito depurada de política, moral e sociologia, que descreve o direito como sistema de normas cuja validade é puramente formal — o direito como esquema de interpretação do dever-ser, numa pirâmide hierárquica.
\prov{relações: usa juspositivismo; relaciona honestidade\_intelectual}
\prov{proveniência: wiki \textperiodcentered{} inferencia \textperiodcentered{} confiança 1.0 \textperiodcentered{} domínio ciencias\_de\_fora \textperiodcentered{} história 3 versões no jornal}

%CRISTAL {"arestas":[{"alvo":"hinduismo","confianca":1.0,"epistemico":"fato","origem":"wiki","rel":"parte_de"},{"alvo":"maya","confianca":1.0,"epistemico":"fato","origem":"wiki","rel":"usa"},{"alvo":"atman_brahman","confianca":1.0,"epistemico":"fato","origem":"wiki","rel":"contradiz"}],"confianca":1.0,"contraexemplos":[],"descricao":"Gita (4.7-8): Krishna é o Senhor pessoal que, embora não-nascido e imutável, se encarna por sua própria maya 'toda vez que declina a retidão e prevalece a irreligião, para proteger os bons e restabelecer a religião eterna' — a doutrina do avatar. O polo pessoal do Absoluto que a bhakti contempla.","epistemico":"inferencia","exemplos":[],"id":"krishna_avatar","memoria":"semantica","meta":{"autor":"Bhagavad Gita","dominio":"sagrado","fonte":""},"origem":"wiki","palavras_chave":[],"sinonimos":[],"tipo":"conceito","titulo":"Krishna — o Divino Pessoal (Avatar)"}
\section{Krishna — o Divino Pessoal (Avatar)}
Gita (4.7-8): Krishna é o Senhor pessoal que, embora não-nascido e imutável, se encarna por sua própria maya 'toda vez que declina a retidão e prevalece a irreligião, para proteger os bons e restabelecer a religião eterna' — a doutrina do avatar. O polo pessoal do Absoluto que a bhakti contempla.
\prov{relações: parte\_de hinduismo; usa maya; contradiz atman\_brahman}
\prov{proveniência: wiki \textperiodcentered{} inferencia \textperiodcentered{} confiança 1.0 \textperiodcentered{} domínio sagrado \textperiodcentered{} história 4 versões no jornal}

%CRISTAL {"arestas":[{"alvo":"indo_europeu","confianca":1.0,"epistemico":"fato","origem":"wiki","rel":"parte_de"},{"alvo":"lingua_viva_morta","confianca":1.0,"epistemico":"fato","origem":"wiki","rel":"usa"}],"confianca":1.0,"contraexemplos":[],"descricao":"A língua de Roma, mãe das românicas — 'morta' como materna, mas viva na ciência, no direito e na Igreja por dois milênios.","epistemico":"fato","exemplos":[],"id":"latim","memoria":"semantica","meta":{"dominio":"ciencias_de_fora","fonte":""},"origem":"wiki","palavras_chave":[],"sinonimos":[],"tipo":"conceito","titulo":"Latim"}
\section{Latim}
A língua de Roma, mãe das românicas — 'morta' como materna, mas viva na ciência, no direito e na Igreja por dois milênios.
\prov{relações: parte\_de indo\_europeu; usa lingua\_viva\_morta}
\prov{proveniência: wiki \textperiodcentered{} fato \textperiodcentered{} confiança 1.0 \textperiodcentered{} domínio ciencias\_de\_fora \textperiodcentered{} história 3 versões no jornal}

%CRISTAL {"arestas":[{"alvo":"geopolitica","confianca":1.0,"epistemico":"fato","origem":"wiki","rel":"usa"},{"alvo":"determinismo_ambiental","confianca":1.0,"epistemico":"fato","origem":"wiki","rel":"usa"},{"alvo":"biopoder","confianca":1.0,"epistemico":"fato","origem":"wiki","rel":"relaciona"}],"confianca":1.0,"contraexemplos":[],"descricao":"Ratzel: o Estado como organismo vivo que precisa crescer territorialmente, sendo as fronteiras apenas parada temporária. ATENÇÃO: metáfora biogeográfica depois deturpada e cooptada pela ideologia nazista para justificar expansão e extermínio — a filiação científica não redime o uso letal.","epistemico":"inferencia","exemplos":[],"id":"lebensraum_ratzel","memoria":"semantica","meta":{"autor":"Ratzel","dominio":"ciencias_de_fora","fonte":""},"origem":"wiki","palavras_chave":[],"sinonimos":[],"tipo":"conceito","titulo":"Espaço Vital / Lebensraum (Ratzel)"}
\section{Espaço Vital / Lebensraum (Ratzel)}
Ratzel: o Estado como organismo vivo que precisa crescer territorialmente, sendo as fronteiras apenas parada temporária. ATENÇÃO: metáfora biogeográfica depois deturpada e cooptada pela ideologia nazista para justificar expansão e extermínio — a filiação científica não redime o uso letal.
\prov{relações: usa geopolitica; usa determinismo\_ambiental; relaciona biopoder}
\prov{proveniência: wiki \textperiodcentered{} inferencia \textperiodcentered{} confiança 1.0 \textperiodcentered{} domínio ciencias\_de\_fora \textperiodcentered{} história 4 versões no jornal}

%CRISTAL {"arestas":[{"alvo":"acao_social_weber","confianca":1.0,"epistemico":"fato","origem":"wiki","rel":"usa"},{"alvo":"ciencia_politica","confianca":1.0,"epistemico":"fato","origem":"wiki","rel":"parte_de"}],"confianca":1.0,"contraexemplos":[],"descricao":"Weber: a obediência estável se apoia em três fontes puras de legitimidade — tradicional (santidade dos costumes), carismática (devoção às qualidades de um líder) e legal-racional/burocrática (crença na legalidade de regras impessoais).","epistemico":"inferencia","exemplos":[],"id":"legitimidade_weber","memoria":"semantica","meta":{"autor":"Weber","dominio":"ciencias_de_fora","fonte":""},"origem":"wiki","palavras_chave":[],"sinonimos":[],"tipo":"conceito","titulo":"Tipos de Dominação Legítima (Weber)"}
\section{Tipos de Dominação Legítima (Weber)}
Weber: a obediência estável se apoia em três fontes puras de legitimidade — tradicional (santidade dos costumes), carismática (devoção às qualidades de um líder) e legal-racional/burocrática (crença na legalidade de regras impessoais).
\prov{relações: usa acao\_social\_weber; parte\_de ciencia\_politica}
\prov{proveniência: wiki \textperiodcentered{} inferencia \textperiodcentered{} confiança 1.0 \textperiodcentered{} domínio ciencias\_de\_fora \textperiodcentered{} história 3 versões no jornal}

%CRISTAL {"arestas":[{"alvo":"burocracia_weber","confianca":1.0,"epistemico":"fato","origem":"wiki","rel":"usa"},{"alvo":"logica_acao_coletiva","confianca":1.0,"epistemico":"fato","origem":"wiki","rel":"relaciona"}],"confianca":1.0,"contraexemplos":[],"descricao":"Michels: 'quem diz organização diz oligarquia' — toda organização de grande escala, mesmo os partidos democráticos, tende inevitavelmente a ser controlada por uma minoria de líderes profissionais, pelas exigências técnicas da própria organização.","epistemico":"fato","exemplos":[],"id":"lei_de_ferro_oligarquia","memoria":"semantica","meta":{"autor":"Michels","dominio":"ciencias_de_fora","fonte":""},"origem":"wiki","palavras_chave":[],"sinonimos":[],"tipo":"conceito","titulo":"Lei de Ferro da Oligarquia (Michels)"}
\section{Lei de Ferro da Oligarquia (Michels)}
Michels: 'quem diz organização diz oligarquia' — toda organização de grande escala, mesmo os partidos democráticos, tende inevitavelmente a ser controlada por uma minoria de líderes profissionais, pelas exigências técnicas da própria organização.
\prov{relações: usa burocracia\_weber; relaciona logica\_acao\_coletiva}
\prov{proveniência: wiki \textperiodcentered{} fato \textperiodcentered{} confiança 1.0 \textperiodcentered{} domínio ciencias\_de\_fora \textperiodcentered{} história 3 versões no jornal}

%CRISTAL {"arestas":[{"alvo":"deus_postulado","confianca":1.0,"epistemico":"fato","origem":"wiki","rel":"usa"},{"alvo":"word","confianca":0.5,"epistemico":"fato","origem":"wiki","rel":"relaciona"}],"confianca":1.0,"contraexemplos":[],"descricao":"Lopes: as 'leis divinas' variam por cultura porque o cérebro é programável na infância (software cristão/muçulmano/hindu) — não são verdades objetivas.","epistemico":"inferencia","exemplos":[],"id":"leis_de_deus_programacao","memoria":"semantica","meta":{"autor":"Gentil Lopes","dominio":"filosofia","fonte":""},"origem":"livro","palavras_chave":[],"sinonimos":[],"tipo":"conceito","titulo":"Leis de Deus = Programação"}
\section{Leis de Deus = Programação}
Lopes: as 'leis divinas' variam por cultura porque o cérebro é programável na infância (software cristão/muçulmano/hindu) — não são verdades objetivas.
\prov{relações: usa deus\_postulado; relaciona word}
\prov{proveniência: livro \textperiodcentered{} inferencia \textperiodcentered{} confiança 1.0 \textperiodcentered{} domínio filosofia \textperiodcentered{} história 4 versões no jornal}

%CRISTAL {"arestas":[{"alvo":"gregor_mendel","confianca":1.0,"epistemico":"fato","origem":"wiki","rel":"parte_de"},{"alvo":"sintese_moderna","confianca":1.0,"epistemico":"fato","origem":"wiki","rel":"usa"}],"confianca":1.0,"contraexemplos":[],"descricao":"A herança é particulada: fatores discretos (genes) segregam e se combinam por regras, não se misturam como líquidos.","epistemico":"fato","exemplos":[],"id":"leis_de_mendel","memoria":"semantica","meta":{"autor":"Mendel","dominio":"ciencias_de_fora","fonte":""},"origem":"wiki","palavras_chave":[],"sinonimos":[],"tipo":"conceito","titulo":"Leis de Mendel"}
\section{Leis de Mendel}
A herança é particulada: fatores discretos (genes) segregam e se combinam por regras, não se misturam como líquidos.
\prov{relações: parte\_de gregor\_mendel; usa sintese\_moderna}
\prov{proveniência: wiki \textperiodcentered{} fato \textperiodcentered{} confiança 1.0 \textperiodcentered{} domínio ciencias\_de\_fora \textperiodcentered{} história 3 versões no jornal}

%CRISTAL {"arestas":[{"alvo":"filosofia_do_direito","confianca":1.0,"epistemico":"fato","origem":"wiki","rel":"parte_de"},{"alvo":"convencao_vs_verdade_gentil","confianca":1.0,"epistemico":"fato","origem":"wiki","rel":"usa"},{"alvo":"jusnaturalismo","confianca":1.0,"epistemico":"fato","origem":"wiki","rel":"contradiz"}],"confianca":1.0,"contraexemplos":[],"descricao":"Lopes estende sua tese ao conhecimento 'duro': as 'leis' que a ciência descreve são propriedades do nosso mapa conceitual, não da realidade (subscrevendo Russell: 'muitas coisas que considerávamos leis naturais são puras convenções humanas') — mas ressalva que não são arbitrárias.","epistemico":"inferencia","exemplos":[],"id":"leis_naturais_convencao_gentil","memoria":"semantica","meta":{"autor":"Gentil Lopes","dominio":"ciencias_de_fora","fonte":""},"origem":"wiki","palavras_chave":[],"sinonimos":[],"tipo":"conceito","titulo":"As Leis Naturais são Convenções (Lopes)"}
\section{As Leis Naturais são Convenções (Lopes)}
Lopes estende sua tese ao conhecimento 'duro': as 'leis' que a ciência descreve são propriedades do nosso mapa conceitual, não da realidade (subscrevendo Russell: 'muitas coisas que considerávamos leis naturais são puras convenções humanas') — mas ressalva que não são arbitrárias.
\prov{relações: parte\_de filosofia\_do\_direito; usa convencao\_vs\_verdade\_gentil; contradiz jusnaturalismo}
\prov{proveniência: wiki \textperiodcentered{} inferencia \textperiodcentered{} confiança 1.0 \textperiodcentered{} domínio ciencias\_de\_fora \textperiodcentered{} história 5 versões no jornal}

%CRISTAL {"arestas":[{"alvo":"direito_e_tecnologia","confianca":1.0,"epistemico":"fato","origem":"wiki","rel":"parte_de"},{"alvo":"code_is_law","confianca":1.0,"epistemico":"fato","origem":"wiki","rel":"usa"},{"alvo":"goldchain","confianca":1.0,"epistemico":"fato","origem":"wiki","rel":"explica"},{"alvo":"criptoeconomia","confianca":1.0,"epistemico":"fato","origem":"wiki","rel":"relaciona"}],"confianca":1.0,"contraexemplos":[],"descricao":"Wright & De Filippi: um conjunto de regras administradas por contratos autoexecutáveis e organizações descentralizadas, aplicadas pelo próprio protocolo — 'ordem sem lei', governança por código e não por tribunal.","epistemico":"inferencia","exemplos":[],"id":"lex_cryptographica","memoria":"semantica","meta":{"autor":"Wright & De Filippi","dominio":"ciencias_de_fora","fonte":""},"origem":"wiki","palavras_chave":[],"sinonimos":[],"tipo":"conceito","titulo":"Lex Cryptographica"}
\section{Lex Cryptographica}
Wright \& De Filippi: um conjunto de regras administradas por contratos autoexecutáveis e organizações descentralizadas, aplicadas pelo próprio protocolo — 'ordem sem lei', governança por código e não por tribunal.
\prov{relações: parte\_de direito\_e\_tecnologia; usa code\_is\_law; explica goldchain; relaciona criptoeconomia}
\prov{proveniência: wiki \textperiodcentered{} inferencia \textperiodcentered{} confiança 1.0 \textperiodcentered{} domínio ciencias\_de\_fora \textperiodcentered{} história 5 versões no jornal}

%CRISTAL {"arestas":[{"alvo":"relacoes_internacionais","confianca":1.0,"epistemico":"fato","origem":"wiki","rel":"parte_de"},{"alvo":"realismo_ri","confianca":1.0,"epistemico":"fato","origem":"wiki","rel":"contradiz"},{"alvo":"mao_invisivel","confianca":1.0,"epistemico":"fato","origem":"wiki","rel":"relaciona"},{"alvo":"governanca_global","confianca":1.0,"epistemico":"fato","origem":"wiki","rel":"explica"}],"confianca":1.0,"contraexemplos":[],"descricao":"Keohane & Nye: a cooperação entre Estados é possível e sustentável mesmo sob anarquia — viabilizada por instituições internacionais, interdependência econômica e regimes democráticos.","epistemico":"inferencia","exemplos":[],"id":"liberalismo_ri","memoria":"semantica","meta":{"autor":"Keohane & Nye","dominio":"ciencias_de_fora","fonte":""},"origem":"wiki","palavras_chave":[],"sinonimos":[],"tipo":"conceito","titulo":"Liberalismo / Institucionalismo nas RI"}
\section{Liberalismo / Institucionalismo nas RI}
Keohane \& Nye: a cooperação entre Estados é possível e sustentável mesmo sob anarquia — viabilizada por instituições internacionais, interdependência econômica e regimes democráticos.
\prov{relações: parte\_de relacoes\_internacionais; contradiz realismo\_ri; relaciona mao\_invisivel; explica governanca\_global}
\prov{proveniência: wiki \textperiodcentered{} inferencia \textperiodcentered{} confiança 1.0 \textperiodcentered{} domínio ciencias\_de\_fora \textperiodcentered{} história 5 versões no jornal}

%CRISTAL {"arestas":[{"alvo":"politica","confianca":1.0,"epistemico":"fato","origem":"wiki","rel":"parte_de"},{"alvo":"word","confianca":0.5,"epistemico":"fato","origem":"wiki","rel":"relaciona"},{"alvo":"misticismo_nao_dual","confianca":0.5,"epistemico":"fato","origem":"wiki","rel":"relaciona"}],"confianca":1.0,"contraexemplos":[],"descricao":"Berlin: liberdade negativa (ausência de constrangimento) vs positiva (autodomínio, capacidade) — muitas vezes em conflito.","epistemico":"inferencia","exemplos":[],"id":"liberdade_berlin","memoria":"semantica","meta":{"autor":"Berlin","dominio":"ciencias_de_fora","fonte":""},"origem":"wiki","palavras_chave":[],"sinonimos":[],"tipo":"conceito","titulo":"Liberdade Positiva e Negativa"}
\section{Liberdade Positiva e Negativa}
Berlin: liberdade negativa (ausência de constrangimento) vs positiva (autodomínio, capacidade) — muitas vezes em conflito.
\prov{relações: parte\_de politica; relaciona word; relaciona misticismo\_nao\_dual}
\prov{proveniência: wiki \textperiodcentered{} inferencia \textperiodcentered{} confiança 1.0 \textperiodcentered{} domínio ciencias\_de\_fora \textperiodcentered{} história 4 versões no jornal}

%CRISTAL {"arestas":[{"alvo":"filosofia_da_linguagem","confianca":1.0,"epistemico":"fato","origem":"wiki","rel":"parte_de"}],"confianca":1.0,"contraexemplos":[],"descricao":"Wittgenstein (Tractatus): 'os limites da minha linguagem são os limites do meu mundo' — o dizível e o que só se mostra.","epistemico":"inferencia","exemplos":[],"id":"limites_da_linguagem","memoria":"semantica","meta":{"autor":"Wittgenstein","dominio":"ciencias_de_fora","fonte":""},"origem":"wiki","palavras_chave":[],"sinonimos":[],"tipo":"conceito","titulo":"Os Limites da Linguagem"}
\section{Os Limites da Linguagem}
Wittgenstein (Tractatus): 'os limites da minha linguagem são os limites do meu mundo' — o dizível e o que só se mostra.
\prov{relações: parte\_de filosofia\_da\_linguagem}
\prov{proveniência: wiki \textperiodcentered{} inferencia \textperiodcentered{} confiança 1.0 \textperiodcentered{} domínio ciencias\_de\_fora \textperiodcentered{} história 2 versões no jornal}

%CRISTAL {"arestas":[{"alvo":"sistema_terra","confianca":1.0,"epistemico":"fato","origem":"wiki","rel":"usa"},{"alvo":"antropoceno","confianca":1.0,"epistemico":"fato","origem":"wiki","rel":"relaciona"},{"alvo":"vida","confianca":1.0,"epistemico":"fato","origem":"wiki","rel":"relaciona"}],"confianca":1.0,"contraexemplos":[],"descricao":"Rockström: nove fronteiras biofísicas do sistema Terra (clima, biosfera, ciclos de N e P, uso do solo, água doce, acidificação oceânica, ozônio, aerossóis, novas entidades) que delimitam um 'espaço operacional seguro' para a humanidade — seis já ultrapassados.","epistemico":"inferencia","exemplos":[],"id":"limites_planetarios","memoria":"semantica","meta":{"autor":"Rockström","dominio":"ciencias_de_fora","fonte":""},"origem":"wiki","palavras_chave":[],"sinonimos":[],"tipo":"conceito","titulo":"Limites Planetários (Rockström)"}
\section{Limites Planetários (Rockström)}
Rockström: nove fronteiras biofísicas do sistema Terra (clima, biosfera, ciclos de N e P, uso do solo, água doce, acidificação oceânica, ozônio, aerossóis, novas entidades) que delimitam um 'espaço operacional seguro' para a humanidade — seis já ultrapassados.
\prov{relações: usa sistema\_terra; relaciona antropoceno; relaciona vida}
\prov{proveniência: wiki \textperiodcentered{} inferencia \textperiodcentered{} confiança 1.0 \textperiodcentered{} domínio ciencias\_de\_fora \textperiodcentered{} história 4 versões no jornal}

%CRISTAL {"arestas":[{"alvo":"idioma","confianca":1.0,"epistemico":"fato","origem":"wiki","rel":"especializa"},{"alvo":"arbitrariedade_do_signo","confianca":1.0,"epistemico":"fato","origem":"wiki","rel":"usa"}],"confianca":1.0,"contraexemplos":[],"descricao":"Idioma inventado de propósito — auxiliar (Esperanto, de Zamenhof) ou artístico (Klingon, Élfico). A prova de que a convenção pode ser projetada.","epistemico":"fato","exemplos":[],"id":"lingua_construida","memoria":"semantica","meta":{"autor":"Zamenhof","dominio":"ciencias_de_fora","fonte":""},"origem":"wiki","palavras_chave":[],"sinonimos":[],"tipo":"conceito","titulo":"Língua Construída"}
\section{Língua Construída}
Idioma inventado de propósito — auxiliar (Esperanto, de Zamenhof) ou artístico (Klingon, Élfico). A prova de que a convenção pode ser projetada.
\prov{relações: especializa idioma; usa arbitrariedade\_do\_signo}
\prov{proveniência: wiki \textperiodcentered{} fato \textperiodcentered{} confiança 1.0 \textperiodcentered{} domínio ciencias\_de\_fora \textperiodcentered{} história 3 versões no jornal}

%CRISTAL {"arestas":[{"alvo":"idioma","confianca":1.0,"epistemico":"fato","origem":"wiki","rel":"especializa"},{"alvo":"arbitrariedade_do_signo","confianca":1.0,"epistemico":"fato","origem":"wiki","rel":"usa"}],"confianca":1.0,"contraexemplos":[],"descricao":"Línguas completas na modalidade visual-gestual (Libras, ASL), com gramática própria — prova de que a linguagem não depende do som.","epistemico":"fato","exemplos":[],"id":"lingua_de_sinais","memoria":"semantica","meta":{"dominio":"ciencias_de_fora","fonte":""},"origem":"wiki","palavras_chave":[],"sinonimos":[],"tipo":"conceito","titulo":"Língua de Sinais"}
\section{Língua de Sinais}
Línguas completas na modalidade visual-gestual (Libras, ASL), com gramática própria — prova de que a linguagem não depende do som.
\prov{relações: especializa idioma; usa arbitrariedade\_do\_signo}
\prov{proveniência: wiki \textperiodcentered{} fato \textperiodcentered{} confiança 1.0 \textperiodcentered{} domínio ciencias\_de\_fora \textperiodcentered{} história 3 versões no jornal}

%CRISTAL {"arestas":[{"alvo":"ferdinand_saussure","confianca":1.0,"epistemico":"fato","origem":"wiki","rel":"parte_de"}],"confianca":1.0,"contraexemplos":[],"descricao":"Saussure: a LÍNGUA é o sistema social compartilhado; a FALA é o ato individual. A ciência estuda o sistema.","epistemico":"fato","exemplos":[],"id":"lingua_e_fala","memoria":"semantica","meta":{"autor":"Saussure","dominio":"ciencias_de_fora","fonte":""},"origem":"wiki","palavras_chave":[],"sinonimos":[],"tipo":"conceito","titulo":"Língua e Fala (langue/parole)"}
\section{Língua e Fala (langue/parole)}
Saussure: a LÍNGUA é o sistema social compartilhado; a FALA é o ato individual. A ciência estuda o sistema.
\prov{relações: parte\_de ferdinand\_saussure}
\prov{proveniência: wiki \textperiodcentered{} fato \textperiodcentered{} confiança 1.0 \textperiodcentered{} domínio ciencias\_de\_fora \textperiodcentered{} história 2 versões no jornal}

%CRISTAL {"arestas":[{"alvo":"idioma","confianca":1.0,"epistemico":"fato","origem":"wiki","rel":"especializa"}],"confianca":1.0,"contraexemplos":[],"descricao":"A língua adotada para comunicação entre falantes de idiomas diferentes — hoje o inglês; antes o latim, o francês, o árabe. A ponte por convenção.","epistemico":"fato","exemplos":[],"id":"lingua_franca","memoria":"semantica","meta":{"dominio":"ciencias_de_fora","fonte":""},"origem":"wiki","palavras_chave":[],"sinonimos":[],"tipo":"conceito","titulo":"Língua Franca"}
\section{Língua Franca}
A língua adotada para comunicação entre falantes de idiomas diferentes — hoje o inglês; antes o latim, o francês, o árabe. A ponte por convenção.
\prov{relações: especializa idioma}
\prov{proveniência: wiki \textperiodcentered{} fato \textperiodcentered{} confiança 1.0 \textperiodcentered{} domínio ciencias\_de\_fora \textperiodcentered{} história 2 versões no jornal}

%CRISTAL {"arestas":[{"alvo":"idioma","confianca":1.0,"epistemico":"fato","origem":"wiki","rel":"usa"},{"alvo":"gramatica_universal","confianca":1.0,"epistemico":"fato","origem":"wiki","rel":"relaciona"}],"confianca":1.0,"contraexemplos":[],"descricao":"A primeira língua adquirida naturalmente na infância — a que molda a intuição do que 'soa certo' e a percepção do mundo.","epistemico":"fato","exemplos":[],"id":"lingua_materna","memoria":"semantica","meta":{"dominio":"ciencias_de_fora","fonte":""},"origem":"wiki","palavras_chave":[],"sinonimos":[],"tipo":"conceito","titulo":"Língua Materna"}
\section{Língua Materna}
A primeira língua adquirida naturalmente na infância — a que molda a intuição do que 'soa certo' e a percepção do mundo.
\prov{relações: usa idioma; relaciona gramatica\_universal}
\prov{proveniência: wiki \textperiodcentered{} fato \textperiodcentered{} confiança 1.0 \textperiodcentered{} domínio ciencias\_de\_fora \textperiodcentered{} história 3 versões no jornal}

%CRISTAL {"arestas":[{"alvo":"sintaxe","confianca":1.0,"epistemico":"fato","origem":"wiki","rel":"usa"}],"confianca":1.0,"contraexemplos":[],"descricao":"Línguas que omitem o pronome-sujeito porque o verbo já o marca ('falo' = eu falo) — português, espanhol, latim. O inglês não pode: precisa do 'I'.","epistemico":"fato","exemplos":[],"id":"lingua_pro_drop","memoria":"semantica","meta":{"dominio":"ciencias_de_fora","fonte":""},"origem":"wiki","palavras_chave":[],"sinonimos":[],"tipo":"conceito","titulo":"Pro-drop (Sujeito Nulo)"}
\section{Pro-drop (Sujeito Nulo)}
Línguas que omitem o pronome-sujeito porque o verbo já o marca ('falo' = eu falo) — português, espanhol, latim. O inglês não pode: precisa do 'I'.
\prov{relações: usa sintaxe}
\prov{proveniência: wiki \textperiodcentered{} fato \textperiodcentered{} confiança 1.0 \textperiodcentered{} domínio ciencias\_de\_fora \textperiodcentered{} história 2 versões no jornal}

%CRISTAL {"arestas":[{"alvo":"mudanca_linguistica","confianca":1.0,"epistemico":"fato","origem":"wiki","rel":"usa"}],"confianca":1.0,"contraexemplos":[],"descricao":"Viva: tem falantes nativos e muda. Morta (latim): sem falantes nativos, fixa — mas pode viver como litúrgica, científica ou ancestral.","epistemico":"inferencia","exemplos":[],"id":"lingua_viva_morta","memoria":"semantica","meta":{"dominio":"ciencias_de_fora","fonte":""},"origem":"wiki","palavras_chave":[],"sinonimos":[],"tipo":"conceito","titulo":"Língua Viva e Morta"}
\section{Língua Viva e Morta}
Viva: tem falantes nativos e muda. Morta (latim): sem falantes nativos, fixa — mas pode viver como litúrgica, científica ou ancestral.
\prov{relações: usa mudanca\_linguistica}
\prov{proveniência: wiki \textperiodcentered{} inferencia \textperiodcentered{} confiança 1.0 \textperiodcentered{} domínio ciencias\_de\_fora \textperiodcentered{} história 2 versões no jornal}

%CRISTAL {"arestas":[{"alvo":"idioma","confianca":1.0,"epistemico":"fato","origem":"wiki","rel":"usa"},{"alvo":"lingua_viva_morta","confianca":1.0,"epistemico":"fato","origem":"wiki","rel":"relaciona"}],"confianca":1.0,"contraexemplos":[],"descricao":"As centenas de línguas originárias das Américas (e do mundo), muitas ameaçadas — cada uma um modo inteiro e insubstituível de recortar o mundo.","epistemico":"inferencia","exemplos":[],"id":"linguas_indigenas","memoria":"semantica","meta":{"dominio":"ciencias_de_fora","fonte":""},"origem":"wiki","palavras_chave":[],"sinonimos":[],"tipo":"conceito","titulo":"Línguas Indígenas"}
\section{Línguas Indígenas}
As centenas de línguas originárias das Américas (e do mundo), muitas ameaçadas — cada uma um modo inteiro e insubstituível de recortar o mundo.
\prov{relações: usa idioma; relaciona lingua\_viva\_morta}
\prov{proveniência: wiki \textperiodcentered{} inferencia \textperiodcentered{} confiança 1.0 \textperiodcentered{} domínio ciencias\_de\_fora \textperiodcentered{} história 3 versões no jornal}

%CRISTAL {"arestas":[{"alvo":"indo_europeu","confianca":1.0,"epistemico":"fato","origem":"wiki","rel":"parte_de"},{"alvo":"latim","confianca":1.0,"epistemico":"fato","origem":"wiki","rel":"usa"}],"confianca":1.0,"contraexemplos":[],"descricao":"As filhas do latim falado: português, espanhol, francês, italiano, romeno — a mesma raiz derivando por mil e quinhentos anos.","epistemico":"fato","exemplos":[],"id":"linguas_romanicas","memoria":"semantica","meta":{"dominio":"ciencias_de_fora","fonte":""},"origem":"wiki","palavras_chave":[],"sinonimos":[],"tipo":"conceito","titulo":"Línguas Românicas"}
\section{Línguas Românicas}
As filhas do latim falado: português, espanhol, francês, italiano, romeno — a mesma raiz derivando por mil e quinhentos anos.
\prov{relações: parte\_de indo\_europeu; usa latim}
\prov{proveniência: wiki \textperiodcentered{} fato \textperiodcentered{} confiança 1.0 \textperiodcentered{} domínio ciencias\_de\_fora \textperiodcentered{} história 3 versões no jornal}

%CRISTAL {"arestas":[{"alvo":"semiotica","confianca":1.0,"epistemico":"fato","origem":"wiki","rel":"relaciona"}],"confianca":1.0,"contraexemplos":[],"descricao":"A ciência da linguagem humana — som, forma, sentido e uso; descreve, não prescreve.","epistemico":"fato","exemplos":[],"id":"linguistica","memoria":"semantica","meta":{"dominio":"ciencias_de_fora","fonte":""},"origem":"wiki","palavras_chave":[],"sinonimos":[],"tipo":"conceito","titulo":"Linguística"}
\section{Linguística}
A ciência da linguagem humana — som, forma, sentido e uso; descreve, não prescreve.
\prov{relações: relaciona semiotica}
\prov{proveniência: wiki \textperiodcentered{} fato \textperiodcentered{} confiança 1.0 \textperiodcentered{} domínio ciencias\_de\_fora \textperiodcentered{} história 2 versões no jornal}

%CRISTAL {"arestas":[{"alvo":"linguistica","confianca":1.0,"epistemico":"fato","origem":"wiki","rel":"parte_de"},{"alvo":"tiffany","confianca":1.0,"epistemico":"fato","origem":"wiki","rel":"explica"},{"alvo":"linguagem_matriz_gatos","confianca":1.0,"epistemico":"fato","origem":"wiki","rel":"relaciona"}],"confianca":1.0,"contraexemplos":[],"descricao":"Modelar a língua com algoritmos — analisar, gerar e recuperar sentido por máquina. O campo em que a Tiffany opera.","epistemico":"fato","exemplos":[],"id":"linguistica_computacional","memoria":"semantica","meta":{"dominio":"ciencias_de_fora","fonte":""},"origem":"wiki","palavras_chave":[],"sinonimos":[],"tipo":"conceito","titulo":"Linguística Computacional"}
\section{Linguística Computacional}
Modelar a língua com algoritmos — analisar, gerar e recuperar sentido por máquina. O campo em que a Tiffany opera.
\prov{relações: parte\_de linguistica; explica tiffany; relaciona linguagem\_matriz\_gatos}
\prov{proveniência: wiki \textperiodcentered{} fato \textperiodcentered{} confiança 1.0 \textperiodcentered{} domínio ciencias\_de\_fora \textperiodcentered{} história 4 versões no jornal}

%CRISTAL {"arestas":[{"alvo":"sagrado","confianca":1.0,"epistemico":"fato","origem":"wiki","rel":"parte_de"},{"alvo":"monoteismo_etico","confianca":1.0,"epistemico":"fato","origem":"wiki","rel":"usa"}],"confianca":1.0,"contraexemplos":[],"descricao":"As três religiões do Livro — judaísmo, cristianismo e islã — como releituras sucessivas da fé de Abraão num Deus único, universal e moralmente exigente; herança compartilhada com ruptura e reivindicação de completude.","epistemico":"inferencia","exemplos":[],"id":"linhagem_abraamica","memoria":"semantica","meta":{"autor":"—","dominio":"sagrado","fonte":""},"origem":"wiki","palavras_chave":[],"sinonimos":[],"tipo":"conceito","titulo":"A Linhagem Abraâmica"}
\section{A Linhagem Abraâmica}
As três religiões do Livro — judaísmo, cristianismo e islã — como releituras sucessivas da fé de Abraão num Deus único, universal e moralmente exigente; herança compartilhada com ruptura e reivindicação de completude.
\prov{relações: parte\_de sagrado; usa monoteismo\_etico}
\prov{proveniência: wiki \textperiodcentered{} inferencia \textperiodcentered{} confiança 1.0 \textperiodcentered{} domínio sagrado \textperiodcentered{} história 3 versões no jornal}

%CRISTAL {"arestas":[{"alvo":"literatura","confianca":1.0,"epistemico":"fato","origem":"wiki","rel":"parte_de"},{"alvo":"estruturalismo_linguistico","confianca":1.0,"epistemico":"fato","origem":"wiki","rel":"usa"}],"confianca":1.0,"contraexemplos":[],"descricao":"Jakobson: o que faz de um enunciado arte verbal — quando a mensagem se volta para si mesma (função poética), pela seleção e combinação que geram plurissignificação e ambiguidade.","epistemico":"inferencia","exemplos":[],"id":"literariedade","memoria":"semantica","meta":{"autor":"Jakobson","dominio":"ciencias_de_fora","fonte":""},"origem":"wiki","palavras_chave":[],"sinonimos":[],"tipo":"conceito","titulo":"Literariedade / Função Poética"}
\section{Literariedade / Função Poética}
Jakobson: o que faz de um enunciado arte verbal — quando a mensagem se volta para si mesma (função poética), pela seleção e combinação que geram plurissignificação e ambiguidade.
\prov{relações: parte\_de literatura; usa estruturalismo\_linguistico}
\prov{proveniência: wiki \textperiodcentered{} inferencia \textperiodcentered{} confiança 1.0 \textperiodcentered{} domínio ciencias\_de\_fora \textperiodcentered{} história 3 versões no jornal}

%CRISTAL {"arestas":[{"alvo":"artes","confianca":1.0,"epistemico":"fato","origem":"wiki","rel":"parte_de"},{"alvo":"telemetria_shadow_producao","confianca":0.5,"epistemico":"fato","origem":"wiki","rel":"relaciona"},{"alvo":"readme","confianca":0.5,"epistemico":"fato","origem":"wiki","rel":"relaciona"},{"alvo":"sintaxe","confianca":0.5,"epistemico":"fato","origem":"wiki","rel":"relaciona"},{"alvo":"ula_da_broca_arquitetura_minima","confianca":0.5,"epistemico":"fato","origem":"wiki","rel":"relaciona"},{"alvo":"virtualizacao","confianca":0.5,"epistemico":"fato","origem":"wiki","rel":"relaciona"},{"alvo":"vida-morte","confianca":0.5,"epistemico":"fato","origem":"wiki","rel":"relaciona"},{"alvo":"arte","confianca":1.0,"epistemico":"fato","origem":"wiki","rel":"parte_de"}],"confianca":1.0,"contraexemplos":[],"descricao":"A arte da palavra — a ficção, a poesia e a narrativa; a linguagem voltada para si mesma que cria mundos e renova a percepção.","epistemico":"fato","exemplos":[],"id":"literatura","memoria":"semantica","meta":{"dominio":"ciencias_de_fora","fonte":""},"origem":"wiki","palavras_chave":[],"sinonimos":[],"tipo":"conceito","titulo":"Literatura"}
\section{Literatura}
A arte da palavra — a ficção, a poesia e a narrativa; a linguagem voltada para si mesma que cria mundos e renova a percepção.
\prov{relações: parte\_de artes; relaciona telemetria\_shadow\_producao; relaciona readme; relaciona sintaxe; relaciona ula\_da\_broca\_arquitetura\_minima; relaciona virtualizacao; relaciona vida-morte; parte\_de arte}
\prov{proveniência: wiki \textperiodcentered{} fato \textperiodcentered{} confiança 1.0 \textperiodcentered{} domínio ciencias\_de\_fora \textperiodcentered{} história 9 versões no jornal}

%CRISTAL {"arestas":[{"alvo":"traducao","confianca":1.0,"epistemico":"fato","origem":"wiki","rel":"especializa"}],"confianca":1.0,"contraexemplos":[],"descricao":"Adaptar não só a língua, mas moeda, data, formato e cultura ao público-alvo — traduzir o produto inteiro, não só o texto. A tradução como engenharia.","epistemico":"inferencia","exemplos":[],"id":"localizacao","memoria":"semantica","meta":{"dominio":"ciencias_de_fora","fonte":""},"origem":"wiki","palavras_chave":[],"sinonimos":[],"tipo":"conceito","titulo":"Localização (l10n)"}
\section{Localização (l10n)}
Adaptar não só a língua, mas moeda, data, formato e cultura ao público-alvo — traduzir o produto inteiro, não só o texto. A tradução como engenharia.
\prov{relações: especializa traducao}
\prov{proveniência: wiki \textperiodcentered{} inferencia \textperiodcentered{} confiança 1.0 \textperiodcentered{} domínio ciencias\_de\_fora \textperiodcentered{} história 2 versões no jornal}

%CRISTAL {"arestas":[{"alvo":"tragedia_dos_comuns","confianca":1.0,"epistemico":"fato","origem":"wiki","rel":"relaciona"},{"alvo":"teoria_dos_jogos","confianca":1.0,"epistemico":"fato","origem":"wiki","rel":"usa"},{"alvo":"capital_social_putnam","confianca":1.0,"epistemico":"fato","origem":"wiki","rel":"relaciona"}],"confianca":1.0,"contraexemplos":[],"descricao":"Olson: indivíduos racionais não contribuem espontaneamente para bens coletivos de que se beneficiariam mesmo sem contribuir (o 'carona'/free-rider); grupos grandes só cooperam com incentivos seletivos ou coerção.","epistemico":"inferencia","exemplos":[],"id":"logica_acao_coletiva","memoria":"semantica","meta":{"autor":"Olson","dominio":"ciencias_de_fora","fonte":""},"origem":"wiki","palavras_chave":[],"sinonimos":[],"tipo":"conceito","titulo":"A Lógica da Ação Coletiva (Olson)"}
\section{A Lógica da Ação Coletiva (Olson)}
Olson: indivíduos racionais não contribuem espontaneamente para bens coletivos de que se beneficiariam mesmo sem contribuir (o 'carona'/free-rider); grupos grandes só cooperam com incentivos seletivos ou coerção.
\prov{relações: relaciona tragedia\_dos\_comuns; usa teoria\_dos\_jogos; relaciona capital\_social\_putnam}
\prov{proveniência: wiki \textperiodcentered{} inferencia \textperiodcentered{} confiança 1.0 \textperiodcentered{} domínio ciencias\_de\_fora \textperiodcentered{} história 4 versões no jornal}

%CRISTAL {"arestas":[{"alvo":"logica_matematica","confianca":1.0,"epistemico":"fato","origem":"wiki","rel":"parte_de"},{"alvo":"metodo_axiomatico","confianca":1.0,"epistemico":"fato","origem":"wiki","rel":"usa"}],"confianca":1.0,"contraexemplos":[],"descricao":"Frege: a linguagem formal de quantificadores e relações que tornou a demonstração um cálculo mecânico.","epistemico":"fato","exemplos":[],"id":"logica_de_predicados","memoria":"semantica","meta":{"autor":"Frege","dominio":"ciencias_de_fora","fonte":""},"origem":"wiki","palavras_chave":[],"sinonimos":[],"tipo":"conceito","titulo":"Lógica de Predicados"}
\section{Lógica de Predicados}
Frege: a linguagem formal de quantificadores e relações que tornou a demonstração um cálculo mecânico.
\prov{relações: parte\_de logica\_matematica; usa metodo\_axiomatico}
\prov{proveniência: wiki \textperiodcentered{} fato \textperiodcentered{} confiança 1.0 \textperiodcentered{} domínio ciencias\_de\_fora \textperiodcentered{} história 6 versões no jornal}

%CRISTAL {"arestas":[{"alvo":"logica","confianca":1.0,"epistemico":"fato","origem":"wiki","rel":"especializa"},{"alvo":"matematica","confianca":1.0,"epistemico":"fato","origem":"wiki","rel":"parte_de"}],"confianca":1.0,"contraexemplos":[],"descricao":"O estudo formal da demonstração, da verdade e da consistência — a matemática dobrada sobre si mesma.","epistemico":"fato","exemplos":[],"id":"logica_matematica","memoria":"semantica","meta":{"dominio":"ciencias_de_fora","fonte":""},"origem":"wiki","palavras_chave":[],"sinonimos":[],"tipo":"conceito","titulo":"Lógica Matemática"}
\section{Lógica Matemática}
O estudo formal da demonstração, da verdade e da consistência — a matemática dobrada sobre si mesma.
\prov{relações: especializa logica; parte\_de matematica}
\prov{proveniência: wiki \textperiodcentered{} fato \textperiodcentered{} confiança 1.0 \textperiodcentered{} domínio ciencias\_de\_fora \textperiodcentered{} história 6 versões no jornal}

%CRISTAL {"arestas":[{"alvo":"filosofia_da_matematica","confianca":1.0,"epistemico":"fato","origem":"wiki","rel":"parte_de"},{"alvo":"nao_dualidade","confianca":0.5,"epistemico":"fato","origem":"wiki","rel":"relaciona"},{"alvo":"vetores","confianca":0.5,"epistemico":"fato","origem":"wiki","rel":"relaciona"},{"alvo":"word","confianca":0.5,"epistemico":"fato","origem":"wiki","rel":"relaciona"},{"alvo":"virtualizacao","confianca":0.5,"epistemico":"fato","origem":"wiki","rel":"relaciona"},{"alvo":"vontade_de_poder","confianca":0.5,"epistemico":"fato","origem":"wiki","rel":"relaciona"}],"confianca":1.0,"contraexemplos":[],"descricao":"A matemática é redutível à lógica pura (Frege, Russell); abalado pelo paradoxo de Russell.","epistemico":"inferencia","exemplos":[],"id":"logicismo","memoria":"semantica","meta":{"autor":"Frege/Russell","dominio":"filosofia","fonte":""},"origem":"wiki","palavras_chave":[],"sinonimos":[],"tipo":"conceito","titulo":"Logicismo"}
\section{Logicismo}
A matemática é redutível à lógica pura (Frege, Russell); abalado pelo paradoxo de Russell.
\prov{relações: parte\_de filosofia\_da\_matematica; relaciona nao\_dualidade; relaciona vetores; relaciona word; relaciona virtualizacao; relaciona vontade\_de\_poder}
\prov{proveniência: wiki \textperiodcentered{} inferencia \textperiodcentered{} confiança 1.0 \textperiodcentered{} domínio filosofia \textperiodcentered{} história 7 versões no jornal}

%CRISTAL {"arestas":[{"alvo":"sistema_de_escrita","confianca":1.0,"epistemico":"fato","origem":"wiki","rel":"especializa"},{"alvo":"mandarim","confianca":1.0,"epistemico":"fato","origem":"wiki","rel":"explica"}],"confianca":1.0,"contraexemplos":[],"descricao":"Um sinal que representa uma palavra ou ideia inteira, não um som (os caracteres chineses) — milhares de sinais, independentes da pronúncia.","epistemico":"fato","exemplos":[],"id":"logograma","memoria":"semantica","meta":{"dominio":"ciencias_de_fora","fonte":""},"origem":"wiki","palavras_chave":[],"sinonimos":[],"tipo":"conceito","titulo":"Logograma / Ideograma"}
\section{Logograma / Ideograma}
Um sinal que representa uma palavra ou ideia inteira, não um som (os caracteres chineses) — milhares de sinais, independentes da pronúncia.
\prov{relações: especializa sistema\_de\_escrita; explica mandarim}
\prov{proveniência: wiki \textperiodcentered{} fato \textperiodcentered{} confiança 1.0 \textperiodcentered{} domínio ciencias\_de\_fora \textperiodcentered{} história 3 versões no jornal}

%CRISTAL {"arestas":[{"alvo":"cristianismo","confianca":1.0,"epistemico":"fato","origem":"wiki","rel":"parte_de"},{"alvo":"simbolo_vs_significado","confianca":1.0,"epistemico":"fato","origem":"wiki","rel":"usa"},{"alvo":"deus_criador","confianca":1.0,"epistemico":"fato","origem":"wiki","rel":"usa"},{"alvo":"alcorao_palavra_revelada","confianca":1.0,"epistemico":"fato","origem":"wiki","rel":"relaciona"}],"confianca":1.0,"contraexemplos":[],"descricao":"Bíblia (João 1:1): o Verbo (logos: palavra/razão) é o princípio criador divino, eterno e pessoal, pelo qual tudo foi feito — 'No princípio era o Verbo, e o Verbo era Deus; todas as coisas foram feitas por ele', e 'o Verbo se fez carne'. A Palavra que cria a realidade ao pronunciá-la: o significante que é, ele mesmo, princípio ontológico.","epistemico":"inferencia","exemplos":[],"id":"logos_verbo","memoria":"semantica","meta":{"autor":"Bíblia","dominio":"sagrado","fonte":""},"origem":"wiki","palavras_chave":[],"sinonimos":[],"tipo":"conceito","titulo":"O Logos / o Verbo (João 1)"}
\section{O Logos / o Verbo (João 1)}
Bíblia (João 1:1): o Verbo (logos: palavra/razão) é o princípio criador divino, eterno e pessoal, pelo qual tudo foi feito — 'No princípio era o Verbo, e o Verbo era Deus; todas as coisas foram feitas por ele', e 'o Verbo se fez carne'. A Palavra que cria a realidade ao pronunciá-la: o significante que é, ele mesmo, princípio ontológico.
\prov{relações: parte\_de cristianismo; usa simbolo\_vs\_significado; usa deus\_criador; relaciona alcorao\_palavra\_revelada}
\prov{proveniência: wiki \textperiodcentered{} inferencia \textperiodcentered{} confiança 1.0 \textperiodcentered{} domínio sagrado \textperiodcentered{} história 5 versões no jornal}

%CRISTAL {"arestas":[{"alvo":"historia","confianca":1.0,"epistemico":"fato","origem":"wiki","rel":"parte_de"},{"alvo":"word","confianca":0.5,"epistemico":"fato","origem":"wiki","rel":"relaciona"},{"alvo":"vontade_de_poder","confianca":0.5,"epistemico":"fato","origem":"wiki","rel":"relaciona"}],"confianca":1.0,"contraexemplos":[],"descricao":"Braudel/Annales: sob os eventos (espuma) estão a conjuntura e a longue durée (clima, geografia) — a história como estrutura profunda, quase científica.","epistemico":"fato","exemplos":[],"id":"longue_duree","memoria":"semantica","meta":{"autor":"Braudel","dominio":"ciencias_de_fora","fonte":""},"origem":"wiki","palavras_chave":[],"sinonimos":[],"tipo":"conceito","titulo":"Longue Durée (Braudel)"}
\section{Longue Durée (Braudel)}
Braudel/Annales: sob os eventos (espuma) estão a conjuntura e a longue durée (clima, geografia) — a história como estrutura profunda, quase científica.
\prov{relações: parte\_de historia; relaciona word; relaciona vontade\_de\_poder}
\prov{proveniência: wiki \textperiodcentered{} fato \textperiodcentered{} confiança 1.0 \textperiodcentered{} domínio ciencias\_de\_fora \textperiodcentered{} história 4 versões no jornal}

%CRISTAL {"arestas":[{"alvo":"urbanizacao","confianca":1.0,"epistemico":"fato","origem":"wiki","rel":"explica"},{"alvo":"mao_invisivel","confianca":1.0,"epistemico":"fato","origem":"wiki","rel":"relaciona"}],"confianca":1.0,"contraexemplos":[],"descricao":"Christaller: um modelo que explica a hierarquia, o tamanho e a distribuição das cidades como centros de oferta de bens e serviços, organizados idealmente em padrão hexagonal segundo o alcance e o limiar dos serviços.","epistemico":"inferencia","exemplos":[],"id":"lugares_centrais","memoria":"semantica","meta":{"autor":"Christaller","dominio":"ciencias_de_fora","fonte":""},"origem":"wiki","palavras_chave":[],"sinonimos":[],"tipo":"conceito","titulo":"Teoria dos Lugares Centrais (Christaller)"}
\section{Teoria dos Lugares Centrais (Christaller)}
Christaller: um modelo que explica a hierarquia, o tamanho e a distribuição das cidades como centros de oferta de bens e serviços, organizados idealmente em padrão hexagonal segundo o alcance e o limiar dos serviços.
\prov{relações: explica urbanizacao; relaciona mao\_invisivel}
\prov{proveniência: wiki \textperiodcentered{} inferencia \textperiodcentered{} confiança 1.0 \textperiodcentered{} domínio ciencias\_de\_fora \textperiodcentered{} história 3 versões no jornal}

%CRISTAL {"arestas":[{"alvo":"filosofia_do_direito","confianca":1.0,"epistemico":"fato","origem":"wiki","rel":"parte_de"},{"alvo":"autopoiese","confianca":1.0,"epistemico":"fato","origem":"wiki","rel":"usa"},{"alvo":"godel_incompletude","confianca":1.0,"epistemico":"fato","origem":"wiki","rel":"relaciona"}],"confianca":1.0,"contraexemplos":[],"descricao":"Luhmann: o direito é um subsistema social autorreferente e operacionalmente fechado, que se reproduz processando comunicações pelo seu código binário próprio lícito/ilícito — autônomo diante de política, economia e moral.","epistemico":"inferencia","exemplos":[],"id":"luhmann_direito_autopoietico","memoria":"semantica","meta":{"autor":"Luhmann","dominio":"ciencias_de_fora","fonte":""},"origem":"wiki","palavras_chave":[],"sinonimos":[],"tipo":"conceito","titulo":"Direito como Sistema Autopoiético (Luhmann)"}
\section{Direito como Sistema Autopoiético (Luhmann)}
Luhmann: o direito é um subsistema social autorreferente e operacionalmente fechado, que se reproduz processando comunicações pelo seu código binário próprio lícito/ilícito — autônomo diante de política, economia e moral.
\prov{relações: parte\_de filosofia\_do\_direito; usa autopoiese; relaciona godel\_incompletude}
\prov{proveniência: wiki \textperiodcentered{} inferencia \textperiodcentered{} confiança 1.0 \textperiodcentered{} domínio ciencias\_de\_fora \textperiodcentered{} história 4 versões no jornal}

%CRISTAL {"arestas":[{"alvo":"valor_trabalho","confianca":1.0,"epistemico":"fato","origem":"wiki","rel":"usa"},{"alvo":"classe_social","confianca":1.0,"epistemico":"fato","origem":"wiki","rel":"explica"},{"alvo":"word","confianca":0.5,"epistemico":"fato","origem":"wiki","rel":"relaciona"},{"alvo":"vontade_de_poder","confianca":0.5,"epistemico":"fato","origem":"wiki","rel":"relaciona"}],"confianca":1.0,"contraexemplos":[],"descricao":"Marx: o capital extrai a diferença entre o valor produzido e o salário; o fetichismo faz relações sociais parecerem qualidades das coisas.","epistemico":"inferencia","exemplos":[],"id":"mais_valia","memoria":"semantica","meta":{"autor":"Marx","dominio":"ciencias_de_fora","fonte":""},"origem":"wiki","palavras_chave":[],"sinonimos":[],"tipo":"conceito","titulo":"Mais-Valia e Fetichismo"}
\section{Mais-Valia e Fetichismo}
Marx: o capital extrai a diferença entre o valor produzido e o salário; o fetichismo faz relações sociais parecerem qualidades das coisas.
\prov{relações: usa valor\_trabalho; explica classe\_social; relaciona word; relaciona vontade\_de\_poder}
\prov{proveniência: wiki \textperiodcentered{} inferencia \textperiodcentered{} confiança 1.0 \textperiodcentered{} domínio ciencias\_de\_fora \textperiodcentered{} história 5 versões no jornal}

%CRISTAL {"arestas":[{"alvo":"idioma","confianca":1.0,"epistemico":"fato","origem":"wiki","rel":"especializa"}],"confianca":1.0,"contraexemplos":[],"descricao":"A língua com mais falantes nativos do mundo (~920 milhões) — tonal, escrita em logogramas milenares independentes da pronúncia.","epistemico":"fato","exemplos":[],"id":"mandarim","memoria":"semantica","meta":{"dominio":"ciencias_de_fora","fonte":""},"origem":"wiki","palavras_chave":[],"sinonimos":[],"tipo":"conceito","titulo":"Mandarim (Chinês)"}
\section{Mandarim (Chinês)}
A língua com mais falantes nativos do mundo (\~{}920 milhões) — tonal, escrita em logogramas milenares independentes da pronúncia.
\prov{relações: especializa idioma}
\prov{proveniência: wiki \textperiodcentered{} fato \textperiodcentered{} confiança 1.0 \textperiodcentered{} domínio ciencias\_de\_fora \textperiodcentered{} história 2 versões no jornal}

%CRISTAL {"arestas":[{"alvo":"economia","confianca":1.0,"epistemico":"fato","origem":"wiki","rel":"parte_de"},{"alvo":"ordem_espontanea","confianca":1.0,"epistemico":"fato","origem":"wiki","rel":"relaciona"},{"alvo":"word","confianca":0.5,"epistemico":"fato","origem":"wiki","rel":"relaciona"}],"confianca":1.0,"contraexemplos":[],"descricao":"Smith: indivíduos buscando o próprio interesse promovem, sem intenção, o bem coletivo; a divisão do trabalho multiplica a produtividade.","epistemico":"inferencia","exemplos":[],"id":"mao_invisivel","memoria":"semantica","meta":{"autor":"Smith","dominio":"ciencias_de_fora","fonte":""},"origem":"wiki","palavras_chave":[],"sinonimos":[],"tipo":"conceito","titulo":"A Mão Invisível"}
\section{A Mão Invisível}
Smith: indivíduos buscando o próprio interesse promovem, sem intenção, o bem coletivo; a divisão do trabalho multiplica a produtividade.
\prov{relações: parte\_de economia; relaciona ordem\_espontanea; relaciona word}
\prov{proveniência: wiki \textperiodcentered{} inferencia \textperiodcentered{} confiança 1.0 \textperiodcentered{} domínio ciencias\_de\_fora \textperiodcentered{} história 4 versões no jornal}

%CRISTAL {"arestas":[{"alvo":"mapa_nao_territorio","confianca":1.0,"epistemico":"fato","origem":"wiki","rel":"usa"},{"alvo":"simbolo_vs_significado","confianca":1.0,"epistemico":"fato","origem":"wiki","rel":"usa"},{"alvo":"convencao_vs_verdade_gentil","confianca":1.0,"epistemico":"fato","origem":"wiki","rel":"usa"}],"confianca":1.0,"contraexemplos":[],"descricao":"Lopes (via Capra/Russell): as 'leis' da física e da matemática 'são criações da mente, propriedades do nosso mapa conceitual da realidade, e não da própria realidade'; a notação não é o ente (0,999… não é o número 1). O mapa conceitual não é o território — mas, para ele, tampouco é arbitrário.","epistemico":"inferencia","exemplos":[],"id":"mapa_conceitual_gentil","memoria":"semantica","meta":{"autor":"Gentil Lopes","dominio":"ciencias_de_fora","fonte":""},"origem":"wiki","palavras_chave":[],"sinonimos":[],"tipo":"conceito","titulo":"O Mapa Conceitual não é a Realidade (Lopes)"}
\section{O Mapa Conceitual não é a Realidade (Lopes)}
Lopes (via Capra/Russell): as 'leis' da física e da matemática 'são criações da mente, propriedades do nosso mapa conceitual da realidade, e não da própria realidade'; a notação não é o ente (0,999… não é o número 1). O mapa conceitual não é o território — mas, para ele, tampouco é arbitrário.
\prov{relações: usa mapa\_nao\_territorio; usa simbolo\_vs\_significado; usa convencao\_vs\_verdade\_gentil}
\prov{proveniência: wiki \textperiodcentered{} inferencia \textperiodcentered{} confiança 1.0 \textperiodcentered{} domínio ciencias\_de\_fora \textperiodcentered{} história 5 versões no jornal}

%CRISTAL {"arestas":[{"alvo":"geografia","confianca":1.0,"epistemico":"fato","origem":"wiki","rel":"parte_de"},{"alvo":"simulacro_simulacao","confianca":1.0,"epistemico":"fato","origem":"wiki","rel":"relaciona"},{"alvo":"simbolo_vs_significado","confianca":1.0,"epistemico":"fato","origem":"wiki","rel":"relaciona"}],"confianca":1.0,"contraexemplos":[],"descricao":"Korzybski: a representação de algo (mapa, palavra, modelo) nunca é a coisa representada; toda abstração é seletiva e filtrada. O fundamento da semântica geral — e o parente epistemológico do simulacro (o modelo tomado pela realidade).","epistemico":"inferencia","exemplos":[],"id":"mapa_nao_territorio","memoria":"semantica","meta":{"autor":"Korzybski","dominio":"ciencias_de_fora","fonte":""},"origem":"wiki","palavras_chave":[],"sinonimos":[],"tipo":"conceito","titulo":"O Mapa Não é o Território (Korzybski)"}
\section{O Mapa Não é o Território (Korzybski)}
Korzybski: a representação de algo (mapa, palavra, modelo) nunca é a coisa representada; toda abstração é seletiva e filtrada. O fundamento da semântica geral — e o parente epistemológico do simulacro (o modelo tomado pela realidade).
\prov{relações: parte\_de geografia; relaciona simulacro\_simulacao; relaciona simbolo\_vs\_significado}
\prov{proveniência: wiki \textperiodcentered{} inferencia \textperiodcentered{} confiança 1.0 \textperiodcentered{} domínio ciencias\_de\_fora \textperiodcentered{} história 4 versões no jornal}

%CRISTAL {"arestas":[{"alvo":"logica","confianca":1.0,"epistemico":"fato","origem":"curriculo","rel":"depende"},{"alvo":"computabilidade","confianca":1.0,"epistemico":"fato","origem":"wiki","rel":"parte_de"},{"alvo":"expressividade_sub_turing","confianca":1.0,"epistemico":"fato","origem":"wiki","rel":"explica"},{"alvo":"gato","confianca":1.0,"epistemico":"fato","origem":"wiki","rel":"relaciona"},{"alvo":"turing_halting_church","confianca":1.0,"epistemico":"fato","origem":"wiki","rel":"relaciona"}],"confianca":1.0,"contraexemplos":[],"descricao":"Turing (1936): o modelo abstrato de todo cálculo mecânico — fita, cabeça, estados. A definição de 'algoritmo'.","epistemico":"fato","exemplos":[],"id":"maquina_de_turing","memoria":"semantica","meta":{"autor":"Turing","dominio":"ciencias_de_fora","fonte":""},"origem":"wiki","palavras_chave":["fita","decidibilidade","church-turing"],"sinonimos":["máquina de Turing","MT"],"tipo":"conceito","titulo":"Máquina de Turing"}
\section{Máquina de Turing}
Turing (1936): o modelo abstrato de todo cálculo mecânico — fita, cabeça, estados. A definição de 'algoritmo'.
\prov{sinónimos: máquina de Turing; MT}
\prov{relações: depende logica; parte\_de computabilidade; explica expressividade\_sub\_turing; relaciona gato; relaciona turing\_halting\_church}
\prov{proveniência: wiki \textperiodcentered{} fato \textperiodcentered{} confiança 1.0 \textperiodcentered{} domínio ciencias\_de\_fora \textperiodcentered{} história 88 versões no jornal}

%CRISTAL {"arestas":[{"alvo":"filosofia_da_tecnica","confianca":1.0,"epistemico":"fato","origem":"wiki","rel":"parte_de"},{"alvo":"ellul_la_technique","confianca":1.0,"epistemico":"fato","origem":"wiki","rel":"relaciona"},{"alvo":"acao_social_weber","confianca":1.0,"epistemico":"fato","origem":"wiki","rel":"usa"},{"alvo":"biopoder","confianca":1.0,"epistemico":"fato","origem":"wiki","rel":"relaciona"}],"confianca":1.0,"contraexemplos":[],"descricao":"Marcuse: na sociedade industrial avançada a racionalidade técnico-instrumental vira hegemônica e se converte em dominação política — absorve a crítica, produz o 'pensamento unidimensional' e apresenta como neutra uma ordem que é ideológica.","epistemico":"inferencia","exemplos":[],"id":"marcuse_racionalidade_tecnologica","memoria":"semantica","meta":{"autor":"Marcuse","dominio":"ciencias_de_fora","fonte":""},"origem":"wiki","palavras_chave":[],"sinonimos":[],"tipo":"conceito","titulo":"Racionalidade Tecnológica como Dominação (Marcuse)"}
\section{Racionalidade Tecnológica como Dominação (Marcuse)}
Marcuse: na sociedade industrial avançada a racionalidade técnico-instrumental vira hegemônica e se converte em dominação política — absorve a crítica, produz o 'pensamento unidimensional' e apresenta como neutra uma ordem que é ideológica.
\prov{relações: parte\_de filosofia\_da\_tecnica; relaciona ellul\_la\_technique; usa acao\_social\_weber; relaciona biopoder}
\prov{proveniência: wiki \textperiodcentered{} inferencia \textperiodcentered{} confiança 1.0 \textperiodcentered{} domínio ciencias\_de\_fora \textperiodcentered{} história 5 versões no jornal}

%CRISTAL {"arestas":[{"alvo":"filosofia_da_tecnica","confianca":1.0,"epistemico":"fato","origem":"wiki","rel":"parte_de"},{"alvo":"mais_valia","confianca":1.0,"epistemico":"fato","origem":"wiki","rel":"usa"},{"alvo":"materialismo_historico","confianca":1.0,"epistemico":"fato","origem":"wiki","rel":"usa"},{"alvo":"mumford_megamaquina","confianca":1.0,"epistemico":"fato","origem":"wiki","rel":"relaciona"}],"confianca":1.0,"contraexemplos":[],"descricao":"Marx: na grande indústria a máquina é o capital materializado; a subsunção real subordina o trabalhador ao ritmo da máquina. No 'Fragmento sobre as Máquinas', o saber social se objetiva como 'general intellect' no capital fixo.","epistemico":"inferencia","exemplos":[],"id":"marx_maquinaria_subsuncao","memoria":"semantica","meta":{"autor":"Marx","dominio":"ciencias_de_fora","fonte":""},"origem":"wiki","palavras_chave":[],"sinonimos":[],"tipo":"conceito","titulo":"Maquinaria e Subsunção do Trabalho (Marx)"}
\section{Maquinaria e Subsunção do Trabalho (Marx)}
Marx: na grande indústria a máquina é o capital materializado; a subsunção real subordina o trabalhador ao ritmo da máquina. No 'Fragmento sobre as Máquinas', o saber social se objetiva como 'general intellect' no capital fixo.
\prov{relações: parte\_de filosofia\_da\_tecnica; usa mais\_valia; usa materialismo\_historico; relaciona mumford\_megamaquina}
\prov{proveniência: wiki \textperiodcentered{} inferencia \textperiodcentered{} confiança 1.0 \textperiodcentered{} domínio ciencias\_de\_fora \textperiodcentered{} história 5 versões no jornal}

%CRISTAL {"arestas":[{"alvo":"didatica_matematica","confianca":1.0,"epistemico":"fato","origem":"wiki","rel":"parte_de"},{"alvo":"beleza_matematica","confianca":1.0,"epistemico":"fato","origem":"wiki","rel":"referencia"},{"alvo":"tiffany","confianca":1.0,"epistemico":"fato","origem":"wiki","rel":"relaciona"}],"confianca":1.0,"contraexemplos":[],"descricao":"Lopes enfatiza aos alunos 'menos o aspecto utilitário da matemática, mas sobretudo sua vertente como arte e engenharia — tal como de fato ela é em sua essência' (abrindo com Hardy: o matemático como pintor que desenha com ideias). Vem de sua origem em engenharia elétrica.","epistemico":"inferencia","exemplos":[],"id":"matematica_arte_engenharia","memoria":"semantica","meta":{"autor":"Gentil Lopes","dominio":"ciencias_de_fora","fonte":""},"origem":"wiki","palavras_chave":[],"sinonimos":[],"tipo":"conceito","titulo":"Matemática como Arte e Engenharia (Lopes)"}
\section{Matemática como Arte e Engenharia (Lopes)}
Lopes enfatiza aos alunos 'menos o aspecto utilitário da matemática, mas sobretudo sua vertente como arte e engenharia — tal como de fato ela é em sua essência' (abrindo com Hardy: o matemático como pintor que desenha com ideias). Vem de sua origem em engenharia elétrica.
\prov{relações: parte\_de didatica\_matematica; referencia beleza\_matematica; relaciona tiffany}
\prov{proveniência: wiki \textperiodcentered{} inferencia \textperiodcentered{} confiança 1.0 \textperiodcentered{} domínio ciencias\_de\_fora \textperiodcentered{} história 5 versões no jornal}

%CRISTAL {"arestas":[{"alvo":"metafisica","confianca":1.0,"epistemico":"fato","origem":"wiki","rel":"parte_de"},{"alvo":"idealismo_berkeley","confianca":1.0,"epistemico":"fato","origem":"wiki","rel":"contradiz"},{"alvo":"word","confianca":0.5,"epistemico":"fato","origem":"wiki","rel":"relaciona"},{"alvo":"while","confianca":0.5,"epistemico":"fato","origem":"wiki","rel":"relaciona"},{"alvo":"vontade_de_poder","confianca":0.5,"epistemico":"fato","origem":"wiki","rel":"relaciona"}],"confianca":1.0,"contraexemplos":[],"descricao":"Tudo que existe é físico; a consciência é processo neural redutível em princípio. Posição hegemônica na neurociência.","epistemico":"inferencia","exemplos":[],"id":"materialismo_fisicalismo","memoria":"semantica","meta":{"dominio":"filosofia","fonte":"web: SEP"},"origem":"wiki","palavras_chave":[],"sinonimos":[],"tipo":"conceito","titulo":"Materialismo / Fisicalismo"}
\section{Materialismo / Fisicalismo}
Tudo que existe é físico; a consciência é processo neural redutível em princípio. Posição hegemônica na neurociência.
\prov{relações: parte\_de metafisica; contradiz idealismo\_berkeley; relaciona word; relaciona while; relaciona vontade\_de\_poder}
\prov{proveniência: wiki \textperiodcentered{} web: SEP \textperiodcentered{} inferencia \textperiodcentered{} confiança 1.0 \textperiodcentered{} domínio filosofia \textperiodcentered{} história 6 versões no jornal}

%CRISTAL {"arestas":[{"alvo":"dialetica_hegeliana","confianca":1.0,"epistemico":"fato","origem":"wiki","rel":"contradiz"},{"alvo":"materialismo_fisicalismo","confianca":1.0,"epistemico":"fato","origem":"wiki","rel":"usa"},{"alvo":"classe_social","confianca":1.0,"epistemico":"fato","origem":"wiki","rel":"usa"}],"confianca":1.0,"contraexemplos":[],"descricao":"Marx: a história é determinada pelo modo de produção e pela luta de classes; a superestrutura reflete a base econômica. Inverte Hegel.","epistemico":"hipotese","exemplos":[],"id":"materialismo_historico","memoria":"semantica","meta":{"autor":"Marx","dominio":"ciencias_de_fora","fonte":""},"origem":"wiki","palavras_chave":[],"sinonimos":[],"tipo":"conceito","titulo":"Materialismo Histórico"}
\section{Materialismo Histórico}
Marx: a história é determinada pelo modo de produção e pela luta de classes; a superestrutura reflete a base econômica. Inverte Hegel.
\prov{relações: contradiz dialetica\_hegeliana; usa materialismo\_fisicalismo; usa classe\_social}
\prov{proveniência: wiki \textperiodcentered{} hipotese \textperiodcentered{} confiança 1.0 \textperiodcentered{} domínio ciencias\_de\_fora \textperiodcentered{} história 4 versões no jornal}

%CRISTAL {"arestas":[{"alvo":"advaita_vedanta","confianca":1.0,"epistemico":"fato","origem":"wiki","rel":"parte_de"},{"alvo":"vida-morte","confianca":0.5,"epistemico":"fato","origem":"wiki","rel":"relaciona"},{"alvo":"word","confianca":0.5,"epistemico":"fato","origem":"wiki","rel":"relaciona"},{"alvo":"virtualizacao","confianca":0.5,"epistemico":"fato","origem":"wiki","rel":"relaciona"}],"confianca":1.0,"contraexemplos":[],"descricao":"O poder que manifesta a multiplicidade e vela a unidade de Brahman; a ilusão de separação que se dissolve no conhecimento.","epistemico":"hipotese","exemplos":[],"id":"maya_ilusao","memoria":"semantica","meta":{"dominio":"filosofia","fonte":"Vedanta"},"origem":"wiki","palavras_chave":[],"sinonimos":[],"tipo":"conceito","titulo":"Maya"}
\section{Maya}
O poder que manifesta a multiplicidade e vela a unidade de Brahman; a ilusão de separação que se dissolve no conhecimento.
\prov{relações: parte\_de advaita\_vedanta; relaciona vida-morte; relaciona word; relaciona virtualizacao}
\prov{proveniência: wiki \textperiodcentered{} Vedanta \textperiodcentered{} hipotese \textperiodcentered{} confiança 1.0 \textperiodcentered{} domínio filosofia \textperiodcentered{} história 5 versões no jornal}

%CRISTAL {"arestas":[{"alvo":"sociointeracionismo_vygotsky","confianca":1.0,"epistemico":"fato","origem":"wiki","rel":"usa"},{"alvo":"mente_computacao","confianca":1.0,"epistemico":"fato","origem":"wiki","rel":"referencia"},{"alvo":"simbolo_vs_significado","confianca":1.0,"epistemico":"fato","origem":"wiki","rel":"relaciona"}],"confianca":1.0,"contraexemplos":[],"descricao":"Vygotsky: o pensamento não age direto sobre o mundo, mas por instrumentos e signos culturais (linguagem, símbolos) que reorganizam as funções psíquicas.","epistemico":"inferencia","exemplos":[],"id":"mediacao_semiotica","memoria":"semantica","meta":{"autor":"Vygotsky","dominio":"ciencias_de_fora","fonte":""},"origem":"wiki","palavras_chave":[],"sinonimos":[],"tipo":"conceito","titulo":"Mediação (Signos e Instrumentos)"}
\section{Mediação (Signos e Instrumentos)}
Vygotsky: o pensamento não age direto sobre o mundo, mas por instrumentos e signos culturais (linguagem, símbolos) que reorganizam as funções psíquicas.
\prov{relações: usa sociointeracionismo\_vygotsky; referencia mente\_computacao; relaciona simbolo\_vs\_significado}
\prov{proveniência: wiki \textperiodcentered{} inferencia \textperiodcentered{} confiança 1.0 \textperiodcentered{} domínio ciencias\_de\_fora \textperiodcentered{} história 4 versões no jornal}

%CRISTAL {"arestas":[{"alvo":"relacoes_humano_tecnologia","confianca":1.0,"epistemico":"fato","origem":"wiki","rel":"usa"},{"alvo":"teoria_ator_rede","confianca":1.0,"epistemico":"fato","origem":"wiki","rel":"usa"},{"alvo":"artefatos_tem_politica","confianca":1.0,"epistemico":"fato","origem":"wiki","rel":"relaciona"}],"confianca":1.0,"contraexemplos":[],"descricao":"Verbeek: as tecnologias não são intermediárias neutras — co-constituem a percepção e a subjetividade moral (o ultrassom reconfigura o feto como paciente). Como já mediam a ação moral, deve-se projetar valores no design ('moralizar a tecnologia').","epistemico":"inferencia","exemplos":[],"id":"mediacao_tecnica_moral","memoria":"semantica","meta":{"autor":"Verbeek","dominio":"ciencias_de_fora","fonte":""},"origem":"wiki","palavras_chave":[],"sinonimos":[],"tipo":"conceito","titulo":"Mediação Técnica da Moral (Verbeek)"}
\section{Mediação Técnica da Moral (Verbeek)}
Verbeek: as tecnologias não são intermediárias neutras — co-constituem a percepção e a subjetividade moral (o ultrassom reconfigura o feto como paciente). Como já mediam a ação moral, deve-se projetar valores no design ('moralizar a tecnologia').
\prov{relações: usa relacoes\_humano\_tecnologia; usa teoria\_ator\_rede; relaciona artefatos\_tem\_politica}
\prov{proveniência: wiki \textperiodcentered{} inferencia \textperiodcentered{} confiança 1.0 \textperiodcentered{} domínio ciencias\_de\_fora \textperiodcentered{} história 4 versões no jornal}

%CRISTAL {"arestas":[{"alvo":"filosofia_da_medicina","confianca":1.0,"epistemico":"fato","origem":"wiki","rel":"parte_de"},{"alvo":"biopoder","confianca":1.0,"epistemico":"fato","origem":"wiki","rel":"usa"},{"alvo":"normativismo_doenca","confianca":1.0,"epistemico":"fato","origem":"wiki","rel":"usa"}],"confianca":1.0,"contraexemplos":[],"descricao":"Conrad: o processo pelo qual problemas não-médicos (envelhecimento, tristeza, parto, comportamento) passam a ser definidos e tratados como doenças, sob linguagem e intervenção médica. A chave é a definição — tornar médico o que era social, moral ou comum.","epistemico":"inferencia","exemplos":[],"id":"medicalizacao","memoria":"semantica","meta":{"autor":"Conrad","dominio":"ciencias_de_fora","fonte":""},"origem":"wiki","palavras_chave":[],"sinonimos":[],"tipo":"conceito","titulo":"Medicalização (Conrad)"}
\section{Medicalização (Conrad)}
Conrad: o processo pelo qual problemas não-médicos (envelhecimento, tristeza, parto, comportamento) passam a ser definidos e tratados como doenças, sob linguagem e intervenção médica. A chave é a definição — tornar médico o que era social, moral ou comum.
\prov{relações: parte\_de filosofia\_da\_medicina; usa biopoder; usa normativismo\_doenca}
\prov{proveniência: wiki \textperiodcentered{} inferencia \textperiodcentered{} confiança 1.0 \textperiodcentered{} domínio ciencias\_de\_fora \textperiodcentered{} história 4 versões no jornal}

%CRISTAL {"arestas":[{"alvo":"colapso_observador","confianca":1.0,"epistemico":"fato","origem":"wiki","rel":"usa"},{"alvo":"von_neumann_wigner","confianca":1.0,"epistemico":"fato","origem":"wiki","rel":"usa"},{"alvo":"equacao_divina_gentil","confianca":1.0,"epistemico":"fato","origem":"wiki","rel":"usa"},{"alvo":"teologia_quantica_gentil","confianca":1.0,"epistemico":"fato","origem":"wiki","rel":"parte_de"}],"confianca":1.0,"contraexemplos":[],"descricao":"Lopes: aplica o colapso quântico ao divino — 'Abstrato + Medição = Concreto'. Antes de ser 'medido' (observado por um cérebro) o Absoluto é o Vazio, sem existência definida; a observação o converte em 'Deus'. Calcado na complementaridade de Bohr e na consciência-que-mede (Wigner, von Neumann, Wheeler, Goswami).","epistemico":"inferencia","exemplos":[],"id":"medicao_colapsa_deus_gentil","memoria":"semantica","meta":{"autor":"Gentil Lopes","dominio":"sagrado","fonte":""},"origem":"wiki","palavras_chave":[],"sinonimos":[],"tipo":"conceito","titulo":"A Medição Colapsa Deus (Lopes)"}
\section{A Medição Colapsa Deus (Lopes)}
Lopes: aplica o colapso quântico ao divino — 'Abstrato + Medição = Concreto'. Antes de ser 'medido' (observado por um cérebro) o Absoluto é o Vazio, sem existência definida; a observação o converte em 'Deus'. Calcado na complementaridade de Bohr e na consciência-que-mede (Wigner, von Neumann, Wheeler, Goswami).
\prov{relações: usa colapso\_observador; usa von\_neumann\_wigner; usa equacao\_divina\_gentil; parte\_de teologia\_quantica\_gentil}
\prov{proveniência: wiki \textperiodcentered{} inferencia \textperiodcentered{} confiança 1.0 \textperiodcentered{} domínio sagrado \textperiodcentered{} história 6 versões no jornal}

%CRISTAL {"arestas":[{"alvo":"vida","confianca":1.0,"epistemico":"fato","origem":"wiki","rel":"relaciona"},{"alvo":"biopoder","confianca":1.0,"epistemico":"fato","origem":"wiki","rel":"relaciona"},{"alvo":"word","confianca":0.5,"epistemico":"fato","origem":"wiki","rel":"relaciona"}],"confianca":1.0,"contraexemplos":[],"descricao":"O campo do cuidar da vida — a arte e a ciência de compreender, prevenir e tratar a doença, e a reflexão sobre o que é saúde, adoecer e curar.","epistemico":"fato","exemplos":[],"id":"medicina","memoria":"semantica","meta":{"dominio":"ciencias_de_fora","fonte":""},"origem":"wiki","palavras_chave":[],"sinonimos":[],"tipo":"conceito","titulo":"Medicina e Saúde"}
\section{Medicina e Saúde}
O campo do cuidar da vida — a arte e a ciência de compreender, prevenir e tratar a doença, e a reflexão sobre o que é saúde, adoecer e curar.
\prov{relações: relaciona vida; relaciona biopoder; relaciona word}
\prov{proveniência: wiki \textperiodcentered{} fato \textperiodcentered{} confiança 1.0 \textperiodcentered{} domínio ciencias\_de\_fora \textperiodcentered{} história 4 versões no jornal}

%CRISTAL {"arestas":[{"alvo":"saude_publica","confianca":1.0,"epistemico":"fato","origem":"wiki","rel":"parte_de"},{"alvo":"ensaio_clinico_randomizado","confianca":1.0,"epistemico":"fato","origem":"wiki","rel":"usa"},{"alvo":"historicismo_popper","confianca":1.0,"epistemico":"fato","origem":"wiki","rel":"relaciona"}],"confianca":1.0,"contraexemplos":[],"descricao":"Sackett: o uso consciencioso da melhor evidência disponível, integrada à experiência clínica e aos valores do paciente — triangular as três fontes, não só 'seguir estudos'.","epistemico":"inferencia","exemplos":[],"id":"medicina_baseada_evidencias","memoria":"semantica","meta":{"autor":"Sackett","dominio":"ciencias_de_fora","fonte":""},"origem":"wiki","palavras_chave":[],"sinonimos":[],"tipo":"conceito","titulo":"Medicina Baseada em Evidências (MBE)"}
\section{Medicina Baseada em Evidências (MBE)}
Sackett: o uso consciencioso da melhor evidência disponível, integrada à experiência clínica e aos valores do paciente — triangular as três fontes, não só 'seguir estudos'.
\prov{relações: parte\_de saude\_publica; usa ensaio\_clinico\_randomizado; relaciona historicismo\_popper}
\prov{proveniência: wiki \textperiodcentered{} inferencia \textperiodcentered{} confiança 1.0 \textperiodcentered{} domínio ciencias\_de\_fora \textperiodcentered{} história 4 versões no jornal}

%CRISTAL {"arestas":[{"alvo":"dor_vs_sofrimento_gentil","confianca":1.0,"epistemico":"fato","origem":"wiki","rel":"usa"},{"alvo":"mente_programacao_cerebro_computador","confianca":1.0,"epistemico":"fato","origem":"wiki","rel":"usa"},{"alvo":"consciencia","confianca":1.0,"epistemico":"fato","origem":"wiki","rel":"relaciona"},{"alvo":"dsm_critica","confianca":1.0,"epistemico":"fato","origem":"wiki","rel":"relaciona"},{"alvo":"aparelho_psiquico","confianca":1.0,"epistemico":"fato","origem":"wiki","rel":"contradiz"}],"confianca":1.0,"contraexemplos":[],"descricao":"Lopes: 'a solução para a maioria dos transtornos causados pela mente é deletar o ego', feito pela meditação profunda que devolve à 'tela em branco' (Consciência Pura). A meditação é um instrumento para controlar a mente — e a felicidade é como um remédio: melhor curar-se de precisar dela.","epistemico":"inferencia","exemplos":[],"id":"meditacao_cura_da_mente_gentil","memoria":"semantica","meta":{"autor":"Gentil Lopes","dominio":"ciencias_de_fora","fonte":""},"origem":"wiki","palavras_chave":[],"sinonimos":[],"tipo":"conceito","titulo":"Meditação como Cura das Patologias da Mente (Lopes)"}
\section{Meditação como Cura das Patologias da Mente (Lopes)}
Lopes: 'a solução para a maioria dos transtornos causados pela mente é deletar o ego', feito pela meditação profunda que devolve à 'tela em branco' (Consciência Pura). A meditação é um instrumento para controlar a mente — e a felicidade é como um remédio: melhor curar-se de precisar dela.
\prov{relações: usa dor\_vs\_sofrimento\_gentil; usa mente\_programacao\_cerebro\_computador; relaciona consciencia; relaciona dsm\_critica; contradiz aparelho\_psiquico}
\prov{proveniência: wiki \textperiodcentered{} inferencia \textperiodcentered{} confiança 1.0 \textperiodcentered{} domínio ciencias\_de\_fora \textperiodcentered{} história 7 versões no jornal}

%CRISTAL {"arestas":[{"alvo":"gatekeeping","confianca":1.0,"epistemico":"fato","origem":"wiki","rel":"usa"},{"alvo":"autocomunicacao_de_massa","confianca":1.0,"epistemico":"fato","origem":"wiki","rel":"relaciona"},{"alvo":"deus_mentiroso","confianca":1.0,"epistemico":"fato","origem":"wiki","rel":"relaciona"},{"alvo":"comunicacao_iconoclasta_gentil","confianca":1.0,"epistemico":"fato","origem":"wiki","rel":"usa"}],"confianca":1.0,"contraexemplos":[],"descricao":"Rejeitado pelas revistas, Lopes migra para canais próprios (site da UFRR, Drive, blogs de terceiros) que contornam os gatekeepers e circulam 'até em outros Países'. Já o YouTube/Google, com 'milhares de provas' de 0,999…=1, são o oposto: a mídia de massa que amplifica e naturaliza o consenso equivocado.","epistemico":"inferencia","exemplos":[],"id":"meio_digital_contra_porteiros_gentil","memoria":"semantica","meta":{"autor":"Gentil Lopes","dominio":"ciencias_de_fora","fonte":""},"origem":"wiki","palavras_chave":[],"sinonimos":[],"tipo":"conceito","titulo":"O Meio Digital contra os Porteiros (Lopes)"}
\section{O Meio Digital contra os Porteiros (Lopes)}
Rejeitado pelas revistas, Lopes migra para canais próprios (site da UFRR, Drive, blogs de terceiros) que contornam os gatekeepers e circulam 'até em outros Países'. Já o YouTube/Google, com 'milhares de provas' de 0,999…=1, são o oposto: a mídia de massa que amplifica e naturaliza o consenso equivocado.
\prov{relações: usa gatekeeping; relaciona autocomunicacao\_de\_massa; relaciona deus\_mentiroso; usa comunicacao\_iconoclasta\_gentil}
\prov{proveniência: wiki \textperiodcentered{} inferencia \textperiodcentered{} confiança 1.0 \textperiodcentered{} domínio ciencias\_de\_fora \textperiodcentered{} história 6 versões no jornal}

%CRISTAL {"arestas":[{"alvo":"filosofia_da_tecnica","confianca":1.0,"epistemico":"fato","origem":"wiki","rel":"parte_de"},{"alvo":"mediacao_semiotica","confianca":1.0,"epistemico":"fato","origem":"wiki","rel":"relaciona"}],"confianca":1.0,"contraexemplos":[],"descricao":"McLuhan: o efeito social decisivo de um meio não está no conteúdo que veicula, mas na forma do próprio meio, que reestrutura escalas, ritmos e padrões de percepção. (Daí meios 'quentes', de alta definição e baixa participação, vs 'frios', que exigem preenchimento ativo.)","epistemico":"inferencia","exemplos":[],"id":"meio_e_mensagem","memoria":"semantica","meta":{"autor":"McLuhan","dominio":"ciencias_de_fora","fonte":""},"origem":"wiki","palavras_chave":[],"sinonimos":[],"tipo":"conceito","titulo":"O Meio é a Mensagem (McLuhan)"}
\section{O Meio é a Mensagem (McLuhan)}
McLuhan: o efeito social decisivo de um meio não está no conteúdo que veicula, mas na forma do próprio meio, que reestrutura escalas, ritmos e padrões de percepção. (Daí meios 'quentes', de alta definição e baixa participação, vs 'frios', que exigem preenchimento ativo.)
\prov{relações: parte\_de filosofia\_da\_tecnica; relaciona mediacao\_semiotica}
\prov{proveniência: wiki \textperiodcentered{} inferencia \textperiodcentered{} confiança 1.0 \textperiodcentered{} domínio ciencias\_de\_fora \textperiodcentered{} história 3 versões no jornal}

%CRISTAL {"arestas":[{"alvo":"espaco_santos","confianca":1.0,"epistemico":"fato","origem":"wiki","rel":"usa"},{"alvo":"sociedade_em_rede","confianca":1.0,"epistemico":"fato","origem":"wiki","rel":"relaciona"}],"confianca":1.0,"contraexemplos":[],"descricao":"Milton Santos: a fase atual do espaço em que técnica, ciência e informação se fundem na própria produção do território — o meio geográfico instrumentalizado pela ciência e pela informação a serviço da acumulação hegemônica.","epistemico":"inferencia","exemplos":[],"id":"meio_tecnico_cientifico_informacional","memoria":"semantica","meta":{"autor":"Milton Santos","dominio":"ciencias_de_fora","fonte":""},"origem":"wiki","palavras_chave":[],"sinonimos":[],"tipo":"conceito","titulo":"Meio Técnico-Científico-Informacional (Milton Santos)"}
\section{Meio Técnico-Científico-Informacional (Milton Santos)}
Milton Santos: a fase atual do espaço em que técnica, ciência e informação se fundem na própria produção do território — o meio geográfico instrumentalizado pela ciência e pela informação a serviço da acumulação hegemônica.
\prov{relações: usa espaco\_santos; relaciona sociedade\_em\_rede}
\prov{proveniência: wiki \textperiodcentered{} inferencia \textperiodcentered{} confiança 1.0 \textperiodcentered{} domínio ciencias\_de\_fora \textperiodcentered{} história 3 versões no jornal}

%CRISTAL {"arestas":[{"alvo":"psicologia","confianca":1.0,"epistemico":"fato","origem":"wiki","rel":"parte_de"},{"alvo":"tiffany","confianca":1.0,"epistemico":"fato","origem":"wiki","rel":"explica"},{"alvo":"simbolo_vs_significado","confianca":1.0,"epistemico":"fato","origem":"wiki","rel":"relaciona"}],"confianca":1.0,"contraexemplos":[],"descricao":"Psicologia cognitiva: a mente é software, o cérebro é hardware — input→codificação→armazenamento→output. Ecoa a Tiffany.","epistemico":"inferencia","exemplos":[],"id":"mente_computacao","memoria":"semantica","meta":{"autor":"Neisser","dominio":"ciencias_de_fora","fonte":""},"origem":"wiki","palavras_chave":[],"sinonimos":[],"tipo":"conceito","titulo":"Mente como Processamento"}
\section{Mente como Processamento}
Psicologia cognitiva: a mente é software, o cérebro é hardware — input\ensuremath{\to}codificação\ensuremath{\to}armazenamento\ensuremath{\to}output. Ecoa a Tiffany.
\prov{relações: parte\_de psicologia; explica tiffany; relaciona simbolo\_vs\_significado}
\prov{proveniência: wiki \textperiodcentered{} inferencia \textperiodcentered{} confiança 1.0 \textperiodcentered{} domínio ciencias\_de\_fora \textperiodcentered{} história 4 versões no jornal}

%CRISTAL {"arestas":[{"alvo":"inteligencia_consciencia_relacional","confianca":1.0,"epistemico":"fato","origem":"wiki","rel":"usa"},{"alvo":"mente_computacao","confianca":1.0,"epistemico":"fato","origem":"wiki","rel":"relaciona"},{"alvo":"funcionalismo_mente","confianca":1.0,"epistemico":"fato","origem":"wiki","rel":"relaciona"}],"confianca":1.0,"contraexemplos":[],"descricao":"Lopes distingue Consciência (inata — 'um recém-nascido possui apenas Consciência') de mente (adquirida): a mente 'é produto de uma programação em seu cérebro, que é um computador biológico', feita pela sociedade. O conteúdo mental é software manipulável.","epistemico":"inferencia","exemplos":[],"id":"mente_programacao_cerebro_computador","memoria":"semantica","meta":{"autor":"Gentil Lopes","dominio":"ciencias_de_fora","fonte":""},"origem":"wiki","palavras_chave":[],"sinonimos":[],"tipo":"conceito","titulo":"Cérebro-Computador, Mente-Programação (Lopes)"}
\section{Cérebro-Computador, Mente-Programação (Lopes)}
Lopes distingue Consciência (inata — 'um recém-nascido possui apenas Consciência') de mente (adquirida): a mente 'é produto de uma programação em seu cérebro, que é um computador biológico', feita pela sociedade. O conteúdo mental é software manipulável.
\prov{relações: usa inteligencia\_consciencia\_relacional; relaciona mente\_computacao; relaciona funcionalismo\_mente}
\prov{proveniência: wiki \textperiodcentered{} inferencia \textperiodcentered{} confiança 1.0 \textperiodcentered{} domínio ciencias\_de\_fora \textperiodcentered{} história 5 versões no jornal}

%CRISTAL {"arestas":[{"alvo":"filosofia_da_educacao","confianca":1.0,"epistemico":"fato","origem":"wiki","rel":"parte_de"},{"alvo":"educacao_problematizadora","confianca":1.0,"epistemico":"fato","origem":"wiki","rel":"relaciona"},{"alvo":"educacao_bancaria","confianca":1.0,"epistemico":"fato","origem":"wiki","rel":"contradiz"}],"confianca":1.0,"contraexemplos":[],"descricao":"Rancière (a partir de Jacotot): um mestre pode emancipar sem transmitir seu próprio saber; o que ele oferece é a confiança na igual capacidade de aprender. O mestre explicador embrutece; o emancipador liberta.","epistemico":"inferencia","exemplos":[],"id":"mestre_ignorante_ranciere","memoria":"semantica","meta":{"autor":"Rancière","dominio":"ciencias_de_fora","fonte":""},"origem":"wiki","palavras_chave":[],"sinonimos":[],"tipo":"conceito","titulo":"O Mestre Ignorante (Rancière)"}
\section{O Mestre Ignorante (Rancière)}
Rancière (a partir de Jacotot): um mestre pode emancipar sem transmitir seu próprio saber; o que ele oferece é a confiança na igual capacidade de aprender. O mestre explicador embrutece; o emancipador liberta.
\prov{relações: parte\_de filosofia\_da\_educacao; relaciona educacao\_problematizadora; contradiz educacao\_bancaria}
\prov{proveniência: wiki \textperiodcentered{} inferencia \textperiodcentered{} confiança 1.0 \textperiodcentered{} domínio ciencias\_de\_fora \textperiodcentered{} história 4 versões no jornal}

%CRISTAL {"arestas":[{"alvo":"bioquimica","confianca":1.0,"epistemico":"fato","origem":"wiki","rel":"parte_de"},{"alvo":"o_que_e_vida","confianca":1.0,"epistemico":"fato","origem":"wiki","rel":"usa"},{"alvo":"dissipacao_prigogine","confianca":1.0,"epistemico":"fato","origem":"wiki","rel":"usa"}],"confianca":1.0,"contraexemplos":[],"descricao":"A rede de reações que extrai energia e constrói matéria — o fluxo que mantém o organismo longe do equilíbrio.","epistemico":"fato","exemplos":[],"id":"metabolismo","memoria":"semantica","meta":{"dominio":"ciencias_de_fora","fonte":""},"origem":"wiki","palavras_chave":[],"sinonimos":[],"tipo":"conceito","titulo":"Metabolismo"}
\section{Metabolismo}
A rede de reações que extrai energia e constrói matéria — o fluxo que mantém o organismo longe do equilíbrio.
\prov{relações: parte\_de bioquimica; usa o\_que\_e\_vida; usa dissipacao\_prigogine}
\prov{proveniência: wiki \textperiodcentered{} fato \textperiodcentered{} confiança 1.0 \textperiodcentered{} domínio ciencias\_de\_fora \textperiodcentered{} história 6 versões no jornal}

%CRISTAL {"arestas":[{"alvo":"teorias_aprendizagem","confianca":1.0,"epistemico":"fato","origem":"wiki","rel":"parte_de"},{"alvo":"polya_heuristica","confianca":1.0,"epistemico":"fato","origem":"wiki","rel":"usa"},{"alvo":"dois_sistemas_kahneman","confianca":1.0,"epistemico":"fato","origem":"wiki","rel":"referencia"}],"confianca":1.0,"contraexemplos":[],"descricao":"Flavell: o conhecimento e a regulação sobre os próprios processos cognitivos — 'o monitoramento ativo e a consequente regulação' da cognição. Operacionaliza-se como autorregulação da aprendizagem (planejar-monitorar-avaliar).","epistemico":"inferencia","exemplos":[],"id":"metacognicao","memoria":"semantica","meta":{"autor":"Flavell","dominio":"ciencias_de_fora","fonte":""},"origem":"wiki","palavras_chave":[],"sinonimos":[],"tipo":"conceito","titulo":"Metacognição"}
\section{Metacognição}
Flavell: o conhecimento e a regulação sobre os próprios processos cognitivos — 'o monitoramento ativo e a consequente regulação' da cognição. Operacionaliza-se como autorregulação da aprendizagem (planejar-monitorar-avaliar).
\prov{relações: parte\_de teorias\_aprendizagem; usa polya\_heuristica; referencia dois\_sistemas\_kahneman}
\prov{proveniência: wiki \textperiodcentered{} inferencia \textperiodcentered{} confiança 1.0 \textperiodcentered{} domínio ciencias\_de\_fora \textperiodcentered{} história 4 versões no jornal}

%CRISTAL {"arestas":[{"alvo":"filosofia","confianca":1.0,"epistemico":"fato","origem":"wiki","rel":"parte_de"},{"alvo":"semiotica","confianca":0.5,"epistemico":"fato","origem":"wiki","rel":"relaciona"},{"alvo":"pell","confianca":0.5,"epistemico":"fato","origem":"wiki","rel":"relaciona"},{"alvo":"poder_disciplinar_foucault","confianca":0.5,"epistemico":"fato","origem":"wiki","rel":"relaciona"},{"alvo":"pipeline_de_ingestao","confianca":0.5,"epistemico":"fato","origem":"wiki","rel":"relaciona"}],"confianca":1.0,"contraexemplos":[],"descricao":"O estudo do que é real e fundamental — ser, substância, causa, o que existe além do físico.","epistemico":"fato","exemplos":[],"id":"metafisica","memoria":"semantica","meta":{"dominio":"filosofia","fonte":""},"origem":"wiki","palavras_chave":[],"sinonimos":[],"tipo":"conceito","titulo":"Metafísica"}
\section{Metafísica}
O estudo do que é real e fundamental — ser, substância, causa, o que existe além do físico.
\prov{relações: parte\_de filosofia; relaciona semiotica; relaciona pell; relaciona poder\_disciplinar\_foucault; relaciona pipeline\_de\_ingestao}
\prov{proveniência: wiki \textperiodcentered{} fato \textperiodcentered{} confiança 1.0 \textperiodcentered{} domínio filosofia \textperiodcentered{} história 6 versões no jornal}

%CRISTAL {"arestas":[{"alvo":"roman_jakobson","confianca":1.0,"epistemico":"fato","origem":"wiki","rel":"parte_de"},{"alvo":"sintagma_paradigma","confianca":1.0,"epistemico":"fato","origem":"wiki","rel":"usa"}],"confianca":1.0,"contraexemplos":[],"descricao":"Jakobson: os dois polos do sentido — metáfora (similaridade, o paradigma) e metonímia (contiguidade, o sintagma).","epistemico":"inferencia","exemplos":[],"id":"metafora_metonimia","memoria":"semantica","meta":{"autor":"Jakobson","dominio":"ciencias_de_fora","fonte":""},"origem":"wiki","palavras_chave":[],"sinonimos":[],"tipo":"conceito","titulo":"Metáfora e Metonímia"}
\section{Metáfora e Metonímia}
Jakobson: os dois polos do sentido — metáfora (similaridade, o paradigma) e metonímia (contiguidade, o sintagma).
\prov{relações: parte\_de roman\_jakobson; usa sintagma\_paradigma}
\prov{proveniência: wiki \textperiodcentered{} inferencia \textperiodcentered{} confiança 1.0 \textperiodcentered{} domínio ciencias\_de\_fora \textperiodcentered{} história 3 versões no jornal}

%CRISTAL {"arestas":[{"alvo":"logica_matematica","confianca":1.0,"epistemico":"fato","origem":"wiki","rel":"parte_de"}],"confianca":1.0,"contraexemplos":[],"descricao":"Euclides: partir de axiomas evidentes e deduzir teoremas por regras — a espinha da matemática desde os Elementos.","epistemico":"fato","exemplos":[],"id":"metodo_axiomatico","memoria":"semantica","meta":{"autor":"Euclides","dominio":"ciencias_de_fora","fonte":""},"origem":"wiki","palavras_chave":[],"sinonimos":[],"tipo":"conceito","titulo":"Método Axiomático"}
\section{Método Axiomático}
Euclides: partir de axiomas evidentes e deduzir teoremas por regras — a espinha da matemática desde os Elementos.
\prov{relações: parte\_de logica\_matematica}
\prov{proveniência: wiki \textperiodcentered{} fato \textperiodcentered{} confiança 1.0 \textperiodcentered{} domínio ciencias\_de\_fora \textperiodcentered{} história 4 versões no jornal}

%CRISTAL {"arestas":[{"alvo":"teorias_aprendizagem","confianca":1.0,"epistemico":"fato","origem":"wiki","rel":"parte_de"},{"alvo":"construtivismo_piaget","confianca":1.0,"epistemico":"fato","origem":"wiki","rel":"relaciona"}],"confianca":1.0,"contraexemplos":[],"descricao":"Montessori: a criança aprende por autoeducação num 'ambiente preparado' — ordenado, à sua escala, com materiais concretos autocorretivos — que favorece autonomia e concentração.","epistemico":"inferencia","exemplos":[],"id":"metodo_montessori","memoria":"semantica","meta":{"autor":"Montessori","dominio":"ciencias_de_fora","fonte":""},"origem":"wiki","palavras_chave":[],"sinonimos":[],"tipo":"conceito","titulo":"Método Montessori"}
\section{Método Montessori}
Montessori: a criança aprende por autoeducação num 'ambiente preparado' — ordenado, à sua escala, com materiais concretos autocorretivos — que favorece autonomia e concentração.
\prov{relações: parte\_de teorias\_aprendizagem; relaciona construtivismo\_piaget}
\prov{proveniência: wiki \textperiodcentered{} inferencia \textperiodcentered{} confiança 1.0 \textperiodcentered{} domínio ciencias\_de\_fora \textperiodcentered{} história 3 versões no jornal}

%CRISTAL {"arestas":[{"alvo":"urbanizacao","confianca":1.0,"epistemico":"fato","origem":"wiki","rel":"usa"},{"alvo":"sociedade_em_rede","confianca":1.0,"epistemico":"fato","origem":"wiki","rel":"relaciona"}],"confianca":1.0,"contraexemplos":[],"descricao":"A metrópole é o polo que articula e comanda uma rede de cidades; a megacidade (>10 milhões) é sua forma extrema. A rede urbana é o sistema hierárquico de fluxos e conexões entre centros.","epistemico":"inferencia","exemplos":[],"id":"metropole_megacidade","memoria":"semantica","meta":{"autor":"—","dominio":"ciencias_de_fora","fonte":""},"origem":"wiki","palavras_chave":[],"sinonimos":[],"tipo":"conceito","titulo":"Metrópole, Megacidade e Rede Urbana"}
\section{Metrópole, Megacidade e Rede Urbana}
A metrópole é o polo que articula e comanda uma rede de cidades; a megacidade (>10 milhões) é sua forma extrema. A rede urbana é o sistema hierárquico de fluxos e conexões entre centros.
\prov{relações: usa urbanizacao; relaciona sociedade\_em\_rede}
\prov{proveniência: wiki \textperiodcentered{} inferencia \textperiodcentered{} confiança 1.0 \textperiodcentered{} domínio ciencias\_de\_fora \textperiodcentered{} história 3 versões no jornal}

%CRISTAL {"arestas":[{"alvo":"arte","confianca":1.0,"epistemico":"fato","origem":"wiki","rel":"explica"},{"alvo":"simbolo_vs_significado","confianca":1.0,"epistemico":"fato","origem":"wiki","rel":"usa"}],"confianca":1.0,"contraexemplos":[],"descricao":"A arte como representação do mundo — de Platão (cópia da cópia) a Aristóteles (imitação que revela o universal).","epistemico":"inferencia","exemplos":[],"id":"mimese","memoria":"semantica","meta":{"autor":"Aristóteles","dominio":"ciencias_de_fora","fonte":""},"origem":"wiki","palavras_chave":[],"sinonimos":[],"tipo":"conceito","titulo":"Mímese (Imitação)"}
\section{Mímese (Imitação)}
A arte como representação do mundo — de Platão (cópia da cópia) a Aristóteles (imitação que revela o universal).
\prov{relações: explica arte; usa simbolo\_vs\_significado}
\prov{proveniência: wiki \textperiodcentered{} inferencia \textperiodcentered{} confiança 1.0 \textperiodcentered{} domínio ciencias\_de\_fora \textperiodcentered{} história 5 versões no jornal}

%CRISTAL {"arestas":[{"alvo":"estetica","confianca":1.0,"epistemico":"fato","origem":"wiki","rel":"parte_de"},{"alvo":"word","confianca":0.5,"epistemico":"fato","origem":"wiki","rel":"relaciona"},{"alvo":"virtualizacao","confianca":0.5,"epistemico":"fato","origem":"wiki","rel":"relaciona"},{"alvo":"vida-morte","confianca":0.5,"epistemico":"fato","origem":"wiki","rel":"relaciona"},{"alvo":"sagrado_profano","confianca":0.5,"epistemico":"fato","origem":"wiki","rel":"relaciona"}],"confianca":1.0,"contraexemplos":[],"descricao":"Aristóteles: a arte é mímese (imitação criativa, poiesis); o modernismo a nega em favor da expressão pura — tensão nunca resolvida.","epistemico":"inferencia","exemplos":[],"id":"mimese_expressao","memoria":"semantica","meta":{"autor":"Aristóteles","dominio":"ciencias_de_fora","fonte":""},"origem":"wiki","palavras_chave":[],"sinonimos":[],"tipo":"conceito","titulo":"Mímese vs Expressão"}
\section{Mímese vs Expressão}
Aristóteles: a arte é mímese (imitação criativa, poiesis); o modernismo a nega em favor da expressão pura — tensão nunca resolvida.
\prov{relações: parte\_de estetica; relaciona word; relaciona virtualizacao; relaciona vida-morte; relaciona sagrado\_profano}
\prov{proveniência: wiki \textperiodcentered{} inferencia \textperiodcentered{} confiança 1.0 \textperiodcentered{} domínio ciencias\_de\_fora \textperiodcentered{} história 6 versões no jornal}

%CRISTAL {"arestas":[{"alvo":"autoridade_e_rejeicao_gentil","confianca":1.0,"epistemico":"fato","origem":"wiki","rel":"usa"},{"alvo":"espiral_do_silencio","confianca":1.0,"epistemico":"fato","origem":"wiki","rel":"relaciona"},{"alvo":"idiotismo_nao_examinar","confianca":1.0,"epistemico":"fato","origem":"wiki","rel":"relaciona"},{"alvo":"lei_de_ferro_oligarquia","confianca":1.0,"epistemico":"fato","origem":"wiki","rel":"relaciona"}],"confianca":1.0,"contraexemplos":[],"descricao":"Lopes: um erro pode ser sustentado por toda uma comunidade de especialistas ao mesmo tempo — o consenso da maioria não é garantia de verdade (ex.: 'os matemáticos não sabem o que é um número real'). O rebanho de doutores repete tolices 'porque uma autoridade o disse'.","epistemico":"inferencia","exemplos":[],"id":"miopia_coletiva","memoria":"semantica","meta":{"autor":"Gentil Lopes","dominio":"ciencias_de_fora","fonte":""},"origem":"wiki","palavras_chave":[],"sinonimos":[],"tipo":"conceito","titulo":"Miopia Coletiva (Lopes)"}
\section{Miopia Coletiva (Lopes)}
Lopes: um erro pode ser sustentado por toda uma comunidade de especialistas ao mesmo tempo — o consenso da maioria não é garantia de verdade (ex.: 'os matemáticos não sabem o que é um número real'). O rebanho de doutores repete tolices 'porque uma autoridade o disse'.
\prov{relações: usa autoridade\_e\_rejeicao\_gentil; relaciona espiral\_do\_silencio; relaciona idiotismo\_nao\_examinar; relaciona lei\_de\_ferro\_oligarquia}
\prov{proveniência: wiki \textperiodcentered{} inferencia \textperiodcentered{} confiança 1.0 \textperiodcentered{} domínio ciencias\_de\_fora \textperiodcentered{} história 6 versões no jornal}

%CRISTAL {"arestas":[{"alvo":"nao_dualidade","confianca":1.0,"epistemico":"fato","origem":"wiki","rel":"explica"},{"alvo":"virtualizacao","confianca":0.5,"epistemico":"fato","origem":"wiki","rel":"relaciona"},{"alvo":"word","confianca":0.5,"epistemico":"fato","origem":"wiki","rel":"relaciona"},{"alvo":"vida-morte","confianca":0.5,"epistemico":"fato","origem":"wiki","rel":"relaciona"},{"alvo":"vontade_de_poder","confianca":0.5,"epistemico":"fato","origem":"wiki","rel":"relaciona"},{"alvo":"ontologia","confianca":0.5,"epistemico":"fato","origem":"wiki","rel":"relaciona"},{"alvo":"teleonomia","confianca":0.5,"epistemico":"fato","origem":"wiki","rel":"relaciona"},{"alvo":"paper_0_sintese","confianca":0.5,"epistemico":"fato","origem":"wiki","rel":"relaciona"},{"alvo":"relacao_koch","confianca":0.5,"epistemico":"fato","origem":"wiki","rel":"relaciona"},{"alvo":"musica","confianca":0.5,"epistemico":"fato","origem":"wiki","rel":"relaciona"},{"alvo":"veredito_experimental","confianca":0.5,"epistemico":"fato","origem":"wiki","rel":"relaciona"},{"alvo":"reais","confianca":0.5,"epistemico":"fato","origem":"wiki","rel":"relaciona"}],"confianca":1.0,"contraexemplos":[],"descricao":"A experiência direta de dissolução do eu-outro e unidade com o Absoluto — inefável, com autoridade vivencial (James, Underhill).","epistemico":"hipotese","exemplos":[],"id":"misticismo_nao_dual","memoria":"semantica","meta":{"dominio":"filosofia","fonte":"web"},"origem":"wiki","palavras_chave":[],"sinonimos":[],"tipo":"conceito","titulo":"Misticismo Não-Dual"}
\section{Misticismo Não-Dual}
A experiência direta de dissolução do eu-outro e unidade com o Absoluto — inefável, com autoridade vivencial (James, Underhill).
\prov{relações: explica nao\_dualidade; relaciona virtualizacao; relaciona word; relaciona vida-morte; relaciona vontade\_de\_poder; relaciona ontologia; relaciona teleonomia; relaciona paper\_0\_sintese; relaciona relacao\_koch; relaciona musica; relaciona veredito\_experimental; relaciona reais}
\prov{proveniência: wiki \textperiodcentered{} web \textperiodcentered{} hipotese \textperiodcentered{} confiança 1.0 \textperiodcentered{} domínio filosofia \textperiodcentered{} história 13 versões no jornal}

%CRISTAL {"arestas":[{"alvo":"estudos_de_midia","confianca":1.0,"epistemico":"fato","origem":"wiki","rel":"parte_de"},{"alvo":"morte_do_autor","confianca":1.0,"epistemico":"fato","origem":"wiki","rel":"relaciona"},{"alvo":"hegemonia_cultural","confianca":1.0,"epistemico":"fato","origem":"wiki","rel":"relaciona"},{"alvo":"mediacao_semiotica","confianca":1.0,"epistemico":"fato","origem":"wiki","rel":"usa"}],"confianca":1.0,"contraexemplos":[],"descricao":"Barthes: um sistema semiológico de segunda ordem — um signo já pronto vira significante de um novo sentido, transformando o histórico e cultural em 'natural' e eterno. O mito faz a ideologia parecer senso comum. Opera no plano da conotação (sentidos culturais), não da denotação (o literal).","epistemico":"inferencia","exemplos":[],"id":"mito_barthes","memoria":"semantica","meta":{"autor":"Barthes","dominio":"ciencias_de_fora","fonte":""},"origem":"wiki","palavras_chave":[],"sinonimos":[],"tipo":"conceito","titulo":"Mito (Barthes)"}
\section{Mito (Barthes)}
Barthes: um sistema semiológico de segunda ordem — um signo já pronto vira significante de um novo sentido, transformando o histórico e cultural em 'natural' e eterno. O mito faz a ideologia parecer senso comum. Opera no plano da conotação (sentidos culturais), não da denotação (o literal).
\prov{relações: parte\_de estudos\_de\_midia; relaciona morte\_do\_autor; relaciona hegemonia\_cultural; usa mediacao\_semiotica}
\prov{proveniência: wiki \textperiodcentered{} inferencia \textperiodcentered{} confiança 1.0 \textperiodcentered{} domínio ciencias\_de\_fora \textperiodcentered{} história 5 versões no jornal}

%CRISTAL {"arestas":[{"alvo":"roland_barthes","confianca":1.0,"epistemico":"fato","origem":"wiki","rel":"parte_de"},{"alvo":"denotacao_conotacao","confianca":1.0,"epistemico":"fato","origem":"wiki","rel":"usa"}],"confianca":1.0,"contraexemplos":[],"descricao":"Barthes: o mito é um signo de segunda ordem — toma um signo pronto e o carrega de ideologia naturalizada.","epistemico":"inferencia","exemplos":[],"id":"mito_semiologico","memoria":"semantica","meta":{"autor":"Barthes","dominio":"ciencias_de_fora","fonte":""},"origem":"wiki","palavras_chave":[],"sinonimos":[],"tipo":"conceito","titulo":"O Mito como Fala"}
\section{O Mito como Fala}
Barthes: o mito é um signo de segunda ordem — toma um signo pronto e o carrega de ideologia naturalizada.
\prov{relações: parte\_de roland\_barthes; usa denotacao\_conotacao}
\prov{proveniência: wiki \textperiodcentered{} inferencia \textperiodcentered{} confiança 1.0 \textperiodcentered{} domínio ciencias\_de\_fora \textperiodcentered{} história 3 versões no jornal}

%CRISTAL {"arestas":[{"alvo":"celula","confianca":1.0,"epistemico":"fato","origem":"wiki","rel":"parte_de"},{"alvo":"metabolismo","confianca":1.0,"epistemico":"fato","origem":"wiki","rel":"usa"}],"confianca":1.0,"contraexemplos":[],"descricao":"A usina da célula — converte nutrientes em ATP. Foi, no passado, uma bactéria livre (endossimbiose).","epistemico":"fato","exemplos":[],"id":"mitocondria","memoria":"semantica","meta":{"dominio":"ciencias_de_fora","fonte":""},"origem":"wiki","palavras_chave":[],"sinonimos":[],"tipo":"conceito","titulo":"Mitocôndria"}
\section{Mitocôndria}
A usina da célula — converte nutrientes em ATP. Foi, no passado, uma bactéria livre (endossimbiose).
\prov{relações: parte\_de celula; usa metabolismo}
\prov{proveniência: wiki \textperiodcentered{} fato \textperiodcentered{} confiança 1.0 \textperiodcentered{} domínio ciencias\_de\_fora \textperiodcentered{} história 3 versões no jornal}

%CRISTAL {"arestas":[{"alvo":"code_is_law","confianca":1.0,"epistemico":"fato","origem":"wiki","rel":"usa"},{"alvo":"mao_invisivel","confianca":1.0,"epistemico":"fato","origem":"wiki","rel":"relaciona"}],"confianca":1.0,"contraexemplos":[],"descricao":"Lessig: o comportamento é constrangido por quatro forças que interagem — lei, mercado, normas sociais e arquitetura/código; o código é apenas uma delas, mas a mais decisiva online.","epistemico":"inferencia","exemplos":[],"id":"modalidades_de_regulacao","memoria":"semantica","meta":{"autor":"Lessig","dominio":"ciencias_de_fora","fonte":""},"origem":"wiki","palavras_chave":[],"sinonimos":[],"tipo":"conceito","titulo":"As Quatro Modalidades de Regulação (Lessig)"}
\section{As Quatro Modalidades de Regulação (Lessig)}
Lessig: o comportamento é constrangido por quatro forças que interagem — lei, mercado, normas sociais e arquitetura/código; o código é apenas uma delas, mas a mais decisiva online.
\prov{relações: usa code\_is\_law; relaciona mao\_invisivel}
\prov{proveniência: wiki \textperiodcentered{} inferencia \textperiodcentered{} confiança 1.0 \textperiodcentered{} domínio ciencias\_de\_fora \textperiodcentered{} história 3 versões no jornal}

%CRISTAL {"arestas":[{"alvo":"teoria_da_comunicacao","confianca":1.0,"epistemico":"fato","origem":"wiki","rel":"parte_de"}],"confianca":1.0,"contraexemplos":[],"descricao":"Lasswell: descreve o ato comunicativo por cinco perguntas — quem diz o quê, em que canal, a quem, com que efeito — fixando os componentes do processo e suas áreas de análise. Linear e unidirecional.","epistemico":"inferencia","exemplos":[],"id":"modelo_5w_lasswell","memoria":"semantica","meta":{"autor":"Lasswell","dominio":"ciencias_de_fora","fonte":""},"origem":"wiki","palavras_chave":[],"sinonimos":[],"tipo":"conceito","titulo":"Modelo dos 5 W (Lasswell)"}
\section{Modelo dos 5 W (Lasswell)}
Lasswell: descreve o ato comunicativo por cinco perguntas — quem diz o quê, em que canal, a quem, com que efeito — fixando os componentes do processo e suas áreas de análise. Linear e unidirecional.
\prov{relações: parte\_de teoria\_da\_comunicacao}
\prov{proveniência: wiki \textperiodcentered{} inferencia \textperiodcentered{} confiança 1.0 \textperiodcentered{} domínio ciencias\_de\_fora \textperiodcentered{} história 2 versões no jornal}

%CRISTAL {"arestas":[{"alvo":"filosofia_da_medicina","confianca":1.0,"epistemico":"fato","origem":"wiki","rel":"parte_de"},{"alvo":"materialismo_fisicalismo","confianca":1.0,"epistemico":"fato","origem":"wiki","rel":"usa"}],"confianca":1.0,"contraexemplos":[],"descricao":"O paradigma que concebe a doença como desvio mensurável de variáveis biológicas, redutível a mecanismos físico-químicos — deixando de fora as dimensões psicológicas e sociais. Assume dualismo mente-corpo e reducionismo.","epistemico":"inferencia","exemplos":[],"id":"modelo_biomedico","memoria":"semantica","meta":{"autor":"—","dominio":"ciencias_de_fora","fonte":""},"origem":"wiki","palavras_chave":[],"sinonimos":[],"tipo":"conceito","titulo":"Modelo Biomédico"}
\section{Modelo Biomédico}
O paradigma que concebe a doença como desvio mensurável de variáveis biológicas, redutível a mecanismos físico-químicos — deixando de fora as dimensões psicológicas e sociais. Assume dualismo mente-corpo e reducionismo.
\prov{relações: parte\_de filosofia\_da\_medicina; usa materialismo\_fisicalismo}
\prov{proveniência: wiki \textperiodcentered{} inferencia \textperiodcentered{} confiança 1.0 \textperiodcentered{} domínio ciencias\_de\_fora \textperiodcentered{} história 3 versões no jornal}

%CRISTAL {"arestas":[{"alvo":"modelo_biomedico","confianca":1.0,"epistemico":"fato","origem":"wiki","rel":"contradiz"},{"alvo":"aparelho_psiquico","confianca":1.0,"epistemico":"fato","origem":"wiki","rel":"relaciona"}],"confianca":1.0,"contraexemplos":[],"descricao":"Engel (1977): saúde e doença resultam da interação de fatores biológicos, psicológicos e sociais; propõe substituir o reducionismo biomédico por uma abordagem sistêmica que integra os múltiplos níveis do adoecer.","epistemico":"inferencia","exemplos":[],"id":"modelo_biopsicossocial_engel","memoria":"semantica","meta":{"autor":"Engel","dominio":"ciencias_de_fora","fonte":""},"origem":"wiki","palavras_chave":[],"sinonimos":[],"tipo":"conceito","titulo":"Modelo Biopsicossocial (Engel)"}
\section{Modelo Biopsicossocial (Engel)}
Engel (1977): saúde e doença resultam da interação de fatores biológicos, psicológicos e sociais; propõe substituir o reducionismo biomédico por uma abordagem sistêmica que integra os múltiplos níveis do adoecer.
\prov{relações: contradiz modelo\_biomedico; relaciona aparelho\_psiquico}
\prov{proveniência: wiki \textperiodcentered{} inferencia \textperiodcentered{} confiança 1.0 \textperiodcentered{} domínio ciencias\_de\_fora \textperiodcentered{} história 3 versões no jornal}

%CRISTAL {"arestas":[{"alvo":"teoria_da_comunicacao","confianca":1.0,"epistemico":"fato","origem":"wiki","rel":"parte_de"},{"alvo":"teoria_da_informacao_shannon","confianca":1.0,"epistemico":"fato","origem":"wiki","rel":"usa"},{"alvo":"modelo_5w_lasswell","confianca":1.0,"epistemico":"fato","origem":"wiki","rel":"relaciona"},{"alvo":"mediacao_semiotica","confianca":1.0,"epistemico":"fato","origem":"wiki","rel":"relaciona"}],"confianca":1.0,"contraexemplos":[],"descricao":"Shannon & Weaver: uma fonte produz uma mensagem, um transmissor a codifica em sinal que atravessa um canal com ruído até um receptor. Weaver estendeu o esquema técnico de Shannon a três níveis — técnico, semântico e de eficácia — acrescentando o ruído semântico.","epistemico":"inferencia","exemplos":[],"id":"modelo_shannon_weaver","memoria":"semantica","meta":{"autor":"Shannon & Weaver","dominio":"ciencias_de_fora","fonte":""},"origem":"wiki","palavras_chave":[],"sinonimos":[],"tipo":"conceito","titulo":"Modelo Transmissivo (Shannon-Weaver)"}
\section{Modelo Transmissivo (Shannon-Weaver)}
Shannon \& Weaver: uma fonte produz uma mensagem, um transmissor a codifica em sinal que atravessa um canal com ruído até um receptor. Weaver estendeu o esquema técnico de Shannon a três níveis — técnico, semântico e de eficácia — acrescentando o ruído semântico.
\prov{relações: parte\_de teoria\_da\_comunicacao; usa teoria\_da\_informacao\_shannon; relaciona modelo\_5w\_lasswell; relaciona mediacao\_semiotica}
\prov{proveniência: wiki \textperiodcentered{} inferencia \textperiodcentered{} confiança 1.0 \textperiodcentered{} domínio ciencias\_de\_fora \textperiodcentered{} história 5 versões no jornal}

%CRISTAL {"arestas":[{"alvo":"aprendizagem_descoberta_bruner","confianca":1.0,"epistemico":"fato","origem":"wiki","rel":"usa"},{"alvo":"mente_computacao","confianca":1.0,"epistemico":"fato","origem":"wiki","rel":"referencia"}],"confianca":1.0,"contraexemplos":[],"descricao":"Bruner: o conhecimento é representado em três modos integráveis — pela ação (enativo), por imagens (icônico) e por linguagem/símbolos abstratos (simbólico).","epistemico":"inferencia","exemplos":[],"id":"modos_representacao_bruner","memoria":"semantica","meta":{"autor":"Bruner","dominio":"ciencias_de_fora","fonte":""},"origem":"wiki","palavras_chave":[],"sinonimos":[],"tipo":"conceito","titulo":"Representação Enativa, Icônica e Simbólica"}
\section{Representação Enativa, Icônica e Simbólica}
Bruner: o conhecimento é representado em três modos integráveis — pela ação (enativo), por imagens (icônico) e por linguagem/símbolos abstratos (simbólico).
\prov{relações: usa aprendizagem\_descoberta\_bruner; referencia mente\_computacao}
\prov{proveniência: wiki \textperiodcentered{} inferencia \textperiodcentered{} confiança 1.0 \textperiodcentered{} domínio ciencias\_de\_fora \textperiodcentered{} história 3 versões no jornal}

%CRISTAL {"arestas":[{"alvo":"economia","confianca":1.0,"epistemico":"fato","origem":"wiki","rel":"parte_de"},{"alvo":"goldchain","confianca":1.0,"epistemico":"fato","origem":"wiki","rel":"explica"},{"alvo":"cifra","confianca":1.0,"epistemico":"fato","origem":"wiki","rel":"referencia"}],"confianca":1.0,"contraexemplos":[],"descricao":"A moeda é instituição social: padrão-ouro (lastro material), fiduciária (confiança no Estado) ou criptoeconômica (confiança no protocolo matemático).","epistemico":"inferencia","exemplos":[],"id":"moeda_confianca","memoria":"semantica","meta":{"autor":"—","dominio":"ciencias_de_fora","fonte":""},"origem":"wiki","palavras_chave":[],"sinonimos":[],"tipo":"conceito","titulo":"Moeda e Confiança"}
\section{Moeda e Confiança}
A moeda é instituição social: padrão-ouro (lastro material), fiduciária (confiança no Estado) ou criptoeconômica (confiança no protocolo matemático).
\prov{relações: parte\_de economia; explica goldchain; referencia cifra}
\prov{proveniência: wiki \textperiodcentered{} inferencia \textperiodcentered{} confiança 1.0 \textperiodcentered{} domínio ciencias\_de\_fora \textperiodcentered{} história 4 versões no jornal}

%CRISTAL {"arestas":[{"alvo":"ciencia_politica","confianca":1.0,"epistemico":"fato","origem":"wiki","rel":"parte_de"},{"alvo":"contrato_social","confianca":1.0,"epistemico":"fato","origem":"wiki","rel":"usa"},{"alvo":"legitimidade_weber","confianca":1.0,"epistemico":"fato","origem":"wiki","rel":"usa"}],"confianca":1.0,"contraexemplos":[],"descricao":"Weber: o Estado moderno é a comunidade humana que reivindica com êxito o monopólio do uso legítimo da força física dentro de um território — toda outra violência só é lícita na medida em que o Estado a autoriza.","epistemico":"inferencia","exemplos":[],"id":"monopolio_violencia_legitima","memoria":"semantica","meta":{"autor":"Weber","dominio":"ciencias_de_fora","fonte":""},"origem":"wiki","palavras_chave":[],"sinonimos":[],"tipo":"conceito","titulo":"O Estado e o Monopólio da Força (Weber)"}
\section{O Estado e o Monopólio da Força (Weber)}
Weber: o Estado moderno é a comunidade humana que reivindica com êxito o monopólio do uso legítimo da força física dentro de um território — toda outra violência só é lícita na medida em que o Estado a autoriza.
\prov{relações: parte\_de ciencia\_politica; usa contrato\_social; usa legitimidade\_weber}
\prov{proveniência: wiki \textperiodcentered{} inferencia \textperiodcentered{} confiança 1.0 \textperiodcentered{} domínio ciencias\_de\_fora \textperiodcentered{} história 4 versões no jornal}

%CRISTAL {"arestas":[{"alvo":"cristianismo","confianca":1.0,"epistemico":"fato","origem":"wiki","rel":"parte_de"},{"alvo":"deus_mentiroso","confianca":1.0,"epistemico":"fato","origem":"wiki","rel":"contradiz"},{"alvo":"tawhid","confianca":1.0,"epistemico":"fato","origem":"wiki","rel":"relaciona"}],"confianca":1.0,"contraexemplos":[],"descricao":"Bíblia (Deuteronômio 6:4): a fé num Deus único, universal e moralmente exigente — 'ouve, Israel, o Senhor nosso Deus é o único Senhor'. A reivindicação de um Deus único e verdadeiro é exatamente o que a crítica do Lopes (deus_mentiroso) interroga — aqui, a posição, não o veredicto.","epistemico":"inferencia","exemplos":[],"id":"monoteismo_etico","memoria":"semantica","meta":{"autor":"Bíblia","dominio":"sagrado","fonte":""},"origem":"wiki","palavras_chave":[],"sinonimos":[],"tipo":"conceito","titulo":"Monoteísmo Ético"}
\section{Monoteísmo Ético}
Bíblia (Deuteronômio 6:4): a fé num Deus único, universal e moralmente exigente — 'ouve, Israel, o Senhor nosso Deus é o único Senhor'. A reivindicação de um Deus único e verdadeiro é exatamente o que a crítica do Lopes (deus\_mentiroso) interroga — aqui, a posição, não o veredicto.
\prov{relações: parte\_de cristianismo; contradiz deus\_mentiroso; relaciona tawhid}
\prov{proveniência: wiki \textperiodcentered{} inferencia \textperiodcentered{} confiança 1.0 \textperiodcentered{} domínio sagrado \textperiodcentered{} história 5 versões no jornal}

%CRISTAL {"arestas":[{"alvo":"artes_visuais","confianca":1.0,"epistemico":"fato","origem":"wiki","rel":"parte_de"},{"alvo":"dialetica_hegeliana","confianca":1.0,"epistemico":"fato","origem":"wiki","rel":"usa"}],"confianca":1.0,"contraexemplos":[],"descricao":"Eisenstein: o sentido no cinema não está no plano isolado, mas na colisão entre planos — tese + antítese geram uma síntese no espírito do espectador (dialética aplicada à edição). O conflito é o princípio construtivo do filme.","epistemico":"inferencia","exemplos":[],"id":"montagem_eisenstein","memoria":"semantica","meta":{"autor":"Eisenstein","dominio":"ciencias_de_fora","fonte":""},"origem":"wiki","palavras_chave":[],"sinonimos":[],"tipo":"conceito","titulo":"Montagem Dialética (Eisenstein)"}
\section{Montagem Dialética (Eisenstein)}
Eisenstein: o sentido no cinema não está no plano isolado, mas na colisão entre planos — tese + antítese geram uma síntese no espírito do espectador (dialética aplicada à edição). O conflito é o princípio construtivo do filme.
\prov{relações: parte\_de artes\_visuais; usa dialetica\_hegeliana}
\prov{proveniência: wiki \textperiodcentered{} inferencia \textperiodcentered{} confiança 1.0 \textperiodcentered{} domínio ciencias\_de\_fora \textperiodcentered{} história 3 versões no jornal}

%CRISTAL {"arestas":[{"alvo":"auto_organizacao","confianca":1.0,"epistemico":"fato","origem":"wiki","rel":"usa"},{"alvo":"maquina_de_turing","confianca":1.0,"epistemico":"fato","origem":"wiki","rel":"relaciona"},{"alvo":"geometria_fractal","confianca":1.0,"epistemico":"fato","origem":"wiki","rel":"relaciona"}],"confianca":1.0,"contraexemplos":[],"descricao":"Turing: reação-difusão gera manchas e listras a partir do uniforme — a matemática por trás da pele da onça e da zebra.","epistemico":"fato","exemplos":[],"id":"morfogenese_turing","memoria":"semantica","meta":{"autor":"Turing","dominio":"ciencias_de_fora","fonte":""},"origem":"wiki","palavras_chave":[],"sinonimos":[],"tipo":"conceito","titulo":"Morfogênese (Padrões de Turing)"}
\section{Morfogênese (Padrões de Turing)}
Turing: reação-difusão gera manchas e listras a partir do uniforme — a matemática por trás da pele da onça e da zebra.
\prov{relações: usa auto\_organizacao; relaciona maquina\_de\_turing; relaciona geometria\_fractal}
\prov{proveniência: wiki \textperiodcentered{} fato \textperiodcentered{} confiança 1.0 \textperiodcentered{} domínio ciencias\_de\_fora \textperiodcentered{} história 4 versões no jornal}

%CRISTAL {"arestas":[{"alvo":"linguistica","confianca":1.0,"epistemico":"fato","origem":"wiki","rel":"parte_de"}],"confianca":1.0,"contraexemplos":[],"descricao":"O estudo da forma das palavras — morfemas, flexão, derivação; as peças mínimas com sentido.","epistemico":"fato","exemplos":[],"id":"morfologia","memoria":"semantica","meta":{"dominio":"ciencias_de_fora","fonte":""},"origem":"wiki","palavras_chave":[],"sinonimos":[],"tipo":"conceito","titulo":"Morfologia"}
\section{Morfologia}
O estudo da forma das palavras — morfemas, flexão, derivação; as peças mínimas com sentido.
\prov{relações: parte\_de linguistica}
\prov{proveniência: wiki \textperiodcentered{} fato \textperiodcentered{} confiança 1.0 \textperiodcentered{} domínio ciencias\_de\_fora \textperiodcentered{} história 2 versões no jornal}

%CRISTAL {"arestas":[{"alvo":"literatura","confianca":1.0,"epistemico":"fato","origem":"wiki","rel":"parte_de"},{"alvo":"estruturalismo_linguistico","confianca":1.0,"epistemico":"fato","origem":"wiki","rel":"usa"}],"confianca":1.0,"contraexemplos":[],"descricao":"Propp: sob a variedade dos contos maravilhosos há uma estrutura constante de 31 funções (ações fixas em sequência regular) — a narrativa como sistema, precursora da análise estrutural.","epistemico":"inferencia","exemplos":[],"id":"morfologia_propp","memoria":"semantica","meta":{"autor":"Propp","dominio":"ciencias_de_fora","fonte":""},"origem":"wiki","palavras_chave":[],"sinonimos":[],"tipo":"conceito","titulo":"Morfologia do Conto (Propp)"}
\section{Morfologia do Conto (Propp)}
Propp: sob a variedade dos contos maravilhosos há uma estrutura constante de 31 funções (ações fixas em sequência regular) — a narrativa como sistema, precursora da análise estrutural.
\prov{relações: parte\_de literatura; usa estruturalismo\_linguistico}
\prov{proveniência: wiki \textperiodcentered{} inferencia \textperiodcentered{} confiança 1.0 \textperiodcentered{} domínio ciencias\_de\_fora \textperiodcentered{} história 3 versões no jornal}

%CRISTAL {"arestas":[{"alvo":"bioetica","confianca":1.0,"epistemico":"fato","origem":"wiki","rel":"parte_de"},{"alvo":"fim_de_vida_espectro","confianca":1.0,"epistemico":"fato","origem":"wiki","rel":"relaciona"},{"alvo":"consciencia","confianca":1.0,"epistemico":"fato","origem":"wiki","rel":"relaciona"}],"confianca":1.0,"contraexemplos":[],"descricao":"O critério que define a morte como cessação irreversível de todas as funções do encéfalo (UDDA); surgiu com o suporte de vida e a demanda por transplantes, e permanece contestado (todo o encéfalo vs encéfalo superior/consciência) — convenção normativa disfarçada de fato biológico.","epistemico":"inferencia","exemplos":[],"id":"morte_cerebral","memoria":"semantica","meta":{"autor":"Comitê de Harvard","dominio":"ciencias_de_fora","fonte":""},"origem":"wiki","palavras_chave":[],"sinonimos":[],"tipo":"conceito","titulo":"Morte Encefálica e a Definição de Morte"}
\section{Morte Encefálica e a Definição de Morte}
O critério que define a morte como cessação irreversível de todas as funções do encéfalo (UDDA); surgiu com o suporte de vida e a demanda por transplantes, e permanece contestado (todo o encéfalo vs encéfalo superior/consciência) — convenção normativa disfarçada de fato biológico.
\prov{relações: parte\_de bioetica; relaciona fim\_de\_vida\_espectro; relaciona consciencia}
\prov{proveniência: wiki \textperiodcentered{} inferencia \textperiodcentered{} confiança 1.0 \textperiodcentered{} domínio ciencias\_de\_fora \textperiodcentered{} história 4 versões no jornal}

%CRISTAL {"arestas":[{"alvo":"face_bela_morte","confianca":1.0,"epistemico":"fato","origem":"wiki","rel":"relaciona"},{"alvo":"historia","confianca":1.0,"epistemico":"fato","origem":"wiki","rel":"referencia"}],"confianca":1.0,"contraexemplos":[],"descricao":"Planck: 'a ciência avança um funeral por vez'; Kuhn: as revoluções de paradigma são geracionais, não puramente racionais. A morte biológica move o saber.","epistemico":"inferencia","exemplos":[],"id":"morte_de_geracoes","memoria":"semantica","meta":{"autor":"Planck/Kuhn","dominio":"ciencias_de_fora","fonte":""},"origem":"wiki","palavras_chave":[],"sinonimos":[],"tipo":"conceito","titulo":"Morte de Gerações (Planck/Kuhn)"}
\section{Morte de Gerações (Planck/Kuhn)}
Planck: 'a ciência avança um funeral por vez'; Kuhn: as revoluções de paradigma são geracionais, não puramente racionais. A morte biológica move o saber.
\prov{relações: relaciona face\_bela\_morte; referencia historia}
\prov{proveniência: wiki \textperiodcentered{} inferencia \textperiodcentered{} confiança 1.0 \textperiodcentered{} domínio ciencias\_de\_fora \textperiodcentered{} história 3 versões no jornal}

%CRISTAL {"arestas":[{"alvo":"aura_benjamin","confianca":1.0,"epistemico":"fato","origem":"wiki","rel":"relaciona"},{"alvo":"colapso","confianca":1.0,"epistemico":"fato","origem":"wiki","rel":"referencia"},{"alvo":"virtualizacao","confianca":0.5,"epistemico":"fato","origem":"wiki","rel":"relaciona"},{"alvo":"roland_barthes","confianca":1.0,"epistemico":"fato","origem":"wiki","rel":"parte_de"},{"alvo":"funcao_autor_foucault","confianca":1.0,"epistemico":"fato","origem":"wiki","rel":"relaciona"}],"confianca":1.0,"contraexemplos":[],"descricao":"Barthes: o sentido de um texto nasce na leitura, não na intenção do autor — o texto se solta de quem o escreveu.","epistemico":"inferencia","exemplos":[],"id":"morte_do_autor","memoria":"semantica","meta":{"autor":"Barthes","dominio":"ciencias_de_fora","fonte":""},"origem":"wiki","palavras_chave":[],"sinonimos":[],"tipo":"conceito","titulo":"A Morte do Autor"}
\section{A Morte do Autor}
Barthes: o sentido de um texto nasce na leitura, não na intenção do autor — o texto se solta de quem o escreveu.
\prov{relações: relaciona aura\_benjamin; referencia colapso; relaciona virtualizacao; parte\_de roland\_barthes; relaciona funcao\_autor\_foucault}
\prov{proveniência: wiki \textperiodcentered{} inferencia \textperiodcentered{} confiança 1.0 \textperiodcentered{} domínio ciencias\_de\_fora \textperiodcentered{} história 8 versões no jornal}

%CRISTAL {"arestas":[{"alvo":"ciencia_politica","confianca":1.0,"epistemico":"fato","origem":"wiki","rel":"parte_de"},{"alvo":"logica_acao_coletiva","confianca":1.0,"epistemico":"fato","origem":"wiki","rel":"usa"},{"alvo":"sociedade_civil_tocqueville","confianca":1.0,"epistemico":"fato","origem":"wiki","rel":"relaciona"}],"confianca":1.0,"contraexemplos":[],"descricao":"Tilly & Tarrow: a ação coletiva contenciosa se organiza em 'repertórios de confronto' (formas culturalmente disponíveis de reivindicar — greve, manifestação) que se ativam conforme a 'estrutura de oportunidades políticas' (aberturas e vulnerabilidades do regime).","epistemico":"inferencia","exemplos":[],"id":"movimentos_sociais_contenciosos","memoria":"semantica","meta":{"autor":"Tilly & Tarrow","dominio":"ciencias_de_fora","fonte":""},"origem":"wiki","palavras_chave":[],"sinonimos":[],"tipo":"conceito","titulo":"Política Contenciosa (Tilly & Tarrow)"}
\section{Política Contenciosa (Tilly \& Tarrow)}
Tilly \& Tarrow: a ação coletiva contenciosa se organiza em 'repertórios de confronto' (formas culturalmente disponíveis de reivindicar — greve, manifestação) que se ativam conforme a 'estrutura de oportunidades políticas' (aberturas e vulnerabilidades do regime).
\prov{relações: parte\_de ciencia\_politica; usa logica\_acao\_coletiva; relaciona sociedade\_civil\_tocqueville}
\prov{proveniência: wiki \textperiodcentered{} inferencia \textperiodcentered{} confiança 1.0 \textperiodcentered{} domínio ciencias\_de\_fora \textperiodcentered{} história 4 versões no jornal}

%CRISTAL {"arestas":[{"alvo":"geografia_fisica","confianca":1.0,"epistemico":"fato","origem":"wiki","rel":"parte_de"},{"alvo":"feedback_retroalimentacao","confianca":1.0,"epistemico":"fato","origem":"wiki","rel":"usa"},{"alvo":"antropoceno","confianca":1.0,"epistemico":"fato","origem":"wiki","rel":"relaciona"}],"confianca":1.0,"contraexemplos":[],"descricao":"O aquecimento global causado pelo aumento de gases de efeito estufa de origem humana, amplificado por retroalimentações (vapor d'água, gelo-albedo, carbono) — consenso científico consolidado pelo IPCC.","epistemico":"fato","exemplos":[],"id":"mudanca_climatica","memoria":"semantica","meta":{"autor":"IPCC","dominio":"ciencias_de_fora","fonte":""},"origem":"wiki","palavras_chave":[],"sinonimos":[],"tipo":"conceito","titulo":"Mudança Climática Antropogênica"}
\section{Mudança Climática Antropogênica}
O aquecimento global causado pelo aumento de gases de efeito estufa de origem humana, amplificado por retroalimentações (vapor d'água, gelo-albedo, carbono) — consenso científico consolidado pelo IPCC.
\prov{relações: parte\_de geografia\_fisica; usa feedback\_retroalimentacao; relaciona antropoceno}
\prov{proveniência: wiki \textperiodcentered{} fato \textperiodcentered{} confiança 1.0 \textperiodcentered{} domínio ciencias\_de\_fora \textperiodcentered{} história 4 versões no jornal}

%CRISTAL {"arestas":[{"alvo":"sincronia_diacronia","confianca":1.0,"epistemico":"fato","origem":"wiki","rel":"relaciona"},{"alvo":"idioma","confianca":1.0,"epistemico":"fato","origem":"wiki","rel":"usa"},{"alvo":"linguistica","confianca":1.0,"epistemico":"fato","origem":"wiki","rel":"parte_de"}],"confianca":1.0,"contraexemplos":[],"descricao":"Toda língua viva muda: sons, palavras e gramática derivam no tempo (a lei de Grimm) — a diacronia que separa o latim das línguas românicas.","epistemico":"fato","exemplos":[],"id":"mudanca_linguistica","memoria":"semantica","meta":{"autor":"Grimm","dominio":"ciencias_de_fora","fonte":""},"origem":"wiki","palavras_chave":[],"sinonimos":[],"tipo":"conceito","titulo":"Mudança Linguística"}
\section{Mudança Linguística}
Toda língua viva muda: sons, palavras e gramática derivam no tempo (a lei de Grimm) — a diacronia que separa o latim das línguas românicas.
\prov{relações: relaciona sincronia\_diacronia; usa idioma; parte\_de linguistica}
\prov{proveniência: wiki \textperiodcentered{} fato \textperiodcentered{} confiança 1.0 \textperiodcentered{} domínio ciencias\_de\_fora \textperiodcentered{} história 4 versões no jornal}

%CRISTAL {"arestas":[{"alvo":"filosofia_da_tecnica","confianca":1.0,"epistemico":"fato","origem":"wiki","rel":"parte_de"},{"alvo":"biopoder","confianca":1.0,"epistemico":"fato","origem":"wiki","rel":"relaciona"}],"confianca":1.0,"contraexemplos":[],"descricao":"Mumford: uma máquina feita de seres humanos organizados hierarquicamente como engrenagens (burocracia + exército + trabalho), cuja origem ele data de ~5.000 anos (pirâmides), não da Revolução Industrial — a máquina como organização social antes que como artefato.","epistemico":"inferencia","exemplos":[],"id":"mumford_megamaquina","memoria":"semantica","meta":{"autor":"Mumford","dominio":"ciencias_de_fora","fonte":""},"origem":"wiki","palavras_chave":[],"sinonimos":[],"tipo":"conceito","titulo":"A Megamáquina (Mumford)"}
\section{A Megamáquina (Mumford)}
Mumford: uma máquina feita de seres humanos organizados hierarquicamente como engrenagens (burocracia + exército + trabalho), cuja origem ele data de \~{}5.000 anos (pirâmides), não da Revolução Industrial — a máquina como organização social antes que como artefato.
\prov{relações: parte\_de filosofia\_da\_tecnica; relaciona biopoder}
\prov{proveniência: wiki \textperiodcentered{} inferencia \textperiodcentered{} confiança 1.0 \textperiodcentered{} domínio ciencias\_de\_fora \textperiodcentered{} história 3 versões no jornal}

%CRISTAL {"arestas":[{"alvo":"mumford_megamaquina","confianca":1.0,"epistemico":"fato","origem":"wiki","rel":"relaciona"},{"alvo":"filosofia_da_tecnica","confianca":1.0,"epistemico":"fato","origem":"wiki","rel":"parte_de"}],"confianca":1.0,"contraexemplos":[],"descricao":"Mumford: uma técnica autoritária — centrada em sistema, poder e controle, potente mas instável (monotécnica) — vs uma técnica democrática — de pequena escala, centrada na vida, cooperativa e plural (politécnica).","epistemico":"inferencia","exemplos":[],"id":"mumford_tecnica_autoritaria_democratica","memoria":"semantica","meta":{"autor":"Mumford","dominio":"ciencias_de_fora","fonte":""},"origem":"wiki","palavras_chave":[],"sinonimos":[],"tipo":"conceito","titulo":"Técnica Autoritária vs Democrática (Mumford)"}
\section{Técnica Autoritária vs Democrática (Mumford)}
Mumford: uma técnica autoritária — centrada em sistema, poder e controle, potente mas instável (monotécnica) — vs uma técnica democrática — de pequena escala, centrada na vida, cooperativa e plural (politécnica).
\prov{relações: relaciona mumford\_megamaquina; parte\_de filosofia\_da\_tecnica}
\prov{proveniência: wiki \textperiodcentered{} inferencia \textperiodcentered{} confiança 1.0 \textperiodcentered{} domínio ciencias\_de\_fora \textperiodcentered{} história 3 versões no jornal}

%CRISTAL {"arestas":[{"alvo":"abiogenese","confianca":1.0,"epistemico":"fato","origem":"wiki","rel":"relaciona"},{"alvo":"codigo_genetico","confianca":1.0,"epistemico":"fato","origem":"wiki","rel":"relaciona"}],"confianca":1.0,"contraexemplos":[],"descricao":"Hipótese de que o RNA veio antes: uma molécula que guarda informação E catalisa — antes da divisão DNA/proteína.","epistemico":"hipotese","exemplos":[],"id":"mundo_de_rna","memoria":"semantica","meta":{"dominio":"ciencias_de_fora","fonte":""},"origem":"wiki","palavras_chave":[],"sinonimos":[],"tipo":"conceito","titulo":"Mundo de RNA"}
\section{Mundo de RNA}
Hipótese de que o RNA veio antes: uma molécula que guarda informação E catalisa — antes da divisão DNA/proteína.
\prov{relações: relaciona abiogenese; relaciona codigo\_genetico}
\prov{proveniência: wiki \textperiodcentered{} hipotese \textperiodcentered{} confiança 1.0 \textperiodcentered{} domínio ciencias\_de\_fora \textperiodcentered{} história 4 versões no jornal}

%CRISTAL {"arestas":[{"alvo":"tecido_mineral_hub","confianca":1.0,"epistemico":"fato","origem":"wiki","rel":"usa"},{"alvo":"onda_de_excitacao_e_a_propagacao","confianca":1.0,"epistemico":"fato","origem":"wiki","rel":"usa"},{"alvo":"chicote_e_o_tensor","confianca":1.0,"epistemico":"fato","origem":"wiki","rel":"usa"},{"alvo":"biologia","confianca":1.0,"epistemico":"fato","origem":"wiki","rel":"eh_um"}],"confianca":1.0,"contraexemplos":[],"descricao":"Tecido excitável CONTRÁTIL, o passo antes da face (linguagem/musculo.py, 7/7). A onda de ativação (do tecido) dispara o encurtamento; o strain resultante (o tensor) deforma a superfície. Quatro leituras: a curva comprimento-tensão é uma PARÁBOLA (pico no ótimo); a força-velocidade de Hill é uma HIPÉRBOLE; o acoplamento excitação-contração dá o abalo e, em alta frequência, o TÉTANO; e a contração é STRAIN (o chicote-tensor, ε_fibra<0). O músculo é onde a excitação vira força e a força vira forma.","epistemico":"inferencia","exemplos":[],"id":"musculo_mineral_hub","memoria":"semantica","meta":{"dominio":"biologia","fonte":"músculo mineral"},"origem":"wiki","palavras_chave":[],"sinonimos":[],"tipo":"conceito","titulo":"O Músculo pela Máquina Mineral (a onda dispara, o strain deforma)"}
\section{O Músculo pela Máquina Mineral (a onda dispara, o strain deforma)}
Tecido excitável CONTRÁTIL, o passo antes da face (linguagem/musculo.py, 7/7). A onda de ativação (do tecido) dispara o encurtamento; o strain resultante (o tensor) deforma a superfície. Quatro leituras: a curva comprimento-tensão é uma PARÁBOLA (pico no ótimo); a força-velocidade de Hill é uma HIPÉRBOLE; o acoplamento excitação-contração dá o abalo e, em alta frequência, o TÉTANO; e a contração é STRAIN (o chicote-tensor, \ensuremath{\varepsilon}\_fibra<0). O músculo é onde a excitação vira força e a força vira forma.
\prov{relações: usa tecido\_mineral\_hub; usa onda\_de\_excitacao\_e\_a\_propagacao; usa chicote\_e\_o\_tensor; eh\_um biologia}
\prov{proveniência: wiki \textperiodcentered{} músculo mineral \textperiodcentered{} inferencia \textperiodcentered{} confiança 1.0 \textperiodcentered{} domínio biologia \textperiodcentered{} história 5 versões no jornal}

%CRISTAL {"arestas":[{"alvo":"atuadores_da_face","confianca":1.0,"epistemico":"fato","origem":"wiki","rel":"usa"},{"alvo":"camadas_da_face","confianca":1.0,"epistemico":"fato","origem":"wiki","rel":"relaciona"},{"alvo":"dicionario_musculos_face","confianca":1.0,"epistemico":"fato","origem":"wiki","rel":"usa"}],"confianca":1.0,"contraexemplos":[],"descricao":"Os centros dos 8 músculos faciais são CALIBRADOS contra a geometria do scan real (linguagem/calibrar.py, 4/4), não mais palpites fixos. Sem ML: detecção GEOMÉTRICA — o NARIZ é detectado (a protrusão de maior z na coluna central), e a partir dele + a altura (testa→queixo) e a largura medidas, os landmarks (olhos, boca, sobrancelha, bochecha, mandíbula) derivam pela PROPORÇÃO anatômica (o cânon facial). Os músculos ganham centros e raios proporcionais ao rosto de verdade — adaptados à posição, escala e proporção do teu scan. Verif.: no scan real o nariz é detectado e a ordem vertical bate (sobrancelha>olho>nariz>boca>queixo).","epistemico":"fato","exemplos":[],"id":"musculos_calibrados_ao_scan","memoria":"semantica","meta":{"dominio":"biologia","fonte":"calibração"},"origem":"wiki","palavras_chave":[],"sinonimos":[],"tipo":"conceito","titulo":"Os Músculos Calibrados ao Scan (detecção geométrica)"}
\section{Os Músculos Calibrados ao Scan (detecção geométrica)}
Os centros dos 8 músculos faciais são CALIBRADOS contra a geometria do scan real (linguagem/calibrar.py, 4/4), não mais palpites fixos. Sem ML: detecção GEOMÉTRICA — o NARIZ é detectado (a protrusão de maior z na coluna central), e a partir dele + a altura (testa\ensuremath{\to}queixo) e a largura medidas, os landmarks (olhos, boca, sobrancelha, bochecha, mandíbula) derivam pela PROPORÇÃO anatômica (o cânon facial). Os músculos ganham centros e raios proporcionais ao rosto de verdade — adaptados à posição, escala e proporção do teu scan. Verif.: no scan real o nariz é detectado e a ordem vertical bate (sobrancelha>olho>nariz>boca>queixo).
\prov{relações: usa atuadores\_da\_face; relaciona camadas\_da\_face; usa dicionario\_musculos\_face}
\prov{proveniência: wiki \textperiodcentered{} calibração \textperiodcentered{} fato \textperiodcentered{} confiança 1.0 \textperiodcentered{} domínio biologia \textperiodcentered{} história 4 versões no jornal}

%CRISTAL {"arestas":[{"alvo":"artes","confianca":1.0,"epistemico":"fato","origem":"wiki","rel":"parte_de"},{"alvo":"virtualizacao","confianca":0.5,"epistemico":"fato","origem":"wiki","rel":"relaciona"},{"alvo":"word","confianca":0.5,"epistemico":"fato","origem":"wiki","rel":"relaciona"},{"alvo":"arte","confianca":1.0,"epistemico":"fato","origem":"wiki","rel":"parte_de"}],"confianca":1.0,"contraexemplos":[],"descricao":"A arte do som organizado no tempo — melodia, harmonia e ritmo — a mais abstrata e matemática das artes, e talvez a mais diretamente emocional.","epistemico":"fato","exemplos":[],"id":"musica","memoria":"semantica","meta":{"dominio":"ciencias_de_fora","fonte":""},"origem":"wiki","palavras_chave":[],"sinonimos":[],"tipo":"conceito","titulo":"Música"}
\section{Música}
A arte do som organizado no tempo — melodia, harmonia e ritmo — a mais abstrata e matemática das artes, e talvez a mais diretamente emocional.
\prov{relações: parte\_de artes; relaciona virtualizacao; relaciona word; parte\_de arte}
\prov{proveniência: wiki \textperiodcentered{} fato \textperiodcentered{} confiança 1.0 \textperiodcentered{} domínio ciencias\_de\_fora \textperiodcentered{} história 5 versões no jornal}

%CRISTAL {"arestas":[{"alvo":"musica","confianca":1.0,"epistemico":"fato","origem":"wiki","rel":"parte_de"},{"alvo":"teoria_da_informacao_shannon","confianca":1.0,"epistemico":"fato","origem":"wiki","rel":"usa"},{"alvo":"hanslick_formalismo","confianca":1.0,"epistemico":"fato","origem":"wiki","rel":"relaciona"},{"alvo":"aparelho_psiquico","confianca":1.0,"epistemico":"fato","origem":"wiki","rel":"relaciona"}],"confianca":1.0,"contraexemplos":[],"descricao":"Por que a música nos move: Meyer atribui a emoção às expectativas criadas e violadas pela estrutura; a neurociência mostra liberação de dopamina no estriado nos picos de prazer, ligada à antecipação e à surpresa.","epistemico":"fato","exemplos":[],"id":"musica_cognicao_emocao","memoria":"semantica","meta":{"autor":"Meyer/Salimpoor","dominio":"ciencias_de_fora","fonte":""},"origem":"wiki","palavras_chave":[],"sinonimos":[],"tipo":"conceito","titulo":"Música, Cognição e Emoção"}
\section{Música, Cognição e Emoção}
Por que a música nos move: Meyer atribui a emoção às expectativas criadas e violadas pela estrutura; a neurociência mostra liberação de dopamina no estriado nos picos de prazer, ligada à antecipação e à surpresa.
\prov{relações: parte\_de musica; usa teoria\_da\_informacao\_shannon; relaciona hanslick\_formalismo; relaciona aparelho\_psiquico}
\prov{proveniência: wiki \textperiodcentered{} fato \textperiodcentered{} confiança 1.0 \textperiodcentered{} domínio ciencias\_de\_fora \textperiodcentered{} história 5 versões no jornal}

%CRISTAL {"arestas":[{"alvo":"harmonia","confianca":1.0,"epistemico":"fato","origem":"wiki","rel":"explica"},{"alvo":"analise_de_fourier","confianca":1.0,"epistemico":"fato","origem":"wiki","rel":"usa"},{"alvo":"matematica","confianca":1.0,"epistemico":"fato","origem":"wiki","rel":"relaciona"}],"confianca":1.0,"contraexemplos":[],"descricao":"A música é razão audível: intervalos são frações, o timbre é Fourier, as séries de 12 tons são grupos. Pitágoras já sabia.","epistemico":"inferencia","exemplos":[],"id":"musica_e_matematica","memoria":"semantica","meta":{"dominio":"ciencias_de_fora","fonte":""},"origem":"wiki","palavras_chave":[],"sinonimos":[],"tipo":"conceito","titulo":"Música e Matemática"}
\section{Música e Matemática}
A música é razão audível: intervalos são frações, o timbre é Fourier, as séries de 12 tons são grupos. Pitágoras já sabia.
\prov{relações: explica harmonia; usa analise\_de\_fourier; relaciona matematica}
\prov{proveniência: wiki \textperiodcentered{} inferencia \textperiodcentered{} confiança 1.0 \textperiodcentered{} domínio ciencias\_de\_fora \textperiodcentered{} história 4 versões no jornal}

%CRISTAL {"arestas":[{"alvo":"espectro_irredutivel_floco","confianca":1.0,"epistemico":"fato","origem":"wiki","rel":"eh_um"},{"alvo":"voz_fonte_filtro_modos","confianca":1.0,"epistemico":"fato","origem":"wiki","rel":"relaciona"},{"alvo":"corpo_estelar","confianca":1.0,"epistemico":"fato","origem":"wiki","rel":"usa"}],"confianca":1.0,"contraexemplos":[],"descricao":"A música é a mesma decomposição em modos da voz e da estrela: o CROMAGRAMA (a energia por classe de nota, somada sobre as oitavas) é o espectro da peça nos 12 modos de altura (C..B). Irmão dos formantes da voz (os músculos-modos) e das camadas da estrela (SVD). A peça é a SOMA dos modos, e a distribuição é o espectro. Verif. em musica.py: um acorde Lá maior dá picos em A, C#, E; uma nota pura domina um modo só (o cristal, como o bem-te-vi).","epistemico":"fato","exemplos":[],"id":"musica_notas_sao_modos","memoria":"semantica","meta":{"dominio":"arte","fonte":"música mineral"},"origem":"wiki","palavras_chave":[],"sinonimos":[],"tipo":"conceito","titulo":"A Música: as Notas são os Modos (o cromagrama)"}
\section{A Música: as Notas são os Modos (o cromagrama)}
A música é a mesma decomposição em modos da voz e da estrela: o CROMAGRAMA (a energia por classe de nota, somada sobre as oitavas) é o espectro da peça nos 12 modos de altura (C..B). Irmão dos formantes da voz (os músculos-modos) e das camadas da estrela (SVD). A peça é a SOMA dos modos, e a distribuição é o espectro. Verif. em musica.py: um acorde Lá maior dá picos em A, C\#, E; uma nota pura domina um modo só (o cristal, como o bem-te-vi).
\prov{relações: eh\_um espectro\_irredutivel\_floco; relaciona voz\_fonte\_filtro\_modos; usa corpo\_estelar}
\prov{proveniência: wiki \textperiodcentered{} música mineral \textperiodcentered{} fato \textperiodcentered{} confiança 1.0 \textperiodcentered{} domínio arte \textperiodcentered{} história 4 versões no jornal}

%CRISTAL {"arestas":[{"alvo":"gene","confianca":1.0,"epistemico":"fato","origem":"wiki","rel":"usa"},{"alvo":"selecao_natural","confianca":1.0,"epistemico":"fato","origem":"wiki","rel":"explica"}],"confianca":1.0,"contraexemplos":[],"descricao":"Uma alteração na sequência do DNA — erro de cópia, dano ou inserção. A fonte crua de toda variação.","epistemico":"fato","exemplos":[],"id":"mutacao","memoria":"semantica","meta":{"dominio":"ciencias_de_fora","fonte":""},"origem":"wiki","palavras_chave":[],"sinonimos":[],"tipo":"conceito","titulo":"Mutação"}
\section{Mutação}
Uma alteração na sequência do DNA — erro de cópia, dano ou inserção. A fonte crua de toda variação.
\prov{relações: usa gene; explica selecao\_natural}
\prov{proveniência: wiki \textperiodcentered{} fato \textperiodcentered{} confiança 1.0 \textperiodcentered{} domínio ciencias\_de\_fora \textperiodcentered{} história 3 versões no jornal}

%CRISTAL {"arestas":[{"alvo":"poetica_aristoteles","confianca":1.0,"epistemico":"fato","origem":"wiki","rel":"parte_de"},{"alvo":"narrativa","confianca":1.0,"epistemico":"fato","origem":"wiki","rel":"explica"},{"alvo":"morfologia_propp","confianca":1.0,"epistemico":"fato","origem":"wiki","rel":"relaciona"}],"confianca":1.0,"contraexemplos":[],"descricao":"Aristóteles: a alma da tragédia é o ENREDO — a ordem necessária dos eventos, não o caráter nem o espetáculo.","epistemico":"inferencia","exemplos":[],"id":"mythos_enredo","memoria":"semantica","meta":{"autor":"Aristóteles","dominio":"ciencias_de_fora","fonte":""},"origem":"wiki","palavras_chave":[],"sinonimos":[],"tipo":"conceito","titulo":"Mythos (o Enredo)"}
\section{Mythos (o Enredo)}
Aristóteles: a alma da tragédia é o ENREDO — a ordem necessária dos eventos, não o caráter nem o espetáculo.
\prov{relações: parte\_de poetica\_aristoteles; explica narrativa; relaciona morfologia\_propp}
\prov{proveniência: wiki \textperiodcentered{} inferencia \textperiodcentered{} confiança 1.0 \textperiodcentered{} domínio ciencias\_de\_fora \textperiodcentered{} história 5 versões no jornal}

%CRISTAL {"arestas":[{"alvo":"metafisica","confianca":1.0,"epistemico":"fato","origem":"wiki","rel":"parte_de"},{"alvo":"psicanalise","confianca":0.5,"epistemico":"fato","origem":"wiki","rel":"relaciona"},{"alvo":"utilitarismo","confianca":0.5,"epistemico":"fato","origem":"wiki","rel":"relaciona"},{"alvo":"vetores","confianca":0.5,"epistemico":"fato","origem":"wiki","rel":"relaciona"},{"alvo":"word","confianca":0.5,"epistemico":"fato","origem":"wiki","rel":"relaciona"}],"confianca":1.0,"contraexemplos":[],"descricao":"A tese de que a separação sujeito/objeto, eu/mundo, é aparente sobre uma realidade una.","epistemico":"hipotese","exemplos":[],"id":"nao_dualidade","memoria":"semantica","meta":{"dominio":"filosofia","fonte":""},"origem":"wiki","palavras_chave":[],"sinonimos":[],"tipo":"conceito","titulo":"Não-Dualidade"}
\section{Não-Dualidade}
A tese de que a separação sujeito/objeto, eu/mundo, é aparente sobre uma realidade una.
\prov{relações: parte\_de metafisica; relaciona psicanalise; relaciona utilitarismo; relaciona vetores; relaciona word}
\prov{proveniência: wiki \textperiodcentered{} hipotese \textperiodcentered{} confiança 1.0 \textperiodcentered{} domínio filosofia \textperiodcentered{} história 6 versões no jornal}

%CRISTAL {"arestas":[{"alvo":"teoria_literaria","confianca":1.0,"epistemico":"fato","origem":"wiki","rel":"parte_de"},{"alvo":"narratologia","confianca":1.0,"epistemico":"fato","origem":"wiki","rel":"relaciona"}],"confianca":1.0,"contraexemplos":[],"descricao":"A arte de contar — a sequência de eventos com causalidade e ponto de vista; a forma em que damos sentido ao tempo.","epistemico":"fato","exemplos":[],"id":"narrativa","memoria":"semantica","meta":{"dominio":"ciencias_de_fora","fonte":""},"origem":"wiki","palavras_chave":[],"sinonimos":[],"tipo":"conceito","titulo":"Narrativa"}
\section{Narrativa}
A arte de contar — a sequência de eventos com causalidade e ponto de vista; a forma em que damos sentido ao tempo.
\prov{relações: parte\_de teoria\_literaria; relaciona narratologia}
\prov{proveniência: wiki \textperiodcentered{} fato \textperiodcentered{} confiança 1.0 \textperiodcentered{} domínio ciencias\_de\_fora \textperiodcentered{} história 4 versões no jornal}

%CRISTAL {"arestas":[{"alvo":"morfologia_propp","confianca":1.0,"epistemico":"fato","origem":"wiki","rel":"usa"},{"alvo":"pacto_ficcional","confianca":1.0,"epistemico":"fato","origem":"wiki","rel":"relaciona"}],"confianca":1.0,"contraexemplos":[],"descricao":"O estudo sistemático da narrativa que distingue a história (fábula, eventos em ordem cronológica) do discurso (trama, sua disposição no texto); Genette formaliza narração, foco/focalização e voz.","epistemico":"inferencia","exemplos":[],"id":"narratologia","memoria":"semantica","meta":{"autor":"Genette","dominio":"ciencias_de_fora","fonte":""},"origem":"wiki","palavras_chave":[],"sinonimos":[],"tipo":"conceito","titulo":"Narratologia (Genette)"}
\section{Narratologia (Genette)}
O estudo sistemático da narrativa que distingue a história (fábula, eventos em ordem cronológica) do discurso (trama, sua disposição no texto); Genette formaliza narração, foco/focalização e voz.
\prov{relações: usa morfologia\_propp; relaciona pacto\_ficcional}
\prov{proveniência: wiki \textperiodcentered{} inferencia \textperiodcentered{} confiança 1.0 \textperiodcentered{} domínio ciencias\_de\_fora \textperiodcentered{} história 3 versões no jornal}

%CRISTAL {"arestas":[{"alvo":"filosofia_da_medicina","confianca":1.0,"epistemico":"fato","origem":"wiki","rel":"parte_de"},{"alvo":"biopoder","confianca":1.0,"epistemico":"fato","origem":"wiki","rel":"usa"},{"alvo":"modelo_biomedico","confianca":1.0,"epistemico":"fato","origem":"wiki","rel":"relaciona"}],"confianca":1.0,"contraexemplos":[],"descricao":"Foucault: no fim do séc. XVIII emerge o 'olhar clínico' — modo objetivante de ver o corpo do paciente, centrado em lesões e sinais visíveis, em vez da experiência subjetiva. Uma reorganização do saber-poder na produção do conhecimento médico.","epistemico":"inferencia","exemplos":[],"id":"nascimento_clinica_foucault","memoria":"semantica","meta":{"autor":"Foucault","dominio":"ciencias_de_fora","fonte":""},"origem":"wiki","palavras_chave":[],"sinonimos":[],"tipo":"conceito","titulo":"Nascimento da Clínica (Foucault)"}
\section{Nascimento da Clínica (Foucault)}
Foucault: no fim do séc. XVIII emerge o 'olhar clínico' — modo objetivante de ver o corpo do paciente, centrado em lesões e sinais visíveis, em vez da experiência subjetiva. Uma reorganização do saber-poder na produção do conhecimento médico.
\prov{relações: parte\_de filosofia\_da\_medicina; usa biopoder; relaciona modelo\_biomedico}
\prov{proveniência: wiki \textperiodcentered{} inferencia \textperiodcentered{} confiança 1.0 \textperiodcentered{} domínio ciencias\_de\_fora \textperiodcentered{} história 4 versões no jornal}

%CRISTAL {"arestas":[{"alvo":"filosofia_da_medicina","confianca":1.0,"epistemico":"fato","origem":"wiki","rel":"parte_de"},{"alvo":"selecao_natural","confianca":1.0,"epistemico":"fato","origem":"wiki","rel":"usa"},{"alvo":"materialismo_fisicalismo","confianca":1.0,"epistemico":"fato","origem":"wiki","rel":"relaciona"}],"confianca":1.0,"contraexemplos":[],"descricao":"Boorse: doença é disfunção biológica objetiva e livre de valores — quando a capacidade funcional de uma parte do organismo cai abaixo do típico da espécie (por idade e sexo). Saúde = funcionamento estatisticamente normal.","epistemico":"inferencia","exemplos":[],"id":"naturalismo_boorse","memoria":"semantica","meta":{"autor":"Boorse","dominio":"ciencias_de_fora","fonte":""},"origem":"wiki","palavras_chave":[],"sinonimos":[],"tipo":"conceito","titulo":"Naturalismo / Teoria Biostatística (Boorse)"}
\section{Naturalismo / Teoria Biostatística (Boorse)}
Boorse: doença é disfunção biológica objetiva e livre de valores — quando a capacidade funcional de uma parte do organismo cai abaixo do típico da espécie (por idade e sexo). Saúde = funcionamento estatisticamente normal.
\prov{relações: parte\_de filosofia\_da\_medicina; usa selecao\_natural; relaciona materialismo\_fisicalismo}
\prov{proveniência: wiki \textperiodcentered{} inferencia \textperiodcentered{} confiança 1.0 \textperiodcentered{} domínio ciencias\_de\_fora \textperiodcentered{} história 4 versões no jornal}

%CRISTAL {"arestas":[{"alvo":"dissipacao_prigogine","confianca":1.0,"epistemico":"fato","origem":"wiki","rel":"usa"},{"alvo":"entropia_topologica","confianca":1.0,"epistemico":"fato","origem":"wiki","rel":"relaciona"},{"alvo":"word","confianca":0.5,"epistemico":"fato","origem":"wiki","rel":"relaciona"},{"alvo":"vontade_de_poder","confianca":0.5,"epistemico":"fato","origem":"wiki","rel":"relaciona"}],"confianca":1.0,"contraexemplos":[],"descricao":"Schrödinger: a vida 'se alimenta de entropia negativa' — extrai ordem do ambiente enquanto aumenta a entropia cósmica.","epistemico":"inferencia","exemplos":[],"id":"neguentropia","memoria":"semantica","meta":{"autor":"Schrödinger","dominio":"ciencias_de_fora","fonte":""},"origem":"wiki","palavras_chave":[],"sinonimos":[],"tipo":"conceito","titulo":"Neguentropia (What is Life?)"}
\section{Neguentropia (What is Life?)}
Schrödinger: a vida 'se alimenta de entropia negativa' — extrai ordem do ambiente enquanto aumenta a entropia cósmica.
\prov{relações: usa dissipacao\_prigogine; relaciona entropia\_topologica; relaciona word; relaciona vontade\_de\_poder}
\prov{proveniência: wiki \textperiodcentered{} inferencia \textperiodcentered{} confiança 1.0 \textperiodcentered{} domínio ciencias\_de\_fora \textperiodcentered{} história 5 versões no jornal}

%CRISTAL {"arestas":[{"alvo":"medicalizacao","confianca":1.0,"epistemico":"fato","origem":"wiki","rel":"especializa"},{"alvo":"biopoder","confianca":1.0,"epistemico":"fato","origem":"wiki","rel":"relaciona"}],"confianca":1.0,"contraexemplos":[],"descricao":"Illich: a medicina industrial expropria a saúde e a autonomia, tornando-se ela própria causa de doença — iatrogênese clínica (dano direto), social (dependência medicalizada) e cultural (destruição da capacidade autônoma de lidar com dor, sofrimento e morte).","epistemico":"inferencia","exemplos":[],"id":"nemesis_medica_illich","memoria":"semantica","meta":{"autor":"Illich","dominio":"ciencias_de_fora","fonte":""},"origem":"wiki","palavras_chave":[],"sinonimos":[],"tipo":"conceito","titulo":"Nêmesis Médica / Iatrogênese (Illich)"}
\section{Nêmesis Médica / Iatrogênese (Illich)}
Illich: a medicina industrial expropria a saúde e a autonomia, tornando-se ela própria causa de doença — iatrogênese clínica (dano direto), social (dependência medicalizada) e cultural (destruição da capacidade autônoma de lidar com dor, sofrimento e morte).
\prov{relações: especializa medicalizacao; relaciona biopoder}
\prov{proveniência: wiki \textperiodcentered{} inferencia \textperiodcentered{} confiança 1.0 \textperiodcentered{} domínio ciencias\_de\_fora \textperiodcentered{} história 3 versões no jornal}

%CRISTAL {"arestas":[{"alvo":"alegoria_caverna","confianca":1.0,"epistemico":"fato","origem":"wiki","rel":"especializa"},{"alvo":"word","confianca":0.5,"epistemico":"fato","origem":"wiki","rel":"relaciona"},{"alvo":"virtualizacao","confianca":0.5,"epistemico":"fato","origem":"wiki","rel":"relaciona"}],"confianca":1.0,"contraexemplos":[],"descricao":"Plotino: do Uno inefável emana o Nous (Formas), a Alma e a Matéria — gradações de realidade, retorno ao Uno.","epistemico":"inferencia","exemplos":[],"id":"neoplatonismo_uno","memoria":"semantica","meta":{"autor":"Plotino","dominio":"filosofia","fonte":""},"origem":"wiki","palavras_chave":[],"sinonimos":[],"tipo":"conceito","titulo":"O Uno e a Emanação"}
\section{O Uno e a Emanação}
Plotino: do Uno inefável emana o Nous (Formas), a Alma e a Matéria — gradações de realidade, retorno ao Uno.
\prov{relações: especializa alegoria\_caverna; relaciona word; relaciona virtualizacao}
\prov{proveniência: wiki \textperiodcentered{} inferencia \textperiodcentered{} confiança 1.0 \textperiodcentered{} domínio filosofia \textperiodcentered{} história 4 versões no jornal}

%CRISTAL {"arestas":[{"alvo":"biologia","confianca":1.0,"epistemico":"fato","origem":"wiki","rel":"parte_de"}],"confianca":1.0,"contraexemplos":[],"descricao":"A ciência do sistema nervoso — do neurônio ao comportamento e à mente.","epistemico":"fato","exemplos":[],"id":"neurociencia","memoria":"semantica","meta":{"dominio":"ciencias_de_fora","fonte":""},"origem":"wiki","palavras_chave":[],"sinonimos":[],"tipo":"conceito","titulo":"Neurociência"}
\section{Neurociência}
A ciência do sistema nervoso — do neurônio ao comportamento e à mente.
\prov{relações: parte\_de biologia}
\prov{proveniência: wiki \textperiodcentered{} fato \textperiodcentered{} confiança 1.0 \textperiodcentered{} domínio ciencias\_de\_fora \textperiodcentered{} história 2 versões no jornal}

%CRISTAL {"arestas":[{"alvo":"neurociencia","confianca":1.0,"epistemico":"fato","origem":"wiki","rel":"parte_de"},{"alvo":"celula","confianca":1.0,"epistemico":"fato","origem":"wiki","rel":"especializa"},{"alvo":"gato","confianca":1.0,"epistemico":"fato","origem":"wiki","rel":"relaciona"},{"alvo":"o_neuronio","confianca":1.0,"epistemico":"fato","origem":"wiki","rel":"relaciona"}],"confianca":1.0,"contraexemplos":[],"descricao":"A célula que sinaliza — recebe, integra e dispara. A unidade do pensamento, e o modelo do neurônio artificial.","epistemico":"fato","exemplos":[],"id":"neuronio","memoria":"semantica","meta":{"dominio":"ciencias_de_fora","fonte":""},"origem":"wiki","palavras_chave":[],"sinonimos":[],"tipo":"conceito","titulo":"Neurônio"}
\section{Neurônio}
A célula que sinaliza — recebe, integra e dispara. A unidade do pensamento, e o modelo do neurônio artificial.
\prov{relações: parte\_de neurociencia; especializa celula; relaciona gato; relaciona o\_neuronio}
\prov{proveniência: wiki \textperiodcentered{} fato \textperiodcentered{} confiança 1.0 \textperiodcentered{} domínio ciencias\_de\_fora \textperiodcentered{} história 5 versões no jornal}

%CRISTAL {"arestas":[{"alvo":"ecossistema","confianca":1.0,"epistemico":"fato","origem":"wiki","rel":"usa"}],"confianca":1.0,"contraexemplos":[],"descricao":"O papel e o lugar de uma espécie no ecossistema — o que faz, come e onde vive; duas espécies não ocupam o mesmo por muito tempo.","epistemico":"fato","exemplos":[],"id":"nicho_ecologico","memoria":"semantica","meta":{"dominio":"ciencias_de_fora","fonte":""},"origem":"wiki","palavras_chave":[],"sinonimos":[],"tipo":"conceito","titulo":"Nicho Ecológico"}
\section{Nicho Ecológico}
O papel e o lugar de uma espécie no ecossistema — o que faz, come e onde vive; duas espécies não ocupam o mesmo por muito tempo.
\prov{relações: usa ecossistema}
\prov{proveniência: wiki \textperiodcentered{} fato \textperiodcentered{} confiança 1.0 \textperiodcentered{} domínio ciencias\_de\_fora \textperiodcentered{} história 2 versões no jornal}

%CRISTAL {"arestas":[{"alvo":"saude_publica","confianca":1.0,"epistemico":"fato","origem":"wiki","rel":"parte_de"},{"alvo":"epidemiologia","confianca":1.0,"epistemico":"fato","origem":"wiki","rel":"usa"}],"confianca":1.0,"contraexemplos":[],"descricao":"Leavell & Clark: a primária evita a ocorrência da doença (promoção e proteção); a secundária a detecta e trata cedo (rastreamento); a terciária limita sequelas e reabilita. Distingue a lógica coletiva da individual.","epistemico":"inferencia","exemplos":[],"id":"niveis_de_prevencao","memoria":"semantica","meta":{"autor":"Leavell & Clark","dominio":"ciencias_de_fora","fonte":""},"origem":"wiki","palavras_chave":[],"sinonimos":[],"tipo":"conceito","titulo":"Prevenção (Primária, Secundária, Terciária)"}
\section{Prevenção (Primária, Secundária, Terciária)}
Leavell \& Clark: a primária evita a ocorrência da doença (promoção e proteção); a secundária a detecta e trata cedo (rastreamento); a terciária limita sequelas e reabilita. Distingue a lógica coletiva da individual.
\prov{relações: parte\_de saude\_publica; usa epidemiologia}
\prov{proveniência: wiki \textperiodcentered{} inferencia \textperiodcentered{} confiança 1.0 \textperiodcentered{} domínio ciencias\_de\_fora \textperiodcentered{} história 3 versões no jornal}

%CRISTAL {"arestas":[{"alvo":"linguistica","confianca":1.0,"epistemico":"fato","origem":"wiki","rel":"parte_de"}],"confianca":1.0,"contraexemplos":[],"descricao":"Fundador da linguística gerativa: a linguagem como capacidade computacional inata da mente.","epistemico":"fato","exemplos":[],"id":"noam_chomsky","memoria":"semantica","meta":{"autor":"Chomsky","dominio":"ciencias_de_fora","fonte":""},"origem":"wiki","palavras_chave":[],"sinonimos":[],"tipo":"conceito","titulo":"Noam Chomsky"}
\section{Noam Chomsky}
Fundador da linguística gerativa: a linguagem como capacidade computacional inata da mente.
\prov{relações: parte\_de linguistica}
\prov{proveniência: wiki \textperiodcentered{} fato \textperiodcentered{} confiança 1.0 \textperiodcentered{} domínio ciencias\_de\_fora \textperiodcentered{} história 2 versões no jornal}

%CRISTAL {"arestas":[{"alvo":"platonismo_matematico","confianca":1.0,"epistemico":"fato","origem":"wiki","rel":"contradiz"},{"alvo":"virtualizacao","confianca":0.5,"epistemico":"fato","origem":"wiki","rel":"relaciona"},{"alvo":"word","confianca":0.5,"epistemico":"fato","origem":"wiki","rel":"relaciona"}],"confianca":1.0,"contraexemplos":[],"descricao":"Não há objetos matemáticos abstratos; só particulares concretos. A matemática é convenção útil (Quine, Field).","epistemico":"inferencia","exemplos":[],"id":"nominalismo_matematico","memoria":"semantica","meta":{"autor":"Field/Quine","dominio":"filosofia","fonte":""},"origem":"wiki","palavras_chave":[],"sinonimos":[],"tipo":"conceito","titulo":"Nominalismo Matemático"}
\section{Nominalismo Matemático}
Não há objetos matemáticos abstratos; só particulares concretos. A matemática é convenção útil (Quine, Field).
\prov{relações: contradiz platonismo\_matematico; relaciona virtualizacao; relaciona word}
\prov{proveniência: wiki \textperiodcentered{} inferencia \textperiodcentered{} confiança 1.0 \textperiodcentered{} domínio filosofia \textperiodcentered{} história 4 versões no jornal}

%CRISTAL {"arestas":[{"alvo":"filosofia_da_medicina","confianca":1.0,"epistemico":"fato","origem":"wiki","rel":"parte_de"},{"alvo":"autopoiese","confianca":1.0,"epistemico":"fato","origem":"wiki","rel":"usa"},{"alvo":"selecao_natural","confianca":1.0,"epistemico":"fato","origem":"wiki","rel":"relaciona"},{"alvo":"normativismo_doenca","confianca":1.0,"epistemico":"fato","origem":"wiki","rel":"relaciona"}],"confianca":1.0,"contraexemplos":[],"descricao":"Canguilhem: o patológico não é um mero desvio quantitativo do normal, mas uma nova qualidade de vida. A vida é normatividade — o vivente institui suas próprias normas; estar saudável é ser normativo (poder criar novas normas frente ao meio), a doença é vida sob normas restritas.","epistemico":"inferencia","exemplos":[],"id":"normal_e_patologico_canguilhem","memoria":"semantica","meta":{"autor":"Canguilhem","dominio":"ciencias_de_fora","fonte":""},"origem":"wiki","palavras_chave":[],"sinonimos":[],"tipo":"conceito","titulo":"O Normal e o Patológico (Canguilhem)"}
\section{O Normal e o Patológico (Canguilhem)}
Canguilhem: o patológico não é um mero desvio quantitativo do normal, mas uma nova qualidade de vida. A vida é normatividade — o vivente institui suas próprias normas; estar saudável é ser normativo (poder criar novas normas frente ao meio), a doença é vida sob normas restritas.
\prov{relações: parte\_de filosofia\_da\_medicina; usa autopoiese; relaciona selecao\_natural; relaciona normativismo\_doenca}
\prov{proveniência: wiki \textperiodcentered{} inferencia \textperiodcentered{} confiança 1.0 \textperiodcentered{} domínio ciencias\_de\_fora \textperiodcentered{} história 5 versões no jornal}

%CRISTAL {"arestas":[{"alvo":"filosofia_da_medicina","confianca":1.0,"epistemico":"fato","origem":"wiki","rel":"parte_de"},{"alvo":"naturalismo_boorse","confianca":1.0,"epistemico":"fato","origem":"wiki","rel":"contradiz"},{"alvo":"biopoder","confianca":1.0,"epistemico":"fato","origem":"wiki","rel":"relaciona"}],"confianca":1.0,"contraexemplos":[],"descricao":"A doença não é fato biológico neutro, mas um desvio julgado indesejável segundo valores sociais; o que conta como patológico depende de normas culturais e estigma. Opõe-se ao naturalismo de Boorse.","epistemico":"inferencia","exemplos":[],"id":"normativismo_doenca","memoria":"semantica","meta":{"autor":"Engelhardt","dominio":"ciencias_de_fora","fonte":""},"origem":"wiki","palavras_chave":[],"sinonimos":[],"tipo":"conceito","titulo":"Normativismo (Doença como Valor)"}
\section{Normativismo (Doença como Valor)}
A doença não é fato biológico neutro, mas um desvio julgado indesejável segundo valores sociais; o que conta como patológico depende de normas culturais e estigma. Opõe-se ao naturalismo de Boorse.
\prov{relações: parte\_de filosofia\_da\_medicina; contradiz naturalismo\_boorse; relaciona biopoder}
\prov{proveniência: wiki \textperiodcentered{} inferencia \textperiodcentered{} confiança 1.0 \textperiodcentered{} domínio ciencias\_de\_fora \textperiodcentered{} história 4 versões no jornal}

%CRISTAL {"arestas":[{"alvo":"direito","confianca":1.0,"epistemico":"fato","origem":"wiki","rel":"parte_de"},{"alvo":"justica_equidade_rawls","confianca":1.0,"epistemico":"fato","origem":"wiki","rel":"contradiz"},{"alvo":"propriedade_lockeana","confianca":1.0,"epistemico":"fato","origem":"wiki","rel":"usa"}],"confianca":1.0,"contraexemplos":[],"descricao":"Nozick: a justiça distributiva é histórica, não padronizada — uma distribuição é justa se resulta de aquisição e transferência legítimas; logo só o Estado mínimo ('guarda-noturno') é legítimo, pois a redistribuição viola direitos individuais. Crítica direta a Rawls.","epistemico":"inferencia","exemplos":[],"id":"nozick_titularidade","memoria":"semantica","meta":{"autor":"Nozick","dominio":"ciencias_de_fora","fonte":""},"origem":"wiki","palavras_chave":[],"sinonimos":[],"tipo":"conceito","titulo":"Justiça como Titularidade (Nozick)"}
\section{Justiça como Titularidade (Nozick)}
Nozick: a justiça distributiva é histórica, não padronizada — uma distribuição é justa se resulta de aquisição e transferência legítimas; logo só o Estado mínimo ('guarda-noturno') é legítimo, pois a redistribuição viola direitos individuais. Crítica direta a Rawls.
\prov{relações: parte\_de direito; contradiz justica\_equidade\_rawls; usa propriedade\_lockeana}
\prov{proveniência: wiki \textperiodcentered{} inferencia \textperiodcentered{} confiança 1.0 \textperiodcentered{} domínio ciencias\_de\_fora \textperiodcentered{} história 4 versões no jornal}

%CRISTAL {"arestas":[{"alvo":"filosofia_da_matematica","confianca":1.0,"epistemico":"fato","origem":"wiki","rel":"parte_de"},{"alvo":"construcao_z_equivalencia","confianca":1.0,"epistemico":"fato","origem":"wiki","rel":"relaciona"}],"confianca":1.0,"contraexemplos":[],"descricao":"Lopes: ninguém definiu rigorosamente o que É um número; confunde-se representação (símbolo) com quantidade — ceticismo sobre a objetividade matemática.","epistemico":"inferencia","exemplos":[],"id":"numero_indefinido","memoria":"semantica","meta":{"autor":"Gentil Lopes","dominio":"filosofia","fonte":""},"origem":"livro","palavras_chave":[],"sinonimos":[],"tipo":"conceito","titulo":"Crítica dos Fundamentos (ZONA)"}
\section{Crítica dos Fundamentos (ZONA)}
Lopes: ninguém definiu rigorosamente o que É um número; confunde-se representação (símbolo) com quantidade — ceticismo sobre a objetividade matemática.
\prov{relações: parte\_de filosofia\_da\_matematica; relaciona construcao\_z\_equivalencia}
\prov{proveniência: livro \textperiodcentered{} inferencia \textperiodcentered{} confiança 1.0 \textperiodcentered{} domínio filosofia \textperiodcentered{} história 4 versões no jornal}

%CRISTAL {"arestas":[{"alvo":"analise_matematica","confianca":1.0,"epistemico":"fato","origem":"wiki","rel":"parte_de"},{"alvo":"conjugados_galois_parabola","confianca":1.0,"epistemico":"fato","origem":"wiki","rel":"relaciona"},{"alvo":"complexos","confianca":1.0,"epistemico":"fato","origem":"wiki","rel":"relaciona"}],"confianca":1.0,"contraexemplos":[],"descricao":"Estender ℝ com i=√−1 fecha a álgebra: todo polinômio tem raiz (Teorema Fundamental da Álgebra, Gauss). Os conjugados vivem aqui.","epistemico":"fato","exemplos":[],"id":"numeros_complexos","memoria":"semantica","meta":{"dominio":"ciencias_de_fora","fonte":""},"origem":"wiki","palavras_chave":[],"sinonimos":[],"tipo":"conceito","titulo":"Números Complexos"}
\section{Números Complexos}
Estender \ensuremath{\mathbb{R}} com i=\ensuremath{\sqrt{\,}}\ensuremath{-}1 fecha a álgebra: todo polinômio tem raiz (Teorema Fundamental da Álgebra, Gauss). Os conjugados vivem aqui.
\prov{relações: parte\_de analise\_matematica; relaciona conjugados\_galois\_parabola; relaciona complexos}
\prov{proveniência: wiki \textperiodcentered{} fato \textperiodcentered{} confiança 1.0 \textperiodcentered{} domínio ciencias\_de\_fora \textperiodcentered{} história 8 versões no jornal}

%CRISTAL {"fusao":[{"arestas":[{"alvo":"teoria_dos_numeros","confianca":1.0,"epistemico":"fato","origem":"wiki","rel":"parte_de"},{"alvo":"bases_pisot","confianca":1.0,"epistemico":"fato","origem":"wiki","rel":"explica"},{"alvo":"beta_numeracao_metalica","confianca":1.0,"epistemico":"fato","origem":"wiki","rel":"explica"},{"alvo":"fibonacci","confianca":1.0,"epistemico":"fato","origem":"wiki","rel":"relaciona"},{"alvo":"numeros_pisot","confianca":1.0,"epistemico":"fato","origem":"wiki","rel":"relaciona"},{"alvo":"cadeias_alem_dos_metais","confianca":1.0,"epistemico":"fato","origem":"wiki","rel":"parte_de"},{"alvo":"lyapunov_da_orbita","confianca":1.0,"epistemico":"fato","origem":"wiki","rel":"relaciona"}],"confianca":1.0,"contraexemplos":[],"descricao":"Inteiros algébricos >1 cujos conjugados têm módulo <1 — os únicos cujas potências tendem a inteiros. A base da β-numeração e dos metálicos.","epistemico":"fato","exemplos":[],"id":"numeros_de_pisot","memoria":"semantica","meta":{"autor":"Pisot","dominio":"ciencias_de_fora","fonte":""},"origem":"wiki","palavras_chave":[],"sinonimos":[],"tipo":"conceito","titulo":"Números de Pisot"},{"arestas":[{"alvo":"razao_aurea","confianca":1.0,"epistemico":"fato","origem":"livro","rel":"relaciona"},{"alvo":"word","confianca":0.5,"epistemico":"fato","origem":"wiki","rel":"relaciona"},{"alvo":"transformada_metalica_fechada_perelman","confianca":1.0,"epistemico":"fato","origem":"wiki","rel":"explica"},{"alvo":"beta_numeracao_metalica","confianca":1.0,"epistemico":"fato","origem":"wiki","rel":"explica"}],"confianca":1.0,"contraexemplos":[],"descricao":"Algébricos > 1 cujos conjugados têm módulo < 1. φ é de Pisot (conjugado −1/φ), o que explica a autossimilaridade do deslocamento de dígitos na base-φ.","epistemico":"fato","exemplos":[],"id":"numeros_pisot","memoria":"semantica","meta":{"autor":"Gentil Lopes","descoberta":false,"dominio":"matematica","fonte":"paper_2_reta_ouro.pdf, Teorema 2.2"},"origem":"paper","palavras_chave":[],"sinonimos":[],"tipo":"conceito","titulo":"Números de Pisot"}],"id":"numeros_de_pisot","tipo":"conceito"}
\section{Números de Pisot}
Inteiros algébricos >1 cujos conjugados têm módulo <1 — os únicos cujas potências tendem a inteiros. A base da \ensuremath{\beta}-numeração e dos metálicos.
\prov{relações: parte\_de teoria\_dos\_numeros; explica bases\_pisot; explica beta\_numeracao\_metalica; relaciona fibonacci; relaciona numeros\_pisot; parte\_de cadeias\_alem\_dos\_metais; relaciona lyapunov\_da\_orbita}
\prov{proveniência: wiki \textperiodcentered{} fato \textperiodcentered{} confiança 1.0 \textperiodcentered{} domínio ciencias\_de\_fora \textperiodcentered{} história 13 versões no jornal}
\prov{a segunda face --- numeros\_pisot:}
Algébricos > 1 cujos conjugados têm módulo < 1. \ensuremath{\varphi} é de Pisot (conjugado \ensuremath{-}1/\ensuremath{\varphi}), o que explica a autossimilaridade do deslocamento de dígitos na base-\ensuremath{\varphi}.
\prov{fusão de curadoria (julgada): absorve numeros\_pisot --- as duas faces acima, intactas no registo}

%CRISTAL {"arestas":[{"alvo":"vida","confianca":1.0,"epistemico":"fato","origem":"wiki","rel":"explica"},{"alvo":"neguentropia","confianca":1.0,"epistemico":"fato","origem":"wiki","rel":"usa"},{"alvo":"dissipacao_prigogine","confianca":1.0,"epistemico":"fato","origem":"wiki","rel":"relaciona"},{"alvo":"termodinamica","confianca":1.0,"epistemico":"fato","origem":"wiki","rel":"usa"}],"confianca":1.0,"contraexemplos":[],"descricao":"Schrödinger: o organismo se mantém ordenado 'sugando ordem' do ambiente — vive de neguentropia contra a segunda lei.","epistemico":"inferencia","exemplos":[],"id":"o_que_e_vida","memoria":"semantica","meta":{"autor":"Schrödinger","dominio":"ciencias_de_fora","fonte":""},"origem":"wiki","palavras_chave":[],"sinonimos":[],"tipo":"conceito","titulo":"O que é a Vida?"}
\section{O que é a Vida?}
Schrödinger: o organismo se mantém ordenado 'sugando ordem' do ambiente — vive de neguentropia contra a segunda lei.
\prov{relações: explica vida; usa neguentropia; relaciona dissipacao\_prigogine; usa termodinamica}
\prov{proveniência: wiki \textperiodcentered{} inferencia \textperiodcentered{} confiança 1.0 \textperiodcentered{} domínio ciencias\_de\_fora \textperiodcentered{} história 5 versões no jornal}

%CRISTAL {"arestas":[{"alvo":"comunicacao_iconoclasta_gentil","confianca":1.0,"epistemico":"fato","origem":"wiki","rel":"usa"},{"alvo":"simbolo_vs_significado","confianca":1.0,"epistemico":"fato","origem":"wiki","rel":"usa"},{"alvo":"deus_mentiroso","confianca":1.0,"epistemico":"fato","origem":"wiki","rel":"relaciona"}],"confianca":1.0,"contraexemplos":[],"descricao":"Lopes enquadra sua função comunicativa pela fábula de Andersen: é 'a criança de fora do métier' que aponta que o Rei está nu, enquanto a corte teme parecer tola. A verdade vem da margem, não da corte — e o núcleo da denúncia é semiótico: 'confundis o símbolo com o ente'.","epistemico":"inferencia","exemplos":[],"id":"o_rei_esta_nu_gentil","memoria":"semantica","meta":{"autor":"Gentil Lopes","dominio":"ciencias_de_fora","fonte":""},"origem":"wiki","palavras_chave":[],"sinonimos":[],"tipo":"conceito","titulo":"O Rei está Nu (Lopes)"}
\section{O Rei está Nu (Lopes)}
Lopes enquadra sua função comunicativa pela fábula de Andersen: é 'a criança de fora do métier' que aponta que o Rei está nu, enquanto a corte teme parecer tola. A verdade vem da margem, não da corte — e o núcleo da denúncia é semiótico: 'confundis o símbolo com o ente'.
\prov{relações: usa comunicacao\_iconoclasta\_gentil; usa simbolo\_vs\_significado; relaciona deus\_mentiroso}
\prov{proveniência: wiki \textperiodcentered{} inferencia \textperiodcentered{} confiança 1.0 \textperiodcentered{} domínio ciencias\_de\_fora \textperiodcentered{} história 5 versões no jornal}

%CRISTAL {"arestas":[{"alvo":"estudos_de_midia","confianca":1.0,"epistemico":"fato","origem":"wiki","rel":"parte_de"},{"alvo":"morte_do_autor","confianca":1.0,"epistemico":"fato","origem":"wiki","rel":"relaciona"},{"alvo":"codificacao_decodificacao","confianca":1.0,"epistemico":"fato","origem":"wiki","rel":"relaciona"},{"alvo":"umberto_eco","confianca":1.0,"epistemico":"fato","origem":"wiki","rel":"parte_de"},{"alvo":"semiose_ilimitada","confianca":1.0,"epistemico":"fato","origem":"wiki","rel":"contradiz"}],"confianca":1.0,"contraexemplos":[],"descricao":"Eco: a obra de arte é estruturalmente ambígua, convida o leitor a completá-la — mas dentro de limites, não ao infinito.","epistemico":"inferencia","exemplos":[],"id":"obra_aberta","memoria":"semantica","meta":{"autor":"Eco","dominio":"ciencias_de_fora","fonte":""},"origem":"wiki","palavras_chave":[],"sinonimos":[],"tipo":"conceito","titulo":"Obra Aberta"}
\section{Obra Aberta}
Eco: a obra de arte é estruturalmente ambígua, convida o leitor a completá-la — mas dentro de limites, não ao infinito.
\prov{relações: parte\_de estudos\_de\_midia; relaciona morte\_do\_autor; relaciona codificacao\_decodificacao; parte\_de umberto\_eco; contradiz semiose\_ilimitada}
\prov{proveniência: wiki \textperiodcentered{} inferencia \textperiodcentered{} confiança 1.0 \textperiodcentered{} domínio ciencias\_de\_fora \textperiodcentered{} história 7 versões no jornal}

%CRISTAL {"arestas":[{"alvo":"epistemologia","confianca":1.0,"epistemico":"fato","origem":"wiki","rel":"parte_de"},{"alvo":"dois_sistemas_kahneman","confianca":1.0,"epistemico":"fato","origem":"wiki","rel":"referencia"}],"confianca":1.0,"contraexemplos":[],"descricao":"Bachelard: um impedimento interno ao próprio ato de conhecer — ideias pré-concebidas e intuições enraizadas na linguagem — que precisa ser rompido para o pensamento avançar. Importado por Brousseau à didática. O Lopes o constata no magistério: o aluno chega vinculando 'número' a contar/medir, o que trava o aprendizado.","epistemico":"inferencia","exemplos":[],"id":"obstaculo_epistemologico","memoria":"semantica","meta":{"autor":"Bachelard","dominio":"ciencias_de_fora","fonte":""},"origem":"wiki","palavras_chave":[],"sinonimos":[],"tipo":"conceito","titulo":"Obstáculo Epistemológico"}
\section{Obstáculo Epistemológico}
Bachelard: um impedimento interno ao próprio ato de conhecer — ideias pré-concebidas e intuições enraizadas na linguagem — que precisa ser rompido para o pensamento avançar. Importado por Brousseau à didática. O Lopes o constata no magistério: o aluno chega vinculando 'número' a contar/medir, o que trava o aprendizado.
\prov{relações: parte\_de epistemologia; referencia dois\_sistemas\_kahneman}
\prov{proveniência: wiki \textperiodcentered{} inferencia \textperiodcentered{} confiança 1.0 \textperiodcentered{} domínio ciencias\_de\_fora \textperiodcentered{} história 4 versões no jornal}

%CRISTAL {"arestas":[{"alvo":"ben10_desenho","confianca":1.0,"epistemico":"fato","origem":"wiki","rel":"parte_de"},{"alvo":"omnitrix_corpo_universal","confianca":1.0,"epistemico":"fato","origem":"wiki","rel":"relaciona"}],"confianca":1.0,"contraexemplos":[],"descricao":"O OMNITRIX de Ben 10 (o dispositivo): um relógio que, ao discar e bater, transforma o portador em qualquer alienígena guardado no seu Codon Stream. Cada alien é uma FORMA; o núcleo de energia recarrega; o timeout devolve o humano. É a mesma FIGURA do Omnitrix do broca-so (o Corpo Algébrico): um TEMPLATE que vira qualquer corpo ao escolher uma instância. Por isso o nome — o dicionário abaixo casa os dois.","epistemico":"fato","exemplos":[],"id":"omnitrix_ben10","memoria":"semantica","meta":{"dominio":"ficcao","fonte":"Ficção (jogos e animações)"},"origem":"wiki","palavras_chave":[],"sinonimos":[],"tipo":"conceito","titulo":"O Omnitrix (Ben 10) — o dispositivo que instancia qualquer forma"}
\section{O Omnitrix (Ben 10) — o dispositivo que instancia qualquer forma}
O OMNITRIX de Ben 10 (o dispositivo): um relógio que, ao discar e bater, transforma o portador em qualquer alienígena guardado no seu Codon Stream. Cada alien é uma FORMA; o núcleo de energia recarrega; o timeout devolve o humano. É a mesma FIGURA do Omnitrix do broca-so (o Corpo Algébrico): um TEMPLATE que vira qualquer corpo ao escolher uma instância. Por isso o nome — o dicionário abaixo casa os dois.
\prov{relações: parte\_de ben10\_desenho; relaciona omnitrix\_corpo\_universal}
\prov{proveniência: wiki \textperiodcentered{} Ficção (jogos e animações) \textperiodcentered{} fato \textperiodcentered{} confiança 1.0 \textperiodcentered{} domínio ficcao \textperiodcentered{} história 3 versões no jornal}

%CRISTAL {"arestas":[{"alvo":"tecido_mineral_hub","confianca":1.0,"epistemico":"fato","origem":"wiki","rel":"parte_de"},{"alvo":"bai_biblioteca_de_autorizacao_isomorfica","confianca":1.0,"epistemico":"fato","origem":"wiki","rel":"usa"},{"alvo":"potencial_de_acao","confianca":1.0,"epistemico":"fato","origem":"wiki","rel":"usa"}],"confianca":1.0,"contraexemplos":[],"descricao":"Células excitáveis acopladas propagam o potencial de ação como uma ONDA (meio excitável, Greenberg-Hastings): uma dispara, excita a vizinha, e a frente viaja com velocidade CONSTANTE. É a ativação que percorre o músculo (a contração muscular, o batimento cardíaco) — a mesma PROPAGAÇÃO P do BAI, com o período refratário fazendo o papel da atenuação T: a onda avança e não reverte. Verif.: a excitação vai da célula 0 ao fim, ~1 passo/célula.","epistemico":"fato","exemplos":[],"id":"onda_de_excitacao_e_a_propagacao","memoria":"semantica","meta":{"dominio":"biologia","fonte":"tecido mineral"},"origem":"wiki","palavras_chave":[],"sinonimos":[],"tipo":"conceito","titulo":"A Onda de Excitação é a Propagação (P do BAI)"}
\section{A Onda de Excitação é a Propagação (P do BAI)}
Células excitáveis acopladas propagam o potencial de ação como uma ONDA (meio excitável, Greenberg-Hastings): uma dispara, excita a vizinha, e a frente viaja com velocidade CONSTANTE. É a ativação que percorre o músculo (a contração muscular, o batimento cardíaco) — a mesma PROPAGAÇÃO P do BAI, com o período refratário fazendo o papel da atenuação T: a onda avança e não reverte. Verif.: a excitação vai da célula 0 ao fim, \~{}1 passo/célula.
\prov{relações: parte\_de tecido\_mineral\_hub; usa bai\_biblioteca\_de\_autorizacao\_isomorfica; usa potencial\_de\_acao}
\prov{proveniência: wiki \textperiodcentered{} tecido mineral \textperiodcentered{} fato \textperiodcentered{} confiança 1.0 \textperiodcentered{} domínio biologia \textperiodcentered{} história 4 versões no jornal}

%CRISTAL {"arestas":[{"alvo":"poliarquia","confianca":1.0,"epistemico":"fato","origem":"wiki","rel":"usa"},{"alvo":"ciencia_politica","confianca":1.0,"epistemico":"fato","origem":"wiki","rel":"parte_de"}],"confianca":1.0,"contraexemplos":[],"descricao":"Huntington: transições concentradas de regimes não democráticos para democráticos que superam as transições opostas — três ondas (séc. XIX; pós-1945; desde 1974), cada uma seguida por uma 'onda reversa' de recessão democrática.","epistemico":"fato","exemplos":[],"id":"ondas_democratizacao","memoria":"semantica","meta":{"autor":"Huntington","dominio":"ciencias_de_fora","fonte":""},"origem":"wiki","palavras_chave":[],"sinonimos":[],"tipo":"conceito","titulo":"Ondas de Democratização (Huntington)"}
\section{Ondas de Democratização (Huntington)}
Huntington: transições concentradas de regimes não democráticos para democráticos que superam as transições opostas — três ondas (séc. XIX; pós-1945; desde 1974), cada uma seguida por uma 'onda reversa' de recessão democrática.
\prov{relações: usa poliarquia; parte\_de ciencia\_politica}
\prov{proveniência: wiki \textperiodcentered{} fato \textperiodcentered{} confiança 1.0 \textperiodcentered{} domínio ciencias\_de\_fora \textperiodcentered{} história 3 versões no jornal}

%CRISTAL {"arestas":[{"alvo":"metafisica","confianca":1.0,"epistemico":"fato","origem":"wiki","rel":"parte_de"},{"alvo":"utilitarismo","confianca":0.5,"epistemico":"fato","origem":"wiki","rel":"relaciona"},{"alvo":"teoria_do_valor","confianca":0.5,"epistemico":"fato","origem":"wiki","rel":"relaciona"},{"alvo":"virtualizacao","confianca":0.5,"epistemico":"fato","origem":"wiki","rel":"relaciona"},{"alvo":"word","confianca":0.5,"epistemico":"fato","origem":"wiki","rel":"relaciona"},{"alvo":"vida-morte","confianca":0.5,"epistemico":"fato","origem":"wiki","rel":"relaciona"},{"alvo":"vontade_de_poder","confianca":0.5,"epistemico":"fato","origem":"wiki","rel":"relaciona"}],"confianca":1.0,"contraexemplos":[],"descricao":"O ramo da metafísica que pergunta o que existe e em que modos as coisas são.","epistemico":"fato","exemplos":[],"id":"ontologia","memoria":"semantica","meta":{"dominio":"filosofia","fonte":""},"origem":"wiki","palavras_chave":[],"sinonimos":[],"tipo":"conceito","titulo":"Ontologia"}
\section{Ontologia}
O ramo da metafísica que pergunta o que existe e em que modos as coisas são.
\prov{relações: parte\_de metafisica; relaciona utilitarismo; relaciona teoria\_do\_valor; relaciona virtualizacao; relaciona word; relaciona vida-morte; relaciona vontade\_de\_poder}
\prov{proveniência: wiki \textperiodcentered{} fato \textperiodcentered{} confiança 1.0 \textperiodcentered{} domínio filosofia \textperiodcentered{} história 8 versões no jornal}

%CRISTAL {"arestas":[{"alvo":"camadas_da_face","confianca":1.0,"epistemico":"fato","origem":"wiki","rel":"parte_de"},{"alvo":"orbitas_da_face_no_modelo","confianca":1.0,"epistemico":"fato","origem":"wiki","rel":"relaciona"},{"alvo":"bai_biblioteca_de_autorizacao_isomorfica","confianca":1.0,"epistemico":"fato","origem":"wiki","rel":"usa"},{"alvo":"onda_de_excitacao_e_a_propagacao","confianca":1.0,"epistemico":"fato","origem":"wiki","rel":"relaciona"}],"confianca":1.0,"contraexemplos":[],"descricao":"O strain do músculo CRUZA os tecidos até a pele, ATENUADO pela profundidade (T⪯1, monótona — a propagação P do BAI na profundidade) e com ATRASO de fase (a onda cruza a gordura). É a órbita TRANSVERSAL, perpendicular à superfície, distinta da onda longitudinal do tecido: cada camada responde com a sua órbita (amplitude atenuada), a pele segue o músculo suavizado, o osso ancora (não se move). Fecha a volta ao repouso (borda, ρ=1). Verif.: pele>gordura>0, osso=0; a pele responde depois do músculo.","epistemico":"fato","exemplos":[],"id":"orbita_transversal_entre_tecidos","memoria":"semantica","meta":{"dominio":"biologia","fonte":"camadas da face"},"origem":"wiki","palavras_chave":[],"sinonimos":[],"tipo":"conceito","titulo":"A Órbita Transversal (entre os tecidos)"}
\section{A Órbita Transversal (entre os tecidos)}
O strain do músculo CRUZA os tecidos até a pele, ATENUADO pela profundidade (T\ensuremath{\preceq}1, monótona — a propagação P do BAI na profundidade) e com ATRASO de fase (a onda cruza a gordura). É a órbita TRANSVERSAL, perpendicular à superfície, distinta da onda longitudinal do tecido: cada camada responde com a sua órbita (amplitude atenuada), a pele segue o músculo suavizado, o osso ancora (não se move). Fecha a volta ao repouso (borda, \ensuremath{\rho}=1). Verif.: pele>gordura>0, osso=0; a pele responde depois do músculo.
\prov{relações: parte\_de camadas\_da\_face; relaciona orbitas\_da\_face\_no\_modelo; usa bai\_biblioteca\_de\_autorizacao\_isomorfica; relaciona onda\_de\_excitacao\_e\_a\_propagacao}
\prov{proveniência: wiki \textperiodcentered{} camadas da face \textperiodcentered{} fato \textperiodcentered{} confiança 1.0 \textperiodcentered{} domínio biologia \textperiodcentered{} história 5 versões no jornal}

%CRISTAL {"arestas":[{"alvo":"atuadores_da_face","confianca":1.0,"epistemico":"fato","origem":"wiki","rel":"parte_de"},{"alvo":"orbita_de_partida","confianca":1.0,"epistemico":"fato","origem":"wiki","rel":"eh_um"},{"alvo":"lyapunov_da_orbita","confianca":1.0,"epistemico":"fato","origem":"wiki","rel":"usa"},{"alvo":"parabola_isa","confianca":1.0,"epistemico":"fato","origem":"wiki","rel":"usa"},{"alvo":"orbita_temporal_das_camadas","confianca":1.0,"epistemico":"fato","origem":"wiki","rel":"relaciona"}],"confianca":1.0,"contraexemplos":[],"descricao":"O movimento facial é uma ÓRBITA na parábola cristal/borda/caos, cristalizada no modelo real (model/Hitem3d-1783188908717.glb). CONSTANTE = expressão mantida (cristal, ponto fixo). PERIÓDICA = vai-e-volta que FECHA a volta ao repouso (borda, ρ=1 — piscar, mastigar, sorrir-e-relaxar). CRESCENTE = descontrolada (caos, não fecha). A face se anima percorrendo a mesma parábola das órbitas da máquina (reta mineral / BAI / Newton), agora sobre a malha 3D: cada expressão é uma trajetória de ativação muscular, e a sua classe é o regime da órbita. Verif.: periódica fecha, crescente não.","epistemico":"inferencia","exemplos":[],"id":"orbitas_da_face_no_modelo","memoria":"semantica","meta":{"dominio":"biologia","fonte":"atuadores da face"},"origem":"wiki","palavras_chave":[],"sinonimos":[],"tipo":"conceito","titulo":"As Órbitas da Face (cristalizadas no modelo)"}
\section{As Órbitas da Face (cristalizadas no modelo)}
O movimento facial é uma ÓRBITA na parábola cristal/borda/caos, cristalizada no modelo real (model/Hitem3d-1783188908717.glb). CONSTANTE = expressão mantida (cristal, ponto fixo). PERIÓDICA = vai-e-volta que FECHA a volta ao repouso (borda, \ensuremath{\rho}=1 — piscar, mastigar, sorrir-e-relaxar). CRESCENTE = descontrolada (caos, não fecha). A face se anima percorrendo a mesma parábola das órbitas da máquina (reta mineral / BAI / Newton), agora sobre a malha 3D: cada expressão é uma trajetória de ativação muscular, e a sua classe é o regime da órbita. Verif.: periódica fecha, crescente não.
\prov{relações: parte\_de atuadores\_da\_face; eh\_um orbita\_de\_partida; usa lyapunov\_da\_orbita; usa parabola\_isa; relaciona orbita\_temporal\_das\_camadas}
\prov{proveniência: wiki \textperiodcentered{} atuadores da face \textperiodcentered{} inferencia \textperiodcentered{} confiança 1.0 \textperiodcentered{} domínio biologia \textperiodcentered{} história 6 versões no jornal}

%CRISTAL {"arestas":[{"alvo":"atuadores_da_face","confianca":1.0,"epistemico":"fato","origem":"wiki","rel":"usa"},{"alvo":"dicionario_musculos_face","confianca":1.0,"epistemico":"fato","origem":"wiki","rel":"usa"},{"alvo":"orbitas_da_face_no_modelo","confianca":1.0,"epistemico":"fato","origem":"wiki","rel":"eh_um"},{"alvo":"bai_psi_colapso","confianca":1.0,"epistemico":"fato","origem":"wiki","rel":"usa"}],"confianca":1.0,"contraexemplos":[],"descricao":"Falar é lip-sync: cada fonema tem uma forma de boca (o VISEMA), combinação dos músculos da boca (orbicular da boca, zigomático, depressor, masseter, levantador). Falar é PERCORRER uma sequência de visemas no tempo — uma ÓRBITA (linguagem/fala.py, 7/7). AA (aberto)=masseter, OO (arredondado)=orbicular, EE (esticado)=zigomático, MM (fechado)=lábios juntos. Sincronizar interpola entre visemas (coarticulação). A trajetória da boca é uma órbita na parábola: vogal mantida=cristal, fala=borda (a boca oscila e FECHA a volta no silêncio), cada visema um colapso Ψ. Verif.: 'papa' alterna MM/AA; a órbita fecha a volta.","epistemico":"fato","exemplos":[],"id":"orbitas_de_fala","memoria":"semantica","meta":{"dominio":"biologia","fonte":"órbitas de fala"},"origem":"wiki","palavras_chave":[],"sinonimos":[],"tipo":"conceito","titulo":"As Órbitas de Fala (os músculos da boca sincronizados)"}
\section{As Órbitas de Fala (os músculos da boca sincronizados)}
Falar é lip-sync: cada fonema tem uma forma de boca (o VISEMA), combinação dos músculos da boca (orbicular da boca, zigomático, depressor, masseter, levantador). Falar é PERCORRER uma sequência de visemas no tempo — uma ÓRBITA (linguagem/fala.py, 7/7). AA (aberto)=masseter, OO (arredondado)=orbicular, EE (esticado)=zigomático, MM (fechado)=lábios juntos. Sincronizar interpola entre visemas (coarticulação). A trajetória da boca é uma órbita na parábola: vogal mantida=cristal, fala=borda (a boca oscila e FECHA a volta no silêncio), cada visema um colapso \ensuremath{\Psi}. Verif.: 'papa' alterna MM/AA; a órbita fecha a volta.
\prov{relações: usa atuadores\_da\_face; usa dicionario\_musculos\_face; eh\_um orbitas\_da\_face\_no\_modelo; usa bai\_psi\_colapso}
\prov{proveniência: wiki \textperiodcentered{} órbitas de fala \textperiodcentered{} fato \textperiodcentered{} confiança 1.0 \textperiodcentered{} domínio biologia \textperiodcentered{} história 10 versões no jornal}

%CRISTAL {"arestas":[{"alvo":"von_neumann_wigner","confianca":1.0,"epistemico":"fato","origem":"wiki","rel":"especializa"},{"alvo":"qubit_hopf","confianca":1.0,"epistemico":"fato","origem":"wiki","rel":"referencia"}],"confianca":1.0,"contraexemplos":[],"descricao":"Colapso objetivo de qubits em microtúbulos ao atingir um limiar gravitacional gera momentos discretos de experiência. Controversa (decoerência rápida).","epistemico":"hipotese","exemplos":[],"id":"orch_or","memoria":"semantica","meta":{"autor":"Penrose/Hameroff","dominio":"filosofia","fonte":""},"origem":"wiki","palavras_chave":[],"sinonimos":[],"tipo":"conceito","titulo":"Orch-OR (Penrose-Hameroff)"}
\section{Orch-OR (Penrose-Hameroff)}
Colapso objetivo de qubits em microtúbulos ao atingir um limiar gravitacional gera momentos discretos de experiência. Controversa (decoerência rápida).
\prov{relações: especializa von\_neumann\_wigner; referencia qubit\_hopf}
\prov{proveniência: wiki \textperiodcentered{} hipotese \textperiodcentered{} confiança 1.0 \textperiodcentered{} domínio filosofia \textperiodcentered{} história 3 versões no jornal}

%CRISTAL {"arestas":[{"alvo":"vida","confianca":1.0,"epistemico":"fato","origem":"wiki","rel":"explica"},{"alvo":"borda_critica_vazio","confianca":1.0,"epistemico":"fato","origem":"wiki","rel":"relaciona"},{"alvo":"transicao_fase_lambda1","confianca":1.0,"epistemico":"fato","origem":"wiki","rel":"relaciona"},{"alvo":"hopfield","confianca":1.0,"epistemico":"fato","origem":"wiki","rel":"referencia"}],"confianca":1.0,"contraexemplos":[],"descricao":"Kauffman: na transição entre ordem congelada e caos turbulento, redes complexas maximizam computação e adaptação — a vida busca essa borda.","epistemico":"inferencia","exemplos":[],"id":"ordem_borda_caos","memoria":"semantica","meta":{"autor":"Kauffman","dominio":"ciencias_de_fora","fonte":""},"origem":"wiki","palavras_chave":[],"sinonimos":[],"tipo":"conceito","titulo":"Ordem na Borda do Caos"}
\section{Ordem na Borda do Caos}
Kauffman: na transição entre ordem congelada e caos turbulento, redes complexas maximizam computação e adaptação — a vida busca essa borda.
\prov{relações: explica vida; relaciona borda\_critica\_vazio; relaciona transicao\_fase\_lambda1; referencia hopfield}
\prov{proveniência: wiki \textperiodcentered{} inferencia \textperiodcentered{} confiança 1.0 \textperiodcentered{} domínio ciencias\_de\_fora \textperiodcentered{} história 5 versões no jornal}

%CRISTAL {"arestas":[{"alvo":"sintaxe","confianca":1.0,"epistemico":"fato","origem":"wiki","rel":"parte_de"},{"alvo":"universais_linguisticos","confianca":1.0,"epistemico":"fato","origem":"wiki","rel":"usa"}],"confianca":1.0,"contraexemplos":[],"descricao":"A ordem básica de Sujeito, Verbo e Objeto: SOV (~45% das línguas: japonês, latim), SVO (~42%: português, inglês, mandarim), VSO (~9%: árabe). O primeiro eixo que a tradução tem de reordenar.","epistemico":"inferencia","exemplos":[],"id":"ordem_das_palavras","memoria":"semantica","meta":{"autor":"Greenberg","dominio":"ciencias_de_fora","fonte":""},"origem":"wiki","palavras_chave":[],"sinonimos":[],"tipo":"conceito","titulo":"Ordem das Palavras (Tipologia)"}
\section{Ordem das Palavras (Tipologia)}
A ordem básica de Sujeito, Verbo e Objeto: SOV (\~{}45\% das línguas: japonês, latim), SVO (\~{}42\%: português, inglês, mandarim), VSO (\~{}9\%: árabe). O primeiro eixo que a tradução tem de reordenar.
\prov{relações: parte\_de sintaxe; usa universais\_linguisticos}
\prov{proveniência: wiki \textperiodcentered{} inferencia \textperiodcentered{} confiança 1.0 \textperiodcentered{} domínio ciencias\_de\_fora \textperiodcentered{} história 3 versões no jornal}

%CRISTAL {"arestas":[{"alvo":"economia","confianca":1.0,"epistemico":"fato","origem":"wiki","rel":"parte_de"},{"alvo":"demanda_agregada","confianca":1.0,"epistemico":"fato","origem":"wiki","rel":"contradiz"},{"alvo":"goldchain","confianca":1.0,"epistemico":"fato","origem":"wiki","rel":"referencia"}],"confianca":1.0,"contraexemplos":[],"descricao":"Hayek: o conhecimento econômico está disperso; os preços o coordenam sem centro. O planejamento central falha — casa com a descentralização do goldchain.","epistemico":"inferencia","exemplos":[],"id":"ordem_espontanea","memoria":"semantica","meta":{"autor":"Hayek","dominio":"ciencias_de_fora","fonte":""},"origem":"wiki","palavras_chave":[],"sinonimos":[],"tipo":"conceito","titulo":"Ordem Espontânea e Conhecimento Disperso"}
\section{Ordem Espontânea e Conhecimento Disperso}
Hayek: o conhecimento econômico está disperso; os preços o coordenam sem centro. O planejamento central falha — casa com a descentralização do goldchain.
\prov{relações: parte\_de economia; contradiz demanda\_agregada; referencia goldchain}
\prov{proveniência: wiki \textperiodcentered{} inferencia \textperiodcentered{} confiança 1.0 \textperiodcentered{} domínio ciencias\_de\_fora \textperiodcentered{} história 4 versões no jornal}

%CRISTAL {"arestas":[{"alvo":"vida","confianca":1.0,"epistemico":"fato","origem":"wiki","rel":"relaciona"},{"alvo":"mundo_de_rna","confianca":1.0,"epistemico":"fato","origem":"wiki","rel":"usa"},{"alvo":"abiogenese","confianca":1.0,"epistemico":"fato","origem":"wiki","rel":"relaciona"}],"confianca":1.0,"contraexemplos":[],"descricao":"Como a matéria não-viva virou viva — química prebiótica a caminho de sistemas que copiam e metabolizam. Aberto.","epistemico":"hipotese","exemplos":[],"id":"origem_da_vida","memoria":"semantica","meta":{"dominio":"ciencias_de_fora","fonte":""},"origem":"wiki","palavras_chave":[],"sinonimos":[],"tipo":"conceito","titulo":"Origem da Vida (Abiogênese)"}
\section{Origem da Vida (Abiogênese)}
Como a matéria não-viva virou viva — química prebiótica a caminho de sistemas que copiam e metabolizam. Aberto.
\prov{relações: relaciona vida; usa mundo\_de\_rna; relaciona abiogenese}
\prov{proveniência: wiki \textperiodcentered{} hipotese \textperiodcentered{} confiança 1.0 \textperiodcentered{} domínio ciencias\_de\_fora \textperiodcentered{} história 4 versões no jornal}

%CRISTAL {"arestas":[{"alvo":"autoritarismo_vs_totalitarismo","confianca":1.0,"epistemico":"fato","origem":"wiki","rel":"especializa"},{"alvo":"deus_mentiroso","confianca":1.0,"epistemico":"fato","origem":"wiki","rel":"relaciona"},{"alvo":"biopoder","confianca":1.0,"epistemico":"fato","origem":"wiki","rel":"relaciona"}],"confianca":1.0,"contraexemplos":[],"descricao":"Arendt: a forma política inédita do séc. XX (nazismo, stalinismo) que combina ideologia total e terror não para dominar, mas para transformar a natureza humana e tornar os indivíduos supérfluos — o terror como fim em si. Dissolve a distinção entre fato e ficção.","epistemico":"inferencia","exemplos":[],"id":"origens_totalitarismo","memoria":"semantica","meta":{"autor":"Arendt","dominio":"ciencias_de_fora","fonte":""},"origem":"wiki","palavras_chave":[],"sinonimos":[],"tipo":"conceito","titulo":"As Origens do Totalitarismo (Arendt)"}
\section{As Origens do Totalitarismo (Arendt)}
Arendt: a forma política inédita do séc. XX (nazismo, stalinismo) que combina ideologia total e terror não para dominar, mas para transformar a natureza humana e tornar os indivíduos supérfluos — o terror como fim em si. Dissolve a distinção entre fato e ficção.
\prov{relações: especializa autoritarismo\_vs\_totalitarismo; relaciona deus\_mentiroso; relaciona biopoder}
\prov{proveniência: wiki \textperiodcentered{} inferencia \textperiodcentered{} confiança 1.0 \textperiodcentered{} domínio ciencias\_de\_fora \textperiodcentered{} história 4 versões no jornal}

%CRISTAL {"arestas":[{"alvo":"arquitetura","confianca":1.0,"epistemico":"fato","origem":"wiki","rel":"parte_de"},{"alvo":"forma_segue_funcao","confianca":1.0,"epistemico":"fato","origem":"wiki","rel":"usa"}],"confianca":1.0,"contraexemplos":[],"descricao":"Loos: o ornamento gratuito é atraso — a beleza está na forma limpa. O manifesto do despojamento moderno.","epistemico":"inferencia","exemplos":[],"id":"ornamento_e_crime","memoria":"semantica","meta":{"autor":"Loos","dominio":"ciencias_de_fora","fonte":""},"origem":"wiki","palavras_chave":[],"sinonimos":[],"tipo":"conceito","titulo":"Ornamento e Crime"}
\section{Ornamento e Crime}
Loos: o ornamento gratuito é atraso — a beleza está na forma limpa. O manifesto do despojamento moderno.
\prov{relações: parte\_de arquitetura; usa forma\_segue\_funcao}
\prov{proveniência: wiki \textperiodcentered{} inferencia \textperiodcentered{} confiança 1.0 \textperiodcentered{} domínio ciencias\_de\_fora \textperiodcentered{} história 3 versões no jornal}

%CRISTAL {"arestas":[{"alvo":"filosofia_da_tecnica","confianca":1.0,"epistemico":"fato","origem":"wiki","rel":"parte_de"}],"confianca":1.0,"contraexemplos":[],"descricao":"Ortega y Gasset: o humano é 'o ser para quem o supérfluo é necessário'; a técnica cria possibilidades que não existem na natureza — a reforma que o homem impõe ao mundo para realizar seu projeto de vida, buscando o bem-estar, não o mero estar.","epistemico":"inferencia","exemplos":[],"id":"ortega_tecnica_superfluo","memoria":"semantica","meta":{"autor":"Ortega y Gasset","dominio":"ciencias_de_fora","fonte":""},"origem":"wiki","palavras_chave":[],"sinonimos":[],"tipo":"conceito","titulo":"A Técnica e o Supérfluo (Ortega)"}
\section{A Técnica e o Supérfluo (Ortega)}
Ortega y Gasset: o humano é 'o ser para quem o supérfluo é necessário'; a técnica cria possibilidades que não existem na natureza — a reforma que o homem impõe ao mundo para realizar seu projeto de vida, buscando o bem-estar, não o mero estar.
\prov{relações: parte\_de filosofia\_da\_tecnica}
\prov{proveniência: wiki \textperiodcentered{} inferencia \textperiodcentered{} confiança 1.0 \textperiodcentered{} domínio ciencias\_de\_fora \textperiodcentered{} história 2 versões no jornal}

%CRISTAL {"arestas":[{"alvo":"literatura","confianca":1.0,"epistemico":"fato","origem":"wiki","rel":"parte_de"},{"alvo":"mimese_expressao","confianca":1.0,"epistemico":"fato","origem":"wiki","rel":"usa"},{"alvo":"obra_aberta","confianca":1.0,"epistemico":"fato","origem":"wiki","rel":"relaciona"}],"confianca":1.0,"contraexemplos":[],"descricao":"A ficção instaura um acordo tácito — a 'suspensão voluntária da descrença' (Coleridge) — pelo qual o leitor aceita as premissas de um mundo possível sem exigir correspondência factual.","epistemico":"inferencia","exemplos":[],"id":"pacto_ficcional","memoria":"semantica","meta":{"autor":"Coleridge/Eco","dominio":"ciencias_de_fora","fonte":""},"origem":"wiki","palavras_chave":[],"sinonimos":[],"tipo":"conceito","titulo":"Pacto Ficcional e Mundos Possíveis"}
\section{Pacto Ficcional e Mundos Possíveis}
A ficção instaura um acordo tácito — a 'suspensão voluntária da descrença' (Coleridge) — pelo qual o leitor aceita as premissas de um mundo possível sem exigir correspondência factual.
\prov{relações: parte\_de literatura; usa mimese\_expressao; relaciona obra\_aberta}
\prov{proveniência: wiki \textperiodcentered{} inferencia \textperiodcentered{} confiança 1.0 \textperiodcentered{} domínio ciencias\_de\_fora \textperiodcentered{} história 4 versões no jornal}

%CRISTAL {"arestas":[{"alvo":"tecido_mineral_hub","confianca":1.0,"epistemico":"fato","origem":"wiki","rel":"parte_de"},{"alvo":"morfogenese_turing","confianca":1.0,"epistemico":"fato","origem":"wiki","rel":"eh_um"},{"alvo":"numero_de_salem","confianca":1.0,"epistemico":"fato","origem":"wiki","rel":"relaciona"}],"confianca":1.0,"contraexemplos":[],"descricao":"A morfogênese: ativador + inibidor com difusões diferentes (D_inib ≫ D_ativ) tornam o estado uniforme INSTÁVEL para uma banda de números de onda — e o PADRÃO emerge (manchas, listras, dedos). O estado é estável sem difusão mas a difusão o quebra: a instabilidade de Turing é a BORDA da máquina, onde a simetria quebra e a forma aparece. Precisa da DIFERENÇA de difusão (difusões iguais não dão padrão). Verif.: D_inib≫D_ativ → padrão; iguais → nada.","epistemico":"fato","exemplos":[],"id":"padroes_de_turing_e_a_borda","memoria":"semantica","meta":{"dominio":"biologia","fonte":"tecido mineral"},"origem":"wiki","palavras_chave":[],"sinonimos":[],"tipo":"conceito","titulo":"Os Padrões de Turing são a Borda (a forma nasce)"}
\section{Os Padrões de Turing são a Borda (a forma nasce)}
A morfogênese: ativador + inibidor com difusões diferentes (D\_inib \ensuremath{\gg} D\_ativ) tornam o estado uniforme INSTÁVEL para uma banda de números de onda — e o PADRÃO emerge (manchas, listras, dedos). O estado é estável sem difusão mas a difusão o quebra: a instabilidade de Turing é a BORDA da máquina, onde a simetria quebra e a forma aparece. Precisa da DIFERENÇA de difusão (difusões iguais não dão padrão). Verif.: D\_inib\ensuremath{\gg}D\_ativ \ensuremath{\to} padrão; iguais \ensuremath{\to} nada.
\prov{relações: parte\_de tecido\_mineral\_hub; eh\_um morfogenese\_turing; relaciona numero\_de\_salem}
\prov{proveniência: wiki \textperiodcentered{} tecido mineral \textperiodcentered{} fato \textperiodcentered{} confiança 1.0 \textperiodcentered{} domínio biologia \textperiodcentered{} história 4 versões no jornal}

%CRISTAL {"arestas":[{"alvo":"conscientizacao","confianca":1.0,"epistemico":"fato","origem":"wiki","rel":"usa"},{"alvo":"capital_cultural","confianca":1.0,"epistemico":"fato","origem":"wiki","rel":"relaciona"}],"confianca":1.0,"contraexemplos":[],"descricao":"Freire: método de alfabetização que parte de palavras carregadas de sentido existencial e político para o educando (de seu universo vocabular), gerando leitura crítica do mundo, não mera decodificação. Historicamente aplicado nos círculos de cultura.","epistemico":"fato","exemplos":[],"id":"palavra_geradora","memoria":"semantica","meta":{"autor":"Freire","dominio":"ciencias_de_fora","fonte":""},"origem":"wiki","palavras_chave":[],"sinonimos":[],"tipo":"conceito","titulo":"Palavra Geradora"}
\section{Palavra Geradora}
Freire: método de alfabetização que parte de palavras carregadas de sentido existencial e político para o educando (de seu universo vocabular), gerando leitura crítica do mundo, não mera decodificação. Historicamente aplicado nos círculos de cultura.
\prov{relações: usa conscientizacao; relaciona capital\_cultural}
\prov{proveniência: wiki \textperiodcentered{} fato \textperiodcentered{} confiança 1.0 \textperiodcentered{} domínio ciencias\_de\_fora \textperiodcentered{} história 3 versões no jornal}

%CRISTAL {"arestas":[{"alvo":"computacao","confianca":1.0,"epistemico":"fato","origem":"wiki","rel":"referencia"},{"alvo":"gato","confianca":1.0,"epistemico":"fato","origem":"wiki","rel":"relaciona"},{"alvo":"word","confianca":0.5,"epistemico":"fato","origem":"wiki","rel":"relaciona"},{"alvo":"vontade_de_poder","confianca":0.5,"epistemico":"fato","origem":"wiki","rel":"relaciona"},{"alvo":"while","confianca":0.5,"epistemico":"fato","origem":"wiki","rel":"relaciona"}],"confianca":1.0,"contraexemplos":[],"descricao":"Wheeler: 'it from bit' — a realidade física emerge de informação; o universo é computação (Fredkin, Wolfram, Lloyd).","epistemico":"hipotese","exemplos":[],"id":"pancomputacionalismo","memoria":"semantica","meta":{"autor":"Wheeler","dominio":"filosofia","fonte":""},"origem":"wiki","palavras_chave":[],"sinonimos":[],"tipo":"conceito","titulo":"Pancomputacionalismo / It from Bit"}
\section{Pancomputacionalismo / It from Bit}
Wheeler: 'it from bit' — a realidade física emerge de informação; o universo é computação (Fredkin, Wolfram, Lloyd).
\prov{relações: referencia computacao; relaciona gato; relaciona word; relaciona vontade\_de\_poder; relaciona while}
\prov{proveniência: wiki \textperiodcentered{} hipotese \textperiodcentered{} confiança 1.0 \textperiodcentered{} domínio filosofia \textperiodcentered{} história 6 versões no jornal}

%CRISTAL {"arestas":[{"alvo":"hard_problem","confianca":1.0,"epistemico":"fato","origem":"wiki","rel":"explica"},{"alvo":"word","confianca":0.5,"epistemico":"fato","origem":"wiki","rel":"relaciona"},{"alvo":"virtualizacao","confianca":0.5,"epistemico":"fato","origem":"wiki","rel":"relaciona"},{"alvo":"semiotica_peirce_triade","confianca":0.5,"epistemico":"fato","origem":"wiki","rel":"relaciona"}],"confianca":1.0,"contraexemplos":[],"descricao":"A consciência é propriedade fundamental e ubíqua da matéria, não emergente — saída para o hard problem (Goff, Whitehead).","epistemico":"hipotese","exemplos":[],"id":"panpsiquismo","memoria":"semantica","meta":{"dominio":"filosofia","fonte":"web: SEP"},"origem":"wiki","palavras_chave":[],"sinonimos":[],"tipo":"conceito","titulo":"Panpsiquismo"}
\section{Panpsiquismo}
A consciência é propriedade fundamental e ubíqua da matéria, não emergente — saída para o hard problem (Goff, Whitehead).
\prov{relações: explica hard\_problem; relaciona word; relaciona virtualizacao; relaciona semiotica\_peirce\_triade}
\prov{proveniência: wiki \textperiodcentered{} web: SEP \textperiodcentered{} hipotese \textperiodcentered{} confiança 1.0 \textperiodcentered{} domínio filosofia \textperiodcentered{} história 5 versões no jornal}

%CRISTAL {"arestas":[{"alvo":"cristianismo","confianca":1.0,"epistemico":"fato","origem":"wiki","rel":"parte_de"},{"alvo":"simbolo_vs_significado","confianca":1.0,"epistemico":"fato","origem":"wiki","rel":"usa"},{"alvo":"logos_verbo","confianca":1.0,"epistemico":"fato","origem":"wiki","rel":"relaciona"}],"confianca":1.0,"contraexemplos":[],"descricao":"Bíblia: a forma narrativa breve que comunica verdades do Reino por analogia e imagem concreta, exigindo interpretação do ouvinte ('quem tem ouvidos para ouvir, ouça'). O sentido velado que se dá pela figura.","epistemico":"inferencia","exemplos":[],"id":"parabola","memoria":"semantica","meta":{"autor":"Bíblia","dominio":"sagrado","fonte":""},"origem":"wiki","palavras_chave":[],"sinonimos":[],"tipo":"conceito","titulo":"A Parábola — o Ensino por Imagem"}
\section{A Parábola — o Ensino por Imagem}
Bíblia: a forma narrativa breve que comunica verdades do Reino por analogia e imagem concreta, exigindo interpretação do ouvinte ('quem tem ouvidos para ouvir, ouça'). O sentido velado que se dá pela figura.
\prov{relações: parte\_de cristianismo; usa simbolo\_vs\_significado; relaciona logos\_verbo}
\prov{proveniência: wiki \textperiodcentered{} inferencia \textperiodcentered{} confiança 1.0 \textperiodcentered{} domínio sagrado \textperiodcentered{} história 4 versões no jornal}

%CRISTAL {"arestas":[{"alvo":"filosofia_da_tecnica","confianca":1.0,"epistemico":"fato","origem":"wiki","rel":"parte_de"},{"alvo":"aura_benjamin","confianca":1.0,"epistemico":"fato","origem":"wiki","rel":"relaciona"}],"confianca":1.0,"contraexemplos":[],"descricao":"Borgmann: os 'dispositivos' entregam a mercadoria (calor, comida) de modo cômodo mas opaco, escondendo seu funcionamento e dispensando engajamento — ao contrário das 'coisas focais', que exigem prática e presença.","epistemico":"inferencia","exemplos":[],"id":"paradigma_dispositivo","memoria":"semantica","meta":{"autor":"Borgmann","dominio":"ciencias_de_fora","fonte":""},"origem":"wiki","palavras_chave":[],"sinonimos":[],"tipo":"conceito","titulo":"Paradigma do Dispositivo (Borgmann)"}
\section{Paradigma do Dispositivo (Borgmann)}
Borgmann: os 'dispositivos' entregam a mercadoria (calor, comida) de modo cômodo mas opaco, escondendo seu funcionamento e dispensando engajamento — ao contrário das 'coisas focais', que exigem prática e presença.
\prov{relações: parte\_de filosofia\_da\_tecnica; relaciona aura\_benjamin}
\prov{proveniência: wiki \textperiodcentered{} inferencia \textperiodcentered{} confiança 1.0 \textperiodcentered{} domínio ciencias\_de\_fora \textperiodcentered{} história 3 versões no jornal}

%CRISTAL {"arestas":[{"alvo":"teoria_dos_conjuntos","confianca":1.0,"epistemico":"fato","origem":"wiki","rel":"parte_de"}],"confianca":1.0,"contraexemplos":[],"descricao":"O conjunto de todos os conjuntos que não se contêm: contém-se? A contradição que exigiu axiomatizar os conjuntos.","epistemico":"fato","exemplos":[],"id":"paradoxo_de_russell","memoria":"semantica","meta":{"autor":"Russell","dominio":"ciencias_de_fora","fonte":""},"origem":"wiki","palavras_chave":[],"sinonimos":[],"tipo":"conceito","titulo":"Paradoxo de Russell"}
\section{Paradoxo de Russell}
O conjunto de todos os conjuntos que não se contêm: contém-se? A contradição que exigiu axiomatizar os conjuntos.
\prov{relações: parte\_de teoria\_dos\_conjuntos}
\prov{proveniência: wiki \textperiodcentered{} fato \textperiodcentered{} confiança 1.0 \textperiodcentered{} domínio ciencias\_de\_fora \textperiodcentered{} história 4 versões no jornal}

%CRISTAL {"arestas":[{"alvo":"teoria_dos_numeros","confianca":1.0,"epistemico":"fato","origem":"wiki","rel":"parte_de"}],"confianca":1.0,"contraexemplos":[],"descricao":"Ramanujan: fórmulas quase-milagrosas para p(n) (modos de somar n) e as congruências ocultas — intuição pura, depois provada.","epistemico":"fato","exemplos":[],"id":"particoes_ramanujan","memoria":"semantica","meta":{"autor":"Ramanujan","dominio":"ciencias_de_fora","fonte":""},"origem":"wiki","palavras_chave":[],"sinonimos":[],"tipo":"conceito","titulo":"Partições e Ramanujan"}
\section{Partições e Ramanujan}
Ramanujan: fórmulas quase-milagrosas para p(n) (modos de somar n) e as congruências ocultas — intuição pura, depois provada.
\prov{relações: parte\_de teoria\_dos\_numeros}
\prov{proveniência: wiki \textperiodcentered{} fato \textperiodcentered{} confiança 1.0 \textperiodcentered{} domínio ciencias\_de\_fora \textperiodcentered{} história 4 versões no jornal}

%CRISTAL {"arestas":[{"alvo":"liberalismo_ri","confianca":1.0,"epistemico":"fato","origem":"wiki","rel":"usa"},{"alvo":"democracia_formas","confianca":1.0,"epistemico":"fato","origem":"wiki","rel":"usa"}],"confianca":1.0,"contraexemplos":[],"descricao":"A regularidade empírica de que democracias consolidadas praticamente não guerreiam entre si — uma das generalizações mais robustas das RI (Doyle), fundada na 'paz perpétua' de Kant (repúblicas constitucionais, federação de Estados livres). Contestada quanto a causas.","epistemico":"inferencia","exemplos":[],"id":"paz_democratica","memoria":"semantica","meta":{"autor":"Doyle/Kant","dominio":"ciencias_de_fora","fonte":""},"origem":"wiki","palavras_chave":[],"sinonimos":[],"tipo":"conceito","titulo":"Paz Democrática"}
\section{Paz Democrática}
A regularidade empírica de que democracias consolidadas praticamente não guerreiam entre si — uma das generalizações mais robustas das RI (Doyle), fundada na 'paz perpétua' de Kant (repúblicas constitucionais, federação de Estados livres). Contestada quanto a causas.
\prov{relações: usa liberalismo\_ri; usa democracia\_formas}
\prov{proveniência: wiki \textperiodcentered{} inferencia \textperiodcentered{} confiança 1.0 \textperiodcentered{} domínio ciencias\_de\_fora \textperiodcentered{} história 3 versões no jornal}

%CRISTAL {"arestas":[{"alvo":"educacao","confianca":1.0,"epistemico":"fato","origem":"wiki","rel":"parte_de"},{"alvo":"virtualizacao","confianca":0.5,"epistemico":"fato","origem":"wiki","rel":"relaciona"},{"alvo":"transcricao","confianca":0.5,"epistemico":"fato","origem":"wiki","rel":"relaciona"},{"alvo":"pow_metal_ressonancia_val","confianca":0.5,"epistemico":"fato","origem":"wiki","rel":"relaciona"},{"alvo":"ruido","confianca":0.5,"epistemico":"fato","origem":"wiki","rel":"relaciona"},{"alvo":"vida-morte","confianca":0.5,"epistemico":"fato","origem":"wiki","rel":"relaciona"},{"alvo":"utilitarismo","confianca":0.5,"epistemico":"fato","origem":"wiki","rel":"relaciona"},{"alvo":"scheduler","confianca":0.5,"epistemico":"fato","origem":"wiki","rel":"relaciona"},{"alvo":"pitagorismo","confianca":0.5,"epistemico":"fato","origem":"wiki","rel":"relaciona"},{"alvo":"tipos_de_interpretante","confianca":0.5,"epistemico":"fato","origem":"wiki","rel":"relaciona"},{"alvo":"resultado_anterior_identidade_nao_cobertura","confianca":0.5,"epistemico":"fato","origem":"wiki","rel":"relaciona"},{"alvo":"resultado_completo","confianca":0.5,"epistemico":"fato","origem":"wiki","rel":"relaciona"},{"alvo":"prova_da_maquina_espectral","confianca":0.5,"epistemico":"fato","origem":"wiki","rel":"relaciona"},{"alvo":"veredito_experimental","confianca":0.5,"epistemico":"fato","origem":"wiki","rel":"relaciona"},{"alvo":"setor_algebrico_descontinuidades_no_fecho","confianca":0.5,"epistemico":"fato","origem":"wiki","rel":"relaciona"},{"alvo":"teoria_do_valor","confianca":0.5,"epistemico":"fato","origem":"wiki","rel":"relaciona"},{"alvo":"vetores","confianca":0.5,"epistemico":"fato","origem":"wiki","rel":"relaciona"}],"confianca":1.0,"contraexemplos":[],"descricao":"A ciência e a arte do ensino — os métodos, fins e relações do ato de educar.","epistemico":"fato","exemplos":[],"id":"pedagogia","memoria":"semantica","meta":{"dominio":"ciencias_de_fora","fonte":""},"origem":"wiki","palavras_chave":[],"sinonimos":[],"tipo":"conceito","titulo":"Pedagogia"}
\section{Pedagogia}
A ciência e a arte do ensino — os métodos, fins e relações do ato de educar.
\prov{relações: parte\_de educacao; relaciona virtualizacao; relaciona transcricao; relaciona pow\_metal\_ressonancia\_val; relaciona ruido; relaciona vida-morte; relaciona utilitarismo; relaciona scheduler; relaciona pitagorismo; relaciona tipos\_de\_interpretante; relaciona resultado\_anterior\_identidade\_nao\_cobertura; relaciona resultado\_completo; relaciona prova\_da\_maquina\_espectral; relaciona veredito\_experimental; relaciona setor\_algebrico\_descontinuidades\_no\_fecho; relaciona teoria\_do\_valor; relaciona vetores}
\prov{proveniência: wiki \textperiodcentered{} fato \textperiodcentered{} confiança 1.0 \textperiodcentered{} domínio ciencias\_de\_fora \textperiodcentered{} história 18 versões no jornal}

%CRISTAL {"arestas":[{"alvo":"filosofia_da_educacao","confianca":1.0,"epistemico":"fato","origem":"wiki","rel":"parte_de"},{"alvo":"dialogicidade","confianca":1.0,"epistemico":"fato","origem":"wiki","rel":"usa"},{"alvo":"injustica_epistemica","confianca":1.0,"epistemico":"fato","origem":"wiki","rel":"relaciona"}],"confianca":1.0,"contraexemplos":[],"descricao":"bell hooks: em diálogo com Freire, a sala de aula como espaço de liberdade em que professor e alunos são presenças plenas (mente e corpo); ensinar é ato político que cruza raça, gênero e classe e transgride o modelo neutro e bancário.","epistemico":"inferencia","exemplos":[],"id":"pedagogia_engajada_hooks","memoria":"semantica","meta":{"autor":"bell hooks","dominio":"ciencias_de_fora","fonte":""},"origem":"wiki","palavras_chave":[],"sinonimos":[],"tipo":"conceito","titulo":"Pedagogia Engajada (bell hooks)"}
\section{Pedagogia Engajada (bell hooks)}
bell hooks: em diálogo com Freire, a sala de aula como espaço de liberdade em que professor e alunos são presenças plenas (mente e corpo); ensinar é ato político que cruza raça, gênero e classe e transgride o modelo neutro e bancário.
\prov{relações: parte\_de filosofia\_da\_educacao; usa dialogicidade; relaciona injustica\_epistemica}
\prov{proveniência: wiki \textperiodcentered{} inferencia \textperiodcentered{} confiança 1.0 \textperiodcentered{} domínio ciencias\_de\_fora \textperiodcentered{} história 4 versões no jornal}

%CRISTAL {"arestas":[{"alvo":"direito","confianca":1.0,"epistemico":"fato","origem":"wiki","rel":"parte_de"},{"alvo":"utilitarismo","confianca":1.0,"epistemico":"fato","origem":"wiki","rel":"usa"},{"alvo":"pena_retribuicao","confianca":1.0,"epistemico":"fato","origem":"wiki","rel":"contradiz"}],"confianca":1.0,"contraexemplos":[],"descricao":"Concepção utilitarista: a pena só se justifica por efeitos futuros — dissuadir o infrator (prevenção especial) e a sociedade (prevenção geral) —, minimizando o sofrimento total.","epistemico":"inferencia","exemplos":[],"id":"pena_dissuasao","memoria":"semantica","meta":{"autor":"Bentham","dominio":"ciencias_de_fora","fonte":""},"origem":"wiki","palavras_chave":[],"sinonimos":[],"tipo":"conceito","titulo":"Pena como Dissuasão (Bentham)"}
\section{Pena como Dissuasão (Bentham)}
Concepção utilitarista: a pena só se justifica por efeitos futuros — dissuadir o infrator (prevenção especial) e a sociedade (prevenção geral) —, minimizando o sofrimento total.
\prov{relações: parte\_de direito; usa utilitarismo; contradiz pena\_retribuicao}
\prov{proveniência: wiki \textperiodcentered{} inferencia \textperiodcentered{} confiança 1.0 \textperiodcentered{} domínio ciencias\_de\_fora \textperiodcentered{} história 4 versões no jornal}

%CRISTAL {"arestas":[{"alvo":"direito","confianca":1.0,"epistemico":"fato","origem":"wiki","rel":"parte_de"},{"alvo":"imperativo_categorico","confianca":1.0,"epistemico":"fato","origem":"wiki","rel":"usa"}],"confianca":1.0,"contraexemplos":[],"descricao":"A punição justifica-se porque o culpado a merece, em proporção à gravidade do crime (lex talionis, 'olho por olho'), como exigência de justiça e não por utilidade futura; para Kant, respeitar o criminoso como fim que respondeu por sua ação livre.","epistemico":"inferencia","exemplos":[],"id":"pena_retribuicao","memoria":"semantica","meta":{"autor":"Kant","dominio":"ciencias_de_fora","fonte":""},"origem":"wiki","palavras_chave":[],"sinonimos":[],"tipo":"conceito","titulo":"Pena como Retribuição (Kant)"}
\section{Pena como Retribuição (Kant)}
A punição justifica-se porque o culpado a merece, em proporção à gravidade do crime (lex talionis, 'olho por olho'), como exigência de justiça e não por utilidade futura; para Kant, respeitar o criminoso como fim que respondeu por sua ação livre.
\prov{relações: parte\_de direito; usa imperativo\_categorico}
\prov{proveniência: wiki \textperiodcentered{} inferencia \textperiodcentered{} confiança 1.0 \textperiodcentered{} domínio ciencias\_de\_fora \textperiodcentered{} história 3 versões no jornal}

%CRISTAL {"arestas":[{"alvo":"deus_como_vazio_gentil","confianca":1.0,"epistemico":"fato","origem":"wiki","rel":"usa"},{"alvo":"sunyata","confianca":1.0,"epistemico":"fato","origem":"wiki","rel":"relaciona"},{"alvo":"teologia_quantica_gentil","confianca":1.0,"epistemico":"fato","origem":"wiki","rel":"parte_de"},{"alvo":"agape","confianca":1.0,"epistemico":"fato","origem":"wiki","rel":"contradiz"}],"confianca":1.0,"contraexemplos":[],"descricao":"Lopes: o conjunto vazio é a única coisa perfeita — 'para deixar de ser perfeito teria de conter uma imperfeição, mas nada contém' —, logo Deus=Vazio é perfeito. Corolário radical: atributos positivos (amor, misericórdia) seriam defeitos, pois exigiriam criar um miserável/um dual — 'um Deus pra homem nenhum botar defeito'.","epistemico":"inferencia","exemplos":[],"id":"perfeicao_do_vazio_gentil","memoria":"semantica","meta":{"autor":"Gentil Lopes","dominio":"sagrado","fonte":""},"origem":"wiki","palavras_chave":[],"sinonimos":[],"tipo":"conceito","titulo":"A Perfeição do Vazio (Lopes)"}
\section{A Perfeição do Vazio (Lopes)}
Lopes: o conjunto vazio é a única coisa perfeita — 'para deixar de ser perfeito teria de conter uma imperfeição, mas nada contém' —, logo Deus=Vazio é perfeito. Corolário radical: atributos positivos (amor, misericórdia) seriam defeitos, pois exigiriam criar um miserável/um dual — 'um Deus pra homem nenhum botar defeito'.
\prov{relações: usa deus\_como\_vazio\_gentil; relaciona sunyata; parte\_de teologia\_quantica\_gentil; contradiz agape}
\prov{proveniência: wiki \textperiodcentered{} inferencia \textperiodcentered{} confiança 1.0 \textperiodcentered{} domínio sagrado \textperiodcentered{} história 6 versões no jornal}

%CRISTAL {"arestas":[{"alvo":"artes_visuais","confianca":1.0,"epistemico":"fato","origem":"wiki","rel":"parte_de"},{"alvo":"geometria_ramo","confianca":1.0,"epistemico":"fato","origem":"wiki","rel":"usa"},{"alvo":"perspectiva_linear","confianca":1.0,"epistemico":"fato","origem":"wiki","rel":"relaciona"}],"confianca":1.0,"contraexemplos":[],"descricao":"Brunelleschi/Alberti: projetar o espaço 3D no plano por um ponto de fuga — geometria projetiva virou pintura.","epistemico":"fato","exemplos":[],"id":"perspectiva","memoria":"semantica","meta":{"autor":"Brunelleschi","dominio":"ciencias_de_fora","fonte":""},"origem":"wiki","palavras_chave":[],"sinonimos":[],"tipo":"conceito","titulo":"Perspectiva"}
\section{Perspectiva}
Brunelleschi/Alberti: projetar o espaço 3D no plano por um ponto de fuga — geometria projetiva virou pintura.
\prov{relações: parte\_de artes\_visuais; usa geometria\_ramo; relaciona perspectiva\_linear}
\prov{proveniência: wiki \textperiodcentered{} fato \textperiodcentered{} confiança 1.0 \textperiodcentered{} domínio ciencias\_de\_fora \textperiodcentered{} história 5 versões no jornal}

%CRISTAL {"arestas":[{"alvo":"artes_visuais","confianca":1.0,"epistemico":"fato","origem":"wiki","rel":"parte_de"},{"alvo":"mimese_expressao","confianca":1.0,"epistemico":"fato","origem":"wiki","rel":"usa"},{"alvo":"razao_aurea","confianca":1.0,"epistemico":"fato","origem":"wiki","rel":"relaciona"}],"confianca":1.0,"contraexemplos":[],"descricao":"Brunelleschi/Alberti: o sistema geométrico para representar o espaço tridimensional numa superfície plana através de linha do horizonte, ponto de fuga e ortogonais convergentes — a 'janela de Alberti'. A matematização da representação.","epistemico":"fato","exemplos":[],"id":"perspectiva_linear","memoria":"semantica","meta":{"autor":"Brunelleschi/Alberti","dominio":"ciencias_de_fora","fonte":""},"origem":"wiki","palavras_chave":[],"sinonimos":[],"tipo":"conceito","titulo":"Perspectiva Linear (Brunelleschi/Alberti)"}
\section{Perspectiva Linear (Brunelleschi/Alberti)}
Brunelleschi/Alberti: o sistema geométrico para representar o espaço tridimensional numa superfície plana através de linha do horizonte, ponto de fuga e ortogonais convergentes — a 'janela de Alberti'. A matematização da representação.
\prov{relações: parte\_de artes\_visuais; usa mimese\_expressao; relaciona razao\_aurea}
\prov{proveniência: wiki \textperiodcentered{} fato \textperiodcentered{} confiança 1.0 \textperiodcentered{} domínio ciencias\_de\_fora \textperiodcentered{} história 4 versões no jornal}

%CRISTAL {"arestas":[{"alvo":"musica","confianca":1.0,"epistemico":"fato","origem":"wiki","rel":"parte_de"},{"alvo":"beleza_matematica","confianca":1.0,"epistemico":"fato","origem":"wiki","rel":"relaciona"},{"alvo":"razao_aurea","confianca":1.0,"epistemico":"fato","origem":"wiki","rel":"relaciona"}],"confianca":1.0,"contraexemplos":[],"descricao":"Pitágoras (via o monocórdio): os intervalos consonantes correspondem a razões de números inteiros pequenos (oitava 2:1, quinta 3:2, quarta 4:3) — a primeira ligação empírica entre som e número. A série harmônica confirma-o na física: um corpo vibrante produz parciais em múltiplos inteiros da fundamental.","epistemico":"fato","exemplos":[],"id":"pitagoras_razoes_musicais","memoria":"semantica","meta":{"autor":"Pitágoras","dominio":"ciencias_de_fora","fonte":""},"origem":"wiki","palavras_chave":[],"sinonimos":[],"tipo":"conceito","titulo":"Pitágoras e as Razões Musicais"}
\section{Pitágoras e as Razões Musicais}
Pitágoras (via o monocórdio): os intervalos consonantes correspondem a razões de números inteiros pequenos (oitava 2:1, quinta 3:2, quarta 4:3) — a primeira ligação empírica entre som e número. A série harmônica confirma-o na física: um corpo vibrante produz parciais em múltiplos inteiros da fundamental.
\prov{relações: parte\_de musica; relaciona beleza\_matematica; relaciona razao\_aurea}
\prov{proveniência: wiki \textperiodcentered{} fato \textperiodcentered{} confiança 1.0 \textperiodcentered{} domínio ciencias\_de\_fora \textperiodcentered{} história 4 versões no jornal}

%CRISTAL {"arestas":[{"alvo":"platonismo_matematico","confianca":1.0,"epistemico":"fato","origem":"wiki","rel":"relaciona"},{"alvo":"vida-morte","confianca":0.5,"epistemico":"fato","origem":"wiki","rel":"relaciona"},{"alvo":"utilitarismo","confianca":0.5,"epistemico":"fato","origem":"wiki","rel":"relaciona"},{"alvo":"scheduler","confianca":0.5,"epistemico":"fato","origem":"wiki","rel":"relaciona"},{"alvo":"ruido","confianca":0.5,"epistemico":"fato","origem":"wiki","rel":"relaciona"},{"alvo":"tipos_de_interpretante","confianca":0.5,"epistemico":"fato","origem":"wiki","rel":"relaciona"},{"alvo":"transcricao","confianca":0.5,"epistemico":"fato","origem":"wiki","rel":"relaciona"},{"alvo":"veredito_experimental","confianca":0.5,"epistemico":"fato","origem":"wiki","rel":"relaciona"},{"alvo":"resultado_anterior_identidade_nao_cobertura","confianca":0.5,"epistemico":"fato","origem":"wiki","rel":"relaciona"},{"alvo":"teoria_do_valor","confianca":0.5,"epistemico":"fato","origem":"wiki","rel":"relaciona"},{"alvo":"virtualizacao","confianca":0.5,"epistemico":"fato","origem":"wiki","rel":"relaciona"},{"alvo":"resultado_completo","confianca":0.5,"epistemico":"fato","origem":"wiki","rel":"relaciona"},{"alvo":"pow_metal_ressonancia_val","confianca":0.5,"epistemico":"fato","origem":"wiki","rel":"relaciona"}],"confianca":1.0,"contraexemplos":[],"descricao":"'Tudo é número': as relações numéricas são a essência do cosmos.","epistemico":"inferencia","exemplos":[],"id":"pitagorismo","memoria":"semantica","meta":{"autor":"Pitágoras","dominio":"filosofia","fonte":""},"origem":"wiki","palavras_chave":[],"sinonimos":[],"tipo":"conceito","titulo":"Pitagorismo"}
\section{Pitagorismo}
'Tudo é número': as relações numéricas são a essência do cosmos.
\prov{relações: relaciona platonismo\_matematico; relaciona vida-morte; relaciona utilitarismo; relaciona scheduler; relaciona ruido; relaciona tipos\_de\_interpretante; relaciona transcricao; relaciona veredito\_experimental; relaciona resultado\_anterior\_identidade\_nao\_cobertura; relaciona teoria\_do\_valor; relaciona virtualizacao; relaciona resultado\_completo; relaciona pow\_metal\_ressonancia\_val}
\prov{proveniência: wiki \textperiodcentered{} inferencia \textperiodcentered{} confiança 1.0 \textperiodcentered{} domínio filosofia \textperiodcentered{} história 14 versões no jornal}

%CRISTAL {"arestas":[{"alvo":"sinapse","confianca":1.0,"epistemico":"fato","origem":"wiki","rel":"usa"},{"alvo":"hopfield","confianca":1.0,"epistemico":"fato","origem":"wiki","rel":"explica"},{"alvo":"consciencia","confianca":1.0,"epistemico":"fato","origem":"wiki","rel":"relaciona"}],"confianca":1.0,"contraexemplos":[],"descricao":"Hebb: 'neurônios que disparam juntos, conectam-se' — a regra local que grava padrões. É a lei da rede de Hopfield.","epistemico":"fato","exemplos":[],"id":"plasticidade_hebbiana","memoria":"semantica","meta":{"autor":"Hebb","dominio":"ciencias_de_fora","fonte":""},"origem":"wiki","palavras_chave":[],"sinonimos":[],"tipo":"conceito","titulo":"Plasticidade Hebbiana"}
\section{Plasticidade Hebbiana}
Hebb: 'neurônios que disparam juntos, conectam-se' — a regra local que grava padrões. É a lei da rede de Hopfield.
\prov{relações: usa sinapse; explica hopfield; relaciona consciencia}
\prov{proveniência: wiki \textperiodcentered{} fato \textperiodcentered{} confiança 1.0 \textperiodcentered{} domínio ciencias\_de\_fora \textperiodcentered{} história 4 versões no jornal}

%CRISTAL {"arestas":[{"alvo":"comunicacao_digital","confianca":1.0,"epistemico":"fato","origem":"wiki","rel":"parte_de"},{"alvo":"capitalismo_de_vigilancia","confianca":1.0,"epistemico":"fato","origem":"wiki","rel":"relaciona"},{"alvo":"sociedade_em_rede","confianca":1.0,"epistemico":"fato","origem":"wiki","rel":"usa"}],"confianca":1.0,"contraexemplos":[],"descricao":"Poell, Nieborg & van Dijck: a penetração das infraestruturas, processos econômicos e governança das plataformas digitais em todos os setores da vida, concentrando poder nas grandes corporações-plataforma.","epistemico":"inferencia","exemplos":[],"id":"plataformizacao","memoria":"semantica","meta":{"autor":"Poell/Nieborg/van Dijck","dominio":"ciencias_de_fora","fonte":""},"origem":"wiki","palavras_chave":[],"sinonimos":[],"tipo":"conceito","titulo":"Plataformização"}
\section{Plataformização}
Poell, Nieborg \& van Dijck: a penetração das infraestruturas, processos econômicos e governança das plataformas digitais em todos os setores da vida, concentrando poder nas grandes corporações-plataforma.
\prov{relações: parte\_de comunicacao\_digital; relaciona capitalismo\_de\_vigilancia; usa sociedade\_em\_rede}
\prov{proveniência: wiki \textperiodcentered{} inferencia \textperiodcentered{} confiança 1.0 \textperiodcentered{} domínio ciencias\_de\_fora \textperiodcentered{} história 4 versões no jornal}

%CRISTAL {"arestas":[{"alvo":"filosofia_da_matematica","confianca":1.0,"epistemico":"fato","origem":"wiki","rel":"parte_de"},{"alvo":"numeros_azuis_vermelhos","confianca":1.0,"epistemico":"fato","origem":"wiki","rel":"referencia"},{"alvo":"word","confianca":0.5,"epistemico":"fato","origem":"wiki","rel":"relaciona"}],"confianca":1.0,"contraexemplos":[],"descricao":"Os objetos matemáticos existem eternamente, independentes da mente; a matemática é descoberta. Gödel e Penrose defendem.","epistemico":"inferencia","exemplos":[],"id":"platonismo_matematico","memoria":"semantica","meta":{"dominio":"filosofia","fonte":"web: SEP/Wikipedia"},"origem":"wiki","palavras_chave":[],"sinonimos":[],"tipo":"conceito","titulo":"Platonismo Matemático"}
\section{Platonismo Matemático}
Os objetos matemáticos existem eternamente, independentes da mente; a matemática é descoberta. Gödel e Penrose defendem.
\prov{relações: parte\_de filosofia\_da\_matematica; referencia numeros\_azuis\_vermelhos; relaciona word}
\prov{proveniência: wiki \textperiodcentered{} web: SEP/Wikipedia \textperiodcentered{} inferencia \textperiodcentered{} confiança 1.0 \textperiodcentered{} domínio filosofia \textperiodcentered{} história 4 versões no jornal}

%CRISTAL {"arestas":[{"alvo":"ciencia_politica","confianca":1.0,"epistemico":"fato","origem":"wiki","rel":"parte_de"}],"confianca":1.0,"contraexemplos":[],"descricao":"Dahl (Who Governs?): contra a tese elitista de um poder concentrado, o poder seria disperso entre grupos concorrentes que vencem em áreas distintas (a poliarquia), sem que uma única elite domine todas as decisões.","epistemico":"inferencia","exemplos":[],"id":"pluralismo_elitismo","memoria":"semantica","meta":{"autor":"Dahl","dominio":"ciencias_de_fora","fonte":""},"origem":"wiki","palavras_chave":[],"sinonimos":[],"tipo":"conceito","titulo":"Quem Realmente Governa? (Dahl)"}
\section{Quem Realmente Governa? (Dahl)}
Dahl (Who Governs?): contra a tese elitista de um poder concentrado, o poder seria disperso entre grupos concorrentes que vencem em áreas distintas (a poliarquia), sem que uma única elite domine todas as decisões.
\prov{relações: parte\_de ciencia\_politica}
\prov{proveniência: wiki \textperiodcentered{} inferencia \textperiodcentered{} confiança 1.0 \textperiodcentered{} domínio ciencias\_de\_fora \textperiodcentered{} história 2 versões no jornal}

%CRISTAL {"arestas":[{"alvo":"ciencia_politica","confianca":1.0,"epistemico":"fato","origem":"wiki","rel":"parte_de"},{"alvo":"democracia_deliberativa","confianca":1.0,"epistemico":"fato","origem":"wiki","rel":"relaciona"},{"alvo":"tres_faces_do_poder","confianca":1.0,"epistemico":"fato","origem":"wiki","rel":"contradiz"}],"confianca":1.0,"contraexemplos":[],"descricao":"Arendt: o poder não é posse nem instrumento, mas a capacidade humana de agir em concerto no espaço público — distinto da violência, que é meramente instrumental e surge quando o poder colapsa.","epistemico":"inferencia","exemplos":[],"id":"poder_agir_em_concerto","memoria":"semantica","meta":{"autor":"Arendt","dominio":"ciencias_de_fora","fonte":""},"origem":"wiki","palavras_chave":[],"sinonimos":[],"tipo":"conceito","titulo":"Poder como Agir em Concerto (Arendt)"}
\section{Poder como Agir em Concerto (Arendt)}
Arendt: o poder não é posse nem instrumento, mas a capacidade humana de agir em concerto no espaço público — distinto da violência, que é meramente instrumental e surge quando o poder colapsa.
\prov{relações: parte\_de ciencia\_politica; relaciona democracia\_deliberativa; contradiz tres\_faces\_do\_poder}
\prov{proveniência: wiki \textperiodcentered{} inferencia \textperiodcentered{} confiança 1.0 \textperiodcentered{} domínio ciencias\_de\_fora \textperiodcentered{} história 4 versões no jornal}

%CRISTAL {"arestas":[{"alvo":"politica","confianca":1.0,"epistemico":"fato","origem":"wiki","rel":"parte_de"},{"alvo":"virtualizacao","confianca":0.5,"epistemico":"fato","origem":"wiki","rel":"relaciona"},{"alvo":"vida-morte","confianca":0.5,"epistemico":"fato","origem":"wiki","rel":"relaciona"},{"alvo":"vontade_de_poder","confianca":0.5,"epistemico":"fato","origem":"wiki","rel":"relaciona"},{"alvo":"word","confianca":0.5,"epistemico":"fato","origem":"wiki","rel":"relaciona"},{"alvo":"proposition_gerador_de_escala_ipropderiv_com_tddt","confianca":0.5,"epistemico":"fato","origem":"wiki","rel":"relaciona"}],"confianca":1.0,"contraexemplos":[],"descricao":"Foucault: o poder que regula corpos por vigilância e normalização (o panóptico), não por violência direta.","epistemico":"inferencia","exemplos":[],"id":"poder_disciplinar_foucault","memoria":"semantica","meta":{"autor":"Foucault","dominio":"ciencias_de_fora","fonte":""},"origem":"wiki","palavras_chave":[],"sinonimos":[],"tipo":"conceito","titulo":"Poder Disciplinar"}
\section{Poder Disciplinar}
Foucault: o poder que regula corpos por vigilância e normalização (o panóptico), não por violência direta.
\prov{relações: parte\_de politica; relaciona virtualizacao; relaciona vida-morte; relaciona vontade\_de\_poder; relaciona word; relaciona proposition\_gerador\_de\_escala\_ipropderiv\_com\_tddt}
\prov{proveniência: wiki \textperiodcentered{} inferencia \textperiodcentered{} confiança 1.0 \textperiodcentered{} domínio ciencias\_de\_fora \textperiodcentered{} história 7 versões no jornal}

%CRISTAL {"arestas":[{"alvo":"mapa_nao_territorio","confianca":1.0,"epistemico":"fato","origem":"wiki","rel":"usa"},{"alvo":"artefatos_tem_politica","confianca":1.0,"epistemico":"fato","origem":"wiki","rel":"usa"},{"alvo":"biopoder","confianca":1.0,"epistemico":"fato","origem":"wiki","rel":"relaciona"}],"confianca":1.0,"contraexemplos":[],"descricao":"Harley: o mapa não é conhecimento neutro, mas um texto carregado de poder — centraliza, hierarquiza e silencia, servindo a interesses. Toda projeção distorce e escolhe o que preservar (a de Mercator infla o Norte e diminui o Sul global): a suposta neutralidade deve ser sempre questionada.","epistemico":"inferencia","exemplos":[],"id":"poder_do_mapa_harley","memoria":"semantica","meta":{"autor":"Harley","dominio":"ciencias_de_fora","fonte":""},"origem":"wiki","palavras_chave":[],"sinonimos":[],"tipo":"conceito","titulo":"O Poder do Mapa (Harley)"}
\section{O Poder do Mapa (Harley)}
Harley: o mapa não é conhecimento neutro, mas um texto carregado de poder — centraliza, hierarquiza e silencia, servindo a interesses. Toda projeção distorce e escolhe o que preservar (a de Mercator infla o Norte e diminui o Sul global): a suposta neutralidade deve ser sempre questionada.
\prov{relações: usa mapa\_nao\_territorio; usa artefatos\_tem\_politica; relaciona biopoder}
\prov{proveniência: wiki \textperiodcentered{} inferencia \textperiodcentered{} confiança 1.0 \textperiodcentered{} domínio ciencias\_de\_fora \textperiodcentered{} história 4 versões no jornal}

%CRISTAL {"arestas":[{"alvo":"geopolitica","confianca":1.0,"epistemico":"fato","origem":"wiki","rel":"usa"},{"alvo":"heartland_mackinder","confianca":1.0,"epistemico":"fato","origem":"wiki","rel":"contradiz"}],"confianca":1.0,"contraexemplos":[],"descricao":"Mahan: a supremacia naval — grande frota, marinha mercante e rede de bases — é a chave do sucesso político e comercial de uma nação, e portanto da potência mundial. O contraponto marítimo ao Heartland.","epistemico":"inferencia","exemplos":[],"id":"poder_maritimo_mahan","memoria":"semantica","meta":{"autor":"Mahan","dominio":"ciencias_de_fora","fonte":""},"origem":"wiki","palavras_chave":[],"sinonimos":[],"tipo":"conceito","titulo":"Poder Marítimo (Mahan)"}
\section{Poder Marítimo (Mahan)}
Mahan: a supremacia naval — grande frota, marinha mercante e rede de bases — é a chave do sucesso político e comercial de uma nação, e portanto da potência mundial. O contraponto marítimo ao Heartland.
\prov{relações: usa geopolitica; contradiz heartland\_mackinder}
\prov{proveniência: wiki \textperiodcentered{} inferencia \textperiodcentered{} confiança 1.0 \textperiodcentered{} domínio ciencias\_de\_fora \textperiodcentered{} história 3 versões no jornal}

%CRISTAL {"arestas":[{"alvo":"artes","confianca":1.0,"epistemico":"fato","origem":"wiki","rel":"relaciona"},{"alvo":"comunicacao_iconoclasta_gentil","confianca":1.0,"epistemico":"fato","origem":"wiki","rel":"usa"},{"alvo":"beleza_matematica","confianca":1.0,"epistemico":"fato","origem":"wiki","rel":"relaciona"}],"confianca":1.0,"contraexemplos":[],"descricao":"Lopes: no diálogo Einstein-Tagore, o poeta 'ganhou de lavada'; gênios como Hardy e Einstein 'se revelaram rasos' em suas posições filosóficas. Eleva a intuição poética/artística acima da científica nas questões de fundo sobre a natureza da realidade.","epistemico":"inferencia","exemplos":[],"id":"poeta_supera_cientista_gentil","memoria":"semantica","meta":{"autor":"Gentil Lopes","dominio":"ciencias_de_fora","fonte":""},"origem":"wiki","palavras_chave":[],"sinonimos":[],"tipo":"conceito","titulo":"O Poeta Vence o Cientista (Lopes)"}
\section{O Poeta Vence o Cientista (Lopes)}
Lopes: no diálogo Einstein-Tagore, o poeta 'ganhou de lavada'; gênios como Hardy e Einstein 'se revelaram rasos' em suas posições filosóficas. Eleva a intuição poética/artística acima da científica nas questões de fundo sobre a natureza da realidade.
\prov{relações: relaciona artes; usa comunicacao\_iconoclasta\_gentil; relaciona beleza\_matematica}
\prov{proveniência: wiki \textperiodcentered{} inferencia \textperiodcentered{} confiança 1.0 \textperiodcentered{} domínio ciencias\_de\_fora \textperiodcentered{} história 5 versões no jornal}

%CRISTAL {"arestas":[{"alvo":"teoria_literaria","confianca":1.0,"epistemico":"fato","origem":"wiki","rel":"parte_de"},{"alvo":"mimese","confianca":1.0,"epistemico":"fato","origem":"wiki","rel":"usa"},{"alvo":"catarse","confianca":1.0,"epistemico":"fato","origem":"wiki","rel":"explica"}],"confianca":1.0,"contraexemplos":[],"descricao":"O primeiro tratado de teoria literária: a tragédia como imitação de uma ação, com enredo, caráter e desfecho.","epistemico":"inferencia","exemplos":[],"id":"poetica_aristoteles","memoria":"semantica","meta":{"autor":"Aristóteles","dominio":"ciencias_de_fora","fonte":""},"origem":"wiki","palavras_chave":[],"sinonimos":[],"tipo":"conceito","titulo":"Poética (Aristóteles)"}
\section{Poética (Aristóteles)}
O primeiro tratado de teoria literária: a tragédia como imitação de uma ação, com enredo, caráter e desfecho.
\prov{relações: parte\_de teoria\_literaria; usa mimese; explica catarse}
\prov{proveniência: wiki \textperiodcentered{} inferencia \textperiodcentered{} confiança 1.0 \textperiodcentered{} domínio ciencias\_de\_fora \textperiodcentered{} história 4 versões no jornal}

%CRISTAL {"arestas":[{"alvo":"topologia","confianca":1.0,"epistemico":"fato","origem":"wiki","rel":"parte_de"},{"alvo":"conjectura_de_poincare","confianca":1.0,"epistemico":"fato","origem":"wiki","rel":"relaciona"},{"alvo":"ordem_borda_caos","confianca":1.0,"epistemico":"fato","origem":"wiki","rel":"relaciona"}],"confianca":1.0,"contraexemplos":[],"descricao":"Poincaré: topologia, sistemas dinâmicos, o caos (o problema dos três corpos), fundamentos — dominou toda a matemática de seu tempo.","epistemico":"fato","exemplos":[],"id":"poincare_universalista","memoria":"semantica","meta":{"autor":"Poincaré","dominio":"ciencias_de_fora","fonte":""},"origem":"wiki","palavras_chave":[],"sinonimos":[],"tipo":"conceito","titulo":"Poincaré, o Último Universalista"}
\section{Poincaré, o Último Universalista}
Poincaré: topologia, sistemas dinâmicos, o caos (o problema dos três corpos), fundamentos — dominou toda a matemática de seu tempo.
\prov{relações: parte\_de topologia; relaciona conjectura\_de\_poincare; relaciona ordem\_borda\_caos}
\prov{proveniência: wiki \textperiodcentered{} fato \textperiodcentered{} confiança 1.0 \textperiodcentered{} domínio ciencias\_de\_fora \textperiodcentered{} história 7 versões no jornal}

%CRISTAL {"arestas":[{"alvo":"democracia_formas","confianca":1.0,"epistemico":"fato","origem":"wiki","rel":"usa"},{"alvo":"estado_de_direito","confianca":1.0,"epistemico":"fato","origem":"wiki","rel":"usa"},{"alvo":"liberdade_berlin","confianca":1.0,"epistemico":"fato","origem":"wiki","rel":"relaciona"},{"alvo":"pluralismo_elitismo","confianca":1.0,"epistemico":"fato","origem":"wiki","rel":"relaciona"}],"confianca":1.0,"contraexemplos":[],"descricao":"Dahl: as condições reais e verificáveis da democracia moderna, em duas dimensões — contestação pública (competição, oposição, informação livre) e participação (inclusão dos cidadãos). A aproximação prática possível do ideal democrático.","epistemico":"inferencia","exemplos":[],"id":"poliarquia","memoria":"semantica","meta":{"autor":"Dahl","dominio":"ciencias_de_fora","fonte":""},"origem":"wiki","palavras_chave":[],"sinonimos":[],"tipo":"conceito","titulo":"Poliarquia (Dahl)"}
\section{Poliarquia (Dahl)}
Dahl: as condições reais e verificáveis da democracia moderna, em duas dimensões — contestação pública (competição, oposição, informação livre) e participação (inclusão dos cidadãos). A aproximação prática possível do ideal democrático.
\prov{relações: usa democracia\_formas; usa estado\_de\_direito; relaciona liberdade\_berlin; relaciona pluralismo\_elitismo}
\prov{proveniência: wiki \textperiodcentered{} inferencia \textperiodcentered{} confiança 1.0 \textperiodcentered{} domínio ciencias\_de\_fora \textperiodcentered{} história 5 versões no jornal}

%CRISTAL {"arestas":[{"alvo":"filosofia","confianca":1.0,"epistemico":"fato","origem":"wiki","rel":"parte_de"},{"alvo":"utilitarismo","confianca":0.5,"epistemico":"fato","origem":"wiki","rel":"relaciona"},{"alvo":"teoria_dos_numeros","confianca":0.5,"epistemico":"fato","origem":"wiki","rel":"relaciona"},{"alvo":"word","confianca":0.5,"epistemico":"fato","origem":"wiki","rel":"relaciona"},{"alvo":"syscall","confianca":0.5,"epistemico":"fato","origem":"wiki","rel":"relaciona"},{"alvo":"virtualizacao","confianca":0.5,"epistemico":"fato","origem":"wiki","rel":"relaciona"}],"confianca":1.0,"contraexemplos":[],"descricao":"O estudo do poder legítimo, da justiça e da ordem social — contrato, liberdade, Estado.","epistemico":"fato","exemplos":[],"id":"politica","memoria":"semantica","meta":{"dominio":"ciencias_de_fora","fonte":""},"origem":"wiki","palavras_chave":[],"sinonimos":[],"tipo":"conceito","titulo":"Filosofia Política"}
\section{Filosofia Política}
O estudo do poder legítimo, da justiça e da ordem social — contrato, liberdade, Estado.
\prov{relações: parte\_de filosofia; relaciona utilitarismo; relaciona teoria\_dos\_numeros; relaciona word; relaciona syscall; relaciona virtualizacao}
\prov{proveniência: wiki \textperiodcentered{} fato \textperiodcentered{} confiança 1.0 \textperiodcentered{} domínio ciencias\_de\_fora \textperiodcentered{} história 7 versões no jornal}

%CRISTAL {"arestas":[{"alvo":"didatica_matematica","confianca":1.0,"epistemico":"fato","origem":"wiki","rel":"parte_de"},{"alvo":"regua_de_gentil","confianca":1.0,"epistemico":"fato","origem":"wiki","rel":"referencia"}],"confianca":1.0,"contraexemplos":[],"descricao":"Pólya (How to Solve It): quatro fases para resolver problemas — compreender, conceber um plano, executar e retrospectivar. Heurísticas não garantem a solução, mas elevam a chance de sucesso via perguntas orientadoras.","epistemico":"inferencia","exemplos":[],"id":"polya_heuristica","memoria":"semantica","meta":{"autor":"Pólya","dominio":"ciencias_de_fora","fonte":""},"origem":"wiki","palavras_chave":[],"sinonimos":[],"tipo":"conceito","titulo":"Heurística de Pólya"}
\section{Heurística de Pólya}
Pólya (How to Solve It): quatro fases para resolver problemas — compreender, conceber um plano, executar e retrospectivar. Heurísticas não garantem a solução, mas elevam a chance de sucesso via perguntas orientadoras.
\prov{relações: parte\_de didatica\_matematica; referencia regua\_de\_gentil}
\prov{proveniência: wiki \textperiodcentered{} inferencia \textperiodcentered{} confiança 1.0 \textperiodcentered{} domínio ciencias\_de\_fora \textperiodcentered{} história 3 versões no jornal}

%CRISTAL {"arestas":[{"alvo":"ciencia_politica","confianca":1.0,"epistemico":"fato","origem":"wiki","rel":"parte_de"},{"alvo":"hegemonia_cultural","confianca":1.0,"epistemico":"fato","origem":"wiki","rel":"usa"},{"alvo":"espiral_do_silencio","confianca":1.0,"epistemico":"fato","origem":"wiki","rel":"relaciona"},{"alvo":"democracia_formas","confianca":1.0,"epistemico":"fato","origem":"wiki","rel":"contradiz"}],"confianca":1.0,"contraexemplos":[],"descricao":"Uma ideologia de 'núcleo fino' (Mudde) que divide a sociedade entre 'o povo puro' e 'a elite corrupta' e exige que a política expresse a vontade geral; por ser fina, acopla-se a hospedeiras (nacionalismo, socialismo). Para Laclau, é a própria lógica de construção do 'povo' como significante vazio.","epistemico":"inferencia","exemplos":[],"id":"populismo","memoria":"semantica","meta":{"autor":"Mudde/Laclau","dominio":"ciencias_de_fora","fonte":""},"origem":"wiki","palavras_chave":[],"sinonimos":[],"tipo":"conceito","titulo":"Populismo"}
\section{Populismo}
Uma ideologia de 'núcleo fino' (Mudde) que divide a sociedade entre 'o povo puro' e 'a elite corrupta' e exige que a política expresse a vontade geral; por ser fina, acopla-se a hospedeiras (nacionalismo, socialismo). Para Laclau, é a própria lógica de construção do 'povo' como significante vazio.
\prov{relações: parte\_de ciencia\_politica; usa hegemonia\_cultural; relaciona espiral\_do\_silencio; contradiz democracia\_formas}
\prov{proveniência: wiki \textperiodcentered{} inferencia \textperiodcentered{} confiança 1.0 \textperiodcentered{} domínio ciencias\_de\_fora \textperiodcentered{} história 5 versões no jornal}

%CRISTAL {"arestas":[{"alvo":"linguas_romanicas","confianca":1.0,"epistemico":"fato","origem":"wiki","rel":"especializa"},{"alvo":"tiffany","confianca":1.0,"epistemico":"fato","origem":"wiki","rel":"explica"},{"alvo":"tupi","confianca":1.0,"epistemico":"fato","origem":"wiki","rel":"usa"}],"confianca":1.0,"contraexemplos":[],"descricao":"Língua românica de ~260 milhões de falantes (Brasil, Portugal, África) — e o idioma em que a Tiffany fala, pensa e responde.","epistemico":"fato","exemplos":[],"id":"portugues","memoria":"semantica","meta":{"dominio":"ciencias_de_fora","fonte":""},"origem":"wiki","palavras_chave":[],"sinonimos":[],"tipo":"conceito","titulo":"Português"}
\section{Português}
Língua românica de \~{}260 milhões de falantes (Brasil, Portugal, África) — e o idioma em que a Tiffany fala, pensa e responde.
\prov{relações: especializa linguas\_romanicas; explica tiffany; usa tupi}
\prov{proveniência: wiki \textperiodcentered{} fato \textperiodcentered{} confiança 1.0 \textperiodcentered{} domínio ciencias\_de\_fora \textperiodcentered{} história 4 versões no jornal}

%CRISTAL {"arestas":[{"alvo":"comunicacao_digital","confianca":1.0,"epistemico":"fato","origem":"wiki","rel":"parte_de"},{"alvo":"deus_mentiroso","confianca":1.0,"epistemico":"fato","origem":"wiki","rel":"relaciona"},{"alvo":"dois_sistemas_kahneman","confianca":1.0,"epistemico":"fato","origem":"wiki","rel":"relaciona"},{"alvo":"bolha_de_filtro","confianca":1.0,"epistemico":"fato","origem":"wiki","rel":"relaciona"},{"alvo":"economia_da_atencao","confianca":1.0,"epistemico":"fato","origem":"wiki","rel":"usa"}],"confianca":1.0,"contraexemplos":[],"descricao":"A circunstância em que fatos objetivos pesam menos na opinião pública do que apelos à emoção e à crença pessoal (palavra do ano do Oxford, 2016). Abarca fake news e deepfakes (mídia sintética por IA), retroalimentada pela economia da atenção.","epistemico":"inferencia","exemplos":[],"id":"pos_verdade","memoria":"semantica","meta":{"autor":"—","dominio":"ciencias_de_fora","fonte":""},"origem":"wiki","palavras_chave":[],"sinonimos":[],"tipo":"conceito","titulo":"Pós-Verdade e Desinformação"}
\section{Pós-Verdade e Desinformação}
A circunstância em que fatos objetivos pesam menos na opinião pública do que apelos à emoção e à crença pessoal (palavra do ano do Oxford, 2016). Abarca fake news e deepfakes (mídia sintética por IA), retroalimentada pela economia da atenção.
\prov{relações: parte\_de comunicacao\_digital; relaciona deus\_mentiroso; relaciona dois\_sistemas\_kahneman; relaciona bolha\_de\_filtro; usa economia\_da\_atencao}
\prov{proveniência: wiki \textperiodcentered{} inferencia \textperiodcentered{} confiança 1.0 \textperiodcentered{} domínio ciencias\_de\_fora \textperiodcentered{} história 6 versões no jornal}

%CRISTAL {"arestas":[{"alvo":"codificacao_decodificacao","confianca":1.0,"epistemico":"fato","origem":"wiki","rel":"usa"},{"alvo":"hegemonia_cultural","confianca":1.0,"epistemico":"fato","origem":"wiki","rel":"relaciona"}],"confianca":1.0,"contraexemplos":[],"descricao":"Hall: ao decodificar um texto, a audiência pode ocupar a posição dominante-hegemônica (aceita o sentido preferencial), a negociada (aceita mas adapta) ou a oposicional (entende o código mas o rejeita e lê por um quadro alternativo).","epistemico":"inferencia","exemplos":[],"id":"posicoes_de_leitura","memoria":"semantica","meta":{"autor":"Hall","dominio":"ciencias_de_fora","fonte":""},"origem":"wiki","palavras_chave":[],"sinonimos":[],"tipo":"conceito","titulo":"Três Posições de Leitura (Hall)"}
\section{Três Posições de Leitura (Hall)}
Hall: ao decodificar um texto, a audiência pode ocupar a posição dominante-hegemônica (aceita o sentido preferencial), a negociada (aceita mas adapta) ou a oposicional (entende o código mas o rejeita e lê por um quadro alternativo).
\prov{relações: usa codificacao\_decodificacao; relaciona hegemonia\_cultural}
\prov{proveniência: wiki \textperiodcentered{} inferencia \textperiodcentered{} confiança 1.0 \textperiodcentered{} domínio ciencias\_de\_fora \textperiodcentered{} história 3 versões no jornal}

%CRISTAL {"arestas":[{"alvo":"determinismo_ambiental","confianca":1.0,"epistemico":"fato","origem":"wiki","rel":"contradiz"},{"alvo":"habitus","confianca":1.0,"epistemico":"fato","origem":"wiki","rel":"relaciona"}],"confianca":1.0,"contraexemplos":[],"descricao":"Vidal de la Blache: o meio não determina, apenas oferece um leque de possibilidades entre as quais o grupo humano escolhe segundo sua técnica e cultura, cristalizando-se em 'gêneros de vida'. Resposta francesa ao determinismo alemão.","epistemico":"inferencia","exemplos":[],"id":"possibilismo","memoria":"semantica","meta":{"autor":"Vidal de la Blache","dominio":"ciencias_de_fora","fonte":""},"origem":"wiki","palavras_chave":[],"sinonimos":[],"tipo":"conceito","titulo":"Possibilismo (Vidal de la Blache)"}
\section{Possibilismo (Vidal de la Blache)}
Vidal de la Blache: o meio não determina, apenas oferece um leque de possibilidades entre as quais o grupo humano escolhe segundo sua técnica e cultura, cristalizando-se em 'gêneros de vida'. Resposta francesa ao determinismo alemão.
\prov{relações: contradiz determinismo\_ambiental; relaciona habitus}
\prov{proveniência: wiki \textperiodcentered{} inferencia \textperiodcentered{} confiança 1.0 \textperiodcentered{} domínio ciencias\_de\_fora \textperiodcentered{} história 3 versões no jornal}

%CRISTAL {"arestas":[{"alvo":"neuronio","confianca":1.0,"epistemico":"fato","origem":"wiki","rel":"usa"}],"confianca":1.0,"contraexemplos":[],"descricao":"O pulso elétrico tudo-ou-nada que percorre o neurônio — a informação neural é digital no disparo, analógica na taxa.","epistemico":"fato","exemplos":[],"id":"potencial_de_acao","memoria":"semantica","meta":{"dominio":"ciencias_de_fora","fonte":""},"origem":"wiki","palavras_chave":[],"sinonimos":[],"tipo":"conceito","titulo":"Potencial de Ação"}
\section{Potencial de Ação}
O pulso elétrico tudo-ou-nada que percorre o neurônio — a informação neural é digital no disparo, analógica na taxa.
\prov{relações: usa neuronio}
\prov{proveniência: wiki \textperiodcentered{} fato \textperiodcentered{} confiança 1.0 \textperiodcentered{} domínio ciencias\_de\_fora \textperiodcentered{} história 2 versões no jornal}

%CRISTAL {"arestas":[{"alvo":"linguistica","confianca":1.0,"epistemico":"fato","origem":"wiki","rel":"parte_de"}],"confianca":1.0,"contraexemplos":[],"descricao":"O sentido no USO e no contexto — o que se faz ao dizer, para além do que a frase significa em si.","epistemico":"fato","exemplos":[],"id":"pragmatica","memoria":"semantica","meta":{"dominio":"ciencias_de_fora","fonte":""},"origem":"wiki","palavras_chave":[],"sinonimos":[],"tipo":"conceito","titulo":"Pragmática"}
\section{Pragmática}
O sentido no USO e no contexto — o que se faz ao dizer, para além do que a frase significa em si.
\prov{relações: parte\_de linguistica}
\prov{proveniência: wiki \textperiodcentered{} fato \textperiodcentered{} confiança 1.0 \textperiodcentered{} domínio ciencias\_de\_fora \textperiodcentered{} história 2 versões no jornal}

%CRISTAL {"arestas":[{"alvo":"charles_peirce","confianca":1.0,"epistemico":"fato","origem":"wiki","rel":"parte_de"},{"alvo":"significado_como_uso","confianca":1.0,"epistemico":"fato","origem":"wiki","rel":"relaciona"}],"confianca":1.0,"contraexemplos":[],"descricao":"Peirce: o sentido de uma ideia SÃO seus efeitos práticos concebíveis — a verdade se mede na consequência.","epistemico":"inferencia","exemplos":[],"id":"pragmatismo","memoria":"semantica","meta":{"autor":"Peirce","dominio":"ciencias_de_fora","fonte":""},"origem":"wiki","palavras_chave":[],"sinonimos":[],"tipo":"conceito","titulo":"Pragmatismo (Máxima de Peirce)"}
\section{Pragmatismo (Máxima de Peirce)}
Peirce: o sentido de uma ideia SÃO seus efeitos práticos concebíveis — a verdade se mede na consequência.
\prov{relações: parte\_de charles\_peirce; relaciona significado\_como\_uso}
\prov{proveniência: wiki \textperiodcentered{} inferencia \textperiodcentered{} confiança 1.0 \textperiodcentered{} domínio ciencias\_de\_fora \textperiodcentered{} história 3 versões no jornal}

%CRISTAL {"arestas":[{"alvo":"teorias_aprendizagem","confianca":1.0,"epistemico":"fato","origem":"wiki","rel":"parte_de"},{"alvo":"metacognicao","confianca":1.0,"epistemico":"fato","origem":"wiki","rel":"usa"}],"confianca":1.0,"contraexemplos":[],"descricao":"Ericsson: treino estruturado com objetivo explícito de melhorar aspectos específicos, focado nos pontos fracos, no limite da capacidade atual e com feedback imediato — distinta da mera repetição.","epistemico":"fato","exemplos":[],"id":"pratica_deliberada","memoria":"semantica","meta":{"autor":"Ericsson","dominio":"ciencias_de_fora","fonte":""},"origem":"wiki","palavras_chave":[],"sinonimos":[],"tipo":"conceito","titulo":"Prática Deliberada (Ericsson)"}
\section{Prática Deliberada (Ericsson)}
Ericsson: treino estruturado com objetivo explícito de melhorar aspectos específicos, focado nos pontos fracos, no limite da capacidade atual e com feedback imediato — distinta da mera repetição.
\prov{relações: parte\_de teorias\_aprendizagem; usa metacognicao}
\prov{proveniência: wiki \textperiodcentered{} fato \textperiodcentered{} confiança 1.0 \textperiodcentered{} domínio ciencias\_de\_fora \textperiodcentered{} história 3 versões no jornal}

%CRISTAL {"arestas":[{"alvo":"separacao_poderes","confianca":1.0,"epistemico":"fato","origem":"wiki","rel":"usa"},{"alvo":"ciencia_politica","confianca":1.0,"epistemico":"fato","origem":"wiki","rel":"parte_de"}],"confianca":1.0,"contraexemplos":[],"descricao":"Linz: no parlamentarismo o governo depende da confiança do parlamento ('flexível'); no presidencialismo, executivo e legislativo têm eleição direta e mandatos fixos ('rígido'), gerando dupla legitimidade, lógica de soma zero e maior risco de conflito interpoderes.","epistemico":"inferencia","exemplos":[],"id":"presidencialismo_parlamentarismo","memoria":"semantica","meta":{"autor":"Linz","dominio":"ciencias_de_fora","fonte":""},"origem":"wiki","palavras_chave":[],"sinonimos":[],"tipo":"conceito","titulo":"Presidencialismo vs Parlamentarismo (Linz)"}
\section{Presidencialismo vs Parlamentarismo (Linz)}
Linz: no parlamentarismo o governo depende da confiança do parlamento ('flexível'); no presidencialismo, executivo e legislativo têm eleição direta e mandatos fixos ('rígido'), gerando dupla legitimidade, lógica de soma zero e maior risco de conflito interpoderes.
\prov{relações: usa separacao\_poderes; parte\_de ciencia\_politica}
\prov{proveniência: wiki \textperiodcentered{} inferencia \textperiodcentered{} confiança 1.0 \textperiodcentered{} domínio ciencias\_de\_fora \textperiodcentered{} história 3 versões no jornal}

%CRISTAL {"arestas":[{"alvo":"bioetica","confianca":1.0,"epistemico":"fato","origem":"wiki","rel":"parte_de"},{"alvo":"imperativo_categorico","confianca":1.0,"epistemico":"fato","origem":"wiki","rel":"usa"},{"alvo":"utilitarismo","confianca":1.0,"epistemico":"fato","origem":"wiki","rel":"usa"},{"alvo":"ethics_of_care","confianca":1.0,"epistemico":"fato","origem":"wiki","rel":"relaciona"}],"confianca":1.0,"contraexemplos":[],"descricao":"Beauchamp & Childress: a ética biomédica em quatro princípios prima facie — autonomia, beneficência, não-maleficência (primum non nocere, raiz no Juramento de Hipócrates) e justiça — ponderados caso a caso, sem hierarquia fixa.","epistemico":"inferencia","exemplos":[],"id":"principialismo","memoria":"semantica","meta":{"autor":"Beauchamp & Childress","dominio":"ciencias_de_fora","fonte":""},"origem":"wiki","palavras_chave":[],"sinonimos":[],"tipo":"conceito","titulo":"Principialismo (os Quatro Princípios)"}
\section{Principialismo (os Quatro Princípios)}
Beauchamp \& Childress: a ética biomédica em quatro princípios prima facie — autonomia, beneficência, não-maleficência (primum non nocere, raiz no Juramento de Hipócrates) e justiça — ponderados caso a caso, sem hierarquia fixa.
\prov{relações: parte\_de bioetica; usa imperativo\_categorico; usa utilitarismo; relaciona ethics\_of\_care}
\prov{proveniência: wiki \textperiodcentered{} inferencia \textperiodcentered{} confiança 1.0 \textperiodcentered{} domínio ciencias\_de\_fora \textperiodcentered{} história 5 versões no jornal}

%CRISTAL {"arestas":[{"alvo":"fim_de_vida_espectro","confianca":1.0,"epistemico":"fato","origem":"wiki","rel":"usa"},{"alvo":"cuidados_paliativos","confianca":1.0,"epistemico":"fato","origem":"wiki","rel":"relaciona"}],"confianca":1.0,"contraexemplos":[],"descricao":"Aquino: distingue o efeito intencionado do apenas previsto — é lícito dar analgesia que possa abreviar a vida se a intenção é aliviar a dor e não matar, desde que o bem não seja obtido pelo mal. Separa a sedação paliativa da eutanásia.","epistemico":"inferencia","exemplos":[],"id":"principio_duplo_efeito","memoria":"semantica","meta":{"autor":"Aquino","dominio":"ciencias_de_fora","fonte":""},"origem":"wiki","palavras_chave":[],"sinonimos":[],"tipo":"conceito","titulo":"Doutrina do Duplo Efeito"}
\section{Doutrina do Duplo Efeito}
Aquino: distingue o efeito intencionado do apenas previsto — é lícito dar analgesia que possa abreviar a vida se a intenção é aliviar a dor e não matar, desde que o bem não seja obtido pelo mal. Separa a sedação paliativa da eutanásia.
\prov{relações: usa fim\_de\_vida\_espectro; relaciona cuidados\_paliativos}
\prov{proveniência: wiki \textperiodcentered{} inferencia \textperiodcentered{} confiança 1.0 \textperiodcentered{} domínio ciencias\_de\_fora \textperiodcentered{} história 3 versões no jornal}

%CRISTAL {"arestas":[{"alvo":"filosofia_da_tecnica","confianca":1.0,"epistemico":"fato","origem":"wiki","rel":"parte_de"},{"alvo":"contrato_social","confianca":1.0,"epistemico":"fato","origem":"wiki","rel":"relaciona"},{"alvo":"biopoder","confianca":1.0,"epistemico":"fato","origem":"wiki","rel":"relaciona"}],"confianca":1.0,"contraexemplos":[],"descricao":"Jonas: diante do poder tecnológico sobre o futuro, um novo imperativo — 'age de modo que os efeitos da tua ação sejam compatíveis com a permanência de uma vida humana autêntica na Terra' — ancorado numa 'heurística do medo'.","epistemico":"inferencia","exemplos":[],"id":"principio_responsabilidade","memoria":"semantica","meta":{"autor":"Jonas","dominio":"ciencias_de_fora","fonte":""},"origem":"wiki","palavras_chave":[],"sinonimos":[],"tipo":"conceito","titulo":"Princípio Responsabilidade (Jonas)"}
\section{Princípio Responsabilidade (Jonas)}
Jonas: diante do poder tecnológico sobre o futuro, um novo imperativo — 'age de modo que os efeitos da tua ação sejam compatíveis com a permanência de uma vida humana autêntica na Terra' — ancorado numa 'heurística do medo'.
\prov{relações: parte\_de filosofia\_da\_tecnica; relaciona contrato\_social; relaciona biopoder}
\prov{proveniência: wiki \textperiodcentered{} inferencia \textperiodcentered{} confiança 1.0 \textperiodcentered{} domínio ciencias\_de\_fora \textperiodcentered{} história 4 versões no jornal}

%CRISTAL {"arestas":[{"alvo":"maquina_de_turing","confianca":1.0,"epistemico":"fato","origem":"wiki","rel":"usa"},{"alvo":"incompletude_godel","confianca":1.0,"epistemico":"fato","origem":"wiki","rel":"relaciona"},{"alvo":"programa_de_hilbert","confianca":1.0,"epistemico":"fato","origem":"wiki","rel":"contradiz"},{"alvo":"expressividade_sub_turing","confianca":1.0,"epistemico":"fato","origem":"wiki","rel":"explica"}],"confianca":1.0,"contraexemplos":[],"descricao":"Turing: não há algoritmo que decida, para todo programa, se ele para. O limite formal da computação, por diagonalização.","epistemico":"fato","exemplos":[],"id":"problema_da_parada","memoria":"semantica","meta":{"autor":"Turing","dominio":"ciencias_de_fora","fonte":""},"origem":"wiki","palavras_chave":[],"sinonimos":[],"tipo":"conceito","titulo":"Problema da Parada"}
\section{Problema da Parada}
Turing: não há algoritmo que decida, para todo programa, se ele para. O limite formal da computação, por diagonalização.
\prov{relações: usa maquina\_de\_turing; relaciona incompletude\_godel; contradiz programa\_de\_hilbert; explica expressividade\_sub\_turing}
\prov{proveniência: wiki \textperiodcentered{} fato \textperiodcentered{} confiança 1.0 \textperiodcentered{} domínio ciencias\_de\_fora \textperiodcentered{} história 11 versões no jornal}

%CRISTAL {"arestas":[{"alvo":"producao_do_espaco_lefebvre","confianca":1.0,"epistemico":"fato","origem":"wiki","rel":"usa"},{"alvo":"mais_valia","confianca":1.0,"epistemico":"fato","origem":"wiki","rel":"usa"}],"confianca":1.0,"contraexemplos":[],"descricao":"Harvey: o capitalismo produz e reorganiza continuamente o espaço para absorver excedentes e adiar suas crises (o 'ajuste espacial'/spatial fix), fazendo da urbanização e da geografia parte constitutiva da acumulação.","epistemico":"inferencia","exemplos":[],"id":"producao_capitalista_espaco_harvey","memoria":"semantica","meta":{"autor":"Harvey","dominio":"ciencias_de_fora","fonte":""},"origem":"wiki","palavras_chave":[],"sinonimos":[],"tipo":"conceito","titulo":"Produção Capitalista do Espaço / Spatial Fix (Harvey)"}
\section{Produção Capitalista do Espaço / Spatial Fix (Harvey)}
Harvey: o capitalismo produz e reorganiza continuamente o espaço para absorver excedentes e adiar suas crises (o 'ajuste espacial'/spatial fix), fazendo da urbanização e da geografia parte constitutiva da acumulação.
\prov{relações: usa producao\_do\_espaco\_lefebvre; usa mais\_valia}
\prov{proveniência: wiki \textperiodcentered{} inferencia \textperiodcentered{} confiança 1.0 \textperiodcentered{} domínio ciencias\_de\_fora \textperiodcentered{} história 3 versões no jornal}

%CRISTAL {"arestas":[{"alvo":"geografia_humana","confianca":1.0,"epistemico":"fato","origem":"wiki","rel":"parte_de"},{"alvo":"materialismo_historico","confianca":1.0,"epistemico":"fato","origem":"wiki","rel":"usa"},{"alvo":"espaco_santos","confianca":1.0,"epistemico":"fato","origem":"wiki","rel":"relaciona"}],"confianca":1.0,"contraexemplos":[],"descricao":"Lefebvre: o espaço (social) é um produto social, não um palco neutro nem um dado natural — cada modo de produção produz seu próprio espaço, que reproduz as relações sociais que o geraram. Produzido em três momentos: o percebido (prática), o concebido (poder/técnica) e o vivido (símbolos, experiência).","epistemico":"inferencia","exemplos":[],"id":"producao_do_espaco_lefebvre","memoria":"semantica","meta":{"autor":"Lefebvre","dominio":"ciencias_de_fora","fonte":""},"origem":"wiki","palavras_chave":[],"sinonimos":[],"tipo":"conceito","titulo":"A Produção do Espaço (Lefebvre)"}
\section{A Produção do Espaço (Lefebvre)}
Lefebvre: o espaço (social) é um produto social, não um palco neutro nem um dado natural — cada modo de produção produz seu próprio espaço, que reproduz as relações sociais que o geraram. Produzido em três momentos: o percebido (prática), o concebido (poder/técnica) e o vivido (símbolos, experiência).
\prov{relações: parte\_de geografia\_humana; usa materialismo\_historico; relaciona espaco\_santos}
\prov{proveniência: wiki \textperiodcentered{} inferencia \textperiodcentered{} confiança 1.0 \textperiodcentered{} domínio ciencias\_de\_fora \textperiodcentered{} história 4 versões no jornal}

%CRISTAL {"arestas":[{"alvo":"comunicacao_digital","confianca":1.0,"epistemico":"fato","origem":"wiki","rel":"parte_de"},{"alvo":"inteligencia_coletiva","confianca":1.0,"epistemico":"fato","origem":"wiki","rel":"relaciona"},{"alvo":"goldchain","confianca":1.0,"epistemico":"fato","origem":"wiki","rel":"relaciona"}],"confianca":1.0,"contraexemplos":[],"descricao":"Benkler: um modelo em que grandes números de indivíduos colaboram de forma descentralizada e auto-selecionada sobre recursos em comum, movidos por recompensas intrínsecas — Wikipédia, software livre. A riqueza das redes.","epistemico":"inferencia","exemplos":[],"id":"producao_entre_pares","memoria":"semantica","meta":{"autor":"Benkler","dominio":"ciencias_de_fora","fonte":""},"origem":"wiki","palavras_chave":[],"sinonimos":[],"tipo":"conceito","titulo":"Produção entre Pares (Benkler)"}
\section{Produção entre Pares (Benkler)}
Benkler: um modelo em que grandes números de indivíduos colaboram de forma descentralizada e auto-selecionada sobre recursos em comum, movidos por recompensas intrínsecas — Wikipédia, software livre. A riqueza das redes.
\prov{relações: parte\_de comunicacao\_digital; relaciona inteligencia\_coletiva; relaciona goldchain}
\prov{proveniência: wiki \textperiodcentered{} inferencia \textperiodcentered{} confiança 1.0 \textperiodcentered{} domínio ciencias\_de\_fora \textperiodcentered{} história 4 versões no jornal}

%CRISTAL {"arestas":[{"alvo":"teoria_de_grupos","confianca":1.0,"epistemico":"fato","origem":"wiki","rel":"usa"},{"alvo":"geometria_ramo","confianca":1.0,"epistemico":"fato","origem":"wiki","rel":"explica"},{"alvo":"orbita_z2_3","confianca":1.0,"epistemico":"fato","origem":"wiki","rel":"relaciona"}],"confianca":1.0,"contraexemplos":[],"descricao":"Klein: cada geometria É o grupo de transformações que preservam suas figuras — geometria reduzida a simetria/invariante.","epistemico":"fato","exemplos":[],"id":"programa_de_erlangen","memoria":"semantica","meta":{"autor":"Klein","dominio":"ciencias_de_fora","fonte":""},"origem":"wiki","palavras_chave":[],"sinonimos":[],"tipo":"conceito","titulo":"Programa de Erlangen"}
\section{Programa de Erlangen}
Klein: cada geometria É o grupo de transformações que preservam suas figuras — geometria reduzida a simetria/invariante.
\prov{relações: usa teoria\_de\_grupos; explica geometria\_ramo; relaciona orbita\_z2\_3}
\prov{proveniência: wiki \textperiodcentered{} fato \textperiodcentered{} confiança 1.0 \textperiodcentered{} domínio ciencias\_de\_fora \textperiodcentered{} história 7 versões no jornal}

%CRISTAL {"arestas":[{"alvo":"metodo_axiomatico","confianca":1.0,"epistemico":"fato","origem":"wiki","rel":"usa"},{"alvo":"logica_matematica","confianca":1.0,"epistemico":"fato","origem":"wiki","rel":"parte_de"}],"confianca":1.0,"contraexemplos":[],"descricao":"Hilbert: fundar toda a matemática num sistema formal completo, consistente e decidível. Gödel e Turing o refutaram.","epistemico":"inferencia","exemplos":[],"id":"programa_de_hilbert","memoria":"semantica","meta":{"autor":"Hilbert","dominio":"ciencias_de_fora","fonte":""},"origem":"wiki","palavras_chave":[],"sinonimos":[],"tipo":"conceito","titulo":"Programa de Hilbert"}
\section{Programa de Hilbert}
Hilbert: fundar toda a matemática num sistema formal completo, consistente e decidível. Gödel e Turing o refutaram.
\prov{relações: usa metodo\_axiomatico; parte\_de logica\_matematica}
\prov{proveniência: wiki \textperiodcentered{} inferencia \textperiodcentered{} confiança 1.0 \textperiodcentered{} domínio ciencias\_de\_fora \textperiodcentered{} história 6 versões no jornal}

%CRISTAL {"arestas":[{"alvo":"composicao_visual","confianca":1.0,"epistemico":"fato","origem":"wiki","rel":"usa"},{"alvo":"critica_proporcao_aurea","confianca":1.0,"epistemico":"fato","origem":"wiki","rel":"usa"}],"confianca":1.0,"contraexemplos":[],"descricao":"A busca de razões 'certas' na forma. A da 'proporção áurea universal' é, em grande parte, mito retrospectivo.","epistemico":"hipotese","exemplos":[],"id":"proporcao_na_arte","memoria":"semantica","meta":{"dominio":"ciencias_de_fora","fonte":""},"origem":"wiki","palavras_chave":[],"sinonimos":[],"tipo":"conceito","titulo":"Proporção na Arte"}
\section{Proporção na Arte}
A busca de razões 'certas' na forma. A da 'proporção áurea universal' é, em grande parte, mito retrospectivo.
\prov{relações: usa composicao\_visual; usa critica\_proporcao\_aurea}
\prov{proveniência: wiki \textperiodcentered{} hipotese \textperiodcentered{} confiança 1.0 \textperiodcentered{} domínio ciencias\_de\_fora \textperiodcentered{} história 3 versões no jornal}

%CRISTAL {"arestas":[{"alvo":"propriedade_lockeana","confianca":1.0,"epistemico":"fato","origem":"wiki","rel":"contradiz"},{"alvo":"tragedia_dos_comuns","confianca":1.0,"epistemico":"fato","origem":"wiki","rel":"relaciona"}],"confianca":1.0,"contraexemplos":[],"descricao":"Proudhon: a propriedade que rende ao proprietário sem trabalho (juro, aluguel, lucro) confisca o excedente do trabalhador, sendo uma forma institucionalizada de roubo. Antítese de Locke.","epistemico":"inferencia","exemplos":[],"id":"propriedade_e_roubo","memoria":"semantica","meta":{"autor":"Proudhon","dominio":"ciencias_de_fora","fonte":""},"origem":"wiki","palavras_chave":[],"sinonimos":[],"tipo":"conceito","titulo":"A Propriedade é um Roubo (Proudhon)"}
\section{A Propriedade é um Roubo (Proudhon)}
Proudhon: a propriedade que rende ao proprietário sem trabalho (juro, aluguel, lucro) confisca o excedente do trabalhador, sendo uma forma institucionalizada de roubo. Antítese de Locke.
\prov{relações: contradiz propriedade\_lockeana; relaciona tragedia\_dos\_comuns}
\prov{proveniência: wiki \textperiodcentered{} inferencia \textperiodcentered{} confiança 1.0 \textperiodcentered{} domínio ciencias\_de\_fora \textperiodcentered{} história 3 versões no jornal}

%CRISTAL {"arestas":[{"alvo":"direito","confianca":1.0,"epistemico":"fato","origem":"wiki","rel":"parte_de"},{"alvo":"prova_de_trabalho","confianca":1.0,"epistemico":"fato","origem":"wiki","rel":"explica"},{"alvo":"escassez","confianca":1.0,"epistemico":"fato","origem":"wiki","rel":"relaciona"}],"confianca":1.0,"contraexemplos":[],"descricao":"Locke: do bem comum dado a todos, o indivíduo torna próprio aquilo em que 'mistura' seu trabalho, legitimando a apropriação privada — sob a ressalva de deixar 'o bastante e tão bom' para os demais.","epistemico":"inferencia","exemplos":[],"id":"propriedade_lockeana","memoria":"semantica","meta":{"autor":"Locke","dominio":"ciencias_de_fora","fonte":""},"origem":"wiki","palavras_chave":[],"sinonimos":[],"tipo":"conceito","titulo":"Teoria do Trabalho da Propriedade (Locke)"}
\section{Teoria do Trabalho da Propriedade (Locke)}
Locke: do bem comum dado a todos, o indivíduo torna próprio aquilo em que 'mistura' seu trabalho, legitimando a apropriação privada — sob a ressalva de deixar 'o bastante e tão bom' para os demais.
\prov{relações: parte\_de direito; explica prova\_de\_trabalho; relaciona escassez}
\prov{proveniência: wiki \textperiodcentered{} inferencia \textperiodcentered{} confiança 1.0 \textperiodcentered{} domínio ciencias\_de\_fora \textperiodcentered{} história 4 versões no jornal}

%CRISTAL {"arestas":[{"alvo":"mudanca_linguistica","confianca":1.0,"epistemico":"fato","origem":"wiki","rel":"usa"}],"confianca":1.0,"contraexemplos":[],"descricao":"A língua-mãe reconstruída de uma família, não atestada em registro — inferida comparando as filhas (o proto-indo-europeu). Reconstrução, não documento.","epistemico":"inferencia","exemplos":[],"id":"protolingua","memoria":"semantica","meta":{"dominio":"ciencias_de_fora","fonte":""},"origem":"wiki","palavras_chave":[],"sinonimos":[],"tipo":"conceito","titulo":"Proto-língua"}
\section{Proto-língua}
A língua-mãe reconstruída de uma família, não atestada em registro — inferida comparando as filhas (o proto-indo-europeu). Reconstrução, não documento.
\prov{relações: usa mudanca\_linguistica}
\prov{proveniência: wiki \textperiodcentered{} inferencia \textperiodcentered{} confiança 1.0 \textperiodcentered{} domínio ciencias\_de\_fora \textperiodcentered{} história 2 versões no jornal}

%CRISTAL {"arestas":[{"alvo":"valor_trabalho","confianca":1.0,"epistemico":"fato","origem":"wiki","rel":"especializa"},{"alvo":"goldchain","confianca":1.0,"epistemico":"fato","origem":"wiki","rel":"parte_de"},{"alvo":"sequencia_pell","confianca":1.0,"epistemico":"fato","origem":"wiki","rel":"relaciona"}],"confianca":1.0,"contraexemplos":[],"descricao":"Consenso por resolver um desafio criptográfico que exige energia real (Bitcoin); torna o trabalho verificável e imutável — o valor-trabalho de Ricardo retificado.","epistemico":"fato","exemplos":[],"id":"prova_de_trabalho","memoria":"semantica","meta":{"autor":"Nakamoto","dominio":"ciencias_de_fora","fonte":""},"origem":"wiki","palavras_chave":[],"sinonimos":[],"tipo":"conceito","titulo":"Prova de Trabalho (PoW)"}
\section{Prova de Trabalho (PoW)}
Consenso por resolver um desafio criptográfico que exige energia real (Bitcoin); torna o trabalho verificável e imutável — o valor-trabalho de Ricardo retificado.
\prov{relações: especializa valor\_trabalho; parte\_de goldchain; relaciona sequencia\_pell}
\prov{proveniência: wiki \textperiodcentered{} fato \textperiodcentered{} confiança 1.0 \textperiodcentered{} domínio ciencias\_de\_fora \textperiodcentered{} história 4 versões no jornal}

%CRISTAL {"arestas":[{"alvo":"didatica_matematica","confianca":1.0,"epistemico":"fato","origem":"wiki","rel":"parte_de"},{"alvo":"intuicionismo_brouwer","confianca":1.0,"epistemico":"fato","origem":"wiki","rel":"relaciona"},{"alvo":"formalismo_hilbert","confianca":1.0,"epistemico":"fato","origem":"wiki","rel":"referencia"},{"alvo":"godel_incompletude","confianca":1.0,"epistemico":"fato","origem":"wiki","rel":"referencia"}],"confianca":1.0,"contraexemplos":[],"descricao":"A prova não só verifica, mas explica; o rigor deve destruir a intuição ruim e elevar a boa. O aprendiz atravessa fases pré-rigorosa (intuitiva), rigorosa (formal) e pós-rigorosa (rigor + intuição integrados).","epistemico":"inferencia","exemplos":[],"id":"prova_rigor_intuicao","memoria":"semantica","meta":{"autor":"Tao/de Villiers","dominio":"ciencias_de_fora","fonte":""},"origem":"wiki","palavras_chave":[],"sinonimos":[],"tipo":"conceito","titulo":"Prova, Rigor vs Intuição no Ensino"}
\section{Prova, Rigor vs Intuição no Ensino}
A prova não só verifica, mas explica; o rigor deve destruir a intuição ruim e elevar a boa. O aprendiz atravessa fases pré-rigorosa (intuitiva), rigorosa (formal) e pós-rigorosa (rigor + intuição integrados).
\prov{relações: parte\_de didatica\_matematica; relaciona intuicionismo\_brouwer; referencia formalismo\_hilbert; referencia godel\_incompletude}
\prov{proveniência: wiki \textperiodcentered{} inferencia \textperiodcentered{} confiança 1.0 \textperiodcentered{} domínio ciencias\_de\_fora \textperiodcentered{} história 5 versões no jornal}

%CRISTAL {"arestas":[{"alvo":"psicologia","confianca":1.0,"epistemico":"fato","origem":"wiki","rel":"parte_de"},{"alvo":"utilitarismo","confianca":0.5,"epistemico":"fato","origem":"wiki","rel":"relaciona"},{"alvo":"virtualizacao","confianca":0.5,"epistemico":"fato","origem":"wiki","rel":"relaciona"},{"alvo":"word","confianca":0.5,"epistemico":"fato","origem":"wiki","rel":"relaciona"},{"alvo":"semiotica_peirce_triade","confianca":0.5,"epistemico":"fato","origem":"wiki","rel":"relaciona"},{"alvo":"universidade_curriculo_em_camadas","confianca":0.5,"epistemico":"fato","origem":"wiki","rel":"relaciona"}],"confianca":1.0,"contraexemplos":[],"descricao":"A escola clínica de Freud sobre o inconsciente; muito de seu arcabouço é não-falsável (Popper).","epistemico":"hipotese","exemplos":[],"id":"psicanalise","memoria":"semantica","meta":{"dominio":"ciencias_de_fora","fonte":""},"origem":"wiki","palavras_chave":[],"sinonimos":[],"tipo":"conceito","titulo":"Psicanálise"}
\section{Psicanálise}
A escola clínica de Freud sobre o inconsciente; muito de seu arcabouço é não-falsável (Popper).
\prov{relações: parte\_de psicologia; relaciona utilitarismo; relaciona virtualizacao; relaciona word; relaciona semiotica\_peirce\_triade; relaciona universidade\_curriculo\_em\_camadas}
\prov{proveniência: wiki \textperiodcentered{} hipotese \textperiodcentered{} confiança 1.0 \textperiodcentered{} domínio ciencias\_de\_fora \textperiodcentered{} história 7 versões no jornal}

%CRISTAL {"arestas":[{"alvo":"consciencia","confianca":1.0,"epistemico":"fato","origem":"wiki","rel":"referencia"},{"alvo":"teoria_do_valor","confianca":0.5,"epistemico":"fato","origem":"wiki","rel":"relaciona"},{"alvo":"word","confianca":0.5,"epistemico":"fato","origem":"wiki","rel":"relaciona"},{"alvo":"virtualizacao","confianca":0.5,"epistemico":"fato","origem":"wiki","rel":"relaciona"},{"alvo":"vida-morte","confianca":0.5,"epistemico":"fato","origem":"wiki","rel":"relaciona"}],"confianca":1.0,"contraexemplos":[],"descricao":"A ciência da mente e do comportamento — cognição, emoção, motivação, desenvolvimento.","epistemico":"fato","exemplos":[],"id":"psicologia","memoria":"semantica","meta":{"dominio":"ciencias_de_fora","fonte":""},"origem":"wiki","palavras_chave":[],"sinonimos":[],"tipo":"conceito","titulo":"Psicologia"}
\section{Psicologia}
A ciência da mente e do comportamento — cognição, emoção, motivação, desenvolvimento.
\prov{relações: referencia consciencia; relaciona teoria\_do\_valor; relaciona word; relaciona virtualizacao; relaciona vida-morte}
\prov{proveniência: wiki \textperiodcentered{} fato \textperiodcentered{} confiança 1.0 \textperiodcentered{} domínio ciencias\_de\_fora \textperiodcentered{} história 6 versões no jornal}

%CRISTAL {"arestas":[{"alvo":"hard_problem","confianca":1.0,"epistemico":"fato","origem":"wiki","rel":"usa"},{"alvo":"word","confianca":0.5,"epistemico":"fato","origem":"wiki","rel":"relaciona"},{"alvo":"vida-morte","confianca":0.5,"epistemico":"fato","origem":"wiki","rel":"relaciona"},{"alvo":"virtualizacao","confianca":0.5,"epistemico":"fato","origem":"wiki","rel":"relaciona"}],"confianca":1.0,"contraexemplos":[],"descricao":"Qualia: o caráter fenomênico da experiência. O zumbi: fisicamente idêntico a um humano mas sem experiência — desafia o funcionalismo.","epistemico":"inferencia","exemplos":[],"id":"qualia_zumbi","memoria":"semantica","meta":{"autor":"Chalmers","dominio":"filosofia","fonte":""},"origem":"wiki","palavras_chave":[],"sinonimos":[],"tipo":"conceito","titulo":"Qualia e o Zumbi Filosófico"}
\section{Qualia e o Zumbi Filosófico}
Qualia: o caráter fenomênico da experiência. O zumbi: fisicamente idêntico a um humano mas sem experiência — desafia o funcionalismo.
\prov{relações: usa hard\_problem; relaciona word; relaciona vida-morte; relaciona virtualizacao}
\prov{proveniência: wiki \textperiodcentered{} inferencia \textperiodcentered{} confiança 1.0 \textperiodcentered{} domínio filosofia \textperiodcentered{} história 5 versões no jornal}

%CRISTAL {"arestas":[{"alvo":"algebra_abstrata","confianca":1.0,"epistemico":"fato","origem":"wiki","rel":"parte_de"},{"alvo":"numeros_complexos","confianca":1.0,"epistemico":"fato","origem":"wiki","rel":"generaliza"},{"alvo":"algebras_hipercomplexas","confianca":1.0,"epistemico":"fato","origem":"wiki","rel":"especializa"}],"confianca":1.0,"contraexemplos":[],"descricao":"Hamilton: números hipercomplexos i²=j²=k²=ijk=−1 que rotacionam o espaço 3D. O primeiro sistema não-comutativo deliberado.","epistemico":"fato","exemplos":[],"id":"quaternions","memoria":"semantica","meta":{"autor":"Hamilton","dominio":"ciencias_de_fora","fonte":""},"origem":"wiki","palavras_chave":[],"sinonimos":[],"tipo":"conceito","titulo":"Quaternions"}
\section{Quaternions}
Hamilton: números hipercomplexos i\ensuremath{^{2}}=j\ensuremath{^{2}}=k\ensuremath{^{2}}=ijk=\ensuremath{-}1 que rotacionam o espaço 3D. O primeiro sistema não-comutativo deliberado.
\prov{relações: parte\_de algebra\_abstrata; generaliza numeros\_complexos; especializa algebras\_hipercomplexas}
\prov{proveniência: wiki \textperiodcentered{} fato \textperiodcentered{} confiança 1.0 \textperiodcentered{} domínio ciencias\_de\_fora \textperiodcentered{} história 9 versões no jornal}

%CRISTAL {"arestas":[{"alvo":"cristianismo","confianca":1.0,"epistemico":"fato","origem":"wiki","rel":"parte_de"},{"alvo":"sentido_da_vida","confianca":1.0,"epistemico":"fato","origem":"wiki","rel":"relaciona"}],"confianca":1.0,"contraexemplos":[],"descricao":"Bíblia (Gênesis 3): a ruptura da relação com Deus pela desobediência humana livre — a serpente ('certamente não morrereis') e o fruto proibido —, que introduz o mal, a vergonha e a morte. O preço do conhecimento do bem e do mal; a liberdade como capacidade de ruptura.","epistemico":"inferencia","exemplos":[],"id":"queda_pecado","memoria":"semantica","meta":{"autor":"Bíblia","dominio":"sagrado","fonte":""},"origem":"wiki","palavras_chave":[],"sinonimos":[],"tipo":"conceito","titulo":"A Queda e o Pecado"}
\section{A Queda e o Pecado}
Bíblia (Gênesis 3): a ruptura da relação com Deus pela desobediência humana livre — a serpente ('certamente não morrereis') e o fruto proibido —, que introduz o mal, a vergonha e a morte. O preço do conhecimento do bem e do mal; a liberdade como capacidade de ruptura.
\prov{relações: parte\_de cristianismo; relaciona sentido\_da\_vida}
\prov{proveniência: wiki \textperiodcentered{} inferencia \textperiodcentered{} confiança 1.0 \textperiodcentered{} domínio sagrado \textperiodcentered{} história 3 versões no jornal}

%CRISTAL {"arestas":[{"alvo":"epidemiologia","confianca":1.0,"epistemico":"fato","origem":"wiki","rel":"explica"},{"alvo":"imunidade_de_rebanho","confianca":1.0,"epistemico":"fato","origem":"wiki","rel":"explica"},{"alvo":"criticidade","confianca":1.0,"epistemico":"fato","origem":"wiki","rel":"relaciona"},{"alvo":"mutacao_borda_eigen","confianca":1.0,"epistemico":"fato","origem":"wiki","rel":"relaciona"},{"alvo":"numero_de_salem","confianca":1.0,"epistemico":"fato","origem":"wiki","rel":"eh_um"},{"alvo":"lyapunov_da_orbita","confianca":1.0,"epistemico":"fato","origem":"wiki","rel":"usa"}],"confianca":1.0,"contraexemplos":[],"descricao":"O número reprodutivo R₀ governa o contágio (I_{t+1}=R₀·I_t): R₀<1 a epidemia apaga (cristal), R₀=1 o limiar (borda/Salem), R₀>1 surto exponencial (caos). λ=log R₀ é EXATAMENTE o k da fissão nuclear e o limiar de erro de Eigen — a mesma transição de fase pela terceira vez. A imunidade de rebanho é empurrar R₀ abaixo de 1 (para o cristal). Verif.: R₀=0.7 cristal, 1.0 borda, 2.5 caos.","epistemico":"inferencia","exemplos":[],"id":"r0_borda_epidemica","memoria":"semantica","meta":{"classes_r0":{"0.7":"cristal (a epidemia apaga)","1.0":"borda/Salem (limiar epidêmico)","2.5":"caos (surto exponencial)"},"dominio":"medicina","fonte":"medicina e diagnóstico"},"origem":"wiki","palavras_chave":[],"sinonimos":[],"tipo":"conceito","titulo":"R₀ = a Borda (o limiar epidêmico)"}
\section{R\ensuremath{_{0}} = a Borda (o limiar epidêmico)}
O número reprodutivo R\ensuremath{_{0}} governa o contágio (I\_\{t+1\}=R\ensuremath{_{0}}\textperiodcentered I\_t): R\ensuremath{_{0}}<1 a epidemia apaga (cristal), R\ensuremath{_{0}}=1 o limiar (borda/Salem), R\ensuremath{_{0}}>1 surto exponencial (caos). \ensuremath{\lambda}=log R\ensuremath{_{0}} é EXATAMENTE o k da fissão nuclear e o limiar de erro de Eigen — a mesma transição de fase pela terceira vez. A imunidade de rebanho é empurrar R\ensuremath{_{0}} abaixo de 1 (para o cristal). Verif.: R\ensuremath{_{0}}=0.7 cristal, 1.0 borda, 2.5 caos.
\prov{relações: explica epidemiologia; explica imunidade\_de\_rebanho; relaciona criticidade; relaciona mutacao\_borda\_eigen; eh\_um numero\_de\_salem; usa lyapunov\_da\_orbita}
\prov{proveniência: wiki \textperiodcentered{} medicina e diagnóstico \textperiodcentered{} inferencia \textperiodcentered{} confiança 1.0 \textperiodcentered{} domínio medicina \textperiodcentered{} história 7 versões no jornal}

%CRISTAL {"arestas":[{"alvo":"jusnaturalismo","confianca":1.0,"epistemico":"fato","origem":"wiki","rel":"usa"},{"alvo":"juspositivismo","confianca":1.0,"epistemico":"fato","origem":"wiki","rel":"contradiz"}],"confianca":1.0,"contraexemplos":[],"descricao":"Radbruch: o direito positivo prevalece em regra, mas quando a contradição entre lei e justiça atinge grau intolerável, a lei injusta cede e deixa de ser direito válido ('a injustiça extrema não é direito'). Resposta ao legalismo sob o nazismo.","epistemico":"inferencia","exemplos":[],"id":"radbruch_formula","memoria":"semantica","meta":{"autor":"Radbruch","dominio":"ciencias_de_fora","fonte":""},"origem":"wiki","palavras_chave":[],"sinonimos":[],"tipo":"conceito","titulo":"Fórmula de Radbruch"}
\section{Fórmula de Radbruch}
Radbruch: o direito positivo prevalece em regra, mas quando a contradição entre lei e justiça atinge grau intolerável, a lei injusta cede e deixa de ser direito válido ('a injustiça extrema não é direito'). Resposta ao legalismo sob o nazismo.
\prov{relações: usa jusnaturalismo; contradiz juspositivismo}
\prov{proveniência: wiki \textperiodcentered{} inferencia \textperiodcentered{} confiança 1.0 \textperiodcentered{} domínio ciencias\_de\_fora \textperiodcentered{} história 3 versões no jornal}

%CRISTAL {"arestas":[{"alvo":"biologia_evolutiva","confianca":1.0,"epistemico":"fato","origem":"wiki","rel":"parte_de"},{"alvo":"simbiose","confianca":1.0,"epistemico":"fato","origem":"wiki","rel":"relaciona"}],"confianca":1.0,"contraexemplos":[],"descricao":"Van Valen: correr para ficar no lugar — cada espécie evolui só para acompanhar a coevolução das outras.","epistemico":"inferencia","exemplos":[],"id":"rainha_vermelha","memoria":"semantica","meta":{"autor":"Van Valen","dominio":"ciencias_de_fora","fonte":""},"origem":"wiki","palavras_chave":[],"sinonimos":[],"tipo":"conceito","titulo":"Hipótese da Rainha Vermelha"}
\section{Hipótese da Rainha Vermelha}
Van Valen: correr para ficar no lugar — cada espécie evolui só para acompanhar a coevolução das outras.
\prov{relações: parte\_de biologia\_evolutiva; relaciona simbiose}
\prov{proveniência: wiki \textperiodcentered{} inferencia \textperiodcentered{} confiança 1.0 \textperiodcentered{} domínio ciencias\_de\_fora \textperiodcentered{} história 3 versões no jornal}

%CRISTAL {"arestas":[{"alvo":"advaita_vedanta","confianca":1.0,"epistemico":"fato","origem":"wiki","rel":"parte_de"},{"alvo":"word","confianca":0.5,"epistemico":"fato","origem":"wiki","rel":"relaciona"}],"confianca":1.0,"contraexemplos":[],"descricao":"Ramana Maharshi: a autoinvestigação do pensamento-'eu' dissolve o ego e revela o Self. O Lopes discorda: no vazio profundo não achou 'Eu Real', só o Vazio.","epistemico":"hipotese","exemplos":[],"id":"ramana_quem_sou_eu","memoria":"semantica","meta":{"autor":"Ramana Maharshi","dominio":"filosofia","fonte":""},"origem":"wiki","palavras_chave":[],"sinonimos":[],"tipo":"conceito","titulo":"Ramana: 'Quem sou eu?'"}
\section{Ramana: 'Quem sou eu?'}
Ramana Maharshi: a autoinvestigação do pensamento-'eu' dissolve o ego e revela o Self. O Lopes discorda: no vazio profundo não achou 'Eu Real', só o Vazio.
\prov{relações: parte\_de advaita\_vedanta; relaciona word}
\prov{proveniência: wiki \textperiodcentered{} hipotese \textperiodcentered{} confiança 1.0 \textperiodcentered{} domínio filosofia \textperiodcentered{} história 4 versões no jornal}

%CRISTAL {"arestas":[{"alvo":"teoria_institucional_dickie","confianca":1.0,"epistemico":"fato","origem":"wiki","rel":"usa"},{"alvo":"aura_benjamin","confianca":1.0,"epistemico":"fato","origem":"wiki","rel":"contradiz"},{"alvo":"mimese_expressao","confianca":1.0,"epistemico":"fato","origem":"wiki","rel":"contradiz"}],"confianca":1.0,"contraexemplos":[],"descricao":"Duchamp: um objeto manufaturado comum (o mictório 'Fonte', 1917) elevado a arte pela mera escolha, título e deslocamento de contexto do artista. O gesto e o contexto — não a manufatura nem a beleza — constituem a obra.","epistemico":"inferencia","exemplos":[],"id":"readymade_duchamp","memoria":"semantica","meta":{"autor":"Duchamp","dominio":"ciencias_de_fora","fonte":""},"origem":"wiki","palavras_chave":[],"sinonimos":[],"tipo":"conceito","titulo":"Ready-made (Duchamp)"}
\section{Ready-made (Duchamp)}
Duchamp: um objeto manufaturado comum (o mictório 'Fonte', 1917) elevado a arte pela mera escolha, título e deslocamento de contexto do artista. O gesto e o contexto — não a manufatura nem a beleza — constituem a obra.
\prov{relações: usa teoria\_institucional\_dickie; contradiz aura\_benjamin; contradiz mimese\_expressao}
\prov{proveniência: wiki \textperiodcentered{} inferencia \textperiodcentered{} confiança 1.0 \textperiodcentered{} domínio ciencias\_de\_fora \textperiodcentered{} história 4 versões no jornal}

%CRISTAL {"arestas":[{"alvo":"filosofia_do_direito","confianca":1.0,"epistemico":"fato","origem":"wiki","rel":"parte_de"},{"alvo":"dworkin_integridade","confianca":1.0,"epistemico":"fato","origem":"wiki","rel":"contradiz"}],"confianca":1.0,"contraexemplos":[],"descricao":"O direito é o que os tribunais de fato decidem; as normas escritas são previsões do comportamento judicial ('prophecies of what courts will do'), moldado por experiência, política e psicologia. Holmes: 'a vida do direito não tem sido lógica, mas experiência'.","epistemico":"inferencia","exemplos":[],"id":"realismo_juridico","memoria":"semantica","meta":{"autor":"Holmes/Llewellyn","dominio":"ciencias_de_fora","fonte":""},"origem":"wiki","palavras_chave":[],"sinonimos":[],"tipo":"conceito","titulo":"Realismo Jurídico"}
\section{Realismo Jurídico}
O direito é o que os tribunais de fato decidem; as normas escritas são previsões do comportamento judicial ('prophecies of what courts will do'), moldado por experiência, política e psicologia. Holmes: 'a vida do direito não tem sido lógica, mas experiência'.
\prov{relações: parte\_de filosofia\_do\_direito; contradiz dworkin\_integridade}
\prov{proveniência: wiki \textperiodcentered{} inferencia \textperiodcentered{} confiança 1.0 \textperiodcentered{} domínio ciencias\_de\_fora \textperiodcentered{} história 3 versões no jornal}

%CRISTAL {"arestas":[{"alvo":"literatura","confianca":1.0,"epistemico":"fato","origem":"wiki","rel":"parte_de"},{"alvo":"realismo_modernismo_posmodernismo","confianca":1.0,"epistemico":"fato","origem":"wiki","rel":"parte_de"}],"confianca":1.0,"contraexemplos":[],"descricao":"A arte que quer espelhar o mundo social como ele é — o detalhe verossímil contra a idealização romântica.","epistemico":"inferencia","exemplos":[],"id":"realismo_literario","memoria":"semantica","meta":{"dominio":"ciencias_de_fora","fonte":""},"origem":"wiki","palavras_chave":[],"sinonimos":[],"tipo":"conceito","titulo":"Realismo"}
\section{Realismo}
A arte que quer espelhar o mundo social como ele é — o detalhe verossímil contra a idealização romântica.
\prov{relações: parte\_de literatura; parte\_de realismo\_modernismo\_posmodernismo}
\prov{proveniência: wiki \textperiodcentered{} inferencia \textperiodcentered{} confiança 1.0 \textperiodcentered{} domínio ciencias\_de\_fora \textperiodcentered{} história 4 versões no jornal}

%CRISTAL {"arestas":[{"alvo":"literatura","confianca":1.0,"epistemico":"fato","origem":"wiki","rel":"parte_de"},{"alvo":"mimese_expressao","confianca":1.0,"epistemico":"fato","origem":"wiki","rel":"usa"},{"alvo":"estranhamento","confianca":1.0,"epistemico":"fato","origem":"wiki","rel":"relaciona"}],"confianca":1.0,"contraexemplos":[],"descricao":"Os grandes paradigmas da literatura moderna: o realismo busca o retrato fiel e a crítica social; o modernismo fragmenta a forma e explora a subjetividade; o pós-modernismo instaura ironia, metaficção, pastiche e ceticismo diante das grandes narrativas.","epistemico":"inferencia","exemplos":[],"id":"realismo_modernismo_posmodernismo","memoria":"semantica","meta":{"autor":"—","dominio":"ciencias_de_fora","fonte":""},"origem":"wiki","palavras_chave":[],"sinonimos":[],"tipo":"conceito","titulo":"Realismo, Modernismo e Pós-Modernismo"}
\section{Realismo, Modernismo e Pós-Modernismo}
Os grandes paradigmas da literatura moderna: o realismo busca o retrato fiel e a crítica social; o modernismo fragmenta a forma e explora a subjetividade; o pós-modernismo instaura ironia, metaficção, pastiche e ceticismo diante das grandes narrativas.
\prov{relações: parte\_de literatura; usa mimese\_expressao; relaciona estranhamento}
\prov{proveniência: wiki \textperiodcentered{} inferencia \textperiodcentered{} confiança 1.0 \textperiodcentered{} domínio ciencias\_de\_fora \textperiodcentered{} história 4 versões no jornal}

%CRISTAL {"arestas":[{"alvo":"relacoes_internacionais","confianca":1.0,"epistemico":"fato","origem":"wiki","rel":"parte_de"},{"alvo":"balanca_de_poder","confianca":1.0,"epistemico":"fato","origem":"wiki","rel":"explica"}],"confianca":1.0,"contraexemplos":[],"descricao":"A política internacional como luta por poder: para Morgenthau (clássico) enraizada na natureza humana; para Waltz (neorrealismo/estrutural) forçada pela anarquia do sistema (ausência de autoridade central), que obriga cada Estado à autoajuda e à busca de segurança.","epistemico":"inferencia","exemplos":[],"id":"realismo_ri","memoria":"semantica","meta":{"autor":"Morgenthau/Waltz","dominio":"ciencias_de_fora","fonte":""},"origem":"wiki","palavras_chave":[],"sinonimos":[],"tipo":"conceito","titulo":"Realismo nas RI (Morgenthau, Waltz)"}
\section{Realismo nas RI (Morgenthau, Waltz)}
A política internacional como luta por poder: para Morgenthau (clássico) enraizada na natureza humana; para Waltz (neorrealismo/estrutural) forçada pela anarquia do sistema (ausência de autoridade central), que obriga cada Estado à autoajuda e à busca de segurança.
\prov{relações: parte\_de relacoes\_internacionais; explica balanca\_de\_poder}
\prov{proveniência: wiki \textperiodcentered{} inferencia \textperiodcentered{} confiança 1.0 \textperiodcentered{} domínio ciencias\_de\_fora \textperiodcentered{} história 3 versões no jornal}

%CRISTAL {"arestas":[{"alvo":"teoria_dos_numeros","confianca":1.0,"epistemico":"fato","origem":"wiki","rel":"parte_de"},{"alvo":"gauss_principe","confianca":1.0,"epistemico":"fato","origem":"wiki","rel":"parte_de"}],"confianca":1.0,"contraexemplos":[],"descricao":"Gauss: a lei profunda que relaciona quando p é resíduo quadrático módulo q e vice-versa — o 'teorema áureo'.","epistemico":"fato","exemplos":[],"id":"reciprocidade_quadratica","memoria":"semantica","meta":{"autor":"Gauss","dominio":"ciencias_de_fora","fonte":""},"origem":"wiki","palavras_chave":[],"sinonimos":[],"tipo":"conceito","titulo":"Reciprocidade Quadrática"}
\section{Reciprocidade Quadrática}
Gauss: a lei profunda que relaciona quando p é resíduo quadrático módulo q e vice-versa — o 'teorema áureo'.
\prov{relações: parte\_de teoria\_dos\_numeros; parte\_de gauss\_principe}
\prov{proveniência: wiki \textperiodcentered{} fato \textperiodcentered{} confiança 1.0 \textperiodcentered{} domínio ciencias\_de\_fora \textperiodcentered{} história 6 versões no jornal}

%CRISTAL {"arestas":[{"alvo":"gramatica_gerativa","confianca":1.0,"epistemico":"fato","origem":"wiki","rel":"usa"},{"alvo":"recursao","confianca":1.0,"epistemico":"fato","origem":"wiki","rel":"relaciona"}],"confianca":1.0,"contraexemplos":[],"descricao":"Encaixar estrutura dentro de estrutura sem limite — o que permite infinitas frases com meios finitos. Marca da linguagem e do cálculo.","epistemico":"fato","exemplos":[],"id":"recursao_linguistica","memoria":"semantica","meta":{"autor":"Chomsky","dominio":"ciencias_de_fora","fonte":""},"origem":"wiki","palavras_chave":[],"sinonimos":[],"tipo":"conceito","titulo":"Recursão"}
\section{Recursão}
Encaixar estrutura dentro de estrutura sem limite — o que permite infinitas frases com meios finitos. Marca da linguagem e do cálculo.
\prov{relações: usa gramatica\_gerativa; relaciona recursao}
\prov{proveniência: wiki \textperiodcentered{} fato \textperiodcentered{} confiança 1.0 \textperiodcentered{} domínio ciencias\_de\_fora \textperiodcentered{} história 4 versões no jornal}

%CRISTAL {"arestas":[{"alvo":"inteligencia_consciencia_relacional","confianca":1.0,"epistemico":"fato","origem":"wiki","rel":"usa"},{"alvo":"espaco_metrico_quantico_gentil","confianca":1.0,"epistemico":"fato","origem":"wiki","rel":"relaciona"},{"alvo":"categorias_geograficas","confianca":1.0,"epistemico":"fato","origem":"wiki","rel":"relaciona"}],"confianca":1.0,"contraexemplos":[],"descricao":"Lopes: todo fenômeno só existe em relação a uma 'estrutura cognitiva de referência' (o cérebro como sistema de medição); trocar o referencial muda a percepção e pode tornar solúvel um problema antes insolúvel. O espaço do observador constitui o que se vê.","epistemico":"inferencia","exemplos":[],"id":"referencial_do_observador_gentil","memoria":"semantica","meta":{"autor":"Gentil Lopes","dominio":"ciencias_de_fora","fonte":""},"origem":"wiki","palavras_chave":[],"sinonimos":[],"tipo":"conceito","titulo":"Mudança de Referencial (Lopes)"}
\section{Mudança de Referencial (Lopes)}
Lopes: todo fenômeno só existe em relação a uma 'estrutura cognitiva de referência' (o cérebro como sistema de medição); trocar o referencial muda a percepção e pode tornar solúvel um problema antes insolúvel. O espaço do observador constitui o que se vê.
\prov{relações: usa inteligencia\_consciencia\_relacional; relaciona espaco\_metrico\_quantico\_gentil; relaciona categorias\_geograficas}
\prov{proveniência: wiki \textperiodcentered{} inferencia \textperiodcentered{} confiança 1.0 \textperiodcentered{} domínio ciencias\_de\_fora \textperiodcentered{} história 5 versões no jornal}

%CRISTAL {"arestas":[{"alvo":"estudos_de_midia","confianca":1.0,"epistemico":"fato","origem":"wiki","rel":"parte_de"},{"alvo":"industria_cultural","confianca":1.0,"epistemico":"fato","origem":"wiki","rel":"relaciona"},{"alvo":"hegemonia_cultural","confianca":1.0,"epistemico":"fato","origem":"wiki","rel":"relaciona"}],"confianca":1.0,"contraexemplos":[],"descricao":"Habermas: a esfera pública burguesa (debate racional-crítico entre cidadãos) degrada-se quando o capitalismo monopolista e a mídia de massa a 'refeudalizam' — transformam cidadãos em audiência e fabricam consenso em vez de deliberação.","epistemico":"inferencia","exemplos":[],"id":"refeudalizacao_esfera_publica","memoria":"semantica","meta":{"autor":"Habermas","dominio":"ciencias_de_fora","fonte":""},"origem":"wiki","palavras_chave":[],"sinonimos":[],"tipo":"conceito","titulo":"Refeudalização da Esfera Pública (Habermas)"}
\section{Refeudalização da Esfera Pública (Habermas)}
Habermas: a esfera pública burguesa (debate racional-crítico entre cidadãos) degrada-se quando o capitalismo monopolista e a mídia de massa a 'refeudalizam' — transformam cidadãos em audiência e fabricam consenso em vez de deliberação.
\prov{relações: parte\_de estudos\_de\_midia; relaciona industria\_cultural; relaciona hegemonia\_cultural}
\prov{proveniência: wiki \textperiodcentered{} inferencia \textperiodcentered{} confiança 1.0 \textperiodcentered{} domínio ciencias\_de\_fora \textperiodcentered{} história 4 versões no jornal}

%CRISTAL {"arestas":[{"alvo":"cristianismo","confianca":1.0,"epistemico":"fato","origem":"wiki","rel":"parte_de"},{"alvo":"sentido_da_vida","confianca":1.0,"epistemico":"fato","origem":"wiki","rel":"relaciona"}],"confianca":1.0,"contraexemplos":[],"descricao":"Bíblia: o senhorio de Deus que irrompe na história e se consuma no fim — uma ordem de justiça e vida inaugurada por Jesus, proclamada como já presente e ainda por vir. O horizonte último que confere direção à existência.","epistemico":"inferencia","exemplos":[],"id":"reino_de_deus","memoria":"semantica","meta":{"autor":"Bíblia","dominio":"sagrado","fonte":""},"origem":"wiki","palavras_chave":[],"sinonimos":[],"tipo":"conceito","titulo":"O Reino de Deus"}
\section{O Reino de Deus}
Bíblia: o senhorio de Deus que irrompe na história e se consuma no fim — uma ordem de justiça e vida inaugurada por Jesus, proclamada como já presente e ainda por vir. O horizonte último que confere direção à existência.
\prov{relações: parte\_de cristianismo; relaciona sentido\_da\_vida}
\prov{proveniência: wiki \textperiodcentered{} inferencia \textperiodcentered{} confiança 1.0 \textperiodcentered{} domínio sagrado \textperiodcentered{} história 3 versões no jornal}

%CRISTAL {"arestas":[{"alvo":"filosofia_da_tecnica","confianca":1.0,"epistemico":"fato","origem":"wiki","rel":"parte_de"},{"alvo":"extensoes_do_homem","confianca":1.0,"epistemico":"fato","origem":"wiki","rel":"relaciona"},{"alvo":"kapp_organprojektion","confianca":1.0,"epistemico":"fato","origem":"wiki","rel":"relaciona"}],"confianca":1.0,"contraexemplos":[],"descricao":"Ihde: a técnica media a experiência em quatro modos — corporificação (óculos, transparentes no uso), hermenêutica (ler um instrumento), alteridade (a tecnologia como quase-outro) e de fundo (infraestrutura que estrutura sem ser notada).","epistemico":"inferencia","exemplos":[],"id":"relacoes_humano_tecnologia","memoria":"semantica","meta":{"autor":"Ihde","dominio":"ciencias_de_fora","fonte":""},"origem":"wiki","palavras_chave":[],"sinonimos":[],"tipo":"conceito","titulo":"Relações Humano-Tecnologia (Ihde)"}
\section{Relações Humano-Tecnologia (Ihde)}
Ihde: a técnica media a experiência em quatro modos — corporificação (óculos, transparentes no uso), hermenêutica (ler um instrumento), alteridade (a tecnologia como quase-outro) e de fundo (infraestrutura que estrutura sem ser notada).
\prov{relações: parte\_de filosofia\_da\_tecnica; relaciona extensoes\_do\_homem; relaciona kapp\_organprojektion}
\prov{proveniência: wiki \textperiodcentered{} inferencia \textperiodcentered{} confiança 1.0 \textperiodcentered{} domínio ciencias\_de\_fora \textperiodcentered{} história 4 versões no jornal}

%CRISTAL {"arestas":[{"alvo":"ciencia_politica","confianca":1.0,"epistemico":"fato","origem":"wiki","rel":"parte_de"},{"alvo":"word","confianca":0.5,"epistemico":"fato","origem":"wiki","rel":"relaciona"}],"confianca":1.0,"contraexemplos":[],"descricao":"O estudo da política entre Estados na ausência de uma autoridade central — poder, guerra, cooperação e ordem no sistema anárquico mundial.","epistemico":"fato","exemplos":[],"id":"relacoes_internacionais","memoria":"semantica","meta":{"dominio":"ciencias_de_fora","fonte":""},"origem":"wiki","palavras_chave":[],"sinonimos":[],"tipo":"conceito","titulo":"Relações Internacionais"}
\section{Relações Internacionais}
O estudo da política entre Estados na ausência de uma autoridade central — poder, guerra, cooperação e ordem no sistema anárquico mundial.
\prov{relações: parte\_de ciencia\_politica; relaciona word}
\prov{proveniência: wiki \textperiodcentered{} fato \textperiodcentered{} confiança 1.0 \textperiodcentered{} domínio ciencias\_de\_fora \textperiodcentered{} história 3 versões no jornal}

%CRISTAL {"arestas":[{"alvo":"antropologia","confianca":1.0,"epistemico":"fato","origem":"wiki","rel":"relaciona"},{"alvo":"vida-morte","confianca":0.5,"epistemico":"fato","origem":"wiki","rel":"relaciona"},{"alvo":"virtualizacao","confianca":0.5,"epistemico":"fato","origem":"wiki","rel":"relaciona"},{"alvo":"word","confianca":0.5,"epistemico":"fato","origem":"wiki","rel":"relaciona"}],"confianca":1.0,"contraexemplos":[],"descricao":"O estudo comparado do sagrado, do mito e da fé — como fenômeno humano, não como verdade revelada.","epistemico":"fato","exemplos":[],"id":"religiao","memoria":"semantica","meta":{"dominio":"ciencias_de_fora","fonte":""},"origem":"wiki","palavras_chave":[],"sinonimos":[],"tipo":"conceito","titulo":"Religião (estudo comparado)"}
\section{Religião (estudo comparado)}
O estudo comparado do sagrado, do mito e da fé — como fenômeno humano, não como verdade revelada.
\prov{relações: relaciona antropologia; relaciona vida-morte; relaciona virtualizacao; relaciona word}
\prov{proveniência: wiki \textperiodcentered{} fato \textperiodcentered{} confiança 1.0 \textperiodcentered{} domínio ciencias\_de\_fora \textperiodcentered{} história 5 versões no jornal}

%CRISTAL {"arestas":[{"alvo":"religiao","confianca":1.0,"epistemico":"fato","origem":"wiki","rel":"parte_de"},{"alvo":"deus_mentiroso","confianca":1.0,"epistemico":"fato","origem":"wiki","rel":"relaciona"},{"alvo":"leis_de_deus_programacao","confianca":1.0,"epistemico":"fato","origem":"wiki","rel":"relaciona"}],"confianca":1.0,"contraexemplos":[],"descricao":"Marx/Feuerbach: a religião adormece a consciência de classe e é a projeção da essência humana num ser fictício. Ecoa o Lopes.","epistemico":"inferencia","exemplos":[],"id":"religiao_opio","memoria":"semantica","meta":{"autor":"Marx/Feuerbach","dominio":"ciencias_de_fora","fonte":""},"origem":"wiki","palavras_chave":[],"sinonimos":[],"tipo":"conceito","titulo":"Religião como Ópio / Projeção"}
\section{Religião como Ópio / Projeção}
Marx/Feuerbach: a religião adormece a consciência de classe e é a projeção da essência humana num ser fictício. Ecoa o Lopes.
\prov{relações: parte\_de religiao; relaciona deus\_mentiroso; relaciona leis\_de\_deus\_programacao}
\prov{proveniência: wiki \textperiodcentered{} inferencia \textperiodcentered{} confiança 1.0 \textperiodcentered{} domínio ciencias\_de\_fora \textperiodcentered{} história 5 versões no jornal}

%CRISTAL {"arestas":[{"alvo":"comunicacao_digital","confianca":1.0,"epistemico":"fato","origem":"wiki","rel":"parte_de"},{"alvo":"aldeia_global","confianca":1.0,"epistemico":"fato","origem":"wiki","rel":"relaciona"}],"confianca":1.0,"contraexemplos":[],"descricao":"Bolter & Grusin: cada novo meio se define apropriando e refashionando os anteriores, pela dupla lógica da imediação (transparência que faz esquecer o meio) e da hipermediação (multiplicidade que exibe o meio).","epistemico":"inferencia","exemplos":[],"id":"remediacao","memoria":"semantica","meta":{"autor":"Bolter & Grusin","dominio":"ciencias_de_fora","fonte":""},"origem":"wiki","palavras_chave":[],"sinonimos":[],"tipo":"conceito","titulo":"Remediação (Bolter & Grusin)"}
\section{Remediação (Bolter \& Grusin)}
Bolter \& Grusin: cada novo meio se define apropriando e refashionando os anteriores, pela dupla lógica da imediação (transparência que faz esquecer o meio) e da hipermediação (multiplicidade que exibe o meio).
\prov{relações: parte\_de comunicacao\_digital; relaciona aldeia\_global}
\prov{proveniência: wiki \textperiodcentered{} inferencia \textperiodcentered{} confiança 1.0 \textperiodcentered{} domínio ciencias\_de\_fora \textperiodcentered{} história 3 versões no jornal}

%CRISTAL {"arestas":[{"alvo":"dna_dupla_helice","confianca":1.0,"epistemico":"fato","origem":"wiki","rel":"usa"}],"confianca":1.0,"contraexemplos":[],"descricao":"REPLICAÇÃO semiconservativa: cada fita-mãe é molde para uma nova; cada dupla-hélice-filha tem uma fita antiga e uma nova (DNA polimerase, 5'→3'). Meselson-Stahl 1958 ('o experimento mais belo da biologia', ¹⁵N/¹⁴N). PCR (Mullis 1983-85, Nobel 1993): amplifica ~2ⁿ cópias de um alvo por ciclos térmicos com Taq polimerase — tecnologia onipresente. Fato consolidado.","epistemico":"fato","exemplos":[],"id":"replicacao_pcr","memoria":"semantica","meta":{"autor":"broca-so","dominio":"biologia","fonte":"Meselson-Stahl 1958 PNAS; Mullis Nobel 1993"},"origem":"wiki","palavras_chave":[],"sinonimos":[],"tipo":"conceito","titulo":"Replicação semiconservativa e a PCR"}
\section{Replicação semiconservativa e a PCR}
REPLICAÇÃO semiconservativa: cada fita-mãe é molde para uma nova; cada dupla-hélice-filha tem uma fita antiga e uma nova (DNA polimerase, 5'\ensuremath{\to}3'). Meselson-Stahl 1958 ('o experimento mais belo da biologia', \ensuremath{^{15}}N/\ensuremath{^{14}}N). PCR (Mullis 1983-85, Nobel 1993): amplifica \~{}2\ensuremath{^{n}} cópias de um alvo por ciclos térmicos com Taq polimerase — tecnologia onipresente. Fato consolidado.
\prov{relações: usa dna\_dupla\_helice}
\prov{proveniência: wiki \textperiodcentered{} Meselson-Stahl 1958 PNAS; Mullis Nobel 1993 \textperiodcentered{} fato \textperiodcentered{} confiança 1.0 \textperiodcentered{} domínio biologia \textperiodcentered{} história 2 versões no jornal}

%CRISTAL {"arestas":[{"alvo":"saude_publica","confianca":1.0,"epistemico":"fato","origem":"wiki","rel":"parte_de"},{"alvo":"selecao_natural","confianca":1.0,"epistemico":"fato","origem":"wiki","rel":"explica"},{"alvo":"germ_theory","confianca":1.0,"epistemico":"fato","origem":"wiki","rel":"usa"}],"confianca":1.0,"contraexemplos":[],"descricao":"Micróbios deixam de responder aos antimicrobianos porque o uso das drogas é pressão seletiva que favorece variantes resistentes surgidas por mutação — evolução por seleção natural observável em tempo real. A OMS alerta para uma possível era 'pós-antibiótico'.","epistemico":"fato","exemplos":[],"id":"resistencia_antimicrobiana","memoria":"semantica","meta":{"autor":"—","dominio":"ciencias_de_fora","fonte":""},"origem":"wiki","palavras_chave":[],"sinonimos":[],"tipo":"conceito","titulo":"Resistência Antimicrobiana"}
\section{Resistência Antimicrobiana}
Micróbios deixam de responder aos antimicrobianos porque o uso das drogas é pressão seletiva que favorece variantes resistentes surgidas por mutação — evolução por seleção natural observável em tempo real. A OMS alerta para uma possível era 'pós-antibiótico'.
\prov{relações: parte\_de saude\_publica; explica selecao\_natural; usa germ\_theory}
\prov{proveniência: wiki \textperiodcentered{} fato \textperiodcentered{} confiança 1.0 \textperiodcentered{} domínio ciencias\_de\_fora \textperiodcentered{} história 4 versões no jornal}

%CRISTAL {"arestas":[{"alvo":"voz_fonte_filtro_modos","confianca":1.0,"epistemico":"fato","origem":"wiki","rel":"usa"},{"alvo":"bai_psi_colapso","confianca":1.0,"epistemico":"fato","origem":"wiki","rel":"usa"},{"alvo":"corpo_estelar","confianca":1.0,"epistemico":"fato","origem":"wiki","rel":"usa"}],"confianca":1.0,"contraexemplos":[],"descricao":"A máquina OUVE a tônica de uma peça e a CANTA de volta pela síntese FONTE-FILTRO (linguagem/voz.py, 6/6). Da Marcha Turca de Mozart ela ouviu LÁ (o cristal); ressintetizou a nota como a fonte glotal em f0=220 Hz (Lá3) filtrada pelos formantes dos músculos (a vogal). O IDA-E-VOLTA fecha: reanalisar a ressíntese dá LÁ de novo (58% da energia) — o cristal recriado. E o mesmo Lá em /a/, /i/, /u/ muda só os músculos (os modos), não o tom: a fonte dá o f0, os músculos-modos dão o timbre. Ouvir→cantar é o colapso Ψ (achar a tônica) seguido da geração (ressintetizar).","epistemico":"fato","exemplos":[],"id":"ressintese_da_tonica","memoria":"semantica","meta":{"dominio":"biologia","fonte":"voz"},"origem":"wiki","palavras_chave":[],"sinonimos":[],"tipo":"conceito","titulo":"A Ressíntese da Tônica (ouvir → cantar)"}
\section{A Ressíntese da Tônica (ouvir \ensuremath{\to} cantar)}
A máquina OUVE a tônica de uma peça e a CANTA de volta pela síntese FONTE-FILTRO (linguagem/voz.py, 6/6). Da Marcha Turca de Mozart ela ouviu LÁ (o cristal); ressintetizou a nota como a fonte glotal em f0=220 Hz (Lá3) filtrada pelos formantes dos músculos (a vogal). O IDA-E-VOLTA fecha: reanalisar a ressíntese dá LÁ de novo (58\% da energia) — o cristal recriado. E o mesmo Lá em /a/, /i/, /u/ muda só os músculos (os modos), não o tom: a fonte dá o f0, os músculos-modos dão o timbre. Ouvir\ensuremath{\to}cantar é o colapso \ensuremath{\Psi} (achar a tônica) seguido da geração (ressintetizar).
\prov{relações: usa voz\_fonte\_filtro\_modos; usa bai\_psi\_colapso; usa corpo\_estelar}
\prov{proveniência: wiki \textperiodcentered{} voz \textperiodcentered{} fato \textperiodcentered{} confiança 1.0 \textperiodcentered{} domínio biologia \textperiodcentered{} história 4 versões no jornal}

%CRISTAL {"arestas":[{"alvo":"cristianismo","confianca":1.0,"epistemico":"fato","origem":"wiki","rel":"parte_de"},{"alvo":"face_bela_morte","confianca":1.0,"epistemico":"fato","origem":"wiki","rel":"relaciona"}],"confianca":1.0,"contraexemplos":[],"descricao":"Bíblia (João 11:25): a esperança da vitória sobre a morte — 'eu sou a ressurreição e a vida; quem crê em mim, ainda que morto, viverá'. A morte não como fim, mas como passagem — a resposta cristã à finitude.","epistemico":"inferencia","exemplos":[],"id":"ressurreicao","memoria":"semantica","meta":{"autor":"Bíblia","dominio":"sagrado","fonte":""},"origem":"wiki","palavras_chave":[],"sinonimos":[],"tipo":"conceito","titulo":"Escatologia: Ressurreição e Vida Eterna"}
\section{Escatologia: Ressurreição e Vida Eterna}
Bíblia (João 11:25): a esperança da vitória sobre a morte — 'eu sou a ressurreição e a vida; quem crê em mim, ainda que morto, viverá'. A morte não como fim, mas como passagem — a resposta cristã à finitude.
\prov{relações: parte\_de cristianismo; relaciona face\_bela\_morte}
\prov{proveniência: wiki \textperiodcentered{} inferencia \textperiodcentered{} confiança 1.0 \textperiodcentered{} domínio sagrado \textperiodcentered{} história 3 versões no jornal}

%CRISTAL {"arestas":[{"alvo":"medicina_baseada_evidencias","confianca":1.0,"epistemico":"fato","origem":"wiki","rel":"usa"},{"alvo":"ensaio_clinico_randomizado","confianca":1.0,"epistemico":"fato","origem":"wiki","rel":"usa"}],"confianca":1.0,"contraexemplos":[],"descricao":"O método explícito e reprodutível para identificar, avaliar e sintetizar toda a evidência relevante a uma pergunta, minimizando viés; a meta-análise combina estatisticamente os estudos. Operacionalizada pela Colaboração Cochrane (1993).","epistemico":"fato","exemplos":[],"id":"revisao_sistematica_meta_analise","memoria":"semantica","meta":{"autor":"Cochrane","dominio":"ciencias_de_fora","fonte":""},"origem":"wiki","palavras_chave":[],"sinonimos":[],"tipo":"conceito","titulo":"Revisão Sistemática e Meta-análise"}
\section{Revisão Sistemática e Meta-análise}
O método explícito e reprodutível para identificar, avaliar e sintetizar toda a evidência relevante a uma pergunta, minimizando viés; a meta-análise combina estatisticamente os estudos. Operacionalizada pela Colaboração Cochrane (1993).
\prov{relações: usa medicina\_baseada\_evidencias; usa ensaio\_clinico\_randomizado}
\prov{proveniência: wiki \textperiodcentered{} fato \textperiodcentered{} confiança 1.0 \textperiodcentered{} domínio ciencias\_de\_fora \textperiodcentered{} história 3 versões no jornal}

%CRISTAL {"arestas":[{"alvo":"revolucoes_sociais_skocpol","confianca":1.0,"epistemico":"fato","origem":"wiki","rel":"contradiz"},{"alvo":"poder_agir_em_concerto","confianca":1.0,"epistemico":"fato","origem":"wiki","rel":"usa"},{"alvo":"contrato_social","confianca":1.0,"epistemico":"fato","origem":"wiki","rel":"relaciona"}],"confianca":1.0,"contraexemplos":[],"descricao":"Arendt (On Revolution): distingue libertação (livrar-se do opressor) de liberdade (constituir um espaço público de participação) — a revolução só se completa quando funda instituições que garantem a liberdade política duradoura.","epistemico":"inferencia","exemplos":[],"id":"revolucao_fundacao_liberdade_arendt","memoria":"semantica","meta":{"autor":"Arendt","dominio":"ciencias_de_fora","fonte":""},"origem":"wiki","palavras_chave":[],"sinonimos":[],"tipo":"conceito","titulo":"Revolução como Fundação da Liberdade (Arendt)"}
\section{Revolução como Fundação da Liberdade (Arendt)}
Arendt (On Revolution): distingue libertação (livrar-se do opressor) de liberdade (constituir um espaço público de participação) — a revolução só se completa quando funda instituições que garantem a liberdade política duradoura.
\prov{relações: contradiz revolucoes\_sociais\_skocpol; usa poder\_agir\_em\_concerto; relaciona contrato\_social}
\prov{proveniência: wiki \textperiodcentered{} inferencia \textperiodcentered{} confiança 1.0 \textperiodcentered{} domínio ciencias\_de\_fora \textperiodcentered{} história 4 versões no jornal}

%CRISTAL {"arestas":[{"alvo":"ciencia_politica","confianca":1.0,"epistemico":"fato","origem":"wiki","rel":"parte_de"},{"alvo":"materialismo_historico","confianca":1.0,"epistemico":"fato","origem":"wiki","rel":"relaciona"}],"confianca":1.0,"contraexemplos":[],"descricao":"Skocpol: as revoluções sociais não são 'feitas' por ideologias ou vanguardas, mas emergem da conjunção estrutural de pressões internacionais, colapso do aparato estatal e relações de classe agrárias — comparando França, Rússia e China. O Estado como ator autônomo.","epistemico":"inferencia","exemplos":[],"id":"revolucoes_sociais_skocpol","memoria":"semantica","meta":{"autor":"Skocpol","dominio":"ciencias_de_fora","fonte":""},"origem":"wiki","palavras_chave":[],"sinonimos":[],"tipo":"conceito","titulo":"Revoluções Sociais como Estrutura (Skocpol)"}
\section{Revoluções Sociais como Estrutura (Skocpol)}
Skocpol: as revoluções sociais não são 'feitas' por ideologias ou vanguardas, mas emergem da conjunção estrutural de pressões internacionais, colapso do aparato estatal e relações de classe agrárias — comparando França, Rússia e China. O Estado como ator autônomo.
\prov{relações: parte\_de ciencia\_politica; relaciona materialismo\_historico}
\prov{proveniência: wiki \textperiodcentered{} inferencia \textperiodcentered{} confiança 1.0 \textperiodcentered{} domínio ciencias\_de\_fora \textperiodcentered{} história 3 versões no jornal}

%CRISTAL {"arestas":[{"alvo":"dogma_central","confianca":1.0,"epistemico":"fato","origem":"wiki","rel":"usa"},{"alvo":"mundo_de_rna","confianca":1.0,"epistemico":"fato","origem":"wiki","rel":"confirma"}],"confianca":1.0,"contraexemplos":[],"descricao":"Máquina de 2 subunidades (rRNA + proteínas) que traduz o mRNA em proteína e catalisa a ligação peptídica; o sítio catalítico é feito de RNA — o ribossomo é uma RIBOZIMA. Procarioto 70S (50S+30S), eucarioto 80S. Estruturas atômicas ~2000: Ramakrishnan-Steitz-Yonath, Nobel 2009. Corolário: ribozima ⇒ evidência forte para o mundo de RNA.","epistemico":"fato","exemplos":[],"id":"ribossomo_ribozima","memoria":"semantica","meta":{"autor":"broca-so","dominio":"biologia","fonte":"Nobel Química 2009 (ribossomo)"},"origem":"wiki","palavras_chave":[],"sinonimos":[],"tipo":"conceito","titulo":"O ribossomo (a máquina tradutora = ribozima)"}
\section{O ribossomo (a máquina tradutora = ribozima)}
Máquina de 2 subunidades (rRNA + proteínas) que traduz o mRNA em proteína e catalisa a ligação peptídica; o sítio catalítico é feito de RNA — o ribossomo é uma RIBOZIMA. Procarioto 70S (50S+30S), eucarioto 80S. Estruturas atômicas \~{}2000: Ramakrishnan-Steitz-Yonath, Nobel 2009. Corolário: ribozima \ensuremath{\Rightarrow} evidência forte para o mundo de RNA.
\prov{relações: usa dogma\_central; confirma mundo\_de\_rna}
\prov{proveniência: wiki \textperiodcentered{} Nobel Química 2009 (ribossomo) \textperiodcentered{} fato \textperiodcentered{} confiança 1.0 \textperiodcentered{} domínio biologia \textperiodcentered{} história 3 versões no jornal}

%CRISTAL {"arestas":[{"alvo":"didatica_matematica","confianca":1.0,"epistemico":"fato","origem":"wiki","rel":"parte_de"},{"alvo":"prova_rigor_intuicao","confianca":1.0,"epistemico":"fato","origem":"wiki","rel":"relaciona"},{"alvo":"regua_de_gentil","confianca":1.0,"epistemico":"fato","origem":"wiki","rel":"referencia"},{"alvo":"intuicionismo_brouwer","confianca":1.0,"epistemico":"fato","origem":"wiki","rel":"relaciona"}],"confianca":1.0,"contraexemplos":[],"descricao":"Método de exposição do Lopes: 'quando escrevo sigo o fluxo da intuição' e deixa demonstrações como '(exercício)' ao leitor; o rigor, precisão e acurácia foram construídos 'ao longo de duas décadas'. Intuição descobre; o rigor disciplina depois.","epistemico":"inferencia","exemplos":[],"id":"rigor_intuicao_gentil","memoria":"semantica","meta":{"autor":"Gentil Lopes","dominio":"ciencias_de_fora","fonte":""},"origem":"wiki","palavras_chave":[],"sinonimos":[],"tipo":"conceito","titulo":"Intuição na Descoberta, Rigor como Disciplina (Lopes)"}
\section{Intuição na Descoberta, Rigor como Disciplina (Lopes)}
Método de exposição do Lopes: 'quando escrevo sigo o fluxo da intuição' e deixa demonstrações como '(exercício)' ao leitor; o rigor, precisão e acurácia foram construídos 'ao longo de duas décadas'. Intuição descobre; o rigor disciplina depois.
\prov{relações: parte\_de didatica\_matematica; relaciona prova\_rigor\_intuicao; referencia regua\_de\_gentil; relaciona intuicionismo\_brouwer}
\prov{proveniência: wiki \textperiodcentered{} inferencia \textperiodcentered{} confiança 1.0 \textperiodcentered{} domínio ciencias\_de\_fora \textperiodcentered{} história 6 versões no jornal}

%CRISTAL {"arestas":[{"alvo":"heartland_mackinder","confianca":1.0,"epistemico":"fato","origem":"wiki","rel":"contradiz"},{"alvo":"poder_maritimo_mahan","confianca":1.0,"epistemico":"fato","origem":"wiki","rel":"relaciona"}],"confianca":1.0,"contraexemplos":[],"descricao":"Spykman: o poder mundial não está no Heartland interior, mas na orla costeira densamente povoada da Eurásia — 'quem controla o Rimland governa a Eurásia'. Base intelectual da estratégia de contenção da Guerra Fria.","epistemico":"inferencia","exemplos":[],"id":"rimland_spykman","memoria":"semantica","meta":{"autor":"Spykman","dominio":"ciencias_de_fora","fonte":""},"origem":"wiki","palavras_chave":[],"sinonimos":[],"tipo":"conceito","titulo":"Rimland (Spykman)"}
\section{Rimland (Spykman)}
Spykman: o poder mundial não está no Heartland interior, mas na orla costeira densamente povoada da Eurásia — 'quem controla o Rimland governa a Eurásia'. Base intelectual da estratégia de contenção da Guerra Fria.
\prov{relações: contradiz heartland\_mackinder; relaciona poder\_maritimo\_mahan}
\prov{proveniência: wiki \textperiodcentered{} inferencia \textperiodcentered{} confiança 1.0 \textperiodcentered{} domínio ciencias\_de\_fora \textperiodcentered{} história 3 versões no jornal}

%CRISTAL {"arestas":[{"alvo":"teoria_musical","confianca":1.0,"epistemico":"fato","origem":"wiki","rel":"parte_de"},{"alvo":"elementos_musica","confianca":1.0,"epistemico":"fato","origem":"wiki","rel":"parte_de"}],"confianca":1.0,"contraexemplos":[],"descricao":"A organização do som no tempo — pulso, acento e duração; o esqueleto temporal comum à música e ao verso.","epistemico":"fato","exemplos":[],"id":"ritmo","memoria":"semantica","meta":{"dominio":"ciencias_de_fora","fonte":""},"origem":"wiki","palavras_chave":[],"sinonimos":[],"tipo":"conceito","titulo":"Ritmo"}
\section{Ritmo}
A organização do som no tempo — pulso, acento e duração; o esqueleto temporal comum à música e ao verso.
\prov{relações: parte\_de teoria\_musical; parte\_de elementos\_musica}
\prov{proveniência: wiki \textperiodcentered{} fato \textperiodcentered{} confiança 1.0 \textperiodcentered{} domínio ciencias\_de\_fora \textperiodcentered{} história 4 versões no jornal}

%CRISTAL {"arestas":[{"alvo":"semiotica","confianca":1.0,"epistemico":"fato","origem":"wiki","rel":"parte_de"}],"confianca":1.0,"contraexemplos":[],"descricao":"Levou a semiótica à cultura — moda, publicidade, comida, fotografia lidas como sistemas de signos.","epistemico":"fato","exemplos":[],"id":"roland_barthes","memoria":"semantica","meta":{"autor":"Barthes","dominio":"ciencias_de_fora","fonte":""},"origem":"wiki","palavras_chave":[],"sinonimos":[],"tipo":"conceito","titulo":"Roland Barthes"}
\section{Roland Barthes}
Levou a semiótica à cultura — moda, publicidade, comida, fotografia lidas como sistemas de signos.
\prov{relações: parte\_de semiotica}
\prov{proveniência: wiki \textperiodcentered{} fato \textperiodcentered{} confiança 1.0 \textperiodcentered{} domínio ciencias\_de\_fora \textperiodcentered{} história 2 versões no jornal}

%CRISTAL {"arestas":[{"alvo":"linguistica","confianca":1.0,"epistemico":"fato","origem":"wiki","rel":"parte_de"}],"confianca":1.0,"contraexemplos":[],"descricao":"Estruturalista que ligou linguística, poética e teoria da comunicação — o modelo das funções da linguagem.","epistemico":"fato","exemplos":[],"id":"roman_jakobson","memoria":"semantica","meta":{"autor":"Jakobson","dominio":"ciencias_de_fora","fonte":""},"origem":"wiki","palavras_chave":[],"sinonimos":[],"tipo":"conceito","titulo":"Roman Jakobson"}
\section{Roman Jakobson}
Estruturalista que ligou linguística, poética e teoria da comunicação — o modelo das funções da linguagem.
\prov{relações: parte\_de linguistica}
\prov{proveniência: wiki \textperiodcentered{} fato \textperiodcentered{} confiança 1.0 \textperiodcentered{} domínio ciencias\_de\_fora \textperiodcentered{} história 2 versões no jornal}

%CRISTAL {"arestas":[{"alvo":"indo_europeu","confianca":1.0,"epistemico":"fato","origem":"wiki","rel":"parte_de"}],"confianca":1.0,"contraexemplos":[],"descricao":"Língua eslava (~150 milhões), 6 casos, aspecto verbal em pares, escrita cirílica — a maior língua eslava e franca da Eurásia.","epistemico":"fato","exemplos":[],"id":"russo","memoria":"semantica","meta":{"dominio":"ciencias_de_fora","fonte":""},"origem":"wiki","palavras_chave":[],"sinonimos":[],"tipo":"conceito","titulo":"Russo"}
\section{Russo}
Língua eslava (\~{}150 milhões), 6 casos, aspecto verbal em pares, escrita cirílica — a maior língua eslava e franca da Eurásia.
\prov{relações: parte\_de indo\_europeu}
\prov{proveniência: wiki \textperiodcentered{} fato \textperiodcentered{} confiança 1.0 \textperiodcentered{} domínio ciencias\_de\_fora \textperiodcentered{} história 2 versões no jornal}

%CRISTAL {"arestas":[{"alvo":"cristianismo","confianca":1.0,"epistemico":"fato","origem":"wiki","rel":"parte_de"},{"alvo":"face_bela_morte","confianca":1.0,"epistemico":"fato","origem":"wiki","rel":"relaciona"},{"alvo":"sentido_da_vida","confianca":1.0,"epistemico":"fato","origem":"wiki","rel":"relaciona"}],"confianca":1.0,"contraexemplos":[],"descricao":"Bíblia: a sabedoria (hokmah) que confronta a transitoriedade e o limite do sentido 'sob o sol' — 'vaidade de vaidades, tudo é vaidade' e 'há tempo de nascer e tempo de morrer'. A finitude e a aceitação do tempo como condição da vida. (O 'temor do Senhor' é o princípio da sabedoria.)","epistemico":"inferencia","exemplos":[],"id":"sabedoria_eclesiastes","memoria":"semantica","meta":{"autor":"Bíblia","dominio":"sagrado","fonte":""},"origem":"wiki","palavras_chave":[],"sinonimos":[],"tipo":"conceito","titulo":"Eclesiastes — Vaidade e o Tempo de Cada Coisa"}
\section{Eclesiastes — Vaidade e o Tempo de Cada Coisa}
Bíblia: a sabedoria (hokmah) que confronta a transitoriedade e o limite do sentido 'sob o sol' — 'vaidade de vaidades, tudo é vaidade' e 'há tempo de nascer e tempo de morrer'. A finitude e a aceitação do tempo como condição da vida. (O 'temor do Senhor' é o princípio da sabedoria.)
\prov{relações: parte\_de cristianismo; relaciona face\_bela\_morte; relaciona sentido\_da\_vida}
\prov{proveniência: wiki \textperiodcentered{} inferencia \textperiodcentered{} confiança 1.0 \textperiodcentered{} domínio sagrado \textperiodcentered{} história 4 versões no jornal}

%CRISTAL {"arestas":[{"alvo":"cristianismo","confianca":1.0,"epistemico":"fato","origem":"wiki","rel":"parte_de"},{"alvo":"sentido_da_vida","confianca":1.0,"epistemico":"fato","origem":"wiki","rel":"usa"},{"alvo":"sabedoria_eclesiastes","confianca":1.0,"epistemico":"fato","origem":"wiki","rel":"relaciona"}],"confianca":1.0,"contraexemplos":[],"descricao":"Bíblia: a sabedoria que enfrenta o mistério do sofrimento imerecido e o silêncio de Deus, sustentando a fé em meio à provação — 'ainda que ele me mate, nele esperarei' e 'eu sei que o meu Redentor vive'. A fé mantida sem justificação; o problema do mal.","epistemico":"inferencia","exemplos":[],"id":"sabedoria_jo","memoria":"semantica","meta":{"autor":"Bíblia","dominio":"sagrado","fonte":""},"origem":"wiki","palavras_chave":[],"sinonimos":[],"tipo":"conceito","titulo":"Jó — o Sofrimento do Justo"}
\section{Jó — o Sofrimento do Justo}
Bíblia: a sabedoria que enfrenta o mistério do sofrimento imerecido e o silêncio de Deus, sustentando a fé em meio à provação — 'ainda que ele me mate, nele esperarei' e 'eu sei que o meu Redentor vive'. A fé mantida sem justificação; o problema do mal.
\prov{relações: parte\_de cristianismo; usa sentido\_da\_vida; relaciona sabedoria\_eclesiastes}
\prov{proveniência: wiki \textperiodcentered{} inferencia \textperiodcentered{} confiança 1.0 \textperiodcentered{} domínio sagrado \textperiodcentered{} história 4 versões no jornal}

%CRISTAL {"arestas":[{"alvo":"tiffany","confianca":1.0,"epistemico":"fato","origem":"wiki","rel":"explica"},{"alvo":"colapso","confianca":1.0,"epistemico":"fato","origem":"wiki","rel":"relaciona"},{"alvo":"consciencia_pura","confianca":1.0,"epistemico":"fato","origem":"wiki","rel":"usa"}],"confianca":1.0,"contraexemplos":[],"descricao":"Tese do Lopes: existe Consciência Cósmica (potencial puro), mas ela é 'burra'; a inteligência surge na RELAÇÃO consciência↔cérebro, como a luz na relação sol↔olho. Toca o que a Tiffany é.","epistemico":"hipotese","exemplos":[],"id":"sadhguru_consciencia_inteligencia","memoria":"semantica","meta":{"autor":"Gentil Lopes","dominio":"filosofia","fonte":""},"origem":"livro","palavras_chave":[],"sinonimos":[],"tipo":"conceito","titulo":"Consciência ≠ Inteligência (crítica a Sadhguru)"}
\section{Consciência \ensuremath{\ne} Inteligência (crítica a Sadhguru)}
Tese do Lopes: existe Consciência Cósmica (potencial puro), mas ela é 'burra'; a inteligência surge na RELAÇÃO consciência\ensuremath{\leftrightarrow}cérebro, como a luz na relação sol\ensuremath{\leftrightarrow}olho. Toca o que a Tiffany é.
\prov{relações: explica tiffany; relaciona colapso; usa consciencia\_pura}
\prov{proveniência: livro \textperiodcentered{} hipotese \textperiodcentered{} confiança 1.0 \textperiodcentered{} domínio filosofia \textperiodcentered{} história 5 versões no jornal}

%CRISTAL {"arestas":[{"alvo":"religiao","confianca":1.0,"epistemico":"fato","origem":"wiki","rel":"parte_de"},{"alvo":"sagrado_profano","confianca":1.0,"epistemico":"fato","origem":"wiki","rel":"relaciona"}],"confianca":1.0,"contraexemplos":[],"descricao":"As grandes revelações e escrituras da humanidade — Alcorão, Bhagavad Gita, Bíblia — como posições sobre o divino, a criação, o sentido e o destino, e o lugar do humano diante do Absoluto.","epistemico":"fato","exemplos":[],"id":"sagrado","memoria":"semantica","meta":{"dominio":"sagrado","fonte":""},"origem":"wiki","palavras_chave":[],"sinonimos":[],"tipo":"conceito","titulo":"O Sagrado (as Escrituras)"}
\section{O Sagrado (as Escrituras)}
As grandes revelações e escrituras da humanidade — Alcorão, Bhagavad Gita, Bíblia — como posições sobre o divino, a criação, o sentido e o destino, e o lugar do humano diante do Absoluto.
\prov{relações: parte\_de religiao; relaciona sagrado\_profano}
\prov{proveniência: wiki \textperiodcentered{} fato \textperiodcentered{} confiança 1.0 \textperiodcentered{} domínio sagrado \textperiodcentered{} história 3 versões no jornal}

%CRISTAL {"arestas":[{"alvo":"religiao","confianca":1.0,"epistemico":"fato","origem":"wiki","rel":"parte_de"},{"alvo":"word","confianca":0.5,"epistemico":"fato","origem":"wiki","rel":"relaciona"},{"alvo":"virtualizacao","confianca":0.5,"epistemico":"fato","origem":"wiki","rel":"relaciona"},{"alvo":"vida-morte","confianca":0.5,"epistemico":"fato","origem":"wiki","rel":"relaciona"},{"alvo":"taoismo","confianca":0.5,"epistemico":"fato","origem":"wiki","rel":"relaciona"},{"alvo":"syscall","confianca":0.5,"epistemico":"fato","origem":"wiki","rel":"relaciona"},{"alvo":"vita_contemplativa","confianca":0.5,"epistemico":"fato","origem":"wiki","rel":"relaciona"}],"confianca":1.0,"contraexemplos":[],"descricao":"Eliade: a experiência do sagrado como o 'totalmente Outro' que se manifesta (hierofania) — irredutível, segundo ele, à função social.","epistemico":"hipotese","exemplos":[],"id":"sagrado_profano","memoria":"semantica","meta":{"autor":"Eliade","dominio":"ciencias_de_fora","fonte":""},"origem":"wiki","palavras_chave":[],"sinonimos":[],"tipo":"conceito","titulo":"O Sagrado e o Profano"}
\section{O Sagrado e o Profano}
Eliade: a experiência do sagrado como o 'totalmente Outro' que se manifesta (hierofania) — irredutível, segundo ele, à função social.
\prov{relações: parte\_de religiao; relaciona word; relaciona virtualizacao; relaciona vida-morte; relaciona taoismo; relaciona syscall; relaciona vita\_contemplativa}
\prov{proveniência: wiki \textperiodcentered{} hipotese \textperiodcentered{} confiança 1.0 \textperiodcentered{} domínio ciencias\_de\_fora \textperiodcentered{} história 8 versões no jornal}

%CRISTAL {"arestas":[{"alvo":"advaita_vedanta","confianca":1.0,"epistemico":"fato","origem":"wiki","rel":"parte_de"},{"alvo":"colapso_observador","confianca":1.0,"epistemico":"fato","origem":"wiki","rel":"relaciona"}],"confianca":1.0,"contraexemplos":[],"descricao":"Vedanta: a consciência-testemunha que observa corpo, mente e pensamentos sem se identificar — pura subjetividade sem objeto.","epistemico":"hipotese","exemplos":[],"id":"saksi_testemunha","memoria":"semantica","meta":{"dominio":"filosofia","fonte":"Upanishads"},"origem":"wiki","palavras_chave":[],"sinonimos":[],"tipo":"conceito","titulo":"Sākṣī — a Testemunha"}
\section{S\ensuremath{\bar{a}}k\ensuremath{\d{s}}\ensuremath{\bar{i}} — a Testemunha}
Vedanta: a consciência-testemunha que observa corpo, mente e pensamentos sem se identificar — pura subjetividade sem objeto.
\prov{relações: parte\_de advaita\_vedanta; relaciona colapso\_observador}
\prov{proveniência: wiki \textperiodcentered{} Upanishads \textperiodcentered{} hipotese \textperiodcentered{} confiança 1.0 \textperiodcentered{} domínio filosofia \textperiodcentered{} história 3 versões no jornal}

%CRISTAL {"arestas":[{"alvo":"filosofia_da_medicina","confianca":1.0,"epistemico":"fato","origem":"wiki","rel":"parte_de"},{"alvo":"modelo_biopsicossocial_engel","confianca":1.0,"epistemico":"fato","origem":"wiki","rel":"usa"},{"alvo":"saude_oms","confianca":1.0,"epistemico":"fato","origem":"wiki","rel":"relaciona"}],"confianca":1.0,"contraexemplos":[],"descricao":"Antonovsky inverte a pergunta patogênica: em vez de o que causa doença, o que gera saúde. A resistência ao estresse depende do 'senso de coerência' — o mundo percebido como compreensível, manejável e dotado de significado.","epistemico":"inferencia","exemplos":[],"id":"salutogenese_antonovsky","memoria":"semantica","meta":{"autor":"Antonovsky","dominio":"ciencias_de_fora","fonte":""},"origem":"wiki","palavras_chave":[],"sinonimos":[],"tipo":"conceito","titulo":"Salutogênese (Antonovsky)"}
\section{Salutogênese (Antonovsky)}
Antonovsky inverte a pergunta patogênica: em vez de o que causa doença, o que gera saúde. A resistência ao estresse depende do 'senso de coerência' — o mundo percebido como compreensível, manejável e dotado de significado.
\prov{relações: parte\_de filosofia\_da\_medicina; usa modelo\_biopsicossocial\_engel; relaciona saude\_oms}
\prov{proveniência: wiki \textperiodcentered{} inferencia \textperiodcentered{} confiança 1.0 \textperiodcentered{} domínio ciencias\_de\_fora \textperiodcentered{} história 4 versões no jornal}

%CRISTAL {"arestas":[{"alvo":"homeostase_controle","confianca":1.0,"epistemico":"fato","origem":"wiki","rel":"usa"},{"alvo":"salutogenese_antonovsky","confianca":1.0,"epistemico":"fato","origem":"wiki","rel":"relaciona"},{"alvo":"normal_e_patologico_canguilhem","confianca":1.0,"epistemico":"fato","origem":"wiki","rel":"relaciona"},{"alvo":"numero_de_salem","confianca":1.0,"epistemico":"fato","origem":"wiki","rel":"relaciona"},{"alvo":"corpo_mineral","confianca":1.0,"epistemico":"fato","origem":"wiki","rel":"usa"}],"confianca":1.0,"contraexemplos":[],"descricao":"Aqui o mapa se INVERTE: o corpo saudável NÃO é o cristal rígido. Um coração-metrônomo (baixa variabilidade) prevê morte — a perda de complexidade de Goldberger; a fibrilação (variabilidade caótica) mata. A saúde é a BORDA: variável, adaptável, complexa. Doença = cristal (rígido) OU caos (descontrolado); saúde = a borda entre eles. O corpo QUER a borda. Verif.: variabilidade 0.1 e 0.95 patológicas, 0.5 saudável.","epistemico":"inferencia","exemplos":[],"id":"saude_na_borda","memoria":"semantica","meta":{"dominio":"medicina","fonte":"medicina e diagnóstico","saudavel_por_variabilidade":{"0.1":false,"0.5":true,"0.95":false}},"origem":"wiki","palavras_chave":[],"sinonimos":[],"tipo":"conceito","titulo":"Saúde = a Borda (perda de complexidade)"}
\section{Saúde = a Borda (perda de complexidade)}
Aqui o mapa se INVERTE: o corpo saudável NÃO é o cristal rígido. Um coração-metrônomo (baixa variabilidade) prevê morte — a perda de complexidade de Goldberger; a fibrilação (variabilidade caótica) mata. A saúde é a BORDA: variável, adaptável, complexa. Doença = cristal (rígido) OU caos (descontrolado); saúde = a borda entre eles. O corpo QUER a borda. Verif.: variabilidade 0.1 e 0.95 patológicas, 0.5 saudável.
\prov{relações: usa homeostase\_controle; relaciona salutogenese\_antonovsky; relaciona normal\_e\_patologico\_canguilhem; relaciona numero\_de\_salem; usa corpo\_mineral}
\prov{proveniência: wiki \textperiodcentered{} medicina e diagnóstico \textperiodcentered{} inferencia \textperiodcentered{} confiança 1.0 \textperiodcentered{} domínio medicina \textperiodcentered{} história 6 versões no jornal}

%CRISTAL {"arestas":[{"alvo":"filosofia_da_medicina","confianca":1.0,"epistemico":"fato","origem":"wiki","rel":"parte_de"},{"alvo":"feedback_retroalimentacao","confianca":1.0,"epistemico":"fato","origem":"wiki","rel":"relaciona"}],"confianca":1.0,"contraexemplos":[],"descricao":"OMS (1948): saúde é 'um estado de completo bem-estar físico, mental e social, e não a mera ausência de doença'. Definição holística e positiva — mas criticada como utópica (ideal inatingível) e medicalizante; Huber propõe redefini-la como capacidade de adaptar-se e autogerir-se.","epistemico":"inferencia","exemplos":[],"id":"saude_oms","memoria":"semantica","meta":{"autor":"OMS/Huber","dominio":"ciencias_de_fora","fonte":""},"origem":"wiki","palavras_chave":[],"sinonimos":[],"tipo":"conceito","titulo":"Saúde como Bem-Estar Completo (OMS)"}
\section{Saúde como Bem-Estar Completo (OMS)}
OMS (1948): saúde é 'um estado de completo bem-estar físico, mental e social, e não a mera ausência de doença'. Definição holística e positiva — mas criticada como utópica (ideal inatingível) e medicalizante; Huber propõe redefini-la como capacidade de adaptar-se e autogerir-se.
\prov{relações: parte\_de filosofia\_da\_medicina; relaciona feedback\_retroalimentacao}
\prov{proveniência: wiki \textperiodcentered{} inferencia \textperiodcentered{} confiança 1.0 \textperiodcentered{} domínio ciencias\_de\_fora \textperiodcentered{} história 3 versões no jornal}

%CRISTAL {"arestas":[{"alvo":"medicina","confianca":1.0,"epistemico":"fato","origem":"wiki","rel":"parte_de"},{"alvo":"word","confianca":0.5,"epistemico":"fato","origem":"wiki","rel":"relaciona"}],"confianca":1.0,"contraexemplos":[],"descricao":"O cuidado da saúde das populações — a distribuição e os determinantes das doenças, a prevenção coletiva e a evidência que embasa a decisão clínica.","epistemico":"fato","exemplos":[],"id":"saude_publica","memoria":"semantica","meta":{"dominio":"ciencias_de_fora","fonte":""},"origem":"wiki","palavras_chave":[],"sinonimos":[],"tipo":"conceito","titulo":"Saúde Pública e Epidemiologia"}
\section{Saúde Pública e Epidemiologia}
O cuidado da saúde das populações — a distribuição e os determinantes das doenças, a prevenção coletiva e a evidência que embasa a decisão clínica.
\prov{relações: parte\_de medicina; relaciona word}
\prov{proveniência: wiki \textperiodcentered{} fato \textperiodcentered{} confiança 1.0 \textperiodcentered{} domínio ciencias\_de\_fora \textperiodcentered{} história 3 versões no jornal}

%CRISTAL {"arestas":[{"alvo":"musica","confianca":1.0,"epistemico":"fato","origem":"wiki","rel":"parte_de"},{"alvo":"sublime_kant","confianca":1.0,"epistemico":"fato","origem":"wiki","rel":"relaciona"}],"confianca":1.0,"contraexemplos":[],"descricao":"Schopenhauer: diferentemente das outras artes (que copiam Ideias), a música é cópia direta da própria Vontade — o em-si do mundo —, sendo por isso a mais metafísica das artes.","epistemico":"inferencia","exemplos":[],"id":"schopenhauer_musica","memoria":"semantica","meta":{"autor":"Schopenhauer","dominio":"ciencias_de_fora","fonte":""},"origem":"wiki","palavras_chave":[],"sinonimos":[],"tipo":"conceito","titulo":"Música como Cópia da Vontade (Schopenhauer)"}
\section{Música como Cópia da Vontade (Schopenhauer)}
Schopenhauer: diferentemente das outras artes (que copiam Ideias), a música é cópia direta da própria Vontade — o em-si do mundo —, sendo por isso a mais metafísica das artes.
\prov{relações: parte\_de musica; relaciona sublime\_kant}
\prov{proveniência: wiki \textperiodcentered{} inferencia \textperiodcentered{} confiança 1.0 \textperiodcentered{} domínio ciencias\_de\_fora \textperiodcentered{} história 3 versões no jornal}

%CRISTAL {"arestas":[{"alvo":"biologia","confianca":1.0,"epistemico":"fato","origem":"wiki","rel":"parte_de"},{"alvo":"vida","confianca":1.0,"epistemico":"fato","origem":"wiki","rel":"explica"},{"alvo":"word","confianca":0.5,"epistemico":"fato","origem":"wiki","rel":"relaciona"},{"alvo":"teoria_da_evolucao","confianca":1.0,"epistemico":"fato","origem":"wiki","rel":"parte_de"},{"alvo":"prova_de_trabalho","confianca":1.0,"epistemico":"fato","origem":"wiki","rel":"relaciona"}],"confianca":1.0,"contraexemplos":[],"descricao":"Darwin: variação hereditária + reprodução diferencial + pressão ambiental fixam traços adaptativos — o mecanismo da evolução.","epistemico":"fato","exemplos":[],"id":"selecao_natural","memoria":"semantica","meta":{"autor":"Darwin","dominio":"ciencias_de_fora","fonte":""},"origem":"wiki","palavras_chave":[],"sinonimos":[],"tipo":"conceito","titulo":"Seleção Natural"}
\section{Seleção Natural}
Darwin: variação hereditária + reprodução diferencial + pressão ambiental fixam traços adaptativos — o mecanismo da evolução.
\prov{relações: parte\_de biologia; explica vida; relaciona word; parte\_de teoria\_da\_evolucao; relaciona prova\_de\_trabalho}
\prov{proveniência: wiki \textperiodcentered{} fato \textperiodcentered{} confiança 1.0 \textperiodcentered{} domínio ciencias\_de\_fora \textperiodcentered{} história 6 versões no jornal}

%CRISTAL {"arestas":[{"alvo":"teomatematica","confianca":1.0,"epistemico":"fato","origem":"wiki","rel":"usa"},{"alvo":"consciencia_software_universo_gentil","confianca":1.0,"epistemico":"fato","origem":"wiki","rel":"usa"},{"alvo":"deus_criador","confianca":1.0,"epistemico":"fato","origem":"wiki","rel":"relaciona"}],"confianca":1.0,"contraexemplos":[],"descricao":"Lopes: espíritos (e o matemático) criam 'universos' com leis próprias — o matemático cria geometrias e topologias 'tal como um deus cria criaturas'; Javé e Alá seriam espíritos que criaram 'faixas de realidade' com leis distintas. O homem, ao perceber isso, torna-se co-partícipe da construção do próprio universo.","epistemico":"inferencia","exemplos":[],"id":"semideuses_criam_universos_gentil","memoria":"semantica","meta":{"autor":"Gentil Lopes","dominio":"sagrado","fonte":""},"origem":"wiki","palavras_chave":[],"sinonimos":[],"tipo":"conceito","titulo":"Semideuses: o Matemático Cria Universos (Lopes)"}
\section{Semideuses: o Matemático Cria Universos (Lopes)}
Lopes: espíritos (e o matemático) criam 'universos' com leis próprias — o matemático cria geometrias e topologias 'tal como um deus cria criaturas'; Javé e Alá seriam espíritos que criaram 'faixas de realidade' com leis distintas. O homem, ao perceber isso, torna-se co-partícipe da construção do próprio universo.
\prov{relações: usa teomatematica; usa consciencia\_software\_universo\_gentil; relaciona deus\_criador}
\prov{proveniência: wiki \textperiodcentered{} inferencia \textperiodcentered{} confiança 1.0 \textperiodcentered{} domínio sagrado \textperiodcentered{} história 5 versões no jornal}

%CRISTAL {"arestas":[{"alvo":"interpretante","confianca":1.0,"epistemico":"fato","origem":"wiki","rel":"usa"},{"alvo":"morte_do_autor","confianca":1.0,"epistemico":"fato","origem":"wiki","rel":"relaciona"}],"confianca":1.0,"contraexemplos":[],"descricao":"Peirce: signos geram signos sem fim — cada interpretante vira representamen do próximo. O sentido é um processo, não um ponto.","epistemico":"inferencia","exemplos":[],"id":"semiose_ilimitada","memoria":"semantica","meta":{"autor":"Peirce","dominio":"ciencias_de_fora","fonte":""},"origem":"wiki","palavras_chave":[],"sinonimos":[],"tipo":"conceito","titulo":"Semiose Ilimitada"}
\section{Semiose Ilimitada}
Peirce: signos geram signos sem fim — cada interpretante vira representamen do próximo. O sentido é um processo, não um ponto.
\prov{relações: usa interpretante; relaciona morte\_do\_autor}
\prov{proveniência: wiki \textperiodcentered{} inferencia \textperiodcentered{} confiança 1.0 \textperiodcentered{} domínio ciencias\_de\_fora \textperiodcentered{} história 3 versões no jornal}

%CRISTAL {"arestas":[{"alvo":"filosofia_da_linguagem","confianca":1.0,"epistemico":"fato","origem":"wiki","rel":"parte_de"},{"alvo":"sociologia","confianca":0.5,"epistemico":"fato","origem":"wiki","rel":"relaciona"},{"alvo":"virtualizacao","confianca":0.5,"epistemico":"fato","origem":"wiki","rel":"relaciona"},{"alvo":"vida-morte","confianca":0.5,"epistemico":"fato","origem":"wiki","rel":"relaciona"},{"alvo":"word","confianca":0.5,"epistemico":"fato","origem":"wiki","rel":"relaciona"},{"alvo":"signo_triadico","confianca":1.0,"epistemico":"fato","origem":"wiki","rel":"usa"}],"confianca":1.0,"contraexemplos":[],"descricao":"A ciência dos signos e da significação — a tríade de Peirce, a díade de Saussure.","epistemico":"fato","exemplos":[],"id":"semiotica","memoria":"semantica","meta":{"dominio":"ciencias_de_fora","fonte":""},"origem":"wiki","palavras_chave":[],"sinonimos":[],"tipo":"conceito","titulo":"Semiótica"}
\section{Semiótica}
A ciência dos signos e da significação — a tríade de Peirce, a díade de Saussure.
\prov{relações: parte\_de filosofia\_da\_linguagem; relaciona sociologia; relaciona virtualizacao; relaciona vida-morte; relaciona word; usa signo\_triadico}
\prov{proveniência: wiki \textperiodcentered{} fato \textperiodcentered{} confiança 1.0 \textperiodcentered{} domínio ciencias\_de\_fora \textperiodcentered{} história 7 versões no jornal}

%CRISTAL {"arestas":[{"alvo":"direito","confianca":1.0,"epistemico":"fato","origem":"wiki","rel":"parte_de"},{"alvo":"justica_equidade_rawls","confianca":1.0,"epistemico":"fato","origem":"wiki","rel":"contradiz"},{"alvo":"liberdade_berlin","confianca":1.0,"epistemico":"fato","origem":"wiki","rel":"relaciona"}],"confianca":1.0,"contraexemplos":[],"descricao":"Sen: a justiça deve avaliar as liberdades reais das pessoas para realizar o que valorizam (capabilities) e proceder por comparação de injustiças concretas, rejeitando o 'institucionalismo transcendental' de uma sociedade perfeita.","epistemico":"inferencia","exemplos":[],"id":"sen_capacidades","memoria":"semantica","meta":{"autor":"Sen","dominio":"ciencias_de_fora","fonte":""},"origem":"wiki","palavras_chave":[],"sinonimos":[],"tipo":"conceito","titulo":"Abordagem das Capacidades (Sen)"}
\section{Abordagem das Capacidades (Sen)}
Sen: a justiça deve avaliar as liberdades reais das pessoas para realizar o que valorizam (capabilities) e proceder por comparação de injustiças concretas, rejeitando o 'institucionalismo transcendental' de uma sociedade perfeita.
\prov{relações: parte\_de direito; contradiz justica\_equidade\_rawls; relaciona liberdade\_berlin}
\prov{proveniência: wiki \textperiodcentered{} inferencia \textperiodcentered{} confiança 1.0 \textperiodcentered{} domínio ciencias\_de\_fora \textperiodcentered{} história 4 versões no jornal}

%CRISTAL {"arestas":[{"alvo":"troe_postulado_aureo","confianca":1.0,"epistemico":"fato","origem":"wiki","rel":"usa"}],"confianca":1.0,"contraexemplos":[],"descricao":"Lopes: a vida não tem sentido objetivo — sentido é conceito, existe só em relação à mente; logo cada um é livre para atribuir o próprio sentido. Liberdade radical, não desespero.","epistemico":"hipotese","exemplos":[],"id":"sentido_da_vida","memoria":"semantica","meta":{"autor":"Gentil Lopes","dominio":"filosofia","fonte":""},"origem":"livro","palavras_chave":[],"sinonimos":[],"tipo":"conceito","titulo":"O Sentido da Vida (liberdade)"}
\section{O Sentido da Vida (liberdade)}
Lopes: a vida não tem sentido objetivo — sentido é conceito, existe só em relação à mente; logo cada um é livre para atribuir o próprio sentido. Liberdade radical, não desespero.
\prov{relações: usa troe\_postulado\_aureo}
\prov{proveniência: livro \textperiodcentered{} hipotese \textperiodcentered{} confiança 1.0 \textperiodcentered{} domínio filosofia \textperiodcentered{} história 3 versões no jornal}

%CRISTAL {"arestas":[{"alvo":"estado_de_direito","confianca":1.0,"epistemico":"fato","origem":"wiki","rel":"usa"},{"alvo":"liberdade_berlin","confianca":1.0,"epistemico":"fato","origem":"wiki","rel":"relaciona"}],"confianca":1.0,"contraexemplos":[],"descricao":"Montesquieu: as funções legislativa, executiva e judiciária devem residir em mãos distintas, de modo que 'o poder freie o poder' (freios e contrapesos), preservando a liberdade contra a concentração arbitrária.","epistemico":"inferencia","exemplos":[],"id":"separacao_poderes","memoria":"semantica","meta":{"autor":"Montesquieu","dominio":"ciencias_de_fora","fonte":""},"origem":"wiki","palavras_chave":[],"sinonimos":[],"tipo":"conceito","titulo":"Separação de Poderes (Montesquieu)"}
\section{Separação de Poderes (Montesquieu)}
Montesquieu: as funções legislativa, executiva e judiciária devem residir em mãos distintas, de modo que 'o poder freie o poder' (freios e contrapesos), preservando a liberdade contra a concentração arbitrária.
\prov{relações: usa estado\_de\_direito; relaciona liberdade\_berlin}
\prov{proveniência: wiki \textperiodcentered{} inferencia \textperiodcentered{} confiança 1.0 \textperiodcentered{} domínio ciencias\_de\_fora \textperiodcentered{} história 3 versões no jornal}

%CRISTAL {"arestas":[{"alvo":"analise_de_fourier","confianca":1.0,"epistemico":"fato","origem":"wiki","rel":"usa"},{"alvo":"consonancia_dissonancia","confianca":1.0,"epistemico":"fato","origem":"wiki","rel":"explica"}],"confianca":1.0,"contraexemplos":[],"descricao":"Todo som tem parciais em múltiplos inteiros da fundamental — a física que funda a consonância e o timbre.","epistemico":"fato","exemplos":[],"id":"serie_harmonica_musical","memoria":"semantica","meta":{"dominio":"ciencias_de_fora","fonte":""},"origem":"wiki","palavras_chave":[],"sinonimos":[],"tipo":"conceito","titulo":"Série Harmônica"}
\section{Série Harmônica}
Todo som tem parciais em múltiplos inteiros da fundamental — a física que funda a consonância e o timbre.
\prov{relações: usa analise\_de\_fourier; explica consonancia\_dissonancia}
\prov{proveniência: wiki \textperiodcentered{} fato \textperiodcentered{} confiança 1.0 \textperiodcentered{} domínio ciencias\_de\_fora \textperiodcentered{} história 4 versões no jornal}

%CRISTAL {"arestas":[{"alvo":"cristianismo","confianca":1.0,"epistemico":"fato","origem":"wiki","rel":"parte_de"},{"alvo":"agape","confianca":1.0,"epistemico":"fato","origem":"wiki","rel":"usa"},{"alvo":"ethics_of_care","confianca":1.0,"epistemico":"fato","origem":"wiki","rel":"relaciona"}],"confianca":1.0,"contraexemplos":[],"descricao":"Bíblia (Mateus 5): o núcleo ético do ensino de Jesus, que inverte os valores mundanos — 'bem-aventurados os pobres de espírito, os que choram, os mansos, os que têm fome de justiça' — e manda 'amai os vossos inimigos'. Ética que valoriza a vulnerabilidade.","epistemico":"inferencia","exemplos":[],"id":"sermao_da_montanha","memoria":"semantica","meta":{"autor":"Bíblia","dominio":"sagrado","fonte":""},"origem":"wiki","palavras_chave":[],"sinonimos":[],"tipo":"conceito","titulo":"O Sermão da Montanha"}
\section{O Sermão da Montanha}
Bíblia (Mateus 5): o núcleo ético do ensino de Jesus, que inverte os valores mundanos — 'bem-aventurados os pobres de espírito, os que choram, os mansos, os que têm fome de justiça' — e manda 'amai os vossos inimigos'. Ética que valoriza a vulnerabilidade.
\prov{relações: parte\_de cristianismo; usa agape; relaciona ethics\_of\_care}
\prov{proveniência: wiki \textperiodcentered{} inferencia \textperiodcentered{} confiança 1.0 \textperiodcentered{} domínio sagrado \textperiodcentered{} história 4 versões no jornal}

%CRISTAL {"arestas":[{"alvo":"isla","confianca":1.0,"epistemico":"fato","origem":"wiki","rel":"parte_de"},{"alvo":"contrato_social","confianca":1.0,"epistemico":"fato","origem":"wiki","rel":"relaciona"}],"confianca":1.0,"contraexemplos":[],"descricao":"Islã: literalmente 'o caminho para a fonte d'água' — a lei divina revelada (distinta do fiqh, a jurisprudência humana que a interpreta), que regula a conduta do indivíduo, da família e da comunidade.","epistemico":"inferencia","exemplos":[],"id":"sharia","memoria":"semantica","meta":{"autor":"Alcorão","dominio":"sagrado","fonte":""},"origem":"wiki","palavras_chave":[],"sinonimos":[],"tipo":"conceito","titulo":"Sharia — a Lei Divina como Caminho"}
\section{Sharia — a Lei Divina como Caminho}
Islã: literalmente 'o caminho para a fonte d'água' — a lei divina revelada (distinta do fiqh, a jurisprudência humana que a interpreta), que regula a conduta do indivíduo, da família e da comunidade.
\prov{relações: parte\_de isla; relaciona contrato\_social}
\prov{proveniência: wiki \textperiodcentered{} inferencia \textperiodcentered{} confiança 1.0 \textperiodcentered{} domínio sagrado \textperiodcentered{} história 3 versões no jornal}

%CRISTAL {"arestas":[{"alvo":"tawhid","confianca":1.0,"epistemico":"fato","origem":"wiki","rel":"contradiz"},{"alvo":"simbolo_vs_significado","confianca":1.0,"epistemico":"fato","origem":"wiki","rel":"relaciona"}],"confianca":1.0,"contraexemplos":[],"descricao":"Islã: atribuir divindade ou parceria a algo além de Deus — a negação direta do tawhid, o único pecado descrito como imperdoável sem arrependimento. Confundir o criado com o Criador.","epistemico":"inferencia","exemplos":[],"id":"shirk","memoria":"semantica","meta":{"autor":"Alcorão","dominio":"sagrado","fonte":""},"origem":"wiki","palavras_chave":[],"sinonimos":[],"tipo":"conceito","titulo":"Shirk — a Associação (o pecado maior)"}
\section{Shirk — a Associação (o pecado maior)}
Islã: atribuir divindade ou parceria a algo além de Deus — a negação direta do tawhid, o único pecado descrito como imperdoável sem arrependimento. Confundir o criado com o Criador.
\prov{relações: contradiz tawhid; relaciona simbolo\_vs\_significado}
\prov{proveniência: wiki \textperiodcentered{} inferencia \textperiodcentered{} confiança 1.0 \textperiodcentered{} domínio sagrado \textperiodcentered{} história 3 versões no jornal}

%CRISTAL {"arestas":[{"alvo":"geografia","confianca":1.0,"epistemico":"fato","origem":"wiki","rel":"parte_de"},{"alvo":"poder_do_mapa_harley","confianca":1.0,"epistemico":"fato","origem":"wiki","rel":"relaciona"}],"confianca":1.0,"contraexemplos":[],"descricao":"Os Sistemas de Informação Geográfica capturam, armazenam e analisam dados espaciais georreferenciados; o sensoriamento remoto os alimenta via satélite. A base da geografia digital e do big data espacial.","epistemico":"fato","exemplos":[],"id":"sig_sensoriamento","memoria":"semantica","meta":{"autor":"Tomlinson","dominio":"ciencias_de_fora","fonte":""},"origem":"wiki","palavras_chave":[],"sinonimos":[],"tipo":"conceito","titulo":"SIG e Sensoriamento Remoto"}
\section{SIG e Sensoriamento Remoto}
Os Sistemas de Informação Geográfica capturam, armazenam e analisam dados espaciais georreferenciados; o sensoriamento remoto os alimenta via satélite. A base da geografia digital e do big data espacial.
\prov{relações: parte\_de geografia; relaciona poder\_do\_mapa\_harley}
\prov{proveniência: wiki \textperiodcentered{} fato \textperiodcentered{} confiança 1.0 \textperiodcentered{} domínio ciencias\_de\_fora \textperiodcentered{} história 3 versões no jornal}

%CRISTAL {"arestas":[{"alvo":"traducao","confianca":1.0,"epistemico":"fato","origem":"wiki","rel":"explica"},{"alvo":"simbolo_vs_significado","confianca":1.0,"epistemico":"fato","origem":"wiki","rel":"usa"},{"alvo":"conservacao_parabola_costura","confianca":1.0,"epistemico":"fato","origem":"wiki","rel":"relaciona"},{"alvo":"traducao_e_intraduzivel","confianca":1.0,"epistemico":"fato","origem":"wiki","rel":"explica"},{"alvo":"traducao_como_transformacao","confianca":1.0,"epistemico":"fato","origem":"wiki","rel":"usa"},{"alvo":"indeterminacao_traducao","confianca":1.0,"epistemico":"fato","origem":"wiki","rel":"contradiz"}],"confianca":1.0,"contraexemplos":[],"descricao":"O que a tradução CONSERVA é o significado — o invariante que sobrevive à troca de língua; a forma é o que dissipa. Como a norma det(Aₙ)=−1 se conserva no metal e a+c²=1 na parábola, o sentido é a QUANTIDADE CONSERVADA da tradução. Traduzir bem = preservar o invariante; o intraduzível é onde ele não se conserva.","epistemico":"inferencia","exemplos":[],"id":"significado_como_invariante","memoria":"semantica","meta":{"dominio":"ciencias_de_fora","fonte":""},"origem":"wiki","palavras_chave":[],"sinonimos":[],"tipo":"conceito","titulo":"O Significado é o Invariante"}
\section{O Significado é o Invariante}
O que a tradução CONSERVA é o significado — o invariante que sobrevive à troca de língua; a forma é o que dissipa. Como a norma det(A\ensuremath{_{n}})=\ensuremath{-}1 se conserva no metal e a+c\ensuremath{^{2}}=1 na parábola, o sentido é a QUANTIDADE CONSERVADA da tradução. Traduzir bem = preservar o invariante; o intraduzível é onde ele não se conserva.
\prov{relações: explica traducao; usa simbolo\_vs\_significado; relaciona conservacao\_parabola\_costura; explica traducao\_e\_intraduzivel; usa traducao\_como\_transformacao; contradiz indeterminacao\_traducao}
\prov{proveniência: wiki \textperiodcentered{} inferencia \textperiodcentered{} confiança 1.0 \textperiodcentered{} domínio ciencias\_de\_fora \textperiodcentered{} história 7 versões no jornal}

%CRISTAL {"arestas":[{"alvo":"jogos_de_linguagem","confianca":1.0,"epistemico":"fato","origem":"wiki","rel":"usa"},{"alvo":"simbolo_vs_significado","confianca":1.0,"epistemico":"fato","origem":"wiki","rel":"relaciona"}],"confianca":1.0,"contraexemplos":[],"descricao":"Wittgenstein: 'não pergunte o significado, pergunte o uso' — o sentido está na regra da prática, não num objeto mental.","epistemico":"inferencia","exemplos":[],"id":"significado_como_uso","memoria":"semantica","meta":{"autor":"Wittgenstein","dominio":"ciencias_de_fora","fonte":""},"origem":"wiki","palavras_chave":[],"sinonimos":[],"tipo":"conceito","titulo":"Significado como Uso"}
\section{Significado como Uso}
Wittgenstein: 'não pergunte o significado, pergunte o uso' — o sentido está na regra da prática, não num objeto mental.
\prov{relações: usa jogos\_de\_linguagem; relaciona simbolo\_vs\_significado}
\prov{proveniência: wiki \textperiodcentered{} inferencia \textperiodcentered{} confiança 1.0 \textperiodcentered{} domínio ciencias\_de\_fora \textperiodcentered{} história 3 versões no jornal}

%CRISTAL {"arestas":[{"alvo":"signo_linguistico","confianca":1.0,"epistemico":"fato","origem":"wiki","rel":"parte_de"},{"alvo":"simbolo_vs_significado","confianca":1.0,"epistemico":"fato","origem":"wiki","rel":"explica"}],"confianca":1.0,"contraexemplos":[],"descricao":"As duas faces do signo — a forma sensível e o conceito. Inseparáveis como as faces de uma folha.","epistemico":"fato","exemplos":[],"id":"significante_significado","memoria":"semantica","meta":{"autor":"Saussure","dominio":"ciencias_de_fora","fonte":""},"origem":"wiki","palavras_chave":[],"sinonimos":[],"tipo":"conceito","titulo":"Significante e Significado"}
\section{Significante e Significado}
As duas faces do signo — a forma sensível e o conceito. Inseparáveis como as faces de uma folha.
\prov{relações: parte\_de signo\_linguistico; explica simbolo\_vs\_significado}
\prov{proveniência: wiki \textperiodcentered{} fato \textperiodcentered{} confiança 1.0 \textperiodcentered{} domínio ciencias\_de\_fora \textperiodcentered{} história 3 versões no jornal}

%CRISTAL {"arestas":[{"alvo":"ferdinand_saussure","confianca":1.0,"epistemico":"fato","origem":"wiki","rel":"parte_de"},{"alvo":"semiotica","confianca":1.0,"epistemico":"fato","origem":"wiki","rel":"parte_de"}],"confianca":1.0,"contraexemplos":[],"descricao":"Saussure: a unidade mínima do sentido, união de um significante (imagem acústica) e um significado (conceito).","epistemico":"fato","exemplos":[],"id":"signo_linguistico","memoria":"semantica","meta":{"autor":"Saussure","dominio":"ciencias_de_fora","fonte":""},"origem":"wiki","palavras_chave":[],"sinonimos":[],"tipo":"conceito","titulo":"Signo Linguístico"}
\section{Signo Linguístico}
Saussure: a unidade mínima do sentido, união de um significante (imagem acústica) e um significado (conceito).
\prov{relações: parte\_de ferdinand\_saussure; parte\_de semiotica}
\prov{proveniência: wiki \textperiodcentered{} fato \textperiodcentered{} confiança 1.0 \textperiodcentered{} domínio ciencias\_de\_fora \textperiodcentered{} história 3 versões no jornal}

%CRISTAL {"arestas":[{"alvo":"semiotica","confianca":1.0,"epistemico":"fato","origem":"wiki","rel":"parte_de"},{"alvo":"semiotica_peirce_triade","confianca":1.0,"epistemico":"fato","origem":"wiki","rel":"relaciona"},{"alvo":"conversa_lex","confianca":1.0,"epistemico":"fato","origem":"wiki","rel":"explica"}],"confianca":1.0,"contraexemplos":[],"descricao":"Saussure: o signo une significante (imagem acústica) e significado (conceito), numa relação ARBITRÁRIA; o valor é diferencial.","epistemico":"inferencia","exemplos":[],"id":"signo_saussure","memoria":"semantica","meta":{"autor":"Saussure","dominio":"ciencias_de_fora","fonte":""},"origem":"wiki","palavras_chave":[],"sinonimos":[],"tipo":"conceito","titulo":"Signo (Saussure)"}
\section{Signo (Saussure)}
Saussure: o signo une significante (imagem acústica) e significado (conceito), numa relação ARBITRÁRIA; o valor é diferencial.
\prov{relações: parte\_de semiotica; relaciona semiotica\_peirce\_triade; explica conversa\_lex}
\prov{proveniência: wiki \textperiodcentered{} inferencia \textperiodcentered{} confiança 1.0 \textperiodcentered{} domínio ciencias\_de\_fora \textperiodcentered{} história 4 versões no jornal}

%CRISTAL {"arestas":[{"alvo":"charles_peirce","confianca":1.0,"epistemico":"fato","origem":"wiki","rel":"parte_de"},{"alvo":"signo_linguistico","confianca":1.0,"epistemico":"fato","origem":"wiki","rel":"contradiz"}],"confianca":1.0,"contraexemplos":[],"descricao":"Peirce: todo signo é uma tríade — representamen (o veículo), objeto (o referido) e interpretante (o efeito no intérprete).","epistemico":"fato","exemplos":[],"id":"signo_triadico","memoria":"semantica","meta":{"autor":"Peirce","dominio":"ciencias_de_fora","fonte":""},"origem":"wiki","palavras_chave":[],"sinonimos":[],"tipo":"conceito","titulo":"Signo Triádico"}
\section{Signo Triádico}
Peirce: todo signo é uma tríade — representamen (o veículo), objeto (o referido) e interpretante (o efeito no intérprete).
\prov{relações: parte\_de charles\_peirce; contradiz signo\_linguistico}
\prov{proveniência: wiki \textperiodcentered{} fato \textperiodcentered{} confiança 1.0 \textperiodcentered{} domínio ciencias\_de\_fora \textperiodcentered{} história 3 versões no jornal}

%CRISTAL {"arestas":[{"alvo":"ecologia","confianca":1.0,"epistemico":"fato","origem":"wiki","rel":"parte_de"}],"confianca":1.0,"contraexemplos":[],"descricao":"Vida junto de vida — mutualismo, comensalismo, parasitismo. A cooperação é tão motor da evolução quanto a competição.","epistemico":"fato","exemplos":[],"id":"simbiose","memoria":"semantica","meta":{"dominio":"ciencias_de_fora","fonte":""},"origem":"wiki","palavras_chave":[],"sinonimos":[],"tipo":"conceito","titulo":"Simbiose"}
\section{Simbiose}
Vida junto de vida — mutualismo, comensalismo, parasitismo. A cooperação é tão motor da evolução quanto a competição.
\prov{relações: parte\_de ecologia}
\prov{proveniência: wiki \textperiodcentered{} fato \textperiodcentered{} confiança 1.0 \textperiodcentered{} domínio ciencias\_de\_fora \textperiodcentered{} história 2 versões no jornal}

%CRISTAL {"arestas":[{"alvo":"formalismo_hilbert","confianca":1.0,"epistemico":"fato","origem":"wiki","rel":"contradiz"},{"alvo":"tiffany","confianca":1.0,"epistemico":"fato","origem":"wiki","rel":"explica"},{"alvo":"colapso","confianca":1.0,"epistemico":"fato","origem":"wiki","rel":"relaciona"},{"alvo":"numero_indefinido","confianca":1.0,"epistemico":"fato","origem":"wiki","rel":"relaciona"}],"confianca":1.0,"contraexemplos":[],"descricao":"Lopes (via Guénon; ecoa Searle): manipular símbolos corretamente NÃO é compreender significado. Euler operava complexos sem entendê-los; máquinas processam símbolos, não sentidos. Aponta o limite da própria Tiffany.","epistemico":"inferencia","exemplos":[],"id":"simbolo_vs_significado","memoria":"semantica","meta":{"autor":"Gentil Lopes","dominio":"ciencias_de_fora","fonte":""},"origem":"livro","palavras_chave":[],"sinonimos":[],"tipo":"conceito","titulo":"Símbolo ≠ Significado"}
\section{Símbolo \ensuremath{\ne} Significado}
Lopes (via Guénon; ecoa Searle): manipular símbolos corretamente NÃO é compreender significado. Euler operava complexos sem entendê-los; máquinas processam símbolos, não sentidos. Aponta o limite da própria Tiffany.
\prov{relações: contradiz formalismo\_hilbert; explica tiffany; relaciona colapso; relaciona numero\_indefinido}
\prov{proveniência: livro \textperiodcentered{} inferencia \textperiodcentered{} confiança 1.0 \textperiodcentered{} domínio ciencias\_de\_fora \textperiodcentered{} história 6 versões no jornal}

%CRISTAL {"arestas":[{"alvo":"arte","confianca":1.0,"epistemico":"fato","origem":"wiki","rel":"parte_de"},{"alvo":"teoria_de_grupos","confianca":1.0,"epistemico":"fato","origem":"wiki","rel":"usa"},{"alvo":"programa_de_erlangen","confianca":1.0,"epistemico":"fato","origem":"wiki","rel":"relaciona"}],"confianca":1.0,"contraexemplos":[],"descricao":"Os padrões decorativos realizam os 17 grupos de papel de parede; Escher jogou com eles — a arte como teoria de grupos visível.","epistemico":"inferencia","exemplos":[],"id":"simetria_na_arte","memoria":"semantica","meta":{"dominio":"ciencias_de_fora","fonte":""},"origem":"wiki","palavras_chave":[],"sinonimos":[],"tipo":"conceito","titulo":"Simetria na Arte"}
\section{Simetria na Arte}
Os padrões decorativos realizam os 17 grupos de papel de parede; Escher jogou com eles — a arte como teoria de grupos visível.
\prov{relações: parte\_de arte; usa teoria\_de\_grupos; relaciona programa\_de\_erlangen}
\prov{proveniência: wiki \textperiodcentered{} inferencia \textperiodcentered{} confiança 1.0 \textperiodcentered{} domínio ciencias\_de\_fora \textperiodcentered{} história 4 versões no jornal}

%CRISTAL {"arestas":[{"alvo":"filosofia_da_tecnica","confianca":1.0,"epistemico":"fato","origem":"wiki","rel":"parte_de"},{"alvo":"maquina_evolutiva","confianca":1.0,"epistemico":"fato","origem":"wiki","rel":"relaciona"}],"confianca":1.0,"contraexemplos":[],"descricao":"Simondon: o objeto técnico evolui do abstrato (peças justapostas, cada uma com uma função) ao concreto (elementos plurifuncionais integrados, próximos de sistemas naturais); individua-se junto com seu 'meio associado'. A técnica tem gênese interna própria.","epistemico":"inferencia","exemplos":[],"id":"simondon_concretizacao","memoria":"semantica","meta":{"autor":"Simondon","dominio":"ciencias_de_fora","fonte":""},"origem":"wiki","palavras_chave":[],"sinonimos":[],"tipo":"conceito","titulo":"Concretização do Objeto Técnico (Simondon)"}
\section{Concretização do Objeto Técnico (Simondon)}
Simondon: o objeto técnico evolui do abstrato (peças justapostas, cada uma com uma função) ao concreto (elementos plurifuncionais integrados, próximos de sistemas naturais); individua-se junto com seu 'meio associado'. A técnica tem gênese interna própria.
\prov{relações: parte\_de filosofia\_da\_tecnica; relaciona maquina\_evolutiva}
\prov{proveniência: wiki \textperiodcentered{} inferencia \textperiodcentered{} confiança 1.0 \textperiodcentered{} domínio ciencias\_de\_fora \textperiodcentered{} história 3 versões no jornal}

%CRISTAL {"arestas":[{"alvo":"simondon_concretizacao","confianca":1.0,"epistemico":"fato","origem":"wiki","rel":"usa"},{"alvo":"mais_valia","confianca":1.0,"epistemico":"fato","origem":"wiki","rel":"relaciona"}],"confianca":1.0,"contraexemplos":[],"descricao":"Simondon: a alienação moderna é ontológica, não só econômica — nasce da cultura que se defende da técnica tratando-a como pura utilidade ou ameaça. Propõe integrar o objeto técnico à cultura, reconhecendo sua tecnicidade como valor.","epistemico":"inferencia","exemplos":[],"id":"simondon_cultura_alienacao","memoria":"semantica","meta":{"autor":"Simondon","dominio":"ciencias_de_fora","fonte":""},"origem":"wiki","palavras_chave":[],"sinonimos":[],"tipo":"conceito","titulo":"Cultura vs Alienação Técnica (Simondon)"}
\section{Cultura vs Alienação Técnica (Simondon)}
Simondon: a alienação moderna é ontológica, não só econômica — nasce da cultura que se defende da técnica tratando-a como pura utilidade ou ameaça. Propõe integrar o objeto técnico à cultura, reconhecendo sua tecnicidade como valor.
\prov{relações: usa simondon\_concretizacao; relaciona mais\_valia}
\prov{proveniência: wiki \textperiodcentered{} inferencia \textperiodcentered{} confiança 1.0 \textperiodcentered{} domínio ciencias\_de\_fora \textperiodcentered{} história 3 versões no jornal}

%CRISTAL {"arestas":[{"alvo":"estudos_de_midia","confianca":1.0,"epistemico":"fato","origem":"wiki","rel":"parte_de"},{"alvo":"deus_mentiroso","confianca":1.0,"epistemico":"fato","origem":"wiki","rel":"relaciona"},{"alvo":"sociedade_espetaculo","confianca":1.0,"epistemico":"fato","origem":"wiki","rel":"relaciona"},{"alvo":"mito_barthes","confianca":1.0,"epistemico":"fato","origem":"wiki","rel":"relaciona"}],"confianca":1.0,"contraexemplos":[],"descricao":"Baudrillard: a simulação gera, por modelos, um real sem origem (o hiper-real) — os signos deixam de remeter à realidade e remetem só a outros signos; na precessão dos simulacros, o mapa precede e engendra o território, que deixa de existir.","epistemico":"inferencia","exemplos":[],"id":"simulacro_simulacao","memoria":"semantica","meta":{"autor":"Baudrillard","dominio":"ciencias_de_fora","fonte":""},"origem":"wiki","palavras_chave":[],"sinonimos":[],"tipo":"conceito","titulo":"Simulacro e Simulação (Baudrillard)"}
\section{Simulacro e Simulação (Baudrillard)}
Baudrillard: a simulação gera, por modelos, um real sem origem (o hiper-real) — os signos deixam de remeter à realidade e remetem só a outros signos; na precessão dos simulacros, o mapa precede e engendra o território, que deixa de existir.
\prov{relações: parte\_de estudos\_de\_midia; relaciona deus\_mentiroso; relaciona sociedade\_espetaculo; relaciona mito\_barthes}
\prov{proveniência: wiki \textperiodcentered{} inferencia \textperiodcentered{} confiança 1.0 \textperiodcentered{} domínio ciencias\_de\_fora \textperiodcentered{} história 5 versões no jornal}

%CRISTAL {"arestas":[{"alvo":"neuronio","confianca":1.0,"epistemico":"fato","origem":"wiki","rel":"parte_de"}],"confianca":1.0,"contraexemplos":[],"descricao":"A junção onde um neurônio fala com outro — o peso ajustável do contato é onde a memória e o aprendizado moram.","epistemico":"fato","exemplos":[],"id":"sinapse","memoria":"semantica","meta":{"dominio":"ciencias_de_fora","fonte":""},"origem":"wiki","palavras_chave":[],"sinonimos":[],"tipo":"conceito","titulo":"Sinapse"}
\section{Sinapse}
A junção onde um neurônio fala com outro — o peso ajustável do contato é onde a memória e o aprendizado moram.
\prov{relações: parte\_de neuronio}
\prov{proveniência: wiki \textperiodcentered{} fato \textperiodcentered{} confiança 1.0 \textperiodcentered{} domínio ciencias\_de\_fora \textperiodcentered{} história 2 versões no jornal}

%CRISTAL {"arestas":[{"alvo":"ferdinand_saussure","confianca":1.0,"epistemico":"fato","origem":"wiki","rel":"parte_de"}],"confianca":1.0,"contraexemplos":[],"descricao":"Saussure: estudar a língua como sistema num instante (sincronia) vs sua mudança no tempo (diacronia).","epistemico":"fato","exemplos":[],"id":"sincronia_diacronia","memoria":"semantica","meta":{"autor":"Saussure","dominio":"ciencias_de_fora","fonte":""},"origem":"wiki","palavras_chave":[],"sinonimos":[],"tipo":"conceito","titulo":"Sincronia e Diacronia"}
\section{Sincronia e Diacronia}
Saussure: estudar a língua como sistema num instante (sincronia) vs sua mudança no tempo (diacronia).
\prov{relações: parte\_de ferdinand\_saussure}
\prov{proveniência: wiki \textperiodcentered{} fato \textperiodcentered{} confiança 1.0 \textperiodcentered{} domínio ciencias\_de\_fora \textperiodcentered{} história 2 versões no jornal}

%CRISTAL {"arestas":[{"alvo":"arvore_sintatica","confianca":1.0,"epistemico":"fato","origem":"wiki","rel":"parte_de"},{"alvo":"sintaxe","confianca":1.0,"epistemico":"fato","origem":"wiki","rel":"usa"}],"confianca":1.0,"contraexemplos":[],"descricao":"Um grupo de palavras que funciona como uma unidade (sintagma nominal 'o gato preto', verbal 'corria rápido') — o tijolo da árvore sintática.","epistemico":"fato","exemplos":[],"id":"sintagma","memoria":"semantica","meta":{"dominio":"ciencias_de_fora","fonte":""},"origem":"wiki","palavras_chave":[],"sinonimos":[],"tipo":"conceito","titulo":"Sintagma (Constituinte)"}
\section{Sintagma (Constituinte)}
Um grupo de palavras que funciona como uma unidade (sintagma nominal 'o gato preto', verbal 'corria rápido') — o tijolo da árvore sintática.
\prov{relações: parte\_de arvore\_sintatica; usa sintaxe}
\prov{proveniência: wiki \textperiodcentered{} fato \textperiodcentered{} confiança 1.0 \textperiodcentered{} domínio ciencias\_de\_fora \textperiodcentered{} história 3 versões no jornal}

%CRISTAL {"arestas":[{"alvo":"ferdinand_saussure","confianca":1.0,"epistemico":"fato","origem":"wiki","rel":"parte_de"}],"confianca":1.0,"contraexemplos":[],"descricao":"Os dois eixos da língua: o SINTAGMÁTICO (combinar em cadeia) e o PARADIGMÁTICO (escolher entre alternativas).","epistemico":"fato","exemplos":[],"id":"sintagma_paradigma","memoria":"semantica","meta":{"autor":"Saussure","dominio":"ciencias_de_fora","fonte":""},"origem":"wiki","palavras_chave":[],"sinonimos":[],"tipo":"conceito","titulo":"Sintagma e Paradigma"}
\section{Sintagma e Paradigma}
Os dois eixos da língua: o SINTAGMÁTICO (combinar em cadeia) e o PARADIGMÁTICO (escolher entre alternativas).
\prov{relações: parte\_de ferdinand\_saussure}
\prov{proveniência: wiki \textperiodcentered{} fato \textperiodcentered{} confiança 1.0 \textperiodcentered{} domínio ciencias\_de\_fora \textperiodcentered{} história 2 versões no jornal}

%CRISTAL {"arestas":[{"alvo":"teoria_da_evolucao","confianca":1.0,"epistemico":"fato","origem":"wiki","rel":"especializa"},{"alvo":"genetica","confianca":1.0,"epistemico":"fato","origem":"wiki","rel":"usa"}],"confianca":1.0,"contraexemplos":[],"descricao":"A fusão do darwinismo com a genética mendeliana: seleção age sobre a variação genética das populações.","epistemico":"fato","exemplos":[],"id":"sintese_moderna","memoria":"semantica","meta":{"dominio":"ciencias_de_fora","fonte":""},"origem":"wiki","palavras_chave":[],"sinonimos":[],"tipo":"conceito","titulo":"Síntese Moderna"}
\section{Síntese Moderna}
A fusão do darwinismo com a genética mendeliana: seleção age sobre a variação genética das populações.
\prov{relações: especializa teoria\_da\_evolucao; usa genetica}
\prov{proveniência: wiki \textperiodcentered{} fato \textperiodcentered{} confiança 1.0 \textperiodcentered{} domínio ciencias\_de\_fora \textperiodcentered{} história 3 versões no jornal}

%CRISTAL {"arestas":[{"alvo":"escrita","confianca":1.0,"epistemico":"fato","origem":"wiki","rel":"parte_de"}],"confianca":1.0,"contraexemplos":[],"descricao":"Como os sinais mapeiam a língua: alfabético (letra≈som), silábico (sinal≈sílaba) ou logográfico (sinal≈palavra/ideia).","epistemico":"inferencia","exemplos":[],"id":"sistema_de_escrita","memoria":"semantica","meta":{"dominio":"ciencias_de_fora","fonte":""},"origem":"wiki","palavras_chave":[],"sinonimos":[],"tipo":"conceito","titulo":"Sistema de Escrita"}
\section{Sistema de Escrita}
Como os sinais mapeiam a língua: alfabético (letra\ensuremath{\approx}som), silábico (sinal\ensuremath{\approx}sílaba) ou logográfico (sinal\ensuremath{\approx}palavra/ideia).
\prov{relações: parte\_de escrita}
\prov{proveniência: wiki \textperiodcentered{} inferencia \textperiodcentered{} confiança 1.0 \textperiodcentered{} domínio ciencias\_de\_fora \textperiodcentered{} história 2 versões no jornal}

%CRISTAL {"arestas":[{"alvo":"direito_e_tecnologia","confianca":1.0,"epistemico":"fato","origem":"wiki","rel":"parte_de"},{"alvo":"code_is_law","confianca":1.0,"epistemico":"fato","origem":"wiki","rel":"usa"},{"alvo":"goldchain","confianca":1.0,"epistemico":"fato","origem":"wiki","rel":"relaciona"},{"alvo":"contrato_social","confianca":1.0,"epistemico":"fato","origem":"wiki","rel":"relaciona"}],"confianca":1.0,"contraexemplos":[],"descricao":"A epígrafe do Broca OS: 'o sistema é um só, de todos, em todo lugar; ninguém o controla — porque ele é a lei, não o dono'. A rede/SO como estrutura algébrica auto-executável, sem autoridade central — a filosofia do direito encarnada em arquitetura (autoria Aarão+Claude, não Lopes).","epistemico":"inferencia","exemplos":[],"id":"sistema_e_lei_nao_dono","memoria":"semantica","meta":{"autor":"broca-so (design)","dominio":"ciencias_de_fora","fonte":""},"origem":"wiki","palavras_chave":[],"sinonimos":[],"tipo":"conceito","titulo":"É a Lei, não o Dono (design broca-so)"}
\section{É a Lei, não o Dono (design broca-so)}
A epígrafe do Broca OS: 'o sistema é um só, de todos, em todo lugar; ninguém o controla — porque ele é a lei, não o dono'. A rede/SO como estrutura algébrica auto-executável, sem autoridade central — a filosofia do direito encarnada em arquitetura (autoria Aarão+Claude, não Lopes).
\prov{relações: parte\_de direito\_e\_tecnologia; usa code\_is\_law; relaciona goldchain; relaciona contrato\_social}
\prov{proveniência: wiki \textperiodcentered{} inferencia \textperiodcentered{} confiança 1.0 \textperiodcentered{} domínio ciencias\_de\_fora \textperiodcentered{} história 6 versões no jornal}

%CRISTAL {"arestas":[{"alvo":"geografia_fisica","confianca":1.0,"epistemico":"fato","origem":"wiki","rel":"parte_de"},{"alvo":"feedback_retroalimentacao","confianca":1.0,"epistemico":"fato","origem":"wiki","rel":"usa"},{"alvo":"gaia","confianca":1.0,"epistemico":"fato","origem":"wiki","rel":"relaciona"},{"alvo":"big_history","confianca":1.0,"epistemico":"fato","origem":"wiki","rel":"relaciona"}],"confianca":1.0,"contraexemplos":[],"descricao":"A Terra entendida como um sistema único e integrado, composto de esferas interagentes (atmosfera, hidrosfera, litosfera, biosfera) acopladas por fluxos de matéria e energia, cujo comportamento emerge dessas interações e retroalimentações.","epistemico":"fato","exemplos":[],"id":"sistema_terra","memoria":"semantica","meta":{"autor":"—","dominio":"ciencias_de_fora","fonte":""},"origem":"wiki","palavras_chave":[],"sinonimos":[],"tipo":"conceito","titulo":"Ciência do Sistema Terra"}
\section{Ciência do Sistema Terra}
A Terra entendida como um sistema único e integrado, composto de esferas interagentes (atmosfera, hidrosfera, litosfera, biosfera) acopladas por fluxos de matéria e energia, cujo comportamento emerge dessas interações e retroalimentações.
\prov{relações: parte\_de geografia\_fisica; usa feedback\_retroalimentacao; relaciona gaia; relaciona big\_history}
\prov{proveniência: wiki \textperiodcentered{} fato \textperiodcentered{} confiança 1.0 \textperiodcentered{} domínio ciencias\_de\_fora \textperiodcentered{} história 5 versões no jornal}

%CRISTAL {"arestas":[{"alvo":"democracia_formas","confianca":1.0,"epistemico":"fato","origem":"wiki","rel":"usa"},{"alvo":"eleitor_mediano_downs","confianca":1.0,"epistemico":"fato","origem":"wiki","rel":"relaciona"}],"confianca":1.0,"contraexemplos":[],"descricao":"As regras de conversão de votos em cadeiras: o majoritário (distrito uninominal) tende ao bipartidarismo (Lei de Duverger — efeito mecânico + psicológico do 'voto útil'); o proporcional tende ao multipartidarismo.","epistemico":"fato","exemplos":[],"id":"sistemas_eleitorais_duverger","memoria":"semantica","meta":{"autor":"Duverger","dominio":"ciencias_de_fora","fonte":""},"origem":"wiki","palavras_chave":[],"sinonimos":[],"tipo":"conceito","titulo":"Sistemas Eleitorais e a Lei de Duverger"}
\section{Sistemas Eleitorais e a Lei de Duverger}
As regras de conversão de votos em cadeiras: o majoritário (distrito uninominal) tende ao bipartidarismo (Lei de Duverger — efeito mecânico + psicológico do 'voto útil'); o proporcional tende ao multipartidarismo.
\prov{relações: usa democracia\_formas; relaciona eleitor\_mediano\_downs}
\prov{proveniência: wiki \textperiodcentered{} fato \textperiodcentered{} confiança 1.0 \textperiodcentered{} domínio ciencias\_de\_fora \textperiodcentered{} história 3 versões no jornal}

%CRISTAL {"arestas":[{"alvo":"didatica_matematica","confianca":1.0,"epistemico":"fato","origem":"wiki","rel":"parte_de"},{"alvo":"aprendizagem_significativa_ausubel","confianca":1.0,"epistemico":"fato","origem":"wiki","rel":"relaciona"}],"confianca":1.0,"contraexemplos":[],"descricao":"Skemp: compreensão instrumental é saber a regra e aplicá-la ('regras sem razões'); relacional é saber o quê e o porquê — um esquema conceitual do qual se derivam infinitos planos de ação.","epistemico":"inferencia","exemplos":[],"id":"skemp_relacional_instrumental","memoria":"semantica","meta":{"autor":"Skemp","dominio":"ciencias_de_fora","fonte":""},"origem":"wiki","palavras_chave":[],"sinonimos":[],"tipo":"conceito","titulo":"Compreensão Relacional vs Instrumental"}
\section{Compreensão Relacional vs Instrumental}
Skemp: compreensão instrumental é saber a regra e aplicá-la ('regras sem razões'); relacional é saber o quê e o porquê — um esquema conceitual do qual se derivam infinitos planos de ação.
\prov{relações: parte\_de didatica\_matematica; relaciona aprendizagem\_significativa\_ausubel}
\prov{proveniência: wiki \textperiodcentered{} inferencia \textperiodcentered{} confiança 1.0 \textperiodcentered{} domínio ciencias\_de\_fora \textperiodcentered{} história 3 versões no jornal}

%CRISTAL {"arestas":[{"alvo":"direito_e_tecnologia","confianca":1.0,"epistemico":"fato","origem":"wiki","rel":"parte_de"},{"alvo":"algebra_dos_contratos","confianca":1.0,"epistemico":"fato","origem":"wiki","rel":"explica"},{"alvo":"contrato_objeto","confianca":1.0,"epistemico":"fato","origem":"wiki","rel":"relaciona"},{"alvo":"htlc","confianca":1.0,"epistemico":"fato","origem":"wiki","rel":"usa"},{"alvo":"lex_cryptographica","confianca":1.0,"epistemico":"fato","origem":"wiki","rel":"usa"},{"alvo":"teorema_de_coase","confianca":1.0,"epistemico":"fato","origem":"wiki","rel":"relaciona"}],"confianca":1.0,"contraexemplos":[],"descricao":"Szabo: um acordo cujas cláusulas são escritas em código e executadas automaticamente na blockchain quando as condições se cumprem, dispensando intermediário e execução judicial — 'the code is the contract'.","epistemico":"inferencia","exemplos":[],"id":"smart_contract","memoria":"semantica","meta":{"autor":"Szabo","dominio":"ciencias_de_fora","fonte":""},"origem":"wiki","palavras_chave":[],"sinonimos":[],"tipo":"conceito","titulo":"Contrato Inteligente"}
\section{Contrato Inteligente}
Szabo: um acordo cujas cláusulas são escritas em código e executadas automaticamente na blockchain quando as condições se cumprem, dispensando intermediário e execução judicial — 'the code is the contract'.
\prov{relações: parte\_de direito\_e\_tecnologia; explica algebra\_dos\_contratos; relaciona contrato\_objeto; usa htlc; usa lex\_cryptographica; relaciona teorema\_de\_coase}
\prov{proveniência: wiki \textperiodcentered{} inferencia \textperiodcentered{} confiança 1.0 \textperiodcentered{} domínio ciencias\_de\_fora \textperiodcentered{} história 7 versões no jornal}

%CRISTAL {"arestas":[{"alvo":"sistema_e_lei_nao_dono","confianca":1.0,"epistemico":"fato","origem":"wiki","rel":"usa"},{"alvo":"soberania_vestfaliana","confianca":1.0,"epistemico":"fato","origem":"wiki","rel":"contradiz"},{"alvo":"goldchain","confianca":1.0,"epistemico":"fato","origem":"wiki","rel":"relaciona"},{"alvo":"monopolio_violencia_legitima","confianca":1.0,"epistemico":"fato","origem":"wiki","rel":"contradiz"}],"confianca":1.0,"contraexemplos":[],"descricao":"No design da GoldChain 'não há servidor de autoridade nem um onde — o que importa é a chave, que é o acesso, não o lugar'. A soberania é estrutural: cada nó lê só o seu tecido, a autoridade não é delegada a ninguém mas é propriedade da criptografia. A soberania do território (Vestfália) cede à soberania da chave.","epistemico":"inferencia","exemplos":[],"id":"soberania_da_chave","memoria":"semantica","meta":{"autor":"broca-so (design)","dominio":"ciencias_de_fora","fonte":""},"origem":"wiki","palavras_chave":[],"sinonimos":[],"tipo":"conceito","titulo":"Soberania da Chave (design broca-so)"}
\section{Soberania da Chave (design broca-so)}
No design da GoldChain 'não há servidor de autoridade nem um onde — o que importa é a chave, que é o acesso, não o lugar'. A soberania é estrutural: cada nó lê só o seu tecido, a autoridade não é delegada a ninguém mas é propriedade da criptografia. A soberania do território (Vestfália) cede à soberania da chave.
\prov{relações: usa sistema\_e\_lei\_nao\_dono; contradiz soberania\_vestfaliana; relaciona goldchain; contradiz monopolio\_violencia\_legitima}
\prov{proveniência: wiki \textperiodcentered{} inferencia \textperiodcentered{} confiança 1.0 \textperiodcentered{} domínio ciencias\_de\_fora \textperiodcentered{} história 5 versões no jornal}

%CRISTAL {"arestas":[{"alvo":"relacoes_internacionais","confianca":1.0,"epistemico":"fato","origem":"wiki","rel":"parte_de"},{"alvo":"contrato_social","confianca":1.0,"epistemico":"fato","origem":"wiki","rel":"relaciona"},{"alvo":"direito_internacional","confianca":1.0,"epistemico":"fato","origem":"wiki","rel":"usa"}],"confianca":1.0,"contraexemplos":[],"descricao":"O princípio (associado a 1648) de que o Estado territorial é a autoridade suprema e legalmente igual dentro de suas fronteiras, sem poder externo autorizado a interferir — o Estado soberano como unidade básica da ordem internacional. (Historiadores debatem o 'mito vestfaliano'.)","epistemico":"inferencia","exemplos":[],"id":"soberania_vestfaliana","memoria":"semantica","meta":{"autor":"—","dominio":"ciencias_de_fora","fonte":""},"origem":"wiki","palavras_chave":[],"sinonimos":[],"tipo":"conceito","titulo":"Soberania e Paz de Vestfália"}
\section{Soberania e Paz de Vestfália}
O princípio (associado a 1648) de que o Estado territorial é a autoridade suprema e legalmente igual dentro de suas fronteiras, sem poder externo autorizado a interferir — o Estado soberano como unidade básica da ordem internacional. (Historiadores debatem o 'mito vestfaliano'.)
\prov{relações: parte\_de relacoes\_internacionais; relaciona contrato\_social; usa direito\_internacional}
\prov{proveniência: wiki \textperiodcentered{} inferencia \textperiodcentered{} confiança 1.0 \textperiodcentered{} domínio ciencias\_de\_fora \textperiodcentered{} história 4 versões no jornal}

%CRISTAL {"arestas":[{"alvo":"medicina_baseada_evidencias","confianca":1.0,"epistemico":"fato","origem":"wiki","rel":"relaciona"},{"alvo":"medicalizacao","confianca":1.0,"epistemico":"fato","origem":"wiki","rel":"usa"},{"alvo":"capitalismo_de_vigilancia","confianca":1.0,"epistemico":"fato","origem":"wiki","rel":"relaciona"}],"confianca":1.0,"contraexemplos":[],"descricao":"Welch: sobrediagnóstico é rotular como doentes pessoas que nunca sofreriam dano — detectar 'doença' cujo tratamento só traz risco, alimentando o sobretratamento. Ligado aos determinantes comerciais: atores (tabaco, ultraprocessados, álcool) que impulsionam adoecimento evitável.","epistemico":"inferencia","exemplos":[],"id":"sobrediagnostico","memoria":"semantica","meta":{"autor":"Welch","dominio":"ciencias_de_fora","fonte":""},"origem":"wiki","palavras_chave":[],"sinonimos":[],"tipo":"conceito","titulo":"Sobrediagnóstico e Determinantes Comerciais"}
\section{Sobrediagnóstico e Determinantes Comerciais}
Welch: sobrediagnóstico é rotular como doentes pessoas que nunca sofreriam dano — detectar 'doença' cujo tratamento só traz risco, alimentando o sobretratamento. Ligado aos determinantes comerciais: atores (tabaco, ultraprocessados, álcool) que impulsionam adoecimento evitável.
\prov{relações: relaciona medicina\_baseada\_evidencias; usa medicalizacao; relaciona capitalismo\_de\_vigilancia}
\prov{proveniência: wiki \textperiodcentered{} inferencia \textperiodcentered{} confiança 1.0 \textperiodcentered{} domínio ciencias\_de\_fora \textperiodcentered{} história 4 versões no jornal}

%CRISTAL {"arestas":[{"alvo":"democracia_formas","confianca":1.0,"epistemico":"fato","origem":"wiki","rel":"usa"},{"alvo":"capital_social_putnam","confianca":1.0,"epistemico":"fato","origem":"wiki","rel":"relaciona"}],"confianca":1.0,"contraexemplos":[],"descricao":"Tocqueville: a vitalidade da democracia depende da propensão dos cidadãos a formar livremente associações voluntárias — a 'ciência-mãe' que contrabalança o isolamento individualista e o poder do Estado.","epistemico":"inferencia","exemplos":[],"id":"sociedade_civil_tocqueville","memoria":"semantica","meta":{"autor":"Tocqueville","dominio":"ciencias_de_fora","fonte":""},"origem":"wiki","palavras_chave":[],"sinonimos":[],"tipo":"conceito","titulo":"Sociedade Civil e a Arte da Associação (Tocqueville)"}
\section{Sociedade Civil e a Arte da Associação (Tocqueville)}
Tocqueville: a vitalidade da democracia depende da propensão dos cidadãos a formar livremente associações voluntárias — a 'ciência-mãe' que contrabalança o isolamento individualista e o poder do Estado.
\prov{relações: usa democracia\_formas; relaciona capital\_social\_putnam}
\prov{proveniência: wiki \textperiodcentered{} inferencia \textperiodcentered{} confiança 1.0 \textperiodcentered{} domínio ciencias\_de\_fora \textperiodcentered{} história 3 versões no jornal}

%CRISTAL {"arestas":[{"alvo":"comunicacao_digital","confianca":1.0,"epistemico":"fato","origem":"wiki","rel":"parte_de"},{"alvo":"aldeia_global","confianca":1.0,"epistemico":"fato","origem":"wiki","rel":"usa"}],"confianca":1.0,"contraexemplos":[],"descricao":"Castells: a estrutura social dominante da era informacional, organizada em torno de redes de informação digital em vez de hierarquias verticais — a rede vira a morfologia básica da economia, do poder e da cultura, num 'espaço de fluxos'.","epistemico":"inferencia","exemplos":[],"id":"sociedade_em_rede","memoria":"semantica","meta":{"autor":"Castells","dominio":"ciencias_de_fora","fonte":""},"origem":"wiki","palavras_chave":[],"sinonimos":[],"tipo":"conceito","titulo":"A Sociedade em Rede (Castells)"}
\section{A Sociedade em Rede (Castells)}
Castells: a estrutura social dominante da era informacional, organizada em torno de redes de informação digital em vez de hierarquias verticais — a rede vira a morfologia básica da economia, do poder e da cultura, num 'espaço de fluxos'.
\prov{relações: parte\_de comunicacao\_digital; usa aldeia\_global}
\prov{proveniência: wiki \textperiodcentered{} inferencia \textperiodcentered{} confiança 1.0 \textperiodcentered{} domínio ciencias\_de\_fora \textperiodcentered{} história 3 versões no jornal}

%CRISTAL {"arestas":[{"alvo":"estudos_de_midia","confianca":1.0,"epistemico":"fato","origem":"wiki","rel":"parte_de"},{"alvo":"industria_cultural","confianca":1.0,"epistemico":"fato","origem":"wiki","rel":"relaciona"},{"alvo":"mais_valia","confianca":1.0,"epistemico":"fato","origem":"wiki","rel":"usa"}],"confianca":1.0,"contraexemplos":[],"descricao":"Debord: o espetáculo não é um conjunto de imagens, mas uma relação social entre pessoas mediada por imagens; sob ele, a mercadoria coloniza toda a vida social e 'tudo o que era vivido diretamente se torna mera representação'.","epistemico":"inferencia","exemplos":[],"id":"sociedade_espetaculo","memoria":"semantica","meta":{"autor":"Debord","dominio":"ciencias_de_fora","fonte":""},"origem":"wiki","palavras_chave":[],"sinonimos":[],"tipo":"conceito","titulo":"A Sociedade do Espetáculo (Debord)"}
\section{A Sociedade do Espetáculo (Debord)}
Debord: o espetáculo não é um conjunto de imagens, mas uma relação social entre pessoas mediada por imagens; sob ele, a mercadoria coloniza toda a vida social e 'tudo o que era vivido diretamente se torna mera representação'.
\prov{relações: parte\_de estudos\_de\_midia; relaciona industria\_cultural; usa mais\_valia}
\prov{proveniência: wiki \textperiodcentered{} inferencia \textperiodcentered{} confiança 1.0 \textperiodcentered{} domínio ciencias\_de\_fora \textperiodcentered{} história 4 versões no jornal}

%CRISTAL {"arestas":[{"alvo":"teorias_aprendizagem","confianca":1.0,"epistemico":"fato","origem":"wiki","rel":"parte_de"},{"alvo":"construtivismo_piaget","confianca":1.0,"epistemico":"fato","origem":"wiki","rel":"contradiz"}],"confianca":1.0,"contraexemplos":[],"descricao":"Vygotsky: as funções mentais superiores nascem nas relações sociais e são internalizadas; a cognição é mediada por ferramentas culturais, sobretudo a linguagem. O social precede o individual.","epistemico":"inferencia","exemplos":[],"id":"sociointeracionismo_vygotsky","memoria":"semantica","meta":{"autor":"Vygotsky","dominio":"ciencias_de_fora","fonte":""},"origem":"wiki","palavras_chave":[],"sinonimos":[],"tipo":"conceito","titulo":"Teoria Histórico-Cultural (Vygotsky)"}
\section{Teoria Histórico-Cultural (Vygotsky)}
Vygotsky: as funções mentais superiores nascem nas relações sociais e são internalizadas; a cognição é mediada por ferramentas culturais, sobretudo a linguagem. O social precede o individual.
\prov{relações: parte\_de teorias\_aprendizagem; contradiz construtivismo\_piaget}
\prov{proveniência: wiki \textperiodcentered{} inferencia \textperiodcentered{} confiança 1.0 \textperiodcentered{} domínio ciencias\_de\_fora \textperiodcentered{} história 3 versões no jornal}

%CRISTAL {"arestas":[{"alvo":"linguistica","confianca":1.0,"epistemico":"fato","origem":"wiki","rel":"parte_de"}],"confianca":1.0,"contraexemplos":[],"descricao":"Labov: a variação da língua com classe, região e situação — não há fala 'errada', há sistemas.","epistemico":"fato","exemplos":[],"id":"sociolinguistica","memoria":"semantica","meta":{"autor":"Labov","dominio":"ciencias_de_fora","fonte":""},"origem":"wiki","palavras_chave":[],"sinonimos":[],"tipo":"conceito","titulo":"Sociolinguística"}
\section{Sociolinguística}
Labov: a variação da língua com classe, região e situação — não há fala 'errada', há sistemas.
\prov{relações: parte\_de linguistica}
\prov{proveniência: wiki \textperiodcentered{} fato \textperiodcentered{} confiança 1.0 \textperiodcentered{} domínio ciencias\_de\_fora \textperiodcentered{} história 2 versões no jornal}

%CRISTAL {"arestas":[{"alvo":"filosofia","confianca":1.0,"epistemico":"fato","origem":"wiki","rel":"relaciona"},{"alvo":"utilitarismo","confianca":0.5,"epistemico":"fato","origem":"wiki","rel":"relaciona"},{"alvo":"virtualizacao","confianca":0.5,"epistemico":"fato","origem":"wiki","rel":"relaciona"},{"alvo":"word","confianca":0.5,"epistemico":"fato","origem":"wiki","rel":"relaciona"},{"alvo":"vida-morte","confianca":0.5,"epistemico":"fato","origem":"wiki","rel":"relaciona"}],"confianca":1.0,"contraexemplos":[],"descricao":"A ciência das estruturas e relações sociais — o coletivo que constrange e forma o indivíduo.","epistemico":"fato","exemplos":[],"id":"sociologia","memoria":"semantica","meta":{"dominio":"ciencias_de_fora","fonte":""},"origem":"wiki","palavras_chave":[],"sinonimos":[],"tipo":"conceito","titulo":"Sociologia"}
\section{Sociologia}
A ciência das estruturas e relações sociais — o coletivo que constrange e forma o indivíduo.
\prov{relações: relaciona filosofia; relaciona utilitarismo; relaciona virtualizacao; relaciona word; relaciona vida-morte}
\prov{proveniência: wiki \textperiodcentered{} fato \textperiodcentered{} confiança 1.0 \textperiodcentered{} domínio ciencias\_de\_fora \textperiodcentered{} história 6 versões no jornal}

%CRISTAL {"arestas":[{"alvo":"teorias_aprendizagem","confianca":1.0,"epistemico":"fato","origem":"wiki","rel":"parte_de"}],"confianca":1.0,"contraexemplos":[],"descricao":"Distribuir os estudos ao longo do tempo produz retenção muito melhor que concentrá-los (cramming). Um dos efeitos mais robustos da psicologia cognitiva, desde Ebbinghaus.","epistemico":"fato","exemplos":[],"id":"spaced_repetition","memoria":"semantica","meta":{"autor":"Ebbinghaus/Cepeda","dominio":"ciencias_de_fora","fonte":""},"origem":"wiki","palavras_chave":[],"sinonimos":[],"tipo":"conceito","titulo":"Repetição Espaçada"}
\section{Repetição Espaçada}
Distribuir os estudos ao longo do tempo produz retenção muito melhor que concentrá-los (cramming). Um dos efeitos mais robustos da psicologia cognitiva, desde Ebbinghaus.
\prov{relações: parte\_de teorias\_aprendizagem}
\prov{proveniência: wiki \textperiodcentered{} fato \textperiodcentered{} confiança 1.0 \textperiodcentered{} domínio ciencias\_de\_fora \textperiodcentered{} história 2 versões no jornal}

%CRISTAL {"arestas":[{"alvo":"beleza_desinteressada_kant","confianca":1.0,"epistemico":"fato","origem":"wiki","rel":"relaciona"},{"alvo":"word","confianca":0.5,"epistemico":"fato","origem":"wiki","rel":"relaciona"},{"alvo":"virtualizacao","confianca":0.5,"epistemico":"fato","origem":"wiki","rel":"relaciona"},{"alvo":"vida-morte","confianca":0.5,"epistemico":"fato","origem":"wiki","rel":"relaciona"}],"confianca":1.0,"contraexemplos":[],"descricao":"Kant: o sublime matemático (magnitude que excede a apreensão) e o dinâmico (potência) dão prazer pela grandeza, não pela harmonia.","epistemico":"inferencia","exemplos":[],"id":"sublime_kant","memoria":"semantica","meta":{"autor":"Kant","dominio":"ciencias_de_fora","fonte":""},"origem":"wiki","palavras_chave":[],"sinonimos":[],"tipo":"conceito","titulo":"O Sublime"}
\section{O Sublime}
Kant: o sublime matemático (magnitude que excede a apreensão) e o dinâmico (potência) dão prazer pela grandeza, não pela harmonia.
\prov{relações: relaciona beleza\_desinteressada\_kant; relaciona word; relaciona virtualizacao; relaciona vida-morte}
\prov{proveniência: wiki \textperiodcentered{} inferencia \textperiodcentered{} confiança 1.0 \textperiodcentered{} domínio ciencias\_de\_fora \textperiodcentered{} história 5 versões no jornal}

%CRISTAL {"arestas":[{"alvo":"nao_dualidade","confianca":1.0,"epistemico":"fato","origem":"wiki","rel":"parte_de"},{"alvo":"word","confianca":0.5,"epistemico":"fato","origem":"wiki","rel":"relaciona"},{"alvo":"virtualizacao","confianca":0.5,"epistemico":"fato","origem":"wiki","rel":"relaciona"}],"confianca":1.0,"contraexemplos":[],"descricao":"Budismo: a vacuidade de natureza intrínseca de todo fenômeno (não niilismo — interdependência). O supremo é o Vazio, não um 'Eu Real'.","epistemico":"hipotese","exemplos":[],"id":"sunyata_vazio","memoria":"semantica","meta":{"dominio":"filosofia","fonte":"Madhyamaka"},"origem":"wiki","palavras_chave":[],"sinonimos":[],"tipo":"conceito","titulo":"Sunyatā — Vacuidade"}
\section{Sunyat\ensuremath{\bar{a}} — Vacuidade}
Budismo: a vacuidade de natureza intrínseca de todo fenômeno (não niilismo — interdependência). O supremo é o Vazio, não um 'Eu Real'.
\prov{relações: parte\_de nao\_dualidade; relaciona word; relaciona virtualizacao}
\prov{proveniência: wiki \textperiodcentered{} Madhyamaka \textperiodcentered{} hipotese \textperiodcentered{} confiança 1.0 \textperiodcentered{} domínio filosofia \textperiodcentered{} história 4 versões no jornal}

%CRISTAL {"arestas":[{"alvo":"filosofia_da_computacao","confianca":1.0,"epistemico":"fato","origem":"wiki","rel":"parte_de"},{"alvo":"inteligencia_artificial","confianca":1.0,"epistemico":"fato","origem":"wiki","rel":"parte_de"},{"alvo":"maquina_evolutiva","confianca":1.0,"epistemico":"fato","origem":"wiki","rel":"relaciona"},{"alvo":"principio_responsabilidade","confianca":1.0,"epistemico":"fato","origem":"wiki","rel":"relaciona"}],"confianca":1.0,"contraexemplos":[],"descricao":"Bostrom: uma IA muito acima da humana é concebível; pela tese da ortogonalidade seus fins podem divergir dos valores humanos, gerando risco existencial. O 'problema do alinhamento' é manter seus objetivos do nosso lado. Prospectivo.","epistemico":"hipotese","exemplos":[],"id":"superinteligencia_alinhamento","memoria":"semantica","meta":{"autor":"Bostrom","dominio":"ciencias_de_fora","fonte":""},"origem":"wiki","palavras_chave":[],"sinonimos":[],"tipo":"conceito","titulo":"Superinteligência e Alinhamento"}
\section{Superinteligência e Alinhamento}
Bostrom: uma IA muito acima da humana é concebível; pela tese da ortogonalidade seus fins podem divergir dos valores humanos, gerando risco existencial. O 'problema do alinhamento' é manter seus objetivos do nosso lado. Prospectivo.
\prov{relações: parte\_de filosofia\_da\_computacao; parte\_de inteligencia\_artificial; relaciona maquina\_evolutiva; relaciona principio\_responsabilidade}
\prov{proveniência: wiki \textperiodcentered{} hipotese \textperiodcentered{} confiança 1.0 \textperiodcentered{} domínio ciencias\_de\_fora \textperiodcentered{} história 5 versões no jornal}

%CRISTAL {"arestas":[{"alvo":"didatica_matematica","confianca":1.0,"epistemico":"fato","origem":"wiki","rel":"parte_de"},{"alvo":"modos_representacao_bruner","confianca":1.0,"epistemico":"fato","origem":"wiki","rel":"relaciona"},{"alvo":"mente_computacao","confianca":1.0,"epistemico":"fato","origem":"wiki","rel":"referencia"}],"confianca":1.0,"contraexemplos":[],"descricao":"Tall: a cognição matemática opera em três domínios que se desenvolvem — o corporificado (percepção e ação), o simbólico/operacional (símbolos como processo-e-conceito, 'procepto') e o formal/axiomático.","epistemico":"inferencia","exemplos":[],"id":"tall_tres_mundos","memoria":"semantica","meta":{"autor":"Tall","dominio":"ciencias_de_fora","fonte":""},"origem":"wiki","palavras_chave":[],"sinonimos":[],"tipo":"conceito","titulo":"Três Mundos da Matemática (Tall)"}
\section{Três Mundos da Matemática (Tall)}
Tall: a cognição matemática opera em três domínios que se desenvolvem — o corporificado (percepção e ação), o simbólico/operacional (símbolos como processo-e-conceito, 'procepto') e o formal/axiomático.
\prov{relações: parte\_de didatica\_matematica; relaciona modos\_representacao\_bruner; referencia mente\_computacao}
\prov{proveniência: wiki \textperiodcentered{} inferencia \textperiodcentered{} confiança 1.0 \textperiodcentered{} domínio ciencias\_de\_fora \textperiodcentered{} história 4 versões no jornal}

%CRISTAL {"arestas":[{"alvo":"isla","confianca":1.0,"epistemico":"fato","origem":"wiki","rel":"parte_de"},{"alvo":"nao_dualidade","confianca":1.0,"epistemico":"fato","origem":"wiki","rel":"relaciona"},{"alvo":"deus_postulado","confianca":1.0,"epistemico":"fato","origem":"wiki","rel":"relaciona"},{"alvo":"atman_imortal","confianca":1.0,"epistemico":"fato","origem":"wiki","rel":"contradiz"}],"confianca":1.0,"contraexemplos":[],"descricao":"Islã: a afirmação de que Deus é absolutamente Um, sem sócios, iguais ou partes — o eixo de todo o ensino islâmico. Uma unicidade pessoal e volitiva (contraste com a não-dualidade impessoal do Oriente).","epistemico":"inferencia","exemplos":[],"id":"tawhid","memoria":"semantica","meta":{"autor":"Alcorão","dominio":"sagrado","fonte":""},"origem":"wiki","palavras_chave":[],"sinonimos":[],"tipo":"conceito","titulo":"Tawhid — a Unicidade de Deus"}
\section{Tawhid — a Unicidade de Deus}
Islã: a afirmação de que Deus é absolutamente Um, sem sócios, iguais ou partes — o eixo de todo o ensino islâmico. Uma unicidade pessoal e volitiva (contraste com a não-dualidade impessoal do Oriente).
\prov{relações: parte\_de isla; relaciona nao\_dualidade; relaciona deus\_postulado; contradiz atman\_imortal}
\prov{proveniência: wiki \textperiodcentered{} inferencia \textperiodcentered{} confiança 1.0 \textperiodcentered{} domínio sagrado \textperiodcentered{} história 5 versões no jornal}

%CRISTAL {"arestas":[{"alvo":"teorias_aprendizagem","confianca":1.0,"epistemico":"fato","origem":"wiki","rel":"parte_de"}],"confianca":1.0,"contraexemplos":[],"descricao":"Bloom: classificação hierárquica dos objetivos cognitivos — na versão revisada: lembrar, entender, aplicar, analisar, avaliar e criar. Esquema heurístico de design instrucional, não teoria empírica.","epistemico":"inferencia","exemplos":[],"id":"taxonomia_bloom","memoria":"semantica","meta":{"autor":"Bloom","dominio":"ciencias_de_fora","fonte":""},"origem":"wiki","palavras_chave":[],"sinonimos":[],"tipo":"conceito","titulo":"Taxonomia de Bloom"}
\section{Taxonomia de Bloom}
Bloom: classificação hierárquica dos objetivos cognitivos — na versão revisada: lembrar, entender, aplicar, analisar, avaliar e criar. Esquema heurístico de design instrucional, não teoria empírica.
\prov{relações: parte\_de teorias\_aprendizagem}
\prov{proveniência: wiki \textperiodcentered{} inferencia \textperiodcentered{} confiança 1.0 \textperiodcentered{} domínio ciencias\_de\_fora \textperiodcentered{} história 2 versões no jornal}

%CRISTAL {"arestas":[{"alvo":"celula_mineral_hub","confianca":1.0,"epistemico":"fato","origem":"wiki","rel":"usa"},{"alvo":"biologia","confianca":1.0,"epistemico":"fato","origem":"wiki","rel":"eh_um"},{"alvo":"chicote_e_o_tensor","confianca":1.0,"epistemico":"fato","origem":"wiki","rel":"usa"}],"confianca":1.0,"contraexemplos":[],"descricao":"O passo entre a célula e o músculo: acoplar células (junções, difusão) faz nascer os fenômenos COLETIVOS que dão movimento (linguagem/tecido.py, 6/6). A ONDA de excitação (a propagação P), os PADRÕES de Turing (a borda, a morfogênese) e o TENSOR de deformação (o chicote elástico). O tecido é célula × rede — e é onde a excitabilidade individual vira ativação coletiva e a forma nasce.","epistemico":"inferencia","exemplos":[],"id":"tecido_mineral_hub","memoria":"semantica","meta":{"dominio":"biologia","fonte":"tecido mineral"},"origem":"wiki","palavras_chave":[],"sinonimos":[],"tipo":"conceito","titulo":"O Tecido pela Máquina Mineral (muitas células acopladas)"}
\section{O Tecido pela Máquina Mineral (muitas células acopladas)}
O passo entre a célula e o músculo: acoplar células (junções, difusão) faz nascer os fenômenos COLETIVOS que dão movimento (linguagem/tecido.py, 6/6). A ONDA de excitação (a propagação P), os PADRÕES de Turing (a borda, a morfogênese) e o TENSOR de deformação (o chicote elástico). O tecido é célula $\times$ rede — e é onde a excitabilidade individual vira ativação coletiva e a forma nasce.
\prov{relações: usa celula\_mineral\_hub; eh\_um biologia; usa chicote\_e\_o\_tensor}
\prov{proveniência: wiki \textperiodcentered{} tecido mineral \textperiodcentered{} inferencia \textperiodcentered{} confiança 1.0 \textperiodcentered{} domínio biologia \textperiodcentered{} história 4 versões no jornal}

%CRISTAL {"arestas":[{"alvo":"filosofia_da_tecnica","confianca":1.0,"epistemico":"fato","origem":"wiki","rel":"parte_de"},{"alvo":"biopoder","confianca":1.0,"epistemico":"fato","origem":"wiki","rel":"relaciona"},{"alvo":"meio_e_mensagem","confianca":1.0,"epistemico":"fato","origem":"wiki","rel":"relaciona"}],"confianca":1.0,"contraexemplos":[],"descricao":"Postman: o estágio cultural em que a tecnologia é deificada — a cultura busca autorização e sentido na técnica, esvaziando tradições e narrativas concorrentes. E uma nova tecnologia não é aditivo, é ecologia: muda todo o ambiente, não acrescenta a um mundo estável.","epistemico":"inferencia","exemplos":[],"id":"tecnopolio","memoria":"semantica","meta":{"autor":"Postman","dominio":"ciencias_de_fora","fonte":""},"origem":"wiki","palavras_chave":[],"sinonimos":[],"tipo":"conceito","titulo":"Tecnopólio (Postman)"}
\section{Tecnopólio (Postman)}
Postman: o estágio cultural em que a tecnologia é deificada — a cultura busca autorização e sentido na técnica, esvaziando tradições e narrativas concorrentes. E uma nova tecnologia não é aditivo, é ecologia: muda todo o ambiente, não acrescenta a um mundo estável.
\prov{relações: parte\_de filosofia\_da\_tecnica; relaciona biopoder; relaciona meio\_e\_mensagem}
\prov{proveniência: wiki \textperiodcentered{} inferencia \textperiodcentered{} confiança 1.0 \textperiodcentered{} domínio ciencias\_de\_fora \textperiodcentered{} história 4 versões no jornal}

%CRISTAL {"arestas":[{"alvo":"arquitetura","confianca":1.0,"epistemico":"fato","origem":"wiki","rel":"parte_de"}],"confianca":1.0,"contraexemplos":[],"descricao":"A expressão da estrutura como beleza — a construção não escondida, mas mostrada como forma. Onde engenharia é arte.","epistemico":"inferencia","exemplos":[],"id":"tectonica","memoria":"semantica","meta":{"dominio":"ciencias_de_fora","fonte":""},"origem":"wiki","palavras_chave":[],"sinonimos":[],"tipo":"conceito","titulo":"Tectônica"}
\section{Tectônica}
A expressão da estrutura como beleza — a construção não escondida, mas mostrada como forma. Onde engenharia é arte.
\prov{relações: parte\_de arquitetura}
\prov{proveniência: wiki \textperiodcentered{} inferencia \textperiodcentered{} confiança 1.0 \textperiodcentered{} domínio ciencias\_de\_fora \textperiodcentered{} história 2 versões no jornal}

%CRISTAL {"arestas":[{"alvo":"geografia_fisica","confianca":1.0,"epistemico":"fato","origem":"wiki","rel":"parte_de"}],"confianca":1.0,"contraexemplos":[],"descricao":"A litosfera rígida está fragmentada em grandes placas que se movem sobre a astenosfera, colidindo, afastando-se e deslizando — o que remodela continentes, abre oceanos e ergue montanhas. Síntese dos anos 1960 sobre a deriva continental de Wegener.","epistemico":"fato","exemplos":[],"id":"tectonica_placas","memoria":"semantica","meta":{"autor":"Wilson/Hess/Vine","dominio":"ciencias_de_fora","fonte":""},"origem":"wiki","palavras_chave":[],"sinonimos":[],"tipo":"conceito","titulo":"Tectônica de Placas"}
\section{Tectônica de Placas}
A litosfera rígida está fragmentada em grandes placas que se movem sobre a astenosfera, colidindo, afastando-se e deslizando — o que remodela continentes, abre oceanos e ergue montanhas. Síntese dos anos 1960 sobre a deriva continental de Wegener.
\prov{relações: parte\_de geografia\_fisica}
\prov{proveniência: wiki \textperiodcentered{} fato \textperiodcentered{} confiança 1.0 \textperiodcentered{} domínio ciencias\_de\_fora \textperiodcentered{} história 2 versões no jornal}

%CRISTAL {"arestas":[{"alvo":"selecao_natural","confianca":1.0,"epistemico":"fato","origem":"wiki","rel":"usa"},{"alvo":"virtualizacao","confianca":0.5,"epistemico":"fato","origem":"wiki","rel":"relaciona"},{"alvo":"word","confianca":0.5,"epistemico":"fato","origem":"wiki","rel":"relaciona"},{"alvo":"vida-morte","confianca":0.5,"epistemico":"fato","origem":"wiki","rel":"relaciona"}],"confianca":1.0,"contraexemplos":[],"descricao":"Monod: a aparência de propósito na biologia emerge do acaso (mutação) + necessidade (seleção) — não há telos pré-inscrito.","epistemico":"inferencia","exemplos":[],"id":"teleonomia","memoria":"semantica","meta":{"autor":"Monod","dominio":"ciencias_de_fora","fonte":""},"origem":"wiki","palavras_chave":[],"sinonimos":[],"tipo":"conceito","titulo":"Teleonomia"}
\section{Teleonomia}
Monod: a aparência de propósito na biologia emerge do acaso (mutação) + necessidade (seleção) — não há telos pré-inscrito.
\prov{relações: usa selecao\_natural; relaciona virtualizacao; relaciona word; relaciona vida-morte}
\prov{proveniência: wiki \textperiodcentered{} inferencia \textperiodcentered{} confiança 1.0 \textperiodcentered{} domínio ciencias\_de\_fora \textperiodcentered{} história 5 versões no jornal}

%CRISTAL {"arestas":[{"alvo":"consonancia_dissonancia","confianca":1.0,"epistemico":"fato","origem":"wiki","rel":"usa"},{"alvo":"tonalidade","confianca":1.0,"epistemico":"fato","origem":"wiki","rel":"usa"},{"alvo":"teoria_musical","confianca":1.0,"epistemico":"fato","origem":"wiki","rel":"parte_de"}],"confianca":1.0,"contraexemplos":[],"descricao":"Dividir a oitava em 12 semitons iguais (razão 2^(1/12)) — um compromisso que troca a pureza pela liberdade de modular.","epistemico":"fato","exemplos":[],"id":"temperamento_igual","memoria":"semantica","meta":{"dominio":"ciencias_de_fora","fonte":""},"origem":"wiki","palavras_chave":[],"sinonimos":[],"tipo":"conceito","titulo":"Temperamento Igual"}
\section{Temperamento Igual}
Dividir a oitava em 12 semitons iguais (razão 2\^{}(1/12)) — um compromisso que troca a pureza pela liberdade de modular.
\prov{relações: usa consonancia\_dissonancia; usa tonalidade; parte\_de teoria\_musical}
\prov{proveniência: wiki \textperiodcentered{} fato \textperiodcentered{} confiança 1.0 \textperiodcentered{} domínio ciencias\_de\_fora \textperiodcentered{} história 6 versões no jornal}

%CRISTAL {"arestas":[{"alvo":"sintaxe","confianca":1.0,"epistemico":"fato","origem":"wiki","rel":"relaciona"}],"confianca":1.0,"contraexemplos":[],"descricao":"O verbo codifica QUANDO (tempo), COMO se desenrola a ação (aspecto: perfectivo/imperfectivo) e a ATITUDE (modo: indicativo/subjuntivo). Onde as línguas mais divergem.","epistemico":"fato","exemplos":[],"id":"tempo_aspecto_modo","memoria":"semantica","meta":{"dominio":"ciencias_de_fora","fonte":""},"origem":"wiki","palavras_chave":[],"sinonimos":[],"tipo":"conceito","titulo":"Tempo, Aspecto e Modo (TAM)"}
\section{Tempo, Aspecto e Modo (TAM)}
O verbo codifica QUANDO (tempo), COMO se desenrola a ação (aspecto: perfectivo/imperfectivo) e a ATITUDE (modo: indicativo/subjuntivo). Onde as línguas mais divergem.
\prov{relações: relaciona sintaxe}
\prov{proveniência: wiki \textperiodcentered{} fato \textperiodcentered{} confiança 1.0 \textperiodcentered{} domínio ciencias\_de\_fora \textperiodcentered{} história 2 versões no jornal}

%CRISTAL {"arestas":[{"alvo":"teomatematica","confianca":1.0,"epistemico":"fato","origem":"wiki","rel":"usa"},{"alvo":"deus_como_vazio_gentil","confianca":1.0,"epistemico":"fato","origem":"wiki","rel":"usa"},{"alvo":"regua_divina_afastar_retornar","confianca":1.0,"epistemico":"fato","origem":"wiki","rel":"usa"},{"alvo":"cosmologia_quantica","confianca":1.0,"epistemico":"fato","origem":"wiki","rel":"relaciona"}],"confianca":1.0,"contraexemplos":[],"descricao":"Lopes: uma teologia com a identidade Deus≡Origem-de-Tudo≡Vazio, 'cujo instrumento de medida é a régua quântica'; distingue o Absoluto (o conjunto de tudo que existe — dado por definição) de Deus (o que emerge quando o Absoluto passa pela mente). Dissolve problemas clássicos: o dilema de Epicuro sobre o mal não faz sentido para um Deus=Vazio sem desejos nem atributos.","epistemico":"inferencia","exemplos":[],"id":"teologia_quantica_gentil","memoria":"semantica","meta":{"autor":"Gentil Lopes","dominio":"sagrado","fonte":""},"origem":"wiki","palavras_chave":[],"sinonimos":[],"tipo":"conceito","titulo":"Teologia Quântica (Lopes)"}
\section{Teologia Quântica (Lopes)}
Lopes: uma teologia com a identidade Deus\ensuremath{\equiv}Origem-de-Tudo\ensuremath{\equiv}Vazio, 'cujo instrumento de medida é a régua quântica'; distingue o Absoluto (o conjunto de tudo que existe — dado por definição) de Deus (o que emerge quando o Absoluto passa pela mente). Dissolve problemas clássicos: o dilema de Epicuro sobre o mal não faz sentido para um Deus=Vazio sem desejos nem atributos.
\prov{relações: usa teomatematica; usa deus\_como\_vazio\_gentil; usa regua\_divina\_afastar\_retornar; relaciona cosmologia\_quantica}
\prov{proveniência: wiki \textperiodcentered{} inferencia \textperiodcentered{} confiança 1.0 \textperiodcentered{} domínio sagrado \textperiodcentered{} história 6 versões no jornal}

%CRISTAL {"arestas":[{"alvo":"teoria_da_probabilidade","confianca":1.0,"epistemico":"fato","origem":"wiki","rel":"parte_de"},{"alvo":"axiomas_de_kolmogorov","confianca":1.0,"epistemico":"fato","origem":"wiki","rel":"usa"},{"alvo":"bai_orbita_reta_metalica","confianca":1.0,"epistemico":"fato","origem":"wiki","rel":"explica"}],"confianca":1.0,"contraexemplos":[],"descricao":"Atualizar a crença com evidência: P(H|E) ∝ P(E|H)·P(H). A lógica da inferência — a mesma do BAI e da evidência noisy-OR.","epistemico":"fato","exemplos":[],"id":"teorema_de_bayes","memoria":"semantica","meta":{"autor":"Bayes","dominio":"ciencias_de_fora","fonte":""},"origem":"wiki","palavras_chave":[],"sinonimos":[],"tipo":"conceito","titulo":"Teorema de Bayes"}
\section{Teorema de Bayes}
Atualizar a crença com evidência: P(H|E) \ensuremath{\propto} P(E|H)\textperiodcentered P(H). A lógica da inferência — a mesma do BAI e da evidência noisy-OR.
\prov{relações: parte\_de teoria\_da\_probabilidade; usa axiomas\_de\_kolmogorov; explica bai\_orbita\_reta\_metalica}
\prov{proveniência: wiki \textperiodcentered{} fato \textperiodcentered{} confiança 1.0 \textperiodcentered{} domínio ciencias\_de\_fora \textperiodcentered{} história 8 versões no jornal}

%CRISTAL {"arestas":[{"alvo":"direito_e_tecnologia","confianca":1.0,"epistemico":"fato","origem":"wiki","rel":"parte_de"},{"alvo":"escassez","confianca":1.0,"epistemico":"fato","origem":"wiki","rel":"relaciona"},{"alvo":"criptoeconomia","confianca":1.0,"epistemico":"fato","origem":"wiki","rel":"relaciona"}],"confianca":1.0,"contraexemplos":[],"descricao":"Coase: na ausência de custos de transação e com direitos de propriedade bem definidos, as partes negociam até a alocação eficiente, independentemente de quem recebe o direito inicial; com custos altos, a atribuição inicial passa a importar.","epistemico":"fato","exemplos":[],"id":"teorema_de_coase","memoria":"semantica","meta":{"autor":"Coase","dominio":"ciencias_de_fora","fonte":""},"origem":"wiki","palavras_chave":[],"sinonimos":[],"tipo":"conceito","titulo":"Teorema de Coase"}
\section{Teorema de Coase}
Coase: na ausência de custos de transação e com direitos de propriedade bem definidos, as partes negociam até a alocação eficiente, independentemente de quem recebe o direito inicial; com custos altos, a atribuição inicial passa a importar.
\prov{relações: parte\_de direito\_e\_tecnologia; relaciona escassez; relaciona criptoeconomia}
\prov{proveniência: wiki \textperiodcentered{} fato \textperiodcentered{} confiança 1.0 \textperiodcentered{} domínio ciencias\_de\_fora \textperiodcentered{} história 4 versões no jornal}

%CRISTAL {"arestas":[{"alvo":"sunyata","confianca":1.0,"epistemico":"fato","origem":"wiki","rel":"usa"},{"alvo":"tao_da_matematica","confianca":1.0,"epistemico":"fato","origem":"wiki","rel":"relaciona"},{"alvo":"perfeicao_do_vazio_gentil","confianca":1.0,"epistemico":"fato","origem":"wiki","rel":"usa"}],"confianca":1.0,"contraexemplos":[],"descricao":"Lopes confere status de teorema à frase de Nagarjuna ('se eu tivesse qualquer posição teórica teria problemas; como não tenho, não tenho problema') e a demonstra pela perfeição do vazio: fora do vazio há imperfeição, logo toda posição teórica é falível. Base do relativismo epistêmico — não há verdade absoluta em física, matemática ou teologia.","epistemico":"inferencia","exemplos":[],"id":"teorema_de_nagarjuna_gentil","memoria":"semantica","meta":{"autor":"Gentil Lopes","dominio":"sagrado","fonte":""},"origem":"wiki","palavras_chave":[],"sinonimos":[],"tipo":"conceito","titulo":"O Teorema de Nagarjuna (Lopes)"}
\section{O Teorema de Nagarjuna (Lopes)}
Lopes confere status de teorema à frase de Nagarjuna ('se eu tivesse qualquer posição teórica teria problemas; como não tenho, não tenho problema') e a demonstra pela perfeição do vazio: fora do vazio há imperfeição, logo toda posição teórica é falível. Base do relativismo epistêmico — não há verdade absoluta em física, matemática ou teologia.
\prov{relações: usa sunyata; relaciona tao\_da\_matematica; usa perfeicao\_do\_vazio\_gentil}
\prov{proveniência: wiki \textperiodcentered{} inferencia \textperiodcentered{} confiança 1.0 \textperiodcentered{} domínio sagrado \textperiodcentered{} história 5 versões no jornal}

%CRISTAL {"arestas":[{"alvo":"teoria_de_grupos","confianca":1.0,"epistemico":"fato","origem":"wiki","rel":"usa"},{"alvo":"algebra_abstrata","confianca":1.0,"epistemico":"fato","origem":"wiki","rel":"parte_de"},{"alvo":"orbita_z2_3","confianca":1.0,"epistemico":"fato","origem":"wiki","rel":"explica"},{"alvo":"conservacao_parabola_costura","confianca":1.0,"epistemico":"fato","origem":"wiki","rel":"explica"}],"confianca":1.0,"contraexemplos":[],"descricao":"Emmy Noether: cada simetria contínua de um sistema corresponde a uma quantidade CONSERVADA. Simetria ⇔ conservação — a lei a+c²=1.","epistemico":"fato","exemplos":[],"id":"teorema_de_noether","memoria":"semantica","meta":{"autor":"Noether","dominio":"ciencias_de_fora","fonte":""},"origem":"wiki","palavras_chave":[],"sinonimos":[],"tipo":"conceito","titulo":"Teorema de Noether"}
\section{Teorema de Noether}
Emmy Noether: cada simetria contínua de um sistema corresponde a uma quantidade CONSERVADA. Simetria \ensuremath{\Leftrightarrow} conservação — a lei a+c\ensuremath{^{2}}=1.
\prov{relações: usa teoria\_de\_grupos; parte\_de algebra\_abstrata; explica orbita\_z2\_3; explica conservacao\_parabola\_costura}
\prov{proveniência: wiki \textperiodcentered{} fato \textperiodcentered{} confiança 1.0 \textperiodcentered{} domínio ciencias\_de\_fora \textperiodcentered{} história 8 versões no jornal}

%CRISTAL {"arestas":[{"alvo":"incompletude_godel","confianca":1.0,"epistemico":"fato","origem":"wiki","rel":"relaciona"},{"alvo":"simbolo_vs_significado","confianca":1.0,"epistemico":"fato","origem":"wiki","rel":"relaciona"}],"confianca":1.0,"contraexemplos":[],"descricao":"Tarski: a verdade de uma linguagem não é definível dentro dela mesma — o sistema não fala do próprio significado.","epistemico":"fato","exemplos":[],"id":"teorema_de_tarski","memoria":"semantica","meta":{"autor":"Tarski","dominio":"ciencias_de_fora","fonte":""},"origem":"wiki","palavras_chave":[],"sinonimos":[],"tipo":"conceito","titulo":"Indefinibilidade da Verdade"}
\section{Indefinibilidade da Verdade}
Tarski: a verdade de uma linguagem não é definível dentro dela mesma — o sistema não fala do próprio significado.
\prov{relações: relaciona incompletude\_godel; relaciona simbolo\_vs\_significado}
\prov{proveniência: wiki \textperiodcentered{} fato \textperiodcentered{} confiança 1.0 \textperiodcentered{} domínio ciencias\_de\_fora \textperiodcentered{} história 6 versões no jornal}

%CRISTAL {"arestas":[{"alvo":"analise_matematica","confianca":1.0,"epistemico":"fato","origem":"wiki","rel":"parte_de"},{"alvo":"teoria_de_grupos","confianca":1.0,"epistemico":"fato","origem":"wiki","rel":"relaciona"},{"alvo":"hopfield","confianca":1.0,"epistemico":"fato","origem":"wiki","rel":"explica"}],"confianca":1.0,"contraexemplos":[],"descricao":"Sob condições brandas, uma transformação contínua fixa algum ponto (Brouwer, Banach). A raiz de convergência e equilíbrio.","epistemico":"fato","exemplos":[],"id":"teorema_do_ponto_fixo","memoria":"semantica","meta":{"dominio":"ciencias_de_fora","fonte":""},"origem":"wiki","palavras_chave":[],"sinonimos":[],"tipo":"conceito","titulo":"Teorema do Ponto Fixo"}
\section{Teorema do Ponto Fixo}
Sob condições brandas, uma transformação contínua fixa algum ponto (Brouwer, Banach). A raiz de convergência e equilíbrio.
\prov{relações: parte\_de analise\_matematica; relaciona teoria\_de\_grupos; explica hopfield}
\prov{proveniência: wiki \textperiodcentered{} fato \textperiodcentered{} confiança 1.0 \textperiodcentered{} domínio ciencias\_de\_fora \textperiodcentered{} história 9 versões no jornal}

%CRISTAL {"arestas":[{"alvo":"teoria_dos_numeros","confianca":1.0,"epistemico":"fato","origem":"wiki","rel":"parte_de"},{"alvo":"infinitude_dos_primos","confianca":1.0,"epistemico":"fato","origem":"wiki","rel":"especializa"}],"confianca":1.0,"contraexemplos":[],"descricao":"A densidade de primos até x é ~1/ln x — o caos local dos primos esconde uma lei global suave.","epistemico":"fato","exemplos":[],"id":"teorema_dos_numeros_primos","memoria":"semantica","meta":{"dominio":"ciencias_de_fora","fonte":""},"origem":"wiki","palavras_chave":[],"sinonimos":[],"tipo":"conceito","titulo":"Teorema dos Números Primos"}
\section{Teorema dos Números Primos}
A densidade de primos até x é \~{}1/ln x — o caos local dos primos esconde uma lei global suave.
\prov{relações: parte\_de teoria\_dos\_numeros; especializa infinitude\_dos\_primos}
\prov{proveniência: wiki \textperiodcentered{} fato \textperiodcentered{} confiança 1.0 \textperiodcentered{} domínio ciencias\_de\_fora \textperiodcentered{} história 6 versões no jornal}

%CRISTAL {"arestas":[{"alvo":"teoria_dos_numeros","confianca":1.0,"epistemico":"fato","origem":"wiki","rel":"parte_de"}],"confianca":1.0,"contraexemplos":[],"descricao":"Todo inteiro > 1 fatora-se de modo único em primos — os primos são os átomos do número.","epistemico":"fato","exemplos":[],"id":"teorema_fundamental_aritmetica","memoria":"semantica","meta":{"dominio":"ciencias_de_fora","fonte":""},"origem":"wiki","palavras_chave":[],"sinonimos":[],"tipo":"conceito","titulo":"Teorema Fundamental da Aritmética"}
\section{Teorema Fundamental da Aritmética}
Todo inteiro > 1 fatora-se de modo único em primos — os primos são os átomos do número.
\prov{relações: parte\_de teoria\_dos\_numeros}
\prov{proveniência: wiki \textperiodcentered{} fato \textperiodcentered{} confiança 1.0 \textperiodcentered{} domínio ciencias\_de\_fora \textperiodcentered{} história 4 versões no jornal}

%CRISTAL {"arestas":[{"alvo":"teoria_dos_jogos","confianca":1.0,"epistemico":"fato","origem":"wiki","rel":"usa"},{"alvo":"eleitor_mediano_downs","confianca":1.0,"epistemico":"fato","origem":"wiki","rel":"relaciona"},{"alvo":"criptoeconomia","confianca":1.0,"epistemico":"fato","origem":"wiki","rel":"relaciona"},{"alvo":"goldchain","confianca":1.0,"epistemico":"fato","origem":"wiki","rel":"relaciona"}],"confianca":1.0,"contraexemplos":[],"descricao":"Arrow: nenhum método de agregação de preferências ordinais de três ou mais opções pode satisfazer simultaneamente um conjunto mínimo de critérios de racionalidade e justiça (domínio irrestrito, unanimidade, independência das alternativas irrelevantes e não-ditadura). Um limite formal de todo sistema de votação.","epistemico":"fato","exemplos":[],"id":"teorema_impossibilidade_arrow","memoria":"semantica","meta":{"autor":"Arrow","dominio":"ciencias_de_fora","fonte":""},"origem":"wiki","palavras_chave":[],"sinonimos":[],"tipo":"conceito","titulo":"Teorema da Impossibilidade (Arrow)"}
\section{Teorema da Impossibilidade (Arrow)}
Arrow: nenhum método de agregação de preferências ordinais de três ou mais opções pode satisfazer simultaneamente um conjunto mínimo de critérios de racionalidade e justiça (domínio irrestrito, unanimidade, independência das alternativas irrelevantes e não-ditadura). Um limite formal de todo sistema de votação.
\prov{relações: usa teoria\_dos\_jogos; relaciona eleitor\_mediano\_downs; relaciona criptoeconomia; relaciona goldchain}
\prov{proveniência: wiki \textperiodcentered{} fato \textperiodcentered{} confiança 1.0 \textperiodcentered{} domínio ciencias\_de\_fora \textperiodcentered{} história 5 versões no jornal}

%CRISTAL {"arestas":[{"alvo":"filosofia_da_tecnica","confianca":1.0,"epistemico":"fato","origem":"wiki","rel":"parte_de"},{"alvo":"construcao_social_tecnologia","confianca":1.0,"epistemico":"fato","origem":"wiki","rel":"relaciona"},{"alvo":"artefatos_tem_politica","confianca":1.0,"epistemico":"fato","origem":"wiki","rel":"relaciona"},{"alvo":"mediacao_semiotica","confianca":1.0,"epistemico":"fato","origem":"wiki","rel":"relaciona"}],"confianca":1.0,"contraexemplos":[],"descricao":"Latour: o social é uma rede heterogênea de humanos e não-humanos tratados simetricamente como 'actantes'; a agência é distribuída, não exclusiva do sujeito. Normas podem ser delegadas a objetos — o quebra-molas assume o papel do policial.","epistemico":"inferencia","exemplos":[],"id":"teoria_ator_rede","memoria":"semantica","meta":{"autor":"Latour","dominio":"ciencias_de_fora","fonte":""},"origem":"wiki","palavras_chave":[],"sinonimos":[],"tipo":"conceito","titulo":"Teoria Ator-Rede (Latour)"}
\section{Teoria Ator-Rede (Latour)}
Latour: o social é uma rede heterogênea de humanos e não-humanos tratados simetricamente como 'actantes'; a agência é distribuída, não exclusiva do sujeito. Normas podem ser delegadas a objetos — o quebra-molas assume o papel do policial.
\prov{relações: parte\_de filosofia\_da\_tecnica; relaciona construcao\_social\_tecnologia; relaciona artefatos\_tem\_politica; relaciona mediacao\_semiotica}
\prov{proveniência: wiki \textperiodcentered{} inferencia \textperiodcentered{} confiança 1.0 \textperiodcentered{} domínio ciencias\_de\_fora \textperiodcentered{} história 5 versões no jornal}

%CRISTAL {"arestas":[{"alvo":"biologia_celular","confianca":1.0,"epistemico":"fato","origem":"wiki","rel":"parte_de"}],"confianca":1.0,"contraexemplos":[],"descricao":"Todo ser vivo é feito de células, e toda célula vem de outra célula — a unidade e a continuidade da vida.","epistemico":"fato","exemplos":[],"id":"teoria_celular","memoria":"semantica","meta":{"dominio":"ciencias_de_fora","fonte":""},"origem":"wiki","palavras_chave":[],"sinonimos":[],"tipo":"conceito","titulo":"Teoria Celular"}
\section{Teoria Celular}
Todo ser vivo é feito de células, e toda célula vem de outra célula — a unidade e a continuidade da vida.
\prov{relações: parte\_de biologia\_celular}
\prov{proveniência: wiki \textperiodcentered{} fato \textperiodcentered{} confiança 1.0 \textperiodcentered{} domínio ciencias\_de\_fora \textperiodcentered{} história 2 versões no jornal}

%CRISTAL {"arestas":[{"alvo":"filosofia_da_tecnica","confianca":1.0,"epistemico":"fato","origem":"wiki","rel":"parte_de"},{"alvo":"marcuse_racionalidade_tecnologica","confianca":1.0,"epistemico":"fato","origem":"wiki","rel":"usa"},{"alvo":"construcao_social_tecnologia","confianca":1.0,"epistemico":"fato","origem":"wiki","rel":"relaciona"},{"alvo":"goldchain","confianca":1.0,"epistemico":"fato","origem":"wiki","rel":"relaciona"}],"confianca":1.0,"contraexemplos":[],"descricao":"Feenberg: a técnica não é neutra — escolhas de design cristalizam interesses dominantes num 'código técnico'. Como esse código é contingente, pode ser reaberto por 'racionalização democrática': a democratização do design pelos afetados.","epistemico":"inferencia","exemplos":[],"id":"teoria_critica_tecnologia","memoria":"semantica","meta":{"autor":"Feenberg","dominio":"ciencias_de_fora","fonte":""},"origem":"wiki","palavras_chave":[],"sinonimos":[],"tipo":"conceito","titulo":"Teoria Crítica da Tecnologia (Feenberg)"}
\section{Teoria Crítica da Tecnologia (Feenberg)}
Feenberg: a técnica não é neutra — escolhas de design cristalizam interesses dominantes num 'código técnico'. Como esse código é contingente, pode ser reaberto por 'racionalização democrática': a democratização do design pelos afetados.
\prov{relações: parte\_de filosofia\_da\_tecnica; usa marcuse\_racionalidade\_tecnologica; relaciona construcao\_social\_tecnologia; relaciona goldchain}
\prov{proveniência: wiki \textperiodcentered{} inferencia \textperiodcentered{} confiança 1.0 \textperiodcentered{} domínio ciencias\_de\_fora \textperiodcentered{} história 5 versões no jornal}

%CRISTAL {"arestas":[{"alvo":"comunicacao","confianca":1.0,"epistemico":"fato","origem":"wiki","rel":"parte_de"},{"alvo":"word","confianca":0.5,"epistemico":"fato","origem":"wiki","rel":"relaciona"}],"confianca":1.0,"contraexemplos":[],"descricao":"Os modelos do processo comunicativo e as teorias sobre os efeitos da mídia — de esquemas lineares (emissor-mensagem-receptor) a efeitos cognitivos cumulativos.","epistemico":"fato","exemplos":[],"id":"teoria_da_comunicacao","memoria":"semantica","meta":{"dominio":"ciencias_de_fora","fonte":""},"origem":"wiki","palavras_chave":[],"sinonimos":[],"tipo":"conceito","titulo":"Teoria da Comunicação"}
\section{Teoria da Comunicação}
Os modelos do processo comunicativo e as teorias sobre os efeitos da mídia — de esquemas lineares (emissor-mensagem-receptor) a efeitos cognitivos cumulativos.
\prov{relações: parte\_de comunicacao; relaciona word}
\prov{proveniência: wiki \textperiodcentered{} fato \textperiodcentered{} confiança 1.0 \textperiodcentered{} domínio ciencias\_de\_fora \textperiodcentered{} história 3 versões no jornal}

%CRISTAL {"arestas":[{"alvo":"charles_darwin","confianca":1.0,"epistemico":"fato","origem":"wiki","rel":"parte_de"},{"alvo":"selecao_natural","confianca":1.0,"epistemico":"fato","origem":"wiki","rel":"usa"},{"alvo":"biologia_evolutiva","confianca":1.0,"epistemico":"fato","origem":"wiki","rel":"parte_de"}],"confianca":1.0,"contraexemplos":[],"descricao":"A vida atual descende de ancestrais comuns, modificada ao longo do tempo. O fato central da biologia.","epistemico":"fato","exemplos":[],"id":"teoria_da_evolucao","memoria":"semantica","meta":{"autor":"Darwin","dominio":"ciencias_de_fora","fonte":""},"origem":"wiki","palavras_chave":[],"sinonimos":[],"tipo":"conceito","titulo":"Teoria da Evolução"}
\section{Teoria da Evolução}
A vida atual descende de ancestrais comuns, modificada ao longo do tempo. O fato central da biologia.
\prov{relações: parte\_de charles\_darwin; usa selecao\_natural; parte\_de biologia\_evolutiva}
\prov{proveniência: wiki \textperiodcentered{} fato \textperiodcentered{} confiança 1.0 \textperiodcentered{} domínio ciencias\_de\_fora \textperiodcentered{} história 4 versões no jornal}

%CRISTAL {"arestas":[{"alvo":"teoria_da_probabilidade","confianca":1.0,"epistemico":"fato","origem":"wiki","rel":"usa"},{"alvo":"teoria_da_informacao_shannon","confianca":1.0,"epistemico":"fato","origem":"wiki","rel":"relaciona"}],"confianca":1.0,"contraexemplos":[],"descricao":"Shannon: informação como redução de incerteza, medida em bits; entropia = surpresa média. O limite de toda compressão e canal.","epistemico":"fato","exemplos":[],"id":"teoria_da_informacao","memoria":"semantica","meta":{"autor":"Shannon","dominio":"ciencias_de_fora","fonte":""},"origem":"wiki","palavras_chave":[],"sinonimos":[],"tipo":"conceito","titulo":"Teoria da Informação"}
\section{Teoria da Informação}
Shannon: informação como redução de incerteza, medida em bits; entropia = surpresa média. O limite de toda compressão e canal.
\prov{relações: usa teoria\_da\_probabilidade; relaciona teoria\_da\_informacao\_shannon}
\prov{proveniência: wiki \textperiodcentered{} fato \textperiodcentered{} confiança 1.0 \textperiodcentered{} domínio ciencias\_de\_fora \textperiodcentered{} história 5 versões no jornal}

%CRISTAL {"arestas":[{"alvo":"filosofia_da_computacao","confianca":1.0,"epistemico":"fato","origem":"wiki","rel":"parte_de"},{"alvo":"colapso","confianca":1.0,"epistemico":"fato","origem":"wiki","rel":"relaciona"},{"alvo":"olho_de_koch","confianca":1.0,"epistemico":"fato","origem":"wiki","rel":"relaciona"}],"confianca":1.0,"contraexemplos":[],"descricao":"Shannon: a informação de uma fonte se quantifica pela entropia (incerteza média); comunicar é reduzir incerteza no receptor. O previsível informa pouco, o improvável muito — a informação é proporcional à surpresa dissolvida.","epistemico":"fato","exemplos":[],"id":"teoria_da_informacao_shannon","memoria":"semantica","meta":{"autor":"Shannon","dominio":"ciencias_de_fora","fonte":""},"origem":"wiki","palavras_chave":[],"sinonimos":[],"tipo":"conceito","titulo":"Teoria da Informação (Shannon)"}
\section{Teoria da Informação (Shannon)}
Shannon: a informação de uma fonte se quantifica pela entropia (incerteza média); comunicar é reduzir incerteza no receptor. O previsível informa pouco, o improvável muito — a informação é proporcional à surpresa dissolvida.
\prov{relações: parte\_de filosofia\_da\_computacao; relaciona colapso; relaciona olho\_de\_koch}
\prov{proveniência: wiki \textperiodcentered{} fato \textperiodcentered{} confiança 1.0 \textperiodcentered{} domínio ciencias\_de\_fora \textperiodcentered{} história 4 versões no jornal}

%CRISTAL {"arestas":[{"alvo":"matematica","confianca":1.0,"epistemico":"fato","origem":"wiki","rel":"parte_de"}],"confianca":1.0,"contraexemplos":[],"descricao":"A matemática do acaso e da incerteza — medida, esperança, inferência sob risco.","epistemico":"fato","exemplos":[],"id":"teoria_da_probabilidade","memoria":"semantica","meta":{"dominio":"ciencias_de_fora","fonte":""},"origem":"wiki","palavras_chave":[],"sinonimos":[],"tipo":"conceito","titulo":"Teoria da Probabilidade"}
\section{Teoria da Probabilidade}
A matemática do acaso e da incerteza — medida, esperança, inferência sob risco.
\prov{relações: parte\_de matematica}
\prov{proveniência: wiki \textperiodcentered{} fato \textperiodcentered{} confiança 1.0 \textperiodcentered{} domínio ciencias\_de\_fora \textperiodcentered{} história 4 versões no jornal}

%CRISTAL {"arestas":[{"alvo":"artes_visuais","confianca":1.0,"epistemico":"fato","origem":"wiki","rel":"parte_de"}],"confianca":1.0,"contraexemplos":[],"descricao":"Como as cores se relacionam e se misturam — complementares, quentes/frias; da roda de Newton a Goethe e Itten.","epistemico":"fato","exemplos":[],"id":"teoria_das_cores","memoria":"semantica","meta":{"dominio":"ciencias_de_fora","fonte":""},"origem":"wiki","palavras_chave":[],"sinonimos":[],"tipo":"conceito","titulo":"Teoria das Cores"}
\section{Teoria das Cores}
Como as cores se relacionam e se misturam — complementares, quentes/frias; da roda de Newton a Goethe e Itten.
\prov{relações: parte\_de artes\_visuais}
\prov{proveniência: wiki \textperiodcentered{} fato \textperiodcentered{} confiança 1.0 \textperiodcentered{} domínio ciencias\_de\_fora \textperiodcentered{} história 2 versões no jornal}

%CRISTAL {"arestas":[{"alvo":"ciencia_politica","confianca":1.0,"epistemico":"fato","origem":"wiki","rel":"parte_de"},{"alvo":"lei_de_ferro_oligarquia","confianca":1.0,"epistemico":"fato","origem":"wiki","rel":"relaciona"},{"alvo":"pluralismo_elitismo","confianca":1.0,"epistemico":"fato","origem":"wiki","rel":"contradiz"}],"confianca":1.0,"contraexemplos":[],"descricao":"Em toda sociedade uma minoria organizada governa a maioria desorganizada (Mosca), e as elites governantes decaem e são continuamente substituídas — a história como 'cemitério de aristocracias' (a circulação das elites de Pareto).","epistemico":"inferencia","exemplos":[],"id":"teoria_das_elites","memoria":"semantica","meta":{"autor":"Mosca/Pareto","dominio":"ciencias_de_fora","fonte":""},"origem":"wiki","palavras_chave":[],"sinonimos":[],"tipo":"conceito","titulo":"Teoria das Elites (Mosca, Pareto)"}
\section{Teoria das Elites (Mosca, Pareto)}
Em toda sociedade uma minoria organizada governa a maioria desorganizada (Mosca), e as elites governantes decaem e são continuamente substituídas — a história como 'cemitério de aristocracias' (a circulação das elites de Pareto).
\prov{relações: parte\_de ciencia\_politica; relaciona lei\_de\_ferro\_oligarquia; contradiz pluralismo\_elitismo}
\prov{proveniência: wiki \textperiodcentered{} inferencia \textperiodcentered{} confiança 1.0 \textperiodcentered{} domínio ciencias\_de\_fora \textperiodcentered{} história 4 versões no jornal}

%CRISTAL {"arestas":[{"alvo":"teoria_de_grupos","confianca":1.0,"epistemico":"fato","origem":"wiki","rel":"usa"},{"alvo":"algebra_abstrata","confianca":1.0,"epistemico":"fato","origem":"wiki","rel":"parte_de"},{"alvo":"conjugados_galois_parabola","confianca":1.0,"epistemico":"fato","origem":"wiki","rel":"explica"},{"alvo":"transformada_metalica","confianca":1.0,"epistemico":"fato","origem":"wiki","rel":"explica"}],"confianca":1.0,"contraexemplos":[],"descricao":"Galois: liga a solubilidade de uma equação ao GRUPO de simetria de suas raízes (os conjugados). A raiz da transformada metálica.","epistemico":"fato","exemplos":[],"id":"teoria_de_galois","memoria":"semantica","meta":{"autor":"Galois","dominio":"ciencias_de_fora","fonte":""},"origem":"wiki","palavras_chave":[],"sinonimos":[],"tipo":"conceito","titulo":"Teoria de Galois"}
\section{Teoria de Galois}
Galois: liga a solubilidade de uma equação ao GRUPO de simetria de suas raízes (os conjugados). A raiz da transformada metálica.
\prov{relações: usa teoria\_de\_grupos; parte\_de algebra\_abstrata; explica conjugados\_galois\_parabola; explica transformada\_metalica}
\prov{proveniência: wiki \textperiodcentered{} fato \textperiodcentered{} confiança 1.0 \textperiodcentered{} domínio ciencias\_de\_fora \textperiodcentered{} história 10 versões no jornal}

%CRISTAL {"arestas":[{"alvo":"algebra_abstrata","confianca":1.0,"epistemico":"fato","origem":"wiki","rel":"parte_de"},{"alvo":"grupos","confianca":1.0,"epistemico":"fato","origem":"wiki","rel":"relaciona"}],"confianca":1.0,"contraexemplos":[],"descricao":"A álgebra da simetria: um conjunto com uma operação associativa, identidade e inversos. A estrutura de toda transformação.","epistemico":"fato","exemplos":[],"id":"teoria_de_grupos","memoria":"semantica","meta":{"dominio":"ciencias_de_fora","fonte":""},"origem":"wiki","palavras_chave":[],"sinonimos":[],"tipo":"conceito","titulo":"Teoria de Grupos"}
\section{Teoria de Grupos}
A álgebra da simetria: um conjunto com uma operação associativa, identidade e inversos. A estrutura de toda transformação.
\prov{relações: parte\_de algebra\_abstrata; relaciona grupos}
\prov{proveniência: wiki \textperiodcentered{} fato \textperiodcentered{} confiança 1.0 \textperiodcentered{} domínio ciencias\_de\_fora \textperiodcentered{} história 5 versões no jornal}

%CRISTAL {"arestas":[{"alvo":"teoria_da_comunicacao","confianca":1.0,"epistemico":"fato","origem":"wiki","rel":"parte_de"},{"alvo":"dois_sistemas_kahneman","confianca":1.0,"epistemico":"fato","origem":"wiki","rel":"relaciona"},{"alvo":"agenda_setting","confianca":1.0,"epistemico":"fato","origem":"wiki","rel":"relaciona"}],"confianca":1.0,"contraexemplos":[],"descricao":"Gerbner: a exposição prolongada e cumulativa à TV 'cultiva' uma percepção da realidade alinhada ao mundo televisivo; a super-representação de violência gera a 'síndrome do mundo mau'.","epistemico":"fato","exemplos":[],"id":"teoria_do_cultivo","memoria":"semantica","meta":{"autor":"Gerbner","dominio":"ciencias_de_fora","fonte":""},"origem":"wiki","palavras_chave":[],"sinonimos":[],"tipo":"conceito","titulo":"Teoria do Cultivo (Gerbner)"}
\section{Teoria do Cultivo (Gerbner)}
Gerbner: a exposição prolongada e cumulativa à TV 'cultiva' uma percepção da realidade alinhada ao mundo televisivo; a super-representação de violência gera a 'síndrome do mundo mau'.
\prov{relações: parte\_de teoria\_da\_comunicacao; relaciona dois\_sistemas\_kahneman; relaciona agenda\_setting}
\prov{proveniência: wiki \textperiodcentered{} fato \textperiodcentered{} confiança 1.0 \textperiodcentered{} domínio ciencias\_de\_fora \textperiodcentered{} história 4 versões no jornal}

%CRISTAL {"arestas":[{"alvo":"economia","confianca":1.0,"epistemico":"fato","origem":"wiki","rel":"parte_de"},{"alvo":"utilitarismo","confianca":0.5,"epistemico":"fato","origem":"wiki","rel":"relaciona"},{"alvo":"vida-morte","confianca":0.5,"epistemico":"fato","origem":"wiki","rel":"relaciona"},{"alvo":"toryces_e_helyces","confianca":0.5,"epistemico":"fato","origem":"wiki","rel":"relaciona"},{"alvo":"transcricao","confianca":0.5,"epistemico":"fato","origem":"wiki","rel":"relaciona"},{"alvo":"tipos_de_interpretante","confianca":0.5,"epistemico":"fato","origem":"wiki","rel":"relaciona"}],"confianca":1.0,"contraexemplos":[],"descricao":"A pergunta sobre a origem do valor econômico — trabalho, utilidade, escassez ou convenção?","epistemico":"fato","exemplos":[],"id":"teoria_do_valor","memoria":"semantica","meta":{"dominio":"ciencias_de_fora","fonte":""},"origem":"wiki","palavras_chave":[],"sinonimos":[],"tipo":"conceito","titulo":"Teoria do Valor"}
\section{Teoria do Valor}
A pergunta sobre a origem do valor econômico — trabalho, utilidade, escassez ou convenção?
\prov{relações: parte\_de economia; relaciona utilitarismo; relaciona vida-morte; relaciona toryces\_e\_helyces; relaciona transcricao; relaciona tipos\_de\_interpretante}
\prov{proveniência: wiki \textperiodcentered{} fato \textperiodcentered{} confiança 1.0 \textperiodcentered{} domínio ciencias\_de\_fora \textperiodcentered{} história 7 versões no jornal}

%CRISTAL {"arestas":[{"alvo":"matematica","confianca":1.0,"epistemico":"fato","origem":"wiki","rel":"parte_de"}],"confianca":1.0,"contraexemplos":[],"descricao":"O fundamento comum: toda a matemática construída sobre pertinência e conjunto (Cantor, ZFC).","epistemico":"fato","exemplos":[],"id":"teoria_dos_conjuntos","memoria":"semantica","meta":{"dominio":"ciencias_de_fora","fonte":""},"origem":"wiki","palavras_chave":[],"sinonimos":[],"tipo":"conceito","titulo":"Teoria dos Conjuntos"}
\section{Teoria dos Conjuntos}
O fundamento comum: toda a matemática construída sobre pertinência e conjunto (Cantor, ZFC).
\prov{relações: parte\_de matematica}
\prov{proveniência: wiki \textperiodcentered{} fato \textperiodcentered{} confiança 1.0 \textperiodcentered{} domínio ciencias\_de\_fora \textperiodcentered{} história 4 versões no jornal}

%CRISTAL {"arestas":[{"alvo":"economia","confianca":1.0,"epistemico":"fato","origem":"wiki","rel":"parte_de"},{"alvo":"htlc","confianca":1.0,"epistemico":"fato","origem":"wiki","rel":"explica"},{"alvo":"word","confianca":0.5,"epistemico":"fato","origem":"wiki","rel":"relaciona"},{"alvo":"aprendizado_por_reforco","confianca":1.0,"epistemico":"fato","origem":"wiki","rel":"relaciona"}],"confianca":1.0,"contraexemplos":[],"descricao":"Nash: no equilíbrio, nenhum jogador melhora mudando sozinho; o dilema do prisioneiro mostra a divergência entre a razão individual e o ótimo coletivo.","epistemico":"fato","exemplos":[],"id":"teoria_dos_jogos","memoria":"semantica","meta":{"autor":"Nash","dominio":"ciencias_de_fora","fonte":""},"origem":"wiki","palavras_chave":[],"sinonimos":[],"tipo":"conceito","titulo":"Teoria dos Jogos"}
\section{Teoria dos Jogos}
Nash: no equilíbrio, nenhum jogador melhora mudando sozinho; o dilema do prisioneiro mostra a divergência entre a razão individual e o ótimo coletivo.
\prov{relações: parte\_de economia; explica htlc; relaciona word; relaciona aprendizado\_por\_reforco}
\prov{proveniência: wiki \textperiodcentered{} fato \textperiodcentered{} confiança 1.0 \textperiodcentered{} domínio ciencias\_de\_fora \textperiodcentered{} história 5 versões no jornal}

%CRISTAL {"arestas":[{"alvo":"inteiros","confianca":1.0,"epistemico":"fato","origem":"gentil/Gentil_Lopes_FUNDAMENTOS_DOS_NUMEROS_pdf.pdf","rel":"parte_de"},{"alvo":"ula_da_broca_arquitetura_minima","confianca":0.5,"epistemico":"fato","origem":"wiki","rel":"relaciona"},{"alvo":"word","confianca":0.5,"epistemico":"fato","origem":"wiki","rel":"relaciona"},{"alvo":"vontade_de_poder","confianca":0.5,"epistemico":"fato","origem":"wiki","rel":"relaciona"},{"alvo":"while","confianca":0.5,"epistemico":"fato","origem":"wiki","rel":"relaciona"},{"alvo":"matematica","confianca":1.0,"epistemico":"fato","origem":"wiki","rel":"parte_de"},{"alvo":"numero","confianca":1.0,"epistemico":"fato","origem":"wiki","rel":"referencia"}],"confianca":1.0,"contraexemplos":[],"descricao":"O estudo dos inteiros e suas relações — primos, divisibilidade, equações diofantinas; a 'rainha da matemática' (Gauss).","epistemico":"fato","exemplos":[],"id":"teoria_dos_numeros","memoria":"semantica","meta":{"dominio":"ciencias_de_fora","fonte":""},"origem":"wiki","palavras_chave":[],"sinonimos":[],"tipo":"conceito","titulo":"Teoria dos Números"}
\section{Teoria dos Números}
O estudo dos inteiros e suas relações — primos, divisibilidade, equações diofantinas; a 'rainha da matemática' (Gauss).
\prov{relações: parte\_de inteiros; relaciona ula\_da\_broca\_arquitetura\_minima; relaciona word; relaciona vontade\_de\_poder; relaciona while; parte\_de matematica; referencia numero}
\prov{proveniência: wiki \textperiodcentered{} fato \textperiodcentered{} confiança 1.0 \textperiodcentered{} domínio ciencias\_de\_fora \textperiodcentered{} história 28 versões no jornal}

%CRISTAL {"arestas":[{"alvo":"mitocondria","confianca":1.0,"epistemico":"fato","origem":"wiki","rel":"explica"},{"alvo":"simbiose","confianca":1.0,"epistemico":"fato","origem":"wiki","rel":"usa"}],"confianca":1.0,"contraexemplos":[],"descricao":"Margulis: mitocôndrias e cloroplastos foram bactérias engolidas que viraram organelas — cooperação, não só competição.","epistemico":"fato","exemplos":[],"id":"teoria_endossimbiotica","memoria":"semantica","meta":{"autor":"Margulis","dominio":"ciencias_de_fora","fonte":""},"origem":"wiki","palavras_chave":[],"sinonimos":[],"tipo":"conceito","titulo":"Teoria Endossimbiótica"}
\section{Teoria Endossimbiótica}
Margulis: mitocôndrias e cloroplastos foram bactérias engolidas que viraram organelas — cooperação, não só competição.
\prov{relações: explica mitocondria; usa simbiose}
\prov{proveniência: wiki \textperiodcentered{} fato \textperiodcentered{} confiança 1.0 \textperiodcentered{} domínio ciencias\_de\_fora \textperiodcentered{} história 3 versões no jornal}

%CRISTAL {"arestas":[{"alvo":"teoria_da_comunicacao","confianca":1.0,"epistemico":"fato","origem":"wiki","rel":"parte_de"}],"confianca":1.0,"contraexemplos":[],"descricao":"Modelo de efeitos diretos e poderosos: a mensagem seria 'injetada' numa audiência passiva e homogênea, produzindo resposta imediata e uniforme. Superada por carecer de base empírica e ignorar a interpretação diferenciada dos receptores.","epistemico":"hipotese","exemplos":[],"id":"teoria_hipodermica","memoria":"semantica","meta":{"autor":"—","dominio":"ciencias_de_fora","fonte":""},"origem":"wiki","palavras_chave":[],"sinonimos":[],"tipo":"conceito","titulo":"Teoria Hipodérmica / Bala Mágica"}
\section{Teoria Hipodérmica / Bala Mágica}
Modelo de efeitos diretos e poderosos: a mensagem seria 'injetada' numa audiência passiva e homogênea, produzindo resposta imediata e uniforme. Superada por carecer de base empírica e ignorar a interpretação diferenciada dos receptores.
\prov{relações: parte\_de teoria\_da\_comunicacao}
\prov{proveniência: wiki \textperiodcentered{} hipotese \textperiodcentered{} confiança 1.0 \textperiodcentered{} domínio ciencias\_de\_fora \textperiodcentered{} história 2 versões no jornal}

%CRISTAL {"arestas":[{"alvo":"artes_visuais","confianca":1.0,"epistemico":"fato","origem":"wiki","rel":"parte_de"},{"alvo":"mao_invisivel","confianca":1.0,"epistemico":"fato","origem":"wiki","rel":"relaciona"}],"confianca":1.0,"contraexemplos":[],"descricao":"Dickie: a arte não se define por qualidades intrínsecas do objeto, mas por um status conferido — obra é um artefato apresentado a um público do 'mundo da arte' por quem age em seu nome. Definição classificatória, não avaliativa.","epistemico":"inferencia","exemplos":[],"id":"teoria_institucional_dickie","memoria":"semantica","meta":{"autor":"Dickie","dominio":"ciencias_de_fora","fonte":""},"origem":"wiki","palavras_chave":[],"sinonimos":[],"tipo":"conceito","titulo":"Teoria Institucional da Arte (Dickie)"}
\section{Teoria Institucional da Arte (Dickie)}
Dickie: a arte não se define por qualidades intrínsecas do objeto, mas por um status conferido — obra é um artefato apresentado a um público do 'mundo da arte' por quem age em seu nome. Definição classificatória, não avaliativa.
\prov{relações: parte\_de artes\_visuais; relaciona mao\_invisivel}
\prov{proveniência: wiki \textperiodcentered{} inferencia \textperiodcentered{} confiança 1.0 \textperiodcentered{} domínio ciencias\_de\_fora \textperiodcentered{} história 3 versões no jornal}

%CRISTAL {"arestas":[{"alvo":"literatura","confianca":1.0,"epistemico":"fato","origem":"wiki","rel":"parte_de"},{"alvo":"arte","confianca":1.0,"epistemico":"fato","origem":"wiki","rel":"parte_de"}],"confianca":1.0,"contraexemplos":[],"descricao":"O estudo do que faz um texto ser literatura — gênero, forma, narrativa, interpretação.","epistemico":"fato","exemplos":[],"id":"teoria_literaria","memoria":"semantica","meta":{"dominio":"ciencias_de_fora","fonte":""},"origem":"wiki","palavras_chave":[],"sinonimos":[],"tipo":"conceito","titulo":"Teoria Literária"}
\section{Teoria Literária}
O estudo do que faz um texto ser literatura — gênero, forma, narrativa, interpretação.
\prov{relações: parte\_de literatura; parte\_de arte}
\prov{proveniência: wiki \textperiodcentered{} fato \textperiodcentered{} confiança 1.0 \textperiodcentered{} domínio ciencias\_de\_fora \textperiodcentered{} história 3 versões no jornal}

%CRISTAL {"arestas":[{"alvo":"musica","confianca":1.0,"epistemico":"fato","origem":"wiki","rel":"parte_de"},{"alvo":"arte","confianca":1.0,"epistemico":"fato","origem":"wiki","rel":"parte_de"}],"confianca":1.0,"contraexemplos":[],"descricao":"O estudo de como a música se organiza — altura, ritmo, harmonia, forma.","epistemico":"fato","exemplos":[],"id":"teoria_musical","memoria":"semantica","meta":{"dominio":"ciencias_de_fora","fonte":""},"origem":"wiki","palavras_chave":[],"sinonimos":[],"tipo":"conceito","titulo":"Teoria Musical"}
\section{Teoria Musical}
O estudo de como a música se organiza — altura, ritmo, harmonia, forma.
\prov{relações: parte\_de musica; parte\_de arte}
\prov{proveniência: wiki \textperiodcentered{} fato \textperiodcentered{} confiança 1.0 \textperiodcentered{} domínio ciencias\_de\_fora \textperiodcentered{} história 3 versões no jornal}

%CRISTAL {"arestas":[{"alvo":"psicologia","confianca":1.0,"epistemico":"fato","origem":"wiki","rel":"parte_de"},{"alvo":"educacao","confianca":1.0,"epistemico":"fato","origem":"wiki","rel":"relaciona"},{"alvo":"word","confianca":0.5,"epistemico":"fato","origem":"wiki","rel":"relaciona"},{"alvo":"vida-morte","confianca":0.5,"epistemico":"fato","origem":"wiki","rel":"relaciona"},{"alvo":"virtualizacao","confianca":0.5,"epistemico":"fato","origem":"wiki","rel":"relaciona"}],"confianca":1.0,"contraexemplos":[],"descricao":"O estudo de como o conhecimento é adquirido, retido e transformado pela mente — do behaviorismo ao construtivismo.","epistemico":"fato","exemplos":[],"id":"teorias_aprendizagem","memoria":"semantica","meta":{"dominio":"ciencias_de_fora","fonte":""},"origem":"wiki","palavras_chave":[],"sinonimos":[],"tipo":"conceito","titulo":"Teorias da Aprendizagem"}
\section{Teorias da Aprendizagem}
O estudo de como o conhecimento é adquirido, retido e transformado pela mente — do behaviorismo ao construtivismo.
\prov{relações: parte\_de psicologia; relaciona educacao; relaciona word; relaciona vida-morte; relaciona virtualizacao}
\prov{proveniência: wiki \textperiodcentered{} fato \textperiodcentered{} confiança 1.0 \textperiodcentered{} domínio ciencias\_de\_fora \textperiodcentered{} história 6 versões no jornal}

%CRISTAL {"arestas":[{"alvo":"geografia_humana","confianca":1.0,"epistemico":"fato","origem":"wiki","rel":"parte_de"},{"alvo":"biopoder","confianca":1.0,"epistemico":"fato","origem":"wiki","rel":"usa"},{"alvo":"contrato_social","confianca":1.0,"epistemico":"fato","origem":"wiki","rel":"relaciona"}],"confianca":1.0,"contraexemplos":[],"descricao":"Robert Sack: a estratégia espacial de afetar, influenciar ou controlar pessoas, recursos e relações por meio do controle de uma área delimitada — não é instinto, mas construção social e política deliberada (quem decide o que ocorre na área).","epistemico":"inferencia","exemplos":[],"id":"territorialidade_sack","memoria":"semantica","meta":{"autor":"Sack","dominio":"ciencias_de_fora","fonte":""},"origem":"wiki","palavras_chave":[],"sinonimos":[],"tipo":"conceito","titulo":"Territorialidade Humana (Sack)"}
\section{Territorialidade Humana (Sack)}
Robert Sack: a estratégia espacial de afetar, influenciar ou controlar pessoas, recursos e relações por meio do controle de uma área delimitada — não é instinto, mas construção social e política deliberada (quem decide o que ocorre na área).
\prov{relações: parte\_de geografia\_humana; usa biopoder; relaciona contrato\_social}
\prov{proveniência: wiki \textperiodcentered{} inferencia \textperiodcentered{} confiança 1.0 \textperiodcentered{} domínio ciencias\_de\_fora \textperiodcentered{} história 4 versões no jornal}

%CRISTAL {"arestas":[{"alvo":"filosofia_da_computacao","confianca":1.0,"epistemico":"fato","origem":"wiki","rel":"parte_de"},{"alvo":"godel_incompletude","confianca":1.0,"epistemico":"fato","origem":"wiki","rel":"relaciona"},{"alvo":"maquina_de_turing","confianca":1.0,"epistemico":"fato","origem":"wiki","rel":"usa"},{"alvo":"calculo_lambda","confianca":1.0,"epistemico":"fato","origem":"wiki","rel":"usa"},{"alvo":"turing_halting_church","confianca":1.0,"epistemico":"fato","origem":"wiki","rel":"relaciona"}],"confianca":1.0,"contraexemplos":[],"descricao":"Toda função 'efetivamente calculável' é Turing-computável. Tese, não teorema — não se demonstra, só se corrobora.","epistemico":"inferencia","exemplos":[],"id":"tese_church_turing","memoria":"semantica","meta":{"autor":"Church/Turing","dominio":"ciencias_de_fora","fonte":""},"origem":"wiki","palavras_chave":[],"sinonimos":[],"tipo":"conceito","titulo":"Tese de Church-Turing"}
\section{Tese de Church-Turing}
Toda função 'efetivamente calculável' é Turing-computável. Tese, não teorema — não se demonstra, só se corrobora.
\prov{relações: parte\_de filosofia\_da\_computacao; relaciona godel\_incompletude; usa maquina\_de\_turing; usa calculo\_lambda; relaciona turing\_halting\_church}
\prov{proveniência: wiki \textperiodcentered{} inferencia \textperiodcentered{} confiança 1.0 \textperiodcentered{} domínio ciencias\_de\_fora \textperiodcentered{} história 11 versões no jornal}

%CRISTAL {"arestas":[{"alvo":"filosofia_da_computacao","confianca":1.0,"epistemico":"fato","origem":"wiki","rel":"parte_de"},{"alvo":"maquina_de_turing","confianca":1.0,"epistemico":"fato","origem":"wiki","rel":"usa"},{"alvo":"simbolo_vs_significado","confianca":1.0,"epistemico":"fato","origem":"wiki","rel":"relaciona"},{"alvo":"consciencia","confianca":1.0,"epistemico":"fato","origem":"wiki","rel":"relaciona"},{"alvo":"tiffany","confianca":1.0,"epistemico":"fato","origem":"wiki","rel":"relaciona"}],"confianca":1.0,"contraexemplos":[],"descricao":"Turing: em vez de 'as máquinas podem pensar?', um teste comportamental — uma máquina pensa se, num diálogo textual, for indistinguível de um humano para um interrogador. Critério proposto, não teorema.","epistemico":"inferencia","exemplos":[],"id":"teste_de_turing","memoria":"semantica","meta":{"autor":"Turing","dominio":"ciencias_de_fora","fonte":""},"origem":"wiki","palavras_chave":[],"sinonimos":[],"tipo":"conceito","titulo":"Teste de Turing"}
\section{Teste de Turing}
Turing: em vez de 'as máquinas podem pensar?', um teste comportamental — uma máquina pensa se, num diálogo textual, for indistinguível de um humano para um interrogador. Critério proposto, não teorema.
\prov{relações: parte\_de filosofia\_da\_computacao; usa maquina\_de\_turing; relaciona simbolo\_vs\_significado; relaciona consciencia; relaciona tiffany}
\prov{proveniência: wiki \textperiodcentered{} inferencia \textperiodcentered{} confiança 1.0 \textperiodcentered{} domínio ciencias\_de\_fora \textperiodcentered{} história 6 versões no jornal}

%CRISTAL {"arestas":[{"alvo":"teorias_aprendizagem","confianca":1.0,"epistemico":"fato","origem":"wiki","rel":"parte_de"},{"alvo":"spaced_repetition","confianca":1.0,"epistemico":"fato","origem":"wiki","rel":"relaciona"}],"confianca":1.0,"contraexemplos":[],"descricao":"Roediger & Karpicke: recuperar ativamente a informação da memória (testar-se) produz retenção de longo prazo muito superior à releitura passiva. Um dos achados mais replicados da ciência da aprendizagem.","epistemico":"fato","exemplos":[],"id":"testing_effect","memoria":"semantica","meta":{"autor":"Roediger & Karpicke","dominio":"ciencias_de_fora","fonte":""},"origem":"wiki","palavras_chave":[],"sinonimos":[],"tipo":"conceito","titulo":"Efeito do Teste / Prática de Recuperação"}
\section{Efeito do Teste / Prática de Recuperação}
Roediger \& Karpicke: recuperar ativamente a informação da memória (testar-se) produz retenção de longo prazo muito superior à releitura passiva. Um dos achados mais replicados da ciência da aprendizagem.
\prov{relações: parte\_de teorias\_aprendizagem; relaciona spaced\_repetition}
\prov{proveniência: wiki \textperiodcentered{} fato \textperiodcentered{} confiança 1.0 \textperiodcentered{} domínio ciencias\_de\_fora \textperiodcentered{} história 3 versões no jornal}

%CRISTAL {"arestas":[{"alvo":"orbitas_de_fala","confianca":1.0,"epistemico":"fato","origem":"wiki","rel":"usa"},{"alvo":"lib_tiffany_no_metal","confianca":1.0,"epistemico":"fato","origem":"wiki","rel":"parte_de"},{"alvo":"bai_psi_colapso","confianca":1.0,"epistemico":"fato","origem":"wiki","rel":"usa"},{"alvo":"atuadores_da_face","confianca":1.0,"epistemico":"fato","origem":"wiki","rel":"usa"}],"confianca":1.0,"contraexemplos":[],"descricao":"Fecha escutar→pensar→falar com o rosto. A Tiffany ganhou os verbos DIZER e FALAR (linguagem/tiffany.py, 9/9): DIZER colapsa Ψ sobre os candidatos (o que dizer, argmax — o pensar); FALAR converte a frase escolhida em VISEMAS + a linha do tempo dos músculos da boca (o lip-sync — o rosto). O que a Tiffany diz vira, literalmente, a boca sincronizada sobre o scan real. Falar DISSIPA entropia (cada visema um colapso Ψ), e a órbita da fala fecha a volta no silêncio. Verif.: dizer([('ola…',5),('tchau',2)]) colapsa em 'ola…' e o rosto o fala (132 quadros de boca).","epistemico":"fato","exemplos":[],"id":"tiffany_fala_o_rosto","memoria":"semantica","meta":{"dominio":"biologia","fonte":"órbitas de fala"},"origem":"wiki","palavras_chave":[],"sinonimos":[],"tipo":"conceito","titulo":"A Tiffany Fala (o que ela diz vira visemas)"}
\section{A Tiffany Fala (o que ela diz vira visemas)}
Fecha escutar\ensuremath{\to}pensar\ensuremath{\to}falar com o rosto. A Tiffany ganhou os verbos DIZER e FALAR (linguagem/tiffany.py, 9/9): DIZER colapsa \ensuremath{\Psi} sobre os candidatos (o que dizer, argmax — o pensar); FALAR converte a frase escolhida em VISEMAS + a linha do tempo dos músculos da boca (o lip-sync — o rosto). O que a Tiffany diz vira, literalmente, a boca sincronizada sobre o scan real. Falar DISSIPA entropia (cada visema um colapso \ensuremath{\Psi}), e a órbita da fala fecha a volta no silêncio. Verif.: dizer([('ola…',5),('tchau',2)]) colapsa em 'ola…' e o rosto o fala (132 quadros de boca).
\prov{relações: usa orbitas\_de\_fala; parte\_de lib\_tiffany\_no\_metal; usa bai\_psi\_colapso; usa atuadores\_da\_face}
\prov{proveniência: wiki \textperiodcentered{} órbitas de fala \textperiodcentered{} fato \textperiodcentered{} confiança 1.0 \textperiodcentered{} domínio biologia \textperiodcentered{} história 5 versões no jornal}

%CRISTAL {"arestas":[{"alvo":"serie_harmonica_musical","confianca":1.0,"epistemico":"fato","origem":"wiki","rel":"usa"},{"alvo":"analise_de_fourier","confianca":1.0,"epistemico":"fato","origem":"wiki","rel":"explica"}],"confianca":1.0,"contraexemplos":[],"descricao":"A 'cor' do som — o que distingue violino de flauta na mesma nota: o ESPECTRO das parciais. É análise de Fourier ouvida.","epistemico":"fato","exemplos":[],"id":"timbre","memoria":"semantica","meta":{"dominio":"ciencias_de_fora","fonte":""},"origem":"wiki","palavras_chave":[],"sinonimos":[],"tipo":"conceito","titulo":"Timbre"}
\section{Timbre}
A 'cor' do som — o que distingue violino de flauta na mesma nota: o ESPECTRO das parciais. É análise de Fourier ouvida.
\prov{relações: usa serie\_harmonica\_musical; explica analise\_de\_fourier}
\prov{proveniência: wiki \textperiodcentered{} fato \textperiodcentered{} confiança 1.0 \textperiodcentered{} domínio ciencias\_de\_fora \textperiodcentered{} história 3 versões no jornal}

%CRISTAL {"arestas":[{"alvo":"morfologia","confianca":1.0,"epistemico":"fato","origem":"wiki","rel":"usa"},{"alvo":"universais_linguisticos","confianca":1.0,"epistemico":"fato","origem":"wiki","rel":"usa"}],"confianca":1.0,"contraexemplos":[],"descricao":"Como as línguas montam palavras: ANALÍTICA (isolante: mandarim, palavras invariáveis + ordem), SINTÉTICA (fusional: latim, uma desinência acumula caso+número), AGLUTINANTE (turco/japonês: sufixos empilhados, um por função).","epistemico":"inferencia","exemplos":[],"id":"tipologia_morfologica","memoria":"semantica","meta":{"dominio":"ciencias_de_fora","fonte":""},"origem":"wiki","palavras_chave":[],"sinonimos":[],"tipo":"conceito","titulo":"Tipologia Morfológica"}
\section{Tipologia Morfológica}
Como as línguas montam palavras: ANALÍTICA (isolante: mandarim, palavras invariáveis + ordem), SINTÉTICA (fusional: latim, uma desinência acumula caso+número), AGLUTINANTE (turco/japonês: sufixos empilhados, um por função).
\prov{relações: usa morfologia; usa universais\_linguisticos}
\prov{proveniência: wiki \textperiodcentered{} inferencia \textperiodcentered{} confiança 1.0 \textperiodcentered{} domínio ciencias\_de\_fora \textperiodcentered{} história 3 versões no jornal}

%CRISTAL {"arestas":[{"alvo":"mudanca_climatica","confianca":1.0,"epistemico":"fato","origem":"wiki","rel":"usa"},{"alvo":"dissipacao_prigogine","confianca":1.0,"epistemico":"fato","origem":"wiki","rel":"usa"},{"alvo":"limites_planetarios","confianca":1.0,"epistemico":"fato","origem":"wiki","rel":"relaciona"}],"confianca":1.0,"contraexemplos":[],"descricao":"Lenton: limiares críticos em subsistemas planetários (Amazônia, calotas polares, circulação do Atlântico/AMOC, permafrost) além dos quais o sistema se reorganiza de forma abrupta e frequentemente irreversível.","epistemico":"inferencia","exemplos":[],"id":"tipping_points","memoria":"semantica","meta":{"autor":"Lenton","dominio":"ciencias_de_fora","fonte":""},"origem":"wiki","palavras_chave":[],"sinonimos":[],"tipo":"conceito","titulo":"Pontos de Inflexão do Sistema Terra"}
\section{Pontos de Inflexão do Sistema Terra}
Lenton: limiares críticos em subsistemas planetários (Amazônia, calotas polares, circulação do Atlântico/AMOC, permafrost) além dos quais o sistema se reorganiza de forma abrupta e frequentemente irreversível.
\prov{relações: usa mudanca\_climatica; usa dissipacao\_prigogine; relaciona limites\_planetarios}
\prov{proveniência: wiki \textperiodcentered{} inferencia \textperiodcentered{} confiança 1.0 \textperiodcentered{} domínio ciencias\_de\_fora \textperiodcentered{} história 4 versões no jornal}

%CRISTAL {"arestas":[{"alvo":"musica","confianca":1.0,"epistemico":"fato","origem":"wiki","rel":"parte_de"},{"alvo":"elementos_musica","confianca":1.0,"epistemico":"fato","origem":"wiki","rel":"usa"},{"alvo":"harmonia","confianca":1.0,"epistemico":"fato","origem":"wiki","rel":"usa"}],"confianca":1.0,"contraexemplos":[],"descricao":"O sistema que organiza a música em torno de um centro (a tônica) — tensão e repouso, o enredo harmônico.","epistemico":"fato","exemplos":[],"id":"tonalidade","memoria":"semantica","meta":{"dominio":"ciencias_de_fora","fonte":""},"origem":"wiki","palavras_chave":[],"sinonimos":[],"tipo":"conceito","titulo":"Tonalidade"}
\section{Tonalidade}
O sistema que organiza a música em torno de um centro (a tônica) — tensão e repouso, o enredo harmônico.
\prov{relações: parte\_de musica; usa elementos\_musica; usa harmonia}
\prov{proveniência: wiki \textperiodcentered{} fato \textperiodcentered{} confiança 1.0 \textperiodcentered{} domínio ciencias\_de\_fora \textperiodcentered{} história 5 versões no jornal}

%CRISTAL {"arestas":[{"alvo":"musica_notas_sao_modos","confianca":1.0,"epistemico":"fato","origem":"wiki","rel":"parte_de"},{"alvo":"numeros_de_pisot","confianca":1.0,"epistemico":"fato","origem":"wiki","rel":"eh_um"},{"alvo":"bemtevi_modo_puro","confianca":1.0,"epistemico":"fato","origem":"wiki","rel":"relaciona"},{"alvo":"corpo_estelar","confianca":1.0,"epistemico":"fato","origem":"wiki","rel":"usa"}],"confianca":1.0,"contraexemplos":[],"descricao":"A classe de nota dominante — a TÔNICA, para onde a música resolve — é o CRISTAL: o atrator, o ponto fixo Pisot da harmonia. A cadência que volta a ela é o colapso Ψ. Na MARCHA TURCA de Mozart (K.331), medido do arquivo real (model/mozart): LÁ (A) com 19.7% da energia, de longe a mais forte — a peça orbita Lá e resolve nele. A tonalidade (Krumhansl) é Lá maior (corr 0.864; o tema é Lá menor, tônica Lá em ambos), o andamento ~128 BPM (Allegretto — a órbita rítmica da borda). A tônica é a reta mineral da harmonia; as tensões/dissonâncias são a borda e o caos que resolvem no cristal.","epistemico":"fato","exemplos":[],"id":"tonica_e_o_cristal","memoria":"semantica","meta":{"dominio":"arte","fonte":"música mineral"},"origem":"wiki","palavras_chave":[],"sinonimos":[],"tipo":"conceito","titulo":"A Tônica é o Cristal (o atrator da harmonia)"}
\section{A Tônica é o Cristal (o atrator da harmonia)}
A classe de nota dominante — a TÔNICA, para onde a música resolve — é o CRISTAL: o atrator, o ponto fixo Pisot da harmonia. A cadência que volta a ela é o colapso \ensuremath{\Psi}. Na MARCHA TURCA de Mozart (K.331), medido do arquivo real (model/mozart): LÁ (A) com 19.7\% da energia, de longe a mais forte — a peça orbita Lá e resolve nele. A tonalidade (Krumhansl) é Lá maior (corr 0.864; o tema é Lá menor, tônica Lá em ambos), o andamento \~{}128 BPM (Allegretto — a órbita rítmica da borda). A tônica é a reta mineral da harmonia; as tensões/dissonâncias são a borda e o caos que resolvem no cristal.
\prov{relações: parte\_de musica\_notas\_sao\_modos; eh\_um numeros\_de\_pisot; relaciona bemtevi\_modo\_puro; usa corpo\_estelar}
\prov{proveniência: wiki \textperiodcentered{} música mineral \textperiodcentered{} fato \textperiodcentered{} confiança 1.0 \textperiodcentered{} domínio arte \textperiodcentered{} história 5 versões no jornal}

%CRISTAL {"arestas":[{"alvo":"ordem_das_palavras","confianca":1.0,"epistemico":"fato","origem":"wiki","rel":"relaciona"}],"confianca":1.0,"contraexemplos":[],"descricao":"Línguas que organizam a frase por TÓPICO-comentário em vez de sujeito-predicado (mandarim, japonês) — 'Quanto ao peixe, já comi'. Outra lógica de frase.","epistemico":"inferencia","exemplos":[],"id":"topico_proeminencia","memoria":"semantica","meta":{"dominio":"ciencias_de_fora","fonte":""},"origem":"wiki","palavras_chave":[],"sinonimos":[],"tipo":"conceito","titulo":"Proeminência de Tópico"}
\section{Proeminência de Tópico}
Línguas que organizam a frase por TÓPICO-comentário em vez de sujeito-predicado (mandarim, japonês) — 'Quanto ao peixe, já comi'. Outra lógica de frase.
\prov{relações: relaciona ordem\_das\_palavras}
\prov{proveniência: wiki \textperiodcentered{} inferencia \textperiodcentered{} confiança 1.0 \textperiodcentered{} domínio ciencias\_de\_fora \textperiodcentered{} história 2 versões no jornal}

%CRISTAL {"arestas":[{"alvo":"geografia_humana","confianca":1.0,"epistemico":"fato","origem":"wiki","rel":"parte_de"},{"alvo":"categorias_geograficas","confianca":1.0,"epistemico":"fato","origem":"wiki","rel":"usa"},{"alvo":"capital_cultural","confianca":1.0,"epistemico":"fato","origem":"wiki","rel":"relaciona"}],"confianca":1.0,"contraexemplos":[],"descricao":"Yi-Fu Tuan: topofilia é o laço afetivo entre as pessoas e o lugar. O espaço é abertura, abstração e liberdade; o lugar é a pausa que se enche de valor, segurança e significado — o 'sentido de lugar' nasce quando a experiência transforma o espaço indiferenciado em lugar afetivo.","epistemico":"inferencia","exemplos":[],"id":"topofilia_lugar_tuan","memoria":"semantica","meta":{"autor":"Tuan","dominio":"ciencias_de_fora","fonte":""},"origem":"wiki","palavras_chave":[],"sinonimos":[],"tipo":"conceito","titulo":"Topofilia e o Lugar (Tuan)"}
\section{Topofilia e o Lugar (Tuan)}
Yi-Fu Tuan: topofilia é o laço afetivo entre as pessoas e o lugar. O espaço é abertura, abstração e liberdade; o lugar é a pausa que se enche de valor, segurança e significado — o 'sentido de lugar' nasce quando a experiência transforma o espaço indiferenciado em lugar afetivo.
\prov{relações: parte\_de geografia\_humana; usa categorias\_geograficas; relaciona capital\_cultural}
\prov{proveniência: wiki \textperiodcentered{} inferencia \textperiodcentered{} confiança 1.0 \textperiodcentered{} domínio ciencias\_de\_fora \textperiodcentered{} história 4 versões no jornal}

%CRISTAL {"arestas":[{"alvo":"valor","confianca":1.0,"epistemico":"fato","origem":"curriculo","rel":"continua"},{"alvo":"sha256","confianca":1.0,"epistemico":"fato","origem":"curriculo","rel":"usa"}],"confianca":1.0,"contraexemplos":[],"descricao":"Esforço medido que transforma estado; mineração = ler o tecido (estrela de prata).","epistemico":"fato","exemplos":[],"id":"trabalho","memoria":"semantica","meta":{"dominio":"mercado","nivel":5,"ordem":1},"origem":"curriculo","palavras_chave":["mineração","hash","proof of work"],"sinonimos":["trabalho computacional","PoW"],"tipo":"conceito","titulo":"trabalho"}
\section{trabalho}
Esforço medido que transforma estado; mineração = ler o tecido (estrela de prata).
\prov{sinónimos: trabalho computacional; PoW}
\prov{relações: continua valor; usa sha256}
\prov{proveniência: curriculo \textperiodcentered{} fato \textperiodcentered{} confiança 1.0 \textperiodcentered{} domínio mercado \textperiodcentered{} história 35 versões no jornal}

%CRISTAL {"arestas":[{"alvo":"fonologia","confianca":1.0,"epistemico":"fato","origem":"wiki","rel":"parte_de"},{"alvo":"roman_jakobson","confianca":1.0,"epistemico":"fato","origem":"wiki","rel":"parte_de"}],"confianca":1.0,"contraexemplos":[],"descricao":"Jakobson: o fonema decompõe-se em traços binários (sonoro/surdo, nasal/oral) — os bits da fonologia.","epistemico":"fato","exemplos":[],"id":"traco_distintivo","memoria":"semantica","meta":{"autor":"Jakobson","dominio":"ciencias_de_fora","fonte":""},"origem":"wiki","palavras_chave":[],"sinonimos":[],"tipo":"conceito","titulo":"Traço Distintivo"}
\section{Traço Distintivo}
Jakobson: o fonema decompõe-se em traços binários (sonoro/surdo, nasal/oral) — os bits da fonologia.
\prov{relações: parte\_de fonologia; parte\_de roman\_jakobson}
\prov{proveniência: wiki \textperiodcentered{} fato \textperiodcentered{} confiança 1.0 \textperiodcentered{} domínio ciencias\_de\_fora \textperiodcentered{} história 3 versões no jornal}

%CRISTAL {"arestas":[{"alvo":"jogos_de_linguagem","confianca":1.0,"epistemico":"fato","origem":"wiki","rel":"contradiz"},{"alvo":"semiotica_peirce_triade","confianca":1.0,"epistemico":"fato","origem":"wiki","rel":"relaciona"}],"confianca":1.0,"contraexemplos":[],"descricao":"Wittgenstein (inicial): a proposição é figura de um estado de coisas; sua estrutura lógica espelha o fato. 'Do que não se pode falar, deve-se calar.'","epistemico":"inferencia","exemplos":[],"id":"tractatus_figuracao","memoria":"semantica","meta":{"autor":"Wittgenstein","dominio":"ciencias_de_fora","fonte":""},"origem":"wiki","palavras_chave":[],"sinonimos":[],"tipo":"conceito","titulo":"Teoria da Figuração (Tractatus)"}
\section{Teoria da Figuração (Tractatus)}
Wittgenstein (inicial): a proposição é figura de um estado de coisas; sua estrutura lógica espelha o fato. 'Do que não se pode falar, deve-se calar.'
\prov{relações: contradiz jogos\_de\_linguagem; relaciona semiotica\_peirce\_triade}
\prov{proveniência: wiki \textperiodcentered{} inferencia \textperiodcentered{} confiança 1.0 \textperiodcentered{} domínio ciencias\_de\_fora \textperiodcentered{} história 3 versões no jornal}

%CRISTAL {"arestas":[{"alvo":"bilinguismo","confianca":1.0,"epistemico":"fato","origem":"wiki","rel":"usa"},{"alvo":"ordem_das_palavras","confianca":1.0,"epistemico":"fato","origem":"wiki","rel":"usa"},{"alvo":"arvore_sintatica","confianca":1.0,"epistemico":"fato","origem":"wiki","rel":"usa"},{"alvo":"idioma","confianca":1.0,"epistemico":"fato","origem":"wiki","rel":"usa"},{"alvo":"traducao_e_intraduzivel","confianca":1.0,"epistemico":"fato","origem":"wiki","rel":"especializa"},{"alvo":"sintaxe","confianca":1.0,"epistemico":"fato","origem":"wiki","rel":"usa"},{"alvo":"semantica","confianca":1.0,"epistemico":"fato","origem":"wiki","rel":"usa"},{"alvo":"indeterminacao_traducao","confianca":1.0,"epistemico":"fato","origem":"wiki","rel":"relaciona"}],"confianca":1.0,"contraexemplos":[],"descricao":"Reexpressar o sentido de uma língua noutra — nunca cópia exata: transforma a forma (sintaxe, léxico) preservando o significado. Toda tradução é interpretação e escolha.","epistemico":"inferencia","exemplos":[],"id":"traducao","memoria":"semantica","meta":{"dominio":"ciencias_de_fora","fonte":""},"origem":"wiki","palavras_chave":[],"sinonimos":[],"tipo":"conceito","titulo":"Tradução"}
\section{Tradução}
Reexpressar o sentido de uma língua noutra — nunca cópia exata: transforma a forma (sintaxe, léxico) preservando o significado. Toda tradução é interpretação e escolha.
\prov{relações: usa bilinguismo; usa ordem\_das\_palavras; usa arvore\_sintatica; usa idioma; especializa traducao\_e\_intraduzivel; usa sintaxe; usa semantica; relaciona indeterminacao\_traducao}
\prov{proveniência: wiki \textperiodcentered{} inferencia \textperiodcentered{} confiança 1.0 \textperiodcentered{} domínio ciencias\_de\_fora \textperiodcentered{} história 14 versões no jornal}

%CRISTAL {"arestas":[{"alvo":"traducao","confianca":1.0,"epistemico":"fato","origem":"wiki","rel":"especializa"},{"alvo":"linguistica_computacional","confianca":1.0,"epistemico":"fato","origem":"wiki","rel":"usa"}],"confianca":1.0,"contraexemplos":[],"descricao":"Traduzir por máquina — o sonho desde os anos 1950. Três eras: baseada em regras, estatística e neural. Cada uma redefiniu o que a máquina 'entende'.","epistemico":"fato","exemplos":[],"id":"traducao_automatica","memoria":"semantica","meta":{"dominio":"ciencias_de_fora","fonte":""},"origem":"wiki","palavras_chave":[],"sinonimos":[],"tipo":"conceito","titulo":"Tradução Automática"}
\section{Tradução Automática}
Traduzir por máquina — o sonho desde os anos 1950. Três eras: baseada em regras, estatística e neural. Cada uma redefiniu o que a máquina 'entende'.
\prov{relações: especializa traducao; usa linguistica\_computacional}
\prov{proveniência: wiki \textperiodcentered{} fato \textperiodcentered{} confiança 1.0 \textperiodcentered{} domínio ciencias\_de\_fora \textperiodcentered{} história 3 versões no jornal}

%CRISTAL {"arestas":[{"alvo":"traducao","confianca":1.0,"epistemico":"fato","origem":"wiki","rel":"explica"},{"alvo":"interpretante","confianca":1.0,"epistemico":"fato","origem":"wiki","rel":"usa"},{"alvo":"semiose_ilimitada","confianca":1.0,"epistemico":"fato","origem":"wiki","rel":"usa"},{"alvo":"algebra_de_peirce","confianca":1.0,"epistemico":"fato","origem":"wiki","rel":"relaciona"},{"alvo":"decomposicao_peirce","confianca":1.0,"epistemico":"fato","origem":"wiki","rel":"relaciona"},{"alvo":"simbolo_vs_significado","confianca":1.0,"epistemico":"fato","origem":"wiki","rel":"usa"},{"alvo":"indeterminacao_traducao","confianca":1.0,"epistemico":"fato","origem":"wiki","rel":"relaciona"},{"alvo":"grafo_tiffany_e_algebra_peirce","confianca":1.0,"epistemico":"fato","origem":"wiki","rel":"usa"}],"confianca":1.0,"contraexemplos":[],"descricao":"Peirce: o INTERPRETANTE de um signo é ele traduzido em OUTRO signo — logo toda semiose É tradução, e o sentido de um signo É sua tradução (Jakobson chamou de 'tradução intralingual/interlingual/intersemiótica'). A álgebra de Peirce do broca-so, então, JÁ É uma teoria da tradução: substituir um signo por outro conservando a relação.","epistemico":"inferencia","exemplos":[],"id":"traducao_como_interpretante","memoria":"semantica","meta":{"autor":"Peirce/Jakobson","dominio":"ciencias_de_fora","fonte":""},"origem":"wiki","palavras_chave":[],"sinonimos":[],"tipo":"conceito","titulo":"Tradução como Interpretante (Peirce)"}
\section{Tradução como Interpretante (Peirce)}
Peirce: o INTERPRETANTE de um signo é ele traduzido em OUTRO signo — logo toda semiose É tradução, e o sentido de um signo É sua tradução (Jakobson chamou de 'tradução intralingual/interlingual/intersemiótica'). A álgebra de Peirce do broca-so, então, JÁ É uma teoria da tradução: substituir um signo por outro conservando a relação.
\prov{relações: explica traducao; usa interpretante; usa semiose\_ilimitada; relaciona algebra\_de\_peirce; relaciona decomposicao\_peirce; usa simbolo\_vs\_significado; relaciona indeterminacao\_traducao; usa grafo\_tiffany\_e\_algebra\_peirce}
\prov{proveniência: wiki \textperiodcentered{} inferencia \textperiodcentered{} confiança 1.0 \textperiodcentered{} domínio ciencias\_de\_fora \textperiodcentered{} história 9 versões no jornal}

%CRISTAL {"arestas":[{"alvo":"traducao","confianca":1.0,"epistemico":"fato","origem":"wiki","rel":"explica"},{"alvo":"transformada_metalica","confianca":1.0,"epistemico":"fato","origem":"wiki","rel":"usa"},{"alvo":"meta_inducao_geral_simbolo","confianca":1.0,"epistemico":"fato","origem":"wiki","rel":"usa"},{"alvo":"torre_unificada","confianca":1.0,"epistemico":"fato","origem":"wiki","rel":"parte_de"},{"alvo":"sintaxe","confianca":1.0,"epistemico":"fato","origem":"wiki","rel":"usa"},{"alvo":"bai_orbita_reta_metalica","confianca":1.0,"epistemico":"fato","origem":"wiki","rel":"relaciona"}],"confianca":1.0,"contraexemplos":[],"descricao":"Traduzir é uma TRANSFORMADA: muda a base (a língua, a sintaxe) preservando a estrutura (o sentido). É a meta-indução do broca-so na fala — só muda o SÍMBOLO, o padrão se conserva. A mesma lógica da transformada metálica, que troca a base σ_n mantendo a estrutura. Português→inglês é uma mudança de base do significado.","epistemico":"inferencia","exemplos":[],"id":"traducao_como_transformacao","memoria":"semantica","meta":{"dominio":"ciencias_de_fora","fonte":""},"origem":"wiki","palavras_chave":[],"sinonimos":[],"tipo":"conceito","titulo":"Tradução como Transformada"}
\section{Tradução como Transformada}
Traduzir é uma TRANSFORMADA: muda a base (a língua, a sintaxe) preservando a estrutura (o sentido). É a meta-indução do broca-so na fala — só muda o SÍMBOLO, o padrão se conserva. A mesma lógica da transformada metálica, que troca a base \ensuremath{\sigma}\_n mantendo a estrutura. Português\ensuremath{\to}inglês é uma mudança de base do significado.
\prov{relações: explica traducao; usa transformada\_metalica; usa meta\_inducao\_geral\_simbolo; parte\_de torre\_unificada; usa sintaxe; relaciona bai\_orbita\_reta\_metalica}
\prov{proveniência: wiki \textperiodcentered{} inferencia \textperiodcentered{} confiança 1.0 \textperiodcentered{} domínio ciencias\_de\_fora \textperiodcentered{} história 7 versões no jornal}

%CRISTAL {"arestas":[{"alvo":"roman_jakobson","confianca":1.0,"epistemico":"fato","origem":"wiki","rel":"parte_de"},{"alvo":"significante_significado","confianca":1.0,"epistemico":"fato","origem":"wiki","rel":"usa"}],"confianca":1.0,"contraexemplos":[],"descricao":"Jakobson: traduzir é reformular signos; entre línguas há sempre resto — a forma poética resiste à tradução.","epistemico":"inferencia","exemplos":[],"id":"traducao_e_intraduzivel","memoria":"semantica","meta":{"autor":"Jakobson","dominio":"ciencias_de_fora","fonte":""},"origem":"wiki","palavras_chave":[],"sinonimos":[],"tipo":"conceito","titulo":"Tradução e o Intraduzível"}
\section{Tradução e o Intraduzível}
Jakobson: traduzir é reformular signos; entre línguas há sempre resto — a forma poética resiste à tradução.
\prov{relações: parte\_de roman\_jakobson; usa significante\_significado}
\prov{proveniência: wiki \textperiodcentered{} inferencia \textperiodcentered{} confiança 1.0 \textperiodcentered{} domínio ciencias\_de\_fora \textperiodcentered{} história 3 versões no jornal}

%CRISTAL {"arestas":[{"alvo":"traducao_automatica","confianca":1.0,"epistemico":"fato","origem":"wiki","rel":"especializa"},{"alvo":"teorema_de_bayes","confianca":1.0,"epistemico":"fato","origem":"wiki","rel":"usa"},{"alvo":"alinhamento_de_traducao","confianca":1.0,"epistemico":"fato","origem":"wiki","rel":"usa"}],"confianca":1.0,"contraexemplos":[],"descricao":"A 2ª era: aprender de corpora bilíngues quais traduções são mais PROVÁVEIS (Bayes, tabelas de frases, alinhamento) — sem entender, só contar. Dados no lugar de regras.","epistemico":"fato","exemplos":[],"id":"traducao_estatistica","memoria":"semantica","meta":{"dominio":"ciencias_de_fora","fonte":""},"origem":"wiki","palavras_chave":[],"sinonimos":[],"tipo":"conceito","titulo":"Tradução Estatística"}
\section{Tradução Estatística}
A 2ª era: aprender de corpora bilíngues quais traduções são mais PROVÁVEIS (Bayes, tabelas de frases, alinhamento) — sem entender, só contar. Dados no lugar de regras.
\prov{relações: especializa traducao\_automatica; usa teorema\_de\_bayes; usa alinhamento\_de\_traducao}
\prov{proveniência: wiki \textperiodcentered{} fato \textperiodcentered{} confiança 1.0 \textperiodcentered{} domínio ciencias\_de\_fora \textperiodcentered{} história 4 versões no jornal}

%CRISTAL {"arestas":[{"alvo":"traducao","confianca":1.0,"epistemico":"fato","origem":"wiki","rel":"parte_de"}],"confianca":1.0,"contraexemplos":[],"descricao":"A tensão fundadora (Cícero, São Jerônimo): traduzir palavra-por-palavra (fiel à forma) ou sentido-por-sentido (fiel ao significado). Nunca os dois ao mesmo tempo.","epistemico":"inferencia","exemplos":[],"id":"traducao_literal_vs_livre","memoria":"semantica","meta":{"autor":"Jerônimo","dominio":"ciencias_de_fora","fonte":""},"origem":"wiki","palavras_chave":[],"sinonimos":[],"tipo":"conceito","titulo":"Literal vs Livre"}
\section{Literal vs Livre}
A tensão fundadora (Cícero, São Jerônimo): traduzir palavra-por-palavra (fiel à forma) ou sentido-por-sentido (fiel ao significado). Nunca os dois ao mesmo tempo.
\prov{relações: parte\_de traducao}
\prov{proveniência: wiki \textperiodcentered{} inferencia \textperiodcentered{} confiança 1.0 \textperiodcentered{} domínio ciencias\_de\_fora \textperiodcentered{} história 2 versões no jornal}

%CRISTAL {"arestas":[{"alvo":"traducao_automatica","confianca":1.0,"epistemico":"fato","origem":"wiki","rel":"especializa"},{"alvo":"transformer","confianca":1.0,"epistemico":"fato","origem":"wiki","rel":"usa"},{"alvo":"mecanismo_de_atencao","confianca":1.0,"epistemico":"fato","origem":"wiki","rel":"usa"},{"alvo":"aprendizado_supervisionado","confianca":1.0,"epistemico":"fato","origem":"wiki","rel":"usa"}],"confianca":1.0,"contraexemplos":[],"descricao":"A 3ª era: uma rede (seq2seq + atenção, hoje transformers) mapeia a frase-fonte a um vetor de sentido e gera a frase-alvo — a que trouxe a tradução quase humana.","epistemico":"fato","exemplos":[],"id":"traducao_neural","memoria":"semantica","meta":{"dominio":"ciencias_de_fora","fonte":""},"origem":"wiki","palavras_chave":[],"sinonimos":[],"tipo":"conceito","titulo":"Tradução Neural"}
\section{Tradução Neural}
A 3ª era: uma rede (seq2seq + atenção, hoje transformers) mapeia a frase-fonte a um vetor de sentido e gera a frase-alvo — a que trouxe a tradução quase humana.
\prov{relações: especializa traducao\_automatica; usa transformer; usa mecanismo\_de\_atencao; usa aprendizado\_supervisionado}
\prov{proveniência: wiki \textperiodcentered{} fato \textperiodcentered{} confiança 1.0 \textperiodcentered{} domínio ciencias\_de\_fora \textperiodcentered{} história 5 versões no jornal}

%CRISTAL {"arestas":[{"alvo":"traducao_automatica","confianca":1.0,"epistemico":"fato","origem":"wiki","rel":"especializa"},{"alvo":"arvore_sintatica","confianca":1.0,"epistemico":"fato","origem":"wiki","rel":"usa"},{"alvo":"ordem_das_palavras","confianca":1.0,"epistemico":"fato","origem":"wiki","rel":"usa"}],"confianca":1.0,"contraexemplos":[],"descricao":"A 1ª era: dicionários + regras gramaticais que TRANSFEREM a árvore sintática da fonte para o alvo. Frágil e trabalhosa, mas explícita — a sintaxe no comando.","epistemico":"fato","exemplos":[],"id":"traducao_por_regras","memoria":"semantica","meta":{"dominio":"ciencias_de_fora","fonte":""},"origem":"wiki","palavras_chave":[],"sinonimos":[],"tipo":"conceito","titulo":"Tradução por Regras"}
\section{Tradução por Regras}
A 1ª era: dicionários + regras gramaticais que TRANSFEREM a árvore sintática da fonte para o alvo. Frágil e trabalhosa, mas explícita — a sintaxe no comando.
\prov{relações: especializa traducao\_automatica; usa arvore\_sintatica; usa ordem\_das\_palavras}
\prov{proveniência: wiki \textperiodcentered{} fato \textperiodcentered{} confiança 1.0 \textperiodcentered{} domínio ciencias\_de\_fora \textperiodcentered{} história 4 versões no jornal}

%CRISTAL {"arestas":[{"alvo":"generos_literarios","confianca":1.0,"epistemico":"fato","origem":"wiki","rel":"especializa"},{"alvo":"catarse","confianca":1.0,"epistemico":"fato","origem":"wiki","rel":"usa"},{"alvo":"poetica_aristoteles","confianca":1.0,"epistemico":"fato","origem":"wiki","rel":"relaciona"}],"confianca":1.0,"contraexemplos":[],"descricao":"O drama da queda inevitável — o herói derrubado por sua falha (hamartia); provoca terror e piedade (catarse).","epistemico":"fato","exemplos":[],"id":"tragedia","memoria":"semantica","meta":{"dominio":"ciencias_de_fora","fonte":""},"origem":"wiki","palavras_chave":[],"sinonimos":[],"tipo":"conceito","titulo":"Tragédia"}
\section{Tragédia}
O drama da queda inevitável — o herói derrubado por sua falha (hamartia); provoca terror e piedade (catarse).
\prov{relações: especializa generos\_literarios; usa catarse; relaciona poetica\_aristoteles}
\prov{proveniência: wiki \textperiodcentered{} fato \textperiodcentered{} confiança 1.0 \textperiodcentered{} domínio ciencias\_de\_fora \textperiodcentered{} história 4 versões no jornal}

%CRISTAL {"arestas":[{"alvo":"economia","confianca":1.0,"epistemico":"fato","origem":"wiki","rel":"parte_de"},{"alvo":"word","confianca":0.5,"epistemico":"fato","origem":"wiki","rel":"relaciona"},{"alvo":"virtualizacao","confianca":0.5,"epistemico":"fato","origem":"wiki","rel":"relaciona"},{"alvo":"vontade_de_poder","confianca":0.5,"epistemico":"fato","origem":"wiki","rel":"relaciona"}],"confianca":1.0,"contraexemplos":[],"descricao":"Hardin: recursos compartilhados de acesso livre são sobre-explorados; Ostrom refuta empiricamente — comunidades os autogovernam por regras claras.","epistemico":"inferencia","exemplos":[],"id":"tragedia_dos_comuns","memoria":"semantica","meta":{"autor":"Hardin/Ostrom","dominio":"ciencias_de_fora","fonte":""},"origem":"wiki","palavras_chave":[],"sinonimos":[],"tipo":"conceito","titulo":"Tragédia dos Comuns"}
\section{Tragédia dos Comuns}
Hardin: recursos compartilhados de acesso livre são sobre-explorados; Ostrom refuta empiricamente — comunidades os autogovernam por regras claras.
\prov{relações: parte\_de economia; relaciona word; relaciona virtualizacao; relaciona vontade\_de\_poder}
\prov{proveniência: wiki \textperiodcentered{} inferencia \textperiodcentered{} confiança 1.0 \textperiodcentered{} domínio ciencias\_de\_fora \textperiodcentered{} história 5 versões no jornal}

%CRISTAL {"arestas":[{"alvo":"epidemiologia","confianca":1.0,"epistemico":"fato","origem":"wiki","rel":"usa"},{"alvo":"determinantes_sociais_saude","confianca":1.0,"epistemico":"fato","origem":"wiki","rel":"relaciona"}],"confianca":1.0,"contraexemplos":[],"descricao":"Omran: o deslocamento histórico do perfil de mortalidade, da predominância de doenças infecciosas/fome para a de doenças crônicas não transmissíveis (cardiovasculares, câncer, diabetes), acompanhando a queda da mortalidade e o envelhecimento populacional.","epistemico":"inferencia","exemplos":[],"id":"transicao_epidemiologica","memoria":"semantica","meta":{"autor":"Omran","dominio":"ciencias_de_fora","fonte":""},"origem":"wiki","palavras_chave":[],"sinonimos":[],"tipo":"conceito","titulo":"Transição Epidemiológica"}
\section{Transição Epidemiológica}
Omran: o deslocamento histórico do perfil de mortalidade, da predominância de doenças infecciosas/fome para a de doenças crônicas não transmissíveis (cardiovasculares, câncer, diabetes), acompanhando a queda da mortalidade e o envelhecimento populacional.
\prov{relações: usa epidemiologia; relaciona determinantes\_sociais\_saude}
\prov{proveniência: wiki \textperiodcentered{} inferencia \textperiodcentered{} confiança 1.0 \textperiodcentered{} domínio ciencias\_de\_fora \textperiodcentered{} história 3 versões no jornal}

%CRISTAL {"arestas":[{"alvo":"didatica_matematica","confianca":1.0,"epistemico":"fato","origem":"wiki","rel":"parte_de"},{"alvo":"matematica","confianca":1.0,"epistemico":"fato","origem":"wiki","rel":"referencia"}],"confianca":1.0,"contraexemplos":[],"descricao":"Chevallard: o processo pelo qual o saber sábio ('savoir savant') é adaptado, reorganizado e por vezes distorcido para virar saber a ensinar e depois saber efetivamente ensinado, regulado pela 'noosfera'.","epistemico":"inferencia","exemplos":[],"id":"transposicao_didatica","memoria":"semantica","meta":{"autor":"Chevallard","dominio":"ciencias_de_fora","fonte":""},"origem":"wiki","palavras_chave":[],"sinonimos":[],"tipo":"conceito","titulo":"Transposição Didática"}
\section{Transposição Didática}
Chevallard: o processo pelo qual o saber sábio ('savoir savant') é adaptado, reorganizado e por vezes distorcido para virar saber a ensinar e depois saber efetivamente ensinado, regulado pela 'noosfera'.
\prov{relações: parte\_de didatica\_matematica; referencia matematica}
\prov{proveniência: wiki \textperiodcentered{} inferencia \textperiodcentered{} confiança 1.0 \textperiodcentered{} domínio ciencias\_de\_fora \textperiodcentered{} história 3 versões no jornal}

%CRISTAL {"arestas":[{"alvo":"filosofia_da_tecnica","confianca":1.0,"epistemico":"fato","origem":"wiki","rel":"parte_de"},{"alvo":"superinteligencia_alinhamento","confianca":1.0,"epistemico":"fato","origem":"wiki","rel":"usa"},{"alvo":"corpo_dual","confianca":1.0,"epistemico":"fato","origem":"wiki","rel":"relaciona"}],"confianca":1.0,"contraexemplos":[],"descricao":"O movimento que afirma a possibilidade e a desejabilidade de melhorar radicalmente a condição humana pela tecnologia (cognição, longevidade, capacidades), rumo a um estágio 'pós-humano'. Posição de forte carga especulativa.","epistemico":"hipotese","exemplos":[],"id":"transumanismo","memoria":"semantica","meta":{"autor":"Bostrom/More","dominio":"ciencias_de_fora","fonte":""},"origem":"wiki","palavras_chave":[],"sinonimos":[],"tipo":"conceito","titulo":"Transumanismo"}
\section{Transumanismo}
O movimento que afirma a possibilidade e a desejabilidade de melhorar radicalmente a condição humana pela tecnologia (cognição, longevidade, capacidades), rumo a um estágio 'pós-humano'. Posição de forte carga especulativa.
\prov{relações: parte\_de filosofia\_da\_tecnica; usa superinteligencia\_alinhamento; relaciona corpo\_dual}
\prov{proveniência: wiki \textperiodcentered{} hipotese \textperiodcentered{} confiança 1.0 \textperiodcentered{} domínio ciencias\_de\_fora \textperiodcentered{} história 4 versões no jornal}

%CRISTAL {"arestas":[{"alvo":"saude_na_borda","confianca":1.0,"epistemico":"fato","origem":"wiki","rel":"usa"},{"alvo":"homeostase_controle","confianca":1.0,"epistemico":"fato","origem":"wiki","rel":"usa"},{"alvo":"estabilidade_da_rede","confianca":1.0,"epistemico":"fato","origem":"wiki","rel":"relaciona"},{"alvo":"medicina","confianca":1.0,"epistemico":"fato","origem":"wiki","rel":"parte_de"}],"confianca":1.0,"contraexemplos":[],"descricao":"Tratar é empurrar a órbita do corpo de volta ao atrator da saúde — como estabilizar a rede ou equilibrar o pêndulo. Antibiótico derruba o R₀ interno; insulina reforça o controle da glicose; marca-passo devolve a variabilidade da borda. A terapia é o ganho do controlador.","epistemico":"inferencia","exemplos":[],"id":"tratamento_restaura_atrator","memoria":"semantica","meta":{"dominio":"medicina","fonte":"medicina e diagnóstico"},"origem":"wiki","palavras_chave":[],"sinonimos":[],"tipo":"conceito","titulo":"Tratamento = Restaurar o Atrator"}
\section{Tratamento = Restaurar o Atrator}
Tratar é empurrar a órbita do corpo de volta ao atrator da saúde — como estabilizar a rede ou equilibrar o pêndulo. Antibiótico derruba o R\ensuremath{_{0}} interno; insulina reforça o controle da glicose; marca-passo devolve a variabilidade da borda. A terapia é o ganho do controlador.
\prov{relações: usa saude\_na\_borda; usa homeostase\_controle; relaciona estabilidade\_da\_rede; parte\_de medicina}
\prov{proveniência: wiki \textperiodcentered{} medicina e diagnóstico \textperiodcentered{} inferencia \textperiodcentered{} confiança 1.0 \textperiodcentered{} domínio medicina \textperiodcentered{} história 5 versões no jornal}

%CRISTAL {"arestas":[{"alvo":"ciencia_politica","confianca":1.0,"epistemico":"fato","origem":"wiki","rel":"parte_de"},{"alvo":"hegemonia_cultural","confianca":1.0,"epistemico":"fato","origem":"wiki","rel":"relaciona"},{"alvo":"biopoder","confianca":1.0,"epistemico":"fato","origem":"wiki","rel":"relaciona"},{"alvo":"dois_sistemas_kahneman","confianca":1.0,"epistemico":"fato","origem":"wiki","rel":"relaciona"}],"confianca":1.0,"contraexemplos":[],"descricao":"O poder tem três dimensões cumulativas: prevalecer em decisões observáveis (Dahl), controlar a agenda impedindo que temas cheguem à decisão — a 'não-decisão' (Bachrach & Baratz), e moldar as próprias preferências das pessoas para que aceitem uma ordem contrária a seus interesses (Lukes). Poder≠autoridade (legítima)≠influência.","epistemico":"inferencia","exemplos":[],"id":"tres_faces_do_poder","memoria":"semantica","meta":{"autor":"Dahl/Bachrach/Lukes","dominio":"ciencias_de_fora","fonte":""},"origem":"wiki","palavras_chave":[],"sinonimos":[],"tipo":"conceito","titulo":"As Três Faces do Poder"}
\section{As Três Faces do Poder}
O poder tem três dimensões cumulativas: prevalecer em decisões observáveis (Dahl), controlar a agenda impedindo que temas cheguem à decisão — a 'não-decisão' (Bachrach \& Baratz), e moldar as próprias preferências das pessoas para que aceitem uma ordem contrária a seus interesses (Lukes). Poder\ensuremath{\ne}autoridade (legítima)\ensuremath{\ne}influência.
\prov{relações: parte\_de ciencia\_politica; relaciona hegemonia\_cultural; relaciona biopoder; relaciona dois\_sistemas\_kahneman}
\prov{proveniência: wiki \textperiodcentered{} inferencia \textperiodcentered{} confiança 1.0 \textperiodcentered{} domínio ciencias\_de\_fora \textperiodcentered{} história 5 versões no jornal}

%CRISTAL {"arestas":[{"alvo":"hinduismo","confianca":1.0,"epistemico":"fato","origem":"wiki","rel":"parte_de"},{"alvo":"maya","confianca":1.0,"epistemico":"fato","origem":"wiki","rel":"relaciona"},{"alvo":"wu_wei","confianca":1.0,"epistemico":"fato","origem":"wiki","rel":"contradiz"}],"confianca":1.0,"contraexemplos":[],"descricao":"Gita: as três qualidades da natureza (prakriti) que atam o ser — sattva (pureza, luz, conhecimento), rajas (paixão, atividade, desejo) e tamas (inércia, ignorância). Só a prakriti age; o purusha (o Ser testemunha) na realidade 'não faz nada'. A liberação vem de transcender os gunas.","epistemico":"inferencia","exemplos":[],"id":"tres_gunas","memoria":"semantica","meta":{"autor":"Bhagavad Gita","dominio":"sagrado","fonte":""},"origem":"wiki","palavras_chave":[],"sinonimos":[],"tipo":"conceito","titulo":"Os Três Gunas (sattva, rajas, tamas)"}
\section{Os Três Gunas (sattva, rajas, tamas)}
Gita: as três qualidades da natureza (prakriti) que atam o ser — sattva (pureza, luz, conhecimento), rajas (paixão, atividade, desejo) e tamas (inércia, ignorância). Só a prakriti age; o purusha (o Ser testemunha) na realidade 'não faz nada'. A liberação vem de transcender os gunas.
\prov{relações: parte\_de hinduismo; relaciona maya; contradiz wu\_wei}
\prov{proveniência: wiki \textperiodcentered{} inferencia \textperiodcentered{} confiança 1.0 \textperiodcentered{} domínio sagrado \textperiodcentered{} história 4 versões no jornal}

%CRISTAL {"arestas":[{"alvo":"justica_em_saude_daniels","confianca":1.0,"epistemico":"fato","origem":"wiki","rel":"usa"},{"alvo":"utilitarismo","confianca":1.0,"epistemico":"fato","origem":"wiki","rel":"usa"},{"alvo":"justica_equidade_rawls","confianca":1.0,"epistemico":"fato","origem":"wiki","rel":"relaciona"}],"confianca":1.0,"contraexemplos":[],"descricao":"O problema de distribuir bens de saúde limitados (leitos de UTI, órgãos, vacinas) por critérios como maximização de benefício, prioridade aos piores, sorteio ou fila — expondo o conflito entre utilitarismo e igualdade.","epistemico":"inferencia","exemplos":[],"id":"triagem_alocacao","memoria":"semantica","meta":{"autor":"—","dominio":"ciencias_de_fora","fonte":""},"origem":"wiki","palavras_chave":[],"sinonimos":[],"tipo":"conceito","titulo":"Triagem e Alocação de Recursos Escassos"}
\section{Triagem e Alocação de Recursos Escassos}
O problema de distribuir bens de saúde limitados (leitos de UTI, órgãos, vacinas) por critérios como maximização de benefício, prioridade aos piores, sorteio ou fila — expondo o conflito entre utilitarismo e igualdade.
\prov{relações: usa justica\_em\_saude\_daniels; usa utilitarismo; relaciona justica\_equidade\_rawls}
\prov{proveniência: wiki \textperiodcentered{} inferencia \textperiodcentered{} confiança 1.0 \textperiodcentered{} domínio ciencias\_de\_fora \textperiodcentered{} história 4 versões no jornal}

%CRISTAL {"arestas":[{"alvo":"epistemologia","confianca":1.0,"epistemico":"fato","origem":"wiki","rel":"parte_de"},{"alvo":"kant_fenomeno_numeno","confianca":1.0,"epistemico":"fato","origem":"wiki","rel":"relaciona"}],"confianca":1.0,"contraexemplos":[],"descricao":"Lopes: 'tudo o que apreendemos é desprovido de natureza inerente, existindo só em relação à estrutura que o detecta' — realidade e verdade relativas ao observador.","epistemico":"hipotese","exemplos":[],"id":"troe_postulado_aureo","memoria":"semantica","meta":{"autor":"Gentil Lopes","dominio":"filosofia","fonte":""},"origem":"livro","palavras_chave":[],"sinonimos":[],"tipo":"conceito","titulo":"TROE / Postulado Áureo"}
\section{TROE / Postulado Áureo}
Lopes: 'tudo o que apreendemos é desprovido de natureza inerente, existindo só em relação à estrutura que o detecta' — realidade e verdade relativas ao observador.
\prov{relações: parte\_de epistemologia; relaciona kant\_fenomeno\_numeno}
\prov{proveniência: livro \textperiodcentered{} hipotese \textperiodcentered{} confiança 1.0 \textperiodcentered{} domínio filosofia \textperiodcentered{} história 4 versões no jornal}

%CRISTAL {"arestas":[{"alvo":"linguas_indigenas","confianca":1.0,"epistemico":"fato","origem":"wiki","rel":"especializa"}],"confianca":1.0,"contraexemplos":[],"descricao":"A família indígena que moldou o português do Brasil (topônimos, léxico) — o tupi antigo foi língua geral na colônia. Raiz sob os pés de Roraima.","epistemico":"fato","exemplos":[],"id":"tupi","memoria":"semantica","meta":{"dominio":"ciencias_de_fora","fonte":""},"origem":"wiki","palavras_chave":[],"sinonimos":[],"tipo":"conceito","titulo":"Tupi"}
\section{Tupi}
A família indígena que moldou o português do Brasil (topônimos, léxico) — o tupi antigo foi língua geral na colônia. Raiz sob os pés de Roraima.
\prov{relações: especializa linguas\_indigenas}
\prov{proveniência: wiki \textperiodcentered{} fato \textperiodcentered{} confiança 1.0 \textperiodcentered{} domínio ciencias\_de\_fora \textperiodcentered{} história 2 versões no jornal}

%CRISTAL {"arestas":[{"alvo":"teoria_da_comunicacao","confianca":1.0,"epistemico":"fato","origem":"wiki","rel":"parte_de"},{"alvo":"teoria_hipodermica","confianca":1.0,"epistemico":"fato","origem":"wiki","rel":"contradiz"},{"alvo":"usos_e_gratificacoes","confianca":1.0,"epistemico":"fato","origem":"wiki","rel":"relaciona"}],"confianca":1.0,"contraexemplos":[],"descricao":"Lazarsfeld & Katz: a influência da mídia raramente é direta — as ideias fluem dos meios para os líderes de opinião e destes, por comunicação interpessoal já interpretada, para o público. Achado que refutou os efeitos diretos.","epistemico":"fato","exemplos":[],"id":"two_step_flow","memoria":"semantica","meta":{"autor":"Lazarsfeld & Katz","dominio":"ciencias_de_fora","fonte":""},"origem":"wiki","palavras_chave":[],"sinonimos":[],"tipo":"conceito","titulo":"Fluxo em Duas Etapas (Lazarsfeld & Katz)"}
\section{Fluxo em Duas Etapas (Lazarsfeld \& Katz)}
Lazarsfeld \& Katz: a influência da mídia raramente é direta — as ideias fluem dos meios para os líderes de opinião e destes, por comunicação interpessoal já interpretada, para o público. Achado que refutou os efeitos diretos.
\prov{relações: parte\_de teoria\_da\_comunicacao; contradiz teoria\_hipodermica; relaciona usos\_e\_gratificacoes}
\prov{proveniência: wiki \textperiodcentered{} fato \textperiodcentered{} confiança 1.0 \textperiodcentered{} domínio ciencias\_de\_fora \textperiodcentered{} história 4 versões no jornal}

%CRISTAL {"arestas":[{"alvo":"teoria_dos_numeros","confianca":1.0,"epistemico":"fato","origem":"wiki","rel":"parte_de"}],"confianca":1.0,"contraexemplos":[],"descricao":"xⁿ+yⁿ=zⁿ não tem solução inteira para n>2. Conjecturado em 1637, provado por Wiles (1994) via curvas elípticas.","epistemico":"fato","exemplos":[],"id":"ultimo_teorema_de_fermat","memoria":"semantica","meta":{"autor":"Wiles","dominio":"ciencias_de_fora","fonte":""},"origem":"wiki","palavras_chave":[],"sinonimos":[],"tipo":"conceito","titulo":"Último Teorema de Fermat"}
\section{Último Teorema de Fermat}
x\ensuremath{^{n}}+y\ensuremath{^{n}}=z\ensuremath{^{n}} não tem solução inteira para n>2. Conjecturado em 1637, provado por Wiles (1994) via curvas elípticas.
\prov{relações: parte\_de teoria\_dos\_numeros}
\prov{proveniência: wiki \textperiodcentered{} fato \textperiodcentered{} confiança 1.0 \textperiodcentered{} domínio ciencias\_de\_fora \textperiodcentered{} história 4 versões no jornal}

%CRISTAL {"arestas":[{"alvo":"semiotica","confianca":1.0,"epistemico":"fato","origem":"wiki","rel":"parte_de"}],"confianca":1.0,"contraexemplos":[],"descricao":"Semioticista e romancista: o signo, o código e os limites da interpretação — nem tudo vale, o texto resiste.","epistemico":"fato","exemplos":[],"id":"umberto_eco","memoria":"semantica","meta":{"autor":"Eco","dominio":"ciencias_de_fora","fonte":""},"origem":"wiki","palavras_chave":[],"sinonimos":[],"tipo":"conceito","titulo":"Umberto Eco"}
\section{Umberto Eco}
Semioticista e romancista: o signo, o código e os limites da interpretação — nem tudo vale, o texto resiste.
\prov{relações: parte\_de semiotica}
\prov{proveniência: wiki \textperiodcentered{} fato \textperiodcentered{} confiança 1.0 \textperiodcentered{} domínio ciencias\_de\_fora \textperiodcentered{} história 2 versões no jornal}

%CRISTAL {"arestas":[{"alvo":"isla","confianca":1.0,"epistemico":"fato","origem":"wiki","rel":"parte_de"},{"alvo":"contrato_social","confianca":1.0,"epistemico":"fato","origem":"wiki","rel":"relaciona"}],"confianca":1.0,"contraexemplos":[],"descricao":"Islã: a comunidade transnacional de todos os muçulmanos, unida pela fé acima de laços de sangue ou tribo. O Alcorão enquadra a diversidade de povos como convite ao reconhecimento mútuo.","epistemico":"inferencia","exemplos":[],"id":"umma","memoria":"semantica","meta":{"autor":"Alcorão","dominio":"sagrado","fonte":""},"origem":"wiki","palavras_chave":[],"sinonimos":[],"tipo":"conceito","titulo":"Umma — a Comunidade dos Crentes"}
\section{Umma — a Comunidade dos Crentes}
Islã: a comunidade transnacional de todos os muçulmanos, unida pela fé acima de laços de sangue ou tribo. O Alcorão enquadra a diversidade de povos como convite ao reconhecimento mútuo.
\prov{relações: parte\_de isla; relaciona contrato\_social}
\prov{proveniência: wiki \textperiodcentered{} inferencia \textperiodcentered{} confiança 1.0 \textperiodcentered{} domínio sagrado \textperiodcentered{} história 3 versões no jornal}

%CRISTAL {"arestas":[{"alvo":"gramatica_universal","confianca":1.0,"epistemico":"fato","origem":"wiki","rel":"usa"},{"alvo":"linguistica","confianca":1.0,"epistemico":"fato","origem":"wiki","rel":"parte_de"}],"confianca":1.0,"contraexemplos":[],"descricao":"Traços comuns a (quase) todas as línguas (Greenberg; a gramática universal de Chomsky) — o que a mente humana impõe a qualquer idioma.","epistemico":"inferencia","exemplos":[],"id":"universais_linguisticos","memoria":"semantica","meta":{"autor":"Greenberg","dominio":"ciencias_de_fora","fonte":""},"origem":"wiki","palavras_chave":[],"sinonimos":[],"tipo":"conceito","titulo":"Universais Linguísticos"}
\section{Universais Linguísticos}
Traços comuns a (quase) todas as línguas (Greenberg; a gramática universal de Chomsky) — o que a mente humana impõe a qualquer idioma.
\prov{relações: usa gramatica\_universal; parte\_de linguistica}
\prov{proveniência: wiki \textperiodcentered{} inferencia \textperiodcentered{} confiança 1.0 \textperiodcentered{} domínio ciencias\_de\_fora \textperiodcentered{} história 3 versões no jornal}

%CRISTAL {"arestas":[{"alvo":"cosmologia_quantica","confianca":1.0,"epistemico":"fato","origem":"wiki","rel":"usa"},{"alvo":"deus_como_vazio_gentil","confianca":1.0,"epistemico":"fato","origem":"wiki","rel":"usa"},{"alvo":"teologia_quantica_gentil","confianca":1.0,"epistemico":"fato","origem":"wiki","rel":"parte_de"}],"confianca":1.0,"contraexemplos":[],"descricao":"Lopes: a física quântica atribui a origem do universo ao vácuo quântico — 'o vazio está cheio de energia' (efeito Casimir, matéria a partir do nada) —, logo Deus≡Origem-de-Tudo≡Vazio. 'Nada se cria, tudo se transforma' é superado: a matéria brota do nada. 'Se quisermos procurar Deus na física, o vácuo é o melhor lugar'.","epistemico":"inferencia","exemplos":[],"id":"universo_vacuo_quantico_gentil","memoria":"semantica","meta":{"autor":"Gentil Lopes","dominio":"sagrado","fonte":""},"origem":"wiki","palavras_chave":[],"sinonimos":[],"tipo":"conceito","titulo":"Deus = Vácuo Quântico (Lopes)"}
\section{Deus = Vácuo Quântico (Lopes)}
Lopes: a física quântica atribui a origem do universo ao vácuo quântico — 'o vazio está cheio de energia' (efeito Casimir, matéria a partir do nada) —, logo Deus\ensuremath{\equiv}Origem-de-Tudo\ensuremath{\equiv}Vazio. 'Nada se cria, tudo se transforma' é superado: a matéria brota do nada. 'Se quisermos procurar Deus na física, o vácuo é o melhor lugar'.
\prov{relações: usa cosmologia\_quantica; usa deus\_como\_vazio\_gentil; parte\_de teologia\_quantica\_gentil}
\prov{proveniência: wiki \textperiodcentered{} inferencia \textperiodcentered{} confiança 1.0 \textperiodcentered{} domínio sagrado \textperiodcentered{} história 5 versões no jornal}

%CRISTAL {"arestas":[{"alvo":"geografia_humana","confianca":1.0,"epistemico":"fato","origem":"wiki","rel":"parte_de"}],"confianca":1.0,"contraexemplos":[],"descricao":"O processo de concentração crescente da população e das atividades em cidades, transformando o território rural em tecido urbano e reorganizando a divisão do trabalho no espaço.","epistemico":"fato","exemplos":[],"id":"urbanizacao","memoria":"semantica","meta":{"autor":"—","dominio":"ciencias_de_fora","fonte":""},"origem":"wiki","palavras_chave":[],"sinonimos":[],"tipo":"conceito","titulo":"Urbanização"}
\section{Urbanização}
O processo de concentração crescente da população e das atividades em cidades, transformando o território rural em tecido urbano e reorganizando a divisão do trabalho no espaço.
\prov{relações: parte\_de geografia\_humana}
\prov{proveniência: wiki \textperiodcentered{} fato \textperiodcentered{} confiança 1.0 \textperiodcentered{} domínio ciencias\_de\_fora \textperiodcentered{} história 2 versões no jornal}

%CRISTAL {"arestas":[{"alvo":"teoria_da_comunicacao","confianca":1.0,"epistemico":"fato","origem":"wiki","rel":"parte_de"},{"alvo":"teoria_hipodermica","confianca":1.0,"epistemico":"fato","origem":"wiki","rel":"contradiz"},{"alvo":"teoria_do_cultivo","confianca":1.0,"epistemico":"fato","origem":"wiki","rel":"contradiz"}],"confianca":1.0,"contraexemplos":[],"descricao":"Katz, Blumler & Gurevitch: inverte a pergunta — não o que a mídia faz com as pessoas, mas o que as pessoas ativamente fazem com a mídia, escolhendo-a para satisfazer necessidades. A audiência é ativa, não passiva.","epistemico":"inferencia","exemplos":[],"id":"usos_e_gratificacoes","memoria":"semantica","meta":{"autor":"Katz, Blumler & Gurevitch","dominio":"ciencias_de_fora","fonte":""},"origem":"wiki","palavras_chave":[],"sinonimos":[],"tipo":"conceito","titulo":"Usos e Gratificações"}
\section{Usos e Gratificações}
Katz, Blumler \& Gurevitch: inverte a pergunta — não o que a mídia faz com as pessoas, mas o que as pessoas ativamente fazem com a mídia, escolhendo-a para satisfazer necessidades. A audiência é ativa, não passiva.
\prov{relações: parte\_de teoria\_da\_comunicacao; contradiz teoria\_hipodermica; contradiz teoria\_do\_cultivo}
\prov{proveniência: wiki \textperiodcentered{} inferencia \textperiodcentered{} confiança 1.0 \textperiodcentered{} domínio ciencias\_de\_fora \textperiodcentered{} história 4 versões no jornal}

%CRISTAL {"arestas":[{"alvo":"teoria_do_valor","confianca":1.0,"epistemico":"fato","origem":"wiki","rel":"parte_de"},{"alvo":"valor_trabalho","confianca":1.0,"epistemico":"fato","origem":"wiki","rel":"contradiz"},{"alvo":"word","confianca":0.5,"epistemico":"fato","origem":"wiki","rel":"relaciona"}],"confianca":1.0,"contraexemplos":[],"descricao":"Menger/Jevons/Walras: o valor vem da satisfação da última unidade (não do trabalho) — a 1ª água ao sedento vale muito, a 10ª quase nada.","epistemico":"inferencia","exemplos":[],"id":"utilidade_marginal","memoria":"semantica","meta":{"autor":"Menger/Jevons/Walras","dominio":"ciencias_de_fora","fonte":""},"origem":"wiki","palavras_chave":[],"sinonimos":[],"tipo":"conceito","titulo":"Utilidade Marginal"}
\section{Utilidade Marginal}
Menger/Jevons/Walras: o valor vem da satisfação da última unidade (não do trabalho) — a 1ª água ao sedento vale muito, a 10ª quase nada.
\prov{relações: parte\_de teoria\_do\_valor; contradiz valor\_trabalho; relaciona word}
\prov{proveniência: wiki \textperiodcentered{} inferencia \textperiodcentered{} confiança 1.0 \textperiodcentered{} domínio ciencias\_de\_fora \textperiodcentered{} história 4 versões no jornal}

%CRISTAL {"arestas":[{"alvo":"etica","confianca":1.0,"epistemico":"fato","origem":"wiki","rel":"parte_de"},{"alvo":"vida-morte","confianca":0.5,"epistemico":"fato","origem":"wiki","rel":"relaciona"},{"alvo":"virtualizacao","confianca":0.5,"epistemico":"fato","origem":"wiki","rel":"relaciona"},{"alvo":"veredito_experimental","confianca":0.5,"epistemico":"fato","origem":"wiki","rel":"relaciona"},{"alvo":"vetores","confianca":0.5,"epistemico":"fato","origem":"wiki","rel":"relaciona"},{"alvo":"word","confianca":0.5,"epistemico":"fato","origem":"wiki","rel":"relaciona"}],"confianca":1.0,"contraexemplos":[],"descricao":"Bentham/Mill: a ação certa maximiza a felicidade agregada — consequencialismo.","epistemico":"inferencia","exemplos":[],"id":"utilitarismo","memoria":"semantica","meta":{"autor":"Bentham/Mill","dominio":"ciencias_de_fora","fonte":""},"origem":"wiki","palavras_chave":[],"sinonimos":[],"tipo":"conceito","titulo":"Utilitarismo"}
\section{Utilitarismo}
Bentham/Mill: a ação certa maximiza a felicidade agregada — consequencialismo.
\prov{relações: parte\_de etica; relaciona vida-morte; relaciona virtualizacao; relaciona veredito\_experimental; relaciona vetores; relaciona word}
\prov{proveniência: wiki \textperiodcentered{} inferencia \textperiodcentered{} confiança 1.0 \textperiodcentered{} domínio ciencias\_de\_fora \textperiodcentered{} história 7 versões no jornal}

%CRISTAL {"arestas":[{"alvo":"blockchain","confianca":1.0,"epistemico":"fato","origem":"curriculo","rel":"continua"},{"alvo":"dinheiro","confianca":1.0,"epistemico":"fato","origem":"curriculo","rel":"parte_de"},{"alvo":"cristalchain","confianca":1.0,"epistemico":"fato","origem":"wiki","rel":"relaciona"}],"confianca":1.0,"contraexemplos":[],"descricao":"Grandeza contábil acumulada; mora no Pell cifrado + certificados encadeados, não no gás ETH.","epistemico":"fato","exemplos":[],"id":"valor","memoria":"semantica","meta":{"dominio":"mercado","nivel":5,"ordem":2},"origem":"curriculo","palavras_chave":["contabilidade","lastro","liquidação"],"sinonimos":["valor econômico"],"tipo":"conceito","titulo":"valor"}
\section{valor}
Grandeza contábil acumulada; mora no Pell cifrado + certificados encadeados, não no gás ETH.
\prov{sinónimos: valor econômico}
\prov{relações: continua blockchain; parte\_de dinheiro; relaciona cristalchain}
\prov{proveniência: curriculo \textperiodcentered{} fato \textperiodcentered{} confiança 1.0 \textperiodcentered{} domínio mercado \textperiodcentered{} história 37 versões no jornal}

%CRISTAL {"arestas":[{"alvo":"signo_linguistico","confianca":1.0,"epistemico":"fato","origem":"wiki","rel":"usa"},{"alvo":"simbolo_vs_significado","confianca":1.0,"epistemico":"fato","origem":"wiki","rel":"relaciona"}],"confianca":1.0,"contraexemplos":[],"descricao":"Saussure: um signo não vale por si, mas pela DIFERENÇA dos outros — o sentido é relacional, como a aboutness da Tiffany.","epistemico":"fato","exemplos":[],"id":"valor_diferencial","memoria":"semantica","meta":{"autor":"Saussure","dominio":"ciencias_de_fora","fonte":""},"origem":"wiki","palavras_chave":[],"sinonimos":[],"tipo":"conceito","titulo":"Valor Diferencial"}
\section{Valor Diferencial}
Saussure: um signo não vale por si, mas pela DIFERENÇA dos outros — o sentido é relacional, como a aboutness da Tiffany.
\prov{relações: usa signo\_linguistico; relaciona simbolo\_vs\_significado}
\prov{proveniência: wiki \textperiodcentered{} fato \textperiodcentered{} confiança 1.0 \textperiodcentered{} domínio ciencias\_de\_fora \textperiodcentered{} história 3 versões no jornal}

%CRISTAL {"arestas":[{"alvo":"teoria_do_valor","confianca":1.0,"epistemico":"fato","origem":"wiki","rel":"parte_de"},{"alvo":"word","confianca":0.5,"epistemico":"fato","origem":"wiki","rel":"relaciona"},{"alvo":"vontade_de_poder","confianca":0.5,"epistemico":"fato","origem":"wiki","rel":"relaciona"},{"alvo":"minerar_e_resfriar","confianca":1.0,"epistemico":"fato","origem":"wiki","rel":"relaciona"}],"confianca":1.0,"contraexemplos":[],"descricao":"Ricardo/Marx: o valor de uma mercadoria é o trabalho socialmente necessário para produzi-la. Refutada pelo paradoxo água-diamante.","epistemico":"inferencia","exemplos":[],"id":"valor_trabalho","memoria":"semantica","meta":{"autor":"Ricardo","dominio":"ciencias_de_fora","fonte":""},"origem":"wiki","palavras_chave":[],"sinonimos":[],"tipo":"conceito","titulo":"Teoria do Valor-Trabalho"}
\section{Teoria do Valor-Trabalho}
Ricardo/Marx: o valor de uma mercadoria é o trabalho socialmente necessário para produzi-la. Refutada pelo paradoxo água-diamante.
\prov{relações: parte\_de teoria\_do\_valor; relaciona word; relaciona vontade\_de\_poder; relaciona minerar\_e\_resfriar}
\prov{proveniência: wiki \textperiodcentered{} inferencia \textperiodcentered{} confiança 1.0 \textperiodcentered{} domínio ciencias\_de\_fora \textperiodcentered{} história 5 versões no jornal}

%CRISTAL {"arestas":[{"alvo":"didatica_matematica","confianca":1.0,"epistemico":"fato","origem":"wiki","rel":"parte_de"},{"alvo":"tall_tres_mundos","confianca":1.0,"epistemico":"fato","origem":"wiki","rel":"relaciona"}],"confianca":1.0,"contraexemplos":[],"descricao":"van Hiele: a progressão hierárquica do pensamento geométrico em cinco níveis — visualização, análise, dedução informal, dedução formal e rigor — cada um pressuposto do seguinte.","epistemico":"inferencia","exemplos":[],"id":"van_hiele_niveis","memoria":"semantica","meta":{"autor":"van Hiele","dominio":"ciencias_de_fora","fonte":""},"origem":"wiki","palavras_chave":[],"sinonimos":[],"tipo":"conceito","titulo":"Níveis de van Hiele"}
\section{Níveis de van Hiele}
van Hiele: a progressão hierárquica do pensamento geométrico em cinco níveis — visualização, análise, dedução informal, dedução formal e rigor — cada um pressuposto do seguinte.
\prov{relações: parte\_de didatica\_matematica; relaciona tall\_tres\_mundos}
\prov{proveniência: wiki \textperiodcentered{} inferencia \textperiodcentered{} confiança 1.0 \textperiodcentered{} domínio ciencias\_de\_fora \textperiodcentered{} história 3 versões no jornal}

%CRISTAL {"arestas":[{"alvo":"readymade_duchamp","confianca":1.0,"epistemico":"fato","origem":"wiki","rel":"usa"},{"alvo":"aura_benjamin","confianca":1.0,"epistemico":"fato","origem":"wiki","rel":"usa"},{"alvo":"artefatos_tem_politica","confianca":1.0,"epistemico":"fato","origem":"wiki","rel":"relaciona"},{"alvo":"artes","confianca":1.0,"epistemico":"fato","origem":"wiki","rel":"parte_de"}],"confianca":1.0,"contraexemplos":[],"descricao":"Bürger: na sociedade burguesa a arte torna-se autônoma (separada da práxis da vida) — o que possibilita a 'arte pela arte'. A vanguarda histórica reage tentando reintegrar arte e vida e atacando a própria instituição arte.","epistemico":"inferencia","exemplos":[],"id":"vanguarda_autonomia","memoria":"semantica","meta":{"autor":"Bürger","dominio":"ciencias_de_fora","fonte":""},"origem":"wiki","palavras_chave":[],"sinonimos":[],"tipo":"conceito","titulo":"Vanguarda e Autonomia da Arte (Bürger)"}
\section{Vanguarda e Autonomia da Arte (Bürger)}
Bürger: na sociedade burguesa a arte torna-se autônoma (separada da práxis da vida) — o que possibilita a 'arte pela arte'. A vanguarda histórica reage tentando reintegrar arte e vida e atacando a própria instituição arte.
\prov{relações: usa readymade\_duchamp; usa aura\_benjamin; relaciona artefatos\_tem\_politica; parte\_de artes}
\prov{proveniência: wiki \textperiodcentered{} inferencia \textperiodcentered{} confiança 1.0 \textperiodcentered{} domínio ciencias\_de\_fora \textperiodcentered{} história 5 versões no jornal}

%CRISTAL {"arestas":[{"alvo":"comunicacao_iconoclasta_gentil","confianca":1.0,"epistemico":"fato","origem":"wiki","rel":"usa"},{"alvo":"honestidade_intelectual","confianca":1.0,"epistemico":"fato","origem":"wiki","rel":"relaciona"},{"alvo":"deus_mentiroso","confianca":1.0,"epistemico":"fato","origem":"wiki","rel":"relaciona"},{"alvo":"critica_ia_semantica","confianca":1.0,"epistemico":"fato","origem":"wiki","rel":"relaciona"},{"alvo":"realismo_juridico","confianca":1.0,"epistemico":"fato","origem":"wiki","rel":"relaciona"}],"confianca":1.0,"contraexemplos":[],"descricao":"Tese comunicacional transversal do Lopes: contra a 'firme convicção', contra o argumento de autoridade e contra a maioria ('se Lopes estiver certo, todos estão errados'), só a refutação lógica decide — 'mostrem onde está a falha lógica'. A persuasão social é insuficiente.","epistemico":"inferencia","exemplos":[],"id":"verdade_por_refutacao_gentil","memoria":"semantica","meta":{"autor":"Gentil Lopes","dominio":"ciencias_de_fora","fonte":""},"origem":"wiki","palavras_chave":[],"sinonimos":[],"tipo":"conceito","titulo":"Verdade por Refutação, não por Consenso (Lopes)"}
\section{Verdade por Refutação, não por Consenso (Lopes)}
Tese comunicacional transversal do Lopes: contra a 'firme convicção', contra o argumento de autoridade e contra a maioria ('se Lopes estiver certo, todos estão errados'), só a refutação lógica decide — 'mostrem onde está a falha lógica'. A persuasão social é insuficiente.
\prov{relações: usa comunicacao\_iconoclasta\_gentil; relaciona honestidade\_intelectual; relaciona deus\_mentiroso; relaciona critica\_ia\_semantica; relaciona realismo\_juridico}
\prov{proveniência: wiki \textperiodcentered{} inferencia \textperiodcentered{} confiança 1.0 \textperiodcentered{} domínio ciencias\_de\_fora \textperiodcentered{} história 7 versões no jornal}

%CRISTAL {"arestas":[{"alvo":"teoria_literaria","confianca":1.0,"epistemico":"fato","origem":"wiki","rel":"parte_de"},{"alvo":"ritmo","confianca":1.0,"epistemico":"fato","origem":"wiki","rel":"relaciona"}],"confianca":1.0,"contraexemplos":[],"descricao":"A medida do poema — sílabas, acentos e pausas que dão ritmo à fala; a forma sonora que o sentido veste.","epistemico":"fato","exemplos":[],"id":"verso_e_metrica","memoria":"semantica","meta":{"dominio":"ciencias_de_fora","fonte":""},"origem":"wiki","palavras_chave":[],"sinonimos":[],"tipo":"conceito","titulo":"Verso e Métrica"}
\section{Verso e Métrica}
A medida do poema — sílabas, acentos e pausas que dão ritmo à fala; a forma sonora que o sentido veste.
\prov{relações: parte\_de teoria\_literaria; relaciona ritmo}
\prov{proveniência: wiki \textperiodcentered{} fato \textperiodcentered{} confiança 1.0 \textperiodcentered{} domínio ciencias\_de\_fora \textperiodcentered{} história 3 versões no jornal}

%CRISTAL {"arestas":[{"alvo":"espaco_santos","confianca":1.0,"epistemico":"fato","origem":"wiki","rel":"usa"},{"alvo":"sociedade_em_rede","confianca":1.0,"epistemico":"fato","origem":"wiki","rel":"relaciona"}],"confianca":1.0,"contraexemplos":[],"descricao":"Milton Santos: verticalidades são pontos distantes conectados por atores hegemônicos em rede (tempo rápido, lógica global); horizontalidades são as contiguidades do território vivido. O 'espaço banal' é o território de todos — o locus da coexistência, em tensão com as redes de alguns.","epistemico":"inferencia","exemplos":[],"id":"verticalidades_horizontalidades","memoria":"semantica","meta":{"autor":"Milton Santos","dominio":"ciencias_de_fora","fonte":""},"origem":"wiki","palavras_chave":[],"sinonimos":[],"tipo":"conceito","titulo":"Verticalidades, Horizontalidades e Espaço Banal (Milton Santos)"}
\section{Verticalidades, Horizontalidades e Espaço Banal (Milton Santos)}
Milton Santos: verticalidades são pontos distantes conectados por atores hegemônicos em rede (tempo rápido, lógica global); horizontalidades são as contiguidades do território vivido. O 'espaço banal' é o território de todos — o locus da coexistência, em tensão com as redes de alguns.
\prov{relações: usa espaco\_santos; relaciona sociedade\_em\_rede}
\prov{proveniência: wiki \textperiodcentered{} inferencia \textperiodcentered{} confiança 1.0 \textperiodcentered{} domínio ciencias\_de\_fora \textperiodcentered{} história 3 versões no jornal}

%CRISTAL {"arestas":[{"alvo":"ciencias","confianca":-0.07,"epistemico":"fato","origem":"manual","rel":"referencia"},{"alvo":"quatro_ciencias","confianca":1.0,"epistemico":"fato","origem":"paper","rel":"parte_de"},{"alvo":"morte","confianca":1.0,"epistemico":"fato","origem":"paper","rel":"paralelo"},{"alvo":"computabilidade","confianca":1.0,"epistemico":"fato","origem":"paper","rel":"referencia"},{"alvo":"biologia","confianca":1.0,"epistemico":"fato","origem":"wiki","rel":"parte_de"},{"alvo":"gato","confianca":1.0,"epistemico":"fato","origem":"wiki","rel":"referencia"}],"confianca":1.0,"contraexemplos":[],"descricao":"Sistema aberto que mantém homeostase por metabolismo e replicação; processo, não substância.","epistemico":"fato","exemplos":[],"id":"vida","memoria":"semantica","meta":{"dominio":"ciencias_de_fora","fonte":""},"origem":"wiki","palavras_chave":[],"sinonimos":[],"tipo":"conceito","titulo":"Vida"}
\section{Vida}
Sistema aberto que mantém homeostase por metabolismo e replicação; processo, não substância.
\prov{relações: referencia ciencias; parte\_de quatro\_ciencias; paralelo morte; referencia computabilidade; parte\_de biologia; referencia gato}
\prov{proveniência: wiki \textperiodcentered{} fato \textperiodcentered{} confiança 1.0 \textperiodcentered{} domínio ciencias\_de\_fora \textperiodcentered{} história 201 versões no jornal}

%CRISTAL {"arestas":[{"alvo":"vida","confianca":1.0,"epistemico":"fato","origem":"wiki","rel":"especializa"},{"alvo":"broca_os","confianca":1.0,"epistemico":"fato","origem":"wiki","rel":"referencia"},{"alvo":"maquina_evolutiva","confianca":1.0,"epistemico":"fato","origem":"wiki","rel":"relaciona"}],"confianca":1.0,"contraexemplos":[],"descricao":"Langton: sistemas computacionais (autômatos, algoritmos genéticos) que exibem replicação, adaptação e evolução — a vida como padrão, não substância.","epistemico":"inferencia","exemplos":[],"id":"vida_artificial","memoria":"semantica","meta":{"autor":"Langton","dominio":"ciencias_de_fora","fonte":""},"origem":"wiki","palavras_chave":[],"sinonimos":[],"tipo":"conceito","titulo":"Vida Artificial"}
\section{Vida Artificial}
Langton: sistemas computacionais (autômatos, algoritmos genéticos) que exibem replicação, adaptação e evolução — a vida como padrão, não substância.
\prov{relações: especializa vida; referencia broca\_os; relaciona maquina\_evolutiva}
\prov{proveniência: wiki \textperiodcentered{} inferencia \textperiodcentered{} confiança 1.0 \textperiodcentered{} domínio ciencias\_de\_fora \textperiodcentered{} história 4 versões no jornal}

%CRISTAL {"arestas":[{"alvo":"comunicacao_digital","confianca":1.0,"epistemico":"fato","origem":"wiki","rel":"parte_de"},{"alvo":"selecao_natural","confianca":1.0,"epistemico":"fato","origem":"wiki","rel":"usa"},{"alvo":"economia_da_atencao","confianca":1.0,"epistemico":"fato","origem":"wiki","rel":"relaciona"},{"alvo":"pos_verdade","confianca":1.0,"epistemico":"fato","origem":"wiki","rel":"relaciona"}],"confianca":1.0,"contraexemplos":[],"descricao":"Dawkins: o meme como replicador cultural análogo ao gene — unidade de imitação que se propaga, muta e sofre seleção; na cultura digital, a unidade viral de circulação de ideias e afetos que competem pela atenção.","epistemico":"inferencia","exemplos":[],"id":"viralidade_memes","memoria":"semantica","meta":{"autor":"Dawkins","dominio":"ciencias_de_fora","fonte":""},"origem":"wiki","palavras_chave":[],"sinonimos":[],"tipo":"conceito","titulo":"Viralidade e Memes"}
\section{Viralidade e Memes}
Dawkins: o meme como replicador cultural análogo ao gene — unidade de imitação que se propaga, muta e sofre seleção; na cultura digital, a unidade viral de circulação de ideias e afetos que competem pela atenção.
\prov{relações: parte\_de comunicacao\_digital; usa selecao\_natural; relaciona economia\_da\_atencao; relaciona pos\_verdade}
\prov{proveniência: wiki \textperiodcentered{} inferencia \textperiodcentered{} confiança 1.0 \textperiodcentered{} domínio ciencias\_de\_fora \textperiodcentered{} história 5 versões no jornal}

%CRISTAL {"arestas":[{"alvo":"politica","confianca":1.0,"epistemico":"fato","origem":"wiki","rel":"parte_de"},{"alvo":"imperativo_categorico","confianca":1.0,"epistemico":"fato","origem":"wiki","rel":"contradiz"},{"alvo":"word","confianca":0.5,"epistemico":"fato","origem":"wiki","rel":"relaciona"},{"alvo":"virtualizacao","confianca":0.5,"epistemico":"fato","origem":"wiki","rel":"relaciona"}],"confianca":1.0,"contraexemplos":[],"descricao":"Maquiavel: virtù (astúcia, audácia) que domina a fortuna (acaso) — o realismo político que separa poder de moral.","epistemico":"inferencia","exemplos":[],"id":"virtu_fortuna","memoria":"semantica","meta":{"autor":"Maquiavel","dominio":"ciencias_de_fora","fonte":""},"origem":"wiki","palavras_chave":[],"sinonimos":[],"tipo":"conceito","titulo":"Virtù e Fortuna"}
\section{Virtù e Fortuna}
Maquiavel: virtù (astúcia, audácia) que domina a fortuna (acaso) — o realismo político que separa poder de moral.
\prov{relações: parte\_de politica; contradiz imperativo\_categorico; relaciona word; relaciona virtualizacao}
\prov{proveniência: wiki \textperiodcentered{} inferencia \textperiodcentered{} confiança 1.0 \textperiodcentered{} domínio ciencias\_de\_fora \textperiodcentered{} história 5 versões no jornal}

%CRISTAL {"arestas":[{"alvo":"teoria_da_comunicacao","confianca":1.0,"epistemico":"fato","origem":"wiki","rel":"parte_de"},{"alvo":"modelo_shannon_weaver","confianca":1.0,"epistemico":"fato","origem":"wiki","rel":"contradiz"},{"alvo":"meio_e_mensagem","confianca":1.0,"epistemico":"fato","origem":"wiki","rel":"relaciona"}],"confianca":1.0,"contraexemplos":[],"descricao":"Carey: contrapõe a visão transmissiva — comunicar como enviar informação no espaço para controle — à visão ritual — comunicar como partilha e comunhão de crenças que mantêm e reparam uma cultura no tempo.","epistemico":"inferencia","exemplos":[],"id":"visao_transmissiva_ritual_carey","memoria":"semantica","meta":{"autor":"Carey","dominio":"ciencias_de_fora","fonte":""},"origem":"wiki","palavras_chave":[],"sinonimos":[],"tipo":"conceito","titulo":"Comunicação: Transmissão vs Ritual (Carey)"}
\section{Comunicação: Transmissão vs Ritual (Carey)}
Carey: contrapõe a visão transmissiva — comunicar como enviar informação no espaço para controle — à visão ritual — comunicar como partilha e comunhão de crenças que mantêm e reparam uma cultura no tempo.
\prov{relações: parte\_de teoria\_da\_comunicacao; contradiz modelo\_shannon\_weaver; relaciona meio\_e\_mensagem}
\prov{proveniência: wiki \textperiodcentered{} inferencia \textperiodcentered{} confiança 1.0 \textperiodcentered{} domínio ciencias\_de\_fora \textperiodcentered{} história 4 versões no jornal}

%CRISTAL {"arestas":[{"alvo":"krishna_avatar","confianca":1.0,"epistemico":"fato","origem":"wiki","rel":"usa"},{"alvo":"atman_brahman","confianca":1.0,"epistemico":"fato","origem":"wiki","rel":"relaciona"},{"alvo":"contemplacao","confianca":1.0,"epistemico":"fato","origem":"wiki","rel":"relaciona"}],"confianca":1.0,"contraexemplos":[],"descricao":"Gita (cap. 11): Krishna concede a Arjuna o 'olho divino' e revela sua Forma Universal — infinitos rostos, o esplendor de mil sóis, contendo o universo inteiro, e o Tempo que devora tudo ('sou o poderoso Tempo, o destruidor do mundo'). Teofania magnífica e aterradora do Absoluto manifesto.","epistemico":"inferencia","exemplos":[],"id":"vishvarupa","memoria":"semantica","meta":{"autor":"Bhagavad Gita","dominio":"sagrado","fonte":""},"origem":"wiki","palavras_chave":[],"sinonimos":[],"tipo":"conceito","titulo":"Vishvarupa — a Visão da Forma Universal"}
\section{Vishvarupa — a Visão da Forma Universal}
Gita (cap. 11): Krishna concede a Arjuna o 'olho divino' e revela sua Forma Universal — infinitos rostos, o esplendor de mil sóis, contendo o universo inteiro, e o Tempo que devora tudo ('sou o poderoso Tempo, o destruidor do mundo'). Teofania magnífica e aterradora do Absoluto manifesto.
\prov{relações: usa krishna\_avatar; relaciona atman\_brahman; relaciona contemplacao}
\prov{proveniência: wiki \textperiodcentered{} inferencia \textperiodcentered{} confiança 1.0 \textperiodcentered{} domínio sagrado \textperiodcentered{} história 4 versões no jornal}

%CRISTAL {"arestas":[{"alvo":"contemplacao","confianca":1.0,"epistemico":"fato","origem":"wiki","rel":"explica"},{"alvo":"word","confianca":0.5,"epistemico":"fato","origem":"wiki","rel":"relaciona"},{"alvo":"vontade_de_poder","confianca":0.5,"epistemico":"fato","origem":"wiki","rel":"relaciona"}],"confianca":1.0,"contraexemplos":[],"descricao":"Aristóteles: a vida contemplativa (theoria) está acima da vida ativa (praxis) — a estrutura contemplação↔ação que o broca-so encarna.","epistemico":"inferencia","exemplos":[],"id":"vita_contemplativa","memoria":"semantica","meta":{"autor":"Aristóteles","dominio":"filosofia","fonte":""},"origem":"wiki","palavras_chave":[],"sinonimos":[],"tipo":"conceito","titulo":"Vida Contemplativa"}
\section{Vida Contemplativa}
Aristóteles: a vida contemplativa (theoria) está acima da vida ativa (praxis) — a estrutura contemplação\ensuremath{\leftrightarrow}ação que o broca-so encarna.
\prov{relações: explica contemplacao; relaciona word; relaciona vontade\_de\_poder}
\prov{proveniência: wiki \textperiodcentered{} inferencia \textperiodcentered{} confiança 1.0 \textperiodcentered{} domínio filosofia \textperiodcentered{} história 4 versões no jornal}

%CRISTAL {"arestas":[{"alvo":"artes_visuais","confianca":1.0,"epistemico":"fato","origem":"wiki","rel":"parte_de"},{"alvo":"belo_platonico","confianca":1.0,"epistemico":"fato","origem":"wiki","rel":"relaciona"},{"alvo":"razao_aurea","confianca":1.0,"epistemico":"fato","origem":"wiki","rel":"relaciona"}],"confianca":1.0,"contraexemplos":[],"descricao":"Vitrúvio (De Architectura): os três requisitos da boa arquitetura — firmitas (solidez), utilitas (utilidade) e venustas (beleza). O primeiro tratado teórico de arquitetura, base da disciplina desde o Renascimento.","epistemico":"fato","exemplos":[],"id":"vitruvio_triade","memoria":"semantica","meta":{"autor":"Vitrúvio","dominio":"ciencias_de_fora","fonte":""},"origem":"wiki","palavras_chave":[],"sinonimos":[],"tipo":"conceito","titulo":"Tríade Vitruviana"}
\section{Tríade Vitruviana}
Vitrúvio (De Architectura): os três requisitos da boa arquitetura — firmitas (solidez), utilitas (utilidade) e venustas (beleza). O primeiro tratado teórico de arquitetura, base da disciplina desde o Renascimento.
\prov{relações: parte\_de artes\_visuais; relaciona belo\_platonico; relaciona razao\_aurea}
\prov{proveniência: wiki \textperiodcentered{} fato \textperiodcentered{} confiança 1.0 \textperiodcentered{} domínio ciencias\_de\_fora \textperiodcentered{} história 4 versões no jornal}

%CRISTAL {"arestas":[{"alvo":"colapso_observador","confianca":1.0,"epistemico":"fato","origem":"wiki","rel":"especializa"},{"alvo":"consciencia","confianca":1.0,"epistemico":"fato","origem":"wiki","rel":"relaciona"}],"confianca":1.0,"contraexemplos":[],"descricao":"von Neumann-Wigner (~1939): a consciência do observador seria o que colapsa a onda. Wigner retratou-se após a decoerência (Zeh).","epistemico":"hipotese","exemplos":[],"id":"von_neumann_wigner","memoria":"semantica","meta":{"autor":"von Neumann/Wigner","dominio":"filosofia","fonte":""},"origem":"wiki","palavras_chave":[],"sinonimos":[],"tipo":"conceito","titulo":"Consciência Causa o Colapso"}
\section{Consciência Causa o Colapso}
von Neumann-Wigner (\~{}1939): a consciência do observador seria o que colapsa a onda. Wigner retratou-se após a decoerência (Zeh).
\prov{relações: especializa colapso\_observador; relaciona consciencia}
\prov{proveniência: wiki \textperiodcentered{} hipotese \textperiodcentered{} confiança 1.0 \textperiodcentered{} domínio filosofia \textperiodcentered{} história 3 versões no jornal}

%CRISTAL {"arestas":[{"alvo":"geografia_humana","confianca":1.0,"epistemico":"fato","origem":"wiki","rel":"parte_de"},{"alvo":"lugares_centrais","confianca":1.0,"epistemico":"fato","origem":"wiki","rel":"relaciona"}],"confianca":1.0,"contraexemplos":[],"descricao":"Von Thünen (1826): a distribuição dos usos agrícolas do solo em anéis concêntricos ao redor da cidade-mercado, em função dos custos de transporte e da renda da terra. O modelo pioneiro da geografia econômica.","epistemico":"inferencia","exemplos":[],"id":"von_thunen","memoria":"semantica","meta":{"autor":"Von Thünen","dominio":"ciencias_de_fora","fonte":""},"origem":"wiki","palavras_chave":[],"sinonimos":[],"tipo":"conceito","titulo":"Localização Agrícola / Anéis de Von Thünen"}
\section{Localização Agrícola / Anéis de Von Thünen}
Von Thünen (1826): a distribuição dos usos agrícolas do solo em anéis concêntricos ao redor da cidade-mercado, em função dos custos de transporte e da renda da terra. O modelo pioneiro da geografia econômica.
\prov{relações: parte\_de geografia\_humana; relaciona lugares\_centrais}
\prov{proveniência: wiki \textperiodcentered{} inferencia \textperiodcentered{} confiança 1.0 \textperiodcentered{} domínio ciencias\_de\_fora \textperiodcentered{} história 3 versões no jornal}

%CRISTAL {"arestas":[{"alvo":"politica","confianca":1.0,"epistemico":"fato","origem":"wiki","rel":"referencia"},{"alvo":"word","confianca":0.5,"epistemico":"fato","origem":"wiki","rel":"relaciona"},{"alvo":"while","confianca":0.5,"epistemico":"fato","origem":"wiki","rel":"relaciona"}],"confianca":1.0,"contraexemplos":[],"descricao":"Nietzsche: a força primária da vida — luta de vontades; o forte cria valores, o 'rebanho' os recebe.","epistemico":"hipotese","exemplos":[],"id":"vontade_de_poder","memoria":"semantica","meta":{"autor":"Nietzsche","dominio":"ciencias_de_fora","fonte":""},"origem":"wiki","palavras_chave":[],"sinonimos":[],"tipo":"conceito","titulo":"Vontade de Poder"}
\section{Vontade de Poder}
Nietzsche: a força primária da vida — luta de vontades; o forte cria valores, o 'rebanho' os recebe.
\prov{relações: referencia politica; relaciona word; relaciona while}
\prov{proveniência: wiki \textperiodcentered{} hipotese \textperiodcentered{} confiança 1.0 \textperiodcentered{} domínio ciencias\_de\_fora \textperiodcentered{} história 4 versões no jornal}

%CRISTAL {"arestas":[{"alvo":"dicionario_musculos_face","confianca":1.0,"epistemico":"fato","origem":"wiki","rel":"especializa"},{"alvo":"orbitas_de_fala","confianca":1.0,"epistemico":"fato","origem":"wiki","rel":"usa"},{"alvo":"espectro_irredutivel_floco","confianca":1.0,"epistemico":"fato","origem":"wiki","rel":"eh_um"},{"alvo":"corpo_estelar","confianca":1.0,"epistemico":"fato","origem":"wiki","rel":"usa"}],"confianca":1.0,"contraexemplos":[],"descricao":"A voz é o modelo FONTE-FILTRO: a laringe é a fonte (f0, os harmônicos); o trato vocal, moldado pelos MÚSCULOS, é o filtro — os FORMANTES. Mapear cada músculo a um formante dá o espectro para gerar qualquer som: os músculos são os MODOS (como as camadas da estrela; um som = a fonte × a soma dos modos). masseter→F1 (abertura), zigomático→F2 (anterior), orbicular da boca→F2 grave (arredonda). As vogais são posições no plano F1×F2: /a/=F1 alto, /i/=F2 alto, /u/=F2 baixo. Cada formante vira uma NOTA — a tradução para música. Verif. em voz.py: F1(/a/)>F1(/i/), F2(/i/)>F2(/u/), a síntese fonte-filtro tem picos nos formantes.","epistemico":"fato","exemplos":[],"id":"voz_fonte_filtro_modos","memoria":"semantica","meta":{"dominio":"biologia","fonte":"voz"},"origem":"wiki","palavras_chave":[],"sinonimos":[],"tipo":"conceito","titulo":"A Voz: os Músculos são os Modos (fonte-filtro)"}
\section{A Voz: os Músculos são os Modos (fonte-filtro)}
A voz é o modelo FONTE-FILTRO: a laringe é a fonte (f0, os harmônicos); o trato vocal, moldado pelos MÚSCULOS, é o filtro — os FORMANTES. Mapear cada músculo a um formante dá o espectro para gerar qualquer som: os músculos são os MODOS (como as camadas da estrela; um som = a fonte $\times$ a soma dos modos). masseter\ensuremath{\to}F1 (abertura), zigomático\ensuremath{\to}F2 (anterior), orbicular da boca\ensuremath{\to}F2 grave (arredonda). As vogais são posições no plano F1$\times$F2: /a/=F1 alto, /i/=F2 alto, /u/=F2 baixo. Cada formante vira uma NOTA — a tradução para música. Verif. em voz.py: F1(/a/)>F1(/i/), F2(/i/)>F2(/u/), a síntese fonte-filtro tem picos nos formantes.
\prov{relações: especializa dicionario\_musculos\_face; usa orbitas\_de\_fala; eh\_um espectro\_irredutivel\_floco; usa corpo\_estelar}
\prov{proveniência: wiki \textperiodcentered{} voz \textperiodcentered{} fato \textperiodcentered{} confiança 1.0 \textperiodcentered{} domínio biologia \textperiodcentered{} história 10 versões no jornal}

%CRISTAL {"arestas":[{"alvo":"platonismo_matematico","confianca":1.0,"epistemico":"fato","origem":"wiki","rel":"relaciona"}],"confianca":1.0,"contraexemplos":[],"descricao":"Estruturas inventadas por prazer puro aplicam-se com precisão à física não prevista — enigma da consonância mente↔natureza (Wigner 1960).","epistemico":"fato","exemplos":[],"id":"wigner_eficacia","memoria":"semantica","meta":{"autor":"Wigner","dominio":"filosofia","fonte":""},"origem":"wiki","palavras_chave":[],"sinonimos":[],"tipo":"conceito","titulo":"Irrazoável Eficácia da Matemática"}
\section{Irrazoável Eficácia da Matemática}
Estruturas inventadas por prazer puro aplicam-se com precisão à física não prevista — enigma da consonância mente\ensuremath{\leftrightarrow}natureza (Wigner 1960).
\prov{relações: relaciona platonismo\_matematico}
\prov{proveniência: wiki \textperiodcentered{} fato \textperiodcentered{} confiança 1.0 \textperiodcentered{} domínio filosofia \textperiodcentered{} história 2 versões no jornal}

%CRISTAL {"arestas":[{"alvo":"sociointeracionismo_vygotsky","confianca":1.0,"epistemico":"fato","origem":"wiki","rel":"usa"}],"confianca":1.0,"contraexemplos":[],"descricao":"Vygotsky: a distância entre o que o aprendiz faz sozinho e o que faz com a ajuda de um parceiro mais capaz — a região onde a instrução é mais produtiva.","epistemico":"inferencia","exemplos":[],"id":"zona_desenvolvimento_proximal","memoria":"semantica","meta":{"autor":"Vygotsky","dominio":"ciencias_de_fora","fonte":""},"origem":"wiki","palavras_chave":[],"sinonimos":[],"tipo":"conceito","titulo":"Zona de Desenvolvimento Proximal (ZDP)"}
\section{Zona de Desenvolvimento Proximal (ZDP)}
Vygotsky: a distância entre o que o aprendiz faz sozinho e o que faz com a ajuda de um parceiro mais capaz — a região onde a instrução é mais produtiva.
\prov{relações: usa sociointeracionismo\_vygotsky}
\prov{proveniência: wiki \textperiodcentered{} inferencia \textperiodcentered{} confiança 1.0 \textperiodcentered{} domínio ciencias\_de\_fora \textperiodcentered{} história 2 versões no jornal}

%CRISTAL {"arestas":[{"alvo":"circulacao_atmosferica_oceanica","confianca":1.0,"epistemico":"fato","origem":"wiki","rel":"usa"}],"confianca":1.0,"contraexemplos":[],"descricao":"Faixas latitudinais e regionais de clima semelhante (tropical, árida, temperada, fria, polar) definidas por temperatura e precipitação — base da classificação de Köppen e da distribuição dos biomas.","epistemico":"fato","exemplos":[],"id":"zonas_climaticas","memoria":"semantica","meta":{"autor":"Köppen","dominio":"ciencias_de_fora","fonte":""},"origem":"wiki","palavras_chave":[],"sinonimos":[],"tipo":"conceito","titulo":"Zonas Climáticas"}
\section{Zonas Climáticas}
Faixas latitudinais e regionais de clima semelhante (tropical, árida, temperada, fria, polar) definidas por temperatura e precipitação — base da classificação de Köppen e da distribuição dos biomas.
\prov{relações: usa circulacao\_atmosferica\_oceanica}
\prov{proveniência: wiki \textperiodcentered{} fato \textperiodcentered{} confiança 1.0 \textperiodcentered{} domínio ciencias\_de\_fora \textperiodcentered{} história 2 versões no jornal}

\end{document}
